\documentclass[11pt,a4paper]{article}
\usepackage[margin=25mm]{geometry}
\usepackage{fontspec}
\tracinglostchars=3
\setmainfont{Latin Modern Roman}
\setsansfont{Latin Modern Sans}
\setmonofont{DejaVu Sans Mono}[Scale=0.78]
\usepackage{amsmath,amssymb,amsthm,mathtools}
\usepackage{mathrsfs}
\usepackage{longtable,booktabs,array,calc,enumitem}
\usepackage{xcolor}
\usepackage[unicode,colorlinks=true,linkcolor=blue!45!black,urlcolor=blue!55!black]{hyperref}
\usepackage{xurl}
\usepackage{newunicodechar}
\newunicodechar{ℓ}{\ensuremath{\ell}}
\newunicodechar{π}{\ensuremath{\pi}}
\newunicodechar{χ}{\ensuremath{\chi}}
\newunicodechar{φ}{\ensuremath{\varphi}}
\newunicodechar{ζ}{\ensuremath{\zeta}}
\newunicodechar{′}{\ensuremath{\prime}}
\newunicodechar{α}{\ensuremath{\alpha}}
\newunicodechar{β}{\ensuremath{\beta}}
\newunicodechar{γ}{\ensuremath{\gamma}}
\newunicodechar{δ}{\ensuremath{\delta}}
\newunicodechar{ε}{\ensuremath{\varepsilon}}
\newunicodechar{θ}{\ensuremath{\theta}}
\newunicodechar{λ}{\ensuremath{\lambda}}
\newunicodechar{μ}{\ensuremath{\mu}}
\newunicodechar{ν}{\ensuremath{\nu}}
\newunicodechar{ξ}{\ensuremath{\xi}}
\newunicodechar{ρ}{\ensuremath{\rho}}
\newunicodechar{σ}{\ensuremath{\sigma}}
\newunicodechar{τ}{\ensuremath{\tau}}
\newunicodechar{ψ}{\ensuremath{\psi}}
\newunicodechar{ω}{\ensuremath{\omega}}
\newunicodechar{Γ}{\ensuremath{\Gamma}}
\newunicodechar{Δ}{\ensuremath{\Delta}}
\newunicodechar{Θ}{\ensuremath{\Theta}}
\newunicodechar{Λ}{\ensuremath{\Lambda}}
\newunicodechar{Ξ}{\ensuremath{\Xi}}
\newunicodechar{Π}{\ensuremath{\Pi}}
\newunicodechar{Σ}{\ensuremath{\Sigma}}
\newunicodechar{Φ}{\ensuremath{\Phi}}
\newunicodechar{Ψ}{\ensuremath{\Psi}}
\newunicodechar{Ω}{\ensuremath{\Omega}}
\newunicodechar{₀}{\ensuremath{_0}}
\newunicodechar{₁}{\ensuremath{_1}}
\newunicodechar{₂}{\ensuremath{_2}}
\newunicodechar{∞}{\ensuremath{\infty}}
\newunicodechar{≤}{\ensuremath{\leq}}
\newunicodechar{≥}{\ensuremath{\geq}}
\newunicodechar{∈}{\ensuremath{\in}}
\newunicodechar{⊗}{\ensuremath{\otimes}}
\newunicodechar{≅}{\ensuremath{\cong}}
\newunicodechar{≃}{\ensuremath{\simeq}}
\newunicodechar{∼}{\ensuremath{\sim}}
\newunicodechar{≠}{\ensuremath{\ne}}
\newunicodechar{→}{\ensuremath{\to}}
\newunicodechar{↔}{\ensuremath{\leftrightarrow}}
\newunicodechar{⊂}{\ensuremath{\subset}}
\newunicodechar{∩}{\ensuremath{\cap}}
\newunicodechar{∪}{\ensuremath{\cup}}
\usepackage{microtype}
\setlength{\emergencystretch}{3em}
\setlength{\parskip}{0.35em}
\setlength{\parindent}{0pt}
\setcounter{tocdepth}{2}
\providecommand{\tightlist}{\setlength{\itemsep}{0pt}\setlength{\parskip}{0pt}}
\newtheorem{theorem}{Theorem}[section]
\newtheorem{proposition}[theorem]{Proposition}
\newtheorem{lemma}[theorem]{Lemma}
\newtheorem{corollary}[theorem]{Corollary}
\theoremstyle{definition}\newtheorem{definition}[theorem]{Definition}
\theoremstyle{remark}\newtheorem{remark}[theorem]{Remark}
\title{ALPOGE\_BRIDGES\\The Riemann hypothesis and Dirichlet GRH\\in Grothendieck's programme}
\author{Research synthesis and proof checks in the present run}
\date{21 November 2026\quad\textnormal{(UTC)}}
\begin{document}
\maketitle
\begin{abstract}
We adopt the supplied Alp\"oge RH and Dirichlet-GRH declarations as
theorems and determine their precise analytic and motivic consequences.
The main work concerns questions and unfinished statements of Grothendieck
and Deligne, using original archive sheets and exact literature locators.
Full proofs are supplied for the deductions and counterexamples established
in this run. Existing results are cited with their operative hypotheses.
The arithmetic consequences include shifted critical lines for abelian
Artin--Tate objects, effective finite Frobenius recognition, the corresponding
Borger realizations, explicit classical consequences, and trace positivity.
Separate cycle-theoretic, Hodge-theoretic and categorical arguments settle
several archive statements without using RH. Counterexamples locate the
information lost by passing to semisimple Euler factors or numerical traces.
Unresolved questions retain explicit gates and attempted approaches.
This working manuscript records ongoing work; its inclusion of a gate
is not a claim that the corresponding problem is solved.
\end{abstract}

\section*{Reading the claims and the check record}
The mathematical input is Alp\"oge (2026), \emph{Remark}, together with
the five named declarations in \nolinkurl{ALPOGE_Main.lean} and their
targets in \nolinkurl{ALPOGE_Statements.lean}. The supplied Mathlib revision
is recorded in the source-scope section. The declarations have been
inspected; the imported proof tree is not in the supplied packet, and
no fresh Lean kernel compilation is claimed. This does not weaken or
replace the input adopted by the task. In this manuscript, ``under the
adopted theorem'' describes the provenance of a consequence; it is not
an additional unproved RH hypothesis within the task's adopted theory.

An identifier denotes a finding established in this run, not a claim
of first publication. We do not claim an exhaustive novelty search.
The check table records actual full-proof readings, primary-source
statement checks, computations, and carried background separately.
Prior intake check reports are never relabelled as checks performed
here. The word ``refuted'' always names the precise statement or
proposed inference for which a counterexample is given.

The strongest arithmetic results concern primitive Dirichlet factors
and integral Tate twists. General Hecke characters over other fields,
nonabelian Artin representations, ramified mixed extensions, integral
lattices, period torsors and monodromy require the additional data
specified in their sections. A proved failure of one reconstruction
rule is not a logical-independence theorem about RH.

\section*{Ranking by the question settled}
\begin{enumerate}
\item \textbf{Original archive statements.} The cote 14 implication and
its proposed converse are addressed over an algebraically closed field
for the precisely defined generated Albanese quotient. All smooth
quasiprojective parameters are treated in every characteristic for the
cohomological implication. CC-ALB01--02 now prove integral equality
of the projective and quasiprojective parameter relations in every
characteristic, using divisibility of the ordinary Albanese kernel
and a uniform alteration degree. The finite-field rational-point reading
has a separate counterexample. Cote 15's abelian--Artin motivic sector is
now treated using the cited 2026 all-CM Hodge theorem and its exact Tate
and standard-conjecture consequences. CC-AG04--06 identify its
coefficient-extended category over the algebraic closure with Weil-germ
gradings, and its protorus with the identity component of the fixed-field
group. CC-AG07--08 give the exact finite transitions and the
Galois-orbit/Schur-index criterion for rational object descent. The
supersingular quaternion obstruction already rules out a rational
neutral fibre functor. Arbitrary-motive generation remains separate. Cote 18's real-Frobenius,
weight-shift and strong-period formulations, and cote 67's capping and
involution claims, receive explicit proofs or counterexamples.
The canonical capping map has an exact real-analytic image but fails
the natural holomorphic formulation. Cote 123's dimension, rational
singularity, nilpotent uniqueness and differential-rank questions have
exact classical answers and checked short deductions. Its strict
nonsimply-laced subregular rank prediction fails: the rank is always
one less than the semisimple rank in every type. The group quotient has a complete
central-character formula; type A has a closed minimal-polynomial
formula. The apposition quadruples and their full automorphism groups are now
classified in every simple complex adjoint type, with exact finite
quotients and orbit counts. Outer-invariant classes exist exactly
outside the even-rank A types. Opposite choices exist in every type;
the all-choice criterion is now fully determined in every simple type.
Besides the stated classical positive cases, it holds in $G_2$; it
fails in $F_4,E_6,E_7,E_8$, with complete exceptional counterexamples.
The recovered cote 18 comparison list has precisely scoped classical
Betti, de Rham, crystalline and integral $p$-adic answers.
Its adelic genericity assertion on physical page 49/footer 48 is
proved by the short deduction BR-18-AD01 from exact local results.
Cote 66's named affine torsion-coefficient comparison cases are
affirmative; its universal degree-zero test is exactly henselianity.
\item \textbf{Deligne and the geometric gate.} The actual problem
Weil II 1.8.15 is answered by the precise admissible-Hodge-theory
theorems. Singular-fibre weights and nearby-cycle weights are computed
separately. The independently checked finite exterior-power criterion reduces the weight--monodromy gate to specified
inertia-invariant trace bounds. The universal stable curve has an explicit
geometric perverse decomposition, and its actual exit functor has the
proved right-Kan direct-image description. The linear theorem of cote
134-8 is established over every scalar ring when its derived functor
commutes canonically with coproducts, in particular in every AB4
abelian category with enough projectives. F-SM-32 constructs the
coherent inverse without AB5, a generator, presentability, an exact
chain-valued realization or functorial projective resolutions. The
remaining original gate is that precise coproduct compatibility;
F-SM-33 further reduces the sufficient condition to countable
compatibility and canonical normalization on all free scalar modules.
F-SM-31 proves that the original tensor and interval conditions cannot
force AB5 or local presentability, even with an exact tensor and a
projective generator. The original projective-generator
corollary is included. The stronger assertion about its original
underived right adjoint retains a separate exact test. Cote 136's
strict equalization assertion is refuted over every nonzero ring;
common sources and terminal ordinary localization are proved over
$\mathbf Q$-algebras and $p$-local rings, hence also over every ring
of positive characteristic. AQ-136-SEMILOCAL extends ordinary terminal
localization to every coefficient ring with finitely many nonunit
rational primes. Strict common sources in that mixed semilocal scope,
and the integral non-strict questions, remain.
\item \textbf{New effective arithmetic access.} AQ-FINITE gives a
finite recognition bound for semisimple abelian Artin--Tate objects;
BR-AT-01 and CC-F01/F01b give the exact shifted analytic factors.
The effective Siegel deduction, abelian-field Chebotarev and prime-search
bounds discharge their original conditional hypotheses.
\item \textbf{Criteria and failed routes.} Trace-class positivity,
moment inequalities, Li and approximation criteria are distinguished
from independent geometric constructions. Explicit mixed-extension,
monodromy, cross-ratio, divided-power and spectral counterexamples
identify the missing information. The 2026 rational Deligne--Drinfeld
theorem is an external result, separate from profinite GT surjectivity.
\end{enumerate}
\tableofcontents
\clearpage
\part{Checked findings and exactly scoped known results}

% BEGIN archive_questions.tex
% Agent contribution: /root/archive_questions.
% AQ identifiers denote findings of this run, not claims of priority.
\section{Archive questions: exact statements and mathematical scope}

\subsection{AQ-FINITE: finite Frobenius recognition of abelian Artin--Tate motives}
\paragraph{Statement.}
Adopt the Dirichlet GRH theorem supplied in \nolinkurl{ALPOGE_Main.lean},
\nolinkurl{LFunction_eq_zero_re}. Let $N\geq1$, put
\[
 Q=\operatorname{lcm}(N,4),\qquad B_Q=(\varphi(Q)\log Q)^2,
\]
and let $M,M'$ be finite direct sums of Tate twists of abelian Artin
motives over $\mathbb Q$, whose Artin representations factor through
$G_N=(\mathbb Z/N\mathbb Z)^\times$. Work initially with a coefficient
number field $E\subset\mathbb C$ splitting $G_N$. Write
\[
 M=\bigoplus_{\chi,n}(\chi\otimes E(n))^{m_{\chi,n}},
 \qquad
 M'=\bigoplus_{\chi,n}(\chi\otimes E(n))^{m'_{\chi,n}}.
\]
Then $M\simeq M'$ if and only if
\[
 \det(1-TF_p\mid M)=\det(1-TF_p\mid M')
 \quad\text{for every prime }p\leq B_Q,\ p\nmid Q.                 \tag{AQ.1}
\]
Here $E(n)$ has Frobenius eigenvalue $p^{-n}$. These finitely many
polynomials recover every multiplicity $m_{\chi,n}$. The same
isomorphism criterion holds with rational coefficients.
The characters are labelled by their Euler-factor values
$\chi(p)$; changing the Artin-reciprocity convention inverts
all character labels and leaves the criterion unchanged.
One may use a single auxiliary $\ell>B_Q$ for all the
displayed Frobenius factors, avoiding the prime $\ell$
among the tested primes.

\paragraph{Known input and exact locator.}
Lamzouri, Li, and Soundararajan, \emph{Conditional bounds for the least
quadratic non-residue and related problems}, arXiv:1309.3595v3,
Corollary 1.2, p.~3, proves under Dirichlet GRH that every reduced residue
class $a\bmod Q$, $Q>3$, contains a prime
$p_a\leq(\varphi(Q)\log Q)^2$. This result is cited, not reproved.

\begin{proof}
The forward direction is functoriality. Inflate all characters to $G_Q$.
For each $a\in G_Q$ choose such a prime $p_a$.
Every eigenvalue at $p_a$ is $\chi(a)p_a^{-n}$, of complex absolute
value $p_a^{-n}$. Since $p_a>1$, different integers $n$ give different
absolute values. Equality in (AQ.1) therefore gives equality,
for every $a$ and $n$, of the sums of eigenvalues of absolute
value $p_a^{-n}$. Multiplying this sum by $p_a^n$ gives
\[
 t_n(a)=\sum_\chi m_{\chi,n}\chi(a).
\]
Finite character orthogonality yields the reconstruction
\[
 m_{\chi,n}=\frac1{|G_Q|}\sum_{a\in G_Q}
             t_n(a)\chi(a)^{-1}.                                \tag{AQ.2}
\]
Indeed, translating $\sum_a\theta(a)$ by an element $b$ with
$\theta(b)\ne1$ makes that sum zero for every nontrivial character
$\theta$; it is $|G_Q|$ for the trivial character. Thus (AQ.2)
gives the same multiplicities for $M'$.
For rational coefficients, the semisimple commutative algebra
$\mathbb Q[G_Q]$ is a product of cyclotomic fields. The
multiplicity of each simple rational module is recovered from
the common multiplicity of its characters over a splitting field.
The rational representations, weight by weight, are therefore
isomorphic as well.
\end{proof}

This is a finite stopping criterion for recognition within the
abelian Artin--Tate part of the programme. It needs every residue
class, or equivalent full character information: values on an
arbitrary generating set do not determine a general character.
It does not reconstruct extension classes, integral lattices,
period comparisons, or nonabelian fundamental groups.

\subsection{AQ-14-FORWARD: the cohomological implication at cote 14, scan 38}
\paragraph{Archive question.}
The sheet labels \emph{Th\'eor\`eme 3.1}, marks its proof
\emph{d\'em. pas faite}, and asks whether $u_Z=0$ implies,
modulo torsion, the vanishing of
\[
 a({}^tZ):H^{2(d-i)+1}(X,\mathbb Q_\ell(d-i))
                     \longrightarrow H^1(A,\mathbb Q_\ell).
                                                               \tag{AQ.3}
\]
Here $d=\dim X$ and $A$ is an abelian parameter variety.
Its margin asks \emph{A-t-on une r\'eciproque}.
The lower display $H^{(n-i)+1}$ is inconsistent with the
degree in the preceding paragraph; (AQ.3) uses the latter.

\paragraph{Precise modern formulation.}
For a smooth projective geometrically connected $X$ over an
algebraically closed field $k$, let $N^i(X)\subset CH^i(X)_{\rm alg}$
be generated by
\[
 \Gamma_*z,\quad
 z\in\ker\bigl(CH_0(T)^0\longrightarrow\operatorname{Alb}_T(k)\bigr),
 \quad\Gamma\in CH^i(T\times X),                                \tag{AQ.4}
\]
where $T$ is smooth projective geometrically connected.
Set $P^i(X)=CH^i(X)_{\rm alg}/N^i(X)$, the generated Albanese
quotient. For $Z\in CH^i(A\times X)$
the map on degree-zero cycles descends to
$u_Z:A(k)\to P^i(X)$, since two lifts of the same point of
$\operatorname{Alb}_A=A$ differ by an Albanese-kernel class.
If $u_Z=0$ integrally, then (AQ.3) is zero for every
$\ell\ne\operatorname{char}k$.
The sole initially supplied sheet 14/038 does not define
$A^i_{\rm alb}(X)$, but the official full cote now resolves
the definition. In its physical PDF page 24 (footer and
manuscript page 23), Definition 1.1 declares a cycle
Albanese-equivalent to zero when it has the form
$Z(v)$ for a smooth connected projective parameter $T$,
a degree-zero zero-cycle $v$ with $S_T(v)=0$, and a
correspondence $Z$ for which the intersection is defined.
The word \emph{projectif} is bracketed in the sheet,
between \emph{un sch\'ema} and \emph{lisse connexe}.
We first prove the projective-parameter formulation.
The subsequent AQ-14-QUASIPROJECTIVE and CC-ALB01--02
establish equality with the quasiprojective relation in every
characteristic. BR-14-SEP in
\nolinkurl{research/root_separated_parameters.tex} further proves
$N_{\rm proj}=N_{\rm aff}=N_{\rm qp}=N_{\rm sep}$ for all
smooth connected separated finite-type parameters over the same
algebraically closed base. Thus the bracketed projectivity
convention does not change this quotient. The nonproper
Albanese and free-zero-cycle conventions remain explicit.
Here $S_T$ is the Albanese summation map.
Proposition 1.2(i), on the same page, says these cycles
form a subgroup containing rational equivalences and
defines the quotient of the Chow group by its image.
Proposition 1.2(v), physical PDF page 25 (footer 24),
says Albanese equivalence implies algebraic equivalence.
Finally Corollary 2.3, physical PDF page 34 (footer 33),
defines $P^i(X)$ itself as the algebraically trivial
subgroup of $A^i_{\rm alb}(X)$. Thus this is the
original quotient in the projective formulation,
not a quotient defined by regular maps to abelian
varieties. The Chow moving lemma makes the finitely
many intersections in (AQ.4) defined; since rational
equivalences are already killed, the cycle-level and
Chow formulations agree. Theorem 2.1, physical PDF
page 31 (footer 30), defines $u_Z$ through the
Albanese summation map. Its image lies in $P^i(X)$,
so vanishing in this subgroup and in the original
total quotient are equivalent. The supplied
\nolinkurl{scan038} uses physical PDF numbering:
it is footer/manuscript page 37, after the cover.
Definition 1.1 itself says only that $k$ is a field.
The preceding handwritten sheets inspected in the full
cote concern other motive conventions and preparatory
cycle calculations; no explicit algebraic-closure
convention was found there. Accordingly the
algebraically closed hypothesis in our theorem is
retained as an explicit mathematical scope condition,
rather than attributed to an unseen original convention.

\paragraph{Primary inputs.}
Bloch (1979), \emph{Torsion algebraic cycles and a theorem of Roitman},
Propositions 3.5 (correspondences), 3.8 (specialization), and
Theorem 4.2 (prime-to-characteristic Roitman), provide the inputs.
They are also explicitly recorded in Achter, Casalaina-Martin,
and Vial, \emph{On the image of the second $\ell$-adic Bloch map},
arXiv:2007.07180v2, Propositions A.23, A.24, A.26, and A.29.
These appendix statements apply in every codimension.

\paragraph{Torsion lemma.}
For Bloch's map
\[
 \lambda_X^i:CH^i(X)[\ell^\infty]\longrightarrow
                    H^{2i-1}(X,\mathbb Q_\ell/\mathbb Z_\ell(i)),
\]
one has
\[
 N^i(X)\cap CH^i(X)[\ell^\infty]\subseteq\ker\lambda_X^i.          \tag{AQ.5}
\]
\begin{proof}
Choose $x$ in the intersection, a finite expression
$x=\sum_j\Gamma_{j*}z_j$ of the form (AQ.4), and $r$ with
$\ell^r x=0$. Spread out this entire expression, its
rational-equivalence certificates, its varieties and
correspondences, their Albanese maps, and the equalities
$\operatorname{alb}(z_j)=0$, over an integral finitely generated
ring over the prime ring. Shrink so that the varieties are
smooth projective, the chosen generic Albanese varieties extend
to abelian schemes $B_j$, and the chosen maps $T_j\to B_j$
extend. It is unnecessary to assert that these are universal
relative Albanese maps. Here the base can be chosen regular:
the finitely generated field of definition has a separating
transcendence basis over its prime field, and a model becomes
smooth over $\mathbb F_p$, or over $\mathbb Z[1/M]$ in
characteristic zero, after a further shrinking. All the
finitely many conditions just listed survive that shrinking.
In characteristic zero invert $\ell$ as well.
Choose a closed point with finite residue field of
characteristic $p\ne\ell$. A regular system of parameters
in its local ring gives a flag of prime ideals with regular
successive quotients. Localizing at consecutive primes
gives discrete valuation rings. Specialization along this
flag, after choosing algebraic closures of the residue
fields, gives geometric special fibres over
$\overline{\mathbb F}_p$, denoted by subscripts $0$.

Every degree-zero zero-cycle on a smooth projective connected
variety over $\overline{\mathbb F}_p$ is torsion: join its
finite support by finitely many projective curves and lift
the differences of points to their normalizations. Every
degree-zero divisor class on those curves is a point of a
Jacobian defined over a finite extension of a finite field,
and has finite order.

The specialized map $T_{j,0}\to B_{j,0}$ factors through
$h_j:\operatorname{Alb}(T_{j,0})\to B_{j,0}$. This map is an
isomorphism on $\ell$-power torsion. Indeed Kummer theory gives
$H^1(T,\mathbb Z_\ell(1))=T_\ell\operatorname{Pic}(T)
=T_\ell\operatorname{Pic}^0(T)_{\rm red}$: the N\'eron--Severi
group is finitely generated and has zero Tate module.
Duality of $\operatorname{Pic}^0_{\rm red}$ and Albanese
therefore gives
$H^1(T,\mathbb Z_\ell)=\operatorname{Hom}(T_\ell
\operatorname{Alb}_T,\mathbb Z_\ell)$, compatibly with maps.
This is also the codimension-one calculation in
Achter--Casalaina-Martin--Vial, Propositions A.25 and A.28.
The generic Albanese map thus induces an isomorphism on
$H^1(-,\mathbb Z_\ell)$, and smooth proper base change
preserves this isomorphism. The latter is obtained
by taking inverse limits of the finite-coefficient base-change
isomorphisms for each variety; no assertion that the Albanese
map is an isomorphism on every finite-coefficient $H^1$ is
needed. Equivalently $T_\ell h_j$ is an isomorphism.
It follows that $h_j$ is an isogeny with no $\ell$-primary
kernel, and hence an isomorphism on $\ell$-power torsion.
Since the specialization of the equality
$\operatorname{alb}(z_j)=0$ says that $z_{j,0}$ maps to zero
in $B_{j,0}$, the $\ell$-primary part of its actual Albanese
class is zero. Prime-to-$p$ Roitman implies that the
$\ell$-primary part of $z_{j,0}$ itself is zero. Thus every
$z_{j,0}$, and hence $x_0$, has order prime to $\ell$.
The equality $\ell^r x_0=0$ also specializes, so $x_0=0$.

Bloch's compatibility with specialization sends
$\lambda_X^i(x)$ to $\lambda_{X_0}^i(x_0)=0$. The
cohomological specialization map is an isomorphism by
smooth proper base change, first for finite coefficients
and then for their direct limit $\mathbb Q_\ell/\mathbb Z_\ell(i)$.
Consequently $\lambda_X^i(x)=0$. Only finitely many data are
spread for each $x$; no simultaneous model for all of $N^i(X)$
has been assumed.
\end{proof}

\begin{proof}[Proof of AQ-14-FORWARD]
Put $g=\dim A$. Roitman identifies
$CH_0(A)[\ell^\infty]$ with $A[\ell^\infty](k)$.
Use the torsion zero-cycle lift of a torsion point under this
isomorphism; it is unnecessary to identify it with the
literal class $[a]-[0]$. For every such torsion cycle $z$,
the premise gives $Z_*z\in N^i(X)\cap CH^i(X)[\ell^\infty]$.
Thus (AQ.5) gives $\lambda_X^i(Z_*z)=0$.
Correspondence compatibility, followed by taking Tate modules
and tensoring with $\mathbb Q_\ell$, makes
\[
 Z_*:H^{2g-1}(A,\mathbb Q_\ell(g))
             \longrightarrow H^{2i-1}(X,\mathbb Q_\ell(i))
\]
zero: the degree-zero Bloch map identifies the source
with $V_\ell A$. Under the perfect Poincar\'e pairings its
transpose is (AQ.3), so (AQ.3) is zero.
\end{proof}

This proves the stated generated-Albanese-quotient formulation
in every characteristic, independently of RH.

\paragraph{AQ-14-RATIONAL-PREMISE.}
The integral premise cannot be replaced by $u_Z\otimes\mathbb Q=0$.
Take an elliptic curve $E/\overline{\mathbb F}_p$ and
$X=A=E$, $Z=\Delta_E-E\times\{0\}$.
Then $P^1(E)=\operatorname{Pic}^0(E)$ and
$u_Z(a)=[a]-[0]$. The group $E(\overline{\mathbb F}_p)$ is
torsion, so the rationalized premise holds. The normalized
diagonal acts as the identity on the nonzero two-dimensional
$H^1(E,\mathbb Q_\ell)$, refuting the weakened assertion.
The phrase \emph{modulo torsion} in the original sheet
qualifies its cohomological conclusion.
For the asserted $P^1$ identification, a divisor
correspondence on $T\times E$ defines a morphism
$T\to\operatorname{Pic}^0(E)$ after choosing a base point,
and hence factors through $\operatorname{Alb}_T$.
Every generator of $N^1(E)$ is therefore zero.

\paragraph{AQ-14-FIELD-SCOPE.}
The algebraically closed base is also substantive.
If the premise is tested only on $k$-rational points
over an arbitrary field, it can fail to detect the
cohomological map. Take $k=\mathbb F_5$ and
$E:y^2=x^3+1$. The affine fibres over
$x=0,1,2,3,4$ have respectively $2,0,2,0,1$
points, so $\#E(\mathbb F_5)=6$, including infinity.
For $Z=6(\Delta_E-E\times\{0\})$, the map
$u_Z$ on $E(\mathbb F_5)$ is zero by Lagrange's
theorem, while the induced rational $\ell$-adic
map is $6\operatorname{id}\ne0$ for every $\ell\ne5$.
The original context's use of divisibility of
$A_Y(k)$ in physical PDF page 34 (footer 33)
is consistent with the algebraically closed
convention used here. This counterexample prevents
an unintended arbitrary-field reading.

\subsection{AQ-14-QUASIPROJECTIVE: removing the parameter qualification in characteristic zero}
Let $k$ be algebraically closed and let $X/k$ be smooth projective.
For a smooth connected parameter variety $T$, write
$Z_0(T)^0$ for the group of finite formal sums
$v=\sum_jn_j[t_j]$ of $k$-points with $\sum_jn_j=0$.
In this paragraph the Albanese of a nonproper $T$ means its
\emph{abelian Albanese}: the universal pointed morphism
$a_T:T\to\operatorname{Alb}_{\mathrm{ab}}(T)$ to an abelian
variety. Put $a_T(v)=\sum_jn_ja_T(t_j)$; this does not depend
on the base point used for the pointed morphism.

For $\Gamma\in\mathrm{CH}^i(X\times T)$ define
\[
 \Gamma(v)=\sum_j n_j\,\iota_{t_j}^{!}\Gamma
                         \in\mathrm{CH}^i(X),
\]
where $\iota_t:X\times\{t\}\hookrightarrow X\times T$
is the regular immersion of the fibre. Thus this is a
sum of fibre restrictions; no nonproper pushforward
$X\times T\to X$ is being used. Let $N^i_{\mathrm{qp}}(X)$
be generated by all such $\Gamma(v)$ with $T$ smooth connected
quasiprojective and $a_T(v)=0$. Let $N^i_{\mathrm{proj}}(X)$
use smooth connected projective $T$ instead.

If every permitted $T$ has a smooth projective compactification,
then
\[
 N^i_{\mathrm{qp}}(X)=N^i_{\mathrm{proj}}(X).                 \tag{AQ.14.1}
\]
In particular this equality holds in characteristic zero.
It also holds if the convention for a nonproper parameter is
the universal semiabelian Albanese, with the kernel imposed
on $Z_0(T)^0$. This last convention must be formulated on
zero-cycles, not automatically on the ordinary Chow group
of the nonproper variety.

\begin{proof}
The inclusion from right to left is immediate, since a
smooth projective parameter is quasiprojective. For the
other inclusion take a generator $\Gamma(v)$ as above.
Choose a smooth projective compactification
$j:T\hookrightarrow\overline T$. Represent $\Gamma$ by
a finite integral linear combination of integral
codimension-$i$ subvarieties of $X\times T$, and take
their closures in $X\times\overline T$. This gives a class
$\overline\Gamma\in\mathrm{CH}^i(X\times\overline T)$
whose restriction is $\Gamma$.

The morphism $a_{\overline T}\circ j:T\to
\operatorname{Alb}(\overline T)$, normalized to send the
chosen base point to zero, factors by universality through
$a_T$. Therefore
$a_{\overline T}(j_*v)=0$. Here $j_*v$ denotes the same
finite formal zero-cycle regarded on $\overline T$;
it does not assert the existence of a pushforward
$j_*:\mathrm{CH}_0(T)\to\mathrm{CH}_0(\overline T)$.
For each point in the support of $v$, refined fibre
restriction commutes with restriction to the open subset
$X\times T$; equivalently, both restrictions are computed
in the same neighbourhood of $X\times\{t_j\}$. Hence
\[
 \overline\Gamma(j_*v)=\Gamma(v).
\]
The left side is a generator allowed in
$N^i_{\mathrm{proj}}(X)$, proving the reverse inclusion.
The same proof works for the semiabelian convention,
because the target $\operatorname{Alb}(\overline T)$ is
also semiabelian, and because the semiabelian Albanese
of a smooth projective parameter is its ordinary
abelian Albanese.

In characteristic zero, embed the quasiprojective $T$
in a projective variety and resolve that projective
closure by a projective resolution that is an
isomorphism over its smooth locus, in particular over
$T$. The resulting $\overline T$ is smooth projective.
This is the strong compactification consequence of
Hironaka's resolution theorem, stated in precisely
this form in Kahn and Sujatha, \emph{Birational geometry
and localisation of categories}, May 22, 2008 author
version, Theorem 1.6.3(a), p.~11.
\end{proof}

Thus AQ-14-FORWARD applies in characteristic zero to the
quotient generated using all smooth quasiprojective
parameters, and AQ-14-CONVERSE over $\mathbb C$ remains
a counterexample for that larger parameter class as
well. The algebraically closed base-field qualification
remains. In positive characteristic this comparison
applies to parameters admitting a smooth projective
compactification; no existence of such compactifications
in all dimensions is asserted here. The next theorem proves the
rational cohomological forward implication without needing that
positive-characteristic equality of relations. AQ-14-ALTERATION gives the intermediate $p$-primary torsion bound; the subsequently checked companion theorems CC-ALB01--02 remove that discrepancy integrally in every characteristic.

\subsection{AQ-14-NONPROPER: the rational cohomological implication for all smooth quasiprojective parameters}
Retain the algebraically closed base $k$ and smooth projective $X$.
Let $N^i_{\mathrm{qp}}(X)$ be the relation defined in
AQ-14-QUASIPROJECTIVE, using the \emph{abelian} Albanese of a
smooth connected quasiprojective parameter and free degree-zero
zero-cycles. Put
$P^i_{\mathrm{qp}}(X)=\mathrm{CH}^i(X)_{\mathrm{alg}}/
N^i_{\mathrm{qp}}(X)$. For every prime
$\ell\ne\operatorname{char}k$, one has
\[
 N^i_{\mathrm{qp}}(X)\cap\mathrm{CH}^i(X)[\ell^\infty]
       \subseteq\ker\lambda_X^i.                            \tag{AQ.14.2}
\]
Consequently the implication of original 14/038 holds for this
quasiprojective-parameter quotient in every characteristic:
if $A$ is abelian, $Z\in\mathrm{CH}^i(A\times X)$, and
$u_Z:A(k)\to P^i_{\mathrm{qp}}(X)$ is zero integrally, then
\[
 a({}^tZ):H^{2(\dim X-i)+1}(X,\mathbb Q_\ell(\dim X-i))
                       \longrightarrow H^1(A,\mathbb Q_\ell)
\]
is zero. The same implication holds for the smaller relation
using the semiabelian Albanese kernel. No equality
$N^i_{\mathrm{qp}}=N^i_{\mathrm{proj}}$ in arbitrary positive
characteristic is asserted.

\paragraph{Inputs and coefficient conventions.}
Write $H_j^S(T,R)$ for Suslin homology: it is the homology of
the complex generated in degree $n$ by integral closed
subvarieties of $T\times\Delta^n$ finite and surjective over
$\Delta^n$, tensored with $R$. Its degree-zero cycles are
finite formal sums of closed points. Write $G_T$ for the
pointed Serre semiabelian Albanese and
\[
 0\longrightarrow S_T\longrightarrow G_T\longrightarrow B_T
       \longrightarrow0
\]
for its torus and abelian quotient. Universality shows that
$B_T$ is the abelian Albanese of $T$: a homomorphism from a
torus to an abelian variety is zero. The notation $[\ell^\infty]$
means the subgroup annihilated by some power of $\ell$.

Geisser, \emph{Rojtman's theorem for normal schemes},
arXiv:1402.1831v2, Theorem 1.1, p.~2, gives, for normal $T$,
\[
 H_0^S(T,\mathbb Z)[\ell^\infty]
       \xrightarrow{\sim}G_T(k)[\ell^\infty],\qquad
 H_1^S(T,\mathbb Z)\otimes\mathbb Q_\ell/\mathbb Z_\ell=0.
                                                               \tag{AQ.14.3}
\]
The degree-zero convention on the left is automatic for
torsion. This prime-to-characteristic theorem does not
require a smooth compactification. The finite-coefficient
perfect pairing used below is Geisser--Schmidt (2016),
Theorem 1.1, journal p.~92; their Theorem 7.1, p.~118, identifies it
with the Suslin--Voevodsky comparison. We invoke its
prime-to-characteristic case, where tame and ordinary
abelian $\ell$-primary covers agree. Kahn (2014),
\emph{Algebraic tori as Nisnevich sheaves with transfers},
Lemma 4.1(b), journal p.~707 (July 4, 2013 author version,
Lemma 3(b), p.~8), states
$\operatorname{Alb}_{\mathrm{sa}}(G)=G$ for a semiabelian $G$.

The integral generic-base-change input is Sun,
arXiv:1511.01393v1, Definition 2.2, p.~3, Lemma 2.3(iii),
pp.~4--5, and Theorem 2.5(i), p.~7, with
$\Lambda=\mathbb Z_\ell$. In Lemma 2.3(iii)'s proof, one
dense open works for \emph{all} finite coefficient levels,
by the finite set of Jordan--H\"older constituents. The
constant $\mathbb Z_\ell$ system has the single simple
constituent $\mathbb F_\ell$. Constructibility and preservation
of stratifiability are part of the stated setup and
Remark 2.5.1(ii). We use this only over a finite-type
$\mathbb F_p$-base, $p\ne\ell$, which satisfies the finite
cohomological-dimension condition (LO).

\begin{proof}
Characteristic zero follows already from
AQ-14-QUASIPROJECTIVE and AQ-14-FORWARD. Thus assume
$\operatorname{char}k=p>0$ and $\ell\ne p$.
We first give the three facts needed for specialization.

\emph{1. A correspondence with proper target descends to
Suslin zero-homology.}
For $\Gamma\in\mathrm{CH}^i(X\times T)$, the fibre-evaluation
map $Z_0(T)\to\mathrm{CH}^i(X)$ kills every Suslin boundary.
Indeed, a generator of a Suslin one-chain is an integral
$W\subset\mathbb A^1\times T$ finite and surjective over
$\mathbb A^1$. Its normalization $C$ is a smooth curve,
finite over $W$ and over $\mathbb A^1$. Let $\overline C$
be its smooth projective completion. Pull $\Gamma$ back
along the morphism $X\times C\to X\times T$; pullback is
well defined between smooth varieties by the graph/refined
Gysin construction. Extend a representing cycle by closure
to a class on $X\times\overline C$.

The function $t/(t-1)$ on $\overline C$ has divisor equal
to the fibre over zero minus the fibre over one, including
multiplicities, and no contribution over infinity. The
normalization projection formula identifies their images
in $Z_0(T)$ with the two faces of $W$. Applying the extended
correspondence to this principal divisor gives zero in
$\mathrm{CH}^i(X)$. Its evaluations at points of $C$ agree
with the original fibre evaluations, since refined pullback
is local around those fibres. This proves the claim and
produces a homomorphism
\[
 \Gamma_*:H_0^S(T,\mathbb Z)\longrightarrow\mathrm{CH}^i(X).
                                                               \tag{AQ.14.4}
\]
The construction is compatible with extension of the base
field and with automorphisms of that field. No pushforward
along the nonproper projection $X\times T\to X$ was used.

\emph{2. The integral $H^1$ of the semiabelian Albanese.}
The homology coefficient sequence and (AQ.14.3) give a
natural isomorphism
\[
 H_1^S(T,\mathbb Q_\ell/\mathbb Z_\ell)
             \simeq G_T(k)[\ell^\infty].                    \tag{AQ.14.5}
\]
For clarity, the coefficient sequence has left term
$H_1^S(T,\mathbb Z)\otimes\mathbb Q_\ell/\mathbb Z_\ell$
and right term
$\operatorname{Tor}_1^{\mathbb Z}(H_0^S(T,\mathbb Z),
\mathbb Q_\ell/\mathbb Z_\ell)=H_0^S(T,\mathbb Z)[\ell^\infty]$;
the left term is zero by (AQ.14.3).

The prime-to-$p$ finite pairings of Geisser--Schmidt are
natural and compatible with the coefficient homomorphisms.
Dualizing the direct system
$H_1^S(T,\mathbb Z/\ell^n)$, whose transition is induced by
$\mathbb Z/\ell^n\hookrightarrow\mathbb Z/\ell^{n+1}$,
identifies its dual with the inverse system
$H^1_{\mathrm{et}}(T,\mathbb Z/\ell^n)$ under reduction.
The pairings are viewed in $\mathbb Q_\ell/\mathbb Z_\ell$
via $1\bmod\ell^n\mapsto\ell^{-n}\bmod\mathbb Z_\ell$;
this makes the transition compatibility explicit. Thus
\[
 H^1_{\mathrm{et}}(T,\mathbb Z_\ell)
  \simeq\operatorname{Hom}\bigl(
       G_T(k)[\ell^\infty],\mathbb Q_\ell/\mathbb Z_\ell\bigr)
  \simeq\operatorname{Hom}_{\mathbb Z_\ell}
       (T_\ell G_T,\mathbb Z_\ell).                         \tag{AQ.14.6}
\]
The left side is the inverse limit of the finite groups;
there is no $\varprojlim^1 H^0$ obstruction, because the
constant $H^0$ transition maps are surjective. All
isomorphisms are functorial. Since $G_T$ is its own Serre
Albanese, the Albanese morphism $a_T:T\to G_T$ induces an
isomorphism on $H^1(-,\mathbb Z_\ell)$. This conclusion
is deliberately integral-adic, not an assertion about
$H^1(-,\mathbb Z/\ell^n)$ for each $n$.

\emph{3. The finite-field vanishing.}
Suppose for the moment that $k=\overline{\mathbb F}_p$.
Every class in $H_0^S(T,\mathbb Z)^0$ is torsion.
To see this, join the finitely many points in a representative
by finitely many integral curves in the connected smooth
quasiprojective variety $T$, and pass to their normalizations.
Such curves can be chosen through any prescribed finite
set by hyperplane sections of sufficiently high degree;
only the resulting curves and their normalizations are
used. Differences of points lift to degree-zero classes
on these smooth open curves. For a smooth curve $C$,
$H_0^S(C,\mathbb Z)^0$ is the group of points of its
Rosenlicht generalized Jacobian, equivalently the curve
case of Geisser, Theorem 1.3, p.~3. Each such Jacobian and
each of its points are defined over some finite extension
of a finite field, and its rational points over that
finite field form a finite group. Each lifted class,
and hence each original degree-zero class, is torsion.

Let $v\in H_0^S(T,\mathbb Z)^0$ have zero image in $B_T(k)$,
and let $v_\ell$ denote its $\ell$-primary component.
Under (AQ.14.3), the class $v_\ell$ corresponds to an
element of $S_T(k)[\ell^\infty]$. Consider the homomorphism
\[
 S_T(k)[\ell^\infty]\longrightarrow
 H_0^S(T,\mathbb Z)[\ell^\infty]
 \xrightarrow{\Gamma_*}\mathrm{CH}^i(X)[\ell^\infty]
 \xrightarrow{\lambda_X^i}
 H^{2i-1}(X,\mathbb Q_\ell/\mathbb Z_\ell(i)).                \tag{AQ.14.7}
\]
All varieties, maps and $\Gamma$ descend to a common finite
field $\mathbb F_q$, after enlarging it if needed. Thus
(AQ.14.7) is Frobenius-equivariant. On
$V_\ell S_T=T_\ell S_T\otimes\mathbb Q_\ell$ geometric
Frobenius has eigenvalues $q^{-1}$ times roots of unity:
a torus splits over a finite extension of $\mathbb F_q$.
On
$H^{2i-1}(X,\mathbb Q_\ell(i))$ its eigenvalues have
complex absolute value $q^{-1/2}$, by Deligne (1974),
Theorem (1.6), printed p.~276, after the indicated Tate
twist. These spectra are disjoint. A linear map
commuting with Frobenius between these two spaces is zero:
apply relatively prime characteristic polynomials and
B\'ezout's identity, or use generalized eigenspaces
after extending scalars. Semisimplicity is not
required for this step.

The rational Tate module of the last group in (AQ.14.7)
is exactly $H^{2i-1}(X,\mathbb Q_\ell(i))$. More explicitly,
that group has maximal divisible subgroup
$H^{2i-1}(X,\mathbb Z_\ell(i))_{\mathrm{free}}
 \otimes\mathbb Q_\ell/\mathbb Z_\ell$ and finite quotient
contained in $H^{2i}(X,\mathbb Z_\ell(i))[\ell^\infty]$.
The domain of (AQ.14.7) is divisible. Its image is therefore
in the maximal divisible subgroup, and a homomorphism
between these finite-corank divisible $\ell$-groups is
zero if its map of rational Tate modules is zero.
Consequently (AQ.14.7) is zero, and
$\lambda_X^i(\Gamma_*v_\ell)=0$.

It follows that over $\overline{\mathbb F}_p$ the conclusion
(AQ.14.2) holds: in a finite sum of generators each
$H_0^S$ class is torsion; its prime-to-$\ell$ component
maps to prime-to-$\ell$ torsion, whereas the $\ell$-primary
component has zero Bloch image by the preceding argument.
Taking the $\ell$-primary component of an expression for
$x\in N^i_{\mathrm{qp}}(X)\cap\mathrm{CH}^i(X)[\ell^\infty]$
gives the claimed vanishing.

\emph{4. Spreading to the finite-field situation.}
Now take arbitrary algebraically closed $k$ of characteristic
$p>0$ and a fixed
$x\in N^i_{\mathrm{qp}}(X)\cap\mathrm{CH}^i(X)[\ell^\infty]$.
Choose a finite expression
$x=\sum_j\Gamma_{j*}v_j$, with $v_j\in Z_0(T_j)^0$
mapping to zero in $B_{T_j}(k)$, and choose $r$ with
$\ell^r x=0$. Spread all varieties, finite zero-cycles,
correspondences, rational-equivalence witnesses for these
equalities, and the pointed maps
$T_j\to G_{T_j}\to B_{T_j}$ over an integral finite-type
$\mathbb F_p$-scheme $S$. After shrinking, $S$ is regular,
$X/S$ is smooth projective, the $T_j/S$ are smooth with
geometrically connected fibres, and the $G_j/S$ are
semiabelian schemes in exact sequences
$0\to S_j\to G_j\to B_j\to0$ with tori $S_j$ and
abelian schemes $B_j$. The sections representing the
zero-cycles remain inside $T_j$, and the relations
$B_j(v_j)=0$ persist. The spread maps need not be
relative universal Albanese morphisms.

Apply Sun's integral generic base change simultaneously to
the finitely many $T_j\to S$ and $G_j\to S$. Their
$R^1(-)_*\mathbb Z_\ell$ are constructible adic sheaves,
and become lisse on a further dense open. By (AQ.14.6)
the map induced by $T_j\to G_j$ is an isomorphism at
a geometric generic point. Indeed its base extension
to $k$ is the chosen Albanese map, and prime-to-$p$
\'etale cohomology is unchanged by extension of
algebraically closed fields. Its kernel and cokernel
vanish on a further dense open, so this is an
isomorphism of integral adic sheaves there. Generic
base change now identifies its stalk at \emph{every}
geometric point of that open with the actual map
on $H^1(-,\mathbb Z_\ell)$ of the fibre. This uses one
open for the full integral system, not separate opens
for infinitely many coefficient levels.

Choose a closed point $s$ in this open; its residue
field is finite. Over $\overline{k(s)}$, the map
$T_{j,s}\to G_{j,s}$ factors by universality as
\[
 T_{j,s}\longrightarrow G_{T_{j,s}}
                  \xrightarrow{h_j}G_{j,s}.
\]
By (AQ.14.6), $T_\ell h_j$ is an isomorphism. A
homomorphism of semiabelian varieties whose Tate module
is an isomorphism is an isogeny with no $\ell$-primary
kernel: a positive-dimensional connected kernel or
quotient has a nonzero Tate module, and a nonzero finite
$\ell$-primary kernel makes the Tate-module image a
proper sublattice. The induced homomorphism of abelian
quotients
\[
 b_j:B_{T_{j,s}}\longrightarrow B_{j,s}
\]
is also an isogeny with no $\ell$-primary kernel.
Indeed an isogeny of semiabelian varieties maps the
torus onto the torus; the kernel of its map of abelian
quotients is a quotient of its finite kernel. The first
assertion follows since the resulting quotient of the
target by the image torus is a quotient of the proper
source abelian variety, and hence has zero toric part.

The equality $B_{j,s}(v_{j,s})=0$ and injectivity of
$b_j$ on $\ell$-primary torsion imply that the
$\ell$-primary component of the actual abelian Albanese
class of $v_{j,s}$ is zero. Although the whole class
need not vanish, this is precisely what step 3 needs:
the $\ell$-primary Suslin class of $v_{j,s}$ corresponds
to the actual torus $S_{T_{j,s}}$. Thus the argument
in step 3 gives $\lambda_{X_s}^i(x_s)=0$.

Finally specialize from a geometric generic point to
$s$ along a regular flag through $s$, as in the complete
flag construction in AQ-14-FORWARD. The finite cycle
and rational-equivalence data make $x_s$ the
specialization of $x$, preserve $\ell^r x_s=0$, and
identify each specialized fibre evaluation with
$\Gamma_{j,s*}v_{j,s}$. Refined Gysin functoriality
provides this compatibility locally at the sections.
Bloch's specialization theorem sends
$\lambda_X^i(x)$ to $\lambda_{X_s}^i(x_s)$, and smooth
proper base change for $X/S$ makes the cohomological
specialization an isomorphism. Cohomological invariance
under algebraically closed extension also identifies
the original $k$-fibre with the geometric generic
fibre for this calculation. Therefore
$\lambda_X^i(x)=0$, proving (AQ.14.2).

For the final cohomological implication, lift each
$\ell$-power torsion point of $A$ to its torsion
zero-cycle class by prime-to-$p$ Roitman. The integral
premise puts its image under $Z$ in the left side of
(AQ.14.2). Bloch correspondence compatibility, followed
by Tate modules and tensoring with $\mathbb Q_\ell$,
therefore kills
$Z_*:H^{2\dim A-1}(A,\mathbb Q_\ell(\dim A))
\to H^{2i-1}(X,\mathbb Q_\ell(i))$.
Poincar\'e duality gives exactly the displayed transpose.
This is the same final argument as AQ-14-FORWARD, now
with the broader relation. For the semiabelian-kernel
convention, its relation is contained in the
abelian-kernel relation, so the torsion vanishing and
the implication follow as well.
\end{proof}

\subsection{AQ-14-ALTERATION: the two parameter relations agree away from the characteristic}
\paragraph{Statement.}
Let $k$ be algebraically closed of characteristic $p>0$,
and let $X/k$ be smooth projective.  Define
$N^i_{\rm proj}(X)\subset N^i_{\rm qp}(X)\subset CH^i(X)$
using the projective and smooth quasiprojective parameters,
respectively, and the ordinary abelian Albanese, exactly as
in AQ-14-QUASIPROJECTIVE.  Then
\[
 N^i_{\rm qp}(X)/N^i_{\rm proj}(X)
 \quad\hbox{is a $p$-primary torsion group.}               \tag{AQ.ALT.1}
\]
In particular the relations agree after tensoring with
$\mathbb Z[1/p]$, hence also with $\mathbb Q$.
The natural map between the two cycle quotients has
$p$-primary kernel and induces an isomorphism on
$\ell$-primary torsion for every $\ell\ne p$.
The same assertions hold if the nonproper relation uses
zero cycles annihilated by the semiabelian Albanese.
Together with AQ-14-QUASIPROJECTIVE this gives equality
with rational coefficients in every characteristic.
This is retained as an intermediate theorem.  The companion
theorems CC-ALB01--02 subsequently prove that this kernel
is zero, by divisibility and the fixed-alteration-degree
quantifier made explicit below.

\paragraph{Exact external inputs.}
The separability refinement is needed here: Illusie and
Temkin (2014), Expos\'e X, Theorem 3.5(iii), printed p.~202,
gives a \emph{separable} projective prime-to-$\ell$ alteration
with regular source over a perfect field.  Their
Theorem 2.1 alone is not the separability statement;
the distinction is explicit on p.~167 and in \S3.5.4,
p.~207.  The other input is the extension of a rational
map from a smooth variety to an abelian variety to a
morphism on the whole source: Milne, \emph{Abelian
Varieties}, version 2.0 (March 16, 2008), Chapter I,
Theorem 3.2, printed p.~17.  We use these as cited
classical inputs.  The moving and trace arguments
needed for their application are proved below.

\begin{proof}
\emph{1. Moving zero cycles inside a dense open using Suslin relations.}
For a smooth connected quasiprojective $T$ and a nonempty
open $U\subset T$, every class in $H^S_0(T,\mathbb Z)$
is represented by a zero cycle supported in $U$.
The replacement can be chosen with the same degree.
We prove this directly on point generators.
There is nothing to prove for a point already in $U$,
or if $\dim T=0$.
Fix $x\in T\setminus U$.  An affine neighborhood of $x$
embeds into affine space.  Successive general hyperplanes
through $x$ cut its dimension down to one and cut
$T\setminus U$ down to dimension zero in that neighborhood.
At each step one avoids finitely many positive-dimensional
components; this is possible because $k$ is infinite and
the linear system has only the point $x$ as its base locus.
An irreducible one-dimensional component through $x$
therefore meets $U$.  Its closure in $T$ is an integral
curve $C_0$ through $x$, still meeting $U$.
Let $C\to C_0$ be its normalization and choose
$x'\in C$ above $x$.  The normalization is finite,
and $C$ is a smooth curve because $k$ is perfect.
Let $\overline C$ be its smooth projective completion,
$D=\overline C\setminus C$, and
$F=C\setminus(C\times_T U)$.
Both $D$ and $F$ are finite sets.  Choose
$y\in C\times_T U$.

For sufficiently large $n$, Riemann--Roch gives a rational
function $f$ on $\overline C$, with poles only at $y$,
whose first-order jet at $x'$ is a uniformizer and whose
value at every point of $D\cup(F\setminus\{x'\})$
is $1$.  Explicitly the required evaluation map is
surjective because the first cohomology of
\[
 \mathcal O_{\overline C}
 \left(ny-2x'-\sum_{z\in D\cup(F\setminus\{x'\})}z\right)
\]
vanishes once its degree exceeds $2g(\overline C)-2$.
Thus $f$ has a simple zero at $x'$ and no zero or pole
at the other points of $D\cup F$.
Its divisor is of the form
\[
 \operatorname{div}(f)=[x']+z-a[y],
 \qquad \operatorname{supp}(z)\subset C\times_T U,
\]
where $z$ is effective and $a>0$.
This divisor is a Suslin boundary on $C$:
the nonconstant morphism
$g=f/(f-1):\overline C\to\mathbb P^1$ is finite,
and $g^{-1}(\mathbb A^1)\subset C$ since $f=1$ on $D$.
Its graph gives a finite correspondence from
$\mathbb A^1$ to $C$; its fibres at $0$ and $1$
differ by $\operatorname{div}(f)$, up to the harmless
choice of boundary sign.  Push this correspondence to
$T$.  It shows that $[x]$ has a same-degree representative
supported in $U$.  Linearity proves the moving assertion.

Both maps used in our relation respect these Suslin
boundaries.  For the correspondence action
$\Gamma_*:Z_0(T)\to CH^i(X)$, this is the proper-target
Suslin descent proved in Step 1 of AQ-14-NONPROPER:
pull to a normalized curve, close on its projective
completion, and evaluate the principal divisor there.
For the ordinary Albanese map it follows directly as well:
the morphism $C\to\operatorname{Alb}_{\rm ab}(T)$
extends to $\overline C$ by properness, and any principal
divisor has zero sum in that abelian variety.
Consequently, if $v$ has degree zero and Albanese sum zero,
we may replace it by a zero cycle $v_U$ supported in $U$
with
\[
 \deg v_U=0,\qquad\operatorname{alb}_T(v_U)=0,
 \qquad \Gamma_*(v_U)=\Gamma_*(v).                \tag{AQ.ALT.2}
\]

\emph{2. The prime-to-$\ell$ transfer.}
Fix a prime $\ell\ne p$ and a generator
$\Gamma_*(v)$ of $N^i_{\rm qp}(X)$.
Choose a projective integral compactification
$T\subset\overline T$.
The cited Theorem 3.5(iii), applied over the perfect field
$k$, gives a regular projective source and a separable
projective prime-to-$\ell$ alteration of $\overline T$.
Choose a connected component that dominates
$\overline T$; write it as
\[
 f:S\longrightarrow\overline T,\qquad
 d=[k(S):k(T)],\qquad (d,\ell)=1.
\]
It is smooth over $k$, projective and integral.
There exists a nonempty open $U\subset T$ over which
$f^{-1}(U)\to U$ is finite \'etale of degree $d$.
Use Step 1 to obtain $v_U$ satisfying (AQ.ALT.2).

Fix a point of $S$ to normalize
$a_S:S\to A_S=\operatorname{Alb}(S)$.
The fibre sum
\[
 \nu:U\longrightarrow A_S,\qquad
 u\longmapsto\sum_{s\in f^{-1}(u)}a_S(s)
\]
is a morphism: on an \'etale cover of $U$ the finite
\'etale cover splits into $d$ sections, their sum is
regular, and these sums descend because they are
invariant under permutations of the sections.
It is a rational map from smooth $T$ to an abelian
variety, so the cited extension theorem gives a morphism
$\overline\nu:T\to A_S$.
After subtracting its value at a base point, universality
of $\operatorname{Alb}_{\rm ab}(T)$ makes this morphism
factor through a homomorphism
$\operatorname{Alb}_{\rm ab}(T)\to A_S$.
Write $v_U=\sum_u n_u[u]$.  Since it has degree and Albanese sum zero,
\[
 w=f_U^*v_U\in Z_0(S)^0,
 \qquad \operatorname{alb}_S(w)
        =\sum_u n_u\nu(u)=0.                       \tag{AQ.ALT.3}
\]

Restrict a cycle representing $\Gamma$ to $X\times U$,
pull it back along the finite \'etale map
$X\times f^{-1}(U)\to X\times U$, and take its
cycle-theoretic closure on $X\times S$.
Call the resulting codimension-$i$ cycle $\Gamma'$.
Fibre Gysin evaluation commutes with this \'etale
pullback and with restriction to the open set.
Each of the $d$ points above $u\in U(k)$ therefore
has the same fibre class $\Gamma_u\in CH^i(X)$.
It follows that
\[
 \Gamma'_*(w)=d\,\Gamma_*(v_U)
             =d\,\Gamma_*(v)
             \in N^i_{\rm proj}(X),               \tag{AQ.ALT.4}
\]
where membership uses the projective parameter $S$
and (AQ.ALT.3).

\emph{3. The quotient and its torsion.}
Every element $c$ of $N^i_{\rm qp}/N^i_{\rm proj}$
is a finite sum of such generators.  For each
$\ell\ne p$, take the product of their degrees in
(AQ.ALT.4); it is prime to $\ell$ and annihilates $c$.
Taking any one $\ell\ne p$ first shows that $c$ is
torsion, of some finite order $n$.  For every
$\ell\ne p$ a prime-to-$\ell$ annihilator forces
$\ell\nmid n$.  Thus $n$ is a power of $p$, proving
(AQ.ALT.1).

Put $Q_{\rm proj}=CH^i(X)/N^i_{\rm proj}$ and
$Q_{\rm qp}=CH^i(X)/N^i_{\rm qp}$, and let
$K=N^i_{\rm qp}/N^i_{\rm proj}$.
The exact sequence
$0\to K\to Q_{\rm proj}\to Q_{\rm qp}\to0$
has $p$-primary torsion kernel.  Multiplication by
$\ell$ is an automorphism of $K$ when $\ell\ne p$:
this holds on every finite cyclic subgroup and hence
on their union.  Therefore the induced map on
$\ell$-primary torsion is injective.  It is surjective
as well: if $y\in Q_{\rm qp}$ is killed by
$\ell^a$, lift it to $x\in Q_{\rm proj}$;
then $\ell^a x\in K$, and choose $z\in K$ with
$\ell^a z=\ell^a x$.  The lift $x-z$ is killed by
$\ell^a$.  This proves the assertion about torsion.
Localization is flat and kills $K$, proving the
claimed equalities after inverting $p$.

Finally, a zero cycle in the semiabelian Albanese
kernel is in the abelian Albanese kernel.  Its relation
is therefore between $N^i_{\rm proj}$ and the relation
just treated.  The same $p$-primary conclusion follows
for that intermediate subgroup and quotient.
\end{proof}

\paragraph{Subsequent integral strengthening: CC-ALB01--02.}
The complete companion proof in
\nolinkurl{research/consequences_criteria_integral_albanese.tex}
proves
\[
                  N^i_{\rm proj}(X)=N^i_{\rm qp}(X)
\]
integrally over every algebraically closed field, for both
of the explicit ordinary and semiabelian parameter conventions.
Its new input is that
$\ker(H^S_0(T,\mathbb Z)^0\to\operatorname{Alb}_{\rm ab}(T)(k))$
is divisible.  CC-ALB01 proves the stronger surjectivity on
$n$-torsion for every $n$, including powers of the characteristic,
using connected isogeny pullbacks, Katz's curve sections and an
explicit connected-component calculation on a Jacobian pullback.
No torsion injectivity is assumed.

The earlier alteration proof has a uniform quantifier:
for fixed $T$ and correspondence, one alteration and one dense
finite-etale open work for every class of that kernel.  Its
single degree $d$ therefore annihilates the entire image in
$CH^i(X)/N^i_{\rm proj}(X)$.  Dividing any kernel class by
that same $d$ makes its image zero.  This removes the earlier
possible $p$-primary discrepancy rather than merely bounding
its exponent.  The full new proof and primary inputs have
received independent complete readings, recorded under
CC-ALB01 and CC-ALB02.

\paragraph{Consequence for the original cohomological question.}
This also gives a second route to the forward implication
of 14/038.  If an $\ell$-primary torsion cycle
$x\in CH^i(X)$ belongs to $N^i_{\rm qp}$, (AQ.ALT.4),
applied to a finite expression for $x$, supplies
$d$ prime to $\ell$ with $dx\in N^i_{\rm proj}$.
AQ-14-FORWARD gives $d\lambda_X^i(x)=0$.
Since $\lambda_X^i(x)$ is $\ell$-primary and
multiplication by $d$ is invertible on that group,
$\lambda_X^i(x)=0$.
The original integral-premise and rational-transpose
argument then applies.  This deduction uses the
projective torsion lemma and the alteration theorem;
the earlier direct Suslin-weight proof is retained
as a separately checked argument.

\subsection{AQ-14-CONVERSE: the generated-quotient converse fails}
For the generated quotient just defined, over $\mathbb C$
the converse fails, even rationally in its target.
The following is an application of Walker (2007), Theorem 5.7(4),
p.~930, and Theorems 6.1 and 6.2, pp.~933--934.
Take a smooth split six-dimensional quadric $W$,
a sufficiently general high-degree smooth hypersurface section
$Y\subset W$, and the two-cycle $\gamma$ supplied by Theorem 6.1
with $r=2$: its singular homology class is zero, whereas
$s\gamma\ne0$ in $L_1H_4(Y,\mathbb Q)$.
The two families of maximal $\mathbb P^3$'s on $W$
give independent middle-dimensional algebraic
homology classes, supplying the rank-at-least-two
hypothesis of Theorem 6.1.
Here $\dim Y=5$. For any complex elliptic curve $E$ put
$X=Y\times E$, so $\dim X=6$, and
\[
 Z=\gamma\times(\Delta_E-E\times\{0\})\in CH^4(E\times X).
\]
Then (AQ.3) is zero but $u_Z$ is nonzero in $P^4(X)$,
even after tensoring with $\mathbb Q$.
\begin{proof}
The cycle class of $Z$ is zero because that of $\gamma$ is zero,
so its cohomological correspondence vanishes.
The needed factorization through the generated Albanese
quotient follows explicitly from Walker's Theorem 5.7(2),(4),
p.~930, and Example 5.5, p.~928. If $d=\dim X$,
$r=d-i$, and $\Gamma\in CH^i(T\times X)$, those statements give
\[
 \Phi^{\rm mor}_{X,r}(\Gamma_*z)
    =\Gamma_{J,*}\Phi^{\rm mor}_{T,0}(z)
    =\Gamma_{J,*}\operatorname{alb}_T(z).
\]
Theorem 5.7(4) applies to every Chow correspondence, not just
an equidimensional representative. The displayed expression
vanishes for each generator (AQ.4). Thus
$\Phi^{\rm mor}_{X,r}$ descends to $P^i(X)$.
The product construction in the proof of Walker's Theorem 6.2
gives a nonzero cylinder
\[
 \gamma\otimes H_1(E,\mathbb Q)\longrightarrow L_2H_5(X,\mathbb Q).
\]
Pairing its image with a dual class on $E$ and pushing to $Y$
gives a nonzero multiple of $s\gamma$, verifying nonvanishing.
This is a morphism of the (ind-)Hodge structures defining
the morphic Jacobian. The source has type $(-1,0),(0,-1)$
after twisting. A nonzero rational morphism from this source
has nonzero tangent map on the associated Jacobians: if the
relevant complex Hodge summand were killed, its conjugate
would also be killed, so the whole morphism would vanish.
More explicitly, strictness of the Hodge filtration gives
a nonzero complex linear map after quotient by $F^0$;
if it were zero, the whole real image would lie in
$F^0\cap\overline{F^0}$, which is zero in weight $-1$.
The underlying Lawson group is countable, as used in
Walker's proof of Theorem 6.2, p.~934. Quotienting the
nonzero vector-space image by its intersection with
this countable period group still gives uncountably
many image points. Torsion in a complex vector space
modulo a countable subgroup is countable, because
its inverse image is the countable union of the
groups $\tfrac1n\Lambda$. Hence
$a\mapsto\gamma\times([a]-[0])$ has nontorsion image
in the morphic Jacobian.
Since it factors through $P^4(X)$, the converse fails
even rationally.
\end{proof}
No unverified effective degree, such as degree 14, is
needed: the primary theorem's sufficiently-large-degree
construction supplies the counterexample. This is a
translation of a known phenomenon to the original
projective quotient of 14/038, independent of RH.
Walker's morphic Jacobian
need not be a finite-dimensional abelian variety.
Consequently this argument does not refute the separate
universality question for the classical Abel--Jacobi
map among finite-dimensional regular abelian targets;
Walker discusses that distinction in \S7.

\subsection{AQ-WEIL-EXIST: the strong conjecture at 15/100}
The source says that its preceding necessary conditions
are sufficient, following the boxed condition
$|\sigma\alpha|=q^{i/2}$. In the unambiguous modern form,
every nonzero algebraic integer with these absolute values
occurs on $H^i$ of a smooth projective scheme over
$\mathbb F_q$. This is true for every $i\geq0$.

The unconditional existence argument is the surjectivity
part of Milne (1994), Proposition 2.6, p.~15 of the author's
manuscript. The whole proposition classifies simple motives
under the Tate conjecture, explicitly assumed at the start
of \S2, p.~12; its full classification must retain that
assumption. Honda (1968), \emph{Isogeny classes of abelian
varieties over finite fields}, Main Theorem, p.~83, is
the unconditional existence input.
\begin{proof}[Short proof of the specified $H^i$ form]
For $i>0$ choose $\beta$ with $\beta^i=\alpha$.
It is integral over the algebraic integers, hence an
algebraic integer, and all its conjugates have absolute
value $q^{1/2}$. Honda's theorem gives an abelian variety
$A/\mathbb F_q$ with eigenvalue $\beta$ on $H^1$.
The K\"unneth summand $H^1(A)^{\otimes i}$ of $H^i(A^i)$
then contains $\beta^i=\alpha$.
For $i=0$, Kronecker's theorem makes $\alpha$ a root of
unity of some order $m$. The permutation action on
$H^0(\operatorname{Spec}\mathbb F_{q^m}
\times_{\mathbb F_q}\overline{\mathbb F}_q,\mathbb Q_\ell)$
contains all $m$th roots of unity. This scheme is smooth
projective. If geometric connectedness were imposed,
the weight-zero assertion would instead allow only $1$.
\end{proof}
The weak assertion of 15/106 follows in this Weil-number
reading, even without a further extension of the finite
field. The claimed identification of the entire motivic
Galois group is separate: realization's full faithfulness
and algebraicity of invariant correspondences are not
proved merely by the existence of each eigenvalue.
Milne's Proposition 2.6 and Remark 2.7 retain Tate precisely
at that point.
More explicitly, his Conjecture 1.14 is the equality
between the dimension of numerical codimension-$r$
cycle classes and the order of the pole at $s=r$.
Proposition 1.15 derives algebraicity of all Frobenius
invariants and semisimplicity; Corollary 1.16 derives
full faithfulness of the $\ell$-adic realization.
These are the additional inputs used for the
injectivity part of Proposition 2.6 and for the
generation assertion of Remark 2.7. The eigenvalue
existence proof above supplies neither of them.
The integral condition also forces $i\geq0$:
for $i<0$ the nonzero integral norm of $\alpha$
would have absolute value less than one.

\paragraph{AQ-TATE-2026: the abelian sector is cleared in current literature.}
The hypothesis required for abelian varieties has now
been supplied in the primary literature. Van Hoften
and Sempliner (2026), Hexagon:2610.00259v2,
Theorem I, p.~1, prove the rational Hodge conjecture
for every complex CM abelian variety. Applying
Milne (1999), \emph{Lefschetz motives and the Tate
conjecture}, Theorem 7.1, journal p.~72, with the definition (0.1)
and finite-model descent in the Introduction, gives
the full Tate conjecture for every abelian variety
over every finite field. Milne's November 9, 2026
author note states this explicitly as Theorem 3.1,
p.~6. His 2002 paper's final author version,
Theorem 3.3(a),(b), p.~8, also gives numerical
equals $\ell$-adic homological equivalence and the
Hodge standard conjecture for abelian varieties
over algebraically closed fields; Remark 3.5,
p.~9, treats good Weil cohomology theories.
These are cited results, not independently reproved
in this module. The older long-version numbering
\emph{Theorem 4.6} used in the 2026 paper must not
be confused with Theorem 3.3 in the specified final
author PDF.

For comparison, Broe (2026), arXiv:2608.28651v1, Theorem 1.1,
printed p.~1, proves for every abelian fourfold
$A/\mathbb F_q$ and $\ell\ne\operatorname{char}\mathbb F_q$ that
\[
 CH^2(A)\otimes\mathbb Q_\ell\longrightarrow
 H^4_{\rm et}(A_{\overline{\mathbb F}_q},\mathbb Q_\ell(2))
 ^{\operatorname{Gal}(\overline{\mathbb F}_q/\mathbb F_q)}
\]
is surjective. With the divisor theorem and hard
Lefschetz this gives Tate in every codimension on $A$.
Corollary 1.2, p.~2, and Corollary 8.2, p.~20,
give agreement of numerical and $\ell$-adic homological
equivalence for abelian fourfolds over arbitrary fields
where $\ell$ is invertible. Corollary 8.1(i), p.~20,
also gives rational equals numerical equivalence with
rational coefficients over a finite field.
These are cited preprint results whose exact statements
were inspected here; their complete proofs are not
claimed independently checked. The fourfold result
alone does not address all powers of a fourfold;
the October--November all-abelian result does.
It therefore clears the realization issue on the
abelian--Artin category below. It does not prove
Tate for arbitrary smooth projective varieties,
or that every numerical motive belongs to that category.
Milne's Remark 7.6(b), journal p.~76, explicitly
records the extension to homological motives of
abelian type.

\paragraph{AQ-ABELIAN-RECOGNITION.}
Fix $k=\mathbb F_q$, $\ell\ne\operatorname{char}k$,
and let $\mathcal A_q$ be the rational numerical
category obtained from the motives of abelian
varieties over $k$, finite \'etale $k$-schemes, and
Tate motives by finite sums, tensor products,
duals, and direct summands. Equivalently, its
objects are direct summands of finite sums
\[
 h^i(A)(n)\otimes h^0(S),\qquad
 A/k\text{ abelian},\quad S/k\text{ finite \'etale},
 \quad n\in\mathbb Z.
\]
We use geometric Frobenius $F$ and
$\mathbb Q_\ell(1)$ with eigenvalue $q^{-1}$.
The summand above has weight $w=i-2n$ and
parity $i\bmod2=w\bmod2$.
Let $R_\ell(M)$ be its ordinary direct-sum
realization, with the weight decomposition remembered.
Then
\[
 \operatorname{Hom}_{\mathcal A_q}(M,N)\otimes_{\mathbb Q}\mathbb Q_\ell
 \xrightarrow{\ \sim\ }
 \operatorname{Hom}_{\mathbb Q_\ell[F]}
       (R_\ell(M),R_\ell(N)).                         \tag{AQ.15}
\]
Moreover, for $M,N\in\mathcal A_q$,
\[
 M\simeq N\quad\Longleftrightarrow\quad
 P_M(t)=P_N(t),\qquad
 P_M(t)=\det(1-tF\mid R_\ell(M)).                    \tag{AQ.16}
\]
The same equivalence holds with the weight-signed
Euler factor
\[
 Z_M(t)=\prod_w
 \det(1-tF\mid R_\ell(M)_w)^{(-1)^{w+1}}            \tag{AQ.17}
\]
in place of $P_M$. These statements are in the
rational numerical category, and concern actual
objects with their full multiplicities.

\begin{proof}
First we explain the scalar extension in (AQ.15).
For an abelian variety $X/k$, numerical and
$\ell$-adic homological equivalence agree by the
cited all-abelian theorem. Let $A^r_{\rm num}(X)_{\mathbb Q}$
denote the space of rational cycles modulo numerical
equivalence. The numerical intersection pairing
with $A^{\dim X-r}_{\rm num}(X)_{\mathbb Q}$ is nondegenerate.
For any finite linearly independent collection
$\alpha_1,\ldots,\alpha_s$ of numerical classes,
one can choose complementary classes
$\beta_1,\ldots,\beta_s$ making
$(\deg(\alpha_i\beta_j))_{i,j}$ invertible over
$\mathbb Q$. Their cohomology classes are then
linearly independent over $\mathbb Q_\ell$, by
pairing against the $\beta_j$. This also bounds
$s$ by the finite Betti number, proving finiteness
and injectivity after extension to $\mathbb Q_\ell$.
The full Tate theorem supplies surjectivity onto
the Frobenius-invariant subspace. Hence
\[
 A^r_{\rm num}(X)_{\mathbb Q}\otimes\mathbb Q_\ell
 \simeq H^{2r}(X_{\bar k},\mathbb Q_\ell(r))^{F=1}.
                                                               \tag{AQ.18}
\]

The cohomological projectors of an abelian variety
are algebraic here without a further conjecture.
For any integer $a\geq2$, multiplication by $a$
acts on $H^j(A)$ as $a^j$. The polynomials
\[
 e_i=\prod_{\substack{0\leq j\leq2\dim A\\j\ne i}}
       \frac{[a]^*-a^j}{a^i-a^j}
\]
give its projectors in numerical correspondences:
their required identities hold in cohomology,
hence numerically by the established equivalence
on $A\times A$. Tensor products are realized
using $A\times B$, and Poincar\'e duality expresses
duals as complementary cohomological degrees
with Tate twists. This proves the stated description
of $\mathcal A_q$.

Consider first $M=h^i(A)(n)$ and $N=h^j(B)(m)$,
and set $a=\dim A$. A correspondence between them
has codimension $a+m-n$ on $A\times B$. Its
relevant K\"unneth component, by Poincar\'e duality,
is
\[
 H^{2a-i}(A)\otimes H^j(B)(a+m-n)
 \simeq\operatorname{Hom}\bigl(H^i(A)(n),H^j(B)(m)\bigr).
\]
It lies in the required cohomological degree
precisely when $j-i=2(m-n)$, equivalently when
the weights agree. In that case (AQ.18) for
$A\times B$, followed by the algebraic projectors,
gives exactly (AQ.15). For unequal weights every
Frobenius-equivariant map is zero: the eigenvalues
have complex absolute values $q^{(i-2n)/2}$ and
$q^{(j-2m)/2}$, so their spectra are disjoint.
The numerical correspondence group is also zero,
since its cohomological component has the wrong
degree and homological equals numerical equivalence.

For finite \'etale factors $S,T$, the schemes
$A\times B\times S\times T$ are disjoint unions
of abelian varieties viewed over finite extensions
of $k$. A Frobenius-invariant class on a component
orbit is determined by a class invariant under the
appropriate power of Frobenius on one component.
The all-abelian Tate theorem over that finite
extension supplies the cycle. Intersection pairings
are multiplied by the nonzero extension degree,
which is invertible over $\mathbb Q$. Thus the
same injection and surjection apply to these
factors. Passing to finite sums and applying
source and target idempotents proves (AQ.15)
for every object of $\mathcal A_q$.

Frobenius acts semisimply on all these realizations.
For $H^1(A)$ this follows from Tate (1966),
Theorem 2(a), printed p.~140: $\mathbb Q[\pi_A]$
is the centre of the semisimple algebra
$\operatorname{End}_k(A)\otimes\mathbb Q$.
It is a product of number fields, so its Frobenius
element has square-free minimal polynomial;
dualizing or inverting preserves this property.
The identity $H^i(A)=\bigwedge^i H^1(A)$ shows
semisimplicity in every degree. Finite \'etale
Frobenius has finite order, and tensor products,
Tate twists and invariant direct summands of
semisimple single-operator representations remain
semisimple in characteristic zero.

Equality in (AQ.16) therefore gives an isomorphism
$R_\ell(M)\simeq R_\ell(N)$ as
$\mathbb Q_\ell[F]$-modules: each irreducible
polynomial in a semisimple operator's characteristic
polynomial records the multiplicity of its simple
$\mathbb Q_\ell[F]$-module. By (AQ.15) this is
an invertible element of
$\operatorname{Hom}_{\mathcal A_q}(M,N)\otimes\mathbb Q_\ell$.
Choose a rational basis $f_1,\ldots,f_s$ of this
Hom space. The polynomial
\[
 D(x_1,\ldots,x_s)=
 \det R_\ell\!\left(\sum_\nu x_\nu f_\nu\right)
\]
over $\mathbb Q_\ell$ is not zero. Since
$\mathbb Q$ is infinite, some rational tuple
has $D\ne0$; this elementary fact follows by
induction on the number of variables, using the
finite number of roots of a nonzero polynomial
in one variable. It gives a rational morphism
$f:M\to N$ with invertible realization.
By full faithfulness its realization's inverse
comes from a morphism in
$\operatorname{Hom}_{\mathcal A_q}(N,M)\otimes\mathbb Q_\ell$.
The equations $gf=1_M$ and $fg=1_N$ form a
finite linear system over $\mathbb Q$ in the
coefficients of $g$. A rational linear system
solvable after extension to $\mathbb Q_\ell$
is solvable over $\mathbb Q$, by row reduction.
Thus $f$ has a rational inverse. The converse
in (AQ.16) is immediate.

Finally a reciprocal root contributed by weight
$w$ has complex absolute value $q^{-w/2}$.
Weights of different parity cannot contribute
the same root. Hence the numerator and denominator
of (AQ.17) are coprime, each with constant term
one. Equality of those rational functions
recovers their numerator and denominator separately,
and their product is $P_M(t)$. This reduces
(AQ.17) to (AQ.16) and completes the proof.
\end{proof}
No essential-surjectivity claim onto all Frobenius
representations is needed here. In particular,
this recognition theorem does not identify an
arbitrary smooth projective motive with an object
of $\mathcal A_q$, does not retain integral lattices,
and does not classify mixed extensions. It is a
genuine strengthened answer to the abelian sector
of cote 15, supplied by the new cycle-algebraicity
theorem rather than by analytic zero location.

\paragraph{AQ-RATIONAL-FIBRE: an explicit neutralization obstruction.}
Even this now-established abelian category need not
admit a fibre functor with values in rational vector
spaces. This is a concrete restriction on any reading
of the archive's motivic Galois group as a neutral
group over $\mathbb Q$. It is not a refutation of a
group defined after scalar extension or of a motivic
gerbe formulation.
\begin{proof}
Consider $E:y^2=x^3-x$ over $\mathbb F_3$, then
base-change it to $k=\mathbb F_9$. It is smooth:
the derivative of $x^3-x$ is $-1$ in characteristic
three. For each of $x=0,1,2$, the right side is
zero, so there is one affine point; with infinity,
$\#E(\mathbb F_3)=4$. Thus its $3$-Frobenius
polynomial on $H^1$ is $T^2+3$. The square of
Frobenius is $-3$ on $H^1$, and injectivity of
the Tate-module action of endomorphisms implies
that the $9$-Frobenius endomorphism is $[-3]$.
Tate (1966), Theorem 2(d), p.~140, identifies
\[
 D=\operatorname{End}_{\mathbb F_9}(E)\otimes\mathbb Q
   =D_{3,\infty},
\]
the quaternion division algebra ramified at $3$
and infinity. Alternatively, the main theorem
on p.~134 gives $\dim_{\mathbb Q}D=4$, since
Frobenius is scalar on a two-dimensional Tate
module; every nonzero endomorphism of an elliptic
curve is an isogeny, so this algebra is a division
algebra.

Put $M=h^1(E)$ in $\mathcal A_9$. Pullback
endomorphisms give an injective algebra map
$D^{\rm op}\to\operatorname{End}(M)$.
By (AQ.15), the target has rational dimension
four, since its scalar extension is
$M_2(\mathbb Q_\ell)$. This map is therefore
an isomorphism, so $\operatorname{End}(M)$ is
a quaternion division algebra.

Use the usual parity-corrected symmetry on
numerical motives, for which the ordinary
$\ell$-adic realization is a symmetric tensor
functor. The categorical trace of $1_M$ is
then $2$: its image is the ordinary trace of
the identity on $H^1(E,\mathbb Q_\ell)$,
and the scalar map $\mathbb Q\to\mathbb Q_\ell$
is injective. A rational fibre functor would
preserve this trace, hence send $M$ to a
two-dimensional rational vector space. Its
faithfulness would embed the four-dimensional
division algebra $\operatorname{End}(M)$ into
$M_2(\mathbb Q)$. Equal dimensions force this
embedding to be an isomorphism. That is
impossible because $M_2(\mathbb Q)$ has
nonzero noninvertible elements. This proves
the asserted obstruction.

It also exhibits an object restriction invisible
in a list of eigenvalues alone. Here
$P_M(t)=(1+3t)^2$. There cannot be an object
$L\in\mathcal A_9$ with $P_L(t)=1+3t$:
otherwise (AQ.16) would give $L\oplus L\simeq M$,
but the former has a nontrivial idempotent
endomorphism and the latter's division algebra
has none. The eigenvalue $-3$ occurs, as
AQ-WEIL-EXIST predicts, but its multiplicity
cannot be discarded in a rational motivic
classification.
\end{proof}

\paragraph{AQ-WEIL-NOTATION: the recovered $Q'$ and the norm-one subgroup.}
The official cote 15 PDF has now been inspected beyond the
supplied excerpts.  Physical page 94, printed footer 93,
opens with the explicit definition
\emph{``$Q'$: cl\^oture tot[alement] r\'eelle de $\mathbb Q$''};
the bracket expands the handwritten abbreviation.
Thus $Q'=\mathbb Q^{\rm tr}$, the field of totally real
algebraic numbers, and the $i$ in $Q'(i)$ means
$\sqrt{-1}$, distinct from the integral weight index.
Physical page 98/footer 97 discusses a totally imaginary
quadratic extension of a totally real field and writes
$\lambda\overline\lambda=q^w$, followed by the
normalization $\lambda=q^{w/2}u$ with
$u\overline u=1$.
Physical page 100/footer 99 contains the boxed correction
requiring every conjugate to have absolute value $q^{w/2}$,
as already preserved in the supplied transcription.
These recovered norm and purity conditions identify the
weight-zero interpretation used in AQ-WEIL-LIMIT.

More precisely, let
\[
 \mathbb Q^{\rm CM}=\bigcup_{L\ {\mathrm{CM}}}L
 \quad\hbox{inside the fixed }\overline{\mathbb Q}\subset\mathbb C.
\]
Then
\[
 \mathbb Q^{\rm CM}=\mathbb Q^{\rm tr}(\sqrt{-1}).
 \tag{AQ.WN.1}
\]
For a prime power $q=p^n$, the group of algebraic
$p$-units with all conjugate absolute values $q^{w/2}$,
for some integer $w$, is exactly
\[
 \left\{\lambda\in\mathbb Q^{\rm tr}(\sqrt{-1})^*:
   \lambda\text{ is a $p$-unit},\quad
   \lambda\overline\lambda=q^w\text{ for some }w\in\mathbb Z
 \right\}.                                                   \tag{AQ.WN.2}
\]
Its weight-zero subgroup is $W_0(p)$, independent of $n$.
Thus the historical group $M(q)$ and its weight-zero part
agree with $W(q)$ and $W_0(p)$ when the explicit
norm/purity conditions on pages 98 and 100 are retained.

\begin{proof}
For a CM field $L$, let $L^+$ be its totally real
subfield.  Choose $b\in L$ with
$\delta=b-\overline b\ne0$.  Then
$\delta^2\in L^+$, and $a=-\delta^2$ is totally positive:
every embedding sends $\delta$ to a nonzero purely
imaginary number.  The algebraic number $\sqrt a$
is totally real, because all its conjugates are real
square roots of positive conjugates of $a$.
Hence
$L=L^+(\delta)\subset\mathbb Q^{\rm tr}(\sqrt{-1})$.
Conversely any element of the latter field belongs to
$F(\sqrt{-1})$ for a finite totally real field $F$,
and $F(\sqrt{-1})$ is a CM field.  This proves (AQ.WN.1),
including the assertion that the displayed union is a field.

Let $\lambda$ have the absolute values in (AQ.WN.2),
and let $L$ be the Galois closure of $\mathbb Q(\lambda)$.
Write $c$ for complex conjugation on $L$.
For every $\sigma\in\operatorname{Gal}(L/\mathbb Q)$,
\[
 c(\sigma\lambda)=q^w/(\sigma\lambda).
\]
For every $\tau$ in the same Galois group,
$c\tau$ and $\tau c$ consequently agree on all the
conjugates $\sigma\lambda$, which generate $L$.
Thus $c$ is central.
If $c=1$, the Galois field $L$ is totally real.
Otherwise $L^+=L^{\langle c\rangle}$ is totally real,
$[L:L^+]=2$, and $L$ has no real embedding, so $L$ is CM.
In either case (AQ.WN.1) places $\lambda$ in
$\mathbb Q^{\rm tr}(\sqrt{-1})$; the norm identity follows
from its specified absolute value at the fixed embedding.

Conversely this field is Galois over $\mathbb Q$, and
complex conjugation fixes $\mathbb Q^{\rm tr}$ and
negates $\sqrt{-1}$.  It therefore commutes with every
embedding.  Applying any embedding to
$\lambda\overline\lambda=q^w$ gives
$|\sigma\lambda|^2=q^w$, as required.
The $p$-unit condition is the same on both sides,
so (AQ.WN.2) follows.  Setting $w=0$ gives the claimed
weight-zero subgroup, independently of $q=p^n$.
\end{proof}

Finite-prime valuations alone must not be substituted
for the norm/purity condition in this identification.
For example $\varepsilon=2+\sqrt3$ and its inverse
$2-\sqrt3$ are algebraic integers, so $\varepsilon$
is a unit at every finite place.  In particular every
finite-prime sum $v_{\mathfrak p}(\varepsilon)+
v_{\overline{\mathfrak p}}(\varepsilon)$ is zero,
for every chosen $p$.  But
$\varepsilon\overline\varepsilon=\varepsilon^2\ne1$
and its two real conjugates have unequal absolute values.
No positive power is a weight-zero Weil number.
This checks the limitation of a \emph{valuation-only}
paraphrase of the margin on 15/102; it does not assert
that the manuscript's terminology \emph{unit\'e absolue}
itself discards archimedean conditions.
The corrected group interpretation uses the explicit
norm condition actually present on the preceding sheet.

\paragraph{AQ-WEIL-LIMIT: the erased question at 15/104.}
The text \emph{Trouve-t-on $P^0/U$ tout entier?}
and its subsequent root-extraction test are inside a
deletion in the supplied transcription. In the standard
Weil-number interpretation its algebraic assertion is true.
Let $W_0(p)$ be the multiplicative group of algebraic
$p$-units whose complex conjugates all have absolute
value one. Then the directed system with transition
maps $u\mapsto u^m$ has limit
\[
 \varinjlim W_0(p)=W_0(p)\otimes_{\mathbb Z}\mathbb Q
                 \simeq W_0(p)/\mu_\infty .
\]
\begin{proof}
For any abelian group $A$, the system indexed by
positive integers, with multiplication by $m$ from
index $n$ to index $mn$, has colimit
$A\otimes\mathbb Q$: the compatible map sends
the class of $a$ at index $n$ to $a/n$.
Every root of an element of $W_0(p)$ is again
a $p$-unit with all conjugate absolute values one.
Thus $W_0(p)$ is divisible, and its torsion subgroup
is exactly $\mu_\infty$. Its quotient by torsion
is uniquely divisible, giving the last isomorphism.
\end{proof}
The recovered original definition $Q'=\mathbb Q^{\rm tr}$
and the preceding explicit norm-one condition identify
this weight-zero interpretation by AQ-WEIL-NOTATION.
The literal deleted calculation therefore has the stated
affirmative answer with its norm/purity conditions retained.
Its identification with the group of arbitrary numerical motives
still requires the separate Tate and generation assertions.
For the coefficient-extended abelian sector, CC-AG04--06 now
complete the categorical base-field passage: finite-field models,
numerical correspondences and idempotents descend, and their
filtered categorical colimit is the category of finite-dimensional
vector spaces graded by Frobenius germs. The precise result is
recorded immediately below; the rational fibre obstruction remains.

\paragraph{Coefficient-extended group: the completed abelian-sector answer.}
The companion proof CC-AG01--03, in
\nolinkurl{research/consequences_criteria_abelian_group.tex},
now strengthens the rational recognition result to a full
category and group description. Fix
$E=\overline{\mathbb Q}_\ell$ and an embedding
$\overline{\mathbb Q}\hookrightarrow E$. Let $W(q)$ be the
group of algebraic $p$-units whose conjugate absolute values
are $q^{w/2}$ for one integer $w$. With the parity-corrected
symmetry and idempotent completion after extending Hom spaces,
the result is
\[
 \operatorname{Kar}(\mathcal A_q\otimes_{\mathbb Q}E)
  \simeq \operatorname{Gr}^{\mathrm{fin}}_{W(q)}(E),\qquad
 \operatorname{Aut}^{\otimes}(\omega_E)=D_E(W(q)).
\]
The target means finite-dimensional vector spaces graded by
$W(q)$, with multiplication of labels as tensor product;
$D_E(M)=\operatorname{Spec}E[M]$ for an abelian group $M$.
Choosing $\sqrt q$ gives
$D_E(W(q))\simeq D_E(W_0(p))\times\mathbb G_{m,E}$,
with Tate character $(1,-2)$. The full proof uses the
Hom-space theorem proved above, Honda occurrence, and
lifting of the individual eigenspace projectors over $E$.
It also proves the distinguished Frobenius has Zariski-dense
powers. The finite-field group is not a protorus, since its
character group contains $\mu_\infty$. This answers the
abelian-sector group question with explicit coefficient and
fibre conventions, while preserving the rational quaternion
obstruction and distinguishing passage to
$\overline{\mathbb F}_p$ from extension of coefficients.
The subsequent complete proof CC-AG04--06, in
\nolinkurl{research/consequences_criteria_abelian_germs.tex},
settles the separate passage to $\overline{\mathbb F}_p$
within this same abelian sector. If $W_p$ denotes the
Frobenius-germ character group, its result is
\[
 \mathcal A_{\overline{\mathbb F}_p,E}
  \simeq\operatorname{Gr}^{\mathrm{fin}}_{W_p}(E),\qquad
 W_p\simeq\mathbb Z\oplus(W_0(p)\otimes_{\mathbb Z}\mathbb Q),
 \qquad\operatorname{Aut}^{\otimes}(\omega_E)=D_E(W_p).
\]
Base extension from any fixed $\mathbb F_q$ corresponds to the
surjection $W(q)\to W_p$ with kernel $\mu_\infty$ and is
essentially surjective. Thus the geometric tensor group is
precisely the identity component of $D_E(W(q))$, is a protorus,
and the quotient is the constant profinite group
$\widehat{\mathbb Z}$. The proof checks descent of varieties,
cycles, numerical equivalence and idempotents, including the
full filtered colimit of Hom spaces after coefficient extension.
There is no remaining abelian-sector categorical germ gate.
Neither coefficient extension nor this base extension proves
that arbitrary numerical motives are of abelian type, and
neither provides the rational neutral fibre functor disproved
above. The recovered historical notation remains as in
AQ-WEIL-NOTATION.

The rational descent requirement is also explicit, rather than an
unanalysed coefficient gate. CC-AG07--08, in
\nolinkurl{research/consequences_criteria_rational_descent.tex},
prove the finite-transition exact sequence
\[
 1\longrightarrow\mu_m\longrightarrow W(q)
  \xrightarrow{\alpha\mapsto\alpha^m}W(q^m)
  \longrightarrow1.
\]
Its cokernel is zero, and finite-field base extension is
essentially surjective after coefficients are extended to $E$.
For either fixed-field labels or Frobenius germs, an $E$-graded
object descends to an actual rational object exactly when its
multiplicities are Galois invariant and, on each Galois orbit
$O$, divisible by the Schur index $d_O$. The index is the least
common multiple of the denominators of Milne's local Brauer
invariants for the associated simple cohomological motive.
The centre for a germ is the stabilized field
$\mathbb Q(\alpha^N)$, with $N$ sufficiently divisible.
The resulting rational isomorphism class is unique, and is
$\bigoplus_O S_O^{\oplus n_O/d_O}$, with endomorphism algebra
$\prod_O M_{n_O/d_O}(D_O)$. Here $D_O$ is the central division
endomorphism algebra of $S_O$; all symbols and local-invariant
sign conventions are specified in the companion proof.
This does not supply a rational neutral fibre functor.
Indeed finite-field base extension need not be essentially
surjective over rational coefficients: the rank-one sign Artin
object over $\mathbb F_{q^2}$ would require a rational rank-one
Frobenius scalar $\alpha$ with $\alpha^2=-1$ to descend to
$\mathbb F_q$. The exact orbit/index test, rather than any
claim of rational neutrality, is the completed rational
abelian-sector coefficient answer.

The rational essential image of base-field extension is now
also exact, by CC-AG09--10 in
\nolinkurl{research/consequences_criteria_rational_basechange.tex}.
For a target Galois orbit $O'$, let $O$ range over source orbits
mapping onto it, and put
\[
 c_O=|O|/|O'|,\qquad r_O=c_Od_O/d_{O'},\qquad
 H_B(O')=\left\{\sum_O a_Or_O:a_O\ge0,
                         \text{ finite support}\right\}.
\]
Each $r_O$ is an integer: the rational simple $S_O$ becomes
$r_O$ copies of $S_{O'}$. A target object's multiplicity of
$S_{O'}$ must, and may, belong to this precise additive monoid.
For a finite extension of degree $m$, the eligible labels are
the finitely many roots of $\prod_{\beta\in O'}(t^m-\beta)$.
For the geometric extension the monoid is finitely generated,
but the proof supplies no effective bound on root-of-unity
orders. The explicit example $q=3$, $K=\mathbb Q(\sqrt{-5})$,
$\beta=2+\sqrt{-5}$ has exactly even multiplicities in the
image of both $\mathbb F_3\to\mathbb F_9$ and
$\mathbb F_3\to\overline{\mathbb F}_3$, for its specified
rank-two target simple. Thus the rational essential image,
its failure of surjectivity, and the extended-coefficient
surjectivity are separately identified.

The last effectiveness limitation is reduced further by CC-AG11--12,
in \nolinkurl{research/consequences_criteria_integral_slope_descent.tex}
and \nolinkurl{research/consequences_criteria_finite_descent.tex}.
For a stabilized germ $\xi=[\pi,b]$ of integral slopes, write
$K=\mathbb Q(\pi^N)$ for its stable centre, and let $h$ generate
the subgroup of integers $n$ for which the prescribed fractional
ideal $J^n$ has a generator with every complex absolute value
$p^{nw/2}$. Then $h\mid b$, a finite polynomial-root test in
$K$ computes $h$, and the geometric multiplicity monoid for
base field $\mathbb F_{p^a}$ is exactly
$\frac{h}{\gcd(h,a)}\mathbb N$.

For every specified target object, no slope restriction is
needed for a finite decision procedure. For a simple target
component $S_\xi^{\oplus r}$ of rank
$D=r d_\xi[K:\mathbb Q]$, define
\[
 R_D=\sum_{e=1}^D e\prod_{j=1}^e
       \left(2\binom ej+1\right),\qquad
 N_D=\operatorname{lcm}(1,\ldots,R_D).
\]
Every possible source simple over $\mathbb F_{p^a}$ has label
among the roots of the explicit rational polynomial
\[
 \prod_{\sigma:K\hookrightarrow\overline{\mathbb Q}}
       \left(X^{bN_D}-\sigma(\pi)^{aN_D}\right).
\]
Factoring it and applying the already proved orbit/Schur-index
formula gives a finite list of positive multiplicities $m_j$.
The finite equation $\sum_j u_jm_j=r$, with $u_j\ge0$,
decides descent and enumerates every source isomorphism class.
The bound is rank-dependent: no uniform cutoff generating
all multiplicity monoids is inferred. This classifies objects
in the established abelian category and does not construct
arbitrary algebraic cycles or establish generation of all
numerical motives by that category. The complete deductions
have author and parent whole-proof readings recorded in their
companion ledgers; this paragraph records their exact scope.

For the prime-field source, CC-AG13 in
\nolinkurl{research/consequences_criteria_prime_field_descent.tex}
now gives the entire geometric monoid in every slope sector.
For a stabilized germ $\xi=[\pi,b]$ with stable centre $K$,
let $h_\xi$ be the positive generator of
\[
 \{n\in\mathbb Z:\ \gamma^b/\pi^n\in\mu(K)
                    \text{ for some }\gamma\in K^\times\}.
\]
Then $h_\xi\mid b$, and the same finite root test over $K$
computes it. The geometric Schur index satisfies
$d_\xi\mid h_\xi$, and the exact multiplicity monoid for
$\mathbb F_p\to\overline{\mathbb F}_p$ is
\[
                    (h_\xi/d_\xi)\mathbb N.
\]
A single source simple realizes its least positive element;
the proof treats real odd-weight germs separately. This whole-monoid
formula is for the prime-field source, whereas the preceding
rank-dependent finite test works for every specified finite
source field. Both statements retain the established abelian,
Artin and Tate numerical category.

The later CC-AG14--16 close the remaining whole-monoid
calculation while retaining its genuine arithmetic complexity.
The exact monoids $\langle2,3\rangle$ and $3\mathbb N$ occur
with the same stabilized centre $\mathbb Q(\sqrt{-23})$,
representative exponent $h=9$, geometric index $d=3$ and
source exponent $a=3$, in characteristics $2$ and $163$
respectively. Thus a cyclic or elementary norm-gcd shortcut
is false. BR-15-PRIME-SUPPORT first confines the needed
root-of-unity orders to primes dividing $h$. The complete
checked proof CC-AG16 in
\nolinkurl{research/consequences_criteria_uniform_descent.tex}
then bounds those tower heights independently of the target
multiplicity. Specifically, for $\xi=[\gamma,h]$, it uses
the splitting field $E$ of $\operatorname{Norm}_{K/\mathbb Q}
(X^h-\gamma^a)$, its degree $M$, and explicit exponents
$B_\ell=c_\ell+v_\ell(M!)+v_\ell(h)$ to form
$N=\prod_{\ell\mid h}\ell^{B_\ell}$. Its fully stated
local constants $c_\ell$ and their proof are in that module.
The finite polynomial
\[
       \operatorname{Norm}_{K/\mathbb Q}
                  (X^N-\gamma^{aN/h})
\]
lists source labels whose actual positive transition
multiplicities generate the entire geometric monoid.
Factoring it and using the local Schur indices computes
minimal generators, gcd, gaps and a conductor relative to
that gcd. The earlier rank-dependent and prime-field
formulas remain valid, but their former tower-height
limitation is now removed in the full established abelian
sector. No practical complexity bound, arbitrary-cycle
construction or rational neutral fibre functor follows.

CC-AG17 further computes the complete profile over all
finite source fields.  For a stabilized germ $\xi=[\gamma,h]$,
its geometric multiplicity monoids satisfy
\[
 H_{a+h}(\xi)=H_a(\xi)=H_{\gcd(a,h)}(\xi),
                   \qquad H_h(\xi)=\mathbb N.
\]
The aligned-label bijection $\alpha\mapsto\alpha\gamma$
between exponents $a$ and $a+h$ preserves the centre,
Schur index and actual geometric multiplicity.  Finite
base extension then gives the gcd reduction.  Thus the
finite divisor table $\{H_d(\xi):d\mid h\}$, computed by
CC-AG16, gives every source exponent.  For a specified
geometric object $Z=\bigoplus_jS_{\xi_j}^{\oplus r_j}$,
one obtains a finite antichain $\mathcal D_Z$ of divisors
of $\operatorname{lcm}_j h_{\xi_j}$ such that
\[
 Z\text{ descends from }\mathbb F_{p^a}
 \quad\Longleftrightarrow\quad
           d\mid a\text{ for some }d\in\mathcal D_Z.
\]
The complete proof is
\nolinkurl{research/consequences_criteria_descent_profile.tex};
it has also been independently read by this module's author.
The profile retains all positive-monoid gaps and concerns
geometric isomorphism, without fixing an additional finite
target-field model.

\paragraph{Concrete spectral targets and the remaining arbitrary-variety step.}
The companion CC-XC11 and CC-XC13 now give, for every proper
smooth finite-field scheme, a unique rational abelian-sector
object with the semisimplification of each of its Frobenius
realizations. The index divisibility is obtained from the
crystalline primary-block determinant and the proper-smooth
alternating-form theorem in odd degree; it is not an assertion
that the arbitrary variety has semisimple Frobenius.
CC-XC14 further gives a terminating construction from a promised
polynomial list: it returns explicit abelian varieties and
finite \'etale schemes together with rational numerical
projector cycles, without Tate twists. Numerical ranks and
complementary intersection bases on abelian products make
endomorphisms and projector tests effective; intersection traces
then test the complete degreewise Frobenius polynomial.
Honda occurrence and the rational classification prove that the
dovetailed search terminates. Smooth projective equations
supply the input list by the stated terminating zeta algorithm;
for proper nonprojective schemes the polynomials are supplied
as spectral input. The complete checked proofs are
\nolinkurl{research/consequences_criteria_proper_parity.tex}
and \nolinkurl{research/consequences_criteria_effective_target.tex}.
The separate cross-correspondences identifying an arbitrary
variety with this fixed target remain the substantive cycle
condition. No general termination of that subsequent search is
asserted.

\subsection{AQ-TWIST: the weight-index correction at 18/009}
The sheet's conventions give
\[
 W_j(V(n))=W_{j+2n}(V),\quad F^j(V(n))=F^{j+n}(V),\quad
 \operatorname{gr}^W_j(V(n))
 =(\operatorname{gr}^W_{j+2n}V)(n).                             \tag{AQ.6}
\]
The last equality follows by taking the quotient of the
first two consecutive weight steps, with the shifted Hodge
filtration. The transcribed unshifted graded formula fails
for $V=\mathbb Q(0)$, $n=-1$, $j=2$: the left side is
$\mathbb Q(-1)$ and the unshifted right side is zero.

\subsection{AQ-REAL: the even-Tate assertion at 18/038 is false}
The sheet claims that $f_{\infty,M}=1$ forces a sum of
tensor powers of $\mathbb Q(2)$.
Let $A$ be the sign Artin motive over $\mathbb R$ and
$L=A\otimes\mathbb Q(-1)$. Its real involution is
$(-1)(-1)=1$, but its weight $2$ is not the weight $-4r$
of any $\mathbb Q(2)^{\otimes r}$. This refutes the claim.
In the semisimple real Artin--Tate category the correct
objects are exactly finite sums of $L^{\otimes r}$,
$r\in\mathbb Z$.
\begin{proof}
The simple objects are $A^\epsilon\otimes\mathbb Q(n)$,
$\epsilon\in\{0,1\}$, with real involution
$(-1)^{\epsilon+n}$. It is trivial exactly when
$\epsilon\equiv n\pmod2$, in which case this object is
$L^{\otimes(-n)}$. Direct-sum decomposition proves both
directions.
\end{proof}
The classification is asserted here only for Artin--Tate
motives. Extending it from Hodge types to all pure motives
would require an additional algebraicity input.

\subsection{AQ-PERIOD-STRONG: a counterexample to 18/011}
The strong transcendence conjecture on 18/011 predicts
$\operatorname{trdeg}_k k(\alpha,\beta)=g+u$.
In the usual reading $T_{\rm Hdg}=\operatorname{gr}_F T_{\rm dR}$
and $Q$ parametrizes normalized splittings, this fails.
Take $E:y^2=x^3-x+1$ over $\mathbb Q$.
Its $j=-6912/23$ is not integral, so it is non-CM;
for $M=h^1(E)$ one has $G=\mathrm{GL}_2$, $g=4$, $u=1$.
Here we use the classical integrality of CM
$j$-invariants, recalled in Habegger (2014),
\S1, p.~1, and the elliptic motivic-group
identification in Kreutz--Shen--Vial,
Proposition 5.14.
Here $\alpha$ is the comparison from Betti to de Rham
cohomology, $T_{\rm Hdg}=\operatorname{gr}_F T_{\rm dR}$,
$\beta$ is the Hodge splitting normalized to induce the
identity on that associated graded, and
$\gamma=(\beta\alpha)^{-1}$, as in the source triangle.
The parabolic preserving its Hodge line has
one-dimensional unipotent radical, explaining $u=1$.
\begin{proof}
Choose rational Betti eigenvectors for the real involution
and a rational de Rham basis with Hodge line $\mathbb Ce_1$.
Write
\[
 P=\begin{pmatrix}a&b\\c&d\end{pmatrix},\quad
 \delta=ad-bc\ne0,\quad C=\operatorname{diag}(1,-1).
\]
The opposite Hodge line is $PCP^{-1}\mathbb Ce_1$,
generated by $(ad+bc,2cd)^t$. Both $c,d$ are nonzero:
otherwise a real-involution eigenline would be the
Hodge line, forcing it to equal its opposite.
Hence the normalized splitting is
\[
 \beta=\begin{pmatrix}1&-(ad+bc)/(2cd)\\0&1\end{pmatrix}.
\]
It is rational in $P$, so $k(P,\beta)=k(P)$ has
transcendence degree at most four, not $g+u=5$.
Moreover
\[
 \gamma=(\beta P)^{-1}
 =\begin{pmatrix}d/\delta&1/(2c)\\-c/\delta&1/(2d)\end{pmatrix},
 \qquad \gamma_{11}\gamma_{22}+\gamma_{12}\gamma_{21}=0.
\]
Thus its transcendence degree is at most three,
also contradicting the predicted value $g=4$.
\end{proof}
This refutation does not settle the ordinary non-CM
elliptic period conjecture or algebraicity of its
Hodge-splitting scalar.

\subsection{AQ-PERIOD-GENERATION: 18/013(b) is a known theorem}
Question (a) asks for period-field transcendence degree
$\dim G$, namely four without CM or two with CM.
Question (b) asks whether the group generated by a
translate of the period locus is all of $G$.
The answer to (b) for one elliptic curve is yes:
Kreutz, Shen, and Vial (2026), \emph{De Rham--Betti classes
with coefficients}, arXiv:2206.08618v3, Proposition 3.1
and Theorem 5.13(ii) identify the generated period torsor
with the motivic torsor; Proposition 5.14 identifies
the elliptic motivic group. The coefficient restriction
in Theorem 5.13(ii) is satisfied for a single elliptic
curve. This known result is cited, not reproved.
It is weaker than density of the period locus.

Question (a) is known for CM curves by Chudnovsky (1980),
\S2A, Theorem 1 and Corollary 2, printed p.~343
in the 1978 Helsinki ICM proceedings. In particular,
a nonzero CM period and $\pi$ are algebraically
independent. The non-CM case remains the full
period-conjecture gate. Question (c) asks whether the
Hodge-splitting scalar can be algebraic without CM.
In the normalization $y^2=4x^3-g_2x-g_3$,
put $\omega=dx/y$, $\xi=x\,dx/y$, take full periods
$\omega_1,\omega_2$ with $\operatorname{Im}
(\omega_2/\omega_1)>0$, and write the periods of
$\xi$ as $-\eta_1,-\eta_2$. The opposite Hodge line
is generated by $\xi+h\omega$. Its periods must be
proportional to $\overline{\omega_1},
\overline{\omega_2}$, which gives
\[
 h=\frac{\eta_1\overline{\omega_2}-\eta_2\overline{\omega_1}}
          {\omega_1\overline{\omega_2}-\omega_2\overline{\omega_1}}
   =\frac{\pi^2}{3\omega_1^2}E_2^*(\tau).
\]
Here $\tau=\omega_2/\omega_1$ and
$E_2^*(\tau)=E_2(\tau)-3/(\pi\operatorname{Im}\tau)$.
The second equality uses
$\eta_1\omega_2-\eta_2\omega_1=2\pi i$ and
$\eta_1\omega_1=\pi^2E_2(\tau)/3$.
The CM conclusions in (a) and (c) can be separated
from the remaining non-CM problem.
\begin{proof}[The CM period field and algebraicity of the splitting]
Take a nonreal CM endomorphism $u$, defined over
$\overline{\mathbb Q}$. In the algebraic de Rham
basis its action has the form
$u^*\omega=u\omega$ and
$u^*\xi=a\omega+\overline u\xi$, with
$a\in\overline{\mathbb Q}$. The opposite Hodge line
is the $\overline u$-eigenline, so
\[
 a+hu=\overline u h,\qquad
 h=\frac{a}{\overline u-u}\in\overline{\mathbb Q}.
\]
The period ratio $\tau$ is imaginary quadratic.
The preceding formula for $h$ and the Legendre
relation give
\[
 \omega_2=\tau\omega_1,\qquad
 \eta_1=h\omega_1+\frac{\pi}{\omega_1\operatorname{Im}\tau},
 \qquad \eta_2=\tau\eta_1-\frac{2\pi i}{\omega_1}.
\]
Thus the full period and quasi-period field is
$\overline{\mathbb Q}(\omega_1,\pi)$; the reverse
inclusion follows from the Legendre relation.
Chudnovsky's cited Corollary 2 makes its
transcendence degree exactly two.
\end{proof}
The remaining non-CM assertion is the nonholomorphic weight-two
Eisenstein-transcendence question attributed to Katz
in Waldschmidt (2004), p.~277, immediately after
Conjecture 3.22. That passage states the same remaining
condition for a non-CM lattice with algebraic invariants.
Its non-CM transcendence is not established by the
supplied proofs. AQ-FINITE cannot supply that missing
algebraic-independence assertion.

\paragraph{AQ-PERIOD-2026: current-source and normalization check.}
The current primary literature does not identify the
three different weight-two functions in these questions.
Kreutz--Shen--Vial, arXiv:2206.08618v3, Introduction
and \S4.2, retain the full period conjecture separately
from their generated-torsor theorem. Mohajer (2026),
arXiv:2611.25399v1, \S1.1, equation (1.3), pp.~2--3,
still formulates the non-CM complex four-period
transcendence-degree assertion as conjectural.
His Theorems 4.2 and 4.4 concern comparison periods
in $B_{\rm dR}$ and their normalized images in
$\mathbb C_p$, respectively. His Proposition 2.1
concerns zeros of the holomorphic $E_2$, whereas
the scalar above involves the almost holomorphic
$E_2^*=E_2-3/(\pi\operatorname{Im}\tau)$.

Sebbar and Besrour (2026), arXiv:2611.14502v1,
\S8.4, Remark 8.11, Proposition 8.12 and
Question 8.13, pp.~32--33, consider the ordinary
$p$-adic Katz form $E^{\rm Katz}_2$. They leave
the existence of a non-CM zero with algebraic
$j$-invariant over $\mathbb Q$ as a question;
their nonvanishing proposition assumes the crystalline
period conjecture of Ancona--Fr\u a\c til\u a.
Algebraicity over a finite extension of $\mathbb Q_p$
does not answer algebraicity over $\mathbb Q$.
Thus these November 2026 results do not supply the
missing complex algebraic-independence assertion in
18/013(a), or the non-CM algebraicity classification of
the complex splitting scalar in 18/013(c). This is a
source-scope conclusion, not a proof of logical
nonimplication from the adopted analytic theorems.

\paragraph{AQ-CM-PERIOD-2026: the new rank-one comparison results.}
Kahn and Otsubo (2026), arXiv:2611.13723v1,
Corollary 3.21, p.~24, show that equal Hecke characters
give equal complex period classes for rank-one abelian-type
$K$-coefficient motives over a number field $k\supset K^{\rm gal}$.
Here the class lies in
$(\mathbb C\otimes_{\mathbb Q}K)^*/(k\otimes_{\mathbb Q}K)^*$
(\S2.4, pp.~6--7).
Remark 3.22, p.~25, relates the converse, up to a finite-order
character, to Rohrlich's conjecture.
For rank-one algebraic motives, equality of crystalline period
classes at almost all places implies equality of Hecke characters;
the converse holds under the preceding abelian-type and base-field
hypotheses (Corollary 3.35, pp.~30--31).
Corollary 3.25, p.~26, proves Mumford--Tate when the rank-one
$K$-motive becomes algebraic after finite extension.
For an abelian CM field $K$ of conductor $m$ and
$k\supset\mathbb Q(\mu_m)$,
Theorems 4.36 and 4.39, pp.~46--47,
refine Gross--Deligne, and Theorem 4.44, p.~49,
gives the corresponding $p$-adic beta-period formula.
These are positive comparison results with specified hypotheses.

For a non-CM elliptic curve, $h^1(E)$ has rational rank two
and endomorphism algebra $\mathbb Q$.  A rank-one structure
over a coefficient field would require a quadratic field to
embed in that endomorphism algebra, which is impossible.
The determinant is a Tate object and records only the
determinant period.  Thus these rank-one statements do not
give the four-period algebraic independence or the
complex $E_2^*$ assertion in 18/013.

\paragraph{AQ-KO-RANK: a correction to an ancillary dimension claim.}
With Definition 2.5, p.~8, of the same paper, the dimension
sentence in Lemma 4.6, p.~34, requires one additional rank:
for $K=\mathbb Q(\mu_m)$ with $m>2$,
\[
 \operatorname{rank}_{\mathbb Z}\Xi(K)
       =\frac{\varphi(m)}2+1,
 \tag{AQ.KO.1}
\]
where the paper prints $\varphi(m)/2$.
\begin{proof}
Put $r=\varphi(m)/2$ and choose one embedding
$\sigma_j$ from each of the $r$ complex-conjugate pairs.
Since $K$ itself is a CM field, the additional CM-subfield
condition in that definition is automatic.
A function $\phi:\Sigma_K\to\mathbb Z$ belongs to
$\Xi(K)$ precisely when all the sums
$\phi(\sigma_j)+\phi(\overline\sigma_j)$ equal
one common integer $w$.
Consequently
\[
 \Xi(K)\longrightarrow\mathbb Z^{r+1},\qquad
 \phi\longmapsto(w,\phi(\sigma_1),\ldots,\phi(\sigma_r))
\]
is an isomorphism of abelian groups: the inverse sends
$(w,a_1,\ldots,a_r)$ to the function with values
$a_j$ and $w-a_j$ on the $j$th pair.
This proves (AQ.KO.1).  The map
$\phi\mapsto2\phi-w$ used in the cited rank argument
has the full group of constant functions as its kernel,
rather than being injective.  Already $m=3$ gives
$\Xi(\mathbb Q(\mu_3))=\mathbb Z^2$.
\end{proof}
This calculation concerns that rank sentence; the
period-comparison results above are cited with their
own hypotheses, and their complete external proofs
are not certified by this source-scope check.

\subsection{AQ-COCYCLE: the image-density statement at 18/059}
Let $P$ be an affine torsor under an algebraic group $G$
over $F\subset\mathbb C$, and $p\in P(\mathbb C)$.
The closure of the image of its automorphism cocycle
is the translated $F$-locus of $p$. Thus the image is
Zariski dense in $G_{\mathbb C}$ exactly when $p$ is
generic in $P$ over $F$.
\begin{proof}
Embed $P$ in affine space over $F$, and let $I$ be the
ideal over $\mathbb C$ of the
$\operatorname{Aut}(\mathbb C/F)$-orbit of $p$.
Each finite-dimensional degree-bounded part of $I$
is stable under the coefficientwise semilinear action.
Its unique reduced-row-echelon basis, relative to an
ordered monomial basis, must therefore have entries
fixed by every automorphism. The fixed field of
$\operatorname{Aut}(\mathbb C/F)$ is $F$: any nontrivial
embedding of $F(a)$ over $F$, for $a\notin F$, extends
first across a transcendence basis and then an algebraic
closure to an automorphism of $\mathbb C$.
Consequently every degree-bounded part has a basis
over $F$, and $I=(I\cap F[\mathbf x])\mathbb C[\mathbf x]$.
An $F$-polynomial vanishes on the orbit exactly when
it vanishes at $p$. Thus $I$ is precisely the ideal of
$Z_F(p)_{\mathbb C}$, proving the density assertion.
The torsor difference map $D_p:P_{\mathbb C}\to G_{\mathbb C}$
is an algebraic isomorphism and the cocycle is
$D_p(\sigma p)$, up to inversion according to convention.
Its closure is consequently $D_p(Z_F(p)_{\mathbb C})$.
\end{proof}
The generated-subgroup version is false for general
torsors. In $G=\mathbb G_m^2$ take $F=\mathbb Q$ and
$p=(t,t+1)$ with $t$ transcendental. Its locus is a curve.
The translated cocycle curve is
$u\mapsto(u/t,(u+1)/(t+1))$.
A character $x^a y^b$ identically one there would have
orders zero at $u=0$ and $u=-1$, giving $a=b=0$.
Thus this proper curve generates the whole torus.
This refutes an abstract torsor inference, without
claiming an arithmetic-motivic counterexample to the
period conjecture.

\subsection{AQ-CM: the hexagonal elliptic curve and a genuine GRH consequence}
The hexagonal real-elliptic classification recorded in SA22,
\S2(A.7), concerns the $j=0$ period lattice and its real
forms. It is not the Grothendieck--Teichm\"uller hexagon.
For example, the latter is the identity
\[
 f(z,x)z^m f(y,z)y^m f(x,y)x^m=1,\qquad
 z=(xy)^{-1},\quad m=(\lambda-1)/2,
\]
in a profinite free group, as in Boavida de Brito, Horel,
and Robertson (2019), Definition 7.2(II). The two uses of
\emph{hexagonal} name different mathematical data.

\paragraph{Hodge obstruction.}
For $E:y^2=x^3+1$, the motive $h^1(E)$ is not Artin--Tate,
even after adjoining its CM field $\mathbb Q(\sqrt{-3})$.
\begin{proof}
An Artin motive has Hodge type $(0,0)$, and $\mathbb Q(n)$
has type $(-n,-n)$ in our convention.
Tensor products, sums, subobjects, and quotients of such
Hodge structures have only diagonal types. But
$H^1(E)$ has nonzero types $(1,0)$ and $(0,1)$.
Base extension to an algebraic number field does not
change these complex Hodge types. This proves the claim.
\end{proof}
Thus the fact that the CM field is abelian does not turn
its weight-one Hecke character into a finite-order
Dirichlet character with an integral Tate twist.
No conclusion about the zeros of $L(h^1(E),s)$ is obtained
merely by substituting that field in the supplied
Dirichlet theorem.

\paragraph{Effective distribution of trace-zero primes.}
Let $a_p=p+1-\#E(\mathbb F_p)$, for $p>3$, and put
$N_0(x)=\#\{3<p\leq x:a_p=0\}$. Then
\[
 a_p=0\ \Longleftrightarrow\ p\equiv2\pmod3,\qquad
 N_0(x)=\tfrac12\operatorname{Li}(x)+O(\sqrt x\log x)            \tag{AQ.7}
\]
under the adopted Dirichlet GRH. The implied constant is
effective.
\begin{proof}
With the quadratic character extended by $\chi_p(0)=0$,
\[
 a_p=-\sum_{x\in\mathbb F_p}\chi_p(x^3+1).
\]
If $p\equiv2\bmod3$, cubing permutes $\mathbb F_p$,
so the sum is $\sum_u\chi_p(u+1)=0$.
If $p\equiv1\bmod3$, set $m=(p-1)/2$.
Euler's criterion and the binomial theorem show
\[
 a_p\equiv-\sum_{j=0}^m {m\choose j}
                 \sum_{x\in\mathbb F_p}x^{3j}
       ={m\choose (p-1)/3}\pmod p.
\]
The $j=0$ sum is $p=0$ modulo $p$.
Among the positive exponents $3j\leq3(p-1)/2$, the
only multiple of $p-1$ is $p-1$, and its power sum is
$-1$ modulo $p$. The displayed binomial coefficient
is nonzero modulo $p$ because all its factorial arguments
are less than $p$. This proves the equivalence.

Now $\zeta_{\mathbb Q(\sqrt{-3})}(s)=
\zeta(s)L(s,\chi_{-3})$, so the adopted theorem gives
GRH for this quadratic Dedekind zeta function.
Lagarias--Odlyzko (1977), Theorem 1.1, equation (1.3),
applied to its nontrivial Frobenius class, gives
\[
 \left|\pi(x;3,2)-\tfrac12\operatorname{Li}(x)\right|
 \leq c_1\left(\tfrac12\sqrt x\log(3x^2)+\log3\right),
 \qquad x>2.
\]
Here $c_1$ is their effective absolute constant.
For $x\geq3$, $N_0(x)=\pi(x;3,2)-1$ because $p=2$
is a bad prime of $E$. This proves (AQ.7).
\end{proof}
This gives an effective arithmetic statement about the
hexagonal curve using the quadratic splitting character.
It does not replace that curve's weight-one Hecke
$L$-function by a Dirichlet $L$-function.

\subsection{AQ-67-RANGE: the global isomorphism in 67/078 is false}
The literal question is \emph{Ces fl\`eches sont elles
des iso --- et si oui, comment expliciter un iso en
sens inverse ?} In the stable nodal $(0,4)$ case,
normalize both three-pointed rational components with
the node at $0$ and the other two marked points at
$1,\infty$. In the tangent-tensor normalization of
the capping construction, a tensor $q$ in the image
has absolute value equal to the product of the two
node-cap conformal radii. Consequently
\[
 |q|\leq16,\qquad q=17\ \hbox{has no preimage}.                  \tag{AQ.8}
\]
If the target is written with dual tangent tensors,
the equivalent statement uses the inverse tensor.
This normalization is fixed by the original sheets:
67/075 calls the propagation datum a nonzero basis of
$T_{x'}\otimes T_{x''}$; 67/078 writes its inverse
as a dual pairing and asks for mutually polar disks
in the two tangent lines. In coordinates, polarity
is $|q^{-1}|r'r''=1$, hence $r'r''=|q|$.
The target in 67/076--077 imposes no bound on this
nonzero tensor. We use the map without additional
marked-point disks in 67/077.

\begin{proof}
For completeness, the Koebe bound needed here follows
from the area argument. If
$F(z)=z+b_0+\sum_{n\geq1}b_nz^{-n}$ is univalent for
$|z|>1$, Green's formula gives the area inside
$F(|z|=r)$ as
$\pi(r^2-\sum_{n\geq1}n|b_n|^2r^{-2n})\geq0$.
Letting $r\downarrow1$ proves
$\sum_{n\geq1}n|b_n|^2\leq1$.
For a univalent $f(z)=z+a_2z^2+\cdots$ in the unit disk,
$g(z)=\sqrt{f(z^2)}=z+(a_2/2)z^3+\cdots$ is an odd
univalent function there. The branch exists because
$f(z^2)/z^2$ never vanishes in the disk. Equality of
two $g$-values forces equality of the squared
arguments under $f$, and oddness then proves
injectivity. Its reciprocal
$1/g(1/z)$ is univalent outside the disk, so the area
inequality gives $|a_2|\leq2$.
If $w\ne0$ is omitted by $f$, the normalized univalent
function $f/(1-f/w)$ has second coefficient $a_2+1/w$.
Thus $|1/w|\leq|a_2|+|a_2+1/w|\leq4$.
This is the Koebe quarter theorem (Duren, 1983,
Theorem 2.3, p.~31).

A node cap is the image of a univalent map $h$ from
the unit disk with $h(0)=0$ and conformal radius
$r=|h'(0)|$. It omits $1$ and $\infty$.
Apply the preceding conclusion to $h/h'(0)$ and
the omitted value $1/h'(0)$ to get $r\leq4$.
Both caps therefore have radii at most four.
The tensor is the product of the two derivative
vectors, up to the unit-modulus boundary identification,
so its norm is at most sixteen. The nonzero tensor
$17$ belongs to the stated unrestricted target,
establishing failure of surjectivity.
\end{proof}
The intake SA22 reports the sharper exact image
$0<|q|<16$, tangent invertibility, and a $\pi_1$-isomorphism.
Its full prior proofs are not in the supplied packet.
The next result independently supplies complete proofs
for the labelled geodesic-cut $(0,4)$ sector, while
AQ-67-NONHOLO separates the real and complex tangent claims.
The elementary obstruction (AQ.8) already answers
67/078 negatively and is independent of RH.

\subsection{AQ-67-EXACT: the canonical radius and the exact labelled $(0,4)$ image}
\paragraph{Scope and normalization.}
Consider the geodesic-cut interpretation of the map without marked-point
disks in cote 67/077. Fix a partition of four labelled punctures into two
pairs and remember the separating geodesic. Its length is $\ell>0$.
Cutting produces two hyperbolic pairs of pants of boundary lengths
$(\ell,0,0)$, where zero denotes a cusp. Cap the geodesic boundary by the
unit disk using its normalized hyperbolic arclength, as prescribed in
67/074. On each resulting three-pointed sphere use the unique coordinate
with cap centre $0$ and the two labelled cusps $1,\infty$.
Let $r(\ell)>0$ be the conformal radius of the cap at zero in that
coordinate. A cap coordinate $w$ has $|w|=1$ at the seam and
$z=r(\ell)w+O(w^2)$ after a rotation of $w$.
The tensor coordinate $q$ of 67/075 then has modulus $r(\ell)^2$.

Put
\[
 x=\frac{\ell}{4\pi},\qquad
 \phi(x)=\operatorname{Im}\bigl(\log\Gamma(1+ix)
                       -\log\Gamma(\tfrac12+ix)\bigr),
\]
where the analytic logarithm of $\Gamma$ on the right half-plane is
real on the positive real axis. Then
\[
 \boxed{\quad r(\ell)=4\exp\!\left(\frac{\phi(x)-\pi/4}{x}\right).\quad}
                                                               \tag{AQ.67.1}
\]
The function $r:(0,\infty)\longrightarrow(0,4)$ is a real-analytic
diffeomorphism. With the orientation of the normalized twist angle
$\theta\in\mathbb R/2\pi\mathbb Z$ chosen so that its increase multiplies
the tangent tensor by $e^{i\theta}$, the capping map is
\[
 (\ell,\theta)\longmapsto r(\ell)^2e^{i\theta}.
                                                               \tag{AQ.67.2}
\]
It is a real-analytic diffeomorphism onto $0<|q|<16$.
Its real tangent map is invertible everywhere, and its composite with
the inclusion $\{0<|q|<16\}\hookrightarrow\mathbb C^\times$ induces an
isomorphism of fundamental groups. These are statements about this
labelled geodesic-cut sector. They do not assert complex holomorphicity
or algebraic \emph{\'etaleness} of the map in 67/077--082, nor do they
settle the full multi-node question.

\paragraph{Geometric inputs and source locators.}
The hyperbolic pair of pants with specified nonnegative boundary
lengths is unique, with zero lengths interpreted as cusps. Its
Fenchel--Nielsen length and twist parametrization is recalled in
Mondello (2011), \S2.2, p.~640. The canonical attachment here is the
infinite flat-cylinder grafting of Mondello, \S5.4,
Theorem 5.4, p.~660; a punctured flat cylinder is conformally the
punctured disk with the arclength parametrization used above.
For an explicit independent source for the projective-connection
calculation below, see F\i rat (2021), equations (3.11), (3.13),
(3.14a), (3.23), and (3.30)--(3.34), arXiv:2102.03936v2,
pp.~14--17 and 24--25. We specialize the two other ends to cusps
and determine the branch directly below; we do not infer the branch
from that paper's integer-valued numerical prescription (3.27).

\begin{proof}
Write $\lambda=\ell/(2\pi)$, so $x=\lambda/2$.
The developing map $D$ of the hyperbolic metric is a local isometry
to a half-plane. Near a geodesic of length $2\pi\lambda$, diagonalize
its hyperbolic monodromy, with fixed points $0,\infty$.
In the arclength cap coordinate its analytic continuation is, up
to a nonzero constant, $w^{i\lambda}$. Indeed, increasing
$\arg w$ by $2\pi$ multiplies this function by $e^{-2\pi\lambda}$,
and along $|w|=1$ its hyperbolic arclength is
$\lambda\,d\arg w$. The same expression supplies the analytic
continuation across the seam. Its Schwarzian derivative therefore
has a double pole at the cap centre with leading coefficient
$(1+\lambda^2)/2$. At a cusp its leading coefficient is $1/2$,
as follows from the local developing coordinate $\log(z-z_0)$.
There are no other poles. On the sphere with centres and cusps
$0,1,\infty$, the resulting Schwarzian is consequently
\[
 T(z)=\frac{1+\lambda^2}{2z^2}
       +\frac1{2(z-1)^2}
       -\frac{1+\lambda^2}{2z(z-1)}.                           \tag{AQ.67.3}
\]
To see uniqueness, any two such projective connections differ by
a quadratic differential with at most simple poles at three points;
this space is $H^0(\mathbb P^1,\mathcal O(-1))=0$.

Set $a=(1+i\lambda)/2$. Direct substitution in
$y''+\tfrac12T(z)y=0$ gives the two solutions
\[
 y_\pm(z)=z^{(1\pm i\lambda)/2}(1-z)^{1/2}
     {}_2F_1\!\left(\frac{1\pm i\lambda}{2},
                    \frac{1\pm i\lambda}{2};
                    1\pm i\lambda;z\right).
\]
The ratio, normalized by its leading term at zero, is
\[
 D(z)=\frac{y_+(z)}{y_-(z)}
     =z^{i\lambda}
       \frac{{}_2F_1(a,a;2a;z)}
            {{}_2F_1(\bar a,\bar a;2\bar a;z)}.                \tag{AQ.67.4}
\]
All developing maps with the prescribed Schwarzian differ by a
M\"obius transformation. The ratio in (AQ.67.4) has the same
diagonal monodromy at zero; hence it is a valid developing map
after changing the target half-plane by that transformation.
Its boundary still contains the two fixed points $0,\infty$;
thus it is a straight half-plane bounded by a line through zero,
rather than a disk with a different boundary circle.

We now fix the phase without an unspecified multiple of $\pi$.
On the real interval $0<z<1$, $y_- =\overline{y_+}$.
Their Wronskian, computed at zero, is
$y_+'y_- -y_+y_-'=i\lambda$. In particular neither solution
vanishes there, and the continuously lifted argument $\Theta(z)$
of $D(z)$, normalized by
$\Theta(z)-\lambda\log z\longrightarrow0$ as $z\downarrow0$,
satisfies
\[
 \Theta'(z)=\frac{\lambda}{|y_+(z)|^2}>0.
                                                               \tag{AQ.67.5}
\]
The zero-balanced hypergeometric asymptotic
\[
 {}_2F_1(a,a;2a;z)
   \sim\frac{\Gamma(2a)}{\Gamma(a)^2}[-\log(1-z)]
       \qquad (z\uparrow1)
\]
is DLMF (15.4.21). Thus
\[
 \lim_{z\uparrow1}\Theta(z)=\vartheta(\lambda),\qquad
 \vartheta(\lambda)=2\operatorname{Im}
       \left(\log\Gamma(1+i\lambda)
                    -2\log\Gamma(\tfrac12+i\lambda/2)\right).
                                                               \tag{AQ.67.6}
\]
This is an equality of lifted arguments. For clarity, the
asymptotic is locally uniform in $\lambda\geq0$. The hypergeometric
factor on $0<z<1$ has no zeros when $\lambda>0$ by the Wronskian,
and at $\lambda=0$ it is positive. Its continuous logarithm,
zero at $z=0$, therefore has limiting phase equal to the indicated
continuous gamma phase, rather than that phase plus a nonzero
multiple of $2\pi$. Equivalently that possible integer is locally
constant in $\lambda$ and is zero at $\lambda=0$.

The pair of pants has its reflection fixing the perpendicular
seams. After capping, this antiholomorphic involution fixes
$0,1,\infty$, so it is complex conjugation. The perpendicular
seam from the cusp at $1$ to the geodesic boundary is the terminal
segment of the real interval $(0,1)$. Its developing image lies
on the unit circle, by (AQ.67.4). The cusp endpoint has phase
$\vartheta$. By (AQ.67.5), approaching that cusp from the core
uses phases less than $\vartheta$. Hence the developing half-plane
of the core is the half-plane with argument interval
$(\vartheta-\pi,\vartheta)$. The axis of the boundary monodromy,
whose fixed points are $0,\infty$, is its ray of argument
$\vartheta-\pi/2$. The seam meets that axis perpendicularly.
Starting from the cusp, this is the first intersection of the
seam with the axis; the further formal intersections obtained
after crossing the boundary line of that half-plane are not in
the hyperbolic core. This selects the lifted phase
$\vartheta-\pi/2$, with no integer ambiguity.

Choose $w=1$ at this perpendicular foot. On a collar of the
geodesic, and thence by analytic continuation into its cap,
\[
 w=\exp\left(\frac{\log D-i(\vartheta-\pi/2)}{i\lambda}\right).
\]
Here the logarithm has the lifted phase just specified.
Its boundary modulus is one and its boundary argument has the
normalized arclength, so this is precisely the prescribed cap
coordinate. At zero, (AQ.67.4) gives
\[
 w=\exp\left(-\frac{\vartheta-\pi/2}{\lambda}\right)z+O(z^2),
 \qquad
 r(\ell)=\exp\left(\frac{\vartheta-\pi/2}{\lambda}\right).
                                                               \tag{AQ.67.7}
\]
The duplication formula, DLMF (5.5.5), yields
$\vartheta=2\lambda\log2+2\phi(\lambda/2)$, proving
(AQ.67.1).

It remains to establish the exact range and the nonvanishing
derivative, rather than only an upper bound. The digamma integral
DLMF (5.9.16) gives, for $x\geq0$,
\[
 \phi'(x)=\operatorname{Re}\bigl(\psi(1+ix)
                            -\psi(\tfrac12+ix)\bigr)
           =\int_0^\infty\frac{\cos(xt)}{e^{t/2}+1}\,dt.
                                                               \tag{AQ.67.8}
\]
Let $f(t)=(e^{t/2}+1)^{-1}$. Then
\[
 f''(t)=\frac{e^{t/2}(e^{t/2}-1)}{4(e^{t/2}+1)^3}>0
               \quad(t>0).
\]
The function and its derivatives decay exponentially. Two
integrations by parts, with their endpoint terms included, give
\[
 \phi'(x)=\frac1{x^2}\int_0^\infty f''(t)(1-\cos xt)\,dt>0
                  \quad(x>0).                                \tag{AQ.67.9}
\]
Stirling's expansion for the same analytic logarithms, DLMF
(5.11.1), or the corresponding gamma-ratio expansion (5.11.13),
gives $\phi(x)\longrightarrow\pi/4$ as $x\to\infty$.
Since $\phi(0)=0$, it follows that
$0<\phi(x)<\pi/4$ for $x>0$. Therefore
\[
 \frac{d}{dx}\log r(4\pi x)
       =\frac{x\phi'(x)+\pi/4-\phi(x)}{x^2}>0.                 \tag{AQ.67.10}
\]
At zero, $\phi'(0)=2\log2$, so
\[
 r(\ell)=16\exp\left(-\frac{\pi^2}{\ell}+O(\ell^2)\right)
              \quad(\ell\downarrow0),
 \qquad r(\ell)\longrightarrow4\quad(\ell\to\infty).
\]
Real analyticity and (AQ.67.10) prove that $r$ is the asserted
diffeomorphism.

Finally Fenchel--Nielsen gluing of the two labelled pairs of
pants has exactly one length $\ell>0$ and one twist modulo
$\ell$, equivalently an angle $\theta$ modulo $2\pi$.
The two derivative vectors each have norm $r(\ell)$, and
the twist contributes their tensor's unit phase. This proves
(AQ.67.2). Its inverse is
$\ell=r^{-1}(\sqrt{|q|})$, $\theta=\arg q$ with the latter
viewed as a circle-valued coordinate. In logarithmic polar
coordinates the real Jacobian is diagonal, with entries
$2r'(\ell)/r(\ell)>0$ and $1$. Both the punctured disk and
$\mathbb C^\times$ deformation retract to the same circle
$|q|=1$; the inclusion therefore induces an isomorphism on
$\pi_1$, taking one full twist to the positive generator.
\end{proof}

\subsection{AQ-67-NONHOLO: the natural complex-analytic formulation fails}
Give the labelled geodesic-cut $(0,4)$ space its natural complex
structure, namely the covering
\[
 \mathcal C=\mathcal T(S_{0,4})/\langle D_\alpha\rangle
\]
of the moduli space of four labelled punctures, where $\alpha$ is
the chosen separating curve and $D_\alpha$ is its Dehn twist.
Equivalently, this is the component of the cover that remembers
the isotopy class of the cut. The two sides and all punctures
remain labelled. Fenchel--Nielsen coordinates upstairs are
$(\ell,\tau)\in\mathbb R_{>0}\times\mathbb R$, with
$D_\alpha(\ell,\tau)=(\ell,\tau+\ell)$; downstairs the
normalized angle is $\theta=2\pi\tau/\ell$ modulo $2\pi$.
The canonical capping map of AQ-67-EXACT is not holomorphic
from this complex manifold to $\mathbb C^\times$.
Thus its everywhere invertible real tangent map does not make
it a complex-analytic or $\mathbb C$-algebraic \emph{\'etale}
map. This refutes that stronger interpretation of 67/077--082
for the natural Teichm\"uller complex structure already in
the labelled $(0,4)$ sector.

\begin{proof}
Suppose it were holomorphic. By AQ-67-EXACT it is a bijection
with nonzero real Jacobian onto the punctured disk
$\mathcal D=\{0<|q|<16\}$. The holomorphic inverse function
theorem then makes it a biholomorphism with $\mathcal D$.
Conjugating the rotations $q\mapsto e^{is}q$ by it gives
holomorphic automorphisms of $\mathcal C$. Their lifts to
$\mathcal T(S_{0,4})$, continuously chosen to equal the identity
at $s=0$, are the explicit global diffeomorphisms
\[
 H_s(\ell,\tau)=
       \left(\ell,\tau+\frac{s\ell}{2\pi}\right),
                      \qquad s\in\mathbb R.                  \tag{AQ.67.11}
\]
This follows directly from the tensor's unit twist phase;
any choice of zero-angle section changes $\theta$ by a function
of $\ell$ and leaves (AQ.67.11) unchanged. Since the quotient
projection is a holomorphic covering, every $H_s$ is
holomorphic as well.

Let $\omega_{\mathrm{WP}}$ and $g_{\mathrm{WP}}$ be the
Weil--Petersson K\"ahler form and metric. Wolpert (1983),
Lemma 4.1, p.~229, says that the arclength twist vector field
$t_\alpha=\partial/\partial\tau$ has Hamiltonian potential
$-\ell$, in his normalization:
\[
 \iota_{t_\alpha}\omega_{\mathrm{WP}}=-d\ell.
\]
The vector field of (AQ.67.11) is
$V=(\ell/2\pi)t_\alpha$, so
\[
 \iota_V\omega_{\mathrm{WP}}
   =-\frac{\ell}{2\pi}d\ell
   =-d\!\left(\frac{\ell^2}{4\pi}\right).
\]
Cartan's formula and $d\omega_{\mathrm{WP}}=0$ imply
$\mathcal L_V\omega_{\mathrm{WP}}=0$; hence every global
$H_s$ preserves $\omega_{\mathrm{WP}}$. A different fixed
normalization of the Weil--Petersson form only multiplies
these identities by a nonzero constant and gives the same
conclusion. A holomorphic map preserving the K\"ahler form
preserves the metric, since that metric is recovered from
the form and the complex structure. Thus every $H_s$ would
be a Weil--Petersson isometry of $\mathcal T(S_{0,4})$.

Brock and Margalit, \emph{Weil--Petersson isometries via the
pants complex}, Main Theorem, p.~2 of arXiv:math/0412499v1,
prove that every such isometry is induced by an element of
the extended mapping class group; their theorem explicitly
includes $S_{0,4}$. That group is countable. On the other
hand (AQ.67.11) gives distinct maps for every two distinct
real numbers $s$, because $\ell>0$ and the upstairs twist
coordinate is real, not taken modulo $\ell$. This is an
uncountable family of isometries, a contradiction.
\end{proof}

\subsection{AQ-67-INVOLUTION: the proposition holds; its corollary fails}
The proposition of 67/091 gives an equivalence between
two-regular modules with involution and triples
$(P,Q,m)$, where $P,Q$ are two-regular and
$m\subset P/2P\oplus Q/2Q$ meets each summand trivially.
This proposition is correct, with morphisms of triples
preserving $m$.
\begin{proof}
For an involution on $M$, let $P=M_+$ and $Q=M_-$.
Two-regularity gives $P\cap Q=0$ and
$2M\subset P\oplus Q\subset M$.
Let $m=2M/(2P\oplus2Q)$. If an element of $2M$
has zero $Q$-coordinate, say $2x\in P$, then
$2(\sigma x-x)=0$, so $x\in P$ and $2x\in2P$.
This proves $m\cap(P/2P)=0$, and the other summand
is identical.

Conversely take the inverse image
$L\subset P\oplus Q$ of $m$, with involution
$(p,q)\mapsto(p,-q)$. It is stable because modulo
two this involution is the identity. It is
two-regular, and the intersection conditions give
$L_+=2P\oplus0$, $L_-=0\oplus2Q$.
Multiplication by two canonically identifies these
eigenspaces with $P,Q$ and recovers the original
$m$ from $2L/(4P\oplus4Q)$.
Starting with $M$, the reconstructed $L$ is $2M$,
canonically isomorphic to $M$ by multiplication
by two. Maps restrict to the eigenspaces, and
maps of triples preserve the indicated inverse
images. These constructions are inverse on maps
as well as objects, proving the equivalence.
\end{proof}

The corollary claiming invertible rank-one eigenspaces
for a nontrivial involution on a rank-two projective
module over a commutative two-regular ring is false
even when the ring is an integral domain. Set
\[
 R=\mathbb Z[x,y,z]/(x^2+yz-1),\quad M=R^2,\qquad
 \sigma=\begin{pmatrix}x&y\\z&-x\end{pmatrix}.                   \tag{AQ.9}
\]
\begin{proof}
The polynomial $f=x^2+yz-1$, viewed in
$\mathbb Z[x,y][z]$, is primitive and linear over
the fraction field; hence it is irreducible in
the unique factorization domain $\mathbb Z[x,y,z]$.
Thus $R$ is a domain, in particular two-regular.
Clearly $\sigma^2=1$ and $\sigma\ne\pm1$.

In $S=\mathbb Z[x,y,z]$ set $A=\sigma-I$ and
$B=\sigma+I$. Then $AB=BA=fI$.
If $\bar A\bar v=0$ over $R$, lift $v$ to $S^2$
and write $Av=fw$. Multiplication by $B$ gives
$fv=fBw$; since $f$ is a non-zero-divisor,
$v=Bw$. Therefore $\ker\bar A=\operatorname{im}\bar B$.
Interchanging $A,B$ proves
$\ker\bar B=\operatorname{im}\bar A$, and hence
\[
 P=\ker(\sigma-I)\simeq\operatorname{coker}(\sigma-I).
\]
At the maximal ideal
$\mathfrak m=(2,x-1,y,z)$ every entry of $A$
vanishes, so
$P/\mathfrak mP\simeq\mathbb F_2^2$.
Generically $A$ has rank one, giving $P$ rank one.
An invertible module would have a one-dimensional
fibre at $\mathfrak m$, a contradiction.
The same argument applies to $Q$.
\end{proof}
Connectedness of $\operatorname{Spec}R$ alone cannot
repair the corollary. This algebraic counterexample
does not involve RH, and cannot be excluded by it.

\subsection{AQ-154-COUNT: what can actually be counted over finite fields}
The supplied 154/120 says
\[
 \chi(\mathcal X)=s_0-s_1+s_2=1-g,\qquad
 2s_1=4s_0,\qquad \chi(\mathcal X)=s_2-s_0.
\]
Its $\mathcal X_2$ is a cellular surface of relative
positions of pseudolines. Neither that sheet nor SA11
supplies it with a scheme over a finite field.
The original Euler-characteristic calculation is
therefore not an assertion about
$\#\mathcal X_2(\mathbb F_q)$.

There is a genuine arithmetic consequence for a
different, explicitly specified object.
Let $L_1,\dots,L_n$ be distinct rational projective
lines in $\mathbb P^2_{\mathbb Q}$, let $U$ be their
complement, and let $S$ contain the finitely many
primes where lines coincide or their incidence
lattice changes. If $v$ runs through their
intersection points and $m_v$ lines meet at $v$, set
\[
 b_2=\sum_v(m_v-1)-n+1.
\]
For every $p\notin S$ and $r\geq1$,
\[
 \#U(\mathbb F_{p^r})=p^{2r}-(n-1)p^r+b_2,\qquad
 Z(U_{\mathbb F_p},t)
  =\frac{(1-pt)^{n-1}}{(1-p^2t)(1-t)^{b_2}}.                  \tag{AQ.10}
\]
Consequently its incomplete arithmetic zeta function is
\[
 \zeta^S(U,s)=
 \frac{\zeta^S(s-2)\zeta^S(s)^{b_2}}
      {\zeta^S(s-1)^{n-1}},\qquad
 \zeta^S(s)=\prod_{p\notin S}(1-p^{-s})^{-1}.                 \tag{AQ.11}
\]
\begin{proof}
Every intersection point is rational because the
lines are rational. Over $\mathbb F_{p^r}$ the sum
of the $n$ line counts is $n(p^r+1)$.
It overcounts an intersection $v$ exactly $m_v-1$
times. Subtracting the corrected union count from
$p^{2r}+p^r+1$ proves the first formula.
Insert it in
$Z(U_{\mathbb F_p},t)=
\exp(\sum_{r\geq1}\#U(\mathbb F_{p^r})t^r/r)$
and use $-\log(1-z)=\sum_{r\geq1}z^r/r$ to get
the second. Multiplying its Euler factors gives
(AQ.11), initially for $\operatorname{Re}s>3$
and then by meromorphic continuation.
\end{proof}
Under the adopted RH, the nontrivial-zero contribution
to the divisor of (AQ.11) is on
$\operatorname{Re}s=\tfrac12,\tfrac32,\tfrac52$:
the middle line gives poles when $n>1$, while the
other two give zeros when their exponents are
positive. Trivial zeros, the zeta pole, and the
explicit omitted Euler factors remain separately
specified by (AQ.11). These shifted divisors do not
give a formula for the pseudoline counts or the
topology of $\mathcal X_2$.

\subsection{AQ-136: recovered admissibility, strict failure, and coefficient-local comparisons}

\paragraph{Original pages and conventions.}
The official cote 136 PDF has alternating mathematical sheets and
unrelated printed reverses.  The relevant physical PDF pages are
111, 113, 115, 117, 119 and 121, with archive footers respectively
110, 112, 114, 116, 118 and 120.  Thus the supplied intake's
``136/117--121'' must not be read as six consecutive mathematical
pages or as the archive-footer numbering.
Physical 111 introduces contravariant functors from simplicial
sets to nonnegatively bigraded commutative differential
$S_k$-algebras with divided powers, and identifies the functors
carrying colimits to limits with simplicial such algebras.
Physical 113--115 requires comparison with constant-coefficient
cohomology after the weight-tail operation $\tau_\infty$.
The coefficient map is $S_k\longrightarrow C_n^{**}$; the sheet
calls this the augmentation.  No map $C_n^{**}\to S_k$ is assumed.

The explicit tests on physical 117 are
\begin{equation}
 \begin{aligned}
 S_k^i&\longrightarrow H^{i,0}(C_n^{**})
  &&\text{is an isomorphism for sufficiently large }i,\\
 H^{i,p}(C_n^{**})&=0
  &&(p>0,\ i\geq\lambda(p)).
 \end{aligned}                                           \tag{136.3}
\end{equation}
The final lines distinguish an initially simplex-dependent
$\lambda_n$ from the desired bound independent of $n$.
On physical 119, let $\pi(i)$ be the smallest $p$ for which the
simplicial module $C_*^{i,p}$ is not flasque, and $+\infty$ if
there is none.  The sheet explains flasqueness by vanishing
simplicial homology, and asks that
\[
                         \pi(i)\longrightarrow+\infty.
                                                               \tag{136.4}
\]
Physical 121 asks for a common admissible source
\[
 C''\longrightarrow C,\qquad C''\longrightarrow C',              \tag{136.a}
\]
and, for parallel morphisms $f,g:C\rightrightarrows C'$, an
admissible $h:C''\to C$ \emph{``qui `\'egalise' ''}, namely
$fh=gh$.  The same sheet also asks whether the groupoid obtained
by inverting all morphisms is connected and simply connected.
The strict equalizer condition is stronger than that groupoid
assertion.  The proofs below exhibit the difference concretely.

\paragraph{Precise category used in the proofs.}
Write $S=\Gamma_k(T)=\bigoplus_{w\geq0}kT^{[w]}$, with
$T^{[a]}T^{[b]}=\binom{a+b}{a}T^{[a+b]}$.
The coefficient weight is $w$, the cochain degree is $p$, and
$d$ has bidegree $(0,1)$.  Let $\mathcal A_k$ consist of simplicial,
nonnegatively bigraded, strictly commutative PD differential
$S$-algebras $C_n^{w,p}$, with morphisms preserving $S$, both
gradings, the differential and divided powers.  Their PD ideal
contains all positive-weight even elements; odd elements square
to zero.  Require (136.3) with bounds uniform in $n$, and
(136.4).  Here flasque means that the normalized simplicial
chain complex is acyclic, including in degree zero.
The associated functor on arbitrary simplicial sets is obtained
by taking compatible sections over all simplices, separately in
each bidegree.  This construction preserves the algebra and PD
operations and is the colimit-to-limit extension on physical 111.
These conventions specify the category rather than presuming
unretrieved prior-run definitions.  The strict counterexample
below needs only the eventual degree-zero case of (136.4), so
it also applies to a larger interpretation with weaker bounds
in (136.3).  The rational assertions concern the explicit
category $\mathcal A_k$ just defined.

\paragraph{AQ-136-COMPARISON: the explicit tests give the cohomological comparison.}
For completeness, the numerical tests used here do imply the
weight-tail comparison on physical 113--115.  More precisely,
for every fixed $d$ there is $W_d$ such that, for every simplicial
set $X$ and $w\geq W_d$, the two natural comparison maps give
\[
 H^j(C^{w,*}(X))\simeq H^j(X;k)\otimes_k S^w
                  \qquad(0\leq j\leq d).                  \tag{136.2}
\]
This assertion uses the uniform-in-simplex form of (136.3).

\begin{proof}
Let $\mathcal P$ be the abelian category of simplicial
$k$-modules.  Sections over $X$ are
$\Gamma_X(F)=\operatorname{Hom}_{\mathcal P}(k[X],F)$.
The normalized chain complex $K=N(k[X])$ has free $k$-modules,
with basis the nondegenerate simplices.  We first check that
if $F$ has acyclic normalization, then it is acyclic for the
right-derived functors of $\Gamma_X$.

Under Dold--Kan this is the assertion
$\operatorname{Ext}^r_{\operatorname{Ch}_{\geq0}(k)}
 (K,N(F))=0$ for $r>0$.
Choose a resolution of $K$ by projective chain complexes,
using sums of projective disks and degree-zero projectives.
Every short exact sequence in this resolution is split in
each chain degree: the first quotient has projective terms,
and its kernel then has projective terms, inductively.
For any nonnegative chain complex $J$ of projective modules,
the cochain complex $\operatorname{Hom}^{\bullet}(J,N(F))$
is acyclic.  Indeed a closed homogeneous map is
null-homotopic by induction on the source degree: subtract
the already chosen homotopy on the boundary, obtain a map
into cycles, and lift it to the next term of the acyclic
target using projectivity of that source term.  This begins
in degree zero and works for every homogeneous degree.

Now a chain map from a kernel in the projective resolution
to $N(F)$ has a graded extension to the preceding projective
complex.  Its defect of commuting with the differential
factors through the quotient and is a closed degree-one
Hom element.  The just-proved Hom acyclicity corrects this
defect.  Thus every such chain map extends as a chain map;
applying Hom to the resolution is exact in positive degrees.
This proves the required derived-section acyclicity.

The same projective-disk resolution gives
$R^j\Gamma_X(\underline{k})=H^j(X;k)$.
One can verify this without identifying a topological site:
on complexes with projective terms, the cohomology of
$\operatorname{Hom}^{\bullet}(K,k[0])$ has degree-zero
functor $\operatorname{Hom}_{\rm Ch}(K,k[0])$ and vanishes
in positive degrees on projective disks and degree-zero
projectives.  Degreewise split exact sequences give its
long exact sequence, so the projective resolution computes
exactly these groups.  The resulting cochain complex is
the usual normalized simplicial cochain complex of $X$.

Choose $w$ so large that (136.3) holds through cochain
degree $d+1$, and that all rows $C_*^{w,p}$ for
$0\leq p\leq d+1$ are acyclic as simplicial modules.
There is one such bound for every $d$, by (136.3)--(136.4).
The first-quadrant hypercohomology spectral sequence for
$R\Gamma_X(C^{w,*})$, applied first to vertical cohomology,
identifies its groups through degree $d$ with
$R^j\Gamma_X(\underline{S^w})$.
Applied instead to the derived sections of each term,
it identifies those same groups with
$H^j(\Gamma_X(C^{w,*}))$, by the acyclicity just proved.
The coefficient augmentation and the map from ordinary
to derived sections are the comparison maps in these two
identifications.  Since $S^w$ is a rank-one free
$k$-module, they give (136.2).
Both spectral sequences converge in these degrees:
all indices are nonnegative and each total degree uses
only finitely many bidegrees.  This proves the claim for
arbitrary $X$, without interchanging a weight-tail limit
with an infinite product.
\end{proof}

\paragraph{Simplicial algebra fact used below.}
Let $L_n=k^{n+1}$ have basis $x_0,\ldots,x_n$.  For
$\theta:[m]\to[n]$, put
\[
 \theta^*x_j=\sum_{\theta(a)=j}x_a.
\]
Then $t_n=\sum_jx_j$ defines an inclusion of the constant
simplicial module $k$ in $L$, and put $Q=L/kt$.
The normalized chain complex of $L$ is the two-term disk
$k\xrightarrow{\pm1}k$ in degrees one and zero: in the
normalization using $\bigcap_{i<n}\ker d_i$, its degree-one
term is $kx_0$, and in degrees at least two the simultaneous
kernels of $d_0$ and $d_1$ are zero.  Thus $L$ is simplicially
homotopy equivalent to zero.  An exact known input for the last
step is the chain-to-simplicial homotopy theorem, Stacks Project,
Lemma 14.29.3, Tag \nolinkurl{01A4}.

For every $r>0$, $\Gamma^r(L)$ and $\operatorname{Sym}^r(L)$
are consequently simplicially contractible: apply the functor
to a homotopy between the identity and zero.  Functorial
preservation of simplicial homotopies is Stacks Project,
Remark 14.26.4, Tag \nolinkurl{019M}; these positive-degree
functors send the zero map to zero.  If $V$ is a simplicially
contractible module and $W$ any simplicial module, then
$V\otimes_k W$ is also contractible.  Indeed the homotopy
maps $h_{n,\alpha}:V_n\to V_n$, indexed by
$\alpha:[n]\to[1]$, give homotopy maps
$h_{n,\alpha}\otimes1$ on $V_n\otimes W_n$, satisfying the
same face and degeneracy identities.  No flatness or
characteristic-zero hypothesis is used in this observation.

\paragraph{AQ-136-DP-MODEL: an explicit admissible integral theory.}
For every commutative ring $k$, define
\[
 B_n=\Gamma_k(x_0,\ldots,x_n)
       \otimes_k\Lambda_k(e_0,\ldots,e_n)/(e_0+\cdots+e_n),
 \quad dx_j=e_j,\quad de_j=0,\quad T\mapsto\sum_jx_j .        \tag{136.5}
\]
Each $x_j$ has bidegree $(1,0)$ and $e_j$ has bidegree
$(1,1)$.  Simplicial operators use the displayed summation
formula on both $x_j$ and $e_j$.  Divided powers of the
coefficient $T$ are sent to divided powers of its sum.
Then $B\in\mathcal A_k$; (136.3) holds in every weight,
and its row $B^{w,p}_*$ is contractible whenever $w>p$.

\begin{proof}
First, all the displayed algebraic structures exist integrally.
The ordinary divided-power polynomial algebra has basis
$\prod_jx_j^{[a_j]}$ and the multiplication above; its precise
universal property is Stacks Project, Lemma 23.5.1,
Tag \nolinkurl{07H5}.  Tensoring it with an exterior algebra
gives the usual strictly graded-commutative PD algebra.
One direct verification is to work first over $\mathbb Z$.
It is a free abelian group inside its rationalization, where
$\gamma_r(z)=z^r/r!$ for positive-weight even $z$.
For an even monomial containing an exterior generator, its
square is zero, so its divided powers of order at least two
vanish.  For a positive polynomial divided-power monomial,
$r!$ divides the factorial coefficient in its $r$th power:
choose one variable occurring positively and count partitions
of its $ra$ occurrences into $r$ unordered blocks of size $a$.
The addition formula for divided powers therefore preserves
the integral lattice.  The identities hold in the
rationalization and hence over $\mathbb Z$; the resulting
integer polynomial identities give the construction over every
$k$, including rings with torsion.  Eliminating $e_0$ gives
an exterior algebra on $e_1,\ldots,e_n$, so the displayed
relation is compatible with this structure.

The differential is explicitly
$d(x_j^{[a]})=x_j^{[a-1]}e_j$ for $a>0$ and extends by the
signed Leibniz rule.  Its square is zero.  It satisfies
$d\gamma_r(z)=\gamma_{r-1}(z)dz$: this follows over
$\mathbb Q$ by differentiation and integrally by the same
free-lattice argument.  The simplicial maps preserve it.
The relation $\sum e_j=0$ says precisely that $dT=0$.
All simplicial identities follow from composition of the
sums over the fibres of order-preserving maps.

At a fixed simplex make the invertible linear change of
variables $(x_0,\ldots,x_n)\mapsto(T,x_1,\ldots,x_n)$.
The PD universal property gives
\[
 B_n\simeq S\otimes_k\bigotimes_{j=1}^n
            \bigl(\Gamma_k(x_j)\otimes\Lambda_k(e_j)\bigr).
                                                               \tag{136.6}
\]
For one factor a contraction onto $k$ is
$h(x^{[a]}e)=x^{[a+1]}$ and $h(x^{[a]})=0$.
The identity $dh+hd=1-\epsilon$ follows on each basis
monomial; $\epsilon$ keeps only the constant term.
The tensor contraction is obtained successively by
$h\otimes1+\epsilon\otimes h$, with the usual sign in the
tensor action of a degree-minus-one map.  Thus the inclusion
$S\to B_n$ is a cochain homotopy equivalence.  The contractions
preserve weight, so $H^{w,0}(B_n)=S^w$ and all
$H^{w,p}(B_n)$ for $p>0$ vanish, uniformly in $n$.

Finally there is a natural identification of simplicial modules
\[
 B_*^{w,p}=\Gamma^{w-p}(L)\otimes_k\Lambda^p(Q),             \tag{136.7}
\]
with zero understood for $w<p$.  If $w>p$, the first factor
is contractible by the preceding simplicial fact, and so is
the tensor product.  Hence (136.4) holds with
$\pi(w)\geq w$.  This proves every assertion.
\end{proof}

\paragraph{AQ-136-STRICT: the original strict equalizing-source conjecture is false.}
If $k\ne0$, there are two parallel morphisms in
$\mathcal A_k$ which admit no equalizing morphism from any
admissible source.  In particular this refutes the strict
question on physical 121 over $\mathbb Z$, without a
rationalization or a proposed lifting lemma.
The supplied intake had reported a strict-form refutation
under the earlier identifier C07; the following is a complete
independent proof from the recovered original conditions,
not a claim to have recovered that absent proof.

\begin{proof}
Use the theory $B$ of (136.5).  Define a differential
$B$-module $M$ by adjoining symbols $\alpha_j$ of bidegree
$(1,0)$ and $\beta_j$ of bidegree $(1,1)$, imposing
$\sum\alpha_j=\sum\beta_j=0$, and putting
$d\alpha_j=\beta_j$, $d\beta_j=0$.
In invariant form
\[
 M=B\otimes_k\bigl(Q\{0\}\xrightarrow{1}Q\{1\}\bigr),
                                                               \tag{136.8}
\]
where the two copies have coefficient weight one and the
braces indicate their cochain degrees.  There is a
weight-preserving degree-zero PD derivation
\[
 \delta:B\longrightarrow M,\qquad
 \delta x_j=\alpha_j,\quad\delta e_j=\beta_j,\quad
 \delta(x_j^{[a]})=x_j^{[a-1]}\alpha_j .                    \tag{136.9}
\]
Extend it by the degree-zero Leibniz rule.  The two sum
relations make $\delta T=0$ and $\delta\sum e_j=0$.
The divided-power addition formula shows that it annihilates
all of $S$ and respects the PD identities.  It commutes
with $d$, as is checked on $x_j,e_j$ and their divided
powers and then by the Leibniz rule.  Everything is natural
in the simplex.

Form the square-zero extension $D=B\ltimes M$, with
$M^2=0$, its induced differential, and divided powers
\[
 \gamma_r(b+m)=\gamma_r(b)+\gamma_{r-1}(b)m
 \quad(r\geq1)                                             \tag{136.10}
\]
for even elements in the PD ideal.  Give $M$ zero higher
divided powers.  Formula (136.10), the addition formula,
and $M^2=0$ verify all the PD axioms; alternatively these
are the terms of degree zero and one in an extra
square-zero variable in the integral construction just
proved.  Differential compatibility follows because
\[
 d(\gamma_{r-1}(b)m)
  =\gamma_{r-2}(b)(db)m+\gamma_{r-1}(b)dm,
\]
with the first term absent for $r=1$, agreeing with
$\gamma_{r-1}(b+m)d(b+m)$ after using $M^2=0$.
Thus
\[
 j_0,j_1:B\rightrightarrows D,
 \qquad j_0(b)=(b,0),\quad j_1(b)=(b,\delta b)              \tag{136.11}
\]
are strict morphisms of differential PD $S$-algebras.

We verify that $D$ is admissible.  The summand $M$ is
vertically contractible: on homogeneous $b\in B$ put
$h(b\beta_j)=(-1)^{\deg b}b\alpha_j$ and
$h(b\alpha_j)=0$.  The relations are preserved, and
$dh+hd=1$ follows by cancelling the two terms involving
$db$.  Consequently $S\to D_n$ has exactly the same
cohomology as $S\to B_n$ in every weight and degree.
Horizontally,
\[
 M_*^{w,p}=
  (B_*^{w-1,p}\otimes Q)\ \oplus\
  (B_*^{w-1,p-1}\otimes Q).                                \tag{136.12}
\]
The second term is zero if $p=0$.  Both terms are
contractible for $w>p+1$, by (136.7) and the tensor
homotopy fact.  Hence $\pi_D(w)\geq\max(w-1,0)$.
No coefficient tail has been discarded: both $B$ and
$D$ retain every $T^{[w]}$ over the original ring $k$.

At every simplex the equalizer of $j_0,j_1$ is exactly
$S$.  To check this, eliminate the zeroth variable as
in (136.6); $M_n$ is free over $B_n$ on
$\alpha_1,\ldots,\alpha_n,\beta_1,\ldots,\beta_n$.
Expand an element of $B_n$ uniquely as
\[
 \sum_{a,I}s_{a,I}x_1^{[a_1]}\cdots x_n^{[a_n]}e_I,
 \qquad s_{a,I}\in S.
\]
The coefficient of $\beta_j$ in its $\delta$ is the
signed exterior partial derivative with respect to $e_j$.
The basis monomials containing $e_j$ map injectively to
distinct monomials after removing it.  Thus $\delta=0$
forces every nonempty $I$ coefficient to vanish.
Then the coefficient of $\alpha_j$ is the divided-power
partial derivative, sending $x_j^{[a_j]}$ to
$x_j^{[a_j-1]}$ with coefficient one.  It forces all
$a_j>0$ coefficients to vanish.  This argument works
over arbitrary $k$; it does not divide by an integer.
The remaining elements are exactly $S$.

Suppose an admissible $E$ and an $S$-morphism $u:E\to B$
equalized (136.11).  By the computed equalizer, $u$
factors through the constant simplicial algebra $S$.
Since it is an $S$-morphism, this factorization retracts
the structure map $S\to E$.  In each weight $w$ it makes
the nonzero constant simplicial module
$kT^{[w]}$ a direct summand of $E_*^{w,0}$.
Therefore $H_0^{\rm simp}(E_*^{w,0})$ contains $k$ for
every $w$.  This contradicts (136.4), which makes that
row acyclic for all sufficiently large $w$.
The contradiction proves the strict assertion false.
\end{proof}

\paragraph{AQ-136-RATIONAL: common sources and terminal localization over a $\mathbb Q$-algebra.}
If $k$ is a $\mathbb Q$-algebra, every finite family of
objects of $\mathcal A_k$ has a common source in
$\mathcal A_k$.  The ordinary category obtained from
$\mathcal A_k$ by inverting every morphism is equivalent
to the terminal groupoid.  This is a statement about
that groupoid, not about all higher homotopy groups of
an infinity-categorical localization.

\begin{proof}
Identify $S=\Gamma_k(T)$ with $k[t]$ by
$T^{[w]}=t^w/w!$.  For an integer $N\geq1$, define
\[
 B_{N,n}=
 \frac{k[t,x_0,\ldots,x_n]\otimes\Lambda(e_0,\ldots,e_n)}
      {(\sum x_j-t^N,\ \sum e_j)},
 \quad dx_j=e_j,\quad dt=de_j=0.                           \tag{136.13}
\]
Put $\operatorname{wt}(t)=1$ and
$\operatorname{wt}(x_j)=\operatorname{wt}(e_j)=N$.
Use the same simplex operators as before, fixing $t$.
Over a $\mathbb Q$-algebra the PD operations on even
positive-weight elements are uniquely $z^r/r!$.
Thus (136.13) is a differential PD $S$-algebra.

Eliminating $x_0,e_0$ gives the polynomial de Rham
complex over $S$ in $n$ variables of weight $N$.
In one variable the homotopy
$h(x^ae)=x^{a+1}/(a+1)$ contracts onto the constants.
Its tensor product contraction proves (136.3) in every
weight, again preserving weight.  For the horizontal
calculation let $A$ be the same construction with
coefficient $U=\sum x_j$ and all $x_j,e_j,U$ of weight
one.  As a simplicial module
$A_*^{m,p}=\operatorname{Sym}^{m-p}(L)\otimes\Lambda^p(Q)$,
so it is contractible if $m>p$.
Before extension of coefficients regrade these weights
by $N$.  We have
\[
 B_N=k[t]\otimes_{k[U]}A,\qquad U\longmapsto t^N,
 \qquad k[t]=\bigoplus_{r=0}^{N-1}t^r k[U].                \tag{136.14}
\]
The last equality is an actual free-module decomposition,
valid over every $k$.  The component of weight $w$
therefore identifies with $A^{m,p}$, where
$w=Nm+r$ and $0\leq r<N$.  It is contractible for
$m>p$.  In particular
$\pi_{B_N}(w)\geq\lfloor w/N\rfloor$, which tends to
infinity.  This proves $B_N\in\mathcal A_k$.

Now fix finitely many admissible theories $C_a$.
Choose $N\geq1$ large enough that every simplicial
module $C_a^{N,0}$ is acyclic, using (136.4).
The coefficient $t^N$ determines a map from the
constant simplicial module $k$ to $C_a^{N,0}$.
It extends along the inclusion $k\to L$ to a
simplicial linear map
\[
                         L\longrightarrow C_a^{N,0}.       \tag{136.15}
\]
Here is the precise lifting reason.  Normalization
sends $k\to L$ to the inclusion of the degree-zero
term into the two-term disk.  Vanishing of $H_0$
says that the coefficient element $t^N$ at simplex
zero is the boundary of an element in normalized
degree one.  Choosing that element gives a chain
map from the disk.  The Dold--Kan equivalence,
Stacks Project Theorem 14.24.3, Tag \nolinkurl{019G},
lifts it to (136.15), retaining the given constant
map.  No lift of an arbitrary cochain complex
through a surjection is being asserted.

Send $x_j$ to the image of its basis vector in
(136.15), and $e_j$ to the vertical differential
of that image.  Their sums are respectively
$t^N$ and zero.  The polynomial/exterior universal
property defines a simplicial differential algebra
map $B_N\to C_a$, and it preserves divided powers
because $k$ contains $\mathbb Q$.  It preserves
weights and $S$ by construction.  These maps give
the claimed simultaneous common source.

It remains to prove terminality after inversion;
strict equalizers cannot be used, by the preceding
counterexample.  For any admissible $D$, let
\[
 P(D)_n=D_n\otimes_k k[s,ds],\quad
 \operatorname{wt}(s)=\operatorname{wt}(ds)=0,
 \quad d s=ds,\quad d(ds)=0.                              \tag{136.16}
\]
It has cochain-degree-zero $s$ and degree-one $ds$.
This is an object of $\mathcal A_k$.  Integration of
polynomials contracts $k[s,ds]$ onto $k$, naturally
in $D$, and so leaves every vertical cohomology
unchanged.  Its row of bidegree $(w,p)$ is the sum
of copies of $D^{w,p}$ and $D^{w,p-1}$ indexed by
powers of $s$.  Normalization and homology commute
with these direct sums; hence (136.4) is retained.
All PD operations exist canonically over $\mathbb Q$.
The maps $j:D\to P(D)$ and
$e_0,e_1:P(D)\to D$ are admissible morphisms and
$e_0j=e_1j=1$.  Consequently $e_0=e_1=j^{-1}$ in
the all-arrow localization.

For any parallel $f,g:C\rightrightarrows D$, choose
a common-source map $a:B_N\to C$ as above.
There is an explicit PD differential map
$H:B_N\to P(D)$ given on each generator by
\[
 H(x_j)=(1-s)fa(x_j)+sga(x_j),\qquad H(e_j)=dH(x_j),
 \qquad H(t)=t.                                           \tag{136.17}
\]
Its sums are $t^N$ and zero, it is natural in the
simplex and preserves weights.  Thus the relations
and all structures are respected.  Its two
evaluations are $fa$ and $ga$.  The equality of
evaluations in the localization, followed by
inversion of $a$, gives $f=g$ there.

Finally any finite zigzag in the localization
involves finitely many objects $C_0,\ldots,C_r$.
Choose a common source $B_M$ and arrows
$a_i:B_M\to C_i$.  For every original arrow
$f:C_i\to C_j$ in the zigzag, the two parallel
arrows $fa_i$ and $a_j$ have the same image in
the localization, by the preceding paragraph.
Hence the image of $f$ is $a_j a_i^{-1}$, and
its inverse is $a_i a_j^{-1}$.  The whole zigzag
telescopes to $a_r a_0^{-1}$.  Applying the same
argument to the union of any two finite zigzags
shows that they agree.  Existence of common sources
makes the groupoid connected; $B_1$ makes it
nonempty.  It therefore has exactly one morphism
between each two objects, proving the assertion.
\end{proof}

\paragraph{AQ-136-INTEGRAL-TAIL: the direct one-generator integral analogue fails.}
The strict counterexample retained every coefficient
$T^{[w]}$ and works over $\mathbb Z$.  There is also a
precise obstruction to extending the rational common-source
construction by replacing $t^N$ with $T^{[N]}$.
For $N\geq2$, form the analogous divided-power theory
$B_N^{\rm PD}$ over $S=\Gamma_{\mathbb Z}(T)$ by
\[
 B_{N,n}^{\rm PD}
  =S\langle x_0,\ldots,x_n\rangle
       \otimes\Lambda(e_0,\ldots,e_n)
          /(\sum x_j-T^{[N]},\ \sum e_j),                   \tag{136.18}
\]
where the first relation means the PD quotient, each
$x_j,e_j$ has weight $N$, and $dx_j=e_j$.
Eliminating the zeroth variables gives a free PD
polynomial de Rham complex over $S$, so vertical
acyclicity is retained.  Nevertheless (136.4) fails:
its degree-zero simplicial homology is nonzero in
arbitrarily large weights.

\begin{proof}
In simplex zero the zero-form algebra is $S$; in
simplex one it is $S\langle x_1\rangle$, with the two
faces evaluating $x_1$ at $T^{[N]}$ and at zero.
Therefore its simplicial degree-zero homology is
\[
 H_0^{\rm simp}((B_N^{\rm PD})^{*,0})=S/J_N,
 \qquad J_N=(\gamma_m(T^{[N]}):m\geq1).                    \tag{136.19}
\]
This is a direct calculation of the cokernel of
$d_0-d_1$, using the $S$-basis $x_1^{[m]}$; it does
not presume flat base change of a simplicial
resolution.  The ideal $J_N$ is the full PD ideal
generated by the chosen coefficient, not merely
its ordinary principal ideal: the iterated-PD and
addition identities show that the displayed
ordinary ideal is PD-stable and is the smallest
such ideal containing $T^{[N]}$.

Let $p>N$ be prime.  Only $mN\leq p$ contributes
in weight $p$.  The corresponding product is
\[
 T^{[p-mN]}\gamma_m(T^{[N]})
   =\frac{p!}{m!(N!)^m(p-mN)!}\,T^{[p]}.                   \tag{136.20}
\]
The coefficient is an integer by the PD identities.
Its numerator has exactly one factor $p$, and
none of the factorials in its denominator has
one: $N<p$, $m\leq p/N<p$, and $p-mN<p$.
Thus it is divisible by $p$.
Every weight-$p$ coefficient of $J_N$ is consequently
a multiple of $p$, so $(S/J_N)^p$ surjects onto
$\mathbb Z/p$ and is nonzero.  Primes $p>N$ are
arbitrarily large, proving the failure of (136.4).
For $N=1$ the PD ideal is the entire positive
coefficient ideal; the counterargument does not
apply.  That value would require a target's
weight-one row to be acyclic, which general
admissibility does not guarantee.
\end{proof}

The rational theorem thus does not extend by the
single-coefficient substitution (136.18), even
though its vertical Poincare lemma is valid.
This is a counterexample to that particular
construction, not to the integral common-source
or integral terminal-localization assertions.
Those two integral assertions remain distinct
open gates in this module.  Their original
defining sheets are now recovered, so the gate
is no longer the absence of the original question.


\paragraph{AQ-136-PLOCAL: common sources over every $\mathbb Z_{(p_0)}$-algebra.}
Fix a prime $p_0$.  If $k$ is any commutative
$\mathbb Z_{(p_0)}$-algebra, every finite family of admissible
strict PD theories in $\mathcal A_k$ has a common admissible
source.  In particular this covers $\mathbb Z_{p_0}$,
$\mathbb F_{p_0}$ and $\mathbb Z/p_0^r$ with their original
coefficients.  Terminality of the $p_0$-local all-arrow groupoid
is proved separately in AQ-136-PLOCAL-LOCALIZATION below.

\begin{proof}
Put $S=\Gamma_k(T)$ and choose $N=p_0^e$, $e\geq0$.
The PD homomorphism
\[
 \rho_N:\Gamma_k(U)\longrightarrow\Gamma_k(T),
                  \qquad U\longmapsto T^{[N]},              \tag{136.21}
\]
becomes a homogeneous map if $U$ is assigned weight $N$.
For every $m\geq0$ it sends the basis element $U^{[m]}$ to
\[
 c_mT^{[mN]},\qquad c_m=\frac{(mN)!}{m!(N!)^m}.
                                                               \tag{136.22}
\]
The integer $c_m$ is a $p_0$-adic unit.  Indeed counting the
multiples of successive powers of $p_0$ in a factorial gives
\[
 v_{p_0}((mp_0^e)!)
  =m\frac{p_0^e-1}{p_0-1}+v_{p_0}(m!),
 \qquad
 v_{p_0}((p_0^e)!)=\frac{p_0^e-1}{p_0-1};
\]
the difference in (136.22) is zero.
For $0\leq a<N$, the integer
$\binom{mN+a}{a}$ is also a $p_0$-adic unit.
For powers $p_0^j$ with $j\leq e$, divisibility of $mN$
by $p_0^j$ makes the corresponding three factorial floors
cancel.  For $j>e$, $a<p_0^e$ implies
\[
 \left\lfloor\frac{mN+a}{p_0^j}\right\rfloor
  =\left\lfloor\frac{m}{p_0^{j-e}}\right\rfloor,
 \qquad \left\lfloor\frac a{p_0^j}\right\rfloor=0,
\]
so those floors also cancel.  These integers are therefore
units in every $\mathbb Z_{(p_0)}$-algebra $k$.

The products
\[
 T^{[a]}\rho_N(U^{[m]})
    =c_m\binom{mN+a}{a}T^{[mN+a]}
\]
are consequently unit multiples of all the distinct basis
elements of $S$.  This proves the actual graded-module
isomorphism
\[
 S=\bigoplus_{a=0}^{N-1}T^{[a]}\Gamma_k(U)                 \tag{136.23}
\]
via $\rho_N$.  In particular $S$ is free, not merely flat,
over the indicated divided-power coefficient algebra.
The claim holds even if $k$ has $p_0$-torsion or nilpotents.

Let $A$ be the admissible integral DP simplex theory
(136.5), with its coefficient named $U$, and regrade
all its weights by $N$.  Form
\[
 B_N^{\rm PD}=S\otimes_{\Gamma_k(U)}A .                    \tag{136.24}
\]
In each simplex this is exactly the PD quotient with
$\sum x_j=T^{[N]}$ and $\sum e_j=0$ from (136.18).
Eliminating $x_0,e_0$ identifies it with the free PD
polynomial de Rham complex over $S$ in the remaining
variables.  The integral contractions in (136.6) prove
vertical cohomology $S$ in degree zero and zero otherwise,
in every weight and every simplex.

For the horizontal calculation, (136.23) is a natural
direct sum decomposition of the whole simplicial module,
since its coefficients are constant simplicially.
The component of weight $w=Nm+a$, $0\leq a<N$, and
cochain degree $r$ is identified, as a simplicial
$k$-module, with $A_*^{m,r}$.  The latter is contractible
for $m>r$, by (136.7).  Thus
$\pi_{B_N^{\rm PD}}(w)\geq\lfloor w/N\rfloor\to\infty$,
and $B_N^{\rm PD}$ is admissible.  The coefficient tail
is complete: every $T^{[w]}$ occurs in the free basis
calculation, with no rationalization.

Given finitely many admissible targets $C_b$, choose
$e$ so large that every row $C_b^{N,0}$ is acyclic.
Condition (136.4) permits this along the unbounded
sequence $N=p_0^e$.  Lift the constant coefficient
$T^{[N]}$ along $k\to L$ to a simplicial linear map
$L\to C_b^{N,0}$, using exactly the normalized-disk
argument of (136.15).  Its coordinate images have
sum $T^{[N]}$.  Send $x_j$ to those images and
$e_j$ to their vertical differentials.  The PD
polynomial universal property, together with
$d\gamma_m(x)=\gamma_{m-1}(x)dx$, gives a strict
simplicial differential PD map
$B_N^{\rm PD}\to C_b$ over $S$.
It preserves both gradings and the two sum relations.
These are the required simultaneous source maps.
\end{proof}

\paragraph{Finite-product and positive-characteristic consequence.}
Common-source existence also holds for a finite product
$k=\prod_{i=1}^r k_i$ of rings covered by AQ-136-PLOCAL.
Indeed the coefficient idempotents split every
$k$-theory, its divided powers and its morphisms into
$r$ component theories over $k_i$.
Both vertical cohomology and normalized simplicial
homology respect this finite direct-product decomposition.
Thus admissibility is equivalent to admissibility of all
components, using the maximum of their finitely many
cutoffs.  Take componentwise common sources and their
finite product; this is the required admissible source
over $k$.  Zero component rings cause no exception.

In particular common-source existence holds for every
commutative ring of positive characteristic.  If its
characteristic is $m=\prod_i p_i^{a_i}>1$, the Chinese
remainder idempotents of $\mathbb Z/m$ give
$k=\prod_i k_i$, each a $\mathbb Z/p_i^{a_i}$-algebra
and hence a $\mathbb Z_{(p_i)}$-algebra.  The preceding
argument applies.  The zero ring is harmless as well.
This finite idempotent decomposition is not available
for $\mathbb Z$; no prime-varying integral coefficient
object is substituted for the single ring in that
remaining question.

\paragraph{AQ-136-PATH-OBSTRUCTION: the unweighted rational interval has no characteristic-$p$ replacement of the same kind.}
Let $k\ne0$ be a commutative ring of characteristic
$p_0>0$.  There is no nonnegative commutative differential
$k$-algebra $R$ with coefficient isomorphism
$k\simeq H^0(R)$, two $k$-algebra augmentations
$\epsilon_0,\epsilon_1:R\to k$, and an element
$s\in R^0$ such that $\epsilon_0(s)=0$ and
$\epsilon_1(s)=1$.

\begin{proof}
Since $R$ is nonnegative, $H^0(R)=\ker(d:R^0\to R^1)$
without a quotient by boundaries.  The Leibniz rule and
characteristic $p_0$ give
$d(s^{p_0})=p_0s^{p_0-1}ds=0$.
The coefficient isomorphism therefore forces
$s^{p_0}$ to be the image of some $c\in k$.
Applying the two augmentations gives both $c=0$ and
$c=1$, contradicting $k\ne0$.
\end{proof}

This last obstruction rules out only such an unweighted
strict interval.  Weighted comparisons and other
localization arguments are not excluded.  In particular
it is consistent with AQ-136-PLOCAL and the terminality
proved in AQ-136-PLOCAL-LOCALIZATION below.  It supplies
no counterexample to integral terminal localization.


\paragraph{AQ-136-COTENSOR: a simplicial path theory over any coefficient ring.}
For any $C\in\mathcal A_k$ over any commutative ring $k$, put
\[
 P_s(C)_n^{w,r}=C^{w,r}(\Delta[n]\times\Delta[1]).             \tag{136.25}
\]
Here $C(X)$ is the compatible-section extension from simplices
used in AQ-136-COMPARISON.  The operations are taken pointwise.
Then $P_s(C)$ belongs to $\mathcal A_k$.  Its constant-path
inclusion $j:C\to P_s(C)$ and two evaluations
$e_0,e_1:P_s(C)\to C$ are strict PD morphisms over $S$,
and $e_0j=e_1j=1$.

\begin{proof}
The construction is a limit of the given bigraded differential
PD algebras, so inherits every operation and the coefficient
map.  Functoriality in $\Delta[n]$ gives its simplicial
structure; projection and the two vertex inclusions give
the displayed maps and their identities.

Fix a cochain bound $d$.  AQ-136-COMPARISON gives a weight
bound $W_d$ independent of the simplicial set.  Each
$\Delta[n]\times\Delta[1]$ is contractible: it is the nerve
of the product ordered set $[n]\times[1]$, which has a
terminal element, or use the usual monotone contraction of
each simplex.  Therefore, for $w\geq W_d$,
\[
 H^0(P_s(C)_n^{w,*})=S^w,\qquad
 H^r(P_s(C)_n^{w,*})=0\quad(1\leq r\leq d),                \tag{136.26}
\]
with the first isomorphism induced by the coefficient map.
The bound is uniform in $n$, as required by (136.3).
There is no interchange of infinitely many separate
simplex-dependent bounds.

For fixed $(w,r)$, write $F=C_*^{w,r}$.  The corresponding
simplicial module in (136.25) is its cotensor
$F^{\Delta[1]}$.  Evaluation at the zero vertex and the
constant map exhibit a simplicial linear homotopy
equivalence
\[
                       F^{\Delta[1]}\simeq F.              \tag{136.27}
\]
Explicitly the order-preserving minimum map
$[1]\times[1]\to[1]$ induces a simplicial map
$\mu:\Delta[1]\times\Delta[1]\to\Delta[1]$
with $\mu(s,0)=0$ and $\mu(s,1)=s$.
Precomposition by this map gives a homotopy between the
constant-path evaluation and the identity on the cotensor.
All its maps are linear, since operations on functions to
$F$ are pointwise.  Thus it induces the same normalized
simplicial homology as $F$, without a flatness hypothesis.
Condition (136.4) for $C$ consequently gives that condition
for $P_s(C)$, with the same cutoff function.  This proves
admissibility and all assertions.
\end{proof}

\paragraph{AQ-136-PLOCAL-LOCALIZATION: the $p_0$-local groupoid is terminal.}
For every commutative $\mathbb Z_{(p_0)}$-algebra $k$, the
ordinary localization of $\mathcal A_k$ at all its morphisms
is equivalent to the terminal groupoid.  The same holds
for a finite product of such coefficient rings and hence
for every ring of positive characteristic.  Together with
AQ-136-PLOCAL this proves both non-strict original
uniqueness assertions in these coefficient ranges.

\begin{proof}
We first show that any two parallel morphisms
$f,g:C\rightrightarrows D$ have the same image in the
all-arrow localization.  Choose $N=p_0^e$ so large that
both $C_*^{N,0}$ and $D_*^{N,0}$ are acyclic.  The proof
of AQ-136-PLOCAL provides an admissible source and map
\[
                         a:B_N^{\rm PD}\longrightarrow C.
\]
Set $F=D_*^{N,0}$.  The coordinate images under $fa$ and
$ga$ define simplicial linear maps
$\phi_0,\phi_1:L\to F$.  On the constant submodule
$kt\subset L$ both send $t=\sum_jx_j$ to the coefficient
$T^{[N]}$.  Their difference thus factors through
\[
             Q=L/kt,\qquad \psi:Q\longrightarrow F.        \tag{136.28}
\]
Normalization is exact, and the disk calculation for $L$
shows that $N(Q)=k[1]$.
The normalized map $N(\psi)$ is consequently a cycle in
$N(F)_1$.  Since $H_1(N(F))=0$, that cycle is the boundary
of an element of $N(F)_2$.  Mapping the generator of
$k[1]$ to this element gives a chain homotopy between
$N(\psi)$ and zero.  The chain-to-simplicial homotopy
statement, Stacks Project Lemma14.29.3/Tag\nolinkurl{01A4},
gives a simplicial linear homotopy from zero to $\psi$.
Equivalently it gives a linear map
$Q\to F^{\Delta[1]}$ with those endpoints.

Compose this map with $L\to Q$ and add the constant-path
map associated with $\phi_0$.  The result is a simplicial
linear map
\[
 \Phi:L\longrightarrow F^{\Delta[1]},\qquad
 e_0\Phi=\phi_0,\quad e_1\Phi=\phi_1,\quad
 \Phi(t)=T^{[N]}\text{ as a constant path}.                 \tag{136.29}
\]
The last equality is exact: the homotopy term factors
through $Q$ and is zero on $kt$, throughout the homotopy.
This relative condition is what preserves the original
coefficient relation.

Apply the free PD polynomial universal property in each
simplex, sending $x_j$ to $\Phi(x_j)$, $e_j$ to
$d\Phi(x_j)$, and every coefficient of $S$ to itself.
Equation (136.29) gives the two relations
$\sum x_j=T^{[N]}$ and $\sum e_j=0$; all operators are
natural in the simplex.  We obtain a strict differential
PD map over $S$
\[
 H:B_N^{\rm PD}\longrightarrow P_s(D),\qquad
                         e_0H=fa,\quad e_1H=ga.            \tag{136.30}
\]
It preserves both gradings, because the coordinate
homotopy was carried out within weight $N$, cochain
degree zero.  AQ-136-COTENSOR proves that its target is
admissible over the same coefficient ring.

Since $e_0j=e_1j=1$, the two evaluations both equal
$j^{-1}$ after inversion of all morphisms.  Thus
$fa=ga$ in the localization, and inversion of $a$
gives $f=g$ there.
Finite common sources exist by AQ-136-PLOCAL.
For any two finite zigzags with the same endpoints,
choose one common source for all their objects.  For
an original arrow $u:C_i\to C_j$, its composite with
the chosen source arrow to $C_i$ and the chosen source
arrow to $C_j$ are parallel; they have the same localized
image by the argument just given.  Every edge is
therefore the corresponding ratio of these source
arrows, and both zigzags telescope to the same ratio.
The localization is nonempty and connected, and has
exactly one arrow between any two objects.

For a finite product of coefficient rings, choose the
sources $B_{N_i}^{\rm PD}$ and the relative coordinate
homotopies separately in each factor.  Their finite
products are admissible and form the same diagrams over
the product ring: finite idempotent decomposition
commutes with compatible sections and hence with
$P_s$.  The identical parallel-arrow and telescope
argument applies.  The Chinese-remainder decomposition
proved in AQ-136-PLOCAL covers every positive-characteristic
ring.  The zero ring again causes no exception.
\end{proof}

The simplicial cotensor varies the argument $X$ of a
theory.  It does not adjoin a weight-zero element with
endpoints zero and one to one cochain algebra; the
hypotheses of AQ-136-PATH-OBSTRUCTION are therefore
not being asserted.  All coefficient divided powers
remain fixed in (136.29)--(136.30).
The two non-strict questions over $\mathbb Z$ remain
unresolved by this argument, since its admissible
source $B_N^{\rm PD}$ uses the one-prime free coefficient
decomposition (136.23).

\paragraph{AQ-136-WEIGHT-ONE: admissibility does not force the lowest row to be acyclic.}
For $N\geq2$ and $k\ne0$, there is an admissible theory
$C$ to which the standard source $B=B_1$ has no
$S$-morphism.  Define, inside (136.5),
\[
 C^{w,r}_n=
 \begin{cases}
 S^w,&w<N,\ r=0,\\
 0,&w<N,\ r>0,\\
 B^{w,r}_n,&w\geq N.
 \end{cases}                                               \tag{136.31}
\]

\begin{proof}
These submodules are preserved by $d$ and simplicial
operators.  Products involving a term of weight at
least $N$ remain in that range; products of lower-weight
coefficients are coefficients.  Divided powers likewise
preserve the subalgebra: they preserve the coefficient
algebra, and multiply a positive weight at least $N$
by a positive integer.  Thus $C$ is a strict differential
PD $S$-subalgebra of $B$.
Its vertical cohomology is $S$ in degree zero and zero
otherwise in every weight.  All sufficiently large
weight rows are the rows of $B$, so (136.4) holds.
Consequently $C$ is admissible.

Its weight-one zero-form row is the nonzero constant
simplicial module $kT$.  A hypothetical $S$-map
$B\to C$ would restrict on the coordinate module to
a simplicial linear map $L\to\underline{kT}$ taking
$t$ to $T$.  On normalization, a chain map from the
disk $k\xrightarrow{\pm1}k$ to $k[0]$ must have zero
degree-zero component, because its composite with
the disk differential is zero.  It cannot take $t$
to nonzero $T$.  This is the required contradiction.
\end{proof}


\paragraph{AQ-136-MULTIWEIGHT: independent coefficient disks in every large weight are still not admissible.}
Fix $N_0\geq1$.  One might try to remove the coefficient-tail
obstruction of AQ-136-INTEGRAL-TAIL by imposing all relations
$\sum_jx_{N,j}=T^{[N]}$ for every $N\geq N_0$, using a
separate divided-power simplex disk for each relation.
This construction has the required vertical cohomology,
but its zero-form row has nonzero horizontal $H_1$ in
arbitrarily large weights, already after extension to
$\mathbb Q$.  It therefore fails condition (136.4), both
over $\mathbb Q$ and over $\mathbb Z$.

\begin{proof}
First define the integral construction precisely.  Put
\[
 R=\Gamma_{\mathbb Z}(U_N:N\geq N_0),\qquad
 S=\Gamma_{\mathbb Z}(T),\qquad
 U_N\longmapsto T^{[N]},\quad \operatorname{wt}(U_N)=N.
                                                               \tag{136.32}
\]
In simplicial degree $n$, let
\[
 A_n=\Gamma_{\mathbb Z}(x_{N,j}:N\geq N_0,\ 0\leq j\leq n)
       \otimes\Lambda(e_{N,j}:N\geq N_0,\ 0\leq j\leq n)
       /(\textstyle\sum_j e_{N,j}:N\geq N_0),                \tag{136.33}
\]
where $d x_{N,j}=e_{N,j}$, $d e_{N,j}=0$, both generators
have weight $N$, and their cochain degrees are zero and
one, respectively.  The coefficient map sends
$U_N$ to $\sum_jx_{N,j}$.  Simplicial operators act on
each $N$-family by the sum-over-fibres rule of (136.5).
All divided powers and differential identities are those
of the integral model AQ-136-DP-MODEL; tensoring the
models and taking the filtered union over finite sets
of indices gives the displayed strict PD algebra.
Define
\[
                         B^{\rm multi}=S\otimes_R A.        \tag{136.34}
\]
Equivalently this is the strict PD quotient with
$\sum_jx_{N,j}=T^{[N]}$ and $\sum_j e_{N,j}=0$ for every
$N\geq N_0$.  Eliminating $x_{N,0}$ and $e_{N,0}$
identifies each simplex with the free divided-power
polynomial de Rham algebra over $S$ on the remaining
variables.  In a fixed weight only finitely many $N$
can occur, and in a fixed simplex there are finitely
many variables for each such $N$.  The weight-preserving
integral disk contractions from (136.6) therefore apply
in each weight, proving vertical cohomology $S$ in
degree zero and zero in every positive degree.

We compute the zero-form horizontal homology.  An
invertible linear change of divided-power variables
replaces $x_{N,0}$ by $U_N=\sum_jx_{N,j}$, independently
for every $N$.  Hence $A_n^0$ is a free $R$-module,
with divided-power monomials in $x_{N,j}$, $j\geq1$,
as basis.  Its augmentation to $\mathbb Z$ sends every
variable to zero.  As a simplicial abelian group, its
weight-$w$ part is the direct sum over the finitely
many tuples $(a_N)$ satisfying $\sum_N N a_N=w$ of
\[
                         \bigotimes_N\Gamma^{a_N}(L).
                                                               \tag{136.35}
\]
For $w>0$, at least one $a_N$ is positive.  That factor
is linearly simplicially contractible by the homotopy
calculation preceding (136.5), and tensoring its
homotopy with the identities contracts the displayed
tensor product.  Weight zero is the constant module
$\mathbb Z$.  Thus the augmented normalized complex
$N(A_*^0)\to\mathbb Z$ is exact.  Each normalized term
is an $R$-linear direct summand of a free term $A_n^0$
by the standard degeneracy decomposition.  It is
therefore a projective $R$-module, so this is a
projective resolution of the augmentation module
$\mathbb Z$ over $R$.

Normalization commutes with arbitrary extension of
scalars here: its degeneracy decomposition is split
and functorial, so tensoring carries the projector
onto the normalized summand to the corresponding
projector after tensoring.  Consequently
\[
 H_i\bigl(N((B^{\rm multi})_*^0)\bigr)
                  =\operatorname{Tor}_i^R(S,\mathbb Z).
                                                               \tag{136.36}
\]
Rationalization is flat.  In characteristic zero the
divided-power polynomial algebra is the polynomial
algebra, with $T^{[N]}=t^N/N!$.  Thus (136.36) after
tensoring with $\mathbb Q$ is computed over
\[
 R_{\mathbb Q}=\mathbb Q[u_N:N\geq N_0],\qquad
 S_{\mathbb Q}=\mathbb Q[t],\qquad
                     u_N\longmapsto a_N=t^N/N!.
                                                               \tag{136.37}
\]
The Koszul complex on the countable polynomial
variables $u_N$ is a free resolution of the
augmentation module $\mathbb Q$.  For completeness,
the Koszul complex on a finite initial set is exact
in positive degrees because adjoining a new
variable tensors the previous resolution with the
two-term complex of multiplication by that variable,
which is injective on the preceding polynomial
quotient.  Its degree-zero homology is the quotient
by the variables in that finite set.  Filtered
unions are exact for modules, so passage through
all finite initial sets gives positive homology
zero and degree-zero quotient $\mathbb Q$.

After tensoring with $\mathbb Q[t]$, we therefore
obtain the Koszul complex $K(a_N:N\geq N_0)$.
Write its degree-one generator as $\xi_N$, with
weight $N$ and $d\xi_N=t^N/N!$.  For $N>N_0$ put
\[
 \eta_N=\xi_N-
              \frac{N_0!}{N!}t^{N-N_0}\xi_{N_0}.
                                                               \tag{136.38}
\]
This is a homogeneous invertible change of basis
of the free direct sum on the generators: both
it and its inverse use only two terms for each
generator.  Moreover $d\eta_N=0$.  It follows that
\[
 K(a_N:N\geq N_0)
  \simeq K(t^{N_0}/N_0!)\otimes_{\mathbb Q[t]}
                 \Lambda_{\mathbb Q[t]}(\eta_N:N>N_0),
                  \qquad d\eta_N=0.                       \tag{136.39}
\]
Multiplication by $t^{N_0}/N_0!$ is injective in
$\mathbb Q[t]$.  The homology of its two-term complex
is therefore $\mathbb Q[t]/(t^{N_0})$ in degree zero.
The second factor in (136.39) is free in every
degree with zero differential, whence
\[
 H_*\bigl(K(a_N:N\geq N_0)\bigr)
  =\mathbb Q[t]/(t^{N_0})\otimes_{\mathbb Q}
                         \Lambda_{\mathbb Q}(\eta_N:N>N_0).
                                                               \tag{136.40}
\]
In particular $1\otimes\eta_N$ is a nonzero
horizontal $H_1$ class of weight $N$ for every
$N>N_0$.  The zero-form row is not horizontally
acyclic in any eventual range of weights, in
contradiction to condition (136.4).  Since flat
rationalization preserves homology, integral
eventual acyclicity would imply rational eventual
acyclicity.  The integral construction fails too.
\end{proof}

The argument applies even when $N_0=1$.  The
obstruction consists of the higher homology of
the redundant, independent coefficient relations;
it is distinct from the prime-weight $H_0$
obstruction in AQ-136-INTEGRAL-TAIL.  It refutes
this free independent-disk construction only.
Adding coherent higher relations or constructing
a different admissible source remains a separate
problem; the integral original conjecture (a)
and integral terminal localization are not
refuted by (136.40).


\paragraph{AQ-136-DERHAM-CELLS: coherent high-weight zero-form cells still cannot give an integral full de Rham source.}
Let $k$ be a commutative ring and put $S=\Gamma_k(T)$,
with $\operatorname{wt}(T)=1$.  Suppose $P_*$ is a
simplicial graded PD $S$-algebra such that in each
simplicial degree it has a homogeneous presentation
\[
                  P_n=S\langle X_n\rangle_{\rm PD},
           \qquad \operatorname{wt}(x)\geq2
                      \quad(x\in X_n).                    \tag{136.41}
\]
The sets of variables may be infinite and may differ
with $n$.  Simplicial operators need not preserve a
chosen generating set: they are arbitrary homogeneous
PD $S$-algebra morphisms.  In particular they can encode
coherent higher relations among the cells.

Write $\Omega^1_{P/S,\mathrm{sp}}$ for the module of
\emph{special} relative PD differentials: it represents
derivations satisfying
\[
                   d\gamma_m(a)=\gamma_{m-1}(a)\,da.
                                                               \tag{136.42}
\]
Thus, in (136.41), it is the free $P_n$-module with
basis $dx$ for $x\in X_n$, of the same weights as
the variables.  This is precisely the one-form row
of the full relative PD polynomial de Rham algebra.
It is not an identification with unrestricted Beck
derivations in the algebraic theory of all PD operations.

For any prime $p$ with $k/pk\ne0$, the following
implication holds:
\[
 [T^{[p]}]=0\ \hbox{in }H_0(N(P_*^p))
 \quad\Longrightarrow\quad
 H_1\bigl(N(\Omega^1_{P_*/S,\mathrm{sp}}{}^{\!p})\bigr)
                                  \ne0.                    \tag{136.43}
\]
Here the superscript $p$ denotes weight, and $N$
denotes simplicial normalization, not a weight bound.
Consequently, if infinitely many rational primes
are nonunits in $k$, no full relative PD de Rham
theory built from (136.41) can satisfy original
condition (136.4).  In particular this rules out
that entire class of constructions over $\mathbb Z$.

\begin{proof}
Fix $p$ with $\bar k=k/pk\ne0$ and reduce the
simplicial algebra modulo $p$.  All its elements
have nonnegative weight, and its weight-zero part
is $\bar k$.  Let $\bar P_+$ denote the ideal of
positive-weight elements and define the weight-$p$
ordinary indecomposables
\[
                  Q_*=(\bar P_+/(\bar P_+)^2)^p.
                                                               \tag{136.44}
\]
This is a simplicial $\bar k$-module.  We emphasize
\emph{ordinary} indecomposables: divided powers of
positive elements are not added to the denominator.
Simplicial maps preserve the ordinary product ideal,
so no invariance of the displayed free variables
under faces is needed.

The divided-power monomial basis gives, in every
simplicial degree, the following basis of $Q_n$:
\[
       z=T^{[p]},\qquad x\in X_n\text{ with }
                                      \operatorname{wt}(x)=p.
                                                               \tag{136.45}
\]
To check this assertion completely, a monomial
involving two distinct factors among the coefficient
variable $T$ and the free variables is a product of
two positive-weight monomials, with coefficient one,
and hence is zero in (136.44).  A monomial involving
only one free variable has the form $x^{[m]}$ with
$m\operatorname{wt}(x)=p$.  Since $p$ is prime and
$\operatorname{wt}(x)\geq2$, this forces $m=1$
and $\operatorname{wt}(x)=p$.  The only remaining
coefficient monomial is $T^{[p]}$.  Any product
of two positive pure coefficient monomials of
total weight $p$ is
\[
           T^{[a]}T^{[p-a]}=\binom pa T^{[p]}=0
                    \quad(0<a<p)                         \tag{136.46}
\]
in characteristic $p$.  Products involving a free
variable cannot contribute to this coefficient
basis vector or to a single free generator.
Thus none of the vectors in (136.45) is killed or
related by the product ideal.  This proves the
asserted free $\bar k$-basis, without assuming that
$\bar k$ is a field or reduced.

There is a canonical simplicial augmentation
$\bar P_*\to\bar k$ sending every positive-weight
element to zero.  Form the simplicial module
\[
 F_*=
 \left(\bar k\otimes_{\bar P_*}
              \Omega^1_{\bar P_*/\bar S,\mathrm{sp}}\right)^p.
                                                               \tag{136.47}
\]
Its degree-$n$ term is free on the symbols $dx$
with $x\in X_n$ of weight $p$.  Composing the
relative derivation with the augmentation kills
products of positive-weight elements by the
Leibniz rule.  It therefore induces a simplicial
linear map $Q_*\to F_*$, taking $x$ to $dx$ and
$z$ to zero, because coefficients in $S$ have
zero relative differential.  The bases just
computed give an exact sequence of simplicial
$\bar k$-modules
\[
 0\longrightarrow\underline{\bar k z}
       \longrightarrow Q_*
       \longrightarrow F_*\longrightarrow0.              \tag{136.48}
\]
The first module is constant: all simplicial
operators fix the original coefficient $T^{[p]}$.

The premise of (136.43) supplies a normalized
one-chain $u\in N_1(P_*^p)$ such that
\[
                            \partial u=T^{[p]}.
                                                               \tag{136.49}
\]
Apply the relative PD derivation to $u$.  It
commutes with all simplicial maps, so
\[
 du\in N_1(\Omega^1_{P_*/S,\mathrm{sp}}{}^{\!p}),
                 \qquad\partial(du)=dT^{[p]}=0.           \tag{136.50}
\]
Reduce $u$ modulo $p$ and project it to $Q_*$;
write $v$ for its normalized image.  Then
$\partial v=z$.  Its image $a$ in $N_1(F_*)$
is a cycle.  The connecting map for (136.48)
sends its homology class to
\[
                         [z]\in H_0(\underline{\bar k z})
                                      =\bar k z,
                                                               \tag{136.51}
\]
which is nonzero.  Hence $[a]\ne0$.
Equivalently, this last step can be checked
without the long exact sequence.  If $a=\partial b$,
exactness of normalization permits lifting
$b\in N_2(F_*)$ to $\widetilde b\in N_2(Q_*)$.
Then $v-\partial\widetilde b$ belongs to
$N_1(\underline{\bar k z})=0$, whereas its
boundary is $z\ne0$, a contradiction.

The class $a$ is exactly the image of $du$ under
the chain map obtained by reduction modulo $p$
and augmentation of the differential module.
If $du$ were a boundary integrally over $k$,
its image would be a boundary in $F_*$.  The
nonzero class in (136.51) therefore proves
that (136.50) is itself a nonzero $H_1$ class.
No flat-base-change assertion for homology or
identification with an abstract cotangent complex
has been used.

Finally, condition (136.4) in cochain degree zero
would make $P_*^w$ acyclic for every sufficiently
large weight $w$, so the premise of (136.43)
would hold at every sufficiently large prime.
If infinitely many primes are nonunits in $k$,
there are arbitrarily large such primes, and
(136.43) contradicts eventual horizontal acyclicity
of the one-form row.  This proves the consequence.
\end{proof}

For a ring, $k/pk\ne0$ is equivalent to $p$ not
being a unit.  Thus this construction obstruction
is consistent with the established $p$-local and
positive-characteristic results: only finitely
many rational primes can remain nonunits in
those cases.  It is also consistent with the
standard integral theory (136.5), whose simplex
presentations have new variables of weight one.

In particular, a relative free simplicial PD
resolution of the truncated coefficient algebra
$S/(T^{[w]}:w\geq N)$, $N\geq2$, that is built
using only new variables of weights at least
$N$, cannot become admissible merely by taking
its full relative special-PD de Rham algebra.
Higher simplicial cells may repair every
zero-form relation; (136.43) then forces the
one-form obstruction.  This rules out that
coherent repair, not all possible coherent
repairs: weight-one variables, additional form
relations, or other non-free theories are
outside the theorem.  The original integral
common-source and localization questions remain
separate.


\paragraph{AQ-136-TENSOR: the flat tensor target already fails.}
The supplied SA08 intake already reports the corresponding
prime-power obstruction in its statement A-S09-26, at
\texttt{sections/A\_hadamard\_image.tex:25}. That prior full
proof is not present among the supplied sources. The proof here
rederives the tensor assertion completely and makes no novelty
claim; it does not certify the intake's additional assertion
about a full PD closure in a rational Hadamard algebra.

Let $k$ be a commutative ring, $S=\Gamma_k(T)$, and let $B$ be
the admissible divided-power simplex theory in AQ-136-DP-MODEL.
Form the actual simplicial PD differential $S$-algebra
\[
                         E=B\otimes_S B.
\]
Both inclusions $B\to E$ preserve the coefficient algebra.
Every simplex of $E$ is free over $S$ in each cochain degree,
and has vertical cohomology $S$ in degree zero and zero in
positive degrees, in every weight. Nevertheless, for every
rational prime $p$ and integer $r\geq1$,
\[
 H_1\bigl(N(E^{p^r,0}_*)\bigr)\simeq k/pk.                 \tag{136.T1}
\]
Consequently $E$ fails the horizontal admissibility condition
over every nonzero ring which is not a $\mathbb Q$-algebra.
In particular, this failure occurs even over the $p$-local
and positive-characteristic rings for which the original
common-source and terminal-groupoid questions are already
answered affirmatively above.

\begin{proof}
In simplex degree $n$, eliminate $x_0$ and $y_0$ using
$\sum x_j=\sum y_j=T$, and eliminate $e_0,f_0$ using
$\sum e_j=\sum f_j=0$. The ordinary tensor has the explicit
PD presentation
\[
 E_n=S\langle x_1,\ldots,x_n,y_1,\ldots,y_n\rangle
       \otimes_k\Lambda_k(e_1,\ldots,e_n,f_1,\ldots,f_n),
 \qquad dx_j=e_j,\quad dy_j=f_j.
\]
Here angle brackets mean free divided-power variables and
all the displayed variables have weight one. Their cochain
degrees are zero for $x_j,y_j$ and one for $e_j,f_j$.
This is a tensor of $2n$ divided-power de Rham disks over $S$.
Each such disk has augmentation to $S$, and its positive
part contracts by
\[
                    h(z^{[a]}\,dz)=z^{[a+1]}.
\]
This is an $S$-linear, weight-preserving contraction.
Tensoring the finite number of contractions proves the
asserted vertical cohomology and simplexwise freeness,
without a flatness hypothesis on $k$.

Put $P_*=B_*^{*,0}$, regarded now as a simplicial graded
$S$-module, and augment it to $k=S/(T^{[1]},T^{[2]},\ldots)$
by killing all positive weights. In simplex degree $n$,
\[
 P_n=\Gamma_k(x_0,\ldots,x_n)
       \simeq S\langle x_1,\ldots,x_n\rangle,
\]
so it is a free graded $S$-module. In weight $w>0$ its
underlying simplicial $k$-module is $\Gamma^w_k(L)$, with
$L_n=k^{n+1}$ and the summation simplicial maps already
specified in AQ-136-DP-MODEL. The explicit linear simplicial
contraction of $L$ induces a contraction of $\Gamma^w(L)$.
Weight zero is the constant $k$. Therefore $N(P_*)\to k$
is a resolution, and its terms are projective graded
$S$-modules: normalization is the image of a degreewise
splitting idempotent on $P_n$.

The zero-form row of $E$ is the diagonal tensor $P_*\otimes_S P_*$.
The Eilenberg--Zilber comparison for flat simplicial modules
identifies its normalized complex, up to quasi-isomorphism,
with
\[
                  N(P_*)\otimes_S N(P_*).
\]
An exact primary formulation of this input is Stacks,
Remark 21.39.11, Tag \nolinkurl{08QD}, following
Lemma 21.39.10. The statement there is first made for the
associated unnormalized complexes; the split degenerate
summands give the normalized formulation used here.
The comparison respects the internal weight grading.
Both factors are nonnegative projective resolutions of
$k$, so their tensor computes $k\otimes_S^{\mathbf L}k$.
One may check the last assertion by the first-quadrant
double-complex filtration: tensoring one resolution with
each projective term of the other preserves its homology,
and each total degree has finitely many bidegrees. Thus
\[
 H_i\bigl(N(E^{*,0}_*)\bigr)
                         \simeq\operatorname{Tor}_i^S(k,k).
                                                               \tag{136.T2}
\]

Let $I=\bigoplus_{w>0}S^w$. Apply $-\otimes_S k$ to
$0\to I\to S\to k\to0$. Since $S$ is projective over
itself and the map $S\otimes_S k\to k\otimes_S k$ is
the identity of $k$, its exact Tor sequence gives
\[
                         \operatorname{Tor}_1^S(k,k)
                           \simeq I/I^2.
\]
For $m>1$, the weight-$m$ component of $I^2$ is generated
inside $kT^{[m]}$ by
\[
              \binom ma T^{[m]},\qquad 1\leq a<m.
\]
If $m=p^r$, every displayed binomial coefficient is divisible
by $p$, because $(1+X)^{p^r}=1+X^{p^r}$ in $\mathbb F_p[X]$.
Their integer greatest common divisor divides
$\binom{p^r}{1}=p^r$. Moreover
\[
 \binom{p^r}{p^{r-1}}
   =p\binom{p^r-1}{p^{r-1}-1}
\]
has $p$-adic valuation exactly one: modulo $p$, all the
coefficients of $(1+X)^{p^r-1}$ are alternately $1,-1$,
as follows by dividing $1+X^{p^r}$ by $1+X$.
The greatest common divisor is therefore exactly $p$.
Bezout's identity shows that these coefficients generate
the ideal $pk$ after any base change to $k$. Hence
$(I/I^2)^{p^r}=k/pk$, proving (136.T1).

If $k$ is a nonzero ring without a $\mathbb Q$-algebra
structure, some rational prime $p$ is not a unit in $k$;
otherwise the map from $\mathbb Z$ would extend to its
localization $\mathbb Q$. For that prime, $k/pk\ne0$.
Equation (136.T1) gives nonzero first simplicial homology
in arbitrarily high weights of the zero-form row. This
contradicts (136.4). \qedhere
\end{proof}

This is a failure of the displayed common-target construction,
including in a simplexwise free and vertically acyclic case.
Replacing a module tensor by its derived tensor does not
establish horizontal admissibility: (136.T2) already is the
derived-module calculation. No model structure on strict PD
cochain algebras, or claim that every possible PD derived
pushout presentation must have this zero-form row, is assumed.
In particular the result neither refutes the original
common-source conjecture nor the integral terminal-groupoid
assertion. Horizontal admissibility is additional information
about an actual simplicial presentation, not a consequence
of its simplexwise vertical derived isomorphism class.


\paragraph{AQ-136-SEMILOCAL: arithmetic path gluing for the ordinary localization.}
Let $P=\{p_1,\ldots,p_s\}$ be a finite nonempty set of rational
primes, and put
\[
 R_P=\mathbb Z[\ell^{-1}:\ell\text{ prime},\ \ell\notin P].
\]
For every commutative $R_P$-algebra $k$, the ordinary groupoid
obtained from $\mathcal A_k$ by inverting every morphism is
equivalent to the terminal groupoid. Together with the rational
case, this proves the assertion for every commutative ring in
which all but finitely many rational primes are units.

This is a theorem about the original localization question.
The present proof constructs common sources after specified
natural arrows to arithmetic path replacements. A subsequent
result, AQ-136-SEMILOCAL-SOURCES below, combines this theorem
with admissible cospan path pullbacks to obtain strict common
sources for the original targets as well. The integral ring $\mathbb Z$, and rings
with infinitely many nonunit rational primes, are not included.
The zero ring is harmless: its unital algebra category has
only the zero algebra. Assume henceforth that $k\ne0$.

\paragraph{1. An exact arithmetic sequence for arbitrary modules.}
Set
\[
 k_i=k\otimes_{R_P}\mathbb Z_{(p_i)},\qquad
 K=k\otimes_{R_P}\mathbb Q.
\]
These are localizations of $k$: to obtain $k_i$, invert the
product of the primes in $P\setminus\{p_i\}$, and to obtain
$K$, invert the product of all primes in $P$.
For every $k$-module $M$, write $M_i=M\otimes_k k_i$ and
$M_0=M\otimes_k K$. There is an exact sequence
\begin{equation}
 \begin{aligned}
 0\longrightarrow M\longrightarrow
       \left(\bigoplus_{i=1}^sM_i\right)\oplus M_0
   &\xrightarrow{\delta} M_0^{\oplus s}\longrightarrow0,\\
 \delta((m_i),m_0)&=(m_i-m_0)_i.
 \end{aligned}                                           \tag{136.S1}
\end{equation}

Here is a direct proof, including the base-change assertion.
For $M=R_P$, the kernel statement is
$\bigcap_i\mathbb Z_{(p_i)}=R_P$ inside $\mathbb Q$.
For surjectivity, given finitely many $a_i\in\mathbb Q$,
choose a common denominator
$D=d\prod_i p_i^{e_i}$ with $d$ prime to every $p_i$.
Write $a_i=b_i/D$. The Chinese remainder theorem gives an
integer $b$ such that $b+b_i$ is divisible by $p_i^{e_i}$
for every $i$. With $a_0=b/D$ we have
$a_i+a_0\in\mathbb Z_{(p_i)}$, as required.
The rightmost module $\mathbb Q^{\oplus s}$ is flat over
$R_P$, so tensoring this short exact sequence with any
$R_P$-module remains exact. Apply this to the underlying
$R_P$-module of $M$. This proves (136.S1) for arbitrary
$M$, with no flatness assumption on $k$ or on $M$.

All localizations used below are flat. Thus vertical and
normalized horizontal homology commute with them, and an
admissible theory localizes to an admissible theory with
the same weight cutoffs. Normalization commutes with these
base changes, either by exactness or by its splitting
idempotents. The coefficient algebra localizes as
$\Gamma_k(T)\otimes_k k_i=\Gamma_{k_i}(T)$.

\paragraph{2. A rational tensor lemma with a uniform weight bound.}
Let $K$ be any $\mathbb Q$-algebra, possibly zero, and let
$N_1,\ldots,N_s$ be positive integers. Over $K$, let $B_{N_i}$
be the weighted simplex theory whose zero-form coordinates
have weight $N_i$ and whose coefficient relation is
$\sum_jx_{i,j}=T^{[N_i]}$. It is the rationalization of
the basic source in AQ-136-PLOCAL. Set
\[
 F=\bigotimes_{i=1}^s{}_{\Gamma_K(T)}B_{N_i}.
\]
Then $F$ is admissible. Its vertical cohomology is exactly
$\Gamma_K(T)$ in degree zero in every weight; for fixed
form degree $q$, its horizontal homology vanishes in all
weights
\[
                    w\ge q\max_iN_i+\sum_iN_i.           \tag{136.S2}
\]

We give the module calculation rather than infer this claim
from a vertical K\"unneth argument. Identify
$\Gamma_K(T)$ with $S_K=K[t]$, by $T^{[a]}=t^a/a!$.
Rescaling the coordinates by $N_i!$ writes the $i$th
coefficient relation as $\sum_jx_{i,j}=t^{N_i}$.
Before this base change, use the coefficient algebra
$K[U_i]$ and relation $U_i=\sum_jx_{i,j}$, where $U_i$
and the coordinates have weight $N_i$. Its form-$q_i$
row is
\[
 \bigoplus_{m\ge0}\operatorname{Sym}^m_K(L)
                         \otimes_K\Lambda^{q_i}_K(Q),
 \qquad Q=L/K\!\left(\sum_jx_j\right),                 \tag{136.S3}
\]
with weight $N_i(m+q_i)$ in the $m$th summand.
The positive-$m$ summands are simplicially contractible:
apply $\operatorname{Sym}^m$ to the explicit linear
contraction of $L$, and tensor the resulting simplicial
homotopy with the identity on $\Lambda^{q_i}Q$.
Hence the augmentation of this row onto
$\Lambda^{q_i}Q$, on which $U_i$ acts as zero, is a
horizontal quasi-isomorphism. Each simplex of the row
is a free $K[U_i]$-module: eliminate $x_{i,0}$ using
$U_i$, and use the exterior basis on the remaining
differentials.

For a fixed tuple $(q_1,\ldots,q_s)$ with sum $q$, take
the tensor of the rows (136.S3), using the independent
coefficient ring $A=K[U_1,\ldots,U_s]$.
The same contraction, on any factor with positive $m$,
shows that it maps by a horizontal quasi-isomorphism to
\[
                V_* =\bigotimes_i\Lambda^{q_i}Q_*.
\]
All of $V_*$ has the single internal weight
$b=\sum_iq_iN_i$, and every normalized term of $V_*$
is projective over $K$, being a direct summand of a
finite free module. The normalized source is a
nonnegative complex of projective $A$-modules. It
therefore computes the derived base change
$S_K\otimes_A^{\mathbf L}N(V_*)$.
Here ordinary base change of the normalized source is
the normalized row after base change; the split
normalization idempotents justify this even though
$A\to S_K$ need not be flat.

Resolve $K=A/(U_1,\ldots,U_s)$ by the Koszul complex
on the variables $U_i$. It is exact over arbitrary
$K$: successive multiplication by a polynomial variable
is injective. Tensoring this resolution with the
nonnegative projective $K$-complex $N(V_*)$ is a
projective $A$-resolution of that complex, by the
first-quadrant filtration. Thus the required derived
base change is computed by
\[
           K(t^{N_1},\ldots,t^{N_s};S_K)
                               \otimes_K N(V_*).         \tag{136.S4}
\]
Reorder the indices so $N_1=\min_iN_i$, and let $\xi_i$
be the Koszul generator of weight $N_i$ and homological
degree one. The invertible basis change
\[
        \eta_i=\xi_i-t^{N_i-N_1}\xi_1\quad(i>1)
\]
has $d\eta_i=0$. Consequently (136.S4) is
quasi-isomorphic to
\[
  S_K/(t^{N_1})\otimes_K
       \Lambda_K(\eta_2,\ldots,\eta_s)\otimes_K N(V_*).
                                                               \tag{136.S5}
\]
The passage to the quotient is valid after tensoring:
$N(V_*)$ is bounded below with projective $K$-terms,
and the exterior factors are free. Every term in
(136.S5) has weight at most
\[
       b+(N_1-1)+\sum_{i>1}N_i
                    \le q\max_iN_i+\sum_iN_i-1.
\]
There are finitely many tuples $(q_i)$ for a fixed $q$,
so this proves (136.S2). The tensor-to-normalized
comparisons used here are the flat Eilenberg--Zilber
comparison of Stacks, Remark 21.39.11, Tag \nolinkurl{08QD};
equivalently one can apply its first-quadrant total
complex form at each finite tensor step.

Finally, simplexwise, eliminate one zero-form coordinate
per factor using $t^{N_i}$. The vertical complex is
$S_K$ tensored with finitely many polynomial de Rham
disks. Integration over the $\mathbb Q$-algebra $K$
contracts each disk onto its constants, preserving
weight. This proves the stated vertical comparison.
Thus $F$ is admissible, with a bound independent of
the simplex dimension. This lemma includes $s=1$.

\paragraph{3. Arithmetic path replacements.}
For a differential $K$-algebra $D$, put
\[
 \operatorname{Path}_{\rm dR}(D)=D\otimes_K K[z,dz],
 \quad |z|=0,\quad |dz|=1,\quad\operatorname{wt}(z)
                  =\operatorname{wt}(dz)=0.
\]
Evaluation at $z=0,1$, $dz=0$ gives $e_0,e_1$.
Integration is a natural contraction of this path
complex onto $D$, so both evaluations induce the
same isomorphism in vertical cohomology. This
construction is used only over the $\mathbb Q$-algebra
$K$. Its positive-weight divided powers are uniquely
specified by $\gamma_m(a)=a^m/m!$ on even elements.

For $C\in\mathcal A_k$, write $C_i=C\otimes_k k_i$
and $C_0=C\otimes_k K$. Define $\mathcal E(C)$,
in every simplex degree, to consist of tuples
\[
       ((c_i)_{i=1}^s,c_0,(\omega_i)_{i=1}^s)
\]
with $c_i\in C_i$, $c_0\in C_0$, and
$\omega_i\in\operatorname{Path}_{\rm dR}(C_0)$,
subject to
\[
                       e_0\omega_i=c_i|_K,
             \qquad    e_1\omega_i=c_0.                   \tag{136.S6}
\]
This is a limit of strict bigraded PD differential
algebras, so it inherits all the required operations.
It is simplicial and functorial in $C$. The coefficient
map takes a coefficient to its localizations and
constant paths. The same prescription gives a natural
morphism
\[
                            \eta_C:C\longrightarrow\mathcal E(C).
                                                               \tag{136.S7}
\]

We verify admissibility explicitly. Forgetting
multiplication gives a short exact sequence of
simplicial cochain modules
\[
0\longrightarrow\mathcal E(C)\longrightarrow D(C)
 \longrightarrow C_0^{\oplus 2s}\longrightarrow0,          \tag{136.S8}
\]
where
$D(C)=(\bigoplus_i C_i)\oplus C_0
\oplus\operatorname{Path}_{\rm dR}(C_0)^{\oplus s}$,
and the last map is the endpoint difference in
(136.S6). It is surjective in every cochain and
simplex degree: $(1-z)a+zb$ gives prescribed
endpoints $(a,b)$. This section is linear and
simplicially natural; it need not be a cochain map.

Fix a cochain cutoff $d$. At sufficiently large weight,
uniformly in simplex degree, every localized vertical
complex has coefficient cohomology in degree zero
and zero cohomology in degrees $1$ through $d$.
The same is true of the rational path complexes.
On degree-zero cohomology, the cokernel of the map
in (136.S8), after eliminating the path constants,
is exactly the cokernel of (136.S1), and its kernel
is $S_k^w$. Thus the long exact vertical cohomology
sequence gives $H^0=S_k^w$ and $H^j=0$ for
$1\le j\le d$. All bounds are independent of
simplex degree. For fixed form degree $q$, the
path row is a direct sum of copies of the two rows
$C_0^{w,q}$ and $C_0^{w,q-1}$. Polynomial coefficients
are free over $K$, so these rows are horizontally
acyclic at large weight. The other rows in (136.S8)
are likewise acyclic, and its horizontal long exact
sequence proves the same for $\mathcal E(C)$.
This proves $\mathcal E(C)\in\mathcal A_k$.

We will also use its compatibility with the simplicial
cotensor $P_s$ of AQ-136-COTENSOR:
\[
                     \mathcal E(P_s C)\simeq P_s\mathcal E(C).
                                                               \tag{136.S9}
\]
For each $n$, the simplicial set $\Delta[n]\times\Delta[1]$
has finitely many nondegenerate simplices. Sections of
a simplicial module on it are therefore computed by
a finite limit of values on simplices. Flat localization
commutes with these limits. Tensoring with $K[z]$ also
commutes with them, since it is flat, and a fixed path
cochain degree is a sum of only the two displayed
form degrees. Limits commute with the finite limits
in (136.S6). These natural identifications preserve
products, the differential, divided powers and all
endpoint maps, proving (136.S9).

\paragraph{4. One local source and homotopies for an entire finite diagram.}
Let a finite directed graph be labelled by objects
$C_v\in\mathcal A_k$ and actual morphisms
$f_e:C_v\to C_{v'}$ on its edges. For each $i$, choose
a power $N_i=p_i^{a_i}$ so large that every zero-form
row $(C_v)_i^{N_i,0}$ is horizontally acyclic.
There is one such power for all vertices, since they
are finite in number. The p-local construction gives
the basic admissible source
$A_i=B_{N_i}(k_i)$ and maps
\[
                         a_{i,v}:A_i\longrightarrow(C_v)_i
                                                               \tag{136.S10}
\]
for every vertex.

For each edge, the two maps $f_ea_{i,v}$ and
$a_{i,v'}$ from this \emph{same} $A_i$ to $(C_{v'})_i$
have a coefficient-preserving simplicial PD homotopy
\[
                  h_{i,e}:A_i\longrightarrow P_s((C_{v'})_i).
                                                               \tag{136.S11}
\]
Here is the simultaneous-choice justification. On the
weight-$N_i$ coordinate module $L$, the difference of
the two coordinate maps kills the constant submodule
$k_i\sum x_j$, because both maps take it to $T^{[N_i]}$.
It factors through $Q=L/k_i\sum x_j$, whose normalized
complex is $k_i[1]$. Horizontal $H_1$ of the target
zero-form row vanishes. The resulting degree-one
cycle is therefore a boundary, giving a chain
nullhomotopy, and Dold--Kan gives a simplicial linear
homotopy. Adding the first coordinate map makes it
relative to the coefficient submodule. The free PD
coordinate presentation extends it, with its
differential, to (136.S11). This is precisely the
construction checked in AQ-136-PLOCAL-LOCALIZATION.
The same row cutoff works for every edge; choices for
the finitely many homotopies may be made separately
after the single common $A_i$ and all maps (136.S10)
have been fixed. The resulting data supply every
edge of the finite graph at once.

\paragraph{5. Gluing that finite diagram and killing every finite loop.}
Form the rational admissible theory
\[
                    F=\bigotimes_{i=1}^s{}_{S_K}(A_i\otimes_{k_i}K),
\]
which is covered by Step 2, and let $j_i:(A_i)_K\to F$
be the canonical inclusions. Define $A$ by the same
path-gluing formula as (136.S6), using local objects
$A_i$, central object $F$, and endpoint maps $j_i$:
\[
 A=\left\{((a_i),a_0,(\omega_i)):
     \begin{array}{l}
       e_0\omega_i=j_i(a_i|_K),\\
       e_1\omega_i=a_0\quad(1\le i\le s)
     \end{array}\right\}.                                 \tag{136.S12}
\]
Here $a_i\in A_i$, $a_0\in F$, and
$\omega_i\in\operatorname{Path}_{\rm dR}(F)$.
The verification in Step 3 applies verbatim to this
diagram of different theories. In detail, its endpoint
difference map is still degreewise surjective; all
local and central vertical cohomology groups are
exactly their coefficient rings, and the induced
maps are the coefficient localizations. Sequence
(136.S1) consequently gives vertical cohomology
exactly $S_k$ in every weight. Its horizontal rows
are eventually acyclic by the finitely many p-local
source cutoffs and the uniform rational bound (136.S2).
Hence $A\in\mathcal A_k$.

For each vertex, multiplication of the rationalized
maps (136.S10) defines a strict map
\[
        b_v:F\longrightarrow(C_v)_0,\qquad
        b_v(u_1\otimes\cdots\otimes u_s)
                   =\prod_i a_{i,v}(u_i).
\]
It is well defined over $S_K$, since all component
maps agree on that coefficient algebra. It preserves
the differential and the gradings; over the rational
algebra it also preserves the uniquely specified
divided powers. Its restriction along $j_i$ is
$(a_{i,v})_K$. Applying it to the central value and
the paths in (136.S12), with (136.S10) on the local
values, gives a strict map
\[
                         \alpha_v:A\longrightarrow\mathcal E(C_v).
                                                               \tag{136.S13}
\]

Likewise multiply the rationalized maps (136.S11)
to obtain $b_e:F\to P_s((C_{v'})_0)$. Its endpoints
are $f_eb_v$ and $b_{v'}$: this follows from
multiplicativity of $f_e$ and of the endpoint maps.
On local values use $h_{i,e}$, on the central value
use $b_e$, and on rational paths use its differential
path extension. This gives
\[
 A\longrightarrow\mathcal E(P_s C_{v'})
                         \simeq P_s\mathcal E(C_{v'}),
\]
whose endpoints are $\mathcal E(f_e)\alpha_v$ and
$\alpha_{v'}$. Thus the single source $A$ supplies
homotopies for every edge of the entire finite graph,
with all its vertex maps fixed.

In the all-arrow localization, the two evaluations
$P_sD\rightrightarrows D$ are equal: they are both
inverse to the constant-path arrow $D\to P_sD$.
All three objects and arrows are admissible by
AQ-136-COTENSOR. Consequently
\[
                   \mathcal E(f_e)\alpha_v=\alpha_{v'}
\]
in that localization. Naturality of (136.S7) now
shows that, with $\beta_v=\eta_{C_v}^{-1}\alpha_v$,
\[
                               f_e\beta_v=\beta_{v'}.
                                                               \tag{136.S14}
\]
Every $\beta_v$ is an isomorphism there, since all
original arrows are inverted.

An arbitrary morphism in the ordinary localization
is represented by a finite word of arrows and
their formal inverses. For any two such words with
the same endpoints, take the finite graph consisting
of all their objects and arrows. Equation (136.S14)
rewrites either word as
$\beta_{\rm end}\beta_{\rm start}^{-1}$, so the words
are equal. Taking two isolated vertices supplies
an isomorphism between any two objects. Finally,
the category is nonempty by AQ-136-DP-MODEL.
It is therefore equivalent to the terminal groupoid.
No higher coherence assertion about an
$\infty$-categorical localization is required or
inferred: the common cone and its edge homotopies
are used to compare arbitrary \emph{finite words}
in the ordinary all-arrow groupoid. \hfill$\square$


\paragraph{AQ-136-TAME: finite-degree repair does not justify a full acyclic closure.}
This finding records an obstruction to a prescribed construction,
not a counterexample to the original common-source or terminal-groupoid
question. It is compatible with the supplied SA08 intake's
statement-level reports about failed delayed bar primitives; their
unavailable full proofs are not certified here.

Let $k$ be any commutative ring, fix a positive weight $m$, and put
\[
 D=\Lambda_k(a)\otimes_k\Gamma_k(b),\qquad
 |a|=1,\quad |b|=2,\quad \operatorname{wt}(a)=\operatorname{wt}(b)=m,
 \qquad da=b,\quad db=0.                                  \tag{136.T1}
\]
Here $|\ |$ denotes form degree, and the positive even part has
its usual divided powers. Then
\[
 H^0(D)=k,\qquad
 H^{2r-1}(D)=k[r],\qquad H^{2r}(D)=k/rk\quad(r\ge1),         \tag{136.T2}
\]
where the two positive-degree groups in (136.T2) have weight
$mr$ and $k[r]=\ker(r:k\to k)$. In particular over
$\mathbb Z[1/d!]$ the augmentation ideal is acyclic through
form degree $2d$. If $p$ is a nonunit prime of a nonzero $k$,
the class of $b^{[p]}$ in form degree $2p$ is nonzero.

There is nevertheless an admissible strict integral theory
$C\in\mathcal A_k$ containing this constant simplicial algebra as a retract
in simplicial PD $k$-algebras. More precisely, let $B$ be
the standard theory of AQ-136-DP-MODEL, let both $B$ and $D$
map to $k$ by killing their positive-weight generators, and set
\[
                             C=B\times_k D.               \tag{136.T3}
\]
Give $C$ the coefficient map
$s\mapsto(s,\epsilon_S(s))$ from $S=\Gamma_k(T)$.
For $k=\mathbb Z$, no map $A\to C$ of strict PD theories
can simultaneously have an element $x\in A_n^{m,1}$ mapping
to $a$ and satisfy
\[
                     H^{2p}(A_n^{mp,*})=0                 \tag{136.T4}
\]
for a prime $p$. Thus a fully vertically acyclic PD replacement
with that prescribed lift cannot be assumed to exist.

\begin{proof}
The monomials $b^{[r]}$ and $ab^{[r]}$ form the indicated
free graded $k$-basis of $D$. The special differential rule
and the signed Leibniz rule give
\[
 d(b^{[r]})=0,\qquad
 d(ab^{[r-1]})=b b^{[r-1]}=r b^{[r]}.
\]
Thus, apart from $k$ in degree zero, the whole cochain complex
is the direct sum of the two-term blocks
\[
 k ab^{[r-1]}\xrightarrow{\ r\ }k b^{[r]},
       \qquad\text{in form degrees }2r-1,2r.
\]
Taking their kernels and cokernels proves (136.T2) over
arbitrary $k$, including rings with torsion or nilpotents.
The PD algebra and differential exist integrally: they are
the usual divided-power polynomial algebra tensored with
one exterior generator, and the displayed formulas satisfy
the special derivation identity. In each block with $r\le d$,
$r$ is a unit over $\mathbb Z[1/d!]$, proving the bounded-degree
claim. A nonunit prime has $k/pk\ne0$, proving the surviving
class assertion.

Both augmentations in (136.T3) preserve divided powers and
differentials. The fibre product has componentwise algebra,
PD and differential operations. The constant simplicial
structure on $D$ and the given structure on $B$ make these
operations simplicial. Its underlying bigraded module is
$B\oplus I_D$, where $I_D$ is the positive-degree augmentation
ideal of $D$; the unit is the common diagonal copy of $k$.
For each fixed positive form degree $q$, $I_D^q$ has only the
single weight $m\lceil q/2\rceil$.
Consequently beyond that weight the simplicial row of $C$
is exactly the row of $B$. Its horizontal homology therefore
vanishes in the required weight tail. Similarly, (136.T2)
shows that in each fixed positive vertical degree the extra
cohomology has finite weight support, independently of the
simplex. In degree zero there is no extra positive-weight
summand and $H^0(C_n)=S$ exactly. These are precisely both
uniform admissibility conditions, so $C\in\mathcal A_k$.
The map $D\to C$, $d\mapsto(\epsilon_D(d),d)$, splits its
projection to $D$.

Now take $k=\mathbb Z$ and a proposed map $f:A\to C$ as in
the statement. The element $dx$ is even of positive weight,
so its divided power $\gamma_p(dx)$ exists. It is closed,
because the special derivation rule gives
$d\gamma_p(dx)=\gamma_{p-1}(dx)d^2x=0$.
Condition (136.T4) would make it a boundary in $A_n$.
But its image under $f$ is $b^{[p]}$, whose class in $C_n$
is the nonzero element $1\bmod p$ of the direct summand
$\mathbb Z/p\mathbb Z$ in (136.T2). A chain map carries
boundaries to boundaries, a contradiction.
\end{proof}

The example has the admissible map
$B\to C$, $b\mapsto(b,\epsilon_B(b))$.
It therefore supplies no obstruction to a common source
or to terminal localization: only the specified acyclic
closure is excluded. The elementary disk calculation also
pinpoints the gap in a degreewise tame construction. Such a
construction still needs strict transition maps between whole
PD algebras and, for each fixed form degree, one eventual weight
bound uniform across stages and simplices. Moreover coefficient
multiplication raises weight without raising form degree, so
leaving a defect above a cutoff at its first weight does not
control its later coefficient multiples. No existence or
nonexistence of the required coherent relative construction
is inferred from this calculation.

\paragraph{AQ-136-PATH-LIFT: the arithmetic path arrow need not admit source lifting.}
Let $P=\{p_1,\ldots,p_s\}$ be nonempty and finite, and let $k$
be a $\mathbb Z_{(P)}$-algebra with
$K=k\otimes_{\mathbb Z}\mathbb Q\ne0$.
Use the exact arithmetic replacement $\mathcal E$ and natural
arrow $\eta_C:C\to\mathcal E(C)$ of AQ-136-SEMILOCAL.
The following proposed source-lifting property is false:
\[
 \begin{gathered}
 F:A\longrightarrow\mathcal E(C),\quad A,C\in\mathcal A_k
 \quad\Longrightarrow\quad
 \text{there exist }A'\in\mathcal A_k,\ h:A'\to A,\ u:A'\to C\\
 \text{ with }Fh=\eta_Cu.
 \end{gathered}                                             \tag{136.PL1}
\]
In particular it fails for $k=\mathbb Q$ and for mixed semilocal
rings such as $\mathbb Z_{(2,3)}$.
This is a failure of a specified way of refining an arrow;
it does not refute common sources or terminal localization.

\begin{proof}
Use the standard $B$ and its square-zero extension $D=B\ltimes M$
from AQ-136-STRICT. To fix every operation needed here,
\[
 B_n=\Gamma_k(x_0,\ldots,x_n)\otimes_k
       \Lambda_k(e_0,\ldots,e_n)/(\textstyle\sum e_j),
 \quad dx_j=e_j,\quad T=\textstyle\sum x_j .                \tag{136.PL2}
\]
The module $M$ is free over $B$ on $\alpha_j$ of form degree
zero and $\beta_j$ of form degree one, subject to
$\sum\alpha_j=\sum\beta_j=0$, with
$d\alpha_j=\beta_j$ and $d\beta_j=0$.
All four kinds of generator have weight one.
The degree-zero relative special derivation
\[
 \delta:B\longrightarrow M,\qquad
 \delta x_j=\alpha_j,\quad\delta e_j=\beta_j               \tag{136.PL3}
\]
commutes with $d$ and is zero on $S$.
The extension has $M^2=0$ and
$\gamma_r(b+m)=\gamma_r(b)+\gamma_{r-1}(b)m$ on its
positive-weight even ideal. Consequently
\[
       j_0(b)=(b,0),\qquad j_1(b)=(b,\delta b)              \tag{136.PL4}
\]
are strict simplicial PD differential $S$-algebra maps.

The earlier full proof verifies $B,D\in\mathcal A_k$.
Its bounds can also be read directly here: $M$ is vertically
contractible by
$h(b\beta_j)=(-1)^{|b|}b\alpha_j$, $h(b\alpha_j)=0$,
so $H^*(D_n)=S$ in every weight. Horizontally, with
$Q=L/k\sum x_j$ as before,
\[
 M_*^{w,q}=(B_*^{w-1,q}\otimes Q)
                 \oplus(B_*^{w-1,q-1}\otimes Q).
\]
Both summands are contractible for $w>q+1$ by the checked
linear contraction of the standard rows. Thus all bounds
are uniform in the simplex and no coefficient tail is removed.

Define a relative derivation $\iota:B\to M$ of form degree
$-1$ and weight zero by
\[
                  \iota(x_j)=0,\qquad\iota(e_j)=\alpha_j.
                                                               \tag{136.PL5}
\]
Equivalently, $\iota=h\delta$ for the graded $B$-linear
contraction of $M$ displayed above. Extend it by the signed
Leibniz and special divided-power rules. These rules preserve the exterior relations: for
example $\iota(e_j^2)=\alpha_je_j-e_j\alpha_j=0$.
They preserve $\sum e_j=0$ because $\sum\alpha_j=0$.
They also preserve every coefficient, and commute with all
simplicial operators, which are the same fibre sums on the
four types of generator. On generators and divided-power
monomials one has
\[
                           d\iota+\iota d=\delta.          \tag{136.PL6}
\]
Both sides are degree-zero special derivations, so this
identity holds on the whole algebra.

For signs, write the rational de Rham path algebra with the
interval factor first,
$K[z,dz]\otimes_K D_K$; the graded switch identifies it
with the path algebra used in AQ-136-SEMILOCAL. Define
\[
 H:B_K\longrightarrow K[z,dz]\otimes_K D_K,\qquad
 H(b)=1\otimes(b,0)+z\otimes(0,\delta b)
                         +dz\otimes(0,\iota b).           \tag{136.PL7}
\]
This formula is homogeneous in form degree and weight.
The square-zero property of $M$ makes it multiplicative:
the $z$ term is the ordinary derivation rule, while the
$dz$ term is the signed rule for $\iota$, together with
$(-1)^{|a|}$ when $1\otimes a$ crosses $dz$.
The unit and coefficients are fixed. Divided powers are
preserved by the same first-order square-zero formula
and the special derivation identities; equivalently, over
$K$ they are uniquely the operations $v\mapsto v^r/r!$.

It is a chain map. For homogeneous $b$, its module component
has differential
\[
 z\otimes\delta(db)+dz\otimes(\delta b-d\iota b)
          =z\otimes\delta(db)+dz\otimes\iota(db),
\]
by (136.PL6), which is the module component of $H(db)$.
The formula is simplicially natural by the generator check
above. Finally its endpoints are $j_{0,K}$ and $j_{1,K}$.

We now obtain an arrow over the original coefficient ring:
\[
 F:B\longrightarrow\mathcal E(D),\qquad
 b\longmapsto
 \bigl((j_0(b_i))_{i=1}^s,\ j_1(b_K),\ (H(b_K))_{i=1}^s\bigr),
                                                               \tag{136.PL8}
\]
where $b_i$ and $b_K$ are the scalar extensions of $b$.
The endpoints in (136.PL7) are exactly the compatibility
conditions defining $\mathcal E(D)$. Thus (136.PL8) is a
strict morphism in the original category.

Suppose (136.PL1) held for this arrow. Every path component
of $Fh=\eta_Du$ would be constant. Evaluating either endpoint
therefore gives
\[
                  j_{0,K}h_K=j_{1,K}h_K,
                    \qquad\text{hence}\quad\delta h_K=0.  \tag{136.PL9}
\]
The kernel of $\delta:B_K\to M_K$ is exactly $S_K$ in every
simplex. Indeed, eliminate $x_0,e_0$ using the coefficient
and sum relations. In the resulting basis
$x_1^{[a_1]}\cdots x_n^{[a_n]}e_I$ over $S_K$, the
$\beta_j$ coefficient of $\delta$ first eliminates every
monomial containing $e_j$. Then its $\alpha_j$ coefficient,
which lowers $a_j>0$ by one with coefficient one, eliminates
every positive $a_j$. The remaining terms are precisely
$S_K$. This argument uses freeness of the displayed basis,
not any field hypothesis on $K$.

Consequently $h_K:A'_K\to B_K$ factors through the constant
simplicial algebra $S_K$. The factorization is an $S_K$-map,
so it retracts $S_K\to A'_K$. In each weight it makes the
nonzero constant module $KT^{[w]}$ a direct summand of the
zero-form row of $A'_K$. That row cannot be acyclic in the
weight tail. On the other hand $A'_K$ is admissible: flat
localization preserves both uniform vertical cohomology
and all normalized horizontal homology. This contradiction
proves failure of (136.PL1).
\end{proof}

Over a nonzero $\mathbb Q$-algebra, common sources and ordinary
terminal localization were separately proved in AQ-136-RATIONAL.
Thus this counterexample cannot be promoted to a negative answer
to either original question. It rules out the general source-lifting
shortcut for the actual natural arithmetic path arrow, including
its mixed semilocal instances; a more restricted lifting assertion
would require a separate proof and narrower hypotheses.

\paragraph{AQ-136-RELATIVE-H1: a relative integral repair through vertical degree one.}
This is a positive partial construction in the original strict PD
category. It does not assert the higher vertical vanishings required
for a fully admissible source.

Let $k$ be a nonzero commutative ring which is torsion-free as a
$\mathbb Z$-module, and put $K=k\otimes_{\mathbb Z}\mathbb Q$.
Let $C_1,\ldots,C_t\in\mathcal A_k$, where $t\ge1$.
Suppose $A$ is a simplicial nonnegatively bigraded strict PD
differential $S_k$-algebra, with maps $a_j:A\to C_j$, such that
\[
 A_0=S_k,\qquad A_n^{0,*}=k[0],\qquad
 H^0((A_K)_n^{w,*})=S_K^w,\quad
 H^1((A_K)_n^{w,*})=0\quad(w\ge W).                       \tag{136.RH1}
\]
The coefficient identification in (136.RH1) is canonical and
holds uniformly in $n$. Enlarge $W\ge1$ so that for every target
and every $n$ one also has
\[
 H^0((C_j)_n^{w,*})=S_k^w,
       \qquad H^1((C_j)_n^{w,*})=0\quad(w\ge W).           \tag{136.RH2}
\]
There is a sequence of strict maps under $A$ and over all targets
whose filtered colimit $A^+$ satisfies
\[
 \begin{gathered}
 A^+_0=S_k,\qquad (A^+)_K\simeq A_K
                   \quad\text{as actual differential PD algebras},\\
 H^0((A^+)_n^{w,*})=S_k^w,\qquad
 H^1((A^+)_n^{w,*})=0\quad(w\ge W,\ n\ge0).
 \end{gathered}                                             \tag{136.RH3}
\]
No horizontal acyclicity is claimed for this repair by itself.
The same single bound $W$ works throughout and in the colimit.
The linear inputs are the Dold--Kan equivalence
(Stacks Project, Theorem~14.24.3, Tag~019G) and the free PD
polynomial universal property (Lemma~23.5.1, Tag~07H5).
The arbitrary-module extension of the latter is constructed
explicitly below; no flatness hypothesis is imported from it.

\begin{proof}
We use normalization $N_r$ in the simplicial direction, with
horizontal differential $\partial:N_r\to N_{r-1}$, and retain
$d$ for the form differential. Normalization is exact and its
individual terms commute with vertical cohomology: each is a
natural direct summand of the corresponding simplicial term.
Thus (136.RH2) gives, for every $r\ge1$ and $w\ge W$,
\[
 H^0(N_r C_j^{w,*})=H^1(N_r C_j^{w,*})=0.
                                                               \tag{136.RH4}
\]
In particular every vertical one-cycle in $N_rC_j$ has a
\emph{unique} zero-form primitive. The corresponding rational
statement holds for $N_r A_K$ by (136.RH1).

We first define one repair step $A\to R(A)$.
In each weight $w\ge W$ define a nonnegative horizontal chain
complex $V^w$ by
\[
 V_0^w=S_k^w,\qquad
 V_r^w=Z^1(N_r A^{w,*})\quad(r\ge1).                       \tag{136.RH5}
\]
For $r\ge2$ its differential is the existing $\partial$ on
one-cycles. It remains to define $\partial:V_1^w\to V_0^w$.
If $u\in Z^1(N_1 A^{w,*})$, rational exactness and the absence
of normalized degree-one vertical $H^0$ give a unique
\[
 v_K\in N_1 A_K^{w,0},\qquad d v_K=u_K.
\]
The image $a_j(u)$ has, by (136.RH4), a unique integral primitive
$y_j\in N_1 C_j^{w,0}$. After scalar extension its value is
$a_{j,K}(v_K)$, by rational uniqueness. Since $A_0^{w,1}=0$,
\[
             d(\partial y_j)=a_j(\partial u)=0.
\]
Consequently $\partial y_j$ is the image of a unique element
$s_j\in S_k^w$, by (136.RH2) and nonnegative form degrees.
Its rationalization is the element $\partial v_K$ of
$A_{K,0}^{w,0}=S_K^w$. The injection $S_k^w\hookrightarrow S_K^w$
therefore makes all $s_j$ equal. Define
\[
                         \partial u=s_j\quad\text{in }V_0^w.
                                                               \tag{136.RH6}
\]
It is $k$-linear: rational primitives and their boundaries are
linear, and the coefficient injection is injective.
For $u\in V_2^w$, the rational primitive of $\partial u$ is
$\partial$ of the primitive of $u$. Its next boundary is zero.
The same injection proves $\partial^2u=0$ integrally.
Thus (136.RH5)--(136.RH6) is a chain complex.

For $w<W$ instead put $V^w=N(A^{w,0})$ unchanged.
There are horizontal chain maps
\[
 f:N(A^{*,0})\longrightarrow V,
             \qquad d_V:V\longrightarrow N(A^{*,1}),      \tag{136.RH7}
\]
where in high weights $f_0$ is the coefficient identity,
$f_r(a)=da$ for $r\ge1$, $(d_V)_0=0$, and $(d_V)_r(u)=u$
for $r\ge1$. In low weights use the identity and the old
vertical differential. The only boundary identity requiring
comment is $r=1$: the rational primitive of $da$ is $a_K$,
so (136.RH6) gives precisely $\partial a\in S_k^w$.
Thus
\[
                         d_Vf=d_A,\qquad d_A d_V=0.
                                                               \tag{136.RH8}
\]

Each target map on zero forms extends to a chain map
$g_j:V\to N(C_j^{*,0})$. In high normalized degrees $r\ge1$,
send $u$ to its unique primitive from (136.RH4); in normalized
degree zero use coefficients. For $r\ge2$ compatibility with
$\partial$ follows from uniqueness of the primitive in degree
$r-1$, and for $r=1$ it is the definition (136.RH6).
The old image of $a\in N_rA^{w,0}$ is itself the primitive
of $a_j(da)$, so
\[
                 g_j f=N(a_j^0),\qquad
                   d_{C_j}g_j=N(a_j^1)d_V.                \tag{136.RH9}
\]
Dold--Kan now gives actual simplicial modules and maps with
all these identities. We continue to denote them by $V,f,d_V,g_j$.

In positive weights take the graded PD algebra pushout
\[
 R(A)=A\amalg_{\Gamma_k(A^{*,0}_{>0})}
                              \Gamma_k(V_{>0}).           \tag{136.RH10}
\]
This formula first denotes a pushout of graded algebras;
the middle algebra is not being given a zero differential
and mapped to $A$ as a differential algebra. Instead define
the differential on the new zero-form generators by
$d[v]=d_V(v)$, and keep the old differential on $A$.
The identities (136.RH8) make every identifying relation
$[f(a)]-a$ stable under this differential. Its PD ideal is
also stable by the special derivation rule. The differential
therefore descends to (136.RH10). Its square vanishes on
old elements and new generators; on divided powers it also
vanishes, using the square-zero rule for the odd element
$d[v]$. This gives a strict PD differential algebra with
the original nonnegative form grading. All constructions
are simplicial and preserve weights and the coefficient map.
The maps (136.RH9) and the free PD universal property give
strict maps $R(A)\to C_j$ extending all $a_j$.
The construction can equivalently be made as the free graded
PD extension by $V_{>0}$ followed by its displayed identifying
PD ideal, so it requires no flatness of $V$ or of a target module.
For an arbitrary module the free PD algebra itself is obtained
from a PD polynomial algebra on a set of module generators by
imposing the PD ideal of all linear relations. This also proves
its universal property and compatibility with scalar extension.

In simplicial degree zero $V_0=S_k=A_0$, so (136.RH10)
changes nothing there. Weight zero also remains $k[0]$.
After extension to $K$, the map $f$ is an isomorphism.
For high weights and $r\ge1$, this is exactly the bijection
\[
 d:N_rA_K^{w,0}\xrightarrow{\sim} Z^1(N_rA_K^{w,*})
\]
from (136.RH1); localization commutes with cycles.
In other weights and normalized degree zero it is already
an identity. Dold--Kan and the free PD universal property,
which commutes with scalar extension, now identify
\[
                           R(A)_K\simeq A_K               \tag{136.RH11}
\]
as differential PD algebras, compatibly with the transition.

The step kills every previous vertical one-cycle in high
weights: in positive normalized degree $r$ the new generator
$[u]$ has differential $u$. In normalized degree zero there
are no one-forms because $A_0=S_k$. It also kills every
previous normalized vertical zero-cycle in degree $r\ge1$,
since its image under $f$ is $da=0$. Consequently the image
of $H^0(A_n^{w,*})\to H^0(R(A)_n^{w,*})$ consists of
coefficients, and the map on $H^1$ is zero, for $w\ge W$.
The normalized direct-sum decomposition and its degeneracies
justify these statements in every simplicial degree, not only
on normalized terms.

Iterate $R$ countably many times and take the filtered colimit.
Every step still satisfies (136.RH1), with the same $W$, by
(136.RH11), and retains $A_0=S_k$ and the target maps.
Filtered colimits of modules are exact and commute with the
finite kernel/cokernel calculations for fixed vertical
cohomology. The same is true of the underlying modules of
these algebras: each operation and each defining identity
has finite arity, including each fixed divided power.
The transition kills every $H^1$ class at the next step and
sends every $H^0$ class to a coefficient. Coefficients remain
injective, since simplicial degeneracies from degree zero have
face retractions and degree zero is still exactly $S_k$.
This proves all assertions of (136.RH3), including the single
uniform bound in all simplices. The rationalized transitions
are isomorphisms compatible with $A_K$, so their colimit is
that actual algebra, as stated.
\end{proof}

\paragraph{AQ-136-PARTIAL-SOURCE: all horizontal rows and the first two vertical conditions.}
For the same torsion-free $k$ and any finite nonempty family
$C_1,\ldots,C_t\in\mathcal A_k$, there is a simplicial strict
PD differential $S_k$-algebra $A$ with strict maps to all $C_j$
such that:
\[
 \begin{gathered}
 A_0=S_k,\qquad A_n^{0,*}=k[0],\qquad H^*(A_{K,n})=S_K
                                   \text{ in every weight},\\
 H^0(A_n^{w,*})=S_k^w,\quad H^1(A_n^{w,*})=0
                        \quad(w\ge W,\ n\ge0),\\
 H_r(N(A_*^{w,q}))=0
                   \quad(w\ge\mu(q),\ r\ge0,\ q\ge0).
 \end{gathered}                                             \tag{136.RH12}
\]
The bounds depend on the given targets, not on the construction
stage or simplex. All positive vertical cohomology groups are
integer-torsion; no vanishing in vertical degree $q\ge2$ is
asserted. In particular (136.RH12) is not yet an object of
$\mathcal A_k$ merely from the conclusions proved here.

\begin{proof}
Choose $W$ from (136.RH2) for the targets. For every form degree
$q$ choose $\mu(q)\ge1$ so that all their normalized horizontal
rows are acyclic at weights $w\ge\mu(q)$.
Begin with the constant coefficient algebra $A^{(0)}=S_k$
and its given maps to all targets. We describe a horizontal
cell step $A\to H(A)$ preserving $A_0=S_k$, weight zero,
and the rational vertical comparison in every weight.

Let $z\in N_r A^{w,q}$ be a horizontal cycle, where
$w\ge\mu(q)$. A free vertical disk $E_q$ has generators
$a,b$ of form degrees $q,q+1$ and $da=b$, $db=0$.
Give them both weight $w$. Let $\mathbb S^r$ be the
simplicial module corresponding to $k[r]$ and let
$\mathbb D^{r+1}$ correspond to the horizontal disk
$k[r+1]\xrightarrow{1}k[r]$. The boundary inclusion
$\mathbb S^r\hookrightarrow\mathbb D^{r+1}$ is split in
each simplicial module degree. The pair $z,dz$ specifies a
map from the free graded PD differential algebra on
$\mathbb S^r\otimes E_q$ to $A$. Attach the free algebra on
$\mathbb D^{r+1}\otimes E_q$ by a pushout along this map.
In normalization the new generator $a$ has
$\partial a=z$, and its partner satisfies
$\partial b=dz$.

For each target the image of $z$ is a boundary in its form-$q$
row. Choose a normalized primitive $y_j$ with
$\partial y_j=a_j(z)$, and send the new generators to
$y_j,dy_j$. This extends the target maps as strict simplicial
PD differential maps. The free universal property checks all
simplicial identities simultaneously. Attach such cells for
all cycles and all triples $(r,w,q)$ satisfying the bound.
These cells form a set, and a simultaneous coproduct and
pushout therefore exist.

No new generator occurs in simplicial degree zero: for $r=0$
the boundary inclusion is already an isomorphism there, and
for $r>0$ both degree-zero modules vanish. Thus $A_0=S_k$
remains exact. Positive weights of the cells retain weight zero.
At a fixed simplex the boundary-module inclusion splits as a
module, and the splitting tensored with $E_q$ respects its
vertical differential. The attached algebra is consequently
an extension of $A_n$ by free vertical PD disks, with their
new generator pairs independent over $A_n$.
After extension to $K$, each such disk is contractible onto
its constants. One explicit proof uses the degree-minus-one
special derivation sending $b$ to $a$ and $a$ to zero:
its commutator with $d$ is multiplication by the positive
monomial arity. Divide by that arity in the $\mathbb Q$-algebra
$K$ to obtain a contraction on the augmentation ideal.
The contraction preserves weight. Tensoring with $A_n$
and taking filtered colimits over finite subsets of the new
generators proves that the attachment preserves rational
vertical cohomology in every weight and simplex.
By construction the transition induces zero on every old
horizontal homology group in the specified tails.

Now alternate one horizontal step $H$ with one repair step
$R$ of the preceding theorem, always using the same $\mu(q)$
and $W$, and take the countable filtered colimit. Every stage
still has simplex zero $S_k$ and exact rational vertical
cohomology, so the repair is always applicable. Every old
horizontal class in the weight tail is killed at the next
horizontal step and remains zero afterwards. Every old
vertical $H^1$ class is killed at the next repair, and every
old vertical $H^0$ class maps to a coefficient there.
Exactness of filtered colimits gives all the assertions of
(136.RH12), with the original fixed bounds.
Finally scalar extension to $K$ commutes with cohomology and
has zero positive vertical cohomology. Its kernel on any
underlying abelian group consists exactly of elements killed
by some nonzero integer. Thus all remaining positive vertical
cohomology is integer-torsion. No common annihilator or higher
vanishing has been deduced.
\end{proof}


\paragraph{AQ-136-HIGHER-COMPLETION: the positive partial axioms do not ensure relative completion.}
The partial source of AQ-136-PARTIAL-SOURCE is not claimed to
admit an arbitrary further relative repair. The following example
shows that such a repair theorem would be false, even without
requiring the completed source to have simplex zero exactly $S$.
It does not refute the original common-source conjecture.

Fix a rational prime $p$, and work over $k=\mathbb Z$.
There is a simplicial strict PD differential $S$-algebra $A$
and two maps $f_0,f_1:A\to B$ to the standard admissible
theory, with the following properties:
\[
 \begin{gathered}
 A_0=S,\qquad A_n^{0,*}=\mathbb Z[0],\\
 H^0(A_n^{w,*})=S^w,\qquad H^1(A_n^{w,*})=0
                                      \quad\text{for all }n,w,\\
 H^*(A_{\mathbb Q,n})=S_{\mathbb Q},\qquad
 H_r(N A_*^{w,q})=0\quad(w>q,\ r\ge0).
 \end{gathered}                                            \tag{136.HC1}
\]
Nevertheless there are no admissible $R$ and strict maps
\[
            A\xrightarrow{g}R
                    \mathrel{\substack{\xrightarrow{j_0}\\[-3pt]
                                      \xrightarrow[j_1]{}}}B,
                 \qquad j_i g=f_i.                       \tag{136.HC2}
\]
Both targets are the same object $B$, so they certainly have
a common admissible source. The obstruction is precisely to
extending the specified maps from this specified partial source.

\begin{proof}
Let $L_n=\mathbb Z^{n+1}$ be the standard coordinate simplicial
module, with coefficient inclusion $\mathbb Z\to L_n$ given
by the sum of the coordinate generators. Put
\[
             Q=L/\mathbb Z,\qquad NQ=\mathbb Z[1].
\]
These are the modules already used in the construction of $B$.
Let $V$ and $U$ be copies of $Q$, placed in form degrees one
and two respectively, and give both weight one. Define
\[
 D=\Lambda_{\mathbb Z}(V)\otimes_{\mathbb Z}\Gamma_{\mathbb Z}(U),
        \qquad d v=p u,\quad d u=0,\qquad
 A=B\otimes_{\mathbb Z}D.                                 \tag{136.HC3}
\]
The correspondence $v\mapsto u$ is the identity between the
two copies of $Q$. Thus the differential is simplicial and
respects every linear relation. Extend it by the strict
graded derivation and special PD rules. Its square is zero;
the odd exterior relations are stable: $d(v^2)=p(uv-vu)=0$,
and the same derivation rule preserves $vv'+v'v=0$.
The usual free graded PD algebra structure supplies divided
powers on positive even elements. The coefficient map to $A$
is its map to the $B$ factor. This verifies the actual strict
PD structure, not merely its underlying bicomplex.

In simplicial degree $n$, the module $Q_n$ is free of rank
$n$. Hence $D_n^0=\mathbb Z$, and
\[
     d:D_n^1=\bigoplus_{i=1}^n\mathbb Z v_i
             \longrightarrow D_n^2
\]
sends the basis $v_i$ to the linearly independent elements
$p u_i$. It is injective. Consequently
\[
                       H^0(D_n)=\mathbb Z,
                              \qquad H^1(D_n)=0.           \tag{136.HC4}
\]
The standard integral PD de Rham contraction of $B_n$ onto
$S$ is $S$-linear and weight-preserving: after solving for
the zeroth coordinate, it is the tensor product of the
one-variable contractions
$x^{[m-1]}dx\mapsto x^{[m]}$ for $m\ge1$.
Tensoring that homotopy with $D_n$ identifies the cohomology
of $A_n$ with that of $S\otimes D_n$.
The module $S$ is free over $\mathbb Z$, so (136.HC4) proves
the asserted integral $H^0$ and $H^1$ in every weight.
Also $Q_0=0$, so $D_0=\mathbb Z$ and $A_0=B_0=S$.
Positive form generators have positive weight, which proves
the weight-zero assertion.

Over $\mathbb Q$ the algebra on each pair $v_i,u_i$ contracts
onto its constants. Indeed
\[
             d(v_i u_i^{[m]})=p(m+1)u_i^{[m+1]},
\]
and $p(m+1)$ is invertible. The tensor product of these
contractions proves the rational vertical assertion.
In particular all positive integral vertical cohomology of
$A_n$ is integer-torsion, as in the partial-source theorem.

For the horizontal assertion, a monomial of form degree $d$
in $D$ has weight $e$ satisfying
\[
                          \lceil d/2\rceil\le e\le d.
                                                               \tag{136.HC5}
\]
There are only finitely many possible $e$ for fixed $d$.
Every summand of $A^{w,q}$ therefore has the form
\[
                 B^{w-e,q-d}\otimes D^{e,d}.
\]
If $w>q$, then $w-e>q-d$. The standard row contraction
of $B^{a,b}$ for $a>b$ applies: its positive coordinate
PD degree is $a-b>0$, and the linear contraction of $L$
induces the contraction of that PD power. Tensoring in
the simplicial direction preserves this acyclicity.
Equivalently, the normalized Eilenberg--Zilber comparison
identifies the diagonal tensor with the tensor of normalized
complexes, whose first factor is contractible. All modules
here are degreewise free; the same flat comparison used in
AQ-136-TENSOR applies (Stacks Project, Remark~21.39.11,
Tag~08QD). Summing the finitely many summands proves the
last assertion of (136.HC1), including horizontal $H_0$.

The standard theory contains the canonical copy of $Q$ as
its closed one-form generators. Define $f_0$ and $f_1$ to
be the identity on $B$, and on $D$ by
\[
       f_0(v)=f_0(u)=0,\qquad f_1(v)=e,
                                  \qquad f_1(u)=0.        \tag{136.HC6}
\]
Here $e$ denotes that canonical closed coordinate
one-form, not a globally chosen primitive on a simplicial
circle. More concretely its normalized generator in
simplicial degree one is $dx$, with coordinate $x$ having
endpoints $0,T$. The maps of copies of $Q$ are simplicial;
$d(dx)=0$ verifies compatibility with $dv=pu$.
The free exterior/PD universal property proves that
(136.HC6) defines strict maps in the stated category.

Let $X=\Delta[1]/\partial\Delta[1]$ be the simplicial circle.
The normalized generators $v$ and $u$ of the copies of $Q$
have both faces zero, so they define elements of $A(X)$,
still denoted $v,u$, satisfying $dv=pu$.
For every $w\ge1$ put
\[
             v_w=T^{[w-1]}v,\qquad
             u_w=T^{[w-1]}u,\qquad dv_w=p u_w.             \tag{136.HC7}
\]
On the edge, the difference $f_1(v_w)-f_0(v_w)$ is the
closed one-form $T^{[w-1]}dx$.
Its class in $H^1(B^{w,*}(X))=S^w$ is
\[
                  T^{[w-1]}T=wT^{[w]}.                    \tag{136.HC8}
\]
For clarity, this normalization can be checked without a
comparison convention: the integral PD antiderivative on
the edge sends $x^{[m-1]}dx$ to $x^{[m]}$.
The difference between its endpoint values identifies
$H^1(B^{w,*}(X))$ with $S^w$; exact circle forms have zero
endpoint difference, and $x^{[w-1]}dx$ has difference
$T^{[w]}$, so the map is an isomorphism. Formula (136.HC8)
follows from the divided-power multiplication rule.

Suppose (136.HC2) existed. The restriction $b:B\to R$ of
$g$ satisfies $j_0b=j_1b=1_B$. The natural comparison theorem
AQ-136-COMPARISON, applied through form degree two, gives,
for every sufficiently large weight $w$,
\[
 H^1(R^{w,*}(X))\simeq S^w,\qquad
 H^2(R^{w,*}(X))=0,\qquad
 H^1(b)\text{ is an isomorphism}.                         \tag{136.HC9}
\]
Thus $H^1(j_0)=H^1(j_1)$, both being the inverse of $H^1(b)$.
Choose such a $w$ not divisible by $p$. Since $u_w$ is
closed, the middle assertion of (136.HC9) provides
$a\in R^{w,1}(X)$ with $da=g(u_w)$. Then
\[
                           c=pa-g(v_w)
\]
is closed. Its two images have the same cohomology class.
The elements $j_i(a)$ are closed because $f_i(u_w)=0$.
Using (136.HC8), the equality of the two classes says
\[
        p\bigl([j_1(a)]-[j_0(a)]\bigr)=wT^{[w]}
                                      \quad\text{in }S^w.
                                                               \tag{136.HC10}
\]
But $S^w=\mathbb ZT^{[w]}$, and $p\nmid w$, which is a
contradiction. Arbitrarily large such weights exist, so
no choice of admissibility cutoff avoids it.
This proves the exact relative-completion obstruction.
\end{proof}

\paragraph{Consequence for the next construction step.}
The hypotheses of the positive partial-source theorem alone
cannot justify a functorial repair killing all higher
vertical torsion while retaining arbitrary prescribed target
maps. Such a construction would have to maintain an additional
compatibility invariant that excludes (136.HC3)--(136.HC6),
or replace the chosen partial source or its maps. No negative
answer to either original integral question is deduced here.


\paragraph{AQ-136-COSPAN: admissible path pullbacks give common sources.}
The supplied SA08 intake reports corresponding statements
A-002 and A-003, but their prior full proofs were not supplied.
The following is a complete proof in the exact category
$\mathcal A_k$ defined from the recovered original sheets.
It uses the already proved coefficient-preserving simplicial
cotensor, rather than a strict equalizer of two composites.
No novelty claim is made.

Let $k$ be any commutative ring. For every cospan of admissible
strict PD theories
\[
                         A\xrightarrow{f}C
                                      \xleftarrow{g}B
                                                               \tag{136.CS1}
\]
there is an admissible $P$ with strict maps $P\to A$ and
$P\to B$. The two resulting maps $P\to C$ are joined by
an actual simplicial cotensor homotopy. Equality of those
two maps is not asserted.

\begin{proof}
We first record the extension property used below. If a
simplicial $k$-module $F$ has acyclic normalized complex,
including $H_0=0$, then for every $n$ the restriction
\[
 \operatorname{Hom}_{s\mathrm{Mod}_k}
       (k[\Delta[n]\times\Delta[1]],F)
   \longrightarrow
 \operatorname{Hom}_{s\mathrm{Mod}_k}
       (k[\Delta[n]\times\partial\Delta[1]],F)
                                                               \tag{136.CS2}
\]
is surjective. Under Dold--Kan, the inclusion of the two
free simplicial modules normalizes to a degreewise split
inclusion of free nonnegative chain complexes. Its quotient
has a basis of relative nondegenerate simplices.
Given the chain map on the subcomplex, extend it in
increasing chain degree. For a new basis element of degree
$r\ge1$, the already assigned image of its boundary is a
cycle in $N_{r-1}F$. It is a boundary by acyclicity, so
choose its primitive in $N_rF$ as the new image. New
degree-zero basis elements, if any, can be assigned zero.
The boundary-square-zero identity verifies the cycle
condition at each step. This constructs a chain extension;
the Dold--Kan equivalence makes it the required simplicial
extension. In the displayed prism all vertices already
belong to the subcomplex. There is no dimension-dependent
condition on $F$: its one acyclicity hypothesis works for
every $n$. These linear facts follow also directly from
Stacks Project, Theorem~14.24.3, Tag~019G.

Let $P_s(C)$ be the cotensor of AQ-136-COTENSOR:
\[
                P_s(C)_n=C(\Delta[n]\times\Delta[1]).
\]
It is admissible over every $k$, with coefficient map,
endpoint maps $e_0,e_1:P_s(C)\to C$ and constant-path map
$C\to P_s(C)$ all strict. Form the actual limit of strict
PD differential algebras
\[
 \begin{gathered}
        P=(A\times B)\times_{C\times C}P_s(C),\\
 (a,b,h)\in P\quad\Longleftrightarrow\quad
                   f(a)=e_0(h),\quad g(b)=e_1(h).
 \end{gathered}                                            \tag{136.CS3}
\]
Its coefficient map sends $s$ to $(s,s,s)$ with constant
path in the third position. Limits carry the operations
componentwise, including divided powers on all positive
even elements; hence this is an actual simplicial strict
PD differential $S_k$-algebra. It has the required strict
projections and the path $h$ between their two composites.
It remains to verify admissibility with uniform cutoffs.

Fix a form degree $q$, and choose $w\ge\mu_C(q)$ so that
$F=C_*^{w,q}$ has acyclic normalization. Formula (136.CS2)
and the two disjoint copies of $\Delta[n]$ identify its
surjective restriction map with
\[
       (e_0,e_1):P_s(C)_n^{w,q}
                           \longrightarrow C_n^{w,q}\oplus C_n^{w,q}.
                                                               \tag{136.CS4}
\]
Consequently, in these weights, the defining kernel of
(136.CS3) gives an exact sequence of simplicial modules
\[
 \begin{gathered}
 0\longrightarrow P_*^{w,q}\longrightarrow
       A_*^{w,q}\oplus B_*^{w,q}\oplus P_s(C)_*^{w,q}
       \xrightarrow{\delta}
       C_*^{w,q}\oplus C_*^{w,q}\longrightarrow0,\\
       \delta(a,b,h)=(f(a)-e_0(h),\,g(b)-e_1(h)).
 \end{gathered}                                            \tag{136.CS5}
\]
The minus signs here describe a map of underlying modules;
the algebra $P$ itself was defined by a limit of algebra
maps in (136.CS3). There is no assertion that $\delta$ is
an algebra homomorphism.

For the horizontal condition, enlarge $w$ so that the
form-$q$ rows of $A,B,C,P_s(C)$ are all acyclic.
Normalization is exact, so the long exact homology sequence
of (136.CS5) proves that the row of $P$ is acyclic, including
degree zero. Its cutoff is independent of the simplicial
degree and of the horizontal homology degree. This proves
the original horizontal condition for each $q$.

For the vertical condition, fix $d\ge1$. Choose one $W_d$
so large that for every $w\ge W_d$ and every simplex $n$,
the vertical cohomology of $A_n,B_n,C_n,P_s(C)_n$ through
degree $d+1$ is canonically $S_k^w$ in degree zero and zero
in positive degrees. Such a uniform bound exists by
admissibility and the already proved cotensor comparison.
Also require $w\ge\mu_C(q)$ for all $0\le q\le d+1$.
The finite maximum of these bounds still works uniformly
in $n$.

Write $D=A_n^{w,*}\oplus B_n^{w,*}\oplus P_s(C)_n^{w,*}$,
$E=C_n^{w,*}\oplus C_n^{w,*}$, and $I=\operatorname{im}\delta$.
The sequence $0\to P_n^{w,*}\to D\to I\to0$ is exact
in all degrees, and $I\to E$ is an isomorphism in degrees
zero through $d+1$ by (136.CS4). Thus it induces an
isomorphism on cohomology through degree $d$.
In degree zero the map $H^0(D)\to H^0(I)=H^0(E)$ is
exactly
\[
             (S_k^w)^3\longrightarrow(S_k^w)^2,
                       \qquad (a,b,h)\longmapsto(a-h,b-h).
                                                               \tag{136.CS6}
\]
It is surjective and has diagonal kernel $S_k^w$.
All positive cohomology of $D$ and $I$ through degree $d$
is zero. Their long exact sequence therefore gives
\[
           H^0(P_n^{w,*})=S_k^w,\qquad
           H^q(P_n^{w,*})=0\quad(1\le q\le d).            \tag{136.CS7}
\]
The first identification is the original coefficient map.
Since the same $W_d$ works in every simplex, these are
precisely the uniform vertical tail conditions.
Together with the horizontal argument they prove
$P\in\mathcal A_k$.
\end{proof}

\paragraph{AQ-136-CONNECTED-SOURCES: the exact connectivity criterion.}
Over every commutative ring $k$, two objects of $\mathcal A_k$
admit a common admissible source if and only if they belong
to the same connected component of its underlying category,
or equivalently of its ordinary all-arrow localization.
Every finite family in one such component has a common
admissible source.

\begin{proof}
Let $A\sim B$ mean that there is an admissible $D$ with
maps $D\to A,D\to B$. This relation is reflexive and
symmetric. If $D$ is a common source of $A,B$ and $E$ a
common source of $B,C$, apply AQ-136-COSPAN to the cospan
$D\to B\leftarrow E$. Its common source $P$ maps to
$A$ through $D$ and to $C$ through $E$. Thus $\sim$ is
transitive. Every arrow $A\to B$ gives $A\sim B$, using
$A$ itself. Hence any finite zigzag gives a common source.
Conversely a common source supplies the zigzag
$A\leftarrow D\to B$. This proves the equivalence.
For a finite family, join a common source for the first
objects to the next object by this equivalence, and compose
its maps with the earlier source maps. Induction gives
one object mapping strictly to every member of the family.
The zero ring is harmless throughout.
\end{proof}

\paragraph{AQ-136-SEMILOCAL-SOURCES: strict common sources for the mixed semilocal case.}
Let $k$ be any commutative ring in which all but finitely
many rational primes are units. Every finite family in
$\mathcal A_k$ has a common admissible strict PD source.
In particular this includes $\mathbb Z_{(2,3)}$ and
every $\mathbb Z_{(P)}$-algebra for a finite set $P$ of
rational primes, with no flatness, reducedness or noetherian
hypothesis.

\begin{proof}
AQ-136-SEMILOCAL proves that the ordinary all-arrow
localization of this very category is terminal. It is
therefore connected. Apply AQ-136-CONNECTED-SOURCES.
The resulting maps to the original targets are strict
maps of admissible PD theories, even though the intermediate
cospan squares commute only by the specified cotensor
homotopies. No equality $Fh=\eta_Cu$ of the kind refuted
by AQ-136-PATH-LIFT has been used or asserted.
\end{proof}

Over $\mathbb Z$, the original common-source conjecture is
now precisely the connectedness question. The proof does
not establish that connectedness, or that every loop in
the integral all-arrow localization is trivial. Those two
integral assertions remain distinct. This result also does
not repair the prescribed partial-source maps excluded by
AQ-136-HIGHER-COMPLETION.


\paragraph{AQ-136-FIRST-SKELETON: a common free integral first skeleton by actual replacements.}
The supplied intake reports a common first-skeleton replacement
as A-004. Its prior full proof is unavailable here. We give a
complete construction using only the original uniform
admissibility tests. This is a reduction of the integral
comparison problem, not a proof of integral connectedness.

Let $k$ be any commutative ring and let
$C_1,\ldots,C_t\in\mathcal A_k$ be a finite nonempty family.
There is an integer $N\ge1$, one $1$-truncated simplicial strict
PD differential $S_k$-algebra $F$, and admissible theories
$C'_j$ with strict maps $C'_j\to C_j$, such that every $C'_j$
has exactly the same first skeleton $F$. Explicitly,
\[
 \begin{gathered}
 F_0=S_k,\qquad
 F_1=S_k\otimes_k
       \Gamma_k(x_m:m\ge N)\otimes_k\Lambda_k(e_m:m\ge N),\\
 |x_m|=(m,0),\quad |e_m|=(m,1),\quad dx_m=e_m,\quad de_m=0,\\
 d_0x_m=0,\quad d_1x_m=T^{[m]},\quad
 d_0e_m=d_1e_m=0,
 \end{gathered}                                             \tag{136.SK1}
\]
with the degeneracy $F_0\to F_1$ given by coefficients.
Interchanging the two endpoint names makes no difference.
The augmentation $S_k\to C'_{j,n}$ still has the original
uniform vertical cohomology comparison, and every horizontal
row is acyclic in its weight tail.

\begin{proof}
\emph{1. Choose the common first skeleton.}
Choose $N$ so that every zero-form row $(C_j)_*^{m,0}$ is
acyclic whenever $m\ge N$. In normalized degree zero its
coefficient $T^{[m]}$ is therefore a boundary. Choose
\[
 y_{j,m}\in (C_j)_1^{m,0},\qquad
       d_0y_{j,m}=0,\quad d_1y_{j,m}=T^{[m]}.
                                                               \tag{136.SK2}
\]
One chooses the sign of a normalized primitive to obtain
these two endpoints, independently of the normalization's
boundary-sign convention. The assignments
$x_m\mapsto y_{j,m}$ and $e_m\mapsto dy_{j,m}$, together with
coefficients in degree zero, define a strict map from $F$
to the first truncation of $C_j$. The free PD and exterior
universal properties verify the algebra and differential
relations, and the two faces and sole degeneracy verify
all truncated simplicial identities.

Every vertical complex $F_1^{w,*}$ has coefficient cohomology
$S_k^w$ in degree zero and zero in higher degrees, in every
weight. This is the integral PD Poincare contraction of
independent even zero-form disks. It is valid for the
infinitely many displayed variables by taking the filtered
union over finite sets; at a fixed weight only finitely many
of their weights can occur. The contraction is $S_k$-linear
and weight-preserving.

\emph{2. Two facts about the first coskeleton.}
For any $1$-truncated simplicial algebra $E$, let
$\operatorname{cosk}_1E$ denote the right Kan extension to
all simplicial degrees, computed by the limits of its
vertex and edge data. Its value at $n$ consists of compatible
sections on the graph $G_n=\operatorname{sk}_1\Delta[n]$:
one value in $E_0$ at each vertex and one in $E_1$ at each
edge, with the specified endpoints. Degenerate edges add
no new variables. Operations are computed componentwise.
This is the finite-limit formula of the Stacks Project,
Lemmas~14.19.2--14.19.4 in Section~14.19, Tag~0AMA;
in particular the first truncation is unchanged.

First, if $M$ is a $1$-truncated simplicial module, then the
normalization of its coskeleton is
\[
 \cdots\longrightarrow0\longrightarrow
       \ker(\partial:N_1M\to M_0)
       \xrightarrow{\;\mathrm{incl}\;}N_1M
       \xrightarrow{\partial}M_0.
                                                               \tag{136.SK3}
\]
The displayed kernel is in degree two. One direct verification
uses triangle edge values: imposing the first two normalized
face conditions leaves exactly one edge value with both
endpoints zero, and the last face gives its inclusion in
$N_1M$. In degree at least three, every edge and vertex lies
in one of the normalized zero faces, so the normalized term
is zero. Thus the only possible horizontal homology of this
coskeleton is $H_0=\operatorname{coker}\partial$.

Second, suppose in a fixed weight the map
$E_1^q\to E_0^q\oplus E_0^q$ is surjective in every degree
needed for a cohomology calculation. Then sections on
$G_n$ are the kernel of the exact incidence sequence
\[
 0\longrightarrow(\operatorname{cosk}_1E)_n
 \longrightarrow E_0^{\,V(G_n)}\oplus E_1^{\,E(G_n)}
 \longrightarrow (E_0\oplus E_0)^{E(G_n)}
 \longrightarrow0.                                      \tag{136.SK4}
\]
The last map subtracts each edge's endpoint values from the
values at its two incident vertices. It is surjective by
choosing the two endpoint values independently on each edge.
If the cohomology of both $E_0,E_1$ is canonically $S_k^w$
in degree zero and zero above, the resulting cohomology of
the kernel is exactly
\[
                 H^r(G_n;k)\otimes_k S_k^w.                \tag{136.SK5}
\]
Indeed the degree-zero incidence complex in (136.SK4) is
the vertex-edge cochain complex of the graph after eliminating
one copy of the scalar for each edge. It has $H^0=S_k^w$
because $G_n$ is connected, its usual graph $H^1$, and no
higher cohomology. All identifications are natural in $E$.
If the hypotheses hold only through degree $d+1$, the same
image-complex truncation as in AQ-136-COSPAN proves (136.SK5)
through degree $d$. Finite products over a graph are finite
direct sums of modules, so no product exactness hypothesis is
needed. The cutoff does not depend on the number of vertices
or edges.

\emph{3. The two coskeleta have the same vertical tail comparison.}
Put $U=\operatorname{cosk}_1F$ and
$V_j=\operatorname{cosk}_1(\operatorname{sk}_1 C_j)$.
The maps of Step~1 induce $U\to V_j$.
For $F$, endpoint surjectivity holds in zero forms for all
$w\ge N$: the coefficient $T^{[w]}$ and the generator
$x_w$ give the two independent endpoint values. In positive
form degrees the endpoint target is zero, since $F_0=S_k$.
Step~1 and (136.SK4)--(136.SK5) therefore compute the vertical
cohomology of $U_n$ exactly as graph cohomology for $w\ge N$.

For $C_j$, acyclicity of its form-$q$ row gives the uniform
endpoint surjectivity $(C_j)_1^{w,q}\to
(C_j)_0^{w,q}\oplus(C_j)_0^{w,q}$ by (136.CS2), with $n=0$.
Its uniform vertical comparison in simplicial degrees zero
and one and the finite-degree version of (136.SK4) give the
same graph-cohomology answer for $V_{j,n}$, in any specified
finite range of form degrees, once $w$ is sufficiently large.
These bounds are independent of $n$.
The map $U\to V_j$ preserves coefficients and hence is the
identity on the incidence cohomology calculation. Consequently
\[
      H^q(U_n^{w,*})\xrightarrow{\sim} H^q(V_{j,n}^{w,*})
                                                               \tag{136.SK6}
\]
holds in every fixed finite range, uniformly in $n$, in a
weight tail. Neither coskeleton is claimed to be admissible:
its graph $H^1$ usually remains nonzero.

\emph{4. The actual replacements and their horizontal test.}
Define the strict PD pullback
\[
                 C'_j=C_j\times_{V_j}U,                    \tag{136.SK7}
\]
using the unit $C_j\to V_j$ of truncation and coskeleton.
For each fixed form degree $q$ and sufficiently large weight,
this unit is degreewise surjective as a map of simplicial
modules. To see this, a compatible family on $G_n$ is a
chain map from $N(k[G_n])$ into the acyclic normalized row
of $C_j$. Extend it across $\Delta[n]$ by the same increasing
degree filling argument as in (136.CS2). The relative
normalized chains have a free basis, and all relative cells
have degree at least two. This proves surjectivity uniformly
in $n$ from the one row cutoff.
Thus in those weights (136.SK7) gives an exact module sequence
\[
       0\longrightarrow (C'_j)_*^{w,q}
       \longrightarrow (C_j)_*^{w,q}\oplus U_*^{w,q}
       \longrightarrow (V_j)_*^{w,q}\longrightarrow0.
                                                               \tag{136.SK8}
\]

The rows of $U$ are acyclic in a uniform weight tail by
(136.SK3). In form degree zero the image of $\partial$
contains $kT^{[w]}$ for every $w\ge N$, by the generator
$x_w$; in positive form degree its degree-zero target is
already zero. Likewise the rows of $V_j$ are acyclic in
the same tails as those of $C_j$, since (136.SK3) has only
the original row's $H_0$. The normalized exact sequence of
(136.SK8) proves horizontal acyclicity for $C'_j$.

\emph{5. The vertical test and the exact first skeleton.}
Fix $d\ge1$. Choose $w$ large enough for (136.SK8) to be
degreewise exact through degree $d+1$, for the coefficient
comparison of $C_j$ through that range, and for (136.SK6)
through that range. The finite maximum of the required
cutoffs is independent of simplex dimension.
Let $I$ be the image of $C_j\oplus U\to V_j$ as a vertical
complex in a fixed simplex and weight. It agrees with $V_j$
through degree $d+1$, so cohomology through degree $d$ is
unchanged. In the long exact sequence of
$0\to C'_j\to C_j\oplus U\to I\to0$, the map contributed
by $U$ is the isomorphism (136.SK6). In degree zero the
matrix is the scalar difference $S_k^w\oplus S_k^w\to S_k^w$,
with diagonal kernel. It follows that $C'_j$ has coefficient
$H^0=S_k^w$ and zero $H^q$ for $1\le q\le d$.
The maps and coefficients in (136.SK7) are strict PD maps,
so this proves $C'_j\in\mathcal A_k$.

Finally the unit $C_j\to\operatorname{cosk}_1\operatorname{sk}_1C_j$
is an isomorphism in simplicial degrees zero and one.
The pullback in those degrees is therefore exactly $F$,
with its displayed faces, degeneracy, products, divided
powers and differential. Thus all the replacements have
the same first skeleton, and the projection in (136.SK7)
is the asserted admissible arrow to the original target.
\end{proof}

The theorem gives actual arrows, so replacing targets in
this way does not change their connected components in the
ordinary all-arrow localization. It does not assert that
arbitrary maps between the replacements fix their common
first skeleton, or that a later prescribed partial model
always has a completion. Higher-simplex compatibility,
including integral secondary periods, remains to be established.


\paragraph{AQ-136-FREE-ZERO-FORMS: a weight-one retraction still obstructs a full free PD de Rham source.}
The earlier AQ-136-DERHAM-CELLS excluded free zero-form
presentations whose new variables all have weight at least two.
The following refinement permits arbitrary variables of weight
one. Its additional hypothesis is automatic for a map to the
actual admissible truncated standard theory described below.

Let $k$ be a commutative ring, $S=\Gamma_k(T)$, and let $P_*$
be a simplicial graded PD $S$-algebra with homogeneous
presentations
\[
                  P_n=S\langle X_n\rangle_{\rm PD},
                \qquad \operatorname{wt}(x)\ge1.
                                                               \tag{136.FK1}
\]
Simplicial operators are arbitrary homogeneous PD $S$-maps.
Suppose its weight-one simplicial module has a coefficient
retraction
\[
          \rho:P_*^1\longrightarrow\underline{kT},
                            \qquad \rho(T)=T.              \tag{136.FK2}
\]
For every prime $p$ with $k/pk\ne0$, one has
\[
 [T^{[p]}]=0\text{ in }H_0(NP_*^p)
       \quad\Longrightarrow\quad
 H_1\!\left(N\Omega^1_{P_*/S,\mathrm{sp}}{}^{\,p}\right)\ne0.
                                                               \tag{136.FK3}
\]
Thus, if infinitely many rational primes are nonunits in $k$,
the full relative special-PD de Rham algebra of such a $P_*$
cannot be admissible. No flatness of Frobenius or of $k/pk$
is assumed.

\begin{proof}
Put $R=k/pk$ and $\bar P=P\otimes_kR$. As in the earlier
proof, form the weight-$p$ ordinary indecomposables and the
augmented special differentials
\[
 Q_*=(\bar P_+/(\bar P_+)^2)^p,\qquad
 F_*=\left(R\otimes_{\bar P_*}
          \Omega^1_{\bar P_*/\Gamma_R(T),\mathrm{sp}}\right)^p.
                                                               \tag{136.FK4}
\]
The augmentation sends all positive weights to zero.
These are simplicial $R$-modules. The special derivation,
followed by augmentation, induces a natural map $Q_*\to F_*$.

In every simplex the PD monomial basis gives the following
basis for $Q_n$:
\[
 T^{[p]},\qquad x^{[p]}\quad(\operatorname{wt}(x)=1),
       \qquad y\quad(\operatorname{wt}(y)=p).               \tag{136.FK5}
\]
Indeed a monomial involving at least two distinct base
variables is a product of two positive-weight monomials
with coefficient one. A pure divided power of a variable
of weight $a$ can have weight $p$ only when $(a,m)=(1,p)$
or $(a,m)=(p,1)$. Products of positive pure powers with total
exponent $p$ have the coefficient $\binom pi=0$ in $R$.
Consequently none of the pure vectors in (136.FK5) is killed
or related by the ordinary product ideal. This argument
works over an arbitrary ring $R$ of characteristic $p$.

The module $F_n$ is free on the $dy$ with
$\operatorname{wt}(y)=p$. The vector $y$ in (136.FK5) maps
to $dy$, while all the other displayed vectors map to zero:
$dT^{[p]}=0$ relatively, and
\[
                  d(x^{[p]})=x^{[p-1]}dx
\]
vanishes after the coefficient augmentation. Thus the kernel
$K_n$ is free on $T^{[p]}$ and the $x^{[p]}$ of weight-one
variables, and there is an exact sequence
\[
                  0\longrightarrow K_*
                      \longrightarrow Q_*
                      \longrightarrow F_*
                      \longrightarrow0.                  \tag{136.FK6}
\]
The point requiring care is the simplicial structure on
$K_*$ when faces mix the displayed weight-one variables.

Let $W_*=\bar P_*^1$ and let
$\mathrm{Fr}:R\to R$ denote $r\mapsto r^p$. Define
\[
              \mathrm{Fr}^*W_*=R\otimes_{\mathrm{Fr},R}W_*,
                                                               \tag{136.FK7}
\]
where the relation is
$a r^p\otimes w=a\otimes rw$.
The map
\[
                  w\longmapsto[\gamma_p(w)]\in Q_*
                                                               \tag{136.FK8}
\]
is additive and Frobenius-semilinear. Additivity follows
from the divided-power sum formula, since each mixed term
is a product of positive-weight elements. Semilinearity is
the identity $\gamma_p(rw)=r^p\gamma_p(w)$.
Consequently (136.FK8) induces a natural simplicial
$R$-linear map
\[
                    \mathrm{Fr}^*W_*\longrightarrow K_*.
                                                               \tag{136.FK9}
\]
The weight-one module $W_n$ is free on $T$ and the
weight-one variables. Formula (136.FK5) shows directly
that (136.FK9) takes a free basis to a free basis, so it is
an isomorphism. No preservation of exact sequences by
$\mathrm{Fr}^*$ is being used.

Reduce (136.FK2) modulo $p$ and apply this tensor functor
to the split pair
\[
          \underline{RT}\longrightarrow W_*
                                \xrightarrow{\bar\rho}\underline{RT}.
                                                               \tag{136.FK10}
\]
Every functor preserves this retraction identity. Also
$R\otimes_{\mathrm{Fr},R}R\simeq R$, by
$a\otimes b\mapsto ab^p$, with inverse $a\mapsto a\otimes1$.
Under (136.FK9), the resulting retraction of $K_*$
is a simplicial map to the constant module $R$ which sends
\[
                           z=T^{[p]}\longmapsto1.
                                                               \tag{136.FK11}
\]
In particular $[z]\ne0$ in $H_0(NK_*)$, because $R\ne0$.

The premise of (136.FK3) gives a normalized chain
$u\in N_1P_*^p$ with $\partial u=T^{[p]}$.
Its reduction and image in $Q_*$ is a normalized chain
$v$ with $\partial v=z$. Its image $a$ in $F_*$
is a cycle. The connecting homomorphism in (136.FK6) sends
$[a]$ to the nonzero class $[z]$, so $[a]\ne0$.
The integral relative differential $du$ is a normalized
one-cycle, since $\partial(du)=dT^{[p]}=0$.
Its reduction and augmented image is exactly $a$.
If $du$ were an integral boundary, $a$ would be a boundary,
a contradiction. This proves (136.FK3) without a flat
homology base-change assertion.

If the zero-form row is acyclic in a weight tail, its
coefficient class is killed at every sufficiently large
prime weight. Infinitely many nonunit primes therefore
give arbitrarily large nonzero one-form $H_1$ groups by
(136.FK3), contrary to horizontal admissibility.
\end{proof}

\paragraph{Actual-target consequence.}
Let $B$ be the standard admissible integral theory and fix
$N\ge2$. Define its strict PD subtheory
\[
 C_N^{w,*}=
 \begin{cases}
  kT^{[w]}[0],&w<N,\\
  B^{w,*},&w\ge N.
 \end{cases}                                               \tag{136.FK12}
\]
The coefficients in low weights are the existing coefficient
submodules of $B$. Products, differential and divided powers
preserve (136.FK12): a product involving a retained high
weight stays high, and positive divided powers multiply
weight by a positive integer. The vertical comparison is
unchanged in every sufficiently large weight, as is each
horizontal row. Thus $C_N\in\mathcal A_k$, exactly as in
AQ-136-WEIGHT-ONE.

Any strict map from the full relative special-PD de Rham
algebra of a $P_*$ satisfying (136.FK1) to $C_N$ induces
on weight-one zero forms the retraction (136.FK2), because
$C_N^{1,0}=\underline{kT}$ and coefficients are preserved.
If $k$ has infinitely many nonunit rational primes, the
source therefore cannot itself be admissible.
In particular, over $\mathbb Z$ even arbitrary additional
weight-one free zero-form cells cannot turn this full
relative-de-Rham construction into a source for $C_N$.

This is a construction-class obstruction. The target
$C_N$ is itself admissible and maps to itself, so it is
not a counterexample to the original common-source question.
The coskeleton replacements of AQ-136-FIRST-SKELETON need
not be degreewise free full de Rham algebras, and are not
excluded. Nonfree form relations and other strict PD
presentations remain available for higher comparisons.


\subsection{AQ-SA08-LIFT: a concrete divided-power lifting obstruction}
The supplied SA08 intake concerns the uniqueness problem
of cote 136/117--121 and a proposed lifting input.
The original cote-136 sheets have now been recovered and
analyzed in AQ-136 above.  The auxiliary lifting claim
is checked separately in its primary source: Caradot--Lin (2026),
arXiv:2511.22614v4, \S2.1, pp.~5--6, and
Proposition 3.10, pp.~19--20.
There \emph{free complex} means free in each degree,
and surjectivity of the PD ideals is degreewise.
Under these hypotheses the asserted lift need not exist,
over any nonzero commutative ring $k$.
\begin{proof}
Use cochain degrees and the square-zero PD dg algebras
\[
 R'=k\oplus(ku\xrightarrow{\,d\,}kv),\qquad
 |u|=-2,\quad |v|=-1,\quad du=v,
 \qquad R=k\oplus k\epsilon,\quad|\epsilon|=-2,\quad d\epsilon=0.
\]
Their augmentation ideals have zero multiplication,
and $\gamma_n=0$ for $n\geq2$ defines divided powers.
All divided-power identities hold: products of two
ideal elements vanish, and the differential identity
$d\gamma_n(u)=\gamma_{n-1}(u)\,du$ is zero for $n\geq2$,
including $u\,du=uv=0$ when $n=2$.
Define $\pi(u)=\epsilon$, $\pi(v)=0$, and $\pi|_k=1$.
It is a PD dg map surjective on the ideals.
Take the termwise free complex $X=k\epsilon$,
concentrated in degree $-2$, and the identity chain
map $X\to k\epsilon\subset R$.
A PD dg algebra lift from the free PD dg algebra
on $X$ would restrict on its generator to a chain
lift $\epsilon\mapsto cu$.
Compatibility with $\pi$ forces $c=1$, whereas
compatibility with the differential forces $cv=0$,
hence $c=0$ in $k$. This contradicts $k\ne0$.
\end{proof}
The notation $X$ here denotes just $k$ placed in
degree $-2$; no resolution or projectivity in the
category of complexes is being assumed.
The failed implication is precisely from a graded
lift to a chain lift. Theorem 3.15 of the same
version, pp.~23--24, assumes a quasi-isomorphism.
Our $\pi$ is not one:
$H^*(R')=k$ but $H^*(R)=k\oplus k\epsilon$.
Consequently this counterexample does not refute
that later theorem.

The unresolved SA08 statements are the actual
integral common-source assertion A-CS, terminality
of the all-arrow localization A-TL, and application
of the SA09 comparison structures A-SA09.
The original common-source and groupoid questions are now
identified precisely. Both common sources for finite families
and the terminal ordinary groupoid are proved for every ring
with only finitely many nonunit rational primes, by
AQ-136-SEMILOCAL and AQ-136-SEMILOCAL-SOURCES. Over every
ring, common-source existence is equivalent to connectedness.
Over $\mathbb Z$, that connectedness and the triviality of
all localization loops remain separate gates. The full prior SA09
comparison structures remain an intake-level source gate;
no absent prior proof is claimed checked here.
Neither a weight grading on a divided-power algebra
nor the adopted analytic zero theorem supplies a
chain lift in the displayed counterexample.
That is a specific false lifting route; it is not
a disproof of the remaining existence or
localization conjectures.




\paragraph{AQ-136-DERIVED-COMPATIBILITY: a uniform necessary invariant for simultaneous completion.}
Let $A$ be any simplicial bigraded strict PD differential
$S_k$-algebra; no flatness or admissibility condition on $A$
is imposed.  Let $C_1,\ldots,C_s\in\mathcal A_k$ and fix
strict maps $f_i:A\to C_i$.  Suppose these maps have a
simultaneous admissible completion:
\[
 A\xrightarrow{g}R\xrightarrow{j_i}C_i,
       \qquad R\in\mathcal A_k,\qquad j_i g=f_i.
                                                               \tag{136.DC1}
\]
There is a necessary condition imposing compatibility
after derived coefficient change, in addition to the
integral comparison.
For every $d\ge0$ there is a weight bound $W_d$ such that,
simultaneously for all $w\ge W_d$, all finite simplicial
sets $X$, all integers $n\ge2$, and all $0\le q\le d$,
the maps
\[
 H^q\bigl(A^{w,*}(X)\otimes^{\mathbf L}_{\mathbb Z}
                         \mathbb Z/n\bigr)
 \longrightarrow \mathcal T^q_n(X;k)\otimes_k S^w
                                                               \tag{136.DC2}
\]
induced by the $f_i$ and the canonical admissible comparisons
are equal for every $i$.  Here the target has the following
precise, theory-independent meaning.  Put
\[
 \mathcal E_X(k)=R\operatorname{Hom}_{s\operatorname{Mod}k}
                         (k[X],\underline{k}),\qquad
 \mathcal T^q_n(X;k)=H^q\bigl(\mathcal E_X(k)
                  \otimes^{\mathbf L}_{\mathbb Z}\mathbb Z/n\bigr).
                                                               \tag{136.DC3}
\]
The underline denotes the constant simplicial module.  In
particular these targets have natural exact sequences
\[
 0\longrightarrow H^q(X;k)/n
 \longrightarrow\mathcal T^q_n(X;k)
 \longrightarrow H^{q+1}(X;k)[n]\longrightarrow0.
                                                               \tag{136.DC4}
\]
No splitting of this sequence is chosen.  In (136.DC2), the
notation with $S^w$ uses its specified free coefficient
generator $T^{[w]}$; equivalently one may replace
$\underline{k}$ in (136.DC3) by $\underline{S^w}$.

The bound $W_d$ is independent of $n$ and $X$.
It can depend on the proposed completion $R$ and on the
finite target family.  Consequently a discrepancy for
arbitrarily large weights, even for one fixed $X$, one
fixed prime modulus and one fixed degree, rules out that
simultaneous completion.  This is a necessary condition,
not a sufficiency criterion for original common sources.

\begin{proof}
We first specify the derived coefficient operation on a
possibly nonflat complex.  For any cochain complex $M$ of
abelian groups define
\[
 D_n(M)^q=M^q\oplus M^{q+1},\qquad
 d(a,b)=(da+n b,-db).
                                                               \tag{136.DC5}
\]
This is the tensor product with the free resolution
$[\mathbb Z\xrightarrow{n}\mathbb Z]$ of $\mathbb Z/n$,
placed in degrees $-1,0$, using the equivalent mapping-cone
sign convention.  A bounded complex of free abelian groups
is $K$-flat: tensoring an acyclic complex with each free
term is exact, and the finite total complex is acyclic by
the finite filtration.  Thus (136.DC5) computes
$M\otimes^{\mathbf L}_{\mathbb Z}\mathbb Z/n$ for every
$M$, with no flatness assumption on its terms.
The mapping-cone long exact sequence gives, naturally,
\[
 0\longrightarrow H^q(M)/n
 \longrightarrow H^q(D_n(M))
 \longrightarrow H^{q+1}(M)[n]\longrightarrow0.
                                                               \tag{136.DC6}
\]
It follows that a map which is an isomorphism on integral
cohomology in degrees $q,q+1$ is an isomorphism on
$H^q(D_n(-))$, for every $n$, with no change to its
cohomological range.  This follows directly from the
commutative diagram of (136.DC6), by the short five lemma.

We use the actual natural comparison construction in
AQ-136-COMPARISON, rather than choosing unrelated
isomorphisms of cohomology groups.  Write
$\Gamma_X(F)=\operatorname{Hom}_{s\operatorname{Mod}k}(k[X],F)$.
For any $C\in\mathcal A_k$ there are natural maps in the
derived category of $k$-modules
\[
 C^{w,*}(X)=\Gamma_X(C^{w,*})
       \longrightarrow R\Gamma_X(C^{w,*})
       \longleftarrow R\Gamma_X(\underline{S^w}).
                                                               \tag{136.DC7}
\]
The right-hand map is induced by the coefficient
augmentation $\underline{S^w}\to C^{w,*}$, and the
left-hand map is the ordinary-to-derived-section map.
Both are therefore natural for every strict
$S$-preserving morphism $C\to C'$.
The complete proof of AQ-136-COMPARISON gives one weight
cutoff, independent of $X$, at which both arrows are
isomorphisms on cohomology through any prescribed finite
degree.  Indeed its first hypercohomology spectral sequence
uses the uniform vertical cohomology bounds, and its
second uses the acyclicity of each required horizontal
row for all higher derived sections.  All total degrees
in that argument have finite diagonals.  Thus this is a
statement about the two displayed maps themselves, not
only the existence of abstract cohomology isomorphisms.
The same proof identifies
$H^q(R\Gamma_X(\underline{k}))$ naturally with $H^q(X;k)$.

For the finite family $R,C_1,\ldots,C_s$, take the maximum
of their comparison cutoffs through degree $d+1$.
Apply $D_n$ to the two arrows in (136.DC7).
By (136.DC6), their cohomology maps through degree $d$
are isomorphisms for every $n\ge2$.  Inverting these
cohomology isomorphisms gives the canonical comparisons
with (136.DC3), and (136.DC6) applied to the coefficient
object proves (136.DC4).
Since every $j_i$ preserves the coefficient map, naturality
of (136.DC7) says that it acts as the identity on the
common target (136.DC3), including after $D_n$.
After precomposition with $D_n(g)$, all the maps in
(136.DC2) are therefore equal.  The weight cutoff was
chosen before $n$ and $X$, so it has the asserted uniformity.

This argument performs derived coefficient change
\emph{after} taking sections.  It never replaces
$A(X)\otimes^{\mathbf L}\mathbb Z/n$ by ordinary reduction
$A(X)/n$, nor identifies it with sections of a
coefficient-changed simplicial algebra.  Such an
interchange could fail for a nonflat source.  In fact
the comparison proof used above works for arbitrary
simplicial sets, so the stated finite-$X$ version needs
no additional finiteness or base-change theorem.
\end{proof}

\paragraph{The earlier obstruction is exactly a failed derived compatibility test.}
Retain the explicit $A$, $f_0,f_1:A\to B$ from
AQ-136-HIGHER-COMPLETION, over $k=\mathbb Z$ and at a
fixed prime $p$.  On the simplicial circle $X$, its
elements $v_w,u_w$ satisfy
$dv_w=p u_w$ and $du_w=0$.  Hence
\[
                 (v_w,-u_w)\in D_p(A^{w,*}(X))^1
                                                               \tag{136.DC8}
\]
is a closed element by (136.DC5).  Its image under $f_0$
is zero.  Under $f_1$ it is
$(T^{[w-1]}dx,0)$, whose class is the image of
$wT^{[w]}$ under the left-hand injection of (136.DC6).
The integral endpoint calculation (136.HC8) gives this
normalization without any coefficient-change assumption.
Since $H^1(B^{w,*}(X))=\mathbb ZT^{[w]}$ and
$H^2(B^{w,*}(X))=0$ in the comparison range, this class
is nonzero exactly when $p\nmid w$.
Such weights are arbitrarily large.  Thus the uniform
necessary condition (136.DC2) already contradicts the
existence of (136.HC2).  This recovers the checked
obstruction with a single mod-$p$ cone class, without
an assertion that the source is flat or that ordinary
mod-$p$ sections compute its derived coefficient change.
For a proposed positive source construction, retaining
only its integral $H^0/H^1$ and rational comparison is
therefore insufficient: the prescribed maps must also
preserve these derived coefficient compatibilities.


\paragraph{AQ-136-CARTAN-SCOPE: what the recovered 1976 primary article proves.}
The entire primary article has been read from its page images:
Cartan, H. (1976). Th\'eories cohomologiques. \emph{Inventiones
Mathematicae, 35}, 261--271. The official GDZ article PDF is
\url{https://gdz.sub.uni-goettingen.de/download/pdf/PPN356556735_0035/LOG_0030.pdf}.
Its first page is a repository cover; physical pages 2--12
are the eleven printed article pages 261--271.

Section 2, p.~263, imposes vertical exactness with constant
degree-zero cycles and horizontal acyclicity in every form
degree. Theorem 1, p.~264, obtains ordinary cohomology.
In Section 4, p.~266, a morphism means a simplicial
$k$-linear cochain map: compatibility with multiplication
is expressly not required. Theorem 2 on that page gives
a unique such integration map to ordinary cochains when
forms vanish above the simplex dimension. Theorem 3,
p.~267, proves multiplicativity on cohomology under a
flatness hypothesis on cycle modules. Section 6, p.~268,
weakens the horizontal axiom. Section 7, pp.~269--271,
constructs the Grothendieck divided-power simplex theory,
obtains $H^n(Gr^{*,r}(X))=H^n(X;\mathbb Z)$ for $r\ge n$
and zero for $r<n$, and identifies its weight-$(r,s)$
cohomology product with $\binom{r+s}{r}$ times cup product.
This is the standard theory $B$ used in this module,
written after eliminating one barycentric coordinate.
The article supplies neither strict PD common sources
nor a terminal-localization theorem for all theories.
Its linear integration theorem cannot be substituted for
either original assertion at cote 136, physical 121.

\paragraph{AQ-136-MULTIPLICATIVE: the canonical comparison preserves the weighted cup product without flatness.}
For every $C\in\mathcal A_k$ and every $D\ge0$, there is
$W_D$ with the following property. For $w,v\ge W_D$,
$p,q\ge0$ and $p+q\le D$, let
\[
 \kappa_w^p:H^p(C^{w,*}(X))\xrightarrow{\ \sim\ }H^p(X;k)
                                                               \tag{136.M1}
\]
denote the canonical comparison, using $T^{[w]}$ to
identify the coefficient line with $k$. Then for every
simplicial set $X$,
\[
 \kappa_{w+v}^{p+q}(\alpha\beta)
   =\binom{w+v}{w}\,
                 \kappa_w^p(\alpha)\smile\kappa_v^q(\beta).
                                                               \tag{136.M2}
\]
No flatness assumption on the terms or cycle modules of
$C$ is needed. The proof uses the underlying cycle
\emph{spaces}, not a possibly nonflat tensor product of
cycle modules. The statement concerns the cohomology
product, and does not produce a strict comparison map.

\begin{proof}
Choose $W_D\ge1$ so that, in every weight at least $W_D$,
the uniform vertical tests hold through degree $D+2$
and all horizontal rows through degree $D+2$ are acyclic.
The finite maximum of their cutoffs exists. It applies
also to $w+v$. We make the comparison, representability,
and normalization steps explicit.

\emph{Relative comparison and connecting maps.}
For a simplicial inclusion $Y\subset X$, put
\[
 C^{w,*}(X,Y)=\ker(C^{w,*}(X)\longrightarrow C^{w,*}(Y)).
\]
The normalized complex of $k[X]/k[Y]$ is free, with
basis the nondegenerate simplices of $X$ not in $Y$.
The complete derived-section proof of AQ-136-COMPARISON
therefore applies to this quotient exactly as it applies
to $k[X]$. It gives canonical relative comparisons with
$H^*(X,Y;k)$ throughout the chosen finite range.

Restriction to $Y$ is surjective in each of the required
cochain degrees. One direct verification is to extend
a normalized chain map to a horizontally acyclic row,
one relative free generator at a time in increasing
chain degree. After the previously prescribed boundary
is mapped, its image is a cycle. Acyclicity, including
$H_0=0$, makes that cycle a boundary; choosing a
preimage defines the value on the new generator.
This proves extension in all degrees. Thus the ordinary
relative short exact sequences are exact in the finite
cochain range needed here. Their comparison with the
derived relative-section triangle shows that the
canonical maps $\kappa$ commute with connecting
homomorphisms. They are normalized in degree zero by
the actual coefficient map $S^w\to C^{w,0}$.

The algebra multiplication also gives relative external
products
\[
 H^p(C^w(X,Y))\otimes H^q(C^v(X',Y'))
 \longrightarrow H^{p+q}
 \bigl(C^{w+v}(X\times X',\,
                Y\times X'\cup X\times Y')\bigr),
                                                               \tag{136.M3}
\]
by pulling back both factors and multiplying. Here and
below the suppressed star denotes the full cochain
complex. The products vanish on the indicated union
because one of their factors vanishes on each piece.
If $a$ is a cocycle on a subspace and $\widetilde a$
is a cochain extension, then its connecting class is
represented by $d\widetilde a$. For a second cocycle $b$,
\[
 d(\widetilde a\,b)=d\widetilde a\,b,
   \qquad
 d(a\,\widetilde b)=(-1)^p a\,d\widetilde b
                  \quad (da=0,\ |a|=p).
                                                               \tag{136.M4}
\]
These formulas prove compatibility of (136.M3) with
connecting maps in the first and second factor, with
the displayed sign. They require no tensor exactness.
The same formulas hold for the ordinary relative cup
and cross products; see Hatcher (2002), p.~218, in
the proof of Theorem 3.15 and the discussion preceding
Theorem 3.18.

\emph{Cycle spaces represent the cohomology.}
Write $Z^r_w=\ker(d:C_*^{w,r}\to C_*^{w,r+1})$.
In the chosen weight range $Z^0_w=\underline{S^w}$
and, for $0\le r<D+1$, there are short exact
sequences of simplicial modules
\[
 0\longrightarrow Z^r_w\longrightarrow C_*^{w,r}
                   \xrightarrow{d}Z^{r+1}_w\longrightarrow0.
                                                               \tag{136.M5}
\]
Each middle object is horizontally acyclic. The homotopy
long exact sequences, beginning with the constant
$Z^0_w$, identify $Z^r_w$ as an Eilenberg--Mac Lane
space $K(S^w,r)$: its only nonzero homotopy group is
$\pi_r=S^w$. This includes all homotopy degrees, since
each middle row is acyclic in every horizontal degree.
Every simplicial abelian group is Kan. The surjection
in (136.M5) is a fibration with contractible total
space, so a map $X\to Z^{r+1}_w$ lifts to that total
space exactly when it is nullhomotopic. The lifting
assertion follows by lifting a nullhomotopy from its
zero endpoint; conversely any lift is nullhomotopic
because its target is contractible. Therefore
\[
 H^r(C^{w,*}(X))=[X,Z^r_w].
                                                               \tag{136.M6}
\]
For $r=0$ the same statement says that closed
zero-forms are locally constant coefficient functions.
This is also the argument of Cartan's Theorem 1,
p.~264, applied only to the finitely many rows in
(136.M5).

The identification of the homotopy group in (136.M5)
is the repeated connecting identification starting
from $S^w$. It agrees with the canonical derived
comparison: both use the coefficient map in degree
zero and commute with the relative connecting maps
just proved. Equivalently, on an oriented relative
$r$-cube both identifications send the repeated
boundary class of a coefficient to that coefficient.
The Eilenberg--Mac Lane representing isomorphism is
determined by this normalization on $\pi_r$: its
fundamental cohomology class corresponds to the
identity on $H_r(K(S^w,r);\mathbb Z)=S^w$.
This last characterization is Hatcher, Theorem 4.57
and its explanation, pp.~393--394. Thus (136.M6)
uses the particular $\kappa$ in (136.M1), rather
than an unspecified additive automorphism.

\emph{The lowest external pairing is forced by coefficients.}
Assume first $p,q>0$. Multiplication gives a pointed
simplicial map
\[
 Z^p_w\times Z^q_v\longrightarrow Z^{p+q}_{w+v}
\]
which vanishes on the wedge, and hence factors through
$Z^p_w\wedge Z^q_v$. The latter smash product has
zero reduced integral homology below degree $p+q$
and
\[
 \widetilde H_{p+q}(Z^p_w\wedge Z^q_v;\mathbb Z)
                         = S^w\otimes_{\mathbb Z}S^v.
                                                               \tag{136.M7}
\]
Indeed each factor has first nonzero reduced integral
homology in its indicated degree, equal to its first
homotopy group (for degree one this is abelianization,
already an abelian group). The reduced integral
K\"unneth sequence has no Tor contribution in degree
$p+q$, and no tensor or Tor contribution below it.
This uses the usual free integral chains of spaces,
not the terms of $C$. Exact primary locators are
Hatcher, Theorem 4.32, p.~366, Theorem 3B.6,
pp.~275--276, and its reduced smash-product version
on p.~276.

By the integral universal coefficient theorem,
\[
 \widetilde H^{p+q}(Z^p_w\wedge Z^q_v;S^{w+v})
   =\operatorname{Hom}_{\mathbb Z}
                      (S^w\otimes_{\mathbb Z}S^v,S^{w+v}).
                                                               \tag{136.M8}
\]
There is no Ext term, since the preceding reduced
homology group vanishes. This is Hatcher, Theorem 3.2,
p.~195. Representability by the target Eilenberg--Mac
Lane space now says that the entire pairing is
determined by its values on cross products of sphere
classes representing $a\in S^w$ and $b\in S^v$.
This argument never uses flatness of these coefficient
groups over $\mathbb Z$, either.

Compute that value using the relative product of
oriented cubes $(I^p,\partial I^p)$ and
$(I^q,\partial I^q)$, with $I=\Delta[1]$.
Here is an explicit connecting induction, including
the relative pairs. Write $u_j(a)$ for the class
corresponding to $a$ in the oriented relative $j$-cube;
for $j=0$ this is the coefficient on a point. For
$j\ge1$ put
\[
 \begin{split}
 K&=I\times I^{j-1},\\
 A&=I\times\partial I^{j-1},\\
 L&=(\partial I\times I^{j-1})\cup A=\partial I^j.
 \end{split}
                                                               \tag{136.M9a}
\]
For $j=1$ the boundary of $I^0$ is empty. The
relative group in degree $j-1$ on $(L,A)$ is two
copies of the group for $(I^{j-1},\partial I^{j-1})$.
This follows directly from compatible sections on
the two disjoint end faces modulo $A$. In the long
exact sequence of the triple $A\subset L\subset K$,
the connecting map sends
$(0,u_{j-1}(a))$ to $u_j(a)$, where the orientation
of $I$ is from its zeroth to its first endpoint.
The same assertion holds in $C$ because the relative
comparison commutes with this connecting map.

Keep the second cube $(I^q,\partial I^q)$ fixed.
In the first-factor connecting step use the product
triple with
$A\times I^q\cup K\times\partial I^q$ and
$L\times I^q\cup K\times\partial I^q$.
The first identity of (136.M4) says that the product
of $\delta(0,u_{j-1}(a))$ with $u_q(b)$ is the
connecting image of
$(0,u_{j-1}(a)\mathbin{\times}u_q(b))$.
This is exactly the same identity and orientation
for ordinary cross product. In the initial case
$j=0$, multiplication by the coefficient $a$ is a
cochain map whose coefficient map is multiplication
in $S$, so naturality of the comparison gives the
claimed value on $u_q(b)$. Induction on $j$ now
proves that the product on the two relative cubes
is ordinary cross product followed by
\[
 S^w\otimes_{\mathbb Z}S^v\longrightarrow S^{w+v},
 \quad (aT^{[w]})\otimes(bT^{[v]})
       \longmapsto ab\binom{w+v}{w}T^{[w+v]}.
                                                               \tag{136.M9}
\]
Surjectivity of restriction, also for these triples,
is the same relative free-chain extension proved
above. Thus each connecting step used here is an
actual cochain calculation for $C$. This proves the required value
on every simple tensor in (136.M7). Such tensors
generate that group, so (136.M8) proves equality of
the universal pairings. Pulling them back along the
two classifying maps and the diagonal of $X$ gives
(136.M2) for $p,q>0$.

If $p=0$, a class is a locally constant coefficient
function. On each connected component multiplication
by its value $aT^{[w]}$ is a natural cochain map
$C^{v,*}\to C^{w+v,*}$ whose map on the coefficient
line is multiplication by the value in (136.M9).
Naturality of the comparison roof therefore gives
(136.M2) in this case. The argument for $q=0$ is
the same, and for $p=q=0$ the formula is simply
the divided-power multiplication in $S$. This treats
every asserted degree, with the single original
weight cutoff $W_D$.
\end{proof}

\paragraph{AQ-136-PD-OPERATIONS: a remaining natural invariant, with its precise torsion limitation.}
Fix $r\ge1$ and $m\ge1$. For every $C\in\mathcal A_k$,
its divided powers induce, for all sufficiently large
$w$, a natural operation
\[
 G_{C,m,w}^{2r}:H^{2r}(X;k)\longrightarrow H^{2rm}(X;k).
                                                               \tag{136.PD1}
\]
The operation is defined using the canonical weight-$w$
and weight-$mw$ comparisons. Strict morphisms of
admissible theories preserve it in their common weight
range. It satisfies
\[
 m!\,G_{C,m,w}^{2r}(a)
       =\frac{(mw)!}{(w!)^m}\,a^m.
                                                               \tag{136.PD2}
\]
Consequently, if multiplication by $m!$ is injective on
$H^{2rm}(X;k)$, then
\[
 G_{C,m,w}^{2r}(a)
       =\frac{(mw)!}{m!(w!)^m}\,a^m.
                                                               \tag{136.PD3}
\]
In particular this holds over a $\mathbb Q$-algebra.
The displayed integer coefficient counts partitions of
$mw$ labelled elements into $m$ unordered blocks of
size $w$, so it is integral for $w\ge1$.

\begin{proof}
The special differential identity shows that
$\gamma_m$ sends closed even forms of bidegree
$(w,2r)$ to closed forms of bidegree $(mw,2rm)$.
It is a simplicial map
$Z^{2r}_w\to Z^{2rm}_{mw}$, preserving zero; it need
not be a linear map. Choose the weight cutoff so that
the two cycle-space descriptions (136.M6) hold.
Two cocycles which define the same cohomology class
are then homotopic as maps from $X$ to $Z^{2r}_w$.
Composing their simplicial homotopy with $\gamma_m$
gives a homotopy in $Z^{2rm}_{mw}$. Hence their
images define the same cohomology class. This proves
well-definedness without presuming that divided powers
of arbitrary boundaries are always boundaries in an
unrestricted PD differential algebra.
Naturality in $X$ is immediate from sections, and
strict PD morphisms commute with the operation.

The identity $m!\gamma_m(x)=x^m$ holds in every PD
algebra. Iterating the product formula (136.M2) gives
the factor
\[
 \prod_{j=1}^{m-1}\binom{(j+1)w}{w}
                          =\frac{(mw)!}{(w!)^m}
\]
for the $m$th power of a weight-$w$ class. Taking a
single large enough cutoff for these finitely many
products proves (136.PD2). Subtract the right-hand
side of (136.PD3) from the operation. Its multiple
by $m!$ is zero, so injectivity proves (136.PD3).
\end{proof}

The ordinary cup-product possibility therefore gives
no component obstruction under the exact admissibility
tests. Higher divided-power operations are well defined
and preserved, but the equation (136.PD2) alone leaves
an $m!$-torsion ambiguity. No pair of admissible theories
with different operations has been constructed here,
and no assertion that this ambiguity is realized is
made. Nor is equality of these operations claimed
sufficient for a strict common source or for vanishing
of all localization loops.

\paragraph{Additional primary reference.}
Hatcher, A. (2002). \emph{Algebraic topology}. Cambridge
University Press. Corrected author PDF:
\url{https://pi.math.cornell.edu/~hatcher/AT/AT.pdf}.
The exact statements used above are Theorem 3.2,
p.~195; the relative connecting-product calculation,
p.~218; Theorem 3B.6 and its reduced relative form,
pp.~275--276; Theorem 4.32, p.~366; and Theorem 4.57,
pp.~393--394. The source statements were inspected;
their full textbook proofs are cited inputs, not
independently reproved here.


\paragraph{AQ-136-PD-RIGIDITY: finitely many local comparisons determine each divided-power operation.}
Let $C,D\in\mathcal A_k$, fix $r\ge1$ and $m\ge1$,
and use the operations of AQ-136-PD-OPERATIONS.
There is $W=W(C,D,r,m)$ such that
\[
           G_{C,m,w}^{2r}(a)=G_{D,m,w}^{2r}(a)
                    \qquad(w\ge W)                       \tag{136.PR1}
\]
for every finite simplicial set $X$ and
$a\in H^{2r}(X;k)$.
If the additive group of $k$ is finitely generated,
then (136.PR1) holds for \emph{every} simplicial set
$X$, with the same independence of the weight cutoff
from $X$. In particular the all-$X$ assertion holds
for the original integral coefficient ring $k=\mathbb Z$.
One may take $D=B$, the explicit standard integral
theory. This asserts equality with its operation,
without replacing that operation in torsion cohomology
by an unproved formula for an ordinary cup power.

\begin{proof}
For $m=1$ the operation is the identity, so assume
$m\ge2$ and write $d=m!$, $s=2r$, and $t=ms$.
For a rational prime $\ell$ let
\[
       k_{(\ell)}=k\otimes_{\mathbb Z}\mathbb Z_{(\ell)},
       \qquad C_{(\ell)}=C\otimes_k k_{(\ell)},
                                                               \tag{136.PR2}
\]
and define $D_{(\ell)}$ likewise. These are flat
localizations. Cycles, cohomology and normalized
simplicial homology commute with them, so the two
localized theories are admissible over $k_{(\ell)}$
with the original weight cutoffs. Divided powers
base change to the specified coefficient algebra
$\Gamma_{k_{(\ell)}}(T)$. If this coefficient ring
is zero, the assertions involving it are vacuous.

By the checked p-local common-source theorem, for
each of the finitely many primes $\ell\mid d$ there
is an admissible theory $R_\ell$ over $k_{(\ell)}$
and strict PD maps
\[
                    C_{(\ell)}\longleftarrow R_\ell
                                      \longrightarrow D_{(\ell)}.
                                                               \tag{136.PR3}
\]
In their common sufficiently large weight range,
naturality of the comparisons and of divided powers
therefore identifies the operations of these two
localized theories on all simplicial sets. There
is one such cutoff for each $\ell\mid d$, and we
take their finite maximum along with the cutoffs
for (136.PD1)--(136.PD2) for $C$ and $D$.

We specify coefficient naturality even on an infinite
simplicial set, without asserting that sections
commute with localization there. The actual map
$C\to C_{(\ell)}$ induces a map of the cycle
Eilenberg--Mac Lane spaces
\[
          Z^s_w(C)\longrightarrow Z^s_w(C_{(\ell)}).
\]
On their only nonzero homotopy groups it is exactly
$k\to k_{(\ell)}$, since normalization and homology
commute with the flat localization. The canonical
cohomology comparison is the coefficient-normalized
representing comparison proved in AQ-136-MULTIPLICATIVE.
Thus the induced map on cohomology for any $X$ is
the ordinary change-of-coefficients map. The same
argument applies in degree $t$ and weight $mw$.
Because the actual coefficient map commutes with
divided powers, it follows that
\[
 G_{C,m,w}^{s}(a)\big|_{k_{(\ell)}}
       =G_{C_{(\ell)},m,w}^{s}
                             (a\big|_{k_{(\ell)}}),
                                                               \tag{136.PR4}
\]
and likewise for $D$. Only these natural maps are
used here; none is claimed to be an isomorphism yet.

Set
\[
                  z=G_{C,m,w}^{s}(a)-G_{D,m,w}^{s}(a)
                               \in H^t(X;k).
\]
The factorial identity (136.PD2), whose right-hand
side is the same weighted cup power for both
theories, gives $dz=0$. Equations (136.PR3)--(136.PR4)
show that the image of $z$ in
$H^t(X;k_{(\ell)})$ is zero for every $\ell\mid d$.

Suppose first that $X$ is finite. Its normalized
integral chain complex is a bounded complex of
finite free modules. Applying $\operatorname{Hom}(-,k)$
and flat localization therefore gives the natural
isomorphism
\[
 H^t(X;k)\otimes_{\mathbb Z}\mathbb Z_{(\ell)}
                         \simeq H^t(X;k_{(\ell)}).
                                                               \tag{136.PR5}
\]
Consequently $z$ has zero localization at every
prime dividing $d$. An element killed by $d$ and
with this property is zero: if its nonzero order
were $e\mid d$, choose a prime $\ell\mid e$.
Vanishing in the $\ell$-localization would imply
$u z=0$ for an integer $u$ prime to $\ell$, whereas
the order $e$ would then divide $u$, a contradiction.
This proves (136.PR1) for all finite $X$, with one
cutoff for all of them.

Now suppose the additive group of $k$ is finitely
generated. Let $Y=K(k,s)$ and let
$\iota\in H^s(Y;k)$ be its fundamental class.
The integral homology groups of $Y$ are finitely
generated in every degree. An exact primary input
is Hatcher (2004), Lemma 1.10, p.~15, applied to
the Serre class of finitely generated abelian
groups. Its proof starts with cyclic groups and
uses the path fibration and Lemma 1.9; in particular
the statement applies to $s=2r\ge2$ here.

For any space whose integral homology in degrees
$t,t-1$ is finitely generated, (136.PR5) still
holds. To see this without localizing an infinite
cochain complex, use the natural universal
coefficient short exact sequence
\[
 0\to\operatorname{Ext}_{\mathbb Z}^1(H_{t-1}(Y;\mathbb Z),k)
  \to H^t(Y;k)
  \to\operatorname{Hom}_{\mathbb Z}(H_t(Y;\mathbb Z),k)\to0.
                                                               \tag{136.PR6}
\]
Each finitely generated homology group has a
resolution by two finite free abelian groups.
Both Hom and Ext in (136.PR6) therefore commute
with localization of $k$. Flatness localizes the
short exact sequence itself, and the natural
diagram with the coefficient-$k_{(\ell)}$ sequence
has isomorphisms at both ends. Its middle map is
an isomorphism by the short five lemma, proving
(136.PR5) for this particular universal space.

Apply the already established local comparison
(136.PR4) to $X=Y$ and $a=\iota$. Its difference
$z$ is $d$-torsion and is now detected by the
localizations (136.PR5). Thus $z=0$ on the
universal fundamental class. Every element of
$H^s(X;k)$ is the pullback of $\iota$ along a map
$X\to K(k,s)$, by representability. Naturality
then proves (136.PR1) for arbitrary $X$. No
additional cutoff depends on the representing
map or on $X$. This proves the all-$X$ assertion.
\end{proof}

\paragraph{AQ-136-LOCALIZATION-CAUTION: torsion cohomology on an arbitrary infinite space need not be detected by all coefficient localizations.}
The finite-space or universal finite-generation step
in the preceding proof is necessary to its reasoning.
Here is an explicit warning, unrelated to a counterexample
to original common-source existence.
Fix a prime $p$, let $A=\mathbb Z_{(p)}$, choose $s\ge2$,
and put $X=K(A,s)$. There is a nonzero class
\[
             z\in H^{s+1}(X;\mathbb Z),\qquad pz=0,
                                                               \tag{136.PR7}
\]
whose image in $H^{s+1}(X;\mathbb Z_{(\ell)})$ is
zero for \emph{every} rational prime $\ell$.

\begin{proof}
The first nonzero reduced integral homology of $X$
is $H_s(X;\mathbb Z)=A$. Universal coefficients
therefore give
\[
 H^s(X;\mathbb Z)=\operatorname{Hom}_{\mathbb Z}(A,\mathbb Z)=0,
 \qquad H^s(X;\mathbb Z/p)=\operatorname{Hom}_{\mathbb Z}
                                      (A,\mathbb Z/p)=\mathbb Z/p.
                                                               \tag{136.PR8}
\]
For the first equality, the image of $1$ under a
homomorphism $A\to\mathbb Z$ must be divisible by
every power of any fixed prime distinct from $p$,
and hence must be zero. Such a homomorphism is
then zero on every fraction. For the second, the
value at $1$ determines the map, and reduction
modulo $p$ extends to $A$ because its denominators
are prime to $p$.

The coefficient sequence
$0\to\mathbb Z\xrightarrow{p}\mathbb Z\to\mathbb Z/p\to0$
has a connecting map whose restriction from the
second group in (136.PR8) is injective, because
the first group is zero. Let $z$ be the image of
the nonzero reduction homomorphism. It is nonzero
and killed by $p$. After replacing integral
coefficients by $\mathbb Z_{(p)}$, that reduction
homomorphism lifts to the identity
$A\to\mathbb Z_{(p)}$ in $H^s$. Its connecting
class is consequently zero. Naturality of the
coefficient long exact sequence proves that the
image of $z$ is zero at $p$. For $\ell\ne p$,
multiplication by $p$ is an automorphism of the
entire coefficient cochain complex with values
in $\mathbb Z_{(\ell)}$ and hence of its cohomology.
The image of a $p$-torsion class there must be
zero. This proves (136.PR7) and the claimed
failure of naive detection.
\end{proof}

\paragraph{Scope of the rigidity result.}
For each fixed divided-power order and form degree,
the common-source results already proved at finitely
many coefficient primes force equality of the
cohomology operations on finite simplicial sets
over every $k$, and on all simplicial sets over
$\mathbb Z$ and more generally over a ring with
finitely generated additive group. This closes
that proposed operation-level distinction in
those scopes. It does not construct an integral
strict common source, a compatible tower of
finite-stage sources, or a contraction of all
ordinary-localization loops.

\paragraph{Exact additional reference.}
Hatcher, A. (2004, January 5). \emph{Spectral sequences
in algebraic topology}, Chapter 1: The Serre spectral
sequence. Author manuscript.
\url{https://pi.math.cornell.edu/~hatcher/SSAT/SSch1.pdf}.
Lemma 1.10 is stated on p.~15; its complete proof on
pp.~15--16 was directly inspected. Lemma 1.9, p.~15, supplies its fibration
step. Universal coefficients and representability are
the previously inspected Hatcher (2002), Theorems 3.2
and 4.57. Those primary results are cited inputs, not
new results of this module.

% END archive_questions.tex

% BEGIN consequences_criteria_integral_albanese.tex
\section{The integral parameter relation in original cote 14}
\label{sec:cc-integral-albanese}

\paragraph{Scope and convention.}
Let $k$ be algebraically closed of characteristic $p>0$ and
let $T/k$ be a smooth connected quasiprojective variety.
Write
\[
 H_T=H^S_0(T,\mathbb Z)^0,\qquad
 A_T=\operatorname{Alb}_{\rm ab}(T),\qquad
 K_T=\ker\bigl(H_T\longrightarrow A_T(k)\bigr).
\]
Here $\operatorname{Alb}_{\rm ab}$ means the ordinary abelian
Albanese: after choosing $t_0\in T(k)$, its pointed morphism
$a:T\to A_T$ is universal for pointed morphisms to
\emph{abelian varieties}.  It is the abelian quotient of the
generalized semiabelian Albanese.  We use the integral Suslin
complex of finite correspondences.  In particular $H_T$ is a
quotient of the group of free degree-zero zero cycles, and
the Albanese homomorphism is the sum of their images.

The new deduction below proves divisibility of $K_T$ and
then removes the possible $p$-primary discrepancy left by
AQ-14-ALTERATION.  It requires no resolution of singularities
or smooth compactification of $T$.  It does not assert
injectivity of the Albanese map on $p$-primary torsion.

\paragraph{Exact cited inputs.}
Katz (1993), Corollary 3.4.3, printed p.~182, gives smooth
irreducible affine hyperplane sections of a smooth
irreducible variety equipped with an unramified morphism to
affine space; the sections can simultaneously preserve
irreducibility of any specified finite \'{e}tale cover.
Its hypotheses are over an algebraically closed field of
characteristic different from the auxiliary coefficient
prime.  We choose that prime different from $p$.
Suslin and Voevodsky (1996), Theorem 3.1, printed p.~65,
identifies the integral zeroth Suslin homology of a smooth
affine curve $C$ with
$\operatorname{Pic}(\overline C,\overline C\setminus C)$,
where the boundary is reduced.
The generalized Albanese universal property is recalled
in Spie{\ss} and Szamuely (2000), \S3, pp.~4--5;
passing to its abelian quotient gives the ordinary
universal property used here.
For the auxiliary curve-through-points construction we
use the Stacks Project, Lemmas 33.47.1 (Tag 0FD5) and
37.32.3 (Tag 0G4F).
Quotients by finite group schemes, their torsor property,
the connected--\'{e}tale sequence, and finiteness and
surjectivity of multiplication on abelian varieties are
recorded in Milne (2017), Theorem 5.28 and Corollaries
5.26--5.27, p.~105, Proposition 5.58, p.~114, and 8.21,
p.~152.  These statements include nonreduced subgroup
schemes; treating only abstract groups of points here
would lose the characteristic-$p$ information.

\subsection{CC-ALB01: divisible ordinary Albanese kernel}

\paragraph{Statement.}
For every positive integer $n$ the map
\[
 H_T[n]\longrightarrow A_T(k)[n]
 \quad\text{is surjective}.                         \tag{CC.ALB.1}
\]
The group $H_T$ is divisible, and its subgroup $K_T$ is
divisible.

\begin{proof}
\emph{1. Curve classes make $H_T$ divisible.}
First, any finite set of points in an integral
quasiprojective variety over $k$ lies on an integral
curve, unless the variety has dimension zero.
Here is the precise Bertini reduction used for this fact.
Take an integral projective closure $Y$, let $F$ be the
finite set, and use sections of
$\mathcal I_F\otimes\mathcal L^{\otimes m}$ for $m$
sufficiently large and $\mathcal L$ ample.
The cited Tag 0FD5 says that this linear system has no
base points outside $F$ and gives an immersion there.
The base locus on $Y$ is exactly $F$, because the ideal
sheaf is generated after a sufficiently high twist.
If $\dim Y\geq2$, this base locus has codimension at least
two.  At a smooth point outside $F$, immersion gives one
section nonzero at that point and two sections whose
local equations have independent differentials.
Their intersection has a codimension-two component
through that point, not contained in the first section.
These verify conditions (3) and (4), as well as (1) and
(2), of the cited Tag 0G4F.  Thus a general member is
geometrically irreducible.  Its reduction contains $F$
and has dimension one less.  Iterating gives an
integral curve through $F$ in the projective closure;
its nonempty intersection with the original variety
is the required integral curve.

Normalize this curve and choose one preimage of each
point of $F$.  Normalization is finite, and a normal
curve over the perfect field $k$ is smooth.  If necessary
remove a finite set disjoint from the chosen preimages
so that the resulting curve $C$ is affine.  Thus every
free degree-zero cycle on $T$ is the image of a
degree-zero cycle on some smooth affine curve $C$.
All residue fields are $k$, so the degree is unchanged.

Let $\overline C$ be the smooth projective completion,
let $D=\overline C\setminus C$ with its reduced structure,
and put $J=\operatorname{Pic}^0(\overline C)$.
The curve calculation in the cited Suslin--Voevodsky
theorem and the elementary exact sequence for line
bundles with a trivialization on $D$ give
\[
 0\longrightarrow (k^\times)^{|D|-1}
 \longrightarrow H^S_0(C,\mathbb Z)^0
 \longrightarrow J(k)\longrightarrow0.             \tag{CC.ALB.2}
\]
Indeed, a line bundle restricts to a trivial line bundle
on the finite reduced set $D$; two choices of
trivialization differ by an element of
$(k^\times)^D/k^\times$.
Both outer groups in (CC.ALB.2) are divisible:
this is immediate for $k^\times$, and multiplication
by $n$ on an abelian variety is surjective.
An extension of divisible groups is divisible: divide
the image in the quotient, lift the division, and
correct the remaining error by division in the kernel.
The functorial pushforward to $H_T$ therefore divides
each degree-zero cycle class by any prescribed $n$.
This proves divisibility of $H_T$.  If $\dim T=0$,
both $H_T$ and $A_T$ are zero and all assertions follow.
We henceforth suppose $\dim T\geq1$.

\emph{2. Albanese universality makes isogeny pullbacks
connected.}
Put $A=A_T$ and normalize $a(t_0)=0$.
Let $\pi:B\to A$ be any finite \'{e}tale isogeny of
abelian varieties.  Its kernel is the constant finite
abelian group scheme associated with a group $G$.
The pullback
\[
 P=T\times_{A,\pi}B
\]
is a finite \'{e}tale $G$-torsor.
Suppose it is disconnected.  Its connected components
are permuted transitively by $G$: each component maps
surjectively to the connected base $T$, and the action
on each geometric fibre is simply transitive.
Translate a chosen component so that it contains
$(t_0,0)$, and let $G_0\subsetneq G$ be its stabilizer.
The chosen component is then a $G_0$-torsor over $T$.
Its map to $B/G_0$ descends to a pointed morphism
\[
 b:T\longrightarrow B/G_0,\qquad \overline\pi b=a,
\]
where $\overline\pi:B/G_0\to A$ is an \'{e}tale isogeny
with nontrivial kernel $G/G_0$.

The target $B/G_0$ is an abelian variety.
The universal property used is precisely that of the
\emph{ordinary abelian} Albanese, so it supplies a
homomorphism $r:A\to B/G_0$ such that $ra=b$.
Uniqueness in that property gives
$\overline\pi r=\operatorname{id}_A$.
This is impossible for a nontrivial isogeny:
$r$ has trivial kernel and closed image of dimension
$\dim A=\dim(B/G_0)$, hence is surjective;
a point of $\ker\overline\pi$ then has the form $r(x)$,
and applying $\overline\pi$ forces $x=0$.
The nontrivial \'{e}tale group $G/G_0$ has nonzero
$k$-points, giving a contradiction.
Consequently $P$ is connected.
As $T$ and $P$ are smooth, they are also irreducible.

\emph{3. Separate the connected part of multiplication.}
Fix $n\geq1$ and let $F_n=A[n]^\circ$ be the connected
component of the identity in the finite group scheme
$A[n]$.  The cited group-scheme quotient results give
the factorization
\[
 A\xrightarrow{r_n}B_n=A/F_n
   \xrightarrow{\pi_n}A,\qquad
 \pi_n r_n=[n]_A.                                  \tag{CC.ALB.3}
\]
The quotient $B_n$ is a smooth connected proper
commutative group, hence an abelian variety.
The map $r_n$ is finite faithfully flat with geometrically
connected, single-point fibres, so it is a universal
homeomorphism.  The kernel of $\pi_n$ is the finite
\'{e}tale group scheme $A[n]/F_n$.  Thus $\pi_n$ is a
finite \'{e}tale isogeny.  These assertions also hold
when this kernel is trivial.
Step 2 shows that $T\times_A B_n$ is connected and
irreducible.

\emph{4. Keep this one cover connected on a curve.}
Choose a nonempty affine open $V\subset T$ and embed
it into affine space.  The restricted cover
$V\times_A B_n$ remains irreducible.  If $\dim V>1$,
apply Katz's Corollary 3.4.3 with $f=0$ and the closed
immersion $V\hookrightarrow\mathbb A^N$ as its
unramified morphism.  A choice in the nonempty
intersection of the two open parameter sets gives a
smooth irreducible affine hyperplane section whose
pullback cover is irreducible.  Iterate until obtaining
a smooth affine curve $C\subset V$ with
$C\times_A B_n$ connected.  If $\dim V=1$, take
$C=V$.  This is a finite sequence of choices for the
single specified $n$; no simultaneous choice for
infinitely many covers is claimed.

Let $\overline C$ be its smooth projective completion.
Properness of $A$ extends $a|_C$ over $\overline C$.
Choose $c_0\in C(k)$ and replace this map by
$a|_{\overline C}-a(c_0)$.  Translation by a lift of
$a(c_0)$ in $B_n(k)$ identifies the old and new
pullbacks, so connectedness is unchanged.
The universal property of the Jacobian gives
\[
 f:J=\operatorname{Jac}(\overline C)\longrightarrow A
\]
whose composite with $c\mapsto[c-c_0]$ is the
normalized map.  The finite \'{e}tale cover
$J\times_A B_n$ is connected as well.
Indeed, if it were disconnected, each of its
nonempty open-and-closed components would map
surjectively to the connected variety $J$.
Their pullbacks to $C$ would be nonempty disjoint
open-and-closed pieces, contradicting the preceding
connectedness.

\emph{5. Compute the components and obtain torsion
surjectivity.}
Consider
\[
 Q=J\times_{A,[n]_A}A
   =\{(x,z):f(x)=[n]_A z\}.
\]
The map $Q\to J\times_A B_n$ induced by $r_n$ is
the base change of the universal homeomorphism
$r_n$ along the projection to $B_n$.  Hence $Q$ is
connected, even though it may be nonreduced.
For each $t\in A[n](k)$, let $L_t$ be the image,
on underlying topological spaces, of the proper
morphism
\[
 J\longrightarrow Q,\qquad
 y\longmapsto\bigl([n]_J y,\ f(y)+t\bigr).           \tag{CC.ALB.4}
\]
Every $L_t$ is closed and connected.
These finitely many sets cover $Q$: given a
$k$-point $(x,z)$, choose $y$ with $[n]_J y=x$;
then $t=z-f(y)\in A[n](k)$.
This covers every closed point, hence all of the
finite-type $k$-scheme $Q$.
Moreover,
\[
 L_t=L_{t'}
 \quad\Longleftrightarrow\quad
 t-t'\in f(J[n](k)),                               \tag{CC.ALB.5}
\]
and any two intersecting $L_t,L_{t'}$ are equal.
For the latter assertion, an intersection contains
a closed point; two preimages $y,y'$ then have
$[n](y-y')=0$ and $f(y-y')=t'-t$.
Conversely translation by any such element of
$J[n](k)$ identifies the two images.
Thus the distinct $L_t$ are a finite disjoint closed
cover and are exactly the connected components.
Connectedness of $Q$ proves
\[
 f:J[n](k)\twoheadrightarrow A[n](k).               \tag{CC.ALB.6}
\]

The exact sequence (CC.ALB.2) for this curve gives a
surjection
$H^S_0(C,\mathbb Z)^0[n]\twoheadrightarrow J[n](k)$:
lift an $n$-torsion point of $J(k)$, then subtract
an $n$-division of the lift's $n$-multiple in the
torus kernel.  Functorial pushforward along $C\to T$
and (CC.ALB.6) prove (CC.ALB.1).
Translations introduced by the choice of $c_0$ vanish
on degree-zero cycles, so this is exactly the original
ordinary Albanese homomorphism.

\emph{6. Divide inside the kernel.}
Let $x\in K_T$ and let $n\geq1$.  By Step 1 choose
$y\in H_T$ with $ny=x$.
Then $\operatorname{alb}_T(y)\in A_T(k)[n]$.
By (CC.ALB.1) choose $z\in H_T[n]$ with the same
Albanese image.  The element $y-z$ belongs to $K_T$
and satisfies $n(y-z)=x$.  This proves divisibility
of $K_T$.
\end{proof}

\subsection{CC-ALB02: integral equality of projective and
quasiprojective parameter relations}

\paragraph{Statement.}
Let $X/k$ be smooth projective over an algebraically
closed field, of arbitrary characteristic.
For a fixed codimension $i$, let
$N^i_{\rm proj}(X)$ be generated in $CH^i(X)$ by
\[
 \Gamma_*v,\qquad
 \Gamma\in CH^i(X\times T),\quad
 v\in Z_0(T)^0,\quad \operatorname{alb}_T(v)=0,
                                                               \tag{CC.ALB.7}
\]
with $T$ smooth connected projective.
Define $N^i_{\rm qp}(X)$ in the same way with smooth
connected quasiprojective $T$, using its ordinary
abelian Albanese.  Then
\[
 N^i_{\rm proj}(X)=N^i_{\rm qp}(X).                 \tag{CC.ALB.8}
\]
The same subgroup results if the quasiprojective
definition uses the generalized semiabelian
Albanese kernel.

\begin{proof}
In characteristic zero this is the already proved
smooth compactification argument of
AQ-14-QUASIPROJECTIVE.  Suppose $\operatorname{char}k=p>0$.
Fix one parameter $T$ and one correspondence $\Gamma$.
The proper-target Suslin descent proved in
AQ-14-NONPROPER makes its action factor through
$H_T$.  We therefore have a homomorphism
\[
 \varphi_\Gamma:K_T\longrightarrow
 CH^i(X)/N^i_{\rm proj}(X).
\]

We make explicit the uniform quantifier in the
already proved alteration argument.
Choose once a projective compactification
$T\subset\overline T$ and a separable projective
alteration $S\to\overline T$ with $S$ smooth
projective, of degree $d>0$.
For example, choose any auxiliary prime $\ell\ne p$
and apply Illusie--Temkin (2014), Expos\'e X,
Theorem 3.5(iii), p.~202, as in AQ-14-ALTERATION.
There is a single nonempty open $U\subset T$ over
which the alteration is finite \'{e}tale of degree $d$.

For every $x\in K_T$, the moving proof in Step 1 of
AQ-14-ALTERATION gives a representative free
degree-zero cycle $v_U$ supported in this same $U$,
with Albanese sum zero and the same Suslin class.
The fibre-sum morphism
$U\to\operatorname{Alb}(S)$ extends across smooth
$T$ and factors, up to translation, through
$\operatorname{Alb}_{\rm ab}(T)$.
Consequently the pullback cycle $w=f_U^*v_U$ has
Albanese sum zero on $S$.
Closing the pullback of $\Gamma|_{X\times U}$
on $X\times S$ gives a correspondence $\Gamma'$
and the fibre Gysin calculation gives
\[
 \Gamma'_*w=d\,\Gamma_*v_U.
\]
The left side lies in $N^i_{\rm proj}(X)$.
Thus the \emph{same} positive integer $d$, fixed
before choosing $x$, satisfies
\[
 d\,\varphi_\Gamma(x)=0\qquad\text{for every }x\in K_T.
                                                               \tag{CC.ALB.9}
\]
Only the moved cycle varies with $x$; the alteration
and its degree do not.

By CC-ALB01, write $x=dy$ with $y\in K_T$.
Equation (CC.ALB.9) yields
$\varphi_\Gamma(x)=d\varphi_\Gamma(y)=0$.
This holds for every $T,\Gamma,x$, so every generator
of $N^i_{\rm qp}(X)$ belongs to $N^i_{\rm proj}(X)$.
The reverse inclusion is immediate, proving
(CC.ALB.8).

Finally the generalized semiabelian Albanese maps
to its ordinary abelian quotient.  Its kernel
relation is therefore contained in the relation
using the ordinary Albanese.  On projective
parameters the generalized Albanese is already
abelian.  The semiabelian parameter subgroup lies
between the two equal groups in (CC.ALB.8), so it
is equal to them.
\end{proof}

\paragraph{What the deduction settles.}
For the free-zero-cycle parameter convention actually
used in the cited cote 14 formulation, the possible
integral $p$-primary discrepancy between smooth
projective and smooth quasiprojective parameters
vanishes.  The earlier away-from-$p$ statement remains
a valid intermediate result.  No statement about
injectivity of the Suslin Albanese map on
$p$-primary torsion, integral motivic projectors,
or resolution of singularities is needed or inferred.

\paragraph{References and inspected locators.}
\begin{itemize}
\item Katz, N. M. (1993). Affine cohomological
transforms, perversity, and monodromy.
\emph{Journal of the American Mathematical Society,
6}(1), 149--222. Corollaries 3.4.2--3.4.3,
pp.~181--182.
\url{https://web.math.princeton.edu/~nmk/old/ACT.pdf}
\item Suslin, A., \& Voevodsky, V. (1996).
Singular homology of abstract algebraic varieties.
\emph{Inventiones Mathematicae, 123}, 61--94.
Theorem 3.1, p.~65, proof pp.~65--66.
\url{https://doi.org/10.1007/BF01232367}
\url{https://www.math.ias.edu/vladimir/sites/math.ias.edu.vladimir/files/singular_homology_published.pdf}
\item Milne, J. S. (2017). \emph{Algebraic groups:
The theory of group schemes of finite type over a
field}. Cambridge University Press.
Theorem 5.28 and Corollaries 5.26--5.27, p.~105;
Proposition 5.58, p.~114; 8.18 and 8.21, p.~152.
\url{https://www.jmilne.org/math/Books/iAG2017.pdf}
\item The Stacks Project Authors. (2026).
\emph{The Stacks Project}. Lemma 33.47.1,
Tag 0FD5, and Lemma 37.32.3, Tag 0G4F.
\url{https://stacks.math.columbia.edu/tag/0FD5}
\url{https://stacks.math.columbia.edu/tag/0G4F}
\item Spie{\ss}, M., \& Szamuely, T. (2000).
\emph{On the Albanese map for smooth quasi-projective
varieties} (arXiv:math/0009017v1). arXiv.
\S3, pp.~4--5, generalized Albanese universal property.
\url{https://arxiv.org/pdf/math/0009017v1}
\item Illusie, L., \& Temkin, M. (2014).
Expos\'e X: Gabber's modification theorem
(log smooth case).
In L. Illusie, Y. Laszlo, \& F. Orgogozo (Eds.),
\emph{Travaux de Gabber sur l'uniformisation locale
et la cohomologie \'{e}tale des sch\'emas
quasi-excellents} (Ast\'erisque 363--364,
pp.~167--212). Soci\'et\'e Math\'ematique de France.
Theorem 3.5(iii), p.~202; the application and
separability were already inspected in
AQ-14-ALTERATION.
\url{https://www.numdam.org/item/AST_2014__363-364__167_0.pdf}
\end{itemize}

% END consequences_criteria_integral_albanese.tex

% BEGIN root_separated_parameters.tex
\section{Removing quasiprojectivity from the cote 14 parameters}

\subsection{BR-14-SEP: reduction to a dense affine parameter}

\paragraph{Original scope.}
Grothendieck's Definition~1.1 in cote~14, physical PDF page~24,
footer~23, places the word \emph{projectif} in brackets before
\emph{lisse connexe}.  The preceding results identify the relation
with or without this qualification when the nonproper parameter is
quasiprojective.  The following deduction removes that remaining
restriction on the parameter scheme.  The algebraically closed
base-field convention is still explicit; it is not attributed to
the handwritten definition.

\paragraph{Statement.}
Let \(k\) be algebraically closed, and let \(T\) be a nonempty smooth
connected separated scheme of finite type over \(k\).  For every
nonempty affine open \(U\subset T\), the ordinary abelian Albanese
of \(T\) exists, and the inclusion induces an isomorphism
\[
              \operatorname{Alb}_{\rm ab}(U)
                    \xrightarrow{\ \sim\ }
              \operatorname{Alb}_{\rm ab}(T).                 \tag{BR.14.1}
\]
Base points may be chosen in \(U(k)\).  Every zero cycle on \(T\)
has a representative of the same degree supported in \(U\), modulo
explicit Suslin one-boundaries.  For \(V=U,T\), write
\(H_V=H^S_0(V,\mathbb Z)^0\) and
\(K_V=\ker(H_V\to\operatorname{Alb}_{\rm ab}(V)(k))\).
Consequently the maps
\[
                    H_U\longrightarrow H_T,\qquad
                    K_U\longrightarrow K_T                  \tag{BR.14.2}
\]
are surjective.  In positive characteristic, CC-ALB01 therefore
gives divisibility of both groups on the right and surjectivity
\[
 H^S_0(T,\mathbb Z)^0[n]\longrightarrow
                    \operatorname{Alb}_{\rm ab}(T)(k)[n]
              \qquad(n\ge1).                                 \tag{BR.14.3}
\]

Let \(X/k\) be smooth projective and geometrically connected.
Define \(N^i_{\rm sep}(X)\) using all such smooth separated
parameters \(T\), all \(\Gamma\in CH^i(T\times X)\), and free
degree-zero zero cycles whose ordinary Albanese sum is zero.
Let \(N^i_{\rm aff}(X)\) be the subgroup using only smooth
connected affine parameters.  Then, integrally and in every
characteristic,
\[
 N^i_{\rm aff}(X)=N^i_{\rm qp}(X)
                =N^i_{\rm sep}(X)=N^i_{\rm proj}(X).           \tag{BR.14.4}
\]
Thus allowing arbitrary smooth separated finite-type parameters
does not enlarge the relation or change the quotient in the
original Theorem~3.1, physical page~38/footer~37.

\paragraph{Exact external inputs.}
Milne (2008), Chapter~I, Theorem~3.2, printed p.~17, extends a
rational map from a nonsingular variety to an abelian variety over
the whole source.  His convention on printed p.~v allows every
geometrically reduced separated finite-type scheme, without
quasiprojectivity.  Lemma~3.6, printed p.~18, says that the union
of the integral curves through a smooth point and nonsingular
there is dense.  Lemma~3.5 on the same page supplies normalization
and smooth completion of curves.  The Gysin maps used below are
the factorization-independent maps of the Stacks Project,
Lemma~42.59.3 (Tag~0FF2) and Definition~42.59.4 (Tag~0FF3);
their functoriality is Lemma~42.59.6 (Tag~0FF5).  Their construction
requires a smooth--immersion factorization, which is supplied
here by the graph of a morphism between smooth separated schemes.

\begin{proof}
\emph{1. The ordinary Albanese is unchanged by the dense open.}
Smoothness and connectedness make \(T\) integral: regular local
rings are domains, so its distinct irreducible components are
disjoint and open.  Thus \(U\) is dense and integral.  Choose
\(u_0\in U(k)\), and write
\(a_U:U\to A_U=\operatorname{Alb}_{\rm ab}(U)\) for the pointed
ordinary Albanese, which exists for the smooth affine variety
\(U\).  Regard \(a_U\) as a rational map on \(T\).  The cited
extension theorem gives a morphism \(a_T:T\to A_U\), still
pointed at \(u_0\).

If \(f:T\to B\) is a pointed morphism to any abelian variety,
its restriction to \(U\) factors through a unique homomorphism
\(h:A_U\to B\).  The two morphisms \(f\) and \(h a_T\)
agree on the dense open \(U\); since \(T\) is reduced and
\(B\) is separated, they agree on \(T\).  Uniqueness follows
from the same universal property on \(U\).  Hence \((A_U,a_T)\)
is the ordinary Albanese of \(T\), proving (BR.14.1).

\emph{2. Moving one point using a curve in an affine neighbourhood.}
It suffices to move \(x\in T(k)\setminus U(k)\).  Choose an
affine neighbourhood \(V\) of \(x\).  The open \(V\cap U\)
is nonempty, because \(T\) is integral.  By the cited curve lemma
there is an integral curve in \(V\), nonsingular at \(x\),
which meets \(U\).  Normalize it and call the resulting smooth
curve \(C\), with its morphism \(\nu:C\to T\) and point
\(x'\) above \(x\).  In particular \(k(x')=k\).  Let
\(\overline C\) be its smooth projective completion.  Put
\[
 D=\overline C\setminus C,\qquad
 F=C\setminus\nu^{-1}(U).
\]
These are finite sets.  Choose \(y\in\nu^{-1}(U)(k)\).
For a sufficiently large integer \(N\), the first cohomology of
\[
 \mathcal O_{\overline C}
 \left(Ny-2x'-\sum_{z\in D\cup(F\setminus\{x'\})}z\right)
\]
vanishes: its degree exceeds \(2g(\overline C)-2\).
The resulting surjective evaluation map gives a rational
function \(f\) with poles only at \(y\), first-order jet a
uniformizer at \(x'\), and value \(1\) at every point of
\(D\cup(F\setminus\{x'\})\).  Thus
\[
       \operatorname{div}(f)=[x']+z-m[y],\qquad
       \operatorname{supp}(z)\subset\nu^{-1}(U),
       \qquad \deg z=m-1.                                  \tag{BR.14.5}
\]
Here \(z\) is effective and \(m>0\): \(f\) is nonconstant,
and it cannot have no poles on the projective curve.

The morphism \(g=f/(f-1):\overline C\to\mathbb P^1\) is
finite.  Its inverse image of \(\mathbb A^1\) lies in \(C\)
because \(f=1\) on \(D\).  Its graph is therefore a finite
correspondence from \(\mathbb A^1\) to \(C\), and the
difference of its fibres at \(0\) and \(1\) is
\(\operatorname{div}(f)\), up to the choice of boundary sign.
Its image under \(\nu\) is a valid finite correspondence to
\(T\).  Indeed, put \(W=g^{-1}(\mathbb A^1)\).  Its graph
\(W\to W\times T\) is closed because \(T\) is separated,
and \(W\times T\to\mathbb A^1\times T\) is finite.
Their composite is finite.  Pushing (BR.14.5) forward shows that \([x]\)
is equal in Suslin homology to \(m[\nu(y)]-\nu_*z\), a
degree-one cycle supported in \(U\).
Linearity proves the moving assertion for every zero cycle,
with its degree unchanged.  No common quasiprojective
neighbourhood of all points of that cycle was used.

\emph{3. Albanese and correspondence values survive this move.}
The morphism \(a_T\nu:C\to A_U\) extends to
\(\overline C\), since the target is proper and the source
is a smooth curve.  A principal divisor on \(\overline C\)
has zero image under its induced map from the Jacobian to
\(A_U\).  Thus (BR.14.5) preserves the ordinary Albanese sum.

For completeness, fix \(\Gamma\in CH^i(T\times X)\).
Pull it back along \(\nu\times1_X\) using the Gysin map;
this exists because both schemes are smooth and separated,
and the graph gives the required factorization.  Choose a
cycle representative on \(C\times X\), and close its
components on \(\overline C\times X\), giving a class
\(\overline\Gamma\).  Its restriction is the chosen pullback.
The usual product with the pullback of
\(\operatorname{div}(f)\) is zero in the Chow group of
\(\overline C\times X\), since a principal divisor is zero
in \(CH^1(\overline C)\).  Push this equality to \(X\);
this projection is proper.  Gysin functoriality and restriction
to the open \(C\) identify its terms with
\[
       \Gamma_{\nu(x')}+
       \sum_{c\in\operatorname{supp}(z)}\operatorname{mult}_c(z)
                   \Gamma_{\nu(c)}-m\Gamma_{\nu(y)}=0.
                                                               \tag{BR.14.6}
\]
Thus the correspondence action also preserves the move.
This argument only uses curves and the projective target;
it does not require a compactification of \(T\).

The Albanese sum descends to Suslin homology by the same
proper-curve argument for an arbitrary Suslin one-chain:
normalize its finite support over \(\mathbb A^1\), complete
the resulting curve, and use the function \(t/(t-1)\).
At each point above infinity this function has value \(1\),
so its divisor is precisely the difference of the finite
fibres.  Step~2 now proves the first surjection in
(BR.14.2).  By (BR.14.1), a lift of a class with zero
Albanese sum has zero Albanese sum already on \(U\).
This proves the second surjection.

\emph{4. Divisibility and equality of the parameter relations.}
In positive characteristic CC-ALB01 applies to the affine
\(U\).  Its two divisible groups map onto those for \(T\),
and a quotient of a divisible abelian group is divisible.
For any \(a\in A_U(k)[n]\), its lift in
\(H^S_0(U,\mathbb Z)^0[n]\) maps to an \(n\)-torsion lift
in \(H^S_0(T,\mathbb Z)^0\).  This proves (BR.14.3).

Now let \(\Gamma_*v\) be a generator of
\(N^i_{\rm sep}(X)\).  The preceding steps replace \(v\)
by a cycle \(v_U\) supported in \(U\), of degree zero,
with zero Albanese sum in \(A_U\), and satisfying
\[
        \Gamma_*v=(\Gamma|_{U\times X})_*v_U.
                                                               \tag{BR.14.7}
\]
The right side belongs to \(N^i_{\rm aff}(X)\).  The reverse
inclusion follows because affine parameters are allowed among
the separated parameters.  Hence
\(N^i_{\rm sep}=N^i_{\rm aff}\).  The evident inclusions
\(N^i_{\rm aff}\subset N^i_{\rm qp}\subset N^i_{\rm sep}\)
give their equality.  The already proved equality
\(N^i_{\rm qp}=N^i_{\rm proj}\), namely
AQ-14-QUASIPROJECTIVE in characteristic zero and CC-ALB02
in positive characteristic, completes (BR.14.4).
\end{proof}

\paragraph{Semiabelian convention and limits.}
If the relation instead uses the kernel of a universal
semiabelian Albanese, every such cycle has zero ordinary
abelian Albanese sum.  Its relation subgroup lies between
the projective-parameter subgroup and \(N^i_{\rm sep}\);
(BR.14.4) makes these endpoints equal.  This observation
does not assert that the semiabelian Albanese itself is
unchanged by removing a divisor.  Nothing here extends
the equality to arbitrary nonclosed base fields, singular
parameters, or an unspecified integral motivic projector.

\paragraph{References.}
Milne, J.~S. (2008, March~16). \emph{Abelian varieties}
(Version~2.0) [Author course notes].
\url{https://www.jmilne.org/math/CourseNotes/AV.pdf}.
Conventions, p.~v; Chapter~I, Theorem~3.2, p.~17,
and Lemmas~3.5--3.6, p.~18.
The Stacks Project Authors. (2026). \emph{The Stacks Project}.
Section~42.59, especially Lemma~42.59.3 (0FF2),
Definition~42.59.4 (0FF3), and Lemma~42.59.6 (0FF5):
\url{https://stacks.math.columbia.edu/tag/0FEZ}.

% END root_separated_parameters.tex

% BEGIN consequences_criteria_albanese_literature.tex
\subsection{CC-ALB-LIT01: exact scope of the nearby Albanese literature}

This source note accompanies CC-ALB01--02 and does not modify their
checked proof. It records overlapping results without making a claim
of priority or novelty. The distinction between the ordinary abelian
Albanese and the generalized semiabelian Albanese is retained throughout.

\paragraph{Binda--Krishna.}
In the inspected author version, arXiv:2202.06808v2, dated 1 September
2022, Section 1.1 defines a reduced modulus pair \((X,D)\) with
\(X\) smooth projective and \(D\) a reduced effective Cartier divisor;
\(U=X\setminus D\). Theorem 1.4, p.~3, with its complete short proof in
Section 3.6, p.~19, identifies the torsion in integral Suslin homology
with the torsion in the generalized semiabelian Albanese over an
algebraically closed field. Thus this theorem includes characteristic
\(p\), but retains the smooth projective compactification hypothesis.
Its characteristic-\(p\) proof uses the modulus comparison in Theorem
1.1 and an earlier torsion theorem for cycles with modulus.

Theorem 1.7, p.~5, restated as Theorem 5.1, p.~25, concerns a different
map: for a reduced projective scheme of pure dimension \(d\) over an
algebraically closed field of characteristic zero, the map from the
lci Chow group of zero-cycles to \(H^{2d}(X,\mathbb Z(d))\) is
surjective with uniquely divisible kernel. It is not a theorem about
the ordinary Albanese kernel of every smooth quasiprojective scheme.
The paper appeared in \emph{Transactions of the American Mathematical
Society} in 2023.

\paragraph{Ghosh--Krishna.}
Theorem 1.7, p.~4, of the inspected arXiv:2012.11249v3, dated 14 January
2022, gives the full torsion isomorphism for smooth quasiprojective
\(U\) over an algebraically closed field under the explicit assumption
that \(U\) admits an open immersion into a smooth projective variety.
The phrase \emph{without resolution} in the surrounding discussion does
not eliminate this stated compactification hypothesis.
Theorem 1.6, on the same page, instead compares the Levine--Weibel
zero-cycle group of a projective variety regular in codimension one
with Suslin homology of its regular locus. Its scope is that specified
regular locus, and does not identify arbitrary smooth open subsets
with that locus. The paper appeared in \emph{Journal of the Institute
of Mathematics of Jussieu} in 2023.

\paragraph{Krishna--Pel\'aez: version-sensitive citation.}
Theorem 9.6, p.~37, occurs in arXiv:1606.05810v1, dated 18 June 2016.
It assumes that a smooth quasiprojective \(U\) of dimension \(d\) over
an algebraically closed field has a normal projective compactification
\(X\) with \(\dim(X\setminus U)\le d-2\). Under that hypothesis it
states a torsion isomorphism for the Albanese map and that the relevant
semiabelian Albanese is abelian. The statement and its short proof
have been inspected; the proof's earlier Theorem 9.3 has not been
independently reconstructed here.

The final arXiv version, 1606.05810v4 dated 1 June 2018, and the
published \emph{Annals of K-Theory} article have eight sections and do
not contain this Theorem 9.6 or its Section 9. Accordingly, this note
does not treat the v1 assertion as a theorem of the published paper.
No reason for the omission, and no assertion that the earlier theorem
is false, is inferred from the version comparison.

\paragraph{Relation to the checked deduction.}
CC-ALB01 proves divisibility of the ordinary abelian Albanese kernel
for every smooth connected quasiprojective variety over an
algebraically closed field of positive characteristic. It uses a
curve calculation, connected pullbacks of a single finite \'etale
isogeny, and finite group schemes; it does not assert injectivity on
\(p\)-primary torsion. CC-ALB02 then combines that divisibility with
the already checked uniform-degree alteration transfer to identify
the projective and quasiprojective parameter relations integrally.
The exact modern results inspected above have different hypotheses
or different maps. This comparison records that difference only;
it is not an exhaustive priority search.

\paragraph{Bibliography and inspected versions.}
Binda, F., \& Krishna, A. (2023).
\emph{Suslin homology via cycles with modulus and applications}.
Transactions of the American Mathematical Society, 376(2), 1445--1473.
\url{https://doi.org/10.1090/tran/8815}.
Inspected author version:
\url{https://arxiv.org/pdf/2202.06808v2}.

Ghosh, M., \& Krishna, A. (2023).
\emph{Zero-cycles on normal projective varieties}.
Journal of the Institute of Mathematics of Jussieu, 22(5), 2241--2295.
\url{https://doi.org/10.1017/S1474748022000032}.
Inspected author version:
\url{https://arxiv.org/pdf/2012.11249v3}.

Krishna, A., \& Pel\'aez, P. (2018).
\emph{The slice spectral sequence for singular schemes and applications}.
Annals of K-Theory, 3(4), 657--708.
\url{https://doi.org/10.2140/akt.2018.3.657}.
Published text:
\url{https://msp.org/akt/2018/3-4/akt-v3-n4-p04-s.pdf}.
Earlier version containing Theorem 9.6:
\url{https://arxiv.org/pdf/1606.05810v1}.
Final author version:
\url{https://arxiv.org/pdf/1606.05810v4}.

% END consequences_criteria_albanese_literature.tex

% BEGIN consequences_criteria_Tate2026.tex
\section{Current primary inputs on CM Hodge, abelian Tate, and the
standard
conjectures}\label{current-primary-inputs-on-cm-hodge-abelian-tate-and-the-standard-conjectures}

Research module \nolinkurl{consequences_criteria}, checked 20--21 November
2026 UTC (literature cutoff: 20 November), GPT-6 Astra Pro; effort not
exposed. This module updates the archive analysis using identified 2026
primary manuscripts. Their major theorems are \textbf{cited literature
inputs}, not theorems independently reproved in this run. The adopted
Alpöge RH and Dirichlet GRH hypotheses play no role in the deductions
below.

\subsection{CC-TATE01. The current all-CM input and its exact
provenance}\label{cc-tate01.-the-current-all-cm-input-and-its-exact-provenance}

\textbf{Cited 2026 preprint theorem.} For every complex abelian variety
of CM type, every rational Hodge class in every codimension is
algebraic. The precise locator is van Hoften and Sempliner, \emph{The
Hodge Conjecture for CM abelian varieties}, \textbf{Theorem I, p.~1},
Hexagon \textbf{2610.00259v2}. This includes powers and products of CM
abelian varieties, because those varieties are themselves of CM type.

The version history dates v1 to 15 October and v2 to \textbf{18 October
2026}. The current record's header displays the original 15 October
date; the downloaded v2 PDF creation and modification timestamps both
read 18 October 2026, 10:11:24 UTC. Thus the edition used here is
specified by the immutable v2 identifier, not merely by the header date.
The named authors are Pol van Hoften and Jack Sempliner. Their §1.3,
p.~4, explains that they give a variant proof inspired by the OpenAI
manuscript and do not certify every detail of that earlier manuscript.
Our cited input is their own Theorem I. Neither a journal publication
nor a complete independent verification of their proof is claimed here.

\textbf{Dependency inspection.} The reduction is Theorem 2.2.1, p.~6;
the algebraic surface criterion is Theorem 2.3.4, pp.~7--8; the required
nonzero cup product is Theorem 4.0.2, p.~12. We read the
surface-criterion proof completely, the cited Albanese statement Theorem
3.1.11 and Corollary 3.1.13, pp.~11--12, and the final argument §4.2.8,
pp.~19--22, including the diagonal-period Lemma 4.2.7, pp.~17--18, and
the archimedean nonvanishing Lemma 4.2.9, pp.~20--21. Those analytic
lemmas retain their cited theta-correspondence inputs; these inputs were
not independently reconstructed. The reduction's reference to Hazama is
recorded as a dependency, including the authors' p.~6 footnote about the
formulation of Hazama's Lemmas 7.11--7.12. This is a dependency read,
not a claim to have checked every underlying paper.

Milne's \textbf{9 November 2026} note, \emph{On the Hodge, Tate, and
standard conjectures for abelian varieties}, is linked from his own
mathematics homepage and available there as
\nolinkurl{math/xnotes/HTS.pdf}. Its Hexagon deposit is
\textbf{2611.00244v1}, submitted 9 November and displayed as published
10 November. Milne states the CM result as Theorem 2.1, p.~4, and
explicitly treats the formerly conditional abelian consequences as
unconditional in §3, p.~6. We read the ten-page note. It is useful
corroboration and context; the numbered primary implications below are
also checked in the original papers.

\subsection{CC-TATE02. Exact abelian consequences of the cited CM
theorem}\label{cc-tate02.-exact-abelian-consequences-of-the-cited-cm-theorem}

Let \(A/\mathbb F_q\) be an abelian variety of arbitrary dimension, let
\(\ell\ne\operatorname{char}\mathbb F_q\), and let
\(N^r(A)_{\mathbb Q}\) denote codimension-\(r\) cycles with rational
coefficients modulo numerical equivalence. The resulting finite-field
statement is

\[
N^r(A)_{\mathbb Q}\otimes_{\mathbb Q}\mathbb Q_\ell
\xrightarrow{\ \sim\ }
H^{2r}_{\mathrm{et}}(A_{\overline{\mathbb F}_q},\mathbb Q_\ell(r))^{G_{\mathbb F_q}}
\qquad\text{for every }r.
\]

In particular, the order of the pole of \(Z(A,t)\) at \(t=q^{-r}\) is
\(\dim_{\mathbb Q}N^r(A)_{\mathbb Q}\). These assertions concern
rational coefficients.

\textbf{Exact original implication.} Milne (1999), \emph{Lefschetz
motives and the Tate conjecture}, \textbf{Theorem 7.1, journal p.~72},
assumes the Hodge conjecture for all complex CM abelian varieties and
proves his Tate statement (0.1) for all abelian varieties over
\(\overline{\mathbb F}_p\). Statement (0.1), \textbf{journal p.~46},
includes both numerical versus homological equivalence and spanning of
the Tate subspace. The following paragraph explicitly gives the
corresponding assertion for each finite-field model and the zeta pole
order. The theorem's proof is organized through Propositions \textbf{7.3
and 7.4, pp.~73--75}. These statements and the complete proofs of these
two propositions were read. The auxiliary theory in §§1--6 is used as
cited literature, not independently reproved. The journal PDF, whose
pagination is 45--76, was inspected; the author's separate erratum
contains an inaccurate bibliographic range 47--81, which is not
propagated here.

\textbf{Finite-field descent check.} The passage from geometric Tate
classes to classes over a given finite field can also be made explicit.
Write an invariant class as a finite \(\mathbb Q_\ell\)-linear
combination of classes of cycles over one common finite extension
\(\mathbb F_{q^m}\). Push each cycle down to \(A\) and divide by \(m\).
On geometric cohomology this is the averaging operator

\[
\frac1m\sum_{j=0}^{m-1}F^j.
\]

It fixes the given \(F\)-invariant class. Thus no enlargement of the
stated ground field is required in the displayed surjectivity assertion.
The numerical-kernel assertion comes from Milne's full Tate statement.
This argument is a formal verification of the field descent, not a new
proof of the Tate input.

\textbf{Standard conjectures.} For every algebraically closed field
\(k\) and every prime \(\ell\ne\operatorname{char}k\), numerical and
\(\ell\)-adic homological equivalence coincide on every abelian variety
over \(k\). For a polarized abelian variety of dimension \(g\), the
primitive rational algebraic classes in codimension \(r\le g/2\) carry
the positive-definite pairing

\[
(a,b)\longmapsto(-1)^r\deg\bigl(a\cdot b\cdot L^{g-2r}\bigr).
\]

The exact inspected locator is Milne (2002), \textbf{Theorem 3.3(a),(b),
pp.~8--9 of the final author manuscript dated 19 September 2001},
\nolinkurl{2002aP.pdf}. Its hypothesis is precisely the CM Hodge theorem in
CC-TATE01. We read the theorem and its full proof, including the
specialization argument to arbitrary algebraically closed \(k\).
\textbf{Corollary 3.4, p.~9}, gives the corresponding polarization of
the category of abelian numerical motives.

For cycles defined over an arbitrary field \(k\), these comparisons
follow by base change to \(\bar k\). Indeed, a complementary cycle over
\(\bar k\) is defined over a finite extension \(k'/k\); the degree
pairing with a cycle defined over \(k\) can be computed using its
pushforward from \(k'\), up to the nonzero extension degree. Thus
numerical triviality over \(k\) is equivalent to numerical triviality
after algebraic closure, with rational coefficients. The \(\ell\)-adic
homological condition is geometric by definition. Positivity on the
rational classes over \(k\) is the restriction of the geometric
positive-definite form. This gives the stated abelian comparison and
positivity over arbitrary fields, for \(\ell\) invertible in the field.

Two source precisions matter. First, van Hoften--Sempliner's §1.1.2
cites \textbf{Theorem 4.6}, the earlier long-version numbering; it
should not be silently represented as the number in the final
manuscript. Second, Milne's abbreviated 2026 display in §3.6 omits the
factor \((-1)^r\); the correct signed form is explicit in the 2002
paper's introduction, p.~2, and §3, p.~6. We use that signed form. The
final manuscript's \textbf{Remark 3.5, p.~9}, concerns an arbitrary
\emph{good} Weil cohomology, where ``good'' means that homological and
numerical equivalence agree on all varieties in the chosen
product-closed class. It does not remove that condition for an
unspecified Weil cohomology theory.

Milne's 2026 note states the finite-field conclusion as \textbf{Theorem
3.1, p.~6}, and the Hodge-standard conclusion as \textbf{Theorem 3.5,
p.~8}. The chain is an application of the cited 2026 theorem to old
conditional theorems. It does not use RH, Dirichlet GRH, a number-field
generalized RH, or an unproved automorphy assertion.

\subsection{CC-TATE03. Rational specialization: a separate 2026
input}\label{cc-tate03.-rational-specialization-a-separate-2026-input}

The Hexagon record \textbf{2610.00226v1}, contributed by \textbf{Keerthi
Madapusi}, is dated \textbf{13 October 2026}. Its actual PDF title is
\emph{Milne's rationality conjecture}, with the byline \textbf{``ChatGPT
and Claude''}; §1.1, p.~3, identifies Madapusi's questions and feedback.
Milne's November note cites this source under Madapusi's name and a
longer descriptive title. The bibliography below records the PDF
attribution and the contributor explicitly.

\textbf{Theorem 1.1, p.~1}, concerns an abelian variety
\(A/\overline{\mathbb Q}\) with good reduction
\(A_0/\overline{\mathbb F}_p\) at a fixed place, a rational Hodge class
\(\gamma\) of codimension \(r\), and a rational Lefschetz class
\(\delta\) of complementary codimension on \(A_0\). The traces

\[
\operatorname{Tr}_v(\gamma_{0,v}\delta_v)
\]

are the images of one rational number for all \(v=\ell\ne p\) and for
the crystalline realization \(v=p\), with the specified coefficient
extension in the latter. Theorem 5.1, pp.~11--12, supplies the
specialized Weil-class dichotomy; §5.1 deduces the main statement. The
lifting input is recorded precisely in \textbf{Proposition A.1,
pp.~13--14}, which invokes Kisin--Zhou, including a quasi-split
replacement and marked tensor transport. These are inspected source
hypotheses and dependency locators; this run does not claim an
independent proof of their lifting arguments.

Combining CC-TATE01 with this manuscript's \textbf{Corollary 6.2,
p.~13}, gives a single cycle \(z_\gamma\in CH^r(A_0)_{\mathbb Q}\)
representing every specialized realization of \(\gamma\). The
corollary's proof was read: CM transport gives
\(\gamma_0=f_0^*\kappa_0\); algebraicity of the CM class \(\kappa\)
permits specialization of a representing cycle; for an actual
homomorphism \(h=Nf_0\), the cycle is \(N^{-2r}h^*Z_0\). This is
algebraicity on the \textbf{special fibre}. It does not assert that
\(\gamma\) is algebraic on a general non-CM characteristic-zero \(A\).

\subsection{CC-TATE04. What this resolves in cote 15, and what it leaves
to
prove}\label{cc-tate04.-what-this-resolves-in-cote-15-and-what-it-leaves-to-prove}

The supplied transcription \nolinkurl{ARCHIVE_15_scan100.tex} asks for
sufficient arithmetic conditions for a Frobenius eigenvalue. Scans 102,
104, and 106 then seek a motivic Galois description and its geometric
version. These are distinct mathematical assertions.

\begin{longtable}[]{@{}
  >{\raggedright\arraybackslash}p{(\columnwidth - 4\tabcolsep) * \real{0.3333}}
  >{\raggedright\arraybackslash}p{(\columnwidth - 4\tabcolsep) * \real{0.3333}}
  >{\raggedright\arraybackslash}p{(\columnwidth - 4\tabcolsep) * \real{0.3333}}@{}}
\toprule\noalign{}
\begin{minipage}[b]{\linewidth}\raggedright
Archive issue
\end{minipage} & \begin{minipage}[b]{\linewidth}\raggedright
Consequence supported by the inspected chain
\end{minipage} & \begin{minipage}[b]{\linewidth}\raggedright
Remaining condition
\end{minipage} \\
\midrule\noalign{}
\endhead
\bottomrule\noalign{}
\endlastfoot
Existence of an individual Weil-integer eigenvalue & The Honda--Tate
construction already supplies the occurrence result analyzed separately
in the AQ module. & An occurrence statement does not identify all
correspondences on another variety with that eigenvalue. \\
Tate classes on abelian varieties and their products over finite fields
& The finite-field rational cycle isomorphism in CC-TATE02 holds in
every dimension and codimension. & The 2026 CM theorem is a cited
primary input, not an independently reproved theorem of this report. \\
A smooth projective variety with a motive of abelian type & Milne
(1999), \textbf{Remark 7.6(b), journal p.~76}, explicitly extends Tate
to a variety whose motive \textbf{modulo homological equivalence} lies
in the Tannakian subcategory generated by abelian varieties; products of
curves are an example. & One must establish this actual motivic
membership for the particular variety. Matching Frobenius characteristic
polynomials alone does not establish it. \\
The geometric motivic group in the abelian category & Milne's
\textbf{Proposition 7.3, pp.~73--74}, gives \(M=P\) for the category of
abelian numerical motives over \(\overline{\mathbb F}_p\), with \(P\)
his Weil-number protorus. & This is a statement about the abelian
category and its Tannakian group, not a neutralization of every motivic
gerbe over \(\mathbb Q\). \\
Primitive positivity and the numerical/homological comparison &
CC-TATE02 discharges these conditions for abelian varieties and the
indicated realizations. & No inspected theorem here proves these
statements for an arbitrary smooth projective variety in positive
characteristic. \\
Integral, mixed, or global analytic strengthening & No such
strengthening follows from the rational pure statements just listed. &
Integral comparison, extension data, and number-field analytic
continuation/zeros require their own hypotheses and constructions. \\
\end{longtable}

\textbf{The group-theoretic application in full.} Let \(M\) be the
fundamental group used in §7 of Milne (1999) for abelian numerical
motives over \(\overline{\mathbb F}_p\). The pre-existing construction
gives \(P\subseteq M\subseteq L\). The CM Hodge theorem identifies the
CM Hodge correspondences with algebraic correspondences, so
specialization gives the additional inclusion \(M\subseteq S\) inside
the common Lefschetz group \(T\). Milne's intersection theorem gives
\(S\cap L=P\). Therefore

\[
P\subseteq M\subseteq S\cap L=P,
\]

and hence \(M=P\). This is exactly the known-result application in
Proposition 7.3; it neither assumes nor proves that all smooth
projective motives belong to this category. Milne defines \(P\) by the
Galois module of germs of Weil numbers in §4, \textbf{author-manuscript
p.~14}. These germs use the power transition maps as the finite field
grows. A convention for homological versus cohomological weights must be
retained when comparing his notation with the archive.

\textbf{Gerbe precision.} Milne's separate \textbf{1999b
addendum/erratum, p.~1}, retracts an optimistic introductory claim that
an intersection of Tannakian groups alone automatically constructs the
desired quotient category. Its stated issue is the nonneutral setting.
It does not retract Theorem 7.1; its §7, pp.~3--4, even supplies
Deligne's alternative proof route. Thus the present all-CM theorem
discharges the hypotheses of the inspected abelian Tate theorem, while
an asserted canonical rational fibre functor or the archive's literal
displayed product decomposition must still be justified on its own
terms.

The AQ module makes both sides of this distinction concrete.

\textbf{AQ-ABELIAN-RECOGNITION} proves that full Frobenius determinants
recognize actual rational numerical objects in its explicitly defined
abelian--Artin category, whereas \textbf{AQ-RATIONAL-FIBRE} gives an
elliptic quaternionic obstruction to a rational fibre functor and to
discarding eigenvalue multiplicities. Both complete proofs were
independently read in this run and passed on their stated inputs.

The safe global status statement is therefore precise: \textbf{the
inspected 2026 literature resolves the abelian Tate/standard inputs and
the stated abelian-type Tate input; the unrestricted archive
identification is not established by those inputs.} No assertion that
every broader version is currently open is needed for this conclusion.

\subsection{CC-TATE05. Fourfold confirmation and revision
control}\label{cc-tate05.-fourfold-confirmation-and-revision-control}

Broe's \emph{The Tate conjecture for abelian fourfolds over finite
fields}, \textbf{arXiv:2608.28651v1, 19 August 2026}, gives
\textbf{Theorem 1.1, p.~1}, for every abelian fourfold over a finite
field and \textbf{Corollary 1.2, p.~2}, for standard D on abelian
fourfolds over arbitrary fields, with \(\ell\) invertible. Corollary
8.2, p.~20, repeats the latter with its proof. The complete argument in
§§7--8 was read, with earlier reduction and lifting results retained as
cited inputs. Its \textbf{Remark 1.3, p.~2}, distinguishes positivity on
the numerical quotient from Grothendieck's original homological
formulation. This verifies and strengthens the previously uninspected
Broe entry in SA22; the later all-CM chain makes the abelian dimension
restriction unnecessary for the Tate/standard conclusions used here.

Ningyi Li's separate \textbf{arXiv:2609.27916v2} is explicitly
\textbf{withdrawn on 8 October 2026}. Its abstract must not be used as a
current theorem citation. The withdrawal does not concern the distinct
fivefold manuscript arXiv:2609.06265v2. That manuscript was examined
only selectively and is not required as an input anywhere in this
module.

\subsection{Bibliography and stable primary
URLs}\label{bibliography-and-stable-primary-urls}

Broe, M. (2026, August 19). \emph{The Tate conjecture for abelian
fourfolds over finite fields} (arXiv:2608.28651v1) {[}Preprint{]}.
\url{https://arxiv.org/pdf/2608.28651v1}

ChatGPT, \& Claude. (2026, October 13). \emph{Milne's rationality
conjecture} (Hexagon:2610.00226v1) {[}Preprint contributed by Keerthi
Madapusi; authorship and supervision explained in §1.1{]}.
\url{https://hexagonmath.org/pdf/2610.00226v1}

Li, N. (2026, October 8). \emph{The Tate conjecture for powers of
abelian fourfolds and the Hodge conjecture for powers of CM fourfolds}
(arXiv:2609.27916v2) {[}Withdrawn preprint record{]}.
\url{https://arxiv.org/abs/2609.27916}

Milne, J. S. (1999). Lefschetz motives and the Tate conjecture.
\emph{Compositio Mathematica, 117}(1), 45--76.
\url{https://doi.org/10.1023/A:1000776613765}

\begin{itemize}
\tightlist
\item
  Inspected journal PDF:
  \url{https://www.cambridge.org/core/services/aop-cambridge-core/content/view/962D833FAAC3219F55A58B193323B52A/S0010437X99000317a.pdf/lefschetz-motives-and-the-tate-conjecture.pdf}
\item
  Author manuscript: \url{https://jmilne.org/math/articles/1999bP.pdf}
\item
  Author addendum/erratum:
  \url{https://jmilne.org/math/articles/add/1999b.pdf}
\end{itemize}

Milne, J. S. (2002). Polarizations and Grothendieck's standard
conjectures. \emph{Annals of Mathematics, 155}(2), 599--610.
\url{https://doi.org/10.2307/3062126}

\begin{itemize}
\tightlist
\item
  Inspected final author manuscript, dated September 19, 2001:
  \url{https://www.jmilne.org/math/articles/2002aP.pdf}
\item
  Publisher metadata: \url{https://annals.math.princeton.edu/2002/155-2/p07}
\item
  Author's separate explanatory addendum:
  \url{https://www.jmilne.org/math/articles/2002a.html}
\end{itemize}

Milne, J. S. (2026, November 9). \emph{On the Hodge, Tate, and standard
conjectures for abelian varieties} (Hexagon:2611.00244v1) {[}Expository
preprint{]}. \url{https://www.jmilne.org/math/xnotes/HTS.pdf}

\begin{itemize}
\tightlist
\item
  Versioned deposit: \url{https://hexagonmath.org/2611.00244v1}
\end{itemize}

van Hoften, P., \& Sempliner, J. (2026, October 18). \emph{The Hodge
conjecture for CM abelian varieties} (Hexagon:2610.00259v2)
{[}Preprint{]}. \url{https://hexagonmath.org/pdf/2610.00259v2}

The source snapshots, extracted text, and selected rendered theorem
pages are preserved in \nolinkurl{research/cc_sources/}. Exact
web-reference identifiers and inspection methods are in the accompanying
citation and check-event ledgers.

% END consequences_criteria_Tate2026.tex

% BEGIN consequences_criteria_abelian_group.tex
\section{An explicit motivic group for the abelian sector after
extending
coefficients}\label{an-explicit-motivic-group-for-the-abelian-sector-after-extending-coefficients}

\textbf{Scope.} This is a complete deduction from the all-abelian Tate
and numerical--homological equivalence inputs recorded in CC-TATE01--02
and the fully proved realization theorem AQ-ABELIAN-RECOGNITION. It
concerns the abelian-generated category over a \textbf{fixed finite
field}, after extending coefficients. It neither classifies arbitrary
numerical motives nor identifies the uncertain historical notation in
cote 15 with a modern group by fiat. Literature cutoff: 20 November
2026; checks: 21 November 2026 UTC. No claim of historical novelty is
made.

\subsection{CC-AG01 --- The abelian category with algebraically closed
coefficients}\label{cc-ag01-the-abelian-category-with-algebraically-closed-coefficients}

Let \(q=p^a\), let \(\ell\ne p\), set \(E=\overline{\mathbf Q}_\ell\),
and fix an embedding \(\iota:\overline{\mathbf Q}\hookrightarrow E\).
Let \(\mathcal A_q\) be the rational numerical category generated by
abelian varieties over \(\mathbf F_q\), finite étale schemes, and Tate
motives under finite sums, tensor products, duals, and direct summands,
exactly as in AQ-ABELIAN-RECOGNITION. Use the parity-corrected symmetry,
so its ordinary direct-sum realization is symmetric monoidal. Define

\[
 \mathcal A_{q,E}
 =\operatorname{Kar}(\mathcal A_q\otimes_{\mathbf Q}E),
\]

where the tensor product first extends every Hom space and
\(\operatorname{Kar}\) then adds images of idempotents. Thus coefficient
extension adds summands; it does not add motives of arbitrary varieties
as generators.

Define the multiplicative abelian group

\[
 W(q)=\left\{\alpha\in\overline{\mathbf Q}^{\,\times}\ \middle|\
 \begin{array}{l}
 \alpha\text{ is a unit at every finite prime away from }p,\\
 \exists w\in\mathbf Z\ \forall\sigma:\ |\sigma(\alpha)|=q^{w/2}
 \end{array}\right\}.
\]

The integer \(w=w_q(\alpha)\) is unique. Let
\(\operatorname{Gr}^{\mathrm{fin}}_{W(q)}(E)\) be the category of
\(W(q)\)-graded \(E\)-vector spaces with finite total dimension. Its
tensor grading is multiplication in \(W(q)\), and its symmetry is the
ordinary interchange of vector factors.

\textbf{Theorem.} Geometric Frobenius gives a symmetric monoidal
equivalence

\[
 \mathcal A_{q,E}\ \xrightarrow{\ \sim\ }\
 \operatorname{Gr}^{\mathrm{fin}}_{W(q)}(E),\qquad
 M\longmapsto
 \bigoplus_{\alpha\in W(q)}
 R_E(M)_{F=\iota(\alpha)}.                         \tag{CC-AG.1}
\]

In particular, every \(\alpha\in W(q)\) labels an invertible rank-one
object \(L_\alpha\), unique up to isomorphism, and every object is a
finite direct sum of these, with arbitrary nonnegative multiplicities.
Tensor product satisfies
\(L_\alpha\otimes L_\beta\simeq L_{\alpha\beta}\). The weight is
\(w_q(\alpha)\), and \(\mathbf Q(1)\) corresponds to the label
\(q^{-1}\). Also

\[
 K_0(\mathcal A_{q,E})\simeq\mathbf Z[W(q)]
\]

as rings, with the classes of actual objects corresponding to the
finitely supported nonnegative sums of group elements.

\textbf{Complete proof.} AQ-ABELIAN-RECOGNITION proves, for objects
\(M,N\) of the rational category,

\[
 \operatorname{Hom}_{\mathcal A_q}(M,N)
          \otimes_{\mathbf Q}\mathbf Q_\ell
 \simeq
 \operatorname{Hom}_{\mathbf Q_\ell[F]}
       (R_\ell M,R_\ell N),                       \tag{CC-AG.2}
\]

and proves semisimplicity of Frobenius. Extend (CC-AG.2) to \(E\). The
commutant on the right is the kernel of the finite linear system
\(TF_M=F_NT\); field extension commutes with this kernel. The left side
becomes \(\operatorname{Hom}_{\mathcal A_q}(M,N)\otimes_{\mathbf Q}E\).
Thus realization remains fully faithful on the category with extended
Hom spaces. For idempotents \(e\in\operatorname{End}(M)\otimes E\) and
\(f\in\operatorname{End}(N)\otimes E\), its Hom space is

\[
 f\bigl(\operatorname{Hom}_{\mathcal A_q}(M,N)\otimes E\bigr)e.
\]

The identical formula holds for maps between their realized images.
Hence realization is fully faithful on the entire idempotent completion.
Frobenius remains semisimple after extension and on every summand.

Its eigenvalues lie in \(W(q)\). Indeed, eigenvalues \(\pi\) on \(H^1\)
of an abelian variety are integral Weil numbers of weight one. Choose an
identification of \(\overline{\mathbf Q}\) with the algebraic numbers in
\(\mathbf C\) when using complex conjugation. Then
\(\pi\overline{\pi}=q\), so integrality of both factors forces every
non-\(p\) finite valuation to vanish. They are therefore \(p\)-units.
The eigenvalues on exterior powers are products of such numbers; Tate
twists multiply them by integral powers of \(q\); finite étale factors
supply roots of unity. These operations, duals and summands preserve the
indicated group of numbers and their integral weights. This proves that
the displayed functor lands in the target category. Since a map between
semisimple Frobenius modules commutes with Frobenius exactly when it
preserves its eigenspaces, the functor is fully faithful into that
category.

It remains to prove that every label occurs and can be split off with
multiplicity one. Let \(\alpha\in W(q)\), of weight \(w\). Choose
\(n\ge0\) large enough that \(\beta=q^n\alpha\) is integral and
\(m=w+2n>0\). Such \(n\) exists: outside \(p\) all valuations are
already zero, and multiplication by \(q^n\) clears the finitely many
negative valuations above \(p\) in a number field containing \(\alpha\).

Choose \(\gamma\in\overline{\mathbf Q}\) with \(\gamma^m=\beta\). It is
integral by transitivity of integrality, and every conjugate has
absolute value \(q^{1/2}\). Honda's Main Theorem (1968, p.~83) therefore
supplies an abelian variety \(B/\mathbf F_q\) for which \(\gamma\) is a
Frobenius eigenvalue on \(H^1\). This is exactly the occurrence argument
already proved in AQ-WEIL-EXIST. The object

\[
 M_\alpha=h^1(B)^{\otimes m}(n)\in\mathcal A_q
\]

has eigenvalue \(\gamma^m q^{-n}=\alpha\), with weight \(m-2n=w\).

Its realized \(\alpha\)-eigenspace is nonzero and Frobenius is scalar
there. Choose a line in that eigenspace and a vector-space complement to
it, including all other eigenspaces. Projection onto the line commutes
with Frobenius. Full faithfulness over \(E\) lifts this rank-one
projector to an idempotent in
\(\operatorname{End}_{\mathcal A_q}(M_\alpha)\otimes E\). The idempotent
completion supplies its image \(L_\alpha\). This is why occurrence alone
would not have sufficed in the rational category: the projector is now
allowed to have extended coefficients.

Every finite-dimensional \(W(q)\)-graded vector space is a finite sum of
these one-dimensional realizations. This proves essential surjectivity.
The tensor-product eigenspace decomposition is compatible with
multiplication of labels, and the corrected symmetry realizes as
ordinary symmetry, so the equivalence is symmetric monoidal. Full
faithfulness and realization give
\(L_\alpha\otimes L_\beta\simeq L_{\alpha\beta}\), in particular
\(L_{\alpha^{-1}}\) is the inverse. The claims about weights and Tate
twists follow from the eigenvalue convention. Finally the target is
semisimple with these simple objects; its Grothendieck ring is freely
generated additively by their classes, with the multiplication just
computed. This proves the ring assertion and the whole theorem.
\(\square\)

\textbf{Comparison with the primary literature.} Milne's Definition 2.5,
author-manuscript pp.~14--15, uses the condition that some \(q^n\alpha\)
is integral. Under the common archimedean absolute-value condition, this
is equivalent to the \(p\)-unit condition above: complex conjugation
gives \(\alpha\overline\alpha=q^w\); away from \(p\), integrality of
both conjugate factors and this identity force both valuations to be
zero. The converse is the denominator-clearing argument in the proof.
Milne's Proposition 2.6, p.~15, contains the occurrence argument, while
his Proposition 3.7, p.~36, gives the broader coefficient-extended
classification under the full Tate conjecture imposed in that part of
the paper. Here (CC-AG.2) is already proved for \(\mathcal A_q\) from
the cited all-abelian theorem; no full Tate conjecture for arbitrary
varieties is used.

\subsection{CC-AG02 --- The split group, its weight factor, and its
Frobenius
point}\label{cc-ag02-the-split-group-its-weight-factor-and-its-frobenius-point}

Write \(\omega_E\) for the underlying-vector-space functor in (CC-AG.1).
For any abelian group \(M\), let

\[
 D_E(M)=\operatorname{Spec}E[M],
\]

with group algebra Hopf structure
\(\Delta([\alpha])=[\alpha]\otimes[\alpha]\).

\textbf{Theorem.} There is a canonical identification, relative to the
chosen realization and \(\iota\),

\[
 \operatorname{Aut}^{\otimes}(\omega_E)
       \simeq D_E(W(q)).                            \tag{CC-AG.3}
\]

Let

\[
 W_0(p)=\left\{\alpha\in\overline{\mathbf Q}^{\,\times}\ \middle|\
 \begin{array}{l}
 \alpha\text{ is a }p\text{-unit},\\
 |\sigma(\alpha)|=1\text{ for every }\sigma
 \end{array}\right\}.
\]

Choosing a square root \(u=\sqrt q\) gives

\[
 W(q)\simeq W_0(p)\times\mathbf Z,\qquad
 \alpha\longmapsto(\alpha u^{-w_q(\alpha)},w_q(\alpha)),
\]

and consequently

\[
 \operatorname{Aut}^{\otimes}(\omega_E)
       \simeq D_E(W_0(p))\times\mathbf G_{m,E}.       \tag{CC-AG.4}
\]

The weight cocharacter is the inclusion of the displayed \(\mathbf G_m\)
factor. The Tate object has character \((1,-2)\). Frobenius is the
distinguished point \(f\) with character values
\(\chi_\alpha(f)=\iota(\alpha)\), and its nonnegative powers are Zariski
dense in this affine group.

\textbf{Complete proof.} Let \(R\) be any commutative \(E\)-algebra. A
tensor automorphism of \(\omega_E\otimes_E R\) acts on each \(L_\alpha\)
by a scalar \(g_\alpha\in R^\times\). Naturality with respect to every
matrix endomorphism of \(L_\alpha^{\oplus r}\), and to the inclusions
and projections of summands, forces its action on every object to be
this scalar on its entire \(\alpha\)-graded summand. Compatibility with
tensor product and the unit says

\[
 g_{\alpha\beta}=g_\alpha g_\beta,\qquad g_1=1.
\]

Conversely any group homomorphism \(W(q)\to R^\times\) defines such a
natural tensor automorphism by those scalars. This correspondence is
functorial in \(R\), and the group algebra represents precisely this
functor. It proves (CC-AG.3).

The function \(w_q:W(q)\to\mathbf Z\) is a group homomorphism with
kernel \(W_0(p)\). The number \(u\) is a \(p\)-unit of weight one, so
\(n\mapsto u^n\) splits this homomorphism. Explicitly, the displayed map
and \((a,n)\mapsto au^n\) are inverse group homomorphisms. Taking group
algebras gives (CC-AG.4). A scalar \(z\) in the weight factor acts on
label \(\alpha\) by \(z^{w_q(\alpha)}\), which proves the cocharacter
statement. Since \(q^{-1}=u^{-2}\), the Tate character is \((1,-2)\).

The map \(\alpha\mapsto\iota(\alpha)\) is multiplicative, so it defines
\(f\in D_E(W(q))(E)\). To prove density, let a regular function

\[
 h=\sum_{j=1}^s c_j[\alpha_j]\in E[W(q)]
\]

with distinct \(\alpha_j\) vanish at every \(f^r\), \(r\ge0\). For
\(r=0,\ldots,s-1\), this gives

\[
 \sum_{j=1}^s c_j\,\iota(\alpha_j)^r=0.
\]

The Vandermonde determinant is nonzero because \(\iota\) is injective
and the labels are distinct. Thus every \(c_j\) is zero. No nonzero
regular function vanishes on these points; their Zariski closure is the
whole affine group. \(\square\)

\subsection{CC-AG03 --- Fixed-field torsion, coefficients, and the
limits of the
description}\label{cc-ag03-fixed-field-torsion-coefficients-and-the-limits-of-the-description}

The character group \(W(q)\) contains all roots of unity \(\mu_\infty\),
precisely its torsion subgroup. Thus \(D_E(W(q))\) is a diagonalizable
affine group of multiplicative type, generally of infinite type; it is
\textbf{not a protorus}. For example, the inclusion
\(\mu_n\subset W(q)\) gives a quotient
\(D_E(W(q))\twoheadrightarrow D_E(\mu_n)\), a nontrivial finite étale
group when \(n>1\). This also shows why the weight-zero factor
\(D_E(W_0(p))\) must not be confused with an identity component: it
retains the same root-of-unity characters. These are the finite-field
Artin contributions.

\textbf{Justification of the quotient assertion.} For a subgroup
\(H\subset M\), choose representatives for the cosets of \(H\) in \(M\).
The group algebra \(E[M]\) is then a free \(E[H]\)-module, with those
representatives as a basis. Its rank is nonzero, so the induced map
\(D_E(M)\to D_E(H)\) is faithfully flat. For \(H=\mu_n\), the abstract
character group is cyclic of order \(n\), whose group algebra is
\(E[t]/(t^n-1)\). In characteristic zero it is finite étale and has
\(n\) geometric points. A protorus, being an inverse limit of tori, is
connected and cannot have this nontrivial finite étale quotient.
\(\square\)

This description is compatible with AQ-RATIONAL-FIBRE. Over
\(\mathbf F_9\), the quaternionic elliptic-curve motive of
AQ-RATIONAL-FIBRE has eigenvalue \(-3\) with multiplicity two, and there
is no rational numerical object having only one copy. After coefficient
extension, (CC-AG.1) supplies the corresponding rank-one summand. Thus
(CC-AG.3) does not assert a neutral \(\mathbf Q\)-valued fibre functor
or remove rational multiplicity restrictions.

There is also a separate base-field operation: passing from a fixed
\(\mathbf F_q\) to \(\overline{\mathbf F}_p\) uses power transitions on
Frobenius characters and their germs. Milne treats this distinction in
\textbf{Proposition 3.2}, p.~34. We have held the finite base field
fixed throughout. The category of arbitrary numerical motives, integral
lattices, mixed extensions, and the archive's uncertain \(Q'\), \(M_0\),
or quotient conventions remain outside (CC-AG.1)--(CC-AG.4).

\subsection{Bibliography and exact
inputs}\label{bibliography-and-exact-inputs}

\begin{itemize}
\tightlist
\item
  Honda, T. (1968). Isogeny classes of abelian varieties over finite
  fields. \emph{Journal of the Mathematical Society of Japan, 20}(1--2),
  83--95. \url{https://doi.org/10.2969/jmsj/02010083}. \textbf{Main Theorem,
  p.~83}; the definition immediately preceding it identifies the
  relevant integral numbers by their conjugate absolute values. Primary
  PDF theorem text inspected through
  \url{https://people.math.rochester.edu/faculty/doug/otherpapers/Honda-Ab.pdf}.
\item
  Milne, J. S. (1994). Motives over finite fields. In U. Jannsen, S. L.
  Kleiman, \& J.-P. Serre (Eds.), \emph{Motives} (Proceedings of
  Symposia in Pure Mathematics, Vol. 55, Part 1, pp.~401--459). American
  Mathematical Society. \url{https://jmilne.org/math/articles/1994aP.pdf}.
  Author-manuscript locators: parity-corrected symmetry, pp.~5--6;
  scalar extension (1.3.2), p.~7; Definition 2.5 and Proposition 2.6,
  pp.~14--15; Proposition 3.2, p.~34; Proposition 3.7, p.~36. The global
  classification in the latter source retains its stated Tate
  hypothesis.
\item
  Internal proved inputs: \textbf{AQ-WEIL-EXIST} and
  \textbf{AQ-ABELIAN-RECOGNITION}, including (AQ.15), in
  research/archive\_questions.tex; the all-abelian Tate and
  numerical--homological input and its exact primary bibliography are
  recorded in \textbf{CC-TATE01--02},
  research/consequences\_criteria\_Tate2026.tex. The major 2026 all-CM
  theorem remains an inspected cited input, not a theorem whose full
  foundational proof has been independently reconstructed in this run.
\end{itemize}

% END consequences_criteria_abelian_group.tex

% BEGIN consequences_criteria_abelian_germs.tex
\section{The abelian sector over the algebraic closure of a finite field}

This section completes the base-field passage left separate in
CC-AG01--03. It is a scoped application of the classical Frobenius-germ
description, using the all-abelian Tate inputs already recorded in
CC-TATE01--02 and AQ-ABELIAN-RECOGNITION. It does not assume the Tate
conjecture for arbitrary smooth projective varieties and makes no
claim of novelty.

Fix \(\bar k=\overline{\mathbb F}_p\), \(\ell\ne p\),
\(E=\overline{\mathbb Q}_\ell\), and
\(\iota:\overline{\mathbb Q}\hookrightarrow E\). Let
\(\mathcal A_{\bar k}\) be the rational numerical category generated by
abelian varieties over \(\bar k\), finite \'etale schemes, and Tate
motives, under finite sums, tensor products, duals, and direct summands.
Set
\[
 \mathcal A_{\bar k,E}
 =\operatorname{Kar}(\mathcal A_{\bar k}\otimes_{\mathbb Q}E).
\]
Use the parity-corrected symmetry, as in CC-AG01. Over \(\bar k\),
finite \'etale motives are sums of the unit.

The primary precedent is Milne (1994), in the corrected author version
dated 2 November 2019: Theorem 1.18(a), p.~11, gives the realization
into germs; Remark 2.26, p.~22, identifies the associated Weil-number
protorus; and Remark 2.29, p.~24, describes the torsion-free germ group.
Those statements occur with the full Tate hypothesis displayed on
p.~11. The proofs below restrict the argument to the established
abelian sector and explicitly check descent and the morphisms.

\subsection{CC-AG04: descent of the category and its morphisms}

Write \(k_a=\mathbb F_{p^a}\), with the positive integers ordered by
divisibility. For \(a\mid b\), the transition functor is extension of
the base field \(k_a\subset k_b\). There are tensor equivalences
\[
 \underset{a}{\operatorname{2-colim}}\mathcal A_{p^a}
 \ \simeq\ \mathcal A_{\bar k},
 \qquad
 \underset{a}{\operatorname{2-colim}}\mathcal A_{p^a,E}
 \ \simeq\ \mathcal A_{\bar k,E}.                 \tag{CC-AG.5}
\]
Here the filtered categorical colimit can be described explicitly:
objects are objects at finite stages, and the morphisms between two
such objects are the filtered colimit of their morphism spaces after
passing to common finite extensions. In particular, for models
\(M_a,N_a\) over \(k_a\),
\[
 \operatorname{Hom}_{\mathcal A_{\bar k}}(M_{\bar k},N_{\bar k})
 =
 \underset{a\mid b}{\operatorname{colim}}
 \operatorname{Hom}_{\mathcal A_{p^b}}(M_b,N_b).   \tag{CC-AG.6}
\]
The same formula holds after extending Hom spaces to \(E\) and for
their idempotent corners. The underlying \(\ell\)-adic realization
therefore exists and is faithful on this abelian category over
\(\bar k\).

\paragraph{Complete proof.}
First check numerical equivalence under extension of the base field.
Let \(X/k\) be smooth projective, with \(k\) finite, and let \(K/k\)
be finite of degree \(d\). Its projection \(f:X_K\to X\) is finite
\'etale. For a rational cycle \(z\) on \(X\) and any complementary
cycle \(v\) on \(X_K\), the projection formula and the definition of
degree give
\[
 d\,\deg_K(f^*z\cdot v)
 =\deg_k f_*(f^*z\cdot v)
 =\deg_k(z\cdot f_*v).                           \tag{CC-AG.7}
\]
Consequently \(z\) numerically trivial implies \(f^*z\) numerically
trivial. Conversely, for a complementary cycle \(u\) on \(X\),
\[
 \deg_K(f^*z\cdot f^*u)=\deg_k(z\cdot u).
\]
The equality follows from compatibility of intersection with flat
pullback and preservation of degree under extension of the ground
field; on a closed point this is the equality of dimensions after
tensoring its residue field with \(K\). Thus numerical triviality is
also reflected by this pullback. These cycle facts are precisely the
projection and pullback formulas of Stacks, Lemmas 42.62.3--42.62.4
(Tags 0FFB--0FFC), and the degree formula 42.41.2 (Tag 0AZ2).

The same equivalence holds for \(\bar k/k\): every complementary
cycle over \(\bar k\) descends to a finite extension, so the finite
calculation applies. Cycle descent is Stacks, Lemma 45.8.2
(Tag 0FH6), or its detailed limit statement 42.67.10 (Tag 0FVQ).
We have proved, in particular, that pullback on rational numerical
cycle spaces is injective.

Every abelian variety over \(\bar k\), together with its group law,
identity, inverse, and a projective embedding, is defined over a
finite subfield. Indeed these are finitely presented schemes,
morphisms, and finitely many equations between their compositions;
Stacks, Lemma 32.10.1 (Tag 01ZM), descends them simultaneously.
Smoothness is detected after the faithfully flat extension to
\(\bar k\), projectivity is supplied by the descended embedding,
and geometric connectedness is detected on the algebraic closure.
Thus the descended object is an abelian variety. Finite tensor
constructions and Tate twists descend as well.

A rational numerical correspondence is represented by a finite
rational combination of cycles. All those cycles descend to one
finite field. Equalities between the resulting numerical
correspondences are reflected at that stage by the injectivity
proved above. In particular, an idempotent over \(\bar k\) descends
as an actual numerical idempotent, not merely as a representative
cycle that might fail its idempotent equation. Taking its image gives
the descended direct summand. This proves descent of every object,
and the same cycle argument proves (CC-AG.6), including injectivity
and compatibility with composition and tensor product. Hence it
proves the first equivalence of (CC-AG.5).

Tensoring a filtered colimit of rational vector spaces with \(E\)
commutes with that colimit. An \(E\)-coefficient correspondence has
only finitely many rational cycle classes in any expression for it,
so its field of definition can again be enlarged to one finite
stage. The previously injective maps remain injective after
tensoring with \(E\). An idempotent equation over \(E\) is therefore
reflected there too. For idempotents \(e,f\), its Hom space is the
corner \(f\operatorname{Hom}e\); the same argument applies to these
corners. This proves the second equivalence.

Finally, numerical and \(\ell\)-homological equivalence agree on
the finite-field abelian products by the cited all-abelian input.
Their agreement over \(\bar k\) follows by descending a cycle and
using numerical injectivity and the unchanged geometric
cohomology. The resulting realization is faithful at each finite
stage, hence on the displayed colimit and its \(E\)-coefficient
idempotent completion. \(\square\)

\subsection{CC-AG05: the actual category of Weil germs}

Use the groups \(W(p^a)\) from CC-AG01 and set
\[
 \mathcal W_p=\underset{a\mid b}{\operatorname{colim}}
 \left(W(p^a),\ \alpha\longmapsto\alpha^{b/a}\right).
                                                        \tag{CC-AG.8}
\]
Write \([\alpha,a]\) for a germ. Explicitly,
\([\alpha,a]=[\beta,b]\) if and only if
\(\alpha^{bN}=\beta^{aN}\) for some positive integer \(N\).
There is a symmetric monoidal equivalence compatible with the
underlying cohomological fibre functor:
\[
 \mathcal A_{\bar k,E}\ \simeq\
 \operatorname{Gr}^{\mathrm{fin}}_{\mathcal W_p}(E).
                                                        \tag{CC-AG.9}
\]
The target consists of finite-dimensional vector spaces graded by
\(\mathcal W_p\). In particular, the germ transition is realized by
the actual numerical category, including all its morphisms.

Choose \(s=\sqrt p\) in \(\overline{\mathbb Q}\). With
\(W_0(p)\) as in CC-AG02, there are isomorphisms
\[
 \mathcal W_p
 \simeq \mathbb Z\oplus\bigl(W_0(p)\otimes_{\mathbb Z}\mathbb Q\bigr)
 \simeq \mathbb Z\oplus\bigl(W_0(p)/\mu_\infty\bigr).
                                                        \tag{CC-AG.10}
\]
The weight is the \(\mathbb Z\)-coordinate, and the Tate object
has coordinates \((-2,0)\), writing the second factor additively.
Consequently, for the underlying-vector-space fibre functor
\(\omega_E\),
\[
 \operatorname{Aut}^{\otimes}(\omega_E)
 \simeq D_E(\mathcal W_p)
 \simeq \mathbb G_{m,E}\times
 D_E\bigl(W_0(p)\otimes_{\mathbb Z}\mathbb Q\bigr). \tag{CC-AG.11}
\]
This group is a protorus.

\paragraph{Complete proof.}
An object defined over \(k_a\) has the semisimple Frobenius
eigenspace decomposition of CC-AG01. Give its \(\alpha\)-eigenspace
the label \([\alpha,a]\), combining eigenspaces with the same germ.
Under \(k_a\subset k_b\), geometric Frobenius becomes its
\(b/a\)-th power. The unchanged geometric cohomology therefore
has exactly the relabeling prescribed in (CC-AG.8).
CC-AG04 makes this construction independent of a chosen model.

We verify full faithfulness, including the maps newly available
over \(\bar k\). Take models of two objects over one common \(k_a\).
A linear map between their realized vector spaces respects germ
labels precisely when each of its nonzero blocks between
Frobenius eigenspaces has eigenvalues \(\alpha,\beta\) with
\((\alpha/\beta)^r=1\) for some positive \(r\). Only finitely many
eigenvalue pairs occur. Choose a common multiple of their
root-of-unity orders. The linear map then commutes with that
power of Frobenius. By the full faithfulness in CC-AG01 over the
corresponding finite extension, it is induced by a unique
correspondence with \(E\)-coefficients at that stage. Conversely,
every correspondence over \(\bar k\) descends to a finite stage
by CC-AG04, so commutes with the corresponding Frobenius power
and preserves the germ labels. Faithfulness was proved in
CC-AG04. This proves the required isomorphism on every Hom space.

Every finite list of germs has representatives at a common finite
stage. CC-AG01 realizes each representative as a rank-one object
there, after extending coefficients and taking the idempotent
completion. Direct sums realize every finite-dimensional graded
vector space. This proves essential surjectivity. Tensor products
multiply eigenvalues and therefore multiply germs. The corrected
symmetry realizes as ordinary interchange, proving (CC-AG.9).

For (CC-AG.10), write an element at stage \(a\) uniquely as
\(\alpha=s^{aw}u\) with \(u\in W_0(p)\). In these coordinates the
transition \(a\to b\) is
\[
 (w,u)\longmapsto(w,u^{b/a}).
\]
Its colimit is \(\mathbb Z\oplus(W_0(p)\otimes\mathbb Q)\):
the second coordinate of the stage-\(a\) element is
\(u\otimes a^{-1}\), in the additive tensor notation for this
abelian group. These maps are compatible and give the standard
localization of an abelian group by all positive integers.

The group \(W_0(p)\) is divisible: any algebraic \(n\)-th root
of a \(p\)-unit all of whose conjugates have absolute value one
has those same two properties. Its torsion is exactly the group
of all roots of unity. Thus \(W_0(p)/\mu_\infty\) is divisible
and torsion-free, and the quotient map identifies it with
\(W_0(p)\otimes\mathbb Q\). This also checks (CC-AG.10) directly.
The Tate eigenvalue at stage \(a\) is \(p^{-a}\), giving the
specified coordinates.

For a finite grading category, the tensor automorphism calculation
in CC-AG02 applies verbatim: naturality forces a scalar on each
label, and tensor compatibility makes the scalars a homomorphism
\(\mathcal W_p\to R^\times\) for every \(E\)-algebra \(R\).
The representing Hopf algebra is \(E[\mathcal W_p]\), proving
(CC-AG.11). The group \(\mathcal W_p\) is torsion-free. Each of
its finitely generated subgroups is free of finite rank, so its
diagonalizable group is the inverse limit of the corresponding
split tori. It is therefore a protorus, as claimed. \(\square\)

\subsection{CC-AG06: base extension is the identity-component inclusion}

Let \(P_{a,E}=D_E(W(p^a))\) be the fixed-field group of CC-AG02,
and let \(P_{\bar k,E}=D_E(\mathcal W_p)\). Extension from \(k_a\)
to \(\bar k\) induces a closed immersion
\[
 P_{\bar k,E}\lhook\joinrel\longrightarrow P_{a,E},
\]
whose image is the identity component. There is an exact sequence
of affine group schemes, with faithfully flat last map,
\[
 1\longrightarrow P_{\bar k,E}\longrightarrow P_{a,E}
 \longrightarrow D_E(\mu_\infty)\longrightarrow1. \tag{CC-AG.12}
\]
With the fixed embedding \(\iota\), the last group is the constant
profinite group scheme \(\widehat{\mathbb Z}_E\). The extension
functor \(\mathcal A_{p^a,E}\to\mathcal A_{\bar k,E}\) is
essentially surjective; its additional morphisms merge exactly the
finite-field labels that differ by a root of unity.

\paragraph{Complete proof.}
The character map \(j_a:W(p^a)\to\mathcal W_p\) is
\(\alpha\mapsto[\alpha,a]\). Its kernel is \(\mu_\infty\):
the germ is the identity if and only if some positive power of
\(\alpha\) is one. It is surjective. Indeed, for any
\([\beta,b]\), choose \(\alpha\) with
\(\alpha^b=\beta^a\). This \(\alpha\) is a \(p\)-unit, and if
\(\beta\) has weight \(w\) relative to \(p^b\), every conjugate
of \(\alpha\) has absolute value \(p^{aw/2}\). Hence
\(\alpha\in W(p^a)\), and the displayed equality identifies
the two germs. Thus
\[
 W(p^a)/\mu_\infty\xrightarrow{\sim}\mathcal W_p. \tag{CC-AG.13}
\]
Under (CC-AG.9), extension of fields pushes a grading forward
along \(j_a\), combining the vector spaces with equal images.
Surjectivity and the rank-one occurrence theorem at that fixed
stage prove essential surjectivity and the assertion about maps.

The contravariant group-algebra construction turns (CC-AG.13)
into the displayed closed immersion. The inclusion
\(\mu_\infty\subset W(p^a)\) gives the last map of (CC-AG.12).
It is faithfully flat because a group algebra is a nonzero free
module over the group algebra of a subgroup, with any set of
coset representatives as basis. Its kernel has coordinate algebra
\[
 E[W(p^a)]\otimes_{E[\mu_\infty]}E
 =E[W(p^a)/\mu_\infty],
\]
which proves exactness, equivalently exactness for the fpqc
topology.

This kernel is connected. Its character group is torsion-free,
so its group algebra is a filtered union of Laurent polynomial
domains, with injective transition maps, and is itself a domain.
Conversely, every connected subset containing the identity in
\(P_{a,E}\) maps to the identity in each finite \'etale quotient
\(D_E(\mu_n)\). It is therefore contained in the kernel.
The kernel is precisely the connected component of the identity.

Finally,
\(D_E(\mu_\infty)=\varprojlim_n D_E(\mu_n)\).
The embedding \(\iota\) identifies
\(D_E(\mu_n)\) with the constant group \(\mathbb Z/n\mathbb Z\)
by \(r:\zeta\mapsto\iota(\zeta)^r\); these identifications are
compatible with divisibility and the reduction transition maps.
This proves the component-group assertion. \(\square\)

For completeness, the distinguished fixed-field Frobenius point
maps to \(1\in\widehat{\mathbb Z}\) in (CC-AG.12), since its
character on a root of unity \(\zeta\) is \(\iota(\zeta)\).
No positive integer power of this universal point lies in the
identity component: its component is that positive integer in
\(\widehat{\mathbb Z}\), which is nonzero. On each specified
finite collection of objects a sufficiently divisible power
does kill the finitely many root-of-unity ratios involved.
Thus the geometric category is correctly described by germs;
one does not obtain a single universal finite-field Frobenius
point inside its protorus.

\paragraph{Exact scope for original cote 15.}
The field/germ transition and the identity-component statement
are now established for actual abelian-generated rational numerical
motives after the displayed coefficient extension. The historical
norm-one group is \(W_0(p)\), as independently recovered in
AQ-WEIL-NOTATION, so (CC-AG.11) gives the precise coefficient-extended
abelian-sector version of its proposed group.
The assertion that \emph{every} numerical motive lies in this sector
is not proved here. The full Tate hypothesis in Milne's general
theory, or the separately established cross-cycle criterion for a
specified variety, remains distinct. Nor does this argument supply
a fibre functor over \(\mathbb Q\), lift numerical projectors to
Chow projectors, or remove the rational quaternion obstruction.

\paragraph{References.}
Milne, J.~S. (1994). Motives over finite fields. In U.~Jannsen,
S.~Kleiman, \& J.-P.~Serre (Eds.), \emph{Motives}
(Proceedings of Symposia in Pure Mathematics, Vol.~55, Part~1,
pp.~401--459). American Mathematical Society.
Corrected author version, 2 November 2019:
\url{https://www.jmilne.org/math/xnotes/Seattle1.pdf}.
Theorem 1.18(a), p.~11; Remark 2.26, p.~22; Remark 2.29, p.~24.

The Stacks Project Authors. (2026). \emph{The Stacks Project}.
Lemma 32.10.1, \url{https://stacks.math.columbia.edu/tag/01ZM};
Lemma 45.8.2, \url{https://stacks.math.columbia.edu/tag/0FH6};
Lemma 42.67.10, \url{https://stacks.math.columbia.edu/tag/0FVQ};
Lemmas 42.62.3--42.62.4,
\url{https://stacks.math.columbia.edu/tag/0FFB},
\url{https://stacks.math.columbia.edu/tag/0FFC};
Lemma 42.41.2, \url{https://stacks.math.columbia.edu/tag/0AZ2}.
The all-abelian Tate dependency and its precise primary versions
remain the separate inspected chain in CC-TATE01--02;
the fixed-field full faithfulness used here is AQ-ABELIAN-RECOGNITION
and CC-AG01.

% END consequences_criteria_abelian_germs.tex

% BEGIN consequences_criteria_rational_descent.tex
\section{Finite transitions and exact rational descent in the abelian sector}

This section sharpens the coefficient and base-field statements in
cote~15.  It uses the all-abelian Tate inputs already recorded in
CC-TATE01--02 and AQ-ABELIAN-RECOGNITION, together with the established
equivalences CC-AG01 and CC-AG04--06.  The arbitrary-variety Tate
conjecture is not an assumption of these scoped conclusions.  The
rational classification applied below is classical: Milne (1994),
corrected author version dated 2 November 2019, Propositions 2.4,
2.6, 2.8, 2.17 and Theorems 2.15, 2.16, 2.18, 2.19.
We prove the precise
consequences for the categories already established here; we do not
reprove the cited classification of central division algebras.

Write \(\mathcal A_q\) for the rational numerical category generated
by abelian varieties, finite \'{e}tale schemes, and Tate motives over
\(\mathbb F_q\), and \(\mathcal A_{\bar k}\) for its counterpart over
\(\bar k=\overline{\mathbb F}_p\).  As before, use the
parity-corrected symmetry, \(E=\overline{\mathbb Q}_\ell\) with
\(\ell\ne p\), and a fixed embedding
\(\iota:\overline{\mathbb Q}\hookrightarrow E\).
Coefficient extension includes idempotent completion.

\subsection{CC-AG07: the exact finite-transition sequence}

For every prime power \(q=p^a\) and positive integer \(m\), there is
an exact sequence of abstract abelian groups with their natural
Galois actions
\[
  1\longrightarrow\mu_m\longrightarrow W(q)
   \xrightarrow{\alpha\mapsto\alpha^m}W(q^m)
   \longrightarrow1.                                      \tag{CC-AG.21}
\]
Here \(\mu_m\) consists of the algebraic \(m\)-th roots of unity;
\(W(q)\) consists of the algebraic \(p\)-units of integral
\(q\)-weight, as in CC-AG01.  In particular, the power map has
\emph{zero cokernel}.

Under the established grading equivalences, extension of the
base field \(\mathbb F_q\subset\mathbb F_{q^m}\) sends a label
\(\alpha\) to \(\alpha^m\).  This functor on the
\(E\)-coefficient categories is essentially surjective.  Its
associated diagonalizable groups have the fpqc exact sequence
\[
  1\longrightarrow P_{q^m,E}\longrightarrow P_{q,E}
   \longrightarrow D_E(\mu_m)\longrightarrow1,
 \qquad P_{q,E}=D_E(W(q)).                                \tag{CC-AG.22}
\]
The last group is finite \'{e}tale of order \(m\).  After the
identifications of component groups in CC-AG06, the first arrow
induces multiplication by \(m\) on
\(\widehat{\mathbb Z}_E\).

\paragraph{Complete proof.}
The power map preserves the defining conditions: if every complex
conjugate of \(\alpha\) has absolute value \(q^{w/2}\), then those
of \(\alpha^m\) have absolute value \((q^m)^{w/2}\); valuations
away from \(p\) remain zero.  Conversely, given
\(\beta\in W(q^m)\) of weight \(w\), choose any algebraic root
\(\alpha\) with \(\alpha^m=\beta\).  Every complex embedding of
\(\mathbb Q(\alpha)\) sends this equation to the corresponding
equation for conjugates.  Therefore
\[
        |\sigma\alpha|^m=|\sigma\beta|=q^{mw/2},
        \qquad |\sigma\alpha|=q^{w/2}.
\]
At any finite place away from \(p\), taking valuations in
\(\alpha^m=\beta\) gives valuation zero for \(\alpha\).
Thus \(\alpha\in W(q)\).  The kernel is precisely \(\mu_m\);
all its elements have weight zero and are \(p\)-units.  The map
commutes with every element of
\(\Gamma=\operatorname{Gal}(\overline{\mathbb Q}/\mathbb Q)\).
This proves (CC-AG.21).

A finite \(W(q^m)\)-graded \(E\)-vector space can now be lifted by
choosing a root label for each of its finitely many nonzero
graded pieces, retaining the vector space and its multiplicity.
The grading equivalences identify this construction with an
object in \(\mathcal A_{q,E}\), whose base extension is the given
object.  No assertion of essential surjectivity before extending
coefficients is used.

For completeness, diagonalizable-group exactness here can be
checked on group algebras.  The surjection \(W(q)\to W(q^m)\)
gives a quotient
\[
       E[W(q)]\longrightarrow E[W(q^m)].
\]
Its kernel is the ideal generated by \(e_\zeta-1\) for
\(\zeta\in\mu_m\), because its quotient identifies precisely
the cosets of \(\mu_m\).  Consequently its spectrum is the
kernel of \(P_{q,E}\to D_E(\mu_m)\).  Moreover
\(E[W(q)]\) is free, and faithfully flat, over \(E[\mu_m]\):
choose one representative for each coset of \(\mu_m\) in
\(W(q)\).  This proves the asserted fpqc exactness.
The finite group algebra \(E[\mu_m]\) is a product of \(m\)
copies of \(E\), since \(E\) is algebraically closed of
characteristic zero; thus its diagonalizable group is finite
\'{e}tale of order \(m\).

The component quotient has character group
\(\mu_\infty\subset W(q)\).  The map on these torsion
character groups is again \(\zeta\mapsto\zeta^m\).
Taking its diagonalizable dual is multiplication by \(m\)
on \(D_E(\mu_\infty)\), which is the identified constant
profinite group \(\widehat{\mathbb Z}_E\).
This also treats \(m=1\). \(\square\)

\subsection{The exact classical rational classification used}

Let \(B\) denote either the fixed field \(\mathbb F_q\) or
\(\bar k\), and put
\[
 \Lambda_B=
 \begin{cases}
      W(q),&B=\mathbb F_q,\\
      \mathcal W_p,&B=\bar k.
 \end{cases}
\]
The group \(\mathcal W_p\) is the Weil-germ group of CC-AG05.
Every \(\Gamma\)-orbit \(O\) in \(\Lambda_B\) is finite.

For a label \(\lambda\), define its centre field
\[
        K_\lambda=
        \overline{\mathbb Q}^{\operatorname{Stab}_\Gamma(\lambda)}.
                                                               \tag{CC-AG.23}
\]
For a fixed-field label \(\alpha\), this is
\(\mathbb Q(\alpha)\).  For a germ
\(\xi=[\alpha,a]\), choose a sufficiently divisible positive
integer \(N\); then
\[
       K_\xi=\mathbb Q(\alpha^N), \qquad
       \xi=[\alpha^N,aN].                                    \tag{CC-AG.24}
\]
To see this stabilization directly, list the finitely many
conjugates of \(\alpha\).  An automorphism fixes its germ
exactly when its image of \(\alpha\), divided by \(\alpha\),
is a root of unity.  Choose \(N\) divisible by the orders of
all such ratios.  An automorphism then fixes \(\alpha^N\)
exactly when it fixes the germ, proving (CC-AG.24).
All larger multiples of this \(N\) give the same field.

Milne's cited classification assigns to each orbit \(O\) a
unique simple rational motive \(S_O\).  Its endomorphism ring
\[
           D_O=\operatorname{End}(S_O)
\]
is a central division algebra over \(K_\lambda\), for any
\(\lambda\in O\).  Here and below End is the endomorphism
algebra of the \emph{cohomological motive}, and the local
Brauer invariant uses Milne's geometric-Frobenius
normalization on printed p.~17.  These conventions fix
the following sign.

Choose \(\pi\in K_\lambda\), \(q_0=p^b\), and weight
\(w\in\mathbb Z\) as follows.  For a fixed-field label take
\((\pi,q_0)=(\alpha,q)\).  For a germ take the stabilized
representative in (CC-AG.24), so
\((\pi,q_0)=(\alpha^N,p^{aN})\).
The local invariants, viewed in \(\mathbb Q/\mathbb Z\), are
\[
 \operatorname{inv}_v(D_O)=
 \begin{cases}
   w/2,&v\text{ real},\\[2pt]
   -\displaystyle\frac{\operatorname{ord}_v(\pi)}
                         {\operatorname{ord}_v(q_0)}
        [K_{\lambda,v}:\mathbb Q_p],&v\mid p,\\[7pt]
   0,&\text{all other places}.
 \end{cases}                                                \tag{CC-AG.25}
\]
The displayed negative sign is the one in Milne's corrected
2019 author text, Theorem~2.16, printed p.~18, carried to
germs by Theorem~2.19, p.~20.  Replacing the cohomological
endomorphism algebra by its opposite reverses all its local
invariants.  Using the opposite normalization for the local
Brauer invariant also reverses their signs.  Neither
operation changes the index used below.

Set
\[
        d_O=\operatorname{lcm}_{v}
            \bigl(\text{denominator of }
                        \operatorname{inv}_v(D_O)\bigr),
                                                               \tag{CC-AG.26}
\]
where the denominator of zero is one.  Then
\[
        [D_O:K_\lambda]=d_O^2,\qquad
        \operatorname{rank}S_O=d_O\,|O|.                       \tag{CC-AG.27}
\]
Only finitely many local invariants are nonzero.
The formulas are independent of taking a further power of a
stabilized germ representative, since numerator and denominator
in its \(p\)-adic slope are both multiplied by the same integer.
Galois conjugation just permutes the places and the labels.

The exact primary locators are as follows: Proposition~2.4,
p.~13, identifies the centre; Proposition~2.6, pp.~13--14,
classifies the simple objects; Proposition~2.8, p.~14, gives
their ranks; Theorem~2.15, p.~18, gives the local-invariant
classification and index; Theorem~2.16, p.~18, computes
the invariants; Proposition~2.17, p.~19, states the integral
Grothendieck-ring image; and Theorems~2.18--2.19 with
Corollary~2.20, p.~20, give the germ versions.

In Milne's paper these results are written under the full Tate
hypothesis.  Their application here is restricted to the
abelian, Artin, and Tate category: the full realization on its
products and summands is the already established all-abelian
input, and occurrence of every Weil label and germ is
CC-AG01 and CC-AG05.  No occurrence or fullness assertion for
an arbitrary smooth projective variety is being inferred.
The classical division-algebra formula (CC-AG.25), rather
than an independent reconstruction of its crystalline
argument, is the cited input in the next deduction.

\subsection{CC-AG08: rational descent as an explicit divisibility test}

Under the established equivalence
\[
       \mathcal A_{B,E}\simeq
                    \mathrm{Gr}^{\mathrm{fin}}_{\Lambda_B}(E),
                                                               \tag{CC-AG.28}
\]
let \(V\) have multiplicity
\(n_\lambda=\dim_E V_\lambda\).
There exists an actual rational object \(M\in\mathcal A_B\)
and an isomorphism \(M_E\simeq V\) if and only if
\[
 \begin{aligned}
 n_{\sigma\lambda}&=n_\lambda
                 &&(\sigma\in\Gamma,\ \lambda\in\Lambda_B),\\
 d_O&\mid n_\lambda
                 &&(\lambda\in O).
 \end{aligned}                                                \tag{CC-AG.29}
\]
When these conditions hold, the rational isomorphism class is
unique, and explicitly
\[
       M\simeq\bigoplus_O S_O^{\oplus n_O/d_O},\qquad
       \operatorname{End}(M)
          \simeq\prod_{O:n_O>0}M_{n_O/d_O}(D_O).               \tag{CC-AG.30}
\]
Here \(n_O\) is the common multiplicity on \(O\).
The zero object corresponds to empty support.

Equivalently, coefficient extension embeds the Grothendieck
ring as
\[
 \begin{split}
 K_0(\mathcal A_B)&\hookrightarrow
                         \mathbb Z[\Lambda_B]^\Gamma,\\
 \operatorname{im}&=
     \left\{\sum_O n_O\!\sum_{\lambda\in O}e_\lambda:
                   d_O\mid n_O,\ \text{finite support}\right\}.
 \end{split}                                                \tag{CC-AG.31}
\]
Its actual-object cone is obtained by requiring \(n_O\ge0\).
The multiplication on the target is convolution in the
group ring, so this also specifies the rational tensor
multiplicities once the \(d_O\) are known.

\paragraph{Complete proof of the scoped deduction.}
The category \(\mathcal A_B\) is semisimple.  The cited
classification identifies its simple isomorphism classes
with the orbits \(O\), and its scalar extension has the
one-dimensional simple labels of (CC-AG.28).
For a simple \(S_O\), scalar extension and the finite
dimensionality of its endomorphism algebra give
\[
 \begin{split}
 \operatorname{End}_{\mathcal A_{B,E}}((S_O)_E)
       &=D_O\otimes_{\mathbb Q}E\\
       &\simeq\prod_{\tau:K_\lambda\hookrightarrow E}M_{d_O}(E).
 \end{split}                                                \tag{CC-AG.32}
\]
The equality of the endomorphism algebra after coefficient
extension is literal in the scalar-extended category, and
idempotent completion does not alter the endomorphisms of
an existing object.  The second equality follows because
the separable centre splits into one \(E\)-factor per
embedding and a central simple algebra splits over an
algebraically closed field.

The centre embeddings correspond to the labels in \(O\);
for germs this uses the stabilized centre in (CC-AG.24).
The associated central idempotents are the projectors onto
those label components.  In the graded-vector-space
category, a component of dimension \(r\) has endomorphism
ring \(M_r(E)\).  Comparing the corresponding factors in
(CC-AG.32) therefore proves
\[
               (S_O)_E
                  \simeq\bigoplus_{\lambda\in O}
                                E_\lambda^{\oplus d_O}.      \tag{CC-AG.33}
\]
This argument identifies each multiplicity, not just the
sum of the dimensions.

Every rational object is a finite direct sum
\(\bigoplus_O S_O^{\oplus a_O}\), with uniquely determined
nonnegative integers \(a_O\).  Equation (CC-AG.33) gives
\(n_\lambda=a_Od_O\) on \(O\), proving the necessity of
(CC-AG.29).  Conversely, if (CC-AG.29) holds, the integers
\(a_O=n_O/d_O\) are nonnegative and have finite support.
The direct sum in (CC-AG.30) exists in \(\mathcal A_B\);
its coefficient extension has exactly the prescribed
graded dimensions.  Finite graded vector spaces with the
same dimensions are isomorphic, so it extends to \(V\).
Uniqueness of the simple multiplicities proves uniqueness
of the rational isomorphism class.  Schur's lemma gives
zero morphisms between different simple classes and
endomorphism algebra \(M_{a_O}(D_O)\) on an isotypic sum,
proving the second formula in (CC-AG.30).

For the Grothendieck ring, the classes of the simple
objects are a free integral basis by semisimplicity.
Equation (CC-AG.33) sends that basis to the disjointly
supported nonzero elements
\(d_O\sum_{\lambda\in O}e_\lambda\).  They are linearly
independent, and their integral span is exactly the
displayed image.  Coefficient extension is a tensor
functor and the grading tensor product multiplies labels,
so the map respects convolution.  Nonnegative simple
multiplicities give precisely the stated actual-object
cone.  This proves all assertions. \(\square\)

\paragraph{Why coefficient extension matters for CC-AG07.}
For \(m=2\), let \(V\) be the one-dimensional sign Artin
motive over \(\mathbb F_{q^2}\); it is the minus summand of
\(h(\operatorname{Spec}\mathbb F_{q^4})\), so its sole
Frobenius label is \(-1\).  It is not the base extension of
any rational object of \(\mathcal A_q\).  Indeed such an
object would have rank one, hence a single label
\(\alpha\in W(q)\); (CC-AG.29) forces that label to be
fixed by \(\Gamma\), so \(\alpha\in\mathbb Q\).
Base extension would require \(\alpha^2=-1\), a
contradiction.  Over \(E\), either label \(i\) or \(-i\)
supplies a preimage, exactly as CC-AG07 asserts.

\subsection{What this resolves, and what it does not identify}

For the established abelian sector, the rational coefficient
condition is now an exact orbitwise arithmetic obstruction,
rather than an unspecified descent requirement.  The
supersingular quaternion example AQ-RATIONAL-FIBRE is its
smallest illustrative case: over \(\overline{\mathbb F}_p\)
the supersingular weight-one germ has centre \(\mathbb Q\),
invariants \(1/2\) at \(p\) and the real place, and index
two.  Its single split label must therefore occur with
even multiplicity in a rational object.

The orbitwise test concerns descent of an \emph{object}.
It does not choose coherent rational vector spaces for all
objects and tensor products, and it does not produce a
\(\mathbb Q\)-valued fibre functor.  In fact,
AQ-RATIONAL-FIBRE already proves that the full rational
abelian category has no such fibre functor, by its
supersingular quaternion endomorphism algebra.  This
nonexistence is a conclusion, not a remaining question.
The nonneutral Tannakian description remains appropriate: Milne,
Corollary~2.25 and Remark~2.26, p.~22, identify its
multiplicative-type band from the natural Galois action on
Weil labels or germs.  A band is not a choice of neutral
fibre functor.  In particular, (CC-AG.29) is not a claim
that the rational category is the ordinary representation
category of that band over \(\mathbb Q\).

Likewise the test does not lift numerical projectors to
Chow projectors.  Applying any of these classifications
to all numerical motives would still require the
arbitrary-motive Tate and realization inputs isolated in
CC-XC01--04 and BR-EFF01.  No such implication follows
from the adopted RH and Dirichlet GRH alone.

\subsection{Primary reference and version convention}

Milne, J.~S. (1994). Motives over finite fields. In
U.~Jannsen, S.~Kleiman, \& J.-P.~Serre (Eds.),
\emph{Motives} (Proceedings of Symposia in Pure Mathematics,
Vol.~55, Part~1, pp.~401--459). American Mathematical
Society. Corrected author version dated November~2, 2019:
\url{https://www.jmilne.org/math/xnotes/Seattle1.pdf}.
All printed page references above refer to that 55-page
version; its theorem numbering agrees with the original
publication.  The sign in (CC-AG.25) was checked directly
in its typeset Theorem~2.16, p.~18.

% END consequences_criteria_rational_descent.tex

% BEGIN consequences_criteria_rational_basechange.tex
\section{The rational essential image of extension of the finite field}

The results here are concrete consequences of the classical
Milne classification applied in CC-AG08, within the established
abelian, Artin, and Tate numerical category.  They are not
claims of new literature results.  In contrast to the split
coefficient statement CC-AG07, they retain rational
coefficients and actual objects, including their positive
multiplicities.

Fix \(q=p^a\).  Let the target field \(B\) be either
\(\mathbb F_{q^m}\), for a positive integer \(m\), or
\(\overline{\mathbb F}_p\).  In the notation of CC-AG08,
extension of the base field acts on split labels by the
surjective, \(\Gamma\)-equivariant map
\[
 f_B:W(q)\longrightarrow\Lambda_B,\qquad
 \alpha\longmapsto
 \begin{cases}
       \alpha^m,& B=\mathbb F_{q^m},\\
       [\alpha,a],& B=\overline{\mathbb F}_p.
 \end{cases}                                                   \tag{CC-AG.34}
\]
Here \(\Gamma=\operatorname{Gal}(\overline{\mathbb Q}/\mathbb Q)\).
For a \(\Gamma\)-orbit \(O\) in the source or target, write
\(S_O\) for its rational simple object and \(d_O\) for its
Schur index, with the exact local-invariant convention of
CC-AG08.

\subsection{CC-AG09: the orbit and index criterion for the essential image}

For a target orbit \(O'\), define
\[
 \mathscr L_B(O')=
   \{O\subset W(q): O\text{ is a }\Gamma\text{-orbit},
                         \ f_B(O)=O'\}.
\]
For each such source orbit, put
\[
 c_O=\frac{|O|}{|O'|},\qquad
 r_O=\frac{c_Od_O}{d_{O'}},\qquad
 H_B(O')=\left\{\sum_{O\in\mathscr L_B(O')}a_Or_O:
       a_O\in\mathbb Z_{\ge0},\ \text{finite support}\right\}.
                                                               \tag{CC-AG.35}
\]
Then \(c_O\) and \(r_O\) are positive integers, and
\[
             (S_O)_B\simeq S_{O'}^{\oplus r_O}.                \tag{CC-AG.36}
\]
In particular, the formula includes the arithmetic divisibility
\(d_{O'}\mid c_Od_O\).

Let an actual rational target object have its unique simple
decomposition
\[
                       M=\bigoplus_{O'}S_{O'}^{\oplus b_{O'}}.
                                                               \tag{CC-AG.37}
\]
It belongs to the essential image of
\(\mathcal A_q\to\mathcal A_B\) if and only if
\[
                         b_{O'}\in H_B(O')
                       \quad\text{for every }O'.              \tag{CC-AG.38}
\]
The rational isomorphism classes of its preimages are exactly
the finite-support nonnegative integer solutions of
\[
                  \sum_{O\in\mathscr L_B(O')}a_Or_O=b_{O'}
                      \quad\text{for every }O'.               \tag{CC-AG.39}
\]
No unspecified object-descent condition remains in this
criterion.

For a finite target extension, \(\mathscr L_B(O')\) is finite:
its source labels are precisely the roots of
\[
                      \prod_{\beta\in O'}(t^m-\beta)
                               \ \in\mathbb Q[t].             \tag{CC-AG.40}
\]
Thus the criterion consists of finitely many field and index
calculations and nonnegative integer equations for a given
target object.  For the geometric target the lift set can be
infinite, but each additive monoid \(H_B(O')\subset\mathbb N\)
is still finitely generated.  This latter assertion does
not supply an effective bound on the orders of roots of
unity that one must inspect.

On Grothendieck groups, the corresponding coefficient condition
is instead divisibility by
\(\gcd\{r_O:O\in\mathscr L_B(O')\}\).  That is the group
completion of (CC-AG.35), rather than a replacement for
the stated actual-object condition.

\paragraph{Complete proof.}
The image of the transitive finite \(\Gamma\)-set \(O\)
under a \(\Gamma\)-equivariant map is the transitive set
\(O'\).  Its fibres all have the same cardinality \(c_O\):
an element of \(\Gamma\) carrying one target label to another
gives a bijection between the corresponding fibres.
Consequently \(c_O=|O|/|O'|\) is a positive integer.
If \(\alpha\in O\), the stabilizer inclusion also identifies
\[
       c_O=[K_\alpha:K_{f_B(\alpha)}].
                                                               \tag{CC-AG.41}
\]
Here \(K_\alpha=\mathbb Q(\alpha)\); for a germ the target
field is its stabilized centre from (CC-AG.24).
The field inclusion is literal for the chosen labels,
because a sufficiently large power representing the germ
lies in \(\mathbb Q(\alpha)\).

After extending coefficients to \(E\), CC-AG08 gives
\((S_O)_E\) multiplicity \(d_O\) on each label of \(O\).
Extension of the base field merges exactly the labels
having the same image under \(f_B\).  Thus \((S_O)_{B,E}\)
has multiplicity \(c_Od_O\) on every label of \(O'\) and
zero elsewhere.
The rational classification in the target says that a
rational object with this support is a sum of copies of
the single simple \(S_{O'}\), whose split multiplicity is
\(d_{O'}\) at each of these labels.  Its number of copies is
therefore \(r_O=c_Od_O/d_{O'}\), a positive integer.
This proves (CC-AG.36) and the asserted index divisibility.

Every rational source object is a unique finite sum
\(\bigoplus_O S_O^{\oplus a_O}\).  Applying (CC-AG.36)
gives exactly (CC-AG.39).  This proves necessity of
(CC-AG.38).  Conversely, for each of the finitely many
target orbits with \(b_{O'}>0\), membership in the monoid
supplies a finite nonnegative expression for \(b_{O'}\).
Taking the corresponding finite sum of source simples
constructs a preimage.  Simples mapping to different
target orbits cannot interfere with these equations, and
every \(r_O\) is positive.  Thus a target coefficient zero
forces all the associated source coefficients to be zero.
Uniqueness of simple decompositions proves the preimage
classification (CC-AG.39), including the zero object.

For a finite extension, every eligible label is an
\(m\)-th root of a label of \(O'\), and every such root
is a Weil \(q\)-number by the root argument in CC-AG07.
The product (CC-AG.40) is Galois invariant and hence has
rational coefficients.  Its roots form a finite set,
partitioned by the source orbits, proving finiteness.
Their centre fields and the local-invariant formula of
CC-AG08 give the stated arithmetic data.

For the geometric target, surjectivity of (CC-AG.34)
ensures that the set of \(r_O\)'s is nonempty.  Let
\(c\) be its least member; it is also the least positive
element of \(H_B(O')\).  For each residue class modulo
\(c\) which meets \(H_B(O')\), let \(h_i\) be the least
member in that class.  Every \(h\in H_B(O')\) is of
the form \(h_i+tc\) with an integer \(t\ge0\).
Hence \(H_B(O')\) is generated by \(c\) and the finitely
many \(h_i\).  Each \(h_i\), by definition of the monoid,
is itself a finite sum of the original \(r_O\)'s.
Taking the union of these finitely many finite supports
shows that finitely many of the actual lift indices
already generate \(H_B(O')\).  The argument proves
existence, without an effective cutoff on the lift set.

Finally, allowing arbitrary integer, rather than
nonnegative integer, source multiplicities gives the
subgroup of \(\mathbb Z\) generated by the \(r_O\)'s.
That subgroup is their greatest common divisor times
\(\mathbb Z\).  Independence of the target orbit
components proves the Grothendieck-group assertion.
\(\square\)

\subsection{CC-AG10: an exact even-multiplicity obstruction over
\(\mathbb F_3\)}

Set
\[
 K=\mathbb Q(s),\quad s^2=-5,\qquad
 \beta=2+s\in W(9),\qquad
 \xi=[\beta,2]\in\mathcal W_3.
                                                               \tag{CC-AG.42}
\]
Let \(T_9\in\mathcal A_9\) and
\(T_{\mathrm{geom}}\in\mathcal A_{\overline{\mathbb F}_3}\)
be the simple rational objects of the orbits of
\(\beta\) and \(\xi\), respectively.  Both have rank two,
endomorphism field \(K\), and Schur index one, and
\((T_9)_{\overline{\mathbb F}_3}\simeq T_{\mathrm{geom}}\).

For every nonnegative integer \(b\), either one of the
objects
\[
                T_9^{\oplus b},\qquad
                T_{\mathrm{geom}}^{\oplus b}
\]
is the base extension of an object in \(\mathcal A_3\)
if and only if \(b\) is even.
Thus the geometric multiplicity monoid in (CC-AG.35)
for this orbit is exactly \(2\mathbb N\).  In particular,
extension of the base field
\(\mathcal A_3\to\mathcal A_{\overline{\mathbb F}_3}\)
is not essentially surjective over rational coefficients.

\paragraph{Complete proof.}
We first establish the small number-field facts explicitly.
The ring of integers of \(K\) is
\(\mathcal O_K=\mathbb Z[s]\).  Indeed, if \(x+ys\)
with \(x,y\in\mathbb Q\) is integral, its trace and norm
are integers, so \(2x\in\mathbb Z\) and
\(x^2+5y^2\in\mathbb Z\).  Their discriminant
\(-20y^2\) is an integer.  Writing \(y\) in lowest
terms shows that its denominator is at most two, because
its denominator squared divides \(20\).
Thus \(x=A/2,\ y=B/2\) for integers \(A,B\), and
\[
                         A^2+5B^2\equiv0\pmod4
\]
forces both \(A\) and \(B\) to be even.  This proves one
inclusion; the reverse inclusion is immediate because
\(s\) is integral.

The two ideals
\[
       \mathfrak p=(3,s-1),\qquad
       \overline{\mathfrak p}=(3,s+1)
\]
are distinct prime ideals with residue field \(\mathbb F_3\).
Reduction of \(s^2+5\) modulo \(3\) gives
\((s-1)(s+1)\) with distinct roots, and the Chinese
remainder theorem gives
\[
                     (3)=\mathfrak p\overline{\mathfrak p}.
\]
Now \(\beta\) vanishes modulo \(\mathfrak p\) and is
equal to \(1\) modulo \(\overline{\mathfrak p}\).
Its field norm is \(4+5=9\).
Unique ideal factorization and the multiplicativity of
the norm therefore give
\[
                           (\beta)=\mathfrak p^2.             \tag{CC-AG.43}
\]
The ideal \(\mathfrak p\) is not principal.  A generator
would be an integer \(x+ys\in\mathcal O_K\) satisfying
\[
                           x^2+5y^2=3,
\]
which has no integral solution: nonzero \(y\) makes the
left side at least five, and \(y=0\) would make three
a square.  Consequently the ideal class of
\(\mathfrak p\) has order exactly two.  The ideal facts
used here are Milne (2020), Theorem~3.20, p.~53, and
Propositions~4.1--4.2, p.~69; no class-number table is
being used.

The two complex conjugates of \(\beta\) have absolute
value \(3\), and (CC-AG.43) shows that it is a \(3\)-unit.
Hence \(\beta\) is a Weil \(9\)-number of weight one.
Every root of unity in \(K\) is an integral unit; its
positive norm is one, and the equation
\(x^2+5y^2=1\) shows that it is \(1\) or \(-1\).
The ratio \(\beta/\overline\beta\) is neither \(1\)
nor \(-1\), since \(\beta\) has both nonzero real and
nonzero imaginary parts.  It is therefore not a root
of unity.  It follows that
\(\mathbb Q(\beta^N)=K\) for every \(N\ge1\):
otherwise the two conjugates would have equal
\(N\)-th powers.  In particular, \(K_\xi=K\), and the
germ orbit has two elements.

The two places above \(3\) have local degree one.
At them the slopes of \(\beta\), relative to \(q_0=9\),
are respectively \(2/2=1\) and \(0/2=0\).
There are no real places.  Formula (CC-AG.25) therefore
makes every local invariant zero, both for \(\beta\)
and for its stabilized germ.  It follows that
\[
       d_\beta=d_\xi=1,\qquad
       D_\beta=D_\xi=K,\qquad
       \operatorname{rank}T_9
         =\operatorname{rank}T_{\mathrm{geom}}=2.
                                                               \tag{CC-AG.44}
\]
The label-to-germ map is a bijection on these two-element
orbits.  Equation (CC-AG.36) now gives the claimed
geometric base change of \(T_9\).

Consider any source orbit which maps to the orbit of
\(\xi\).  Replacing a chosen label \(\alpha\in W(3)\)
by a Galois conjugate if necessary, arrange that
\([\alpha,1]=\xi\).  By the definition of a germ,
for some integer \(N\ge1\),
\[
                            \alpha^{2N}=\beta^N.             \tag{CC-AG.45}
\]
Put \(L=\mathbb Q(\alpha)\).
Since \(\mathbb Q(\beta^N)=K\), this equation proves
the required inclusion \(K\subset L\).
It also proves that \(\alpha\) is integral: it is a
root of a monic polynomial with coefficient
\(\beta^N\), an algebraic integer.
Passing to fractional ideals in \(\mathcal O_L\),
(CC-AG.43)--(CC-AG.45) give
\[
                      (\alpha)^{2N}
                           =(\mathfrak p\mathcal O_L)^{2N}.
\]
The group of fractional ideals is free abelian on the
prime ideals, so its power maps are injective and
\[
                            (\alpha)=\mathfrak p\mathcal O_L.
                                                               \tag{CC-AG.46}
\]
The relative ideal norm sends this equality to
\[
             \bigl(N_{L/K}(\alpha)\bigr)
                        =\mathfrak p^{[L:K]}.                 \tag{CC-AG.47}
\]
This is exactly Proposition~4.1(a),(c) of Milne (2020),
p.~69; the extension \(L/K\) need not be Galois.
Since the class of \(\mathfrak p\) has order two,
(CC-AG.47) forces \([L:K]\) to be even.

For clarity, the source Schur index is also exactly one.
The field \(L\) contains the imaginary quadratic field
\(K\), so it has no real embeddings.  At a place \(v\)
above \(3\), let \(e_v,f_v\) be its ramification index
and residue degree.  The normalized integral valuation
satisfies
\(\operatorname{ord}_v(3)=e_v\) and
\([L_v:\mathbb Q_3]=e_vf_v\).
Since the source field is \(\mathbb F_3\), its invariant
from (CC-AG.25) is
\[
        -\frac{\operatorname{ord}_v(\alpha)}{e_v}
                          e_vf_v
          =-\operatorname{ord}_v(\alpha)f_v
          \equiv0\pmod{\mathbb Z}.
                                                               \tag{CC-AG.48}
\]
All other invariants vanish as well.  Hence \(d_\alpha=1\).
The transition multiplicity is therefore
\[
        r_O=\frac{[L:\mathbb Q]}{[K:\mathbb Q]}
                              \frac{d_\alpha}{d_\xi}
             =[L:K],
                                                               \tag{CC-AG.49}
\]
which is even.  Every rational source object is a sum
of such simples, proving that a geometric preimage can
exist only in even multiplicity.

Finally, choose \(\alpha\) with \(\alpha^2=\beta\).
It is a Weil \(3\)-number by the root calculation
of CC-AG07.  The field \(L=\mathbb Q(\alpha)\)
contains \(K\).  The element \(\beta\) is not a
square in \(K\): a square root in \(K\) would be
integral and would have positive norm three, contrary
to the equation already considered.  Thus
\([L:K]=2\).  Equivalently, the minimal polynomial
of \(\alpha\) is the irreducible polynomial
\[
                              t^4-4t^2+9.
\]
The index computation (CC-AG.48) applies, so \(d_\alpha=1\),
and its simple rational object has rank four.
Equation (CC-AG.36) gives two copies of
\(T_{\mathrm{geom}}\).  Finite extension to \(\mathbb F_9\)
already gives two copies of \(T_9\), since the power map
has orbit fibre size two and both relevant indices
are one.  Taking arbitrary finite sums realizes every
even multiplicity.
An odd-multiplicity finite-field preimage would give an
odd-multiplicity geometric preimage after further base
extension, which has just been excluded.
This proves both asserted equivalences and
\(H_{\overline{\mathbb F}_3}(O_\xi)=2\mathbb N\).
\(\square\)

\subsection{Scope and primary references}

The criterion describes the exact rational essential image
within \(\mathcal A_q\).  It does not assert that every
smooth projective variety belongs to that category, that
numerical projectors lift to Chow projectors, or that the
rational category has a neutral fibre functor.  The latter
has already been disproved by AQ-RATIONAL-FIBRE.  The
arbitrary-variety Tate gate remains the one separated in
CC-XC01--04 and BR-EFF01.

\begin{itemize}
\item Milne, J.~S. (1994). Motives over finite fields. In
U.~Jannsen, S.~Kleiman, \& J.-P.~Serre (Eds.),
\emph{Motives} (Proceedings of Symposia in Pure Mathematics,
Vol.~55, Part~1, pp.~401--459). American Mathematical
Society. Corrected author version dated November~2, 2019:
\url{https://www.jmilne.org/math/xnotes/Seattle1.pdf}.
The classification, rank, invariant, and scalar-extension
inputs are the exact statements listed in CC-AG08.
\item Milne, J.~S. (2020, July~19).
\emph{Algebraic number theory} (Version~3.08)
[Author course notes].
\url{https://www.jmilne.org/math/CourseNotes/ANT.pdf}.
Exact ideal inputs: Theorem~3.20, printed p.~53;
Proposition~4.1(a),(c) and Proposition~4.2(a),
printed p.~69.  These are ideal factorization and norm
identities, rather than a precomputed assertion about
the example's ideal class.
\end{itemize}

% END consequences_criteria_rational_basechange.tex

% BEGIN consequences_criteria_integral_slope_descent.tex
\section{A finite rational descent test for integral-slope germs}

This section makes the geometric multiplicity criterion CC-AG09
effective in a specified sector.  It is a deduction from the
checked Milne classification, elementary ideal theory, and
Kronecker's root-of-unity criterion.  The category remains the
abelian, Artin, and Tate numerical category.  No novelty claim
or extension to arbitrary motives is made.

\subsection{CC-AG11: a cyclic multiplicity monoid with a finite root test}

Let \(\xi\in\mathcal W_p\) have weight \(w\in\mathbb Z\).
Choose a stabilized representative
\[
          \xi=[\pi,b],\qquad \pi\in K,\quad b\ge1,\qquad
          K=K_\xi=\mathbb Q(\pi^N)\quad(N\ge1).
                                                               \tag{CC-AG.50}
\]
Such a representative exists by (CC-AG.24).
For each place \(v\mid p\) of \(K\), normalize its
valuation to have image \(\mathbb Z\), put
\(e_v=\operatorname{ord}_v(p)\), and define
\[
       \lambda_v=\frac{\operatorname{ord}_v(\pi)}{b e_v}.
                                                               \tag{CC-AG.51}
\]
Assume throughout this theorem that every
\(\lambda_v\) is an integer.  Define the fractional ideal
\[
                   J=\prod_{v\mid p}\mathfrak p_v^{e_v\lambda_v}
                             \quad\text{of }K.                \tag{CC-AG.52}
\]
Negative exponents are allowed.  In particular,
\((\pi)=J^b\).

Let
\[
 \mathscr E_\xi=
 \left\{n\in\mathbb Z:\ \text{there is }\gamma\in K^\times
       \text{ with }(\gamma)=J^n
       \text{ and }|\sigma\gamma|=p^{nw/2}
                    \text{ for every }\sigma:K\hookrightarrow\mathbb C
 \right\}.
                                                               \tag{CC-AG.53}
\]
Then \(\mathscr E_\xi=h\mathbb Z\) for a positive integer
\(h\) dividing \(b\).
This integer can be found by the finite test
\[
 h=\min\left\{n>0:n\mid b,\quad
       X^b-\zeta\pi^n\text{ has a root in }K
                      \text{ for some }\zeta\in\mu(K)\right\}.
                                                               \tag{CC-AG.54}
\]
Here \(\mu(K)\) is the finite group of roots of unity in \(K\).
Thus only divisors of a given exponent and roots of finitely
many specified polynomials over a fixed number field are
required.

Let \(S_\xi\) denote the rational simple object of the
\(\Gamma\)-orbit of \(\xi\) over
\(\overline{\mathbb F}_p\).  For every \(a\ge1\), its
exact multiplicity monoid for descent to \(\mathbb F_{p^a}\)
is
\[
 H_{\overline{\mathbb F}_p}(O_\xi)
                 =\frac{h}{\gcd(h,a)}\,\mathbb N.
                                                               \tag{CC-AG.55}
\]
Equivalently, for every \(r\ge0\),
\[
 S_\xi^{\oplus r}\text{ is a base extension from }
             \mathcal A_{p^a}
       \quad\Longleftrightarrow\quad
                  \frac{h}{\gcd(h,a)}\mid r.                 \tag{CC-AG.56}
\]
Every relevant target and source simple has Schur index one.
One source simple already realizes the least positive
multiplicity.  More precisely, choose \(\gamma\) in
(CC-AG.53) for \(n=h\), and put \(g=\gcd(h,a)\).
Any algebraic root
\[
                     \alpha^{h/g}=\gamma^{a/g}               \tag{CC-AG.57}
\]
is a Weil \(p^a\)-number of weight \(w\), has germ \(\xi\),
and satisfies
\[
                [\mathbb Q(\alpha):K]=h/g.                  \tag{CC-AG.58}
\]
This gives an actual finite-degree construction of a
minimal preimage.

\paragraph{Complete proof.}
The set in (CC-AG.53) contains zero and is closed under
addition and negation, by multiplying and inverting its
witnesses.  The element \(b\) belongs to it, witnessed
by \(\pi\), since (CC-AG.51)--(CC-AG.52) give
\((\pi)=J^b\) and \(\pi\) has the stated Weil absolute
values.  Hence \(\mathscr E_\xi\) is a nonzero subgroup
of \(\mathbb Z\), equal to \(h\mathbb Z\) with \(h\ge1\)
and \(h\mid b\).

We verify the finite test in both directions.
If \(\gamma\) witnesses \(n\in\mathscr E_\xi\), then
\[
                              u=\gamma^b/\pi^n
                                                               \tag{CC-AG.59}
\]
has trivial fractional ideal, and every complex conjugate
has absolute value one.  Thus \(u\) is an algebraic
integer and a unit.  Kronecker's criterion makes it
a root of unity, and it lies in \(K\), so
\(u\in\mu(K)\).  Conversely, suppose
\(\gamma\in K\) satisfies
\(\gamma^b=\zeta\pi^n\) for \(\zeta\in\mu(K)\).
It is nonzero.  Taking fractional ideals and using
their torsion-free group gives
\[
                    (\gamma)^b=J^{bn}
                       \quad\Longrightarrow\quad
                    (\gamma)=J^n.
 \]
Taking the absolute value at every complex embedding
gives \(|\sigma\gamma|=p^{nw/2}\).
Therefore \(\gamma\) witnesses membership in
\(\mathscr E_\xi\).  Since its positive generator is a
divisor of \(b\), testing the positive divisors of \(b\)
is sufficient and proves (CC-AG.54).

The precise root-of-unity input is Milne (2020),
Proposition~5.5 and Corollary~5.6, printed p.~87.
The elementary reason is that the minimal polynomials
of the powers of an algebraic integer whose conjugates
have absolute value one have uniformly bounded degree
and integer coefficients bounded by elementary symmetric
functions of numbers of modulus one.  Only finitely
many polynomials, and hence finitely many powers, can
occur.  Two powers coincide, which proves torsion.
The same bounded-polynomial argument proves the
finiteness of \(\mu(K)\), and supplies an explicit finite
way to determine it if it is not already known.  Put
\[
 R_D=\sum_{e=1}^{D}e\prod_{j=1}^{e}
                    \left(2\binom ej+1\right),
 \qquad D=[K:\mathbb Q],\qquad
 M_D=\operatorname{lcm}(1,2,\ldots,R_D).
\]
There are at most \(R_D\) complex numbers whose monic
minimal polynomials have degree at most \(D\) and the
coefficient bounds just described.  All the distinct
powers of a root of unity in \(K\) are among them, so
its order is at most \(R_D\).  Consequently \(\mu(K)\)
is exactly the set of roots in \(K\) of
\(X^{M_D}-1\).  This deliberately large bound requires
no search without a stopping condition.
Formula (CC-AG.54) is therefore a finite collection
of exact number-field root questions.

We next justify the Schur-index assertion, including
the archimedean places.  If \(K\) has a real embedding,
the Weil condition on \(\pi\) implies
\(\pi^2=p^{bw}\): its image under that real embedding
has square \(p^{bw}\), and the embedding is injective.
The stabilized condition (CC-AG.50) then gives
\(K=\mathbb Q(\pi^2)=\mathbb Q\).
Its single \(p\)-adic slope is \(w/2\), so the
integral-slope assumption forces \(w\) to be even.
Otherwise \(K\) has no real embeddings.
Consequently every real-place invariant of \(S_\xi\)
vanishes.  Its \(p\)-adic invariant at \(v\) is
\[
                       -\lambda_v[K_v:\mathbb Q_p],
\]
an integer, and all other invariants vanish by
(CC-AG.25).  Hence \(d_\xi=1\).

Take any source label \(\alpha\in W(p^a)\) mapping to
the orbit of \(\xi\).  Conjugating it if necessary,
assume that \([\alpha,a]=[\pi,b]\).  Some \(N\ge1\)
then satisfies
\[
                           \alpha^{bN}=\pi^{aN}.             \tag{CC-AG.60}
\]
Set \(L=\mathbb Q(\alpha)\) and \(d=[L:K]\).
The stabilized condition gives
\(K=\mathbb Q(\pi^{aN})\subset L\), so this degree
is well-defined.
Taking ideals in \(L\) in (CC-AG.60) and using
torsion-freeness gives
\[
                              (\alpha)=J^a\mathcal O_L.
                                                               \tag{CC-AG.61}
\]
At a place \(z\) of \(L\) above a place \(v\mid p\)
of \(K\), the corresponding normalized slope of
\(\alpha\), relative to \(p^a\), is exactly
\(\lambda_v\); this follows either from
(CC-AG.60) or from (CC-AG.61), accounting for the
ramification index in \(L/K\).
It is an integer.  If \(L\) has a real embedding,
its restriction makes \(K\) real, so \(w\) is even
by the preceding paragraph.
Thus the local-invariant formula also gives
\(d_\alpha=1\).

Now apply the relative field norm to \(\alpha\).
The ideal norm identities, which hold without a
Galois assumption, give
\[
                  \bigl(N_{L/K}(\alpha)\bigr)=J^{ad}.
                                                               \tag{CC-AG.62}
\]
For each embedding \(\sigma:K\hookrightarrow\mathbb C\),
the conjugate of this norm is the product over the
\(d\) embeddings of \(L\) extending \(\sigma\).
Each factor has absolute value \(p^{aw/2}\).
Therefore the norm is a witness for
\(ad\in\mathscr E_\xi=h\mathbb Z\).
It follows that
\[
                           h/\gcd(h,a)\mid d.                 \tag{CC-AG.63}
\]
The ideal and field norm formulas used here are Milne
(2020), Proposition~4.1(a),(c), p.~69, and
Proposition~2.19, p.~31, respectively.
Because both indices are one, (CC-AG.35)--(CC-AG.41)
identify the transition multiplicity of this source
simple with \(d\).
Thus every allowed multiplicity is divisible by
\(h/\gcd(h,a)\).

It remains to construct that least positive
multiplicity.  Choose \(\gamma\) witnessing exponent
\(h\) and let \(\alpha\) satisfy (CC-AG.57).
Every conjugate of \(\alpha\) has absolute value
\(p^{aw/2}\).  Valuations away from \(p\) vanish,
because those of \(\gamma\) vanish, so
\(\alpha\in W(p^a)\).
Equation (CC-AG.59) for \(n=h\) shows that
\([\gamma,h]=[\pi,b]=\xi\).
Moreover (CC-AG.57), raised to the \(g\)-th power,
is \(\alpha^h=\gamma^a\), so
\([\alpha,a]=[\gamma,h]=\xi\).

The germ centre is \(K\).  Since \(\gamma\in K\),
every positive power of \(\gamma\) generates \(K\):
its field contains the eventual stabilized germ
centre and is contained in \(K\).
Consequently
\(\gamma^{a/g}=\alpha^{h/g}\) proves
\(K\subset\mathbb Q(\alpha)\).
The defining equation bounds this extension degree
above by \(h/g\), whereas (CC-AG.63) says it is a
positive multiple of \(h/g\).
It must therefore equal \(h/g\), proving
(CC-AG.58).  Its source simple realizes exactly
\(h/g\) copies of \(S_\xi\).
Direct sums realize all nonnegative multiples, and
the preceding divisibility excludes every other
multiplicity.  Equations (CC-AG.55)--(CC-AG.56)
follow. \(\square\)

\paragraph{Ordinary imaginary-quadratic special case.}
Suppose \(K\) is imaginary quadratic, \(w=1\), and the
germ is ordinary, so that \(p\) splits and the two
slopes are zero and one.  Let \(\mathfrak p\) be the
prime where the slope is one.
Then \(J=\mathfrak p\), and
\[
                           h=\operatorname{ord}([\mathfrak p])
                                  \quad\text{in }\operatorname{Cl}(K).
                                                               \tag{CC-AG.64}
\]
Indeed, a Weil generator of \(J^n\) makes \(J^n\)
principal.  Conversely, if
\(J^n=(\gamma)\), then
the positive field norm of \(\gamma\) is \(p^n\);
the two complex conjugates therefore each have
absolute value \(p^{n/2}\).  Thus this generator
satisfies (CC-AG.53).  The two exponent subgroups
are identical, proving (CC-AG.64).
The exact multiplicity for descent to \(\mathbb F_{p^a}\)
is consequently
\[
             \frac{\operatorname{ord}([\mathfrak p])}
                  {\gcd(\operatorname{ord}([\mathfrak p]),a)}.
\]
For the germ of CC-AG10, this order is two and
\(a=1\), recovering its exact even-multiplicity
obstruction.  For \(a=2\) the least multiplicity
is one.

\paragraph{Exact boundary.}
The integral-slope hypothesis makes every Schur index
in this proof one.  With fractional slopes the local
indices can vary with the source centre field, and
the transition multiplicity is the full expression
\([K_\alpha:K_\xi]d_\alpha/d_\xi\).
Neither the cyclic-monoid formula (CC-AG.55) nor the
root-degree construction is asserted in that remaining
case.  This result removes the previously unbounded
root-order search precisely in the stated sector.
It does not provide new cross-correspondences for an
arbitrary variety, so the separate arbitrary-motive
Tate gate is unchanged.

\subsection{Primary references}

\begin{itemize}
\item Milne, J.~S. (1994). Motives over finite fields. In
U.~Jannsen, S.~Kleiman, \& J.-P.~Serre (Eds.),
\emph{Motives} (Proceedings of Symposia in Pure Mathematics,
Vol.~55, Part~1, pp.~401--459). American Mathematical
Society. Corrected author version dated November~2, 2019:
\url{https://www.jmilne.org/math/xnotes/Seattle1.pdf}.
Exact classification and local-index locators are
Theorems~2.15--2.16, p.~18, and Theorems~2.18--2.19,
p.~20, with the geometric-Frobenius normalization
on p.~17.
\item Milne, J.~S. (2020, July~19).
\emph{Algebraic number theory} (Version~3.08)
[Author course notes].
\url{https://www.jmilne.org/math/CourseNotes/ANT.pdf}.
Exact inputs: Proposition~2.19, p.~31;
Theorem~3.20, p.~53; Proposition~4.1(a),(c),
p.~69; Proposition~5.5 and Corollary~5.6, p.~87.
\end{itemize}

% END consequences_criteria_integral_slope_descent.tex

% BEGIN consequences_criteria_finite_descent.tex
\section{A finite decision procedure for each specified rational descent}

\subsection{CC-AG12: the rank bound removes the root-order search}

Retain the established rational abelian, Artin, and Tate
numerical categories and the exact Milne orbit and
Schur-index classification of CC-AG08.  Fix a source
field \(\mathbb F_{p^a}\), with \(a\ge1\).
No integral-slope hypothesis is imposed in this section.

First consider one rational simple geometric object
\(S_\xi\), with stabilized representative and centre
\[
         \xi=[\pi,b],\qquad b\ge1,\qquad
         K=K_\xi=\mathbb Q(\pi^N)\quad(N\ge1).
                                                               \tag{CC-AG.65}
\]
Let \(d_\xi\) be its Schur index, and let \(r\ge1\)
be a specified target multiplicity.  Its rank is
\[
                         D=r\,d_\xi[K:\mathbb Q].
                                                               \tag{CC-AG.66}
\]
Define the explicit integers
\[
 R_D=\sum_{e=1}^{D}e\prod_{j=1}^{e}
                           \left(2\binom ej+1\right),
       \qquad N_D=\operatorname{lcm}(1,2,\ldots,R_D),
                                                               \tag{CC-AG.67}
\]
and the monic rational polynomial
\[
 P_{\xi,a,r}(X)=
   \prod_{\sigma:K\hookrightarrow\overline{\mathbb Q}}
       \left(X^{bN_D}-\sigma(\pi)^{aN_D}\right)
                            \quad\in\mathbb Q[X].             \tag{CC-AG.68}
\]
Factor it into irreducible polynomials over \(\mathbb Q\).
Let \(O_1,\ldots,O_t\) be their root orbits, and choose
\(\alpha_j\in O_j\).  Every \(\alpha_j\) is a Weil
\(p^a\)-number whose germ belongs to \(O_\xi\).
Compute \(d_j\) from the full local-invariant formula
(CC-AG.25), and put
\[
        m_j=
          \frac{[\mathbb Q(\alpha_j):\mathbb Q]}
               {[K:\mathbb Q]}\,
          \frac{d_j}{d_\xi}.
                                                               \tag{CC-AG.69}
\]
These are positive integers.

Then the exact rational descent criterion is the finite
nonnegative-integer problem
\[
 \begin{split}
 S_\xi^{\oplus r}\text{ descends to }\mathcal A_{p^a}
   \quad\Longleftrightarrow\quad
   &\text{there are }u_1,\ldots,u_t\in\mathbb Z_{\ge0}\\
   &\text{such that }\sum_{j=1}^{t}u_jm_j=r.
 \end{split}                                                  \tag{CC-AG.70}
\]
Factors with \(m_j>r\) can be discarded.
Moreover all rational isomorphism classes of preimages,
without omission, are
\[
                         \bigoplus_{j=1}^{t}
                                    S_{O_j}^{\oplus u_j}
                                                               \tag{CC-AG.71}
\]
as the finite solutions of (CC-AG.70) vary.

For any specified geometric object with finitely many
simple isotypic components, perform this finite test on
each component.  The results combine independently.
Thus actual rational descent of every specified object
in the established abelian sector is decidable by finite
algebraic calculations, and there are only finitely many
rational isomorphism classes of its preimages over a
fixed finite field.  For a zero component the preimage
is zero and no polynomial is required.

\paragraph{Complete proof.}
We first recall the elementary bound in (CC-AG.67).
If a root of unity \(\zeta\) has degree at most \(D\),
then all of its powers are algebraic integers of degree
at most \(D\), with every conjugate of absolute value
one.  A monic minimal polynomial of degree \(e\le D\)
has its \(j\)-th coefficient bounded in absolute value
by \(\binom ej\), because that coefficient is an
elementary symmetric function of roots of modulus one.
There are at most
\(\prod_{j=1}^{e}(2\binom ej+1)\) possible integer
coefficient lists in degree \(e\), each with at most
\(e\) distinct roots.  Therefore the set of all such
algebraic numbers has cardinality at most \(R_D\).
The distinct powers of \(\zeta\) lie in that set, so
its order is at most \(R_D\) and divides \(N_D\).
This is the bounded-polynomial argument of Milne
(2020), Proposition~5.5 and Corollary~5.6, p.~87,
with its bound made explicit.

Suppose a source object descends to
\(S_\xi^{\oplus r}\).  Decompose that source into
rational simples.  By CC-AG09, every source simple
which occurs has a label \(\alpha\in W(p^a)\)
whose germ belongs to \(O_\xi\), and its geometric
base change is a positive number of copies of
\(S_\xi\).  In particular its rank is at most \(D\).
Conjugate the label if necessary to arrange
\([\alpha,a]=[\pi,b]\), and let
\(L=\mathbb Q(\alpha)\).
The source rank formula gives
\[
                      [L:\mathbb Q]
                       \le d_\alpha[L:\mathbb Q]
                       \le D.
                                                               \tag{CC-AG.72}
\]
Germ equality gives
\(\alpha^{bN}=\pi^{aN}\) for some \(N\ge1\).
Since the representative is stabilized,
\(K=\mathbb Q(\pi^{aN})\subset L\).
Consequently
\[
                              \theta=\alpha^b/\pi^a
\]
is a root of unity \emph{in \(L\)}.
Its degree over \(\mathbb Q\) is at most \(D\).
The bound just proved implies \(\theta^{N_D}=1\),
and hence
\[
                          \alpha^{bN_D}=\pi^{aN_D}.
                                                               \tag{CC-AG.73}
\]
Thus every possible source simple has its entire
Galois orbit among the roots of (CC-AG.68).
This is the step that makes the search finite even
when the original fibre over a germ contains
infinitely many root-of-unity twists.

We also verify that every root listed in (CC-AG.68)
has the correct label conditions.  The polynomial
has rational coefficients because the product is
Galois invariant.  If \(\alpha\) is a root, then
for one embedding \(\sigma\) it satisfies
\[
               \alpha^{bN_D}=\sigma(\pi)^{aN_D}.
                                                               \tag{CC-AG.74}
\]
Stabilization implies
\(\sigma K=\mathbb Q(\sigma(\pi)^{aN_D})
                  \subset\mathbb Q(\alpha)\).
For every complex embedding of \(\mathbb Q(\alpha)\),
taking absolute values in (CC-AG.74) gives
\(|\alpha'|=p^{aw/2}\), where \(w\) is the weight
of \(\xi\).
At every finite place away from \(p\), the same
equation gives valuation zero.
Thus \(\alpha\in W(p^a)\).
Moreover (CC-AG.74) gives
\([\alpha,a]=[\sigma\pi,b]=\sigma\xi\).
All candidate roots therefore have the correct
weight, \(p\)-unit condition, and germ orbit;
no unproved enlargement of the category is involved.

Distinct embeddings of \(K\) give distinct values
of \(\sigma(\pi)^{aN_D}\), since that power
generates \(K\).  Each binomial in (CC-AG.68)
is separable in characteristic zero, and its
constant term is nonzero.  The product is consequently
separable.  Its irreducible factors thus give a
finite list of distinct source simple classes.

Apply the exact base-change formula CC-AG09 to
each of those classes.  Its image is
\(S_\xi^{\oplus m_j}\), with the full index
ratio (CC-AG.69); this proves the asserted positive
integrality of \(m_j\).
Necessity of (CC-AG.70) now follows from the
completeness of the candidate list.
Conversely, each nonnegative solution constructs
the actual rational object (CC-AG.71), and its
base change is exactly \(S_\xi^{\oplus r}\).
Uniqueness of simple decompositions identifies
distinct solutions with distinct rational
isomorphism classes.
Since every \(m_j>0\), each \(u_j\le r/m_j\).
Only finitely many solutions exist, and factors
with \(m_j>r\) cannot occur.

Finally, a source simple has isotypic geometric
image in one target orbit, so source summands
for different target orbits cannot contribute
to one another.  Decomposing a general target
into its finitely many isotypic components
reduces the problem to the preceding tests
and their finite Cartesian product of solution
sets.  The zero target has only the zero
preimage because all nonzero simples have
positive rank.  This proves all assertions.
\(\square\)

\paragraph{Scope of effectivity.}
The input is the finite orbit and index data of a
specified object.  The output is its complete list
of rational simple-label multiplicities over the
specified source field.  Exact factorization and
the local-invariant computations are finite
number-field calculations; the bound is not
intended as an efficient algorithm.
No explicit algebraic cycle representatives for
the numerical projectors are claimed.

The integer \(D\), and hence the polynomial cutoff,
depends on the target multiplicity.  The theorem
does not give one cutoff for an entire geometric
multiplicity monoid, and does not claim that all
such monoids are cyclic.  In particular it keeps
the fractional-slope Schur indices in the formula.
It does resolve the formerly unbounded root-order
search for every specified object in every slope
sector.  The separate arbitrary-variety
cross-correspondence and full Tate conditions
are unchanged.

\subsection{Primary references}

\begin{itemize}
\item Milne, J.~S. (1994). Motives over finite fields.
In U.~Jannsen, S.~Kleiman, \& J.-P.~Serre (Eds.),
\emph{Motives} (Proceedings of Symposia in Pure
Mathematics, Vol.~55, Part~1, pp.~401--459).
American Mathematical Society.
Corrected author version dated November~2, 2019:
\url{https://www.jmilne.org/math/xnotes/Seattle1.pdf}.
The precise classification, rank, centre, and
local-index inputs are those inspected and
recorded in CC-AG08, especially Proposition~2.8,
p.~14, Theorems~2.15--2.16, p.~18, and
Theorems~2.18--2.19, p.~20.
\item Milne, J.~S. (2020, July~19).
\emph{Algebraic number theory} (Version~3.08)
[Author course notes].
\url{https://www.jmilne.org/math/CourseNotes/ANT.pdf}.
Proposition~5.5 and Corollary~5.6, printed p.~87.
\end{itemize}

% END consequences_criteria_finite_descent.tex

% BEGIN consequences_criteria_prime_field_descent.tex
\section{The exact geometric multiplicity monoid over a prime field}

This section completes the geometric rational essential-image test
from the prime field in every slope sector.  It is a consequence of
the inspected Milne classification and the transition formula
CC-AG09.  The category is still the rational numerical category
generated by abelian varieties, finite étale schemes and Tate
twists.  No statement about the missing cross-correspondences of
an arbitrary variety is used or implied.

\subsection{CC-AG13: a computable generator in all slopes}

Fix a prime \(p\).  Let \(\xi\in\mathcal W_p\) have weight
\(w\in\mathbb Z\), and choose a stabilized representative
\[
 \xi=[\pi,b],\qquad b\ge1,\qquad
 K=K_\xi=\mathbb Q(\pi^N)\quad\text{for every }N\ge1.
                                                        \tag{CC-AG.80}
\]
Write \(d_\xi\) for the geometric Schur index of its rational
simple object \(S_\xi\).  In particular, its realization has
dimension \(d_\xi[K:\mathbb Q]\).
Define the exponent subgroup directly by
\[
 \mathscr E_\xi=
 \{n\in\mathbb Z:\ \text{there exists }\gamma\in K^\times
                    \text{ with }\gamma^b/\pi^n\in\mu(K)\}.
                                                        \tag{CC-AG.81}
\]
This definition uses no integrality assumption on the normalized
\(p\)-slopes and no fractional ideal with fractional exponents.
Then \(\mathscr E_\xi=h_\xi\mathbb Z\), where \(h_\xi>0\) and
\(h_\xi\mid b\), and
\[
 h_\xi=\min\{n>0:n\mid b,\ 
            X^b-\zeta\pi^n\text{ has a root in }K
                          \text{ for some }\zeta\in\mu(K)\}.
                                                        \tag{CC-AG.82}
\]
The positive integer \(h_\xi\) is independent of the chosen
stabilized representative: equivalently, it is the least positive
exponent \(n\) at which \(\xi\) has a representative in \(K\).

The geometric index divides this exponent:
\[
                              d_\xi\mid h_\xi.
                                                        \tag{CC-AG.83}
\]
The exact multiplicity monoid for base extension from
\(\mathbb F_p\) to \(\overline{\mathbb F}_p\) is
\[
                 H_{\mathbb F_p}(O_\xi)
                      =\frac{h_\xi}{d_\xi}\,\mathbb N.
                                                        \tag{CC-AG.84}
\]
Thus a rational geometric object
\[
                         \bigoplus_\xi S_\xi^{\oplus r_\xi}
\]
with finitely many distinct Galois orbits of labels descends
to \(\mathcal A_p\) if and only if
\[
                        h_\xi/d_\xi\mid r_\xi
                            \quad\text{for every }\xi.
                                                        \tag{CC-AG.85}
\]
Both integers in this test are effectively determined by the
finite equations (CC-AG.82) and the local-invariant formula
of CC-AG08.

If \(K\) has no real embeddings, choose a witness \(\gamma\)
for \(h_\xi\) in (CC-AG.81).  Every root
\[
                             \alpha^{h_\xi}=\gamma
                                                        \tag{CC-AG.86}
\]
is a weight-\(w\) Weil \(p\)-number with germ \(\xi\),
has centre degree
\[
                           [\mathbb Q(\alpha):K]=h_\xi,
                                                        \tag{CC-AG.87}
\]
and has source Schur index one.  Its rational simple object
base changes to \(S_\xi^{\oplus h_\xi/d_\xi}\), realizing the
least positive multiplicity by a single simple object.

\paragraph{Complete proof.}
Multiplication and inversion of witnesses in (CC-AG.81) show
that \(\mathscr E_\xi\) is a subgroup of \(\mathbb Z\).
It contains \(b\), with witness \(\pi\), so it is a nonzero
subgroup \(h_\xi\mathbb Z\), with \(h_\xi>0\) and \(h_\xi\mid b\).
Testing its positive divisors of \(b\) gives (CC-AG.82).

For \(n>0\), a witness satisfies
\(\gamma^b=\zeta\pi^n\).  At every complex embedding it
therefore has absolute value \(p^{nw/2}\); at every finite
place away from \(p\) its valuation is zero.  It is a
weight-\(w\) Weil \(p^n\)-number.  After raising the equation
to the order of \(\zeta\), it gives
\([\gamma,n]=[\pi,b]\).  Conversely, this equality of germs
means that \(\gamma^{bN}=\pi^{nN}\) for some \(N\ge1\).
Consequently \(\gamma^b/\pi^n\) is a root of unity in \(K\).
This proves the intrinsic representative description, and
hence independence in (CC-AG.82).

This is a finite exact test over a fixed number field.
For completeness, one explicit way to determine \(\mu(K)\)
is the bounded-coefficient construction already proved
in CC-AG11--12.  Put \(D=[K:\mathbb Q]\) and
\[
 R_D=\sum_{e=1}^{D}e\prod_{j=1}^{e}
                  \left(2\binom ej+1\right),\qquad
 N_D=\operatorname{lcm}(1,\ldots,R_D).
\]
The minimal polynomial of every power of a root of unity
in \(K\) has degree at most \(D\) and its coefficient of
degree \(e-j\) has absolute value at most \(\binom ej\).
There are at most \(R_D\) possible roots of all these
bounded polynomials.  Hence the order of every root of
unity in \(K\) is at most \(R_D\), and \(\mu(K)\) is precisely
the set of roots in \(K\) of \(X^{N_D}-1\).  Factoring this
polynomial and the finitely many polynomials in
(CC-AG.82) terminates.  This deliberately coarse bound is
used only for finiteness, not for computational efficiency.

First suppose that \(K\) has no real embeddings.
Let \(\gamma\) witness exponent \(h=h_\xi\), and let
\(v\mid p\) be a place of \(K\), with normalized integral
valuation \(\operatorname{ord}_v\), ramification index \(e_v\)
and residue degree \(f_v\).  The geometric local invariant,
computed using the representative \([\gamma,h]\), is
\[
 -\frac{\operatorname{ord}_v(\gamma)}
              {h e_v}[K_v:\mathbb Q_p]
       =-\frac{\operatorname{ord}_v(\gamma)f_v}{h}
                                             \quad\bmod\mathbb Z.
                                                        \tag{CC-AG.88}
\]
Its denominator divides \(h\).  There are no real-place
invariants, and all remaining invariants are zero.
The Schur index is their least common denominator, so
\(d_\xi\mid h\), proving (CC-AG.83) in this case.
This uses the geometric-Frobenius convention in the
corrected Milne manuscript, p.~17 and Theorems~2.15--2.16
and 2.18--2.19, pp.~18 and 20.

Take any source Weil \(p\)-number \(\alpha\) whose germ
orbit is \(O_\xi\).  After conjugation, assume
\([\alpha,1]=[\pi,b]\), and put \(L=\mathbb Q(\alpha)\).
For some \(N\ge1\),
\[
                              \alpha^{bN}=\pi^N.
                                                        \tag{CC-AG.89}
\]
The stabilization in (CC-AG.80) gives \(K\subset L\).
The element \(\theta=\alpha^b/\pi\) lies in \(\mu(L)\).
Set \(d=[L:K]\).  Taking the relative field norm gives
\[
 \frac{N_{L/K}(\alpha)^b}{\pi^d}
                         =N_{L/K}(\theta)\in\mu(K).
                                                        \tag{CC-AG.90}
\]
The norm of a root of unity is a root of unity because
the norm is multiplicative.  No Galois hypothesis on
\(L/K\) is required.  Thus \(d\in\mathscr E_\xi\), and
\[
                                     h\mid d.
                                                        \tag{CC-AG.91}
\]

The field \(L\) has no real embedding, since one would
restrict to a real embedding of \(K\).  At a place
\(z\mid p\) of \(L\), the local invariant of the source
simple over the prime field is
\[
 -\frac{\operatorname{ord}_z(\alpha)}
                 {\operatorname{ord}_z(p)}[L_z:\mathbb Q_p]
               =-\operatorname{ord}_z(\alpha)f_z
                                             \quad\bmod\mathbb Z.
                                                        \tag{CC-AG.92}
\]
It is zero.  The other invariants are also zero, and
there are no real places; hence the source index
\(d_\alpha\) is one.  By CC-AG09 its base-change
multiplicity is
\[
                   [L:K]\frac{d_\alpha}{d_\xi}
                         =\frac{d}{d_\xi}
                    \in\frac{h}{d_\xi}\mathbb N.
                                                        \tag{CC-AG.93}
\]
This proves the necessary divisibility for every
source simple and thus for every positive direct sum.

For attainment, take a root as in (CC-AG.86).
All complex absolute values are \(p^{w/2}\), and all
valuations away from \(p\) vanish, so \(\alpha\in W(p)\).
The exact equation gives
\([\alpha,1]=[\gamma,h]=\xi\).
Every positive power of \(\gamma\) generates \(K\).
Indeed, if \(\gamma^b=\zeta\pi^h\), raising to the
order of \(\zeta\) and then to any positive power
shows that the field of that power of \(\gamma\)
contains a positive power of \(\pi\), which generates
\(K\); the reverse inclusion is immediate.
In particular \(K\subset\mathbb Q(\alpha)\).
The degree of this extension is at most \(h\), by
(CC-AG.86), and is a positive multiple of \(h\), by
(CC-AG.91).  It is exactly \(h\).
Equation (CC-AG.92) gives source index one, so
(CC-AG.93) realizes \(h/d_\xi\).
Direct sums realize all its nonnegative multiples.
This proves (CC-AG.84) when \(K\) has no real embeddings.

It remains to handle every possible real-place case.
If \(K\) has a real embedding \(\sigma\), then
\(\sigma(\pi)^2=p^{bw}\) by the Weil absolute-value
condition.  Injectivity of \(\sigma\) gives
\(\pi^2=p^{bw}\).  Stabilization therefore implies
\[
                             K=\mathbb Q(\pi^2)=\mathbb Q.
                                                        \tag{CC-AG.94}
\]
The germ is uniquely determined by \(w\); signs are
roots of unity and disappear in the germ.
If \(w\) is even, it has the stage-one representative
\(p^{w/2}\in\mathbb Q\), so \(h=1\).  Its real and
\(p\)-adic invariants are integral, giving \(d_\xi=1\).
The corresponding rational one-dimensional Tate
source realizes multiplicity one.

If \(w\) is odd, a rational representative at exponent
\(n\) would have absolute value \(p^{nw/2}\).
This is rational precisely when \(n\) is even.
The representative \(p^w\) at exponent two exists,
including when \(w<0\); hence \(h=2\).
The real invariant \(w/2\) and the \(p\)-adic invariant
\(-w/2\) have denominator two, so \(d_\xi=2\).
Choose instead the equivalent stage-two representative
\(\gamma=-p^w\), and choose
\[
                                  \alpha^2=-p^w.
                                                        \tag{CC-AG.95}
\]
Its field is imaginary quadratic, because \(p^w>0\)
and \(w\) is odd; its degree over \(\mathbb Q\) is two.
It is a weight-\(w\) Weil \(p\)-number with germ \(\xi\).
Equation (CC-AG.92) and absence of real places give
source index one.  The transition multiplicity is
\(2/2=1=h/d_\xi\).
Every nonnegative integer is now realized.  Other
source labels with real centre may have index two;
this does not affect the exact monoid \(\mathbb N\).
Thus (CC-AG.83)--(CC-AG.84) hold also in the real case.

Finally, simple rational objects are classified by
Galois orbits of labels.  Base extension of one simple
has only the one target orbit occurring in CC-AG09.
Semisimple decomposition and positivity therefore
separate distinct target orbits.  Applying the proved
single-orbit criterion to each proves (CC-AG.85).
The zero object has only zero multiplicities and is
included. \(\square\)

\paragraph{Scope relative to the preceding results.}
CC-AG11 gives \(h/\gcd(h,a)\) in the integral-slope sector
over every finite source field.  CC-AG13 gives the entire
all-slope geometric monoid over the prime field, with the
geometric index \(d_\xi\) retained.  Over
\(\mathbb F_{p^a}\) for \(a>1\), the source invariant is
\(-\operatorname{ord}_z(\alpha)f_z/a\); it need not vanish.
No all-slope cyclic formula for those source fields is
asserted here.  CC-AG12 still supplies the finite exact
decision and all preimage isomorphism classes for each
specified target object in that general case.

\subsection{Primary references and exact conventions}

\begin{itemize}
\item Milne, J.~S. (1994). Motives over finite fields. In
U.~Jannsen, S.~Kleiman, \& J.-P.~Serre (Eds.),
\emph{Motives} (Proceedings of Symposia in Pure Mathematics,
Vol.~55, Part~1, pp.~401--459). American Mathematical Society.
Corrected author version dated November~2, 2019:
\url{https://www.jmilne.org/math/xnotes/Seattle1.pdf}.
Geometric-Frobenius invariant convention, p.~17;
Theorems~2.15--2.16, p.~18; Theorems~2.18--2.19,
p.~20.  These inputs are applied only within the
established abelian--Artin--Tate category.
\item Milne, J.~S. (2020, July~19).
\emph{Algebraic number theory} (Version~3.08)
[Author course notes].
\url{https://www.jmilne.org/math/CourseNotes/ANT.pdf}.
Relative field norm, Proposition~2.19,
statement p.~31 and proof pp.~31--32;
bounded-polynomial root-of-unity argument,
Proposition~5.5 and Corollary~5.6, p.~87.
\end{itemize}

% END consequences_criteria_prime_field_descent.tex

% BEGIN consequences_criteria_noncyclic_descent.tex
\section{An exact noncyclic rational descent
monoid}\label{an-exact-noncyclic-rational-descent-monoid}

\textbf{Scope and conventions.} Work in the rational numerical
abelian--Artin--Tate categories \(\mathcal A_q\) already established in
AQ.15 and CC-AG01. The simple-object, rank and Schur-index formulas are
the inspected Milne classification used in CC-AG08--09, with its
geometric-Frobenius local-invariant convention. Only the denominators of
the invariants enter the indices below. The example is a complete
deduction from those inputs and elementary number-field calculations; no
novelty claim is made.

\subsection{CC-AG14 --- Two copies and three copies descend, but one
copy does
not}\label{cc-ag14-two-copies-and-three-copies-descend-but-one-copy-does-not}

Put

\[
 s^2=-23,\qquad K=\mathbf Q(s),\qquad
 \omega=\frac{1+s}{2},\qquad
 \eta=1+\omega=\frac{3+s}{2},
\]

\[
 \gamma=8\bar\eta=12-4s,\qquad
 \xi=[\gamma,9]\in\mathcal W_2,\qquad
 \beta=\gamma^3\in W(2^{27}).
 \tag{CC-AG.96}
\]

Let \(S_\xi\) be the geometric rational simple object of the orbit of
\(\xi\), and let \(T_\beta\in\mathcal A_{2^{27}}\) be the rational
simple object of the orbit of \(\beta\). Both have centre \(K\), Schur
index \(3\), and rank \(6\), and
\((T_\beta)_{\bar{\mathbf F}_2}\simeq S_\xi\).

For descent from the specified field \(\mathbf F_8\), both exact
multiplicity monoids are

\[
 \begin{aligned}
 \{r\ge0:S_\xi^{\oplus r}
          \text{ comes from }\mathcal A_8\}
 &=\langle2,3\rangle_{\mathbf N},\\
 \{r\ge0:T_\beta^{\oplus r}
          \text{ comes from }\mathcal A_8\}
 &=\langle2,3\rangle_{\mathbf N}\\
 &=\{0\}\cup\{r\in\mathbf Z:r\ge2\}.
 \end{aligned}
 \tag{CC-AG.97}
\]

The second base extension is the finite extension
\(\mathbf F_8\subset\mathbf F_{2^{27}}\) of degree \(9\).

Thus these monoids are \textbf{not cyclic}, although their group
completions are \(\mathbf Z\). In particular a gcd or virtual-object
criterion does not decide descent of actual objects. The stabilized
exponent and geometric index are

\[
 h_\xi=9,\qquad d_\xi=3.
 \tag{CC-AG.98}
\]

The general necessary norm bound \((h_\xi/d_\xi)/\gcd(h_\xi/d_\xi,3)=1\)
is not sufficient in this example.

Here are two source simples realizing the generators. Choose
\(\alpha_0^3=\gamma\), let \(\zeta_9\) be a primitive ninth root of
unity, and put \(\alpha_1=\alpha_0\zeta_9\).

\begin{longtable}[]{@{}
  >{\raggedright\arraybackslash}p{(\columnwidth - 8\tabcolsep) * \real{0.1579}}
  >{\raggedleft\arraybackslash}p{(\columnwidth - 8\tabcolsep) * \real{0.2105}}
  >{\raggedleft\arraybackslash}p{(\columnwidth - 8\tabcolsep) * \real{0.2105}}
  >{\raggedleft\arraybackslash}p{(\columnwidth - 8\tabcolsep) * \real{0.2105}}
  >{\raggedleft\arraybackslash}p{(\columnwidth - 8\tabcolsep) * \real{0.2105}}@{}}
\toprule\noalign{}
\begin{minipage}[b]{\linewidth}\raggedright
Source label over \(\mathbf F_8\)
\end{minipage} & \begin{minipage}[b]{\linewidth}\raggedleft
Degree of its centre over \(K\)
\end{minipage} & \begin{minipage}[b]{\linewidth}\raggedleft
Source Schur index
\end{minipage} & \begin{minipage}[b]{\linewidth}\raggedleft
Rank
\end{minipage} & \begin{minipage}[b]{\linewidth}\raggedleft
Multiplicity after either stated base extension
\end{minipage} \\
\midrule\noalign{}
\endhead
\bottomrule\noalign{}
\endlastfoot
\(\alpha_0\) & \(3\) & \(3\) & \(18\) & \(3\) \\
\(\alpha_1\) & \(6\) & \(1\) & \(12\) & \(2\) \\
\end{longtable}

Their irreducible minimal polynomials over \(\mathbf Q\) are
respectively

\[
 P_0(t)=t^6-24t^3+512,
 \qquad
 P_1(t)=t^{12}+24t^9+64t^6+12288t^3+262144.
 \tag{CC-AG.99}
\]

\textbf{Complete proof.}

\subsubsection{1. The quadratic ideal calculation and stabilized
germ}\label{the-quadratic-ideal-calculation-and-stabilized-germ}

First,

\[
 \mathcal O_K=\mathbf Z[\omega],\qquad
 \omega^2-\omega+6=0,\qquad
 N_{K/\mathbf Q}(x+y\omega)=x^2+xy+6y^2.
 \tag{CC-AG.100}
\]

For completeness, if \(a+bs\), \(a,b\in\mathbf Q\), is integral, its
trace and norm are integers. Thus \(2a\in\mathbf Z\), and
\((2a)^2-4N(a+bs)=-92b^2\) is integral. The denominator of \(b\) in
lowest terms is at most \(2\), because its square divides \(92\). Write
\(a=A/2,b=B/2\). Integrality of the norm says
\(A^2+23B^2\equiv0\pmod4\), so \(A\) and \(B\) have the same parity.
This is exactly an element of \(\mathbf Z[\omega]\). Conversely
\(\omega\) satisfies the displayed monic polynomial, so that ring
consists of algebraic integers. The norm formula follows by
multiplication with the conjugate.

Modulo \(2\), this polynomial is \(t(t-1)\). Consequently

\[
 \mathfrak p=(2,\omega-1),\qquad
 \bar{\mathfrak p}=(2,\omega),\qquad
 (2)=\mathfrak p\bar{\mathfrak p},
 \qquad N\mathfrak p=N\bar{\mathfrak p}=2.
 \tag{CC-AG.101}
\]

The primes are distinct, and the two completions of \(K\) above \(2\)
are \(\mathbf Q_2\). One can see the latter directly by lifting the two
simple roots of \(t^2-t+6\) modulo \(2\); the simple-root lifting
statement is Milne (2020), Proposition 7.31, p.~120.

Now \(N(\eta)=8\), while \(\eta\) vanishes modulo \(\mathfrak p\) and is
\(1\) modulo \(\bar{\mathfrak p}\). Unique ideal factorization and the
ideal norm give

\[
 (\eta)=\mathfrak p^3.
 \tag{CC-AG.102}
\]

The ideal \(\mathfrak p\) is not principal. A generator would have
positive norm \(2\), but

\[
 x^2+xy+6y^2=(x+y/2)^2+\frac{23}{4}y^2=2
\]

has no integral solution: if \(y\ne0\), the second term alone exceeds
\(2\); if \(y=0\), it would say \(x^2=2\). Therefore the class of
\(\mathfrak p\) has order exactly \(3\). No class-number table is
needed.

It follows that

\[
 (\gamma)=\mathfrak p^3\bar{\mathfrak p}^{\,6},
 \qquad \gamma\bar\gamma=512=2^9.
 \tag{CC-AG.103}
\]

Hence \(\gamma\) is an integral weight-one Weil \(2^9\)-number. The
ratio \(\gamma/\bar\gamma\) has nonzero valuations at the two primes
above \(2\), and is therefore not a root of unity. Since \(K\) is
quadratic, no positive power of \(\gamma\) is rational: otherwise its
two conjugates would have equal powers. Thus

\[
 \mathbf Q(\gamma^N)=K\qquad(N\ge1),
 \tag{CC-AG.104}
\]

so \(K\) is the stabilized germ centre. The two normalized slopes are
\(1/3\) and \(2/3\). Their local degrees over \(\mathbf Q_2\) are one,
and \(K\) has no real place. Milne's invariant formula therefore gives
local invariants \(-1/3,-2/3\) modulo \(\mathbf Z\), with all others
zero. Their least common denominator is \(3\), so \(d_\xi=3\). The same
calculation applies to \(\beta=\gamma^3\) relative to \(2^{27}\), giving
\(d_\beta=3\). The rank is \(3[K:\mathbf Q]=6\). Their label-to-germ
orbit map is bijective, so CC-AG09 gives
\((T_\beta)_{\bar{\mathbf F}_2}\simeq S_\xi\).

We also compute \(h_\xi\) directly. If \(\delta\in K\) represents
\(\xi\) at exponent \(n>0\), then \(\delta^{9N}=\gamma^{nN}\) for some
\(N\ge1\). Taking normalized valuations gives

\[
 \operatorname{ord}_{\mathfrak p}(\delta)=n/3,\qquad
 \operatorname{ord}_{\bar{\mathfrak p}}(\delta)=2n/3,
 \qquad
 (\delta)=\mathfrak p^{n/3}\bar{\mathfrak p}^{\,2n/3}.
 \tag{CC-AG.105}
\]

First \(3\mid n\). Since \([\bar{\mathfrak p}]=[\mathfrak p]^{-1}\),
principality of this ideal forces \([\mathfrak p]^{-n/3}=1\), and hence
\(9\mid n\). The representative \(\gamma\) at exponent \(9\) exists,
proving \(h_\xi=9\).

\subsubsection{2. Necessary degrees and indices for every possible
geometric
lift}\label{necessary-degrees-and-indices-for-every-possible-geometric-lift}

Consider any source label \(\alpha\in W(8)\) whose germ orbit is that of
\(\xi\). Conjugate it if necessary so that \([\alpha,3]=[\gamma,9]\).
For some \(N\ge1\),

\[
 \alpha^{9N}=\gamma^{3N}.
 \tag{CC-AG.106}
\]

Put \(L=\mathbf Q(\alpha)\). By (CC-AG.104), the right side generates
\(K\), so \(K\subset L\). Write \(c=[L:K]\).

Taking the field norm from \(L\) to \(K\) in (CC-AG.106) shows that
\(N_{L/K}(\alpha)\) is a representative of \(\xi\) at exponent \(3c\).
Equivalently its \(9N\)-th power equals \(\gamma^{3cN}\). By the
computed exponent condition,

\[
 9\mid3c,\qquad\text{so }3\mid c.
 \tag{CC-AG.107}
\]

This norm argument does not assume that \(L/K\) is Galois.

There are no real embeddings of \(L\), since it contains the imaginary
quadratic field \(K\). For each place \(z\mid2\) of \(L\), with
normalized valuation and residue degree \(f_z\), the source local
invariant relative to \(q=8\) is

\[
 -\frac{\operatorname{ord}_z(\alpha)}
          {\operatorname{ord}_z(8)}[L_z:\mathbf Q_2]
 =-\frac{\operatorname{ord}_z(\alpha) f_z}{3}
 \quad\bmod\mathbf Z.
 \tag{CC-AG.108}
\]

All remaining invariants vanish. Thus

\[
 d_\alpha\in\{1,3\},\qquad
 r_\alpha=\frac{c\,d_\alpha}{3},
 \tag{CC-AG.109}
\]

by the exact transition formula CC-AG09. A multiplicity-one source
simple would therefore have

\[
 [L:K]=3,\qquad [L:\mathbf Q]=6,\qquad d_\alpha=1.
 \tag{CC-AG.110}
\]

Moreover (CC-AG.106) gives the root-of-unity statement itself:

\[
 \theta=\alpha^3/\gamma\in L,\qquad
 \theta^{3N}=1.
 \tag{CC-AG.111}
\]

In particular \(\theta\), not only its cube, is a root of unity.

\subsubsection{\texorpdfstring{3. Roots of unity in a degree-six field
containing
\(K\)}{3. Roots of unity in a degree-six field containing K}}\label{roots-of-unity-in-a-degree-six-field-containing-k}

If \(L\) has degree \(6\) over \(\mathbf Q\) and contains
\(K=\mathbf Q(\sqrt{-23})\), then

\[
 \mu(L)=\{1,-1\}.
 \tag{CC-AG.112}
\]

To prove this, let a root of unity in \(L\) have order \(n\). The
cyclotomic degree formula, Milne (2020), Theorem 6.4(a), p.~99, gives
\(\varphi(n)\mid6\). The complete list is

\[
 n\in\{1,2,3,4,6,7,9,14,18\}.
 \tag{CC-AG.113}
\]

Here is a finite justification of completeness. Every prime \(l\mid n\)
has \(l-1\mid\varphi(n)\mid6\), so \(l\in\{2,3,7\}\). Their exponents
are at most \(2,2,1\), respectively, because the factors
\(l^{e-1}(l-1)\) must divide \(6\). Checking these finitely many
combinations gives exactly (CC-AG.113).

Orders \(3,6,9,18\) would put \(\mathbf Q(\sqrt{-3})\) in \(L\), and
order \(4\) would put \(\mathbf Q(i)\) in \(L\). Either is a quadratic
field distinct from \(K\), so its compositum with \(K\) would have
degree \(4\), impossible inside a field of degree \(6\).

For orders \(7,14\), \(L\) would contain \(\mathbf Q(\zeta_7)\), itself
of degree \(6\), and hence equal that field. Its Galois group is cyclic
of order \(6\): it injects into \((\mathbf Z/7\mathbf Z)^\times\), and
the two groups have the same order. Thus it has a unique quadratic
subfield. This subfield is \(\mathbf Q(\sqrt{-7})\), as may be verified
without a table: put \(u=\zeta_7+\zeta_7^2+\zeta_7^4\); then
\(\bar u=\zeta_7^3+\zeta_7^5+\zeta_7^6\), \(u+\bar u=-1\), and
\(u^2=u+2\bar u=-u-2\). Hence \(u^2+u+2=0\), with discriminant \(-7\).
This quadratic field is not \(K\), another contradiction. Only orders
\(1,2\) remain, proving (CC-AG.112).

\subsubsection{4. The local cube obstruction excludes multiplicity
one}\label{the-local-cube-obstruction-excludes-multiplicity-one}

At either completion \(K_v=\mathbf Q_2\), (CC-AG.103) gives
\(\operatorname{ord}_2(\gamma)=3\) or \(6\). Every unit of
\(\mathbf Z_2\) is a cube: its cube-root equation has the simple root
\(1\) modulo \(2\), so Proposition 7.31 applies. Therefore \(\gamma\)
has a cube root \(b_v\in\mathbf Q_2\), with valuation \(1\) or \(2\),
respectively.

Suppose a multiplicity-one lift existed. By (CC-AG.110)--(CC-AG.112),
\(\alpha^3=\gamma\) or \(-\gamma\). Replacing \(\alpha\) by \(-\alpha\)
in the latter case preserves its centre, valuations and index, and gives
\(\alpha^3=\gamma\). Its minimal polynomial over \(K\) has degree \(3\)
and is \(t^3-\gamma\). At either place \(v\mid2\), it factors as

\[
 t^3-\gamma=(t-b_v)(t^2+b_vt+b_v^2).
 \tag{CC-AG.114}
\]

In particular the centre \(L\) has a place above \(v\) of local degree
one, where the valuation of \(\alpha\) is \(1\), respectively \(2\). The
source invariant (CC-AG.108) is then \(-1/3\), respectively \(-2/3\),
and is nonzero. Hence \(d_\alpha=3\), contradicting (CC-AG.110).

This excludes every possible simple lift of multiplicity one, including
arbitrary root-of-unity twists. An actual source object is a positive
sum of source simples, each contributing a positive integer to exactly
one target orbit. It therefore cannot give one copy of \(S_\xi\) unless
such a simple lift exists. We have proved

\[
 1\notin H_{\bar{\mathbf F}_2}(O_\xi).
 \tag{CC-AG.115}
\]

\subsubsection{5. A multiplicity-three
lift}\label{a-multiplicity-three-lift}

Choose \(\alpha_0\) with \(\alpha_0^3=\gamma\). It is integral and all
its conjugates have absolute value \(2^{3/2}\), so it is a weight-one
Weil \(8\)-number. Its germ is \(\xi\). Equation (CC-AG.104) gives
\(K\subset\mathbf Q(\alpha_0)\), while its defining polynomial bounds
the relative degree by \(3\). The norm constraint (CC-AG.107) forces
that degree to be exactly \(3\).

The local argument (CC-AG.114) gives \(d_{\alpha_0}=3\). More
explicitly, after scaling by \(b_v\) its quadratic factor is
\(t^2+t+1\), irreducible modulo \(2\), so the local factors have degrees
\(1\) and \(2\); already the degree-one factor proves the nontrivial
index. Formula (CC-AG.109) yields

\[
 r_{\alpha_0}=3\cdot3/3=3.
 \tag{CC-AG.116}
\]

Its centre degree over \(\mathbf Q\) is \(6\). Taking the norm of
\(t^3-\gamma\) from \(K\) to \(\mathbf Q\), using
\(\gamma+\bar\gamma=24\) and \(\gamma\bar\gamma=512\), gives \(P_0\) in
(CC-AG.99). The degree just proved establishes its irreducibility.

\subsubsection{6. A multiplicity-two
lift}\label{a-multiplicity-two-lift}

Set \(\alpha_1=\alpha_0\zeta_9\), with \(\zeta_9\) primitive. It is
again an integral weight-one Weil \(8\)-number, and

\[
 \alpha_1^9=\gamma^3,\qquad
 \alpha_1^3/\gamma=\zeta_3.
 \tag{CC-AG.117}
\]

The first equation and stabilization give
\(K\subset L_1=\mathbf Q(\alpha_1)\); the second puts \(\zeta_3\) in
\(L_1\). Since \(K\ne\mathbf Q(\sqrt{-3})\), the relative degree
\(c_1=[L_1:K]\) is even. It is divisible by \(3\) by (CC-AG.107). On the
other hand, \(\alpha_1\) is a root of

\[
 Q_1(t)=t^6+\gamma t^3+\gamma^2\in K[t],
 \tag{CC-AG.118}
\]

so \(c_1\le6\). Consequently \(c_1=6\), and \(Q_1\) is its minimal
polynomial over \(K\).

At each \(v\mid2\), use the cube root \(b_v\) above. Then

\[
 Q_1(b_vt)=b_v^6(t^6+t^3+1)=b_v^6\Phi_9(t).
 \tag{CC-AG.119}
\]

The reduction of \(\Phi_9(t)\) modulo \(2\) is irreducible of degree
\(6\). Indeed, its roots in characteristic \(2\) have exact order \(9\):
the identity \((t^3-1)\Phi_9(t)=t^9-1\) has no repeated roots, and
excludes order \(1\) or \(3\). The least positive integer \(e\) with
\(9\mid2^e-1\) is \(6\), as \(2^3\equiv-1\pmod9\). Therefore every root
has degree \(6\) over \(\mathbf F_2\), proving irreducibility.

A monic polynomial with irreducible reduction remains irreducible over
\(\mathbf Q_2\), so (CC-AG.119) proves that \(L_1\) has a single place
of local degree \(6\) over each \(v\mid2\). Its normalized slopes are
still \(1/3,2/3\), because of the germ relation. The two source local
invariants are therefore

\[
 -\tfrac13\cdot6=-2,\qquad
 -\tfrac23\cdot6=-4
 \quad\bmod\mathbf Z,
\]

both zero. There are no real places and all other local invariants
vanish. Thus \(d_{\alpha_1}=1\), and

\[
 r_{\alpha_1}=6\cdot1/3=2.
 \tag{CC-AG.120}
\]

Taking the norm of \(Q_1(t)\) from \(K\) to \(\mathbf Q\) gives
\(P_1(t)\) in (CC-AG.99). Its degree \(12=[L_1:\mathbf Q]\) proves its
irreducibility. The polynomial expansion is also checked by the
accompanying exact rational script; the irreducibility argument is the
field-degree proof, not an unreported numerical computation.

\subsubsection{7. Exact geometric and finite-target
monoids}\label{exact-geometric-and-finite-target-monoids}

We have realized multiplicities \(2\) and \(3\), and excluded \(1\) for
the geometric target. Every integer \(r\ge2\) is a nonnegative
combination of \(2\) and \(3\): use only \(2\)'s for even \(r\), and
\(3\) followed by \(2\)'s for odd \(r\ge3\). Direct sums give all these
multiplicities. This proves the geometric equality in (CC-AG.97).

For the finite target \(T_\beta\), (CC-AG.117) and
\(\alpha_0^9=\gamma^3\) show that both source simples already map to its
orbit under the ninth-power map. Their source-to-target multiplicities
are still \(3\) and \(2\), since the target centre and index are \(K\)
and \(3\). Conversely, a finite-field multiplicity-one lift would remain
one geometric copy after further base extension, contradicting
(CC-AG.115). Hence the same exact monoid holds for the finite target.

There is also a direct finite-orbit verification. The polynomial whose
roots are all possible finite-target source labels is

\[
 \begin{aligned}
 (t^9-\gamma^3)(t^9-\bar\gamma^3)
 &=P_0(t)P_1(t)\\
 &=t^{18}+23040t^9+134217728.
 \end{aligned}
 \tag{CC-AG.121}
\]

The two factors are the irreducible polynomials already proved above, so
there are exactly two source orbits, with multiplicities \(3\) and
\(2\). This proves the finite-target assertion independently of any
search cutoff.

Finally, a cyclic submonoid of \(\mathbf N\) containing both \(2\) and
\(3\) would have generator dividing their gcd, hence generator \(1\),
and would contain \(1\). Thus the monoids in (CC-AG.97) are not cyclic.
Their group completions are \(\mathbf Z\), since \(3-2=1\). This
finishes the proof. \(\square\)

\subsubsection{Classical abelian-variety corollary, independent of the
additional motivic
inputs}\label{classical-abelian-variety-corollary-independent-of-the-additional-motivic-inputs}

There is a geometrically simple abelian threefold \(A/\mathbf F_{2^9}\)
with Frobenius characteristic polynomial

\[
 f_A(t)=(t^2-24t+512)^3,
 \tag{CC-AG.122}
\]

such that, for every integer \(r\ge0\),

\[
 \begin{aligned}
 A_{\bar{\mathbf F}_2}^{r}
  \text{ is isogenous to }B_{\bar{\mathbf F}_2}
  \text{ for some }B/\mathbf F_8
 &\quad\Longleftrightarrow\quad r\in\langle2,3\rangle_{\mathbf N},\\
 A_{\mathbf F_{2^{27}}}^{r}
  \text{ is }\mathbf F_{2^{27}}\text{-isogenous to }
        B_{\mathbf F_{2^{27}}}
  \text{ for some }B/\mathbf F_8
 &\quad\Longleftrightarrow\quad r\in\langle2,3\rangle_{\mathbf N}.
 \end{aligned}
 \tag{CC-AG.123}
\]

The first assertion concerns geometric isogeny descent; the second
concerns isogeny descent to a specified finite target. In particular,
the square and cube of \(A\) have both kinds of descent, while \(A\) has
neither. The zero-dimensional case \(r=0\) means the zero abelian
variety.

\textbf{Proof using only classical Honda--Tate and the explicit
arithmetic above.} Tate's account of Honda--Tate, Theorem 1(i)--(ii),
printed p.~96, classifies simple abelian varieties up to isogeny by
integral Weil numbers, describes their endomorphism division algebras,
and gives \(2\dim B=[\mathbf Q(\pi_B):\mathbf Q]d_B\). Its Remark 2 on
the same page gives the Frobenius characteristic polynomial as the
\(d_B\)-th power of the minimal polynomial. Here Tate uses the
endomorphism algebra of the abelian variety, whose \(p\)-invariant has
the opposite sign to the motivic convention used earlier; the Schur
indices, being least common denominators, are identical. No result about
higher-codimension cycles on abelian varieties is needed in this
corollary.

The integral Weil number \(\gamma\) therefore gives a simple threefold
\(A/\mathbf F_{2^9}\), with (CC-AG.122). For every \(N\ge1\), the number
\(\gamma^N\) has centre \(K\) by (CC-AG.104), and the normalized local
invariants relative to \(2^{9N}\) have the same denominators as those of
\(\gamma\). Its index remains \(3\). The simple abelian variety
associated to \(\gamma^N\) consequently has dimension \(3\), and the
same Frobenius characteristic polynomial as \(A_{\mathbf F_{2^{9N}}}\).
Tate (1966), Theorem 1(c), printed p.~139, makes these varieties
isogenous. Thus \(A\) stays simple over every finite extension. Any
abelian subvariety over \(\bar{\mathbf F}_2\) is defined over a finite
extension, since its defining equations and group maps use only finitely
many algebraic coefficients. Therefore \(A\) is geometrically simple.

Let \(B_9,B_6\) be the simple abelian varieties over \(\mathbf F_8\)
associated to \(\alpha_0,\alpha_1\), respectively. The subscripts
indicate their dimensions. The centre degrees and indices established in
Sections 5--6 give

\[
 \dim B_9=9,\qquad \dim B_6=6,\qquad
 f_{B_9}=P_0^3,\qquad f_{B_6}=P_1.
 \tag{CC-AG.124}
\]

Under extension from \(\mathbf F_8\) to \(\mathbf F_{2^{27}}\),
Frobenius becomes its ninth power. Since
\(\alpha_0^9=\alpha_1^9=\beta=\gamma^3\), all the resulting eigenvalues
lie in the two-element conjugacy orbit of \(\beta\). The two eigenvalues
occur with the same multiplicity: they are Galois conjugate and the
characteristic polynomial is rational. The dimensions therefore give
characteristic polynomials \(m_\beta(t)^9\) and \(m_\beta(t)^6\),
respectively, where \(m_\beta(t)=(t-\beta)(t-\bar\beta)\). The variety
\(A_{\mathbf F_{2^{27}}}\) has characteristic polynomial
\(m_\beta(t)^3\). Applying Tate's isogeny criterion gives actual
finite-field isogenies

\[
 (B_9)_{\mathbf F_{2^{27}}}
       \sim_{\mathbf F_{2^{27}}} A_{\mathbf F_{2^{27}}}^{3},
 \qquad
 (B_6)_{\mathbf F_{2^{27}}}
       \sim_{\mathbf F_{2^{27}}} A_{\mathbf F_{2^{27}}}^{2}.
 \tag{CC-AG.125}
\]

Products of these source varieties supply every \(r\ge2\) in both
assertions of (CC-AG.123).

For completeness, the impossibility of a single geometric copy can also
be read entirely in the classical isogeny category. If \(C/\mathbf F_8\)
were geometrically isogenous to \(A\), then \(C\) would be a
geometrically simple threefold, hence simple over \(\mathbf F_8\). Let
\(\alpha\) be its Weil label and \(L=\mathbf Q(\alpha)\). An isogeny
over the algebraic closure is defined over some finite common extension,
so, after conjugating \(\alpha\), its germ is \(\xi\). Sections 2--4
then apply to its centre and classical Schur index: \(K\subset L\),
\(3\mid c=[L:K]\), and \(d_\alpha\in\{1,3\}\). The classical dimension
formula and \(\dim C=3\) require \(c d_\alpha=3\), hence
\(c=3,d_\alpha=1\). But the complete root-of-unity and local cube
argument in Sections 3--4 forces \(d_\alpha=3\), a contradiction. This
argument uses only the Honda--Tate index formula, not the
numerical-motive transition theorem. It also excludes a finite-target
isogeny, which would imply a geometric one. This proves (CC-AG.123).
\(\square\)

\subsubsection{Consequence for the remaining descent
criterion}\label{consequence-for-the-remaining-descent-criterion}

CC-AG11's integral-slope cyclic formula and CC-AG13's prime-field cyclic
formula retain their stated hypotheses. A universal cyclic formula for
an arbitrary specified field \(\mathbf F_{p^a}\) is false, even for a
weight-one orbit with imaginary quadratic stabilized centre. The
positive multiplicity equations of CC-AG09 and the finite fixed-object
procedure of CC-AG12 cannot in general be replaced by one divisibility
condition. This is an exact obstruction, not a claim that those existing
finite decision procedures fail.

\subsection{Primary references and exact
locators}

\begin{itemize}
\item
  Honda, T. (1968). Isogeny classes of abelian varieties over finite
  fields. \emph{Journal of the Mathematical Society of Japan, 20}(1--2),
  83--95. \url{https://doi.org/10.2969/jmsj/02010083} . Exact locator: Main
  Theorem, p.~83. The occurrence statement and its integral Weil-number
  definition were previously inspected in AQ-WEIL-EXIST and CC-AG01.
  Primary scan:
  \url{https://people.math.rochester.edu/faculty/doug/otherpapers/Honda-Ab.pdf}
  .
\item
  Milne, J. S. (1994). Motives over finite fields. In U. Jannsen, S.
  Kleiman, \& J.-P. Serre (Eds.), \emph{Motives} (Proceedings of
  Symposia in Pure Mathematics, Vol. 55, Part 1, pp.~401--459). American
  Mathematical Society. \textbf{Version used:} corrected author
  manuscript, November 2, 2019.
  \url{https://www.jmilne.org/math/xnotes/Seattle1.pdf} . Exact locators:
  invariant convention p.~17; Theorem 2.15(c) and Theorem 2.16, p.~18;
  Theorems 2.18--2.19, p.~20. These are the checked classification
  inputs, applied in the abelian-generated category as in CC-AG08.
\item
  Milne, J. S. (2020, July 19). \emph{Algebraic number theory} (Version
  3.08) {[}Author course notes{]}.
  \url{https://www.jmilne.org/math/CourseNotes/ANT.pdf} . Exact locators:
  field norm, Proposition 2.19, pp.~31--32; ideal factorization, Theorem
  3.20, p.~53; ideal norms, Propositions 4.1--4.2, p.~69; cyclotomic
  degree, Theorem 6.4(a), p.~99; simple-root lifting over a complete
  discrete valuation ring, Proposition 7.31 and its proof, p.~120. The
  class order, root-of-unity exclusion, local factors, indices and exact
  monoids in this example are proved explicitly above.
\item
  Tate, J. (1966). Endomorphisms of abelian varieties over finite
  fields. \emph{Inventiones Mathematicae, 2}, 134--144.
  \url{https://doi.org/10.1007/BF01404549} . Exact locator: Theorem 1(c),
  p.~139, equivalence of isogeny and equality of Frobenius
  characteristic polynomials. Primary scan directly inspected:
  \url{https://pazuki.perso.math.cnrs.fr/index_fichiers/Tate66.pdf} .
\item
  Tate, J. (1971). Classes d'isogénie des variétés abéliennes sur un
  corps fini (d'après T. Honda). In \emph{Séminaire Bourbaki: Vol.
  1968/69, exposés 347--363} (Exposé 352, pp.~95--110). Springer-Verlag.
  \url{https://www.numdam.org/book-part/SB_1968-1969__11__95_0/} . Exact
  locator: Theorem 1(i)--(ii), p.~96, equations (1)--(2), and Remark 2
  on p.~96; talk delivered November 1968, published 1971. Primary scan
  directly inspected:
  \url{https://www.numdam.org/article/SB_1968-1969__11__95_0.pdf} .
\end{itemize}

% END consequences_criteria_noncyclic_descent.tex

% BEGIN consequences_criteria_local_descent_obstruction.tex
\section{A local obstruction to the gcd of all descent
multiplicities}\label{a-local-obstruction-to-the-gcd-of-all-descent-multiplicities}

\subsection{CC-AG15 --- Identical exponent and index data can give a
different
gcd}\label{cc-ag15-identical-exponent-and-index-data-can-give-a-different-gcd}

The finite fixed-object decision procedure CC-AG12 remains valid. This
section rules out a proposed shortcut for the \emph{whole} geometric
multiplicity monoid. It gives an exact local obstruction which survives
every root-of-unity twist, with no bound on the degree of a possible
lift. The motivic formulation uses only the previously specified
classification within the abelian--Artin--Tate sector. The final
abelian-variety formulation uses classical Honda--Tate alone.

\subsubsection{A finite local criterion}\label{a-finite-local-criterion}

Let a geometric simple label have stabilized representative

\[
 \xi=[\gamma,9],\qquad
 K=\mathbf Q(\gamma^N)\quad(N\ge1).
 \tag{CC-AG.126}
\]

Assume that \(K\) has no real embeddings, \(3\nmid[K:\mathbf Q]\), and
its minimal representative exponent and geometric Schur index are

\[
 h_\xi=9,\qquad d_\xi=3.
 \tag{CC-AG.127}
\]

Suppose there is a place \(v\mid p\) of \(K\) with all three properties:

\begin{enumerate}
\def\labelenumi{\arabic{enumi}.}
\tightlist
\item
  The geometric division-algebra invariant at \(v\) has exact order
  \(3\) in \(\mathbf Q/\mathbf Z\).
\item
  The element \(\gamma\) is a cube in \(K_v\).
\item
  The field \(K_v\) contains all twenty-seven-th roots of unity.
\end{enumerate}

Then the exact geometric multiplicity monoid for descent from
\(\mathbf F_{p^3}\) is

\[
 H_{\mathbf F_{p^3}}(O_\xi)=3\mathbf N.
 \tag{CC-AG.128}
\]

For the simple finite target with label \(\beta=\gamma^3\) over
\(\mathbf F_{p^{27}}\), the exact monoid for base extension from
\(\mathbf F_{p^3}\) is also \(3\mathbf N\).

\textbf{Complete proof of the criterion.} Let \(\alpha\in W(p^3)\) have
germ \(\xi\), after Galois conjugation if needed. For some \(N\ge1\),

\[
 \alpha^{9N}=\gamma^{3N}.
 \tag{CC-AG.129}
\]

Stabilization gives \(K\subset L=\mathbf Q(\alpha)\). Write \(c=[L:K]\).
Taking the relative norm in (CC-AG.129) shows that \(N_{L/K}(\alpha)\)
represents the germ at exponent \(3c\). By \(h_\xi=9\),

\[
 3\mid c.
 \tag{CC-AG.130}
\]

The field \(L\) has no real embeddings. At any place \(z\mid p\), its
finite-source invariant is

\[
 -\frac{\operatorname{ord}_z(\alpha)}
           {\operatorname{ord}_z(p^3)}[L_z:\mathbf Q_p]
 =-\frac{\operatorname{ord}_z(\alpha)f_z}{3}
                                      \quad\bmod\mathbf Z,
 \tag{CC-AG.131}
\]

where the valuation is normalized to have value group \(\mathbf Z\) and
\(f_z\) is the residue degree over \(\mathbf F_p\). All other invariants
vanish. Thus its source index \(d_\alpha\) is either \(1\) or \(3\), and
the exact geometric multiplicity is

\[
 r_\alpha=\frac{c d_\alpha}{3}.
 \tag{CC-AG.132}
\]

Suppose \(3\nmid r_\alpha\). Equations (CC-AG.130)--(CC-AG.132) force

\[
 v_3(c)=1,\qquad d_\alpha=1.
 \tag{CC-AG.133}
\]

The element \(\theta=\alpha^3/\gamma\in L\) is a root of unity, since
\(\theta^{3N}=1\). Factor it uniquely as \(\theta=\theta_3\theta_{3'}\),
where the two factors have respectively power-of-three and
prime-to-three orders. Both factors are integral powers of \(\theta\),
so belong to \(L\).

Write the order of \(\theta_3\) as \(3^m\). If \(m\ge1\), the cyclotomic
degree formula gives \(2\cdot3^{m-1}\mid[L:\mathbf Q]\). Since
\(3\nmid[K:\mathbf Q]\), (CC-AG.133) implies \(m\le2\). Thus

\[
 \theta_3^9=1.
 \tag{CC-AG.134}
\]

If \(n\) is the order of \(\theta_{3'}\), choose an integer \(u\) with
\(3u\equiv1\pmod n\), and put \(\rho=\theta_{3'}^u\in L\). Then
\(\rho^3=\theta_{3'}\). Define

\[
 \alpha'=\alpha/\rho,\qquad L'=\mathbf Q(\alpha').
\]

It satisfies

\[
 (\alpha')^3=\gamma\theta_3,
 \qquad (\alpha')^{27}=\gamma^9.
 \tag{CC-AG.135}
\]

Consequently \(K\subset L'\subset L\), by stabilization, and \(\alpha'\)
has the same Weil conditions and the same germ as \(\alpha\):
multiplication by a root of unity changes neither. The norm argument
(CC-AG.130) therefore gives \(3\mid[L':K]\). Because \([L':K]\) divides
\(c\), it follows that

\[
 v_3([L':K])=1,\qquad 3\nmid[L:L'].
 \tag{CC-AG.136}
\]

Now use the specified place \(v\). Choose \(b\in K_v\) with
\(b^3=\gamma\). The roots of \(X^{27}-\gamma^9\) are exactly \(b\zeta\),
with \(\zeta^{27}=1\), and all lie in \(K_v\) by assumption. They are
distinct, since the characteristic is zero. The minimal polynomial of
\(\alpha'\) over \(K\) divides this polynomial, so it splits into
distinct linear factors over \(K_v\). Hence

\[
 L'\otimes_K K_v\simeq K_v^{[L':K]}.
 \tag{CC-AG.137}
\]

This is the completion-factor correspondence of Milne (2020),
Propositions 8.1--8.2, p.~136; it also follows directly by applying the
Chinese remainder theorem to the split minimal polynomial.

Choose one resulting place \(w\) of \(L'\), so \(L'_w=K_v\). The same
completion decomposition for \(L/L'\), and equality of vector-space
dimensions, give

\[
 \sum_{z\mid w}[L_z:L'_w]=[L:L'].
 \tag{CC-AG.138}
\]

The right side is prime to \(3\). Therefore at least one summand is
prime to \(3\). For that \(z\), the source invariant equals
\([L_z:K_v]\) times the geometric invariant at \(v\). Indeed the germ
equation makes their normalized slopes equal, and the local degree
multiplies in the tower. Its order is consequently still \(3\). This
contradicts \(d_\alpha=1\) in (CC-AG.133).

We have proved that every source simple contributes a multiple of \(3\).
Positive direct sums do too, since the base-change decomposition has one
geometric orbit for each source simple and no cancellation. Thus the
geometric monoid is contained in \(3\mathbf N\).

For attainment, take any root \(\alpha_0^3=\gamma\). It is a Weil
\(p^3\)-number of the correct weight and germ. Stabilization gives
\(K\subset\mathbf Q(\alpha_0)\); its relative degree is at most \(3\),
while (CC-AG.130) makes it divisible by \(3\). It is exactly \(3\). At
\(v\), the polynomial \(X^3-\gamma\) splits completely, using the local
cube root and \(\mu_3\subset K_v\). Thus there is a local degree-one
place above \(v\), at which the source invariant has order \(3\).
Equation (CC-AG.131) bounds the index by \(3\), so it is exactly \(3\).
Its multiplicity is \(3\cdot3/3=3\). Direct sums give every nonnegative
multiple of \(3\), proving (CC-AG.128).

For the finite target, stabilization and the unchanged normalized slopes
give centre \(K\), index \(3\), and geometric base change equal to the
one simple object \(S_\xi\). A finite-target source multiplicity must
therefore belong to the geometric monoid. Conversely
\(\alpha_0^9=\gamma^3=\beta\), so the source just constructed already
gives three copies of this finite target. Direct sums prove the asserted
finite-target equality. This completes the proof of the criterion.
\(\square\)

\subsubsection{A completely explicit
instance}\label{a-completely-explicit-instance}

Take

\[
 p=163,\qquad s^2=-23,\qquad K=\mathbf Q(s),\qquad
 \omega=(1+s)/2,
\]

\[
 \eta=818+399s=419+798\omega,
 \qquad \gamma=p^3\bar\eta,
 \qquad \xi=[\gamma,9].
 \tag{CC-AG.139}
\]

Then every hypothesis of the criterion holds. In particular

\[
 h_\xi=9,\qquad d_\xi=3,\qquad
 H_{\mathbf F_{163^3}}(O_\xi)=3\mathbf N.
 \tag{CC-AG.140}
\]

\textbf{Complete verification.} The integral basis and norm formula for
\(K\) were proved in CC-AG14: \(\mathcal O_K=\mathbf Z[\omega]\) and
\(N(x+y\omega)=x^2+xy+6y^2\). Direct calculation gives

\[
 N(\eta)=818^2+23\cdot399^2=4330747=163^3,
 \qquad \operatorname{Tr}(\eta)=1636\equiv6\pmod{163}.
 \tag{CC-AG.141}
\]

The two roots of \(X^2+23\) modulo \(163\) are \(125\) and \(38\). At
\(s=125\), \(\eta\) vanishes; at \(s=38\), it equals \(6\). Thus \(p\)
splits as \(\mathfrak p\bar{\mathfrak p}\), with
\(\mathfrak p=(p,s-125)\), and

\[
 (\eta)=\mathfrak p^3,
 \qquad (\gamma)=\mathfrak p^3\bar{\mathfrak p}^{\,6},
 \qquad \gamma\bar\gamma=p^9.
 \tag{CC-AG.142}
\]

The two roots are simple, so their Hensel lifts also identify both
completions with \(\mathbf Q_{163}\).

The ideal \(\mathfrak p\) is not principal. A generator would have norm
\(163\). Writing its coordinates in the integral basis would give

\[
 (2x+y)^2=652-23y^2.
\]

It suffices to test \(|y|\le5\), since larger absolute values give a
negative right side. For \(|y|=0,1,2,3,4,5\), the values are

\[
 652,\quad629,\quad560,\quad445,\quad284,\quad77,
\]

none of which is a square. The class of \(\mathfrak p\) therefore has
exact order \(3\), by (CC-AG.142).

Exactly as can also be checked directly here, the valuations of
\(\gamma/\bar\gamma\) are nonzero. No positive power of \(\gamma\) is
rational, so \(K\) is its stabilized centre. A representative in \(K\)
at exponent \(n\) would have ideal
\(\mathfrak p^{n/3}\bar{\mathfrak p}^{2n/3}\). Integrality of its
valuations requires \(3\mid n\), and the exact order-three class then
requires \(9\mid n\). The displayed \(\gamma\) realizes exponent \(9\),
so \(h_\xi=9\).

The slopes at the two \(p\)-adic places are \(1/3\) and \(2/3\), with
local degree one. The geometric invariants have exact order \(3\) at
both, and all others vanish, since \(K\) is imaginary quadratic. Thus
\(d_\xi=3\). Also \(3\nmid[K:\mathbf Q]=2\).

It remains to check the local cube and roots-of-unity conditions, not
merely the ideal valuations. At \(\mathfrak p\), the unit
\(\gamma/p^3=\bar\eta\) has residue \(6\). At \(\bar{\mathfrak p}\), the
unit \(\gamma/p^6=1/\eta\) has residue \(1/6\). The explicit congruence

\[
 35^3\equiv6\pmod{163}
 \tag{CC-AG.143}
\]

and simple-root Hensel lifting show that these units are cubes in their
respective \(\mathbf Q_{163}\) completions. Since their extracted powers
of \(p\) have exponents \(3\) and \(6\), \(\gamma\) is itself a cube at
both places.

Finally,

\[
 6^{27}\equiv1\pmod{163},\qquad
 6^9\equiv58\not\equiv1\pmod{163}.
 \tag{CC-AG.144}
\]

Thus \(6\) has exact order \(27\) modulo \(163\). The root of
\(X^{27}-1\) at \(6\) is simple. Its Hensel lift in \(\mathbf Q_{163}\)
still has order exactly \(27\), since its reduction does. Its powers
give all twenty-seven-th roots of unity. The criterion applies, proving
(CC-AG.140) and its finite-target counterpart. All numerical identities
and the finite nonsquare test are also checked by the accompanying exact
integer script. \(\square\)

\subsubsection{What this rules out, and the classical threefold
version}\label{what-this-rules-out-and-the-classical-threefold-version}

Put \(t=h_\xi/d_\xi\). The norm argument always supplies the necessary
divisibility \(t/\gcd(t,a)\mid r\) for source exponent \(a\). In this
example,

\[
 t=3,\qquad a=3,\qquad
 \frac{t}{\gcd(t,a)}=1,\qquad
 \gcd H_{\mathbf F_{163^3}}(O_\xi)=3.
 \tag{CC-AG.145}
\]

Consequently the proposed equality between that gcd and the necessary
norm divisor is false. Moreover CC-AG14 has the same values
\(h=9,d=3,a=3\), and the same stabilized quadratic field \(K\), but its
monoid has gcd \(1\). These integer and centre data alone do not
determine the answer: the local root-of-unity and root conditions
matter. This does not invalidate the finite test CC-AG12, whose full
local-index computation retains those conditions.

The element \(\gamma\) in (CC-AG.139) is an integral weight-one Weil
\(163^9\)-number. Classical Honda--Tate therefore gives a simple abelian
threefold \(A/\mathbf F_{163^9}\), with Frobenius characteristic
polynomial

\[
 \bigl(X^2-1636\cdot163^3 X+163^9\bigr)^3.
 \tag{CC-AG.146}
\]

It is geometrically simple: every positive power of \(\gamma\) still has
centre \(K\) and index \(3\), so the argument using Tate's isogeny
criterion in the classical corollary to CC-AG14 applies at every finite
extension. Let \(I_{\rm geom}\) be the set of integers \(r\ge0\) for
which \(A_{\bar{\mathbf F}_{163}}^r\) is isogenous to the base extension
of an abelian variety over \(\mathbf F_{163^3}\). Let \(I_{27}\) be the
analogous set requiring an isogeny already over \(\mathbf F_{163^{27}}\)
to \(A_{\mathbf F_{163^{27}}}^r\). Then

\[
 I_{\rm geom}=I_{27}=3\mathbf N.
 \tag{CC-AG.147}
\]

Here \(r\ge0\), and the zero case has its usual meaning. To see that no
motivic assumption is used in (CC-AG.147), decompose any source abelian
variety up to isogeny into simple factors. A factor with this geometric
isogeny type has Weil label \(\alpha\) and dimension
\([L:\mathbf Q]d_\alpha/2\), hence contributes \(c d_\alpha/3\) copies
of the geometric threefold. The complete arbitrary-twist argument above
proves that every such number is divisible by \(3\), using only
classical local indices. Conversely, \(\alpha_0^3=\gamma\) gives a
simple nine-dimensional source variety whose ninth Frobenius power has
the characteristic polynomial of \(A_{\mathbf F_{163^{27}}}^{3}\).
Tate's finite-field isogeny criterion realizes the required isogeny, and
products give all multiples. Thus both assertions are consequences of
classical Honda--Tate and the explicit arithmetic, independent of any
higher-codimension Tate or all-CM Hodge input.

\subsection{Primary references and exact
locators}

\begin{itemize}
\tightlist
\item
  Milne, J. S. (1994). Motives over finite fields. In U. Jannsen, S.
  Kleiman, \& J.-P. Serre (Eds.), \emph{Motives} (Proceedings of
  Symposia in Pure Mathematics, Vol. 55, Part 1, pp.~401--459). American
  Mathematical Society. Version used: corrected author manuscript,
  November 2, 2019. \url{https://www.jmilne.org/math/xnotes/Seattle1.pdf} .
  Exact locators: geometric-Frobenius convention p.~17; Theorems
  2.15--2.16, p.~18; Theorems 2.18--2.19, p.~20. Reused only on the
  established abelian-generated category.
\item
  Milne, J. S. (2020, July 19). \emph{Algebraic number theory} (Version
  3.08) {[}Author course notes{]}.
  \url{https://www.jmilne.org/math/CourseNotes/ANT.pdf} . Exact locators:
  relative field norm, Proposition 2.19, pp.~31--32; cyclotomic degree,
  Theorem 6.4(a), p.~99; simple-root Hensel lifting, Proposition 7.31,
  p.~120; completions and factorization, Propositions 8.1--8.2 and the
  proof of 8.2, p.~136.
\item
  Tate, J. (1966). Endomorphisms of abelian varieties over finite
  fields. \emph{Inventiones Mathematicae, 2}, 134--144.
  \url{https://doi.org/10.1007/BF01404549} . Exact locator: Theorem 1(c),
  p.~139. Primary scan:
  \url{https://pazuki.perso.math.cnrs.fr/index_fichiers/Tate66.pdf} .
\item
  Tate, J. (1971). Classes d'isogénie des variétés abéliennes sur un
  corps fini (d'après T. Honda). In \emph{Séminaire Bourbaki: Vol.
  1968/69, exposés 347--363} (Exposé 352, pp.~95--110). Springer-Verlag.
  \url{https://www.numdam.org/article/SB_1968-1969__11__95_0.pdf} .
  Exact locator: Theorem 1(i)--(ii), equations (1)--(2), and Remark 2,
  p.~96. The classical endomorphism-algebra invariant has the opposite
  sign to the stated motivic convention, with identical index
  denominators.
\end{itemize}

% END consequences_criteria_local_descent_obstruction.tex

% BEGIN root_descent_prime_support.tex
\section{Only finitely many cyclotomic prime towers are needed for geometric descent}

\paragraph{BR-15-PRIME-SUPPORT.}
Work in the rational numerical abelian--Artin--Tate categories and with
the exact transition formulas of CC-AG08--09. Let $\xi$ be a geometric
Weil germ of weight $w$, let $K$ be its stabilized centre, and let
$d_\xi$ be its geometric Schur index. Choose a representative
\[
                  \xi=[\gamma,h],\qquad \gamma\in K^\times,
\]
at the least positive exponent $h=h_\xi$ of CC-AG13. Thus
$\mathbb Q(\gamma^n)=K$ for every $n\ge1$. Fix the source field
$\mathbb F_{p^a}$, and let $H_a(\xi)$ be the additive monoid of
multiplicities of the geometric simple $S_\xi$ obtained by base change
from that field. For a source Weil label $\alpha$ with
$[\alpha,a]=\xi$, put
\[
 L_\alpha=\mathbb Q(\alpha),\qquad
 r_\alpha=[L_\alpha:K]\,d_\alpha/d_\xi.
                                                               \tag{PS.1}
\]
The positive integer $r_\alpha$ is the transition multiplicity already
proved in CC-AG09.

Then $H_a(\xi)$ is generated by the $r_\alpha$ for which
\[
       \alpha^h=\gamma^a\theta,
       \qquad \theta\in\mu_{h\text{-primary}}.
                                                               \tag{PS.2}
\]
Here the subscript means that the order has no prime divisor outside
those of $h$. In fact every source label can be replaced by a label
satisfying (PS.2), whose centre is a subfield of the original centre
and whose transition multiplicity
\emph{divides} the original multiplicity. No finite bound on the
exponents of the primes dividing $h$ is asserted.

The assertion concerns a geometric target. Removing a root-of-unity
factor need not preserve a prescribed finite-field target label, so
no such finite-target assertion is included.

\begin{proof}
Take any source simple contributing to the orbit of $\xi$, and conjugate
its label so that $[\alpha,a]=[\gamma,h]$. There is an $N\ge1$ with
\[
                      \alpha^{hN}=\gamma^{aN}.
                                                               \tag{PS.3}
\]
Consequently $K\subset L=\mathbb Q(\alpha)$, because the right side
generates $K$, and
\[
                         \theta=\alpha^h/\gamma^a\in\mu(L).
\]
Write uniquely $\theta=\theta_h\theta_0$, where the order of
$\theta_h$ has only prime divisors dividing $h$, and the order $m$ of
$\theta_0$ is coprime to $h$. Both factors are powers of $\theta$ and
therefore belong to $L$. Choose an integer $e$ with $he\equiv1\pmod m$
(with $\theta_0=1$ if $m=1$), and put
\[
                \zeta=\theta_0^e,\qquad \alpha'=\alpha/\zeta.
                                                               \tag{PS.4}
\]
Then $\zeta^h=\theta_0$, so $(\alpha')^h=\gamma^a\theta_h$.
Multiplication by a root of unity preserves the Weil absolute values,
weight, and the condition of being a unit away from $p$. Thus $\alpha'$
is a valid label over the same source field, and its geometric germ
is still $\xi$. Let $L'=\mathbb Q(\alpha')$. We have $L'\subset L$,
and some positive power of $(\alpha')^h$ is a positive power of
$\gamma^a$; stabilization gives $K\subset L'$.

Let $d'$ and $d$ be the source Schur indices attached to $\alpha'$
and $\alpha$, respectively. We prove directly from the explicit local
invariants that
\[
                              d'\mid[L:L']d.
                                                               \tag{PS.5}
\]
This avoids replacing actual multiplicities by a virtual-object gcd.
The classification input is precisely the formula in CC-AG08: at a
finite place above $p$ the invariant is the negative normalized slope
times the local degree over $\mathbb Q_p$; at a real place it is $w/2$
modulo $\mathbb Z$; at the remaining places it is zero. The Schur index
is the least common denominator of these finitely many invariants.
Write $b'_v$ and $b_z$ for the two sets of invariants.

If $z$ lies over a $p$-adic place $v$ of $L'$, the quotient
$\alpha/\alpha'$ is a root of unity and has valuation zero. Thus
\[
 \begin{split}
 b_z
   &=-\frac{\operatorname{ord}_v(\alpha')}{
                    \operatorname{ord}_v(p^a)}[L_z:\mathbb Q_p]\\
   &=[L_z:L'_v]b'_v\pmod{\mathbb Z}.
 \end{split}                                                   \tag{PS.6}
\]
At the other finite places the same congruence holds with both sides
zero. At an archimedean place over a complex place it also holds with
both sides zero. For a real place extending a real place the local
degree is one and the weights agree. For a complex place over a real
place the local degree is two, which sends $w/2$ to the integer $w$.
Thus the congruence holds at every place, including the real-centre case.

By the definition of $d$, each $d b_z$ is zero in $\mathbb Q/\mathbb Z$.
Sum (PS.6), multiplied by $d$, over the places $z$ above a fixed $v$.
The local degrees sum to the global degree $[L:L']$: this follows by
tensoring the $L'$-vector space $L$ with its completion and using the
product decomposition into the completions of $L$. Consequently
\[
                              [L:L']d b'_v=0
                           \quad\text{in }\mathbb Q/\mathbb Z.
                                                               \tag{PS.7}
\]
It holds for every $v$. Taking the least common denominator of the
$b'_v$ gives (PS.5). The completion decomposition is Milne (2020),
\emph{Algebraic number theory}, Propositions~8.1--8.2, p.~136;
its real and complex version is the same scalar-extension dimension
identity. No Galois assumption on either field extension is needed.

Applying (PS.1) to the two labels now gives
\[
                 \frac{r_\alpha}{r_{\alpha'}}
                          =\frac{[L:L']d}{d'}\in\mathbb Z_{>0}.
                                                               \tag{PS.8}
\]
Thus every original generator of $H_a(\xi)$ is a positive integer
multiple of one with the restricted root-of-unity support (PS.2).
Conversely every restricted label is an allowed source label, so its
multiplicity already belongs to $H_a(\xi)$. Taking finite positive sums
proves equality of the two monoids.
\end{proof}

\paragraph{Exact remaining gate and source scope.}
The unrestricted geometric search over all root-of-unity orders reduces
to towers at the finitely many primes dividing the computable integer
$h_\xi$. The new statement does not bound the height in those towers,
compute a whole-monoid conductor, or assert that the monoid is cyclic.
CC-AG14 already disproves the last assertion. CC-AG12's separate
rank-dependent finite procedure remains valid for every specified object.
The division-algebra and transition inputs are Milne (1994), corrected
author manuscript of 2 November 2019, pp.~17--20, particularly
Theorems~2.15(c), 2.16 and 2.18--2.19, in the established abelian sector;
the proof above supplies the root-of-unity reduction and actual positive
multiplicity divisibility in full.

% END root_descent_prime_support.tex

% BEGIN consequences_criteria_uniform_descent.tex
\section{A uniform finite computation of each geometric descent
monoid}\label{a-uniform-finite-computation-of-each-geometric-descent-monoid}

\subsection{CC-AG16 --- A conductor bound independent of the target
multiplicity}\label{cc-ag16-a-conductor-bound-independent-of-the-target-multiplicity}

Work in the established rational numerical abelian--Artin--Tate
categories, with the exact simple-object and local-index formulas of
CC-AG08--09. The theorem below computes the entire geometric
multiplicity monoid over a specified finite source field. It retains the
full positive multiplicity data; the noncyclic monoid of CC-AG14 and the
stronger local obstruction of CC-AG15 are included.

Choose a stabilized representative at the least positive exponent:

\[
 \xi=[\gamma,h],\qquad
 K=\mathbf Q(\gamma^n)\quad(n\ge1),\qquad h=h_\xi.
 \tag{CC-AG.148}
\]

Let \(d_\xi\) be the geometric Schur index, and fix the source field
\(\mathbf F_{p^a}\), with \(a\ge1\). Write \(H_a(\xi)\) for the additive
monoid of multiplicities of the geometric simple \(S_\xi\) obtained by
base extension from that source.

If \(K\) has a real embedding, then \(K=\mathbf Q\) and
\(H_a(\xi)=\mathbf N\), as follows from the real-centre cases of CC-AG13
and subsequent base extension. Suppose henceforth that \(K\) has no real
embeddings.

Let \(E/\mathbf Q\) be the splitting field of

\[
 \operatorname{Norm}_{K/\mathbf Q}(X^h-\gamma^a),
 \qquad M=[E:\mathbf Q],\qquad S=\{\ell:\ell\text{ prime},\ \ell\mid h\}.
 \tag{CC-AG.149}
\]

For \(\ell\in S\), put

\[
 \varepsilon_\ell=
 \begin{cases}1&\ell\ne2,\\2&\ell=2,\end{cases}
 \qquad
 b_\ell=\max\bigl(\{v_\ell(r-1):r\in S\setminus\{p,\ell\}\}\cup\{0\}\bigr),
 \tag{CC-AG.150}
\]

\[
 c_\ell=
 \begin{cases}
 \varepsilon_\ell,&\ell=p,\\
 v_\ell(p^{\ell-1}-1)+b_\ell,&\ell\ne p,\ \ell\ne2,\\
 v_2(p^2-1)+b_2,&\ell=2\ne p,
 \end{cases}
 \qquad
 B_\ell=c_\ell+v_\ell(M!)+v_\ell(h).
 \tag{CC-AG.151}
\]

Define

\[
 N=\prod_{\ell\in S}\ell^{B_\ell}.
 \tag{CC-AG.152}
\]

Empty products mean \(N=1\). Since \(B_\ell\ge v_\ell(h)\), we have
\(h\mid N\). Form the finite rational polynomial

\[
 P_{\xi,a}(X)=
 \operatorname{Norm}_{K/\mathbf Q}
       \left(X^N-\gamma^{aN/h}\right).
 \tag{CC-AG.153}
\]

Factor it over \(\mathbf Q\), and let \(\alpha_1,\ldots,\alpha_t\)
represent its distinct irreducible root orbits. Compute each source
Schur index \(d_j\) from the specified local-invariant formula, and put

\[
 m_j=\frac{[\mathbf Q(\alpha_j):K]d_j}{d_\xi}
     =\frac{[\mathbf Q(\alpha_j):\mathbf Q]d_j}
            {[K:\mathbf Q]d_\xi}.
 \tag{CC-AG.154}
\]

Here an embedding of \(K\) into each source centre is chosen through its
corresponding factor of the norm; the second formula is independent of
that choice. All \(m_j\) are positive integers, and

\[
                         H_a(\xi)=
                  \langle m_1,\ldots,m_t\rangle_{\mathbf N}.
 \tag{CC-AG.155}
\]

Thus the finite polynomial cutoff depends on the germ and the source
field, and \textbf{does not depend on a target multiplicity}. It
computes the exact monoid, its minimal generators, its gcd, its gaps
among multiples of that gcd, and a conductor relative to those
multiples. The displayed bounds are deliberately coarse. No practical
complexity estimate is asserted.

\subsubsection{1. Elementary principal-unit
facts}\label{elementary-principal-unit-facts}

For a prime \(\ell\), write \[
 U_\ell(s)=1+\ell^s\mathbf Z_\ell\qquad(s\ge\varepsilon_\ell).
\] We shall use two explicit elementary facts:

\[
 \overline{\langle u\rangle}=U_\ell(s)
 \quad\text{if }v_\ell(u-1)=s\ge\varepsilon_\ell,
 \qquad
 U_\ell(s)^m=U_\ell(s+v_\ell(m)).
 \tag{CC-AG.156}
\]

Here the first closure is in the \(\ell\)-adic topology, and the second
expression is the image of the \(m\)-th power map, for any positive
integer \(m\).

For a proof, binomial expansion gives \[
 v_\ell(u^n-1)=s+v_\ell(n)\qquad(n\ge1).
 \tag{CC-AG.157}
\] When \(\ell\nmid n\), the linear term has valuation \(s\) and every
higher term has larger valuation. For \(n=\ell\), the linear term has
valuation \(s+1\), the intermediate terms have larger valuation, and the
final term has larger valuation for \(s\ge1\) if \(\ell\) is odd and for
\(s\ge2\) if \(\ell=2\). Iterating proves the general formula.
Consequently the order of \(u\) modulo \(U_\ell(s+j)\) is exactly
\(\ell^j\), which is the order of that whole quotient. The generated
cyclic subgroup is therefore dense. The same argument applied to
\(u^m\), together with compactness of \(U_\ell(s)\), proves the
power-map formula in (CC-AG.156).

\subsubsection{2. A computable cyclotomic subgroup that fixes the
auxiliary
field}\label{a-computable-cyclotomic-subgroup-that-fixes-the-auxiliary-field}

Let \(\mu_{S^\infty}\) denote the roots of unity whose orders have all
prime divisors in \(S\). The cyclotomic character identifies

\[
 \operatorname{Gal}\bigl(\mathbf Q(\mu_{S^\infty})/\mathbf Q\bigr)
       =\prod_{\ell\in S}\mathbf Z_\ell^\times.
 \tag{CC-AG.158}
\]

This is the inverse limit of the usual finite cyclotomic Galois groups:
their orders and their faithful action on a primitive root give the full
groups \((\mathbf Z/n\mathbf Z)^\times\), and the Chinese remainder
theorem separates the prime factors. The finite degree input is Milne
(2020), Theorem 6.4(a), p.~99.

Put \[
                    C_\ell=c_\ell+v_\ell(M!).
 \tag{CC-AG.159}
\]

The following two assertions use the same constants:

\begin{itemize}
\tightlist
\item
  In \(E(\mu_{S^\infty})/E\), every cyclotomic action which is the
  identity at the primes other than \(\ell\), and is any specified
  element of \(U_\ell(C_\ell)\) at \(\ell\), is realized by an
  automorphism fixing \(E\).
\item
  For every \(p\)-adic completion \(E_w\), the identical assertion holds
  in \(E_w(\mu_{S^\infty})/E_w\).
\end{itemize}

We prove both assertions, including the independence from the exponents
in every other cyclotomic tower.

For the global assertion, the cyclotomic actions which fix \(E\) form a
subgroup of (CC-AG.158) of index \[
 [E\cap\mathbf Q(\mu_{S^\infty}):\mathbf Q]\le M.
\] Its quotient is finite abelian of order at most \(M\), so it is
killed by \(M!\). Thus the \(M!\)-th powers of all independently
supported principal units occur. By (CC-AG.156), the subgroup contains
the independently supported \(U_\ell(\varepsilon_\ell+v_\ell(M!))\).
Since \(c_\ell\ge\varepsilon_\ell\), this proves the stated global
assertion.

For the local assertion, first work over \(\mathbf Q_p\). Roots of unity
of order prime to \(p\) lie in the maximal unramified extension;
Frobenius acts on them by \(\zeta\mapsto\zeta^p\). Hence the image of
the local cyclotomic character on the factors indexed by
\(S\setminus\{p\}\) is the closure of the cyclic subgroup generated by
\(p\) in \(\prod_{r\in S\setminus\{p\}}\mathbf Z_r^\times\). The
\(p\)-power roots, when \(p\in S\), give the full independent factor
\(\mathbf Z_p^\times\): their cyclotomic extensions are totally
ramified, whereas the preceding extensions are unramified. The exact
primary inputs are Milne (2020), Proposition 6.2(a),(c), pp.~96--97,
Proposition 7.50 and Corollary 7.51, pp.~127--128, and Example 7.54,
p.~129. Equivalently, the prime-to-\(p\) assertion follows by
simple-root lifting from finite residue fields, and total ramification
follows from the Eisenstein polynomial for \(\zeta_{p^k}-1\).

We now isolate one prime \(\ell\ne p\) in that procyclic Frobenius
image. The continuous power map from the profinite integers sends \[
 t\in\widehat{\mathbf Z}=\prod_q\mathbf Z_q
          \quad\text{to}\quad p^t.
\] Take all coordinates \(t_q=0\) for \(q\ne\ell\). In a factor
\(\mathbf Z_r^\times\) with \(r\ne\ell\), the pro-\(r\) part is thereby
killed. Its remaining \(\ell\)-primary finite part has order dividing
\(\ell^{v_\ell(r-1)}\). This statement also covers \(r=2\ne\ell\), when
that \(\ell\)-primary part is trivial.

If \(\ell\) is odd, further take
\(t_\ell\in\ell^{b_\ell}\mathbf Z_\ell\). This kills all the other
factors. In its own \(\ell\)-factor, the pro-\(\ell\) component of \(p\)
generates the same closed subgroup as \(p^{\ell-1}\), since \(\ell-1\)
is an \(\ell\)-adic unit. By (CC-AG.156), its image after the additional
power is exactly \[
 U_\ell\bigl(v_\ell(p^{\ell-1}-1)+b_\ell\bigr)=U_\ell(c_\ell).
\] If \(\ell=2\ne p\), use \(t_2\in2^{b_2+1}\mathbf Z_2\). The
additional factor \(2\) kills the sign in \(\mathbf Z_2^\times\); the
powers of \(p^2\) generate \(U_2(v_2(p^2-1))\). The resulting
independently supported image is therefore \(U_2(c_2)\). Finally, for
\(\ell=p\), the independent totally ramified factor already contains
\(U_p(\varepsilon_p)=U_p(c_p)\).

Thus the local cyclotomic image over \(\mathbf Q_p\) contains each
independently supported \(U_\ell(c_\ell)\). Passing to \(E_w\) imposes a
subgroup of index \[
 [E_w\cap\mathbf Q_p(\mu_{S^\infty}):\mathbf Q_p]
       \le [E_w:\mathbf Q_p]\le M.
\] Again the quotient is finite abelian and is killed by \(M!\).
Applying (CC-AG.156) gives every independently supported
\(U_\ell(C_\ell)\), as claimed. This proof works at once for every \(w\)
and for all heights in the other prime towers.

\subsubsection{3. An orbit-degree lemma with a specified lower-conductor
label}\label{an-orbit-degree-lemma-with-a-specified-lower-conductor-label}

We isolate the exact group argument needed below. Let \(F\) be a
characteristic-zero field, let \(E_0/F\) be finite Galois, and let
\(x\in E_0^\times\) satisfy \(x^h\in F\). Suppose that in \[
                         E_0(\mu_{S^\infty})/E_0
\] every independently supported \(U_\ell(C)\) action occurs, for a
fixed \(\ell\in S\). Put \(H=v_\ell(h)\). Fix roots of unity at the
other primes in \(S\), with product \(\eta\) of order prime to \(\ell\).
Choose compatible primitive roots
\(\zeta_{\ell^{k+1}}^\ell=\zeta_{\ell^k}\), and define

\[
 x_k=x\zeta_{\ell^k}\eta,\qquad
 F_k=E_0(\zeta_{\ell^k},\eta).
 \tag{CC-AG.160}
\]

Then for every \(k\ge C+H\),

\[
 [F(x_{k+1}):F]=\ell[F(x_k):F].
 \tag{CC-AG.161}
\]

To prove it, write \(\Gamma_k=\operatorname{Gal}(F_k/F)\). Projection to
\(\operatorname{Gal}(E_0/F)\) is injective on the stabilizer of \(x_k\).
Indeed an automorphism which fixes \(E_0\) and \(x_k\) fixes the product
\(\zeta_{\ell^k}\eta\); its coprime primary components generate the
cyclotomic part of \(F_k\), so it fixes \(F_k\).

For \(\sigma\in\operatorname{Gal}(E_0/F)\), set
\(\lambda_\sigma=\sigma(x)/x\in\mu_h\). Decompose it into its primary
parts. A lift fixes \(x_k\) precisely when its action on each cyclotomic
component compensates the corresponding component of \(\lambda_\sigma\).
All conditions at primes other than \(\ell\) are independent of \(k\).
At \(\ell\), write \[
 (\lambda_\sigma)_\ell
        =\zeta_{\ell^k}^{\,j\ell^{k-H}}.
\] The required exponent of the cyclotomic action is \[
                       u_k=1-j\ell^{k-H}.
 \tag{CC-AG.162}
\] For \(k\ge C+H\) this is an element of \(U_\ell(C)\).

More explicitly, a lift at level \(k\) extends to the infinite Galois
extension. Its full \(\ell\)-adic cyclotomic action is congruent to
\(u_k\) modulo \(\ell^k\), hence lies in \(U_\ell(C)\). Multiply that
lift by the independently supported inverse action which fixes \(E_0\)
and all the other cyclotomic towers. This produces a lift of the same
\(\sigma\) whose action on the entire \(\ell\)-tower is the identity and
whose other required actions are unchanged. Conversely, from any such
lift with trivial \(\ell\)-action, multiplication by the independently
supported action (CC-AG.162) produces a lift fixing \(x_k\). Thus the
image of the stabilizer in \(\operatorname{Gal}(E_0/F)\), and
consequently its cardinality, is independent of \(k\) throughout this
range. This argument is uniform in \(\eta\) and works for every choice
of the primitive root.

For \(k\ge C\), the extension \(F_{k+1}/F_k\) has degree exactly
\(\ell\). Its degree is at most \(\ell\) by adjoining a root of
\(T^\ell-\zeta_{\ell^k}\). The independently supported subgroup
\(U_\ell(k)\subset U_\ell(C)\) acts trivially on \(F_k\) and has an
image of order \(\ell\) on \(\zeta_{\ell^{k+1}}\), proving the reverse
inequality. Therefore \(|\Gamma_{k+1}|=\ell|\Gamma_k|\). The
orbit-stabilizer formula, with the constant stabilizer just proved,
yields (CC-AG.161).

\subsubsection{4. Simultaneous global and local
recurrences}\label{simultaneous-global-and-local-recurrences}

Fix a root \(x=\alpha_0\) of \(X^h-\gamma^a\) in \(E\). The field \(E\)
in (CC-AG.149) contains \(\mu_h\), since quotients of the roots of one
binomial factor give these roots of unity. Every positive power of
\(\gamma\) generates \(K\), so \(K\subset\mathbf Q(\alpha_0)\).

Consider any label \[
 \alpha_k=\alpha_0\zeta_{\ell^k}\eta,
\] where \(\eta\) has order supported on \(S\setminus\{\ell\}\), and let
\[
 \alpha_{k+1}=\alpha_0\zeta_{\ell^{k+1}}\eta,
 \qquad \zeta_{\ell^{k+1}}^\ell=\zeta_{\ell^k}.
 \tag{CC-AG.163}
\] Both labels are Weil \(p^a\)-numbers with germ \(\xi\), and their
centres contain \(K\). In particular the lower label is specified
explicitly; it is not an unspecified field of a related degree.

Apply Section 3 first with \(F=K\), \(E_0=E\), and \(C=C_\ell\) of
(CC-AG.159). Since \(B_\ell=C_\ell+v_\ell(h)\), for \(k\ge B_\ell\) it
gives \[
 [\mathbf Q(\alpha_{k+1}):K]
            =\ell[\mathbf Q(\alpha_k):K].
 \tag{CC-AG.164}
\]

Now fix any embedding of the common Galois number field \[
                 E(\zeta_{\ell^{k+1}},\eta)
\] into \(\overline{\mathbf Q}_p\). Its restriction determines a place
\(v\) of \(K\) and a completion \(E_w\). Use Section 3 with \(F=K_v\),
\(E_0=E_w\), and the embedded conjugate of \(\alpha_0\). The completion
\(E_w/K_v\) is Galois, because \(E/\mathbf Q\) is Galois. Section 2
supplies the same \(U_\ell(C_\ell)\) subgroup. We obtain, for this
paired embedding,

\[
 [K_v(\alpha_{k+1}):K_v]
                =\ell[K_v(\alpha_k):K_v].
 \tag{CC-AG.165}
\]

The notation in (CC-AG.165) means the embedded elements. These fields
are exactly the corresponding completions of the two source centres.
Every embedding of either centre into \(\overline{\mathbf Q}_p\) extends
to the displayed common Galois field. Thus every local place at the
upper level and every local place at the lower level occurs in such a
pair. No local place is discarded by the comparison.

The quotient of the two labels in a pair is a root of unity, so their
normalized \(p\)-adic slopes are identical. Formula (CC-AG.165)
therefore multiplies the source local invariant by \(\ell\). Since all
centres contain the field \(K\) with no real embeddings, there are no
real-place invariants. All finite invariants away from \(p\) vanish.
Consequently, if \(d_k,d_{k+1}\) are the two source Schur indices, the
\emph{whole sets} of local invariants, with repetitions immaterial, give

\[
                  d_{k+1}=\frac{d_k}{\gcd(d_k,\ell)}.
 \tag{CC-AG.166}
\]

Indeed multiplying a finite set of rational classes of least common
denominator \(d_k\) by \(\ell\) changes that denominator to
\(d_k/\gcd(d_k,\ell)\). Combining (CC-AG.164)--(CC-AG.166) with the
exact transition multiplicity formula gives

\[
 \frac{r_{\alpha_{k+1}}}{r_{\alpha_k}}
                =\frac{\ell}{\gcd(d_k,\ell)}
                    \ \in\ \{1,\ell\}.
 \tag{CC-AG.167}
\]

Thus every step above the cutoff replaces a generator by a positive
integer multiple of the specified lower-conductor generator. There is no
need to assume that the two source centres are nested, or to wait until
an \(\ell\)-primary index has vanished.

\subsubsection{5. Reduction to the finite norm
polynomial}\label{reduction-to-the-finite-norm-polynomial}

The already proved prime-support lemma BR-15-PRIME-SUPPORT says that
every source label can be replaced by one whose \[
                         \alpha^h/\gamma^a
\] has order supported on \(S\), and whose positive transition
multiplicity divides the original multiplicity. This is a reduction for
a geometric target; it does not purport to preserve a prescribed
finite-field target label.

For such a restricted label, \(\zeta=\alpha/\alpha_0\) is itself a root
of unity supported on \(S\). Its \(h\)-th power is the indicated root of
unity. A prime outside \(S\) dividing the order of \(\zeta\) would still
divide the order of its \(h\)-th power, which is impossible. Thus all
the restricted labels have the form used in (CC-AG.163).

If the \(\ell\)-primary component of \(\zeta\) has order \(\ell^k\) with
\(k>B_\ell\), keep every other primary component fixed and replace that
component by its \(\ell\)-th power. Equation (CC-AG.167), applied with
lower level \(k-1\ge B_\ell\), shows that the old multiplicity is either
the same as, or \(\ell\) times, the new one. Repeating finitely many
times for each prime in \(S\) produces a label \(\alpha_0\zeta'\) with
\((\zeta')^N=1\), whose multiplicity divides that of the original source
label.

Because \(h\mid N\), \[
       \alpha_0^N=(\alpha_0^h)^{N/h}=\gamma^{aN/h}.
 \tag{CC-AG.168}
\] The reduced labels, with their conjugates, are exactly the roots of
(CC-AG.153). Conversely every root of that polynomial is a valid source
Weil number with the required germ orbit. If
\(\alpha^N=\sigma(\gamma)^{aN/h}\), taking complex absolute values gives
weight \(w\) relative to \(p^a\), taking valuations away from \(p\)
gives zero, and raising to \(h\) gives germ equality. Also
\(\sigma(K)\subset\mathbf Q(\alpha)\), by stabilization.

The polynomial lies in \(\mathbf Q[X]\) by the norm construction. It is
separable: each binomial has distinct roots, and the constants belonging
to different embeddings of \(K\) are distinct because \(\gamma^{aN/h}\)
generates \(K\). Its factors therefore give finitely many distinct
source simple classes. Their multiplicities (CC-AG.154) are positive
integers by the exact transition theorem.

We have proved that every unrestricted source generator is a positive
integer multiple of one of these finitely many generators. Conversely
each of the listed generators occurs as a source simple. Positivity and
semisimple decomposition prove (CC-AG.155).

If \(h=1\), the same conclusion holds with \(S\) empty and \(N=1\): the
prime-support step reduces every label to \(\gamma^a\). No tower
argument is then needed. This completes the proof of the uniform finite
cutoff for nonreal \(K\).

\subsubsection{6. Real centres, arithmetic effectivity, and exact monoid
output}\label{real-centres-arithmetic-effectivity-and-exact-monoid-output}

For completeness, a stabilized centre with a real embedding is
\(\mathbf Q\). If the weight \(w\) is even, the germ is the Tate germ
represented by \(p^{w/2}\), and a single copy is defined over the prime
field. If \(w\) is odd, the representative \(-p^w\) at exponent \(2\),
followed by a square root, gives an imaginary quadratic Weil
\(p\)-number of index one whose geometric image is one copy of the
rank-two geometric simple. These are exactly the fully proved
real-centre cases of CC-AG13. Base extension of this prime-field source
to \(\mathbf F_{p^a}\) still realizes one copy. Hence
\(H_a(\xi)=\mathbf N\) for every \(a\), as claimed.

All constants in the nonreal case are finite and effectively determined.
The minimal exponent and a witness \(\gamma\) are obtained by the finite
root equations of CC-AG13. A splitting field of the finite polynomial
(CC-AG.149), and its degree \(M\), are found by finite number-field
calculations. The valuations in (CC-AG.150)--(CC-AG.151) are valuations
of explicit nonzero integers; \(v_\ell(M!)\) can be computed as
\(\sum_{j\ge1}\lfloor M/\ell^j\rfloor\). The norm polynomial
(CC-AG.153), its rational factorization, and the local-index
computations are therefore finite exact operations.

Here is also a direct finite procedure for the additional monoid data
asserted after (CC-AG.155). Let \(g=\gcd(m_1,\ldots,m_t)\), put
\(a_j=m_j/g\), and let \(m=\min_j a_j\). On the finite set
\(\mathbf Z/m\mathbf Z\), give each directed edge \[
                     x\longrightarrow x+a_j
\] the positive integer cost \(a_j\). All vertices are reachable from
\(0\): the nonnegative sums in a finite group form a subgroup, and the
gcd of the \(a_j\) is one. Finite shortest-path calculation therefore
gives a least cost \(A_x\) for every residue \(x\). A cost is the
corresponding nonnegative integer sum, so \(A_x\) is precisely the least
element of the divided monoid in that residue class.

An integer \(n\ge0\) is in the divided monoid if and only if \[
                          n\ge A_{n\bmod m}.
 \tag{CC-AG.169}
\] Necessity is minimality. For sufficiency, the difference is a
nonnegative multiple of the generator \(m\). Hence the conductor is \[
                         \max_x A_x-m+1,
 \tag{CC-AG.170}
\] with value \(0\) when \(m=1\). Multiplication by \(g\) gives the
conductor relative to \(g\mathbf N\) and all gaps among those multiples.
Every indecomposable positive element of the original monoid belongs to
the original finite generator list, because any expression using at
least two positive summands decomposes it. First remove repeated
integers from the list. Testing each remaining generator against the
bounded positive combinations of the others then gives the minimal
generator list. Thus the theorem computes all the claimed monoid data.
\(\square\)

\subsubsection{Scope}\label{scope}

The cutoff lists enough source classes to generate all possible
multiplicities. It does \textbf{not} say that every source simple
belongs to that finite list. The earlier rank-dependent procedure
CC-AG12 remains the statement that enumerates every preimage class of a
specified target object. For a geometric object with finitely many
simple isotypic components, the present theorem gives one finite
positive multiplicity system for each orbit, valid for all target
multiplicities in that orbit.

The whole construction takes place in the established
abelian--Artin--Tate sector and does not construct arbitrary-variety
cross-correspondences. It therefore closes the uniform
coefficient-descent calculation in that sector while leaving the stated
arbitrary-motive/full-Tate gate explicit. No cyclicity or norm-only gcd
formula is asserted; CC-AG14--15 show why those shortcuts are
inadequate.

\subsection{Primary references and exact
locators}

\begin{itemize}
\tightlist
\item
  Milne, J. S. (1994). Motives over finite fields. In U. Jannsen, S.
  Kleiman, \& J.-P. Serre (Eds.), \emph{Motives} (Proceedings of
  Symposia in Pure Mathematics, Vol. 55, Part 1, pp.~401--459). American
  Mathematical Society. Version used: corrected author manuscript,
  November 2, 2019. \url{https://www.jmilne.org/math/xnotes/Seattle1.pdf} .
  Exact local-invariant and simple-object inputs: convention p.~17;
  Theorems 2.15--2.16, p.~18; Theorems 2.18--2.19, p.~20, applied only
  within the established abelian sector.
\item
  Milne, J. S. (2020, July 19). \emph{Algebraic number theory} (Version
  3.08) {[}Author course notes{]}.
  \url{https://www.jmilne.org/math/CourseNotes/ANT.pdf} . Exact locators:
  prime-power cyclotomic degree and total ramification, Proposition
  6.2(a),(c), pp.~96--97; cyclotomic degree, Theorem 6.4(a), p.~99;
  simple-root lifting, Proposition 7.31, p.~120; unramified extensions
  and roots of unity, Proposition 7.50 and Corollary 7.51, pp.~127--128;
  Frobenius description, Example 7.54, p.~129; completion decomposition,
  Propositions 8.1--8.2, p.~136. The principal-unit estimates,
  independently supported deep-subgroup bound, orbit-degree lemma,
  simultaneous local recurrence, conductor cutoff and exact
  positive-monoid deduction are proved in full above.
\end{itemize}

% END consequences_criteria_uniform_descent.tex

% BEGIN consequences_criteria_descent_profile.tex
\section{The complete finite-field profile of geometric rational
descent}\label{the-complete-finite-field-profile-of-geometric-rational-descent}

\subsection{CC-AG17 --- Periodicity, divisor tables, and all source
fields}\label{cc-ag17-periodicity-divisor-tables-and-all-source-fields}

Use the established rational abelian--Artin--Tate categories and the
exact simple-object transition formulas of CC-AG08--09. Let

\[
 \xi=[\gamma,h],\qquad
 h=h_\xi,\qquad K=\mathbf Q(\gamma^n)\quad(n\ge1),
 \tag{CC-AG.171}
\]

be a stabilized geometric germ of weight \(w\), represented at its least
positive exponent. Let \(H_a(\xi)\) be its geometric multiplicity monoid
for source \(\mathbf F_{p^a}\), for \(a\ge1\).

Then

\[
 H_{a+h}(\xi)=H_a(\xi),\qquad
 H_a(\xi)=H_{\gcd(a,h)}(\xi),\qquad H_h(\xi)=\mathbf N.
 \tag{CC-AG.172}
\]

More generally,

\[
 \gcd(a,h)\mid\gcd(b,h)
          \quad\Longrightarrow\quad H_a(\xi)\subseteq H_b(\xi).
 \tag{CC-AG.173}
\]

The periodic equality has an explicit bijection on \textbf{all} source
simple-label orbits, preserving the actual geometric multiplicity. After
aligning a source label with the germ \(\xi\), that bijection is

\[
                 \alpha\longmapsto\alpha\gamma
       \quad\text{from source exponent }a\text{ to }a+h,
 \tag{CC-AG.174}
\]

with inverse \(\beta\mapsto\beta/\gamma\). These two labels have exactly
the same centre and Schur index. The equality involving \(\gcd(a,h)\) is
an equality of multiplicity monoids; it does not assert this particular
bijection of simple orbits between arbitrary equal-gcd exponents.

Consequently the entire monoid profile over all finite source fields is
determined by the finite table

\[
                         \{H_d(\xi):d\mid h\},
 \tag{CC-AG.175}
\]

which CC-AG16 computes by finite algebraic calculations. This uses the
full positive monoids, including their gaps, rather than replacing them
by gcd divisibility tests.

\subsubsection{Complete proof}

We first make the orbit alignment in (CC-AG.174) precise. The stabilizer
of the germ \(\xi\) in
\(\operatorname{Gal}(\overline{\mathbf Q}/\mathbf Q)\) is exactly
\(\operatorname{Gal}(\overline{\mathbf Q}/K)\). One inclusion is
immediate because \(\gamma\in K\). For the other, if \(\sigma\xi=\xi\),
then some positive powers of the two stage-\(h\) representatives agree:
\[
                         \sigma(\gamma)^N=\gamma^N
\] for a positive integer \(N\). Stabilization says that this power
generates \(K\), so \(\sigma\) fixes \(K\) pointwise, and in particular
fixes \(\gamma\).

Every source Galois orbit whose geometric germ orbit is \(O_\xi\) has a
representative with germ exactly \(\xi\). Two such aligned
representatives in the same full Galois orbit differ by an automorphism
fixing \(K\), by the preceding stabilizer calculation. Thus these full
source orbits correspond precisely to the
\(\operatorname{Gal}(\overline{\mathbf Q}/K)\)-orbits of aligned labels.

Take an aligned source label \(\alpha\in W(p^a)\). For some \(N\ge1\),

\[
                         \alpha^{hN}=\gamma^{aN}.
 \tag{CC-AG.176}
\]

It follows that \(K\subset\mathbf Q(\alpha)\). Put
\(\beta=\alpha\gamma\). At every complex embedding, \[
 |\beta|=p^{aw/2}p^{hw/2}=p^{(a+h)w/2},
\] and its valuations away from \(p\) vanish. Therefore \(\beta\) is a
Weil \(p^{a+h}\)-number of weight \(w\). Equation (CC-AG.176) gives

\[
                        \beta^{hN}=\gamma^{(a+h)N},
 \tag{CC-AG.177}
\]

so its germ is again \(\xi\), and \(K\subset\mathbf Q(\beta)\). Since
\(\gamma\in K\), the inclusions in both directions now give

\[
                       \mathbf Q(\alpha)=\mathbf Q(\beta)=L.
 \tag{CC-AG.178}
\]

Conversely, if \(\beta\) is an aligned label at source exponent \(a+h\),
division by \(\gamma\) gives the correct Weil absolute values, the unit
condition away from \(p\), and equation (CC-AG.176) at exponent \(a>0\).
It is therefore the inverse operation. Multiplication and division are
equivariant for automorphisms fixing \(K\), so the aligned-label
bijection induces the claimed bijection on all source Galois orbits.

The source indices also agree. At a \(p\)-adic place \(z\) of the common
centre \(L\), equation (CC-AG.176) implies \[
 \frac{\operatorname{ord}_z(\alpha)}{a}
       =\frac{\operatorname{ord}_z(\gamma)}{h}.
\] Hence

\[
 \frac{\operatorname{ord}_z(\beta)}{a+h}
       =\frac{\operatorname{ord}_z(\alpha)}{a}.
 \tag{CC-AG.179}
\]

The normalized slopes, local field degrees and therefore the \(p\)-adic
Brauer invariants in the source-index formula are the same. At every
real place the weight is the same, so the invariant \(w/2\) is
unchanged; at the remaining places both invariants vanish. Their least
common denominators are equal:

\[
                              d_\beta=d_\alpha.
 \tag{CC-AG.180}
\]

Using the exact transition formula, \[
 r_\beta=[L:K]d_\beta/d_\xi
        =[L:K]d_\alpha/d_\xi=r_\alpha.
 \tag{CC-AG.181}
\] The bijection of simple orbits therefore preserves every positive
multiplicity sum, giving \(H_{a+h}(\xi)=H_a(\xi)\). It also gives a
bijection of isomorphism classes of all source objects with geometric
image \(S_\xi^{\oplus r}\), for every fixed \(r\), by transporting their
simple summands and multiplicities.

For the divisor assertions, finite-field base extension supplies

\[
                         H_a(\xi)\subseteq H_{ua}(\xi)
                       \qquad(u\in\mathbf Z_{>0}),
 \tag{CC-AG.182}
\]

because the subsequent geometric base extension is the same. If
\(\gcd(a,h)\mid\gcd(b,h)\), the congruence \(au\equiv b\pmod h\) has an
integer solution. Adding a positive multiple of \(h/\gcd(a,h)\) makes
\(u\) positive. Periodicity equates the monoids at the two positive
congruent exponents \(au\) and \(b\). Combining with (CC-AG.182) proves
(CC-AG.173). Taking \(b=\gcd(a,h)\), and then using the reverse
inclusion since \(b\mid a\), gives \(H_a(\xi)=H_{\gcd(a,h)}(\xi)\).

Finally the finite label \(\gamma\) over \(\mathbf F_{p^h}\) has centre
\(K\). Its local invariants are exactly the geometric invariants of the
stabilized germ, including the real-place convention, so its finite
index is \(d_\xi\). The transition multiplicity is \(1\). Direct sums
give \(H_h(\xi)=\mathbf N\). This proves all assertions in
(CC-AG.172)--(CC-AG.175). \(\square\)

\subsubsection{The field-of-definition profile of a specified geometric
object}\label{the-field-of-definition-profile-of-a-specified-geometric-object}

Let

\[
                  Z=\bigoplus_{j=1}^s S_{\xi_j}^{\oplus r_j}
 \tag{CC-AG.183}
\]

be a geometric rational object in the established sector, with distinct
simple germ orbits and positive multiplicities. Put

\[
                              H=\operatorname{lcm}(h_{\xi_1},\ldots,h_{\xi_s}).
 \tag{CC-AG.184}
\]

There is an effectively computable finite antichain \(\mathcal D_Z\) of
divisors of \(H\), with respect to divisibility, such that for every
\(a\ge1\),

\[
 Z\text{ comes from }\mathcal A_{p^a}
       \quad\Longleftrightarrow\quad
                 d\mid a\text{ for some }d\in\mathcal D_Z.
 \tag{CC-AG.185}
\]

Here ``comes from'' means isomorphism after geometric base extension; it
does not choose or preserve a finite target-field model.

\textbf{Proof and finite construction.} For every divisor \(e\mid H\),
test the finite collection of conditions

\[
                r_j\in H_{\gcd(e,h_{\xi_j})}(\xi_j)
                         \qquad(1\le j\le s).
 \tag{CC-AG.186}
\]

These are decidable either from the whole-monoid tables of CC-AG16 or
directly from the earlier specified-object procedure CC-AG12. By
positivity and the one-orbit image of each source simple, their
conjunction is exactly descent of \(Z\) from \(\mathbf F_{p^e}\).

The set of successful divisors is nonempty: \(e=H\) succeeds by
\(H_{h_{\xi_j}}(\xi_j)=\mathbf N\). It is upward closed under
divisibility, by finite-field base extension. Define \(\mathcal D_Z\) to
be its minimal elements. This is a finite, effectively computable
antichain.

For any \(a\ge1\), set \(e=\gcd(a,H)\). Then
\(\gcd(e,h_{\xi_j})=\gcd(a,h_{\xi_j})\) for every \(j\). The first part
of the theorem shows that descent at \(a\) is equivalent to success of
the divisor \(e\). Since the successful set is upward closed in the
finite divisor poset, this is equivalent to \(d\mid e\) for some
\(d\in\mathcal D_Z\). Because every such \(d\) divides \(H\), this is
equivalent to \(d\mid a\), proving (CC-AG.185).

For the zero object use \(H=1\) and \(\mathcal D_Z=\{1\}\), giving
descent from every source field. No assumption that the antichain has
only one element is needed or made. \(\square\)

\subsubsection{The complete profiles of the two explicit threefold
germs}\label{the-complete-profiles-of-the-two-explicit-threefold-germs}

For each of the examples CC-AG14 and CC-AG15, \(h=9\) and \(d_\xi=3\).
Their entire source-field profiles are therefore:

\begin{longtable}[]{@{}
  >{\raggedright\arraybackslash}p{(\columnwidth - 4\tabcolsep) * \real{0.3333}}
  >{\raggedright\arraybackslash}p{(\columnwidth - 4\tabcolsep) * \real{0.3333}}
  >{\raggedright\arraybackslash}p{(\columnwidth - 4\tabcolsep) * \real{0.3333}}@{}}
\toprule\noalign{}
\begin{minipage}[b]{\linewidth}\raggedright
Condition on the source exponent \(a\)
\end{minipage} & \begin{minipage}[b]{\linewidth}\raggedright
CC-AG14, characteristic \(2\)
\end{minipage} & \begin{minipage}[b]{\linewidth}\raggedright
CC-AG15, characteristic \(163\)
\end{minipage} \\
\midrule\noalign{}
\endhead
\bottomrule\noalign{}
\endlastfoot
\(3\nmid a\) & \(3\mathbf N\) & \(3\mathbf N\) \\
\(3\mid a,\ 9\nmid a\) & \(\langle2,3\rangle_{\mathbf N}\) &
\(3\mathbf N\) \\
\(9\mid a\) & \(\mathbf N\) & \(\mathbf N\) \\
\end{longtable}

Indeed the three divisors of \(h\) are \(1,3,9\). At \(1\), CC-AG13
gives \((h/d_\xi)\mathbf N=3\mathbf N\); at \(3\), CC-AG14 and CC-AG15
give the two different monoids displayed; at \(9\), (CC-AG.172) gives
\(\mathbf N\). The gcd reduction proves every entry for all positive
\(a\).

\subsubsection{Source scope and what has been
computed}\label{source-scope-and-what-has-been-computed}

For a single germ, only the divisor table of its computable exponent
\(h_\xi\) is needed to know its multiplicity monoid over every finite
source field. For a specified geometric object, (CC-AG.185) computes its
entire geometric field-of-definition profile as a finite union of
divisibility conditions. A fixed source's actual preimage isomorphism
classes remain those of CC-AG12; the periodic orbit bijection
(CC-AG.174) transports them between source exponents differing by
\(h_\xi\).

The proofs use the inspected Milne classification and transition theorem
in the established abelian--Artin--Tate sector. They do not supply
correspondences placing an arbitrary numerical motive in that sector,
and do not make a finite-target descent statement by changing a
Frobenius label.

\subsection{Primary reference}

Milne, J. S. (1994). Motives over finite fields. In U. Jannsen, S.
Kleiman, \& J.-P. Serre (Eds.), \emph{Motives} (Proceedings of Symposia
in Pure Mathematics, Vol. 55, Part 1, pp.~401--459). American
Mathematical Society. Version used: corrected author manuscript,
November 2, 2019. \url{https://www.jmilne.org/math/xnotes/Seattle1.pdf} .
Exact classification and local-invariant inputs: convention p.~17;
Theorems 2.15--2.16, p.~18; Theorems 2.18--2.19, p.~20. The orbit
alignment, periodic bijection, divisor-table reduction and finite
field-of-definition profile are proved in full above.

% END consequences_criteria_descent_profile.tex

% BEGIN root_class_number_descent.tex
\section{A class-number shortcut to finite-field descent fails}
\label{root-class-number-descent}

\textbf{Scope.} This source check concerns one exact assertion in
Thakur's January 2020 preprint, not the other results of that paper.
It uses the classical abelian-variety corollary of CC-AG14; no
additional theorem about higher-codimension abelian cycles, and no
adopted analytic Riemann hypothesis, enters the counterexample.

\subsection{BR-15-CLASSBOUND --- Refutation of the stated class-number bound}

Thakur (2020, Definition 2.3 and Proposition 2.9, p.~6,
arXiv:2001.00210v1) defines a model by isogeny after base extension
and asserts that a type $\mathrm{IV}(1,d)$ abelian variety has such a
model over $\mathbf F_{p^{h_K}}$, where $h_K$ is the class number of
the centre $K$ of its geometric endomorphism algebra. The assertion
is false, including if one weakens the demanded isogeny to a geometric
isogeny.

\textbf{Counterexample and complete deduction.} Let
\[
 K=\mathbf Q(s),\quad s^2=-23,\quad
 \omega=(1+s)/2,\quad \eta=1+\omega,\quad
 \gamma=8\bar\eta=12-4s.
\]
By the fully proved classical corollary of CC-AG14, there is a
geometrically simple abelian threefold $A/\mathbf F_{2^9}$ with
Frobenius polynomial
\[
                 (T^2-24T+512)^3,
\]
and no abelian variety over $\mathbf F_8$ is geometrically isogenous
to $A$. The corollary proves this using the classical Honda--Tate
classification and its local-index formula: every possible source
label giving a single geometric copy has degree three over $K$ and
would need index one, whereas the complete root-of-unity and local
cube calculation forces index three. Thus its nonexistence conclusion
allows every source label and every root-of-unity twist; it is not
only a failure to take a particular cube root of $\gamma$.

We verify the remaining class-number hypothesis explicitly. The ring
of integers is $\mathcal O_K=\mathbf Z[\omega]$, its discriminant is
$-23$, and
\[
 N_{K/\mathbf Q}(x+y\omega)=x^2+xy+6y^2
             =(x+y/2)^2+23y^2/4.
\]
The maximality of this quadratic order follows because its
discriminant is squarefree: the square of the index of any larger
order would divide $23$. Put
\[
 \mathfrak p=(2,\omega-1),\quad
 \bar{\mathfrak p}=(2,\omega),\quad
 \mathfrak q=(3,\omega),\quad
 \bar{\mathfrak q}=(3,\omega-1).
\]
The polynomial of $\omega$ is $T^2-T+6$; its distinct roots modulo
$2$ and modulo $3$ give precisely these primes, each with residue
degree one. In particular
$(2)=\mathfrak p\bar{\mathfrak p}$ and
$(3)=\mathfrak q\bar{\mathfrak q}$.
The norm of $\eta$ is $8$, and its residues at the two primes above
$2$ show $(\eta)=\mathfrak p^3$. No element has norm $2$: if $y\ne0$
the displayed quadratic form is at least $23/4$, and if $y=0$ it is
a square. Hence $\mathfrak p$ is not principal and its class has
order exactly three. Also $N(\omega)=6$, and its residues give
\[
                  (\omega)=\bar{\mathfrak p}\mathfrak q,
             \qquad [\mathfrak q]=[\mathfrak p].
\]
Minkowski's ideal-class bound (Milne, 2020, Theorem 4.3, p.~70)
gives an integral representative of norm at most
$2\sqrt{23}/\pi<4$ for every class. An ideal of norm $2$ or $3$ is
prime, so the above four primes, together with the unit ideal,
exhaust these representatives. Their classes are
$1,[\mathfrak p],[\mathfrak p]^{-1}$. Therefore $h_K=3$.

The endomorphism division algebra of $A$ has centre $K$, degree
three, and is ramified only at the two primes above $2$, with
invariants of denominators three. These statements are part of the
Honda--Tate computation in CC-AG14. They also describe the geometric
endomorphism algebra: over every finite extension the centre and
index remain unchanged, and the injective maps between the
endomorphism algebras therefore have the same finite dimension and
are isomorphisms. Every geometric endomorphism descends to a finite
extension. A polarization gives the required involution of the
second kind. Thus $A$ satisfies Thakur's Definition 2.1, p.~3, for
type $\mathrm{IV}(1,3)$.

The asserted bound would now give a model over
$\mathbf F_{2^{h_K}}=\mathbf F_8$, which is a subfield of the
displayed field of definition $\mathbf F_{2^9}$. Any isogeny over
$\mathbf F_{2^9}$ would give a geometric isogeny, contradicting
CC-AG14. This proves the refutation in the source's own model
convention. $\square$

\subsection{Location of the failed inference}

The proof of the cited Proposition 2.9 chooses a generator $\pi_0$
of $\mathfrak p^{h_K}$ and assigns its associated abelian variety
dimension $d$ and the original division algebra. When $p$ splits
in the imaginary quadratic field, this generator instead has slopes
$1$ and $0$. The classical local-index formula gives index one;
its centre is $K$, since its two valuations are different. Its
Honda--Tate variety is therefore an ordinary elliptic curve.
For $d>1$, it cannot supply the asserted variety. This is the
precise failed inference; the argument above independently
refutes the conclusion rather than relying only on a gap in that
proof.

\subsection{BR-15-SLOPEBOUND --- The corrected geometric construction}

For clarity, the slope-sensitive construction in Thakur (2020,
Proposition 2.4, pp.~4--5) has a different exponent. Here is the
exact form needed for geometric descent.

Let $A$ be a geometrically simple abelian variety over a finite field
of characteristic $p$ whose geometric endomorphism division algebra has imaginary
quadratic centre $K$ and degree $d>1$. Suppose $p$ splits in $K$,
and write its invariant at one prime $\mathfrak p$ as $j/d$ with
$0<j<d$ and $\gcd(j,d)=1$. Let $t$ be the order of
$[\mathfrak p]$ in $\operatorname{Cl}(K)$. Then the geometric
isogeny class of $A$ has a representative over
$\mathbf F_{p^{td}}$, hence also over $\mathbf F_{p^{h_Kd}}$.
This is a sufficient geometric field, not a claimed smallest field
and not a prescribed finite-field Frobenius descent.

\textbf{Complete proof of this form.} Choose a finite field
$\mathbf F_{p^N}$ defining $A$ and a finite $\mathbf Q$-basis of
its geometric endomorphism algebra. Such a field exists because each
basis element, after clearing a denominator, is a morphism defined
over some finite extension. Thus the rational endomorphism algebra
over this field is already the full geometric algebra.
Frobenius $\pi$ commutes with those endomorphisms, so $\pi\in K$;
the Honda--Tate centre formula gives $\mathbf Q(\pi)=K$.
At the split primes the local degree is one. Its slopes lie in
$[0,1]$ and represent the nonzero invariants $j/d$ and
$(d-j)/d$, so they equal those fractions themselves. In particular
$d\mid N$ and $\dim A=d$.

Choose $\eta\in\mathcal O_K$ with $(\eta)=\mathfrak p^t$.
Its norm is $p^t$. Put
\[
                  \rho=\eta^j\bar\eta^{\,d-j}.
\]
It is an integral Weil $p^{td}$ number with the same two slopes.
The fractions cannot both be $1/2$: if they were, the two valuations
of $\pi$ would be equal, and $\pi^2/p^N$ would be a unit of
the imaginary quadratic field. Its finite order would make a power
of $\pi$ rational, contradicting the unchanged geometric centre
$K$ over every finite extension where all endomorphisms are already
defined. Hence $j\ne d-j$, and $\rho$ and all its positive powers
have unequal valuations and centre $K$. Their index is $d$.
Honda--Tate therefore gives an abelian variety $A_0$ of dimension
$d$ over $\mathbf F_{p^{td}}$ with Frobenius $\rho$.

The elements $\rho^N$ and $\pi^{td}$ have identical valuations
at every prime of $K$, so their ratio is a unit. The unit group of
an imaginary quadratic field is finite (equivalently, its lattice
has only finitely many points on the unit circle). If this ratio
has order $e$, then
\[
                         \rho^{Ne}=\pi^{tde}.
\]
Over $\mathbf F_{p^{Ntde}}$ the two Frobenius eigenvalues and their
conjugates thus agree, with the same multiplicity $d$. Tate's
isogeny criterion makes the two base extensions isogenous. This
gives the geometric representative $A_0$. Finally $t\mid h_K$,
so extension of $A_0$ to $\mathbf F_{p^{h_Kd}}$ is available.
$\square$

The exact finite divisor profile of CC-AG17 remains stronger than
this sufficient construction. The present source correction is not
counted as a reduction of the remaining arbitrary-variety cycle gate.

\subsection{Primary locators}

Thakur, S. (2020). \emph{Isogenies of certain abelian varieties over
finite fields with $p$-ranks zero} (arXiv:2001.00210v1, January 1).
\url{https://arxiv.org/abs/2001.00210v1} . Definition 2.1, p.~3;
Proposition 2.4, pp.~4--5; Definition 2.3 and Proposition 2.9 with
its proof, p.~6. The refutation refers exactly to this version.

Milne, J. S. (2020). \emph{Algebraic number theory} (Version 3.08,
July 19). Theorem 4.3, p.~70, ideal-class Minkowski bound.
\url{https://www.jmilne.org/math/CourseNotes/ANT.pdf} .

Tate, J. (1969). Classes d'isog\'enie des vari\'et\'es ab\'eliennes
sur un corps fini (d'apr\`es T.~Honda). \emph{S\'eminaire Bourbaki},
Expos\'e 352, pp.~95--110. Theorem 1(i)--(ii) and Remark 2,
p.~96. \url{https://www.numdam.org/article/SB_1968-1969__11__95_0.pdf} .

Tate, J. (1966). Endomorphisms of abelian varieties over finite
fields. \emph{Inventiones Mathematicae, 2}, 134--144.
Theorem 1(c), p.~139, Frobenius-polynomial isogeny criterion.
\url{https://doi.org/10.1007/BF01404549} .

% END root_class_number_descent.tex

% BEGIN consequences_criteria_cross_cycles.tex
\section{A finite cross-correspondence criterion for a specified
variety}\label{a-finite-cross-correspondence-criterion-for-a-specified-variety}

\textbf{Scope and input.} Fix \(k=\mathbf F_q\), a prime
\(\ell\ne\operatorname{char}k\), and \(E=\overline{\mathbf Q}_\ell\),
with the algebraic-number embedding used in CC-AG01. This section uses
the proved abelian-sector equivalence CC-AG01, and therefore the
inspected cited all-abelian Tate and numerical--homological inputs
underlying AQ.15. It gives an exact remaining criterion for a specified
smooth projective variety; it does not assert that its cycle condition
always holds. Literature cutoff: 20 November 2026; proof checks: 21
November 2026 UTC.

\subsection{The homological category used
here}\label{the-homological-category-used-here}

Let \(\mathcal H_{\ell,E}(k)\) be the category of pure motives formed
from \textbf{\(E\)-linear algebraic correspondences modulo the kernel of
their \(E\)-valued \(\ell\)-adic cycle class}, followed by idempotent
completion and inversion of the Tate object. Use the usual correction of
commutativity by cohomological parity. Its realization \(R\) on the
ordinary direct sum of cohomology is faithful. Indeed, Poincaré duality
and the Künneth formula identify the class of a correspondence with its
graded linear action; being zero under the action is exactly the
equivalence relation imposed.

This coefficient convention matters. For arbitrary varieties we do not
identify this category with a naive scalar extension of rational
homological Hom spaces without proving scalar injectivity. Nor do we
begin by realizing arbitrary numerical motives. On the abelian--Artin
subcategory, the numerical--homological equality and scalar-injectivity
proof in AQ.15 do identify it with \(\mathcal A_{q,E}\); hence that
subcategory embeds fully faithfully in \(\mathcal H_{\ell,E}(k)\).

All cohomological projectors needed here are algebraic in the
homological category. The characteristic polynomials of Frobenius on
distinct cohomological degrees are relatively prime by their different
weights. The Chinese remainder theorem produces polynomials in the
Frobenius correspondence acting as identity in one degree and zero in
the others. Their idempotent identities hold after realization and thus
in this category. Frobenius semisimplicity is not required for this
construction: the full characteristic polynomials already annihilate the
relevant generalized eigenspaces. This is the finite-field construction
discussed by Milne (1994), author-manuscript pp.~5--6.

\subsection{CC-XC01 --- An abelian target realizing the
semisimplification}\label{cc-xc01-an-abelian-target-realizing-the-semisimplification}

Let \(X/k\) be nonempty, smooth and projective of pure dimension \(d\),
and put

\[
 H=h_{\ell,E}(X),\qquad V=R(H)=\bigoplus_i H^i_{\mathrm{et}}(X_{\bar k},E).
\]

There exist finitely many \(Y_\nu/k\), each an abelian variety or a
finite étale scheme, and an idempotent \(e\) on

\[
 T=\bigoplus_{\nu=1}^s h(Y_\nu),\qquad M=(T,e)\in\mathcal A_{q,E},
\]

such that \(W=R(M)\) is isomorphic, as a Frobenius module, to
\(V^{\mathrm{ss}}\). In particular \(\dim_E W=\dim_E V\). No
semisimplicity assumption on \(V\) is made.

\textbf{Proof.} Deligne's Weil theorem (1974, Theorem 1.6, p.~276) gives
integral algebraic Frobenius eigenvalues of weight \(i\) on \(H^i(X)\).
Their products with their complex conjugates are \(q^i\), so they are
\(p\)-units. Thus every eigenvalue belongs to \(W(q)\), independently of
possible Jordan blocks.

For an eigenvalue \(\alpha\) in a positive degree \(i\), choose
\(\gamma\) with \(\gamma^i=\alpha\). As in AQ-WEIL-EXIST, \(\gamma\) is
an integral weight-one Weil number. Honda's Main Theorem gives an
abelian variety \(B_\alpha/k\) with eigenvalue \(\gamma\) on \(H^1\).
Then \(H^i(B_\alpha^i)\) contains the Künneth summand
\(H^1(B_\alpha)^{\otimes i}\) and hence the eigenvalue \(\alpha\). The
eigenvalues in degree zero are roots of unity, since Frobenius permutes
the finitely many geometric connected components of \(X\). A root of
unity of order \(r\) occurs in
\(H^0(\operatorname{Spec}\mathbf F_{q^r})\).

Use enough copies of these varieties to accommodate every algebraic
multiplicity in the characteristic polynomial of \(V\). This gives a
finite \(T\) whose semisimple realization contains the desired
eigenspaces with at least those multiplicities. Project onto the desired
subspaces. This projector commutes with Frobenius; AQ.15 after
coefficient extension lifts it to \(e\). Its image has exactly
\(V^{\mathrm{ss}}\) as realization. No Tate twist was necessary in this
construction, because all degrees of \(X\) are nonnegative and its
eigenvalues are integral. \(\square\)

\subsection{CC-XC02 --- A determinant condition on one projected
cross-cycle
map}\label{cc-xc02-a-determinant-condition-on-one-projected-cross-cycle-map}

Fix \(T,e,M,W\) as in CC-XC01. A codimension-\(d\) cycle \(z_\nu\) on
\(X\times Y_\nu\) acts in degree zero from \(h(X)\) to \(h(Y_\nu)\).
Define the projected cross-cycle map

\[
 \begin{split}
 \kappa:\ 
 \bigoplus_{\nu=1}^s
 \bigl(\mathrm{CH}^d(X\times Y_\nu)_{\mathbf Q}
              \otimes_{\mathbf Q}E\bigr)
 &\longrightarrow \operatorname{Hom}_{E[F]}(V,W),\\
 (z_\nu)_\nu&\longmapsto R(e)\circ R((z_\nu)_\nu).
 \end{split}                                      \tag{CC-XC.1}
\]

Write \(I=\operatorname{im}\kappa\). It is a finite-dimensional
\(E\)-linear subspace of \(\operatorname{Hom}_E(V,W)\). Choose bases of
the equal-dimensional spaces \(V,W\); the condition
\(\det|_I\not\equiv0\) is independent of those choices.

\textbf{Theorem.} The following conditions are equivalent:

\begin{enumerate}
\def\labelenumi{\arabic{enumi}.}
\tightlist
\item
  \(H\) is isomorphic in \(\mathcal H_{\ell,E}(k)\) to an object of
  \(\mathcal A_{q,E}\).
\item
  \(H\simeq M\) in that category.
\item
  \(I\) contains an invertible linear map.
\item
  The determinant polynomial is not identically zero on \(I\).
\end{enumerate}

When they hold, Frobenius on \(V\) is automatically semisimple. Thus the
determinant condition incorporates both semisimplicity and the needed
cross-correspondence algebraicity.

\textbf{Complete proof.} By the definition of the homological category
and the idempotent completion,

\[
 I=R\bigl(\operatorname{Hom}_{\mathcal H_{\ell,E}}(H,M)\bigr).
                                                        \tag{CC-XC.2}
\]

The codimension in (CC-XC.1) is \(d=\dim X\), independently of
\(\dim Y_\nu\): the action

\[
 a\longmapsto(p_{Y_\nu})_*\bigl(p_X^*a\smile\operatorname{cl}(z_\nu)\bigr)
\]

adds \(2d\) and then subtracts \(2d\) from cohomological degree, with
cancelling Tate twists. The cycle is defined over \(k\), so its action
commutes with Frobenius. Different cohomological degrees have disjoint
Frobenius spectra, so a Frobenius-equivariant map between these
realizations preserves the degrees.

Suppose (3). By (CC-XC.2), choose \(f:H\to M\) with invertible
realization \(\phi:V\to W\). Its Frobenius equivariance shows that \(V\)
is semisimple, since \(W\) is. We must still prove that \(f\) has an
algebraic inverse; faithfulness alone would not establish this.

Poincaré duality gives an isomorphism

\[
 \delta_H:H\xrightarrow{\sim}H^\vee(-d).
\]

Transfer its realization across \(\phi\) to the Frobenius-equivariant
isomorphism

\[
 b=(\phi^\vee(-d))^{-1}\,R(\delta_H)\,\phi^{-1}:
       W\longrightarrow W^\vee(-d).               \tag{CC-XC.3}
\]

Both \(M\) and \(M^\vee(-d)\) belong to \(\mathcal A_{q,E}\). Its full
realization theorem, AQ.15, therefore lifts \(b\) to an algebraic
morphism \(\delta_M:M\to M^\vee(-d)\). This step requires fullness only
inside the already established abelian sector.

Now form the algebraic morphism

\[
 g=\delta_H^{-1}\circ f^\vee(-d)\circ\delta_M:
          M\longrightarrow H.                    \tag{CC-XC.4}
\]

Using (CC-XC.3), its realization is

\[
 R(g)=R(\delta_H)^{-1}\,
       \phi^\vee(-d)\,
       (\phi^\vee(-d))^{-1}\,
       R(\delta_H)\,\phi^{-1}
       =\phi^{-1}.
\]

Faithfulness of the homological realization now gives \(gf=1_H\) and
\(fg=1_M\). Thus (3) implies (2). This argument needs neither a separate
surjectivity assumption for reverse cross-correspondences nor
numerical--homological equivalence for \(X\).

Condition (2) implies (1), and an isomorphism in (2) has invertible
realization, so it implies (3). If (1) holds, say \(H\simeq M'\) with
\(M'\in\mathcal A_{q,E}\), then \(V\) is semisimple. Hence
\(R(M')\simeq V\simeq V^{\mathrm{ss}}\simeq W\). Full faithfulness
within \(\mathcal A_{q,E}\) lifts this Frobenius-module isomorphism and
its inverse to \(M'\simeq M\); therefore (1) implies (2).

Finally (3) implies (4). Conversely, a nonzero polynomial on a
finite-dimensional vector space over the infinite field \(E\) is nonzero
at some point: after choosing coordinates this follows by induction on
the number of variables, using the finitely many roots of a nonzero
one-variable polynomial. Apply this to \(\det|_I\) to deduce (3) from
(4). This proves all equivalences and the semisimplicity conclusion.
\(\square\)

\textbf{Finite witnesses.} The image \(I\) is spanned by the images of
algebraic cycles and is finite-dimensional. Consequently it has a finite
spanning collection of such images. Equivalently, the condition is
witnessed by finitely many cycles \(z_1,\ldots,z_t\) for which

\[
 \det\!\left(\sum_{j=1}^t x_j\kappa(z_j)\right)
\]

is not the zero polynomial. This is an existence criterion; it supplies
no effective method or complexity bound for finding the cycles.

\subsection{CC-XC03 --- The semisimple version as a finite projected
Tate
condition}\label{cc-xc03-the-semisimple-version-as-a-finite-projected-tate-condition}

\textbf{Corollary.} If Frobenius on \(V\) is semisimple, then the
conditions in CC-XC02 are also equivalent to the surjectivity of
\(\kappa\).

\textbf{Proof.} Under semisimplicity, \(V\simeq W\) as Frobenius
modules. Surjectivity then puts such an isomorphism in \(I\), so CC-XC02
applies. Conversely, if \(H\simeq M\), transport any
Frobenius-equivariant map \(V\to W\) across that isomorphism to an
endomorphism of \(W\). Full faithfulness on \(M\) supplies its algebraic
endomorphism. Composing gives the required morphism \(H\to M\), and
(CC-XC.2) places its realization in \(I\). Thus \(\kappa\) is onto.
\(\square\)

One sufficient hypothesis for this projected surjectivity is the usual
Tate surjectivity in \textbf{codimension \(d\)} on each of the finitely
many \(X\times Y_\nu\):

\[
 \mathrm{CH}^d(X\times Y_\nu)_{\mathbf Q}
       \otimes_{\mathbf Q}\mathbf Q_\ell
 \twoheadrightarrow
 H^{2d}((X\times Y_\nu)_{\bar k},\mathbf Q_\ell(d))^{F=1}.
                                                        \tag{CC-XC.5}
\]

Indeed, Poincaré duality and Künneth identify this invariant space with
the Frobenius-equivariant maps from \(H^*(X)\) to \(H^*(Y_\nu)\). Extend
scalars to \(E\), sum, and project by \(e\). Every map to \(W=eR(T)\) is
obtained by projecting its composition with the inclusion
\(W\hookrightarrow R(T)\), so (CC-XC.5) implies the surjectivity of
(CC-XC.1). This remains a condition on a specified finite family and
specified codimension. It requires no Tate assertion for all varieties
or all codimensions. The projected condition can be weaker than full
surjectivity in (CC-XC.5).

Without semisimplicity, surjectivity onto
\(\operatorname{Hom}_{E[F]}(V,W)\) alone would not guarantee an
invertible map: all such maps could have nontrivial kernel. The
determinant criterion is the exact formulation retaining that issue.

\subsection{CC-XC04 --- Tate and numerical--homological equality after
the
criterion}\label{cc-xc04-tate-and-numericalhomological-equality-after-the-criterion}

\textbf{Corollary.} Suppose the equivalent conditions of CC-XC02 hold
for \(X\). Let \(k'/k\) be any finite extension, let \(n\ge0\), let
\(A/k'\) be an abelian variety and \(S/k'\) finite étale, and put

\[
 Z=X_{k'}^{\,n}\times A\times S.
\]

For the fixed prime \(\ell\), Frobenius on
\(H^*(Z_{\bar k},\mathbf Q_\ell)\) is semisimple; the rational Tate
cycle map is surjective in every codimension; and numerical equivalence
and \(\ell\)-adic homological equivalence agree on rational algebraic
cycles of \(Z\).

\textbf{Complete proof.} If \(Z\) is empty, all assertions are vacuous,
so assume it is nonempty. Base change of algebraic correspondences
preserves the isomorphism \(h_{\ell,E}(X)\simeq M\) proved in CC-XC02.
Its base change remains an abelian--Artin motive over \(k'\). Tensoring,
and then multiplying by \(h(A)\) and \(h(S)\), gives an isomorphism of
\(h_{\ell,E}(Z)\) with an object of \(\mathcal A_{q',E}\), where
\(q'=\#k'\). AQ.15 applies over every finite field. It proves
semisimplicity after extension to \(E\), and semisimplicity over
\(\mathbf Q_\ell\) follows because a characteristic-zero operator is
semisimple if and only if its extension to an algebraic closure is.

For codimension \(r\), fullness on the abelian sector identifies
algebraic morphisms \(\mathbf 1\to h_{\ell,E}(Z)(r)\) with the
Frobenius-invariant vectors of its realization. The weight decomposition
shows that only cohomological degree \(2r\) contributes. Consequently

\[
 \operatorname{im}\!\left(
 \mathrm{CH}^r(Z)_{\mathbf Q}\otimes_{\mathbf Q}E
 \longrightarrow H^{2r}(Z_{\bar k},E(r))
 \right)
 =H^{2r}(Z_{\bar k},E(r))^{F=1}.                 \tag{CC-XC.6}
\]

Here is the descent to the usual \(\mathbf Q_\ell\)-valued Tate
statement. Let \(C_\ell\) denote the image of
\(\mathrm{CH}^r(Z)_{\mathbf Q}\otimes_{\mathbf Q}\mathbf Q_\ell\) in
\(U=H^{2r}(Z_{\bar k},\mathbf Q_\ell(r))^{F=1}\). Flat scalar extension
commutes both with that image and with the kernel defining \(F=1\)
invariants. Equation (CC-XC.6) therefore says
\(C_\ell\otimes_{\mathbf Q_\ell}E=U\otimes_{\mathbf Q_\ell}E\). Faithful
flatness of a field extension gives \(U/C_\ell=0\). This argument uses
the image subspace; it does not presume injectivity after scalar
extension of rational homological classes.

Write \(D=\dim Z\). The Poincaré pairing

\[
 H^{2r}(Z_{\bar k},\mathbf Q_\ell(r))
 \times
 H^{2D-2r}(Z_{\bar k},\mathbf Q_\ell(D-r))
 \longrightarrow\mathbf Q_\ell
\]

is perfect and Frobenius invariant. Its restriction to the two invariant
subspaces is perfect: after diagonalizing semisimple Frobenius over
\(E\), eigenvalues \(\lambda,\mu\) can pair nontrivially only when
\(\lambda\mu=1\); the \(1\)-eigenspace is therefore paired perfectly
with the other \(1\)-eigenspace. Nondegeneracy descends to
\(\mathbf Q_\ell\).

If a rational cycle \(z\) is numerically trivial but
\(a=\operatorname{cl}_\ell(z)\ne0\), this invariant perfect pairing
provides an invariant complementary class \(b\) with
\(\langle a,b\rangle\ne0\). The already proved rational Tate
surjectivity expresses \(b\) as a finite \(\mathbf Q_\ell\)-linear
combination of classes of rational complementary cycles. Numerical
triviality makes every corresponding intersection degree zero,
contradicting the chosen pairing. Thus numerical triviality implies
homological triviality. The converse follows directly from compatibility
of the cycle class with intersection degrees. This proves the assertion
on rational cycles without assuming numerical realization at the start.
\(\square\)

The conclusion is for the specified prime \(\ell\) and this class of
products and base extensions. It does not assert the determinant
condition for arbitrary \(X\), or a comparison with another prime
without an additional argument.

\subsection{Relation to Milne and the remaining
question}\label{relation-to-milne-and-the-remaining-question}

Milne (1994), \textbf{Proposition 1.15 and Corollary 1.16},
author-manuscript pp.~11--12, derive semisimplicity,
numerical--homological equality and full realization from the full Tate
conjecture. \textbf{Proposition 2.6 and Remark 2.7}, p.~15, then give
the global abelian/Artin-generation conclusion under that hypothesis.
Those inspected statements do not themselves state the particular
one-sided finite-target criterion above. Here we prove that criterion
using only the already settled abelian-sector realization and the
specified cycle condition for \(X\). No historical priority claim is
made.

For a particular \(X\), the unresolved requirement has therefore been
made concrete: construct finitely many cross-correspondences whose
projected realization has nonzero determinant. Pure numerical Frobenius
eigenvalues alone do not establish this. The adopted RH and Dirichlet
GRH do not provide these cycles or prove the determinant condition.

\subsection{Sources and locators}\label{sources-and-locators}

\begin{itemize}
\tightlist
\item
  Deligne, P. (1974). La conjecture de Weil: I. \emph{Publications
  Mathématiques de l'IHÉS, 43}, 273--307. \textbf{Theorem 1.6, p.~276}.
  \url{https://www.numdam.org/item/PMIHES_1974__43__273_0.pdf}.
\item
  Honda, T. (1968). Isogeny classes of abelian varieties over finite
  fields. \emph{Journal of the Mathematical Society of Japan, 20}(1--2),
  83--95. \url{https://doi.org/10.2969/jmsj/02010083}. Main Theorem, p.~83;
  primary PDF:
  \url{https://people.math.rochester.edu/faculty/doug/otherpapers/Honda-Ab.pdf}.
\item
  Milne, J. S. (1994). Motives over finite fields. In U. Jannsen, S. L.
  Kleiman, \& J.-P. Serre (Eds.), \emph{Motives} (Proceedings of
  Symposia in Pure Mathematics, Vol. 55, Part 1, pp.~401--459). American
  Mathematical Society. \url{https://jmilne.org/math/articles/1994aP.pdf}.
  Author-manuscript pp.~5--6 for finite-field projectors, corrected
  symmetry and duality; pp.~11--12 for Proposition 1.15 and Corollary
  1.16; p.~15 for Proposition 2.6 and Remark 2.7.
\item
  Proved internal inputs: \textbf{CC-AG01} in
  \nolinkurl{research/consequences_criteria_abelian_group.tex};
  \textbf{AQ.15 / AQ-ABELIAN-RECOGNITION} and \textbf{AQ-WEIL-EXIST} in
  \nolinkurl{research/archive_questions.tex}. Their exact all-abelian primary
  inputs and limits are recorded in CC-TATE01--02 and the corresponding
  source ledger.
\end{itemize}

% END consequences_criteria_cross_cycles.tex

% BEGIN consequences_criteria_cycle_certificates.tex
\section{Finite algebraic certificates for abelian-sector membership}

The criterion CC-XC02 was stated using the image of an
\(E\)-valued cycle-class map.  The following results replace its
positive witnesses by rational cycles with a finite verifier using
intersection numbers.  They also show that this membership condition
is independent of the prime \(\ell\ne p\).  The search supplied below
is one-sided: it halts exactly on positive instances.  It supplies
neither a bound on the necessary cycle degrees nor a decision
procedure for negative instances.

Use the rational homological category defined by the kernel of
the \(\ell\)-adic cycle class, its \(E\)-coefficient counterpart
from CC-XC02, and the established abelian--Artin--Tate category
\(\mathcal A_q\).  The latter's realization theorem is the
inspected input AQ.15.  An isomorphism of an arbitrary numerical
motive is not assumed to have a realization before it is
constructed.

\subsection{CC-XC05: a finite determinant verifier from intersection traces}

Let \(k=\mathbb F_q\), with characteristic \(p\), and let
\(X/k\) be smooth projective and geometrically integral of
dimension \(d\).  Write
\[
 P_i(t)=\det(t-F_X\mid H^i(X_{\bar k},\mathbb Q_\ell))
       \in\mathbb Z[t],\qquad b_i=\deg P_i.
                                                        \tag{CC-XC.C10}
\]
These polynomials are independent of \(\ell\ne p\).
They can be computed from the equations of \(X\).
Here \(F_X\) denotes the Frobenius correspondence with
the usual Weil eigenvalues, as in CC-XC01.

A candidate certificate consists of finitely many abelian
varieties or finite étale \(k\)-schemes \(Y_\nu\), integers
\(n_\nu\), and rational cycles
\[
 \begin{split}
 f_\nu&\in\mathrm{CH}^{d+n_\nu}(X\times Y_\nu)_{\mathbb Q},\\
 g_\nu&\in\mathrm{CH}^{\dim Y_\nu-n_\nu}
                                  (Y_\nu\times X)_{\mathbb Q}.
 \end{split}                                            \tag{CC-XC.C11}
\]
The cycles are supplied as finite rational linear combinations
of subvarieties, rather than as unspecified equivalence classes.
A negative codimension supplies only the zero cycle.
They define morphisms
\[
 f:h(X)\longrightarrow T=\bigoplus_\nu h(Y_\nu)(n_\nu),
 \qquad g:T\longrightarrow h(X),\qquad a=gf.
                                                        \tag{CC-XC.C12}
\]
The composition \(a\) has degree zero on \(X\).

There is a finite algorithm computing a rational number
\(\mathfrak D_X(f,g)\) such that, for every \(\ell\ne p\),
\[
 \mathfrak D_X(f,g)
       =\prod_i\det\!\left(a\mid
                         H^i(X_{\bar k},\mathbb Q_\ell)\right).
                                                        \tag{CC-XC.C13}
\]
Consequently its nonvanishing is independent of \(\ell\).
The algorithm uses the zeta function of \(X\), rational
polynomial arithmetic and finitely many intersection numbers
of explicitly described cycles.  It needs neither a test for
homological equivalence nor an algorithm for computing an
\(\ell\)-adic cycle-class map.

\paragraph{Complete proof and verifier.}
Poonen, Testa, and van Luijk (2015), Proposition~8.1,
author-manuscript p.~10, computes \(Z(X,t)\) unconditionally
over a finite field.  Their proof uses Katz's computable
upper bound for the sum of Betti numbers, counts points over
finitely many extensions and reconstructs the resulting
rational function.  Deligne's Theorem~1.6 (1974, p.~276)
puts the roots in distinct absolute-value classes according
to \(i\).  Thus the rational zeta function determines each
\(P_i\), not merely their alternating product.  This is the
same separation used in Poonen--Testa--van Luijk,
Proposition~8.2, p.~10.  Factoring over \(\mathbb Q\) and
distinguishing the finitely many possible root absolute
values is a finite calculation.

Discard indices with \(b_i=0\) for the moment.
The polynomials \(P_i\) are pairwise coprime: roots in
different degrees have absolute values \(q^{i/2}\) and
therefore cannot coincide.  By the Chinese remainder
algorithm choose \(u_i(t)\in\mathbb Q[t]\) with
\[
 u_i\equiv1\pmod {P_i},\qquad
 u_i\equiv0\pmod {P_j}\quad(j\ne i).
 \]
Define the explicit rational correspondence
\[
                              \pi_i=u_i(F_X).
                                                        \tag{CC-XC.C14}
\]
Cayley--Hamilton for the full characteristic polynomial
\(P_j\), including its multiplicities, shows that
\(\pi_i\) acts as the identity on \(H^i\) and as zero
on every other degree.  This argument permits arbitrary
Jordan blocks.  Neither Frobenius semisimplicity nor
idempotence modulo rational equivalence is being assumed.
The existence of these common algebraic Künneth projectors
is the classical Katz--Messing Theorem~2 (1974, p.~76).
The explicit CRT calculation here fixes the rational
cycles used by the verifier.

We record the trace calculation including its sign.
For any rational degree-zero correspondence \(c\) on \(X\),
Poincaré duality, Künneth and the compatibility of cycle
classes with products and degrees give
\[
 \deg(\Delta_X\cdot c)
                  =\sum_j(-1)^j\operatorname{Tr}(c\mid H^j(X)).
                                                        \tag{CC-XC.C15}
\]
One can see the formula directly.  For each \(j\), choose
a basis \(e_{j,s}\) of \(H^j\), with dual classes
\(e^\vee_{j,s}\in H^{2d-j}(d)\) normalized by
\(\int_X e_{j,s}\cup e^\vee_{j,t}=\delta_{st}\).
Under Künneth, the correspondence class is
\(\sum_{j,s}e^\vee_{j,s}\otimes c(e_{j,s})\).
Pulling back to the diagonal and integrating gives
\[
 \sum_{j,s}\int_X e^\vee_{j,s}\cup c(e_{j,s})
                    =\sum_j(-1)^j\operatorname{Tr}(c\mid H^j).
 \]
The last sign is graded commutativity.  Degree compatibility
identifies the integral with the intersection number on the
left of (CC-XC.C15).  The underlying diagonal and degree
identities, with these signs, are also given in Stacks
Project, Lemma~45.9.7, Tag~0FHH, and Lemma~45.10.3,
Tag~0FHR.  Thus the calculation applies to rational cycles
and requires no isolated-fixed-point or transversality
hypothesis on the original correspondence.

For \(1\le r\le b_i\), compute
\[
 s_{i,r}=(-1)^i\deg\bigl(\Delta_X\cdot(\pi_i\circ a^r)\bigr)
                         \in\mathbb Q.                  \tag{CC-XC.C16}
\]
Equation (CC-XC.C15) and the projector action show that these
are exactly \(\operatorname{Tr}(a_i^r)\), where \(a_i\)
denotes the action of \(a\) on \(H^i\).
Starting with \(e_{i,0}=1\), use Newton's recursion
\[
 m e_{i,m}=
      \sum_{r=1}^{m}(-1)^{r-1}e_{i,m-r}s_{i,r},
                 \qquad 1\le m\le b_i.                 \tag{CC-XC.C17}
\]
The result is \(e_{i,b_i}=\det(a_i)\).
For \(b_i=0\), set this determinant equal to one.
Their product is the required \(\mathfrak D_X(f,g)\).
This proves (CC-XC.C13), including independence of \(\ell\).

Every operation in this verifier is finite and algebraic.
Composition of correspondences uses pullback, intersection
and proper pushforward on products.  Intersection numbers
are computable by Poonen--Testa--van Luijk, Lemma~8.7,
p.~11: move a cycle by an enumerated rational equivalence
to a proper transverse intersection and compute its degree.
Their Remark~8.18, p.~14, records the same effective
intersection operation in the Chow ring.  Applied successively
on the smooth projective products, this computes a
representative for each required composition, or equivalently
the degree of the corresponding iterated intersection.
Rational coefficients are handled by clearing denominators
and bilinearity.  Pullbacks along product projections and
proper pushforwards of the resulting cycles are computable
from their defining ideals and generic degrees.
If a product has several geometric components, these
operations can instead be made after a finite extension
splitting the components, computing degrees over that extension.  Such a degree equals
the original degree over the ground field.  If cycles are
instead returned by proper trace, that pushforward multiplies
degrees by the extension degree, which must then be divided
out.  These are equivalent ways to perform the computation.
No bound on the time required for the moving search is
claimed, but the cited lemma proves its termination.
\(\square\)

\subsection{CC-XC06: rational certificates, prime independence and an
actual summand}

For \(X\) as above, the following conditions are equivalent.
\begin{enumerate}
\item For some \(\ell\ne p\), the \(E=\overline{\mathbb Q}_\ell\)
homological motive \(h_{\ell,E}(X)\) satisfies the equivalent
conditions of CC-XC02.
\item There is a rational candidate (CC-XC.C11) with
\(\mathfrak D_X(f,g)\ne0\).
\item There is such a candidate and rational correspondence
formulas defining a splitting
\[
 h_{\ell,\mathbb Q}(X)\longrightarrow T
                  \longrightarrow h_{\ell,\mathbb Q}(X)
 \]
whose composite is the identity, simultaneously for every
\(\ell\ne p\).
\end{enumerate}
When they hold, one rational numerical motive
\(M\in\mathcal A_q\), given by a rational idempotent on \(T\),
is isomorphic to \(h_{\mathrm{num}}(X)\); the same cycle
formulas give its homological realization as \(h(X)\) for
every \(\ell\ne p\).  In particular CC-XC02 is independent
of the initially selected prime and coefficient extension.
This is a statement modulo homological or numerical
equivalence, not a lift of the idempotent to Chow motives.

\paragraph{Complete proof.}
Suppose (1).  An object of \(\mathcal A_{q,E}\) is a summand
of a finite \(T_E\) as in (CC-XC.C12).  Compose an isomorphism
from \(h_{\ell,E}(X)\) with its inclusion into \(T_E\), and
compose the corresponding projection with the inverse
isomorphism.  This gives \(E\)-linear algebraic
correspondences \(f_E,g_E\) with
\[
                                  g_E f_E=1
                                                        \tag{CC-XC.C18}
\]
in the \(E\)-homological category.
Each is a finite \(E\)-linear combination of rational
cycles.  Choose finite lists \(f_1,\ldots,f_u\) and
\(g_1,\ldots,g_v\) of such rational cycles, preserving
the finitely many components and codimensions in
(CC-XC.C11), and form
\[
 f(\mathbf x)=\sum_jx_jf_j,\qquad
 g(\mathbf y)=\sum_ky_kg_k.
 \]
Composition and the traces in (CC-XC.C16) are polynomial
in the variables.  Newton's recursion uses only rational
coefficients.  Therefore
\[
        D(\mathbf x,\mathbf y)
               =\mathfrak D_X(f(\mathbf x),g(\mathbf y))
                         \in\mathbb Q[\mathbf x,\mathbf y].
                                                        \tag{CC-XC.C19}
\]
At the \(E\)-coefficients of (CC-XC.C18) every determinant
equals one, so \(D\) is not the zero polynomial.
A nonzero polynomial over the infinite field \(\mathbb Q\)
cannot vanish on every rational tuple: induction on the
number of variables proves this from the finite number of
roots of a one-variable polynomial.  Choose a rational
tuple at which \(D\ne0\).  It supplies (2).
This argument concerns actual finite cycle combinations
and their intersection traces; it does not assume
injectivity of a scalar-extension map on rational
homological Hom spaces.

Suppose (2), and let \(a=gf\).  The recursion in CC-XC05
also computes the rational characteristic polynomial
\[
 \chi(t)=\prod_i\det(t-a_i)
        =t^N+c_1t^{N-1}+\cdots+c_N\in\mathbb Q[t],
        \qquad N=\sum_i b_i,\quad c_N\ne0.
 \]
The same polynomial is obtained for every \(\ell\ne p\).
Set
\[
 v(t)=-c_N^{-1}
        (t^{N-1}+c_1t^{N-2}+\cdots+c_{N-1}).
 \]
Then \(t v(t)=1-c_N^{-1}\chi(t)\).
Cayley--Hamilton gives
\[
                       a\,v(a)=v(a)\,a=1
                                                        \tag{CC-XC.C20}
\]
after every such realization, hence in each rational
homological category.  Define the rational cycle
formulas
\[
                         j=f,\qquad r=v(a)g,\qquad e=jr.
                                                        \tag{CC-XC.C21}
\]
They satisfy \(rj=1\) and \(e^2=e\) in all those categories,
and give \(h(X)\simeq(T,e)\).  This proves (3).
The same identities hold modulo numerical equivalence,
because homologically trivial rational cycles are
numerically trivial by compatibility with intersection
degrees.  Thus the rational numerical projector \(e\)
on the abelian--Artin--Tate object \(T\) defines
\(M=(T,e)\in\mathcal A_q\) and the same \(j,r\)
give \(h_{\mathrm{num}}(X)\simeq M\).
Within this established sector numerical and homological
equivalence agree, so its realization is precisely the
already exhibited summand for every \(\ell\).
This constructs, rather than presupposes, the realization
needed for this particular numerical object.

Finally, (3) makes \(h_{\ell,E}(X)\) an abelian-sector
summand for every \(\ell\), and therefore implies (1)
by CC-XC02.  All implications are proved. \(\square\)

\paragraph{All-prime consequence for products.}
For a certified \(X\), every finite base extension and
every product \(X^m\times A\times S\), with \(A\) abelian
and \(S\) finite étale, has semisimple Frobenius,
surjective rational Tate cycle maps in every codimension,
and numerical equivalence equal to \(\ell\)-adic
homological equivalence for every \(\ell\ne p\).
Indeed, apply the same base-change and tensor constructions
to (CC-XC.C21), and apply AQ.15 in its abelian sector.
The scalar descent and invariant complementary-pairing
argument are exactly those fully proved in CC-XC04,
now applicable to every prime by the common certificate.
No crystalline assertion is added here.

\subsection{CC-XC07: an exact one-sided search for certificates}

There is an algorithm on equations for \(X/k\) which halts
if and only if the equivalent conditions of CC-XC06 hold.
When it halts it returns the finitely many varieties,
rational cycles and nonzero rational determinant of a
certificate, together with the splitting formulas
(CC-XC.C20)--(CC-XC.C21).  Each returned certificate can
be checked by an independent finite computation.
If membership fails, the search runs forever and makes
no negative declaration.

\paragraph{Complete proof.}
Finite projective models and cycles over a finite field
have enumerable encodings by homogeneous polynomial
ideals and finite integer or rational coefficient lists;
this is the encoding used in Poonen--Testa--van Luijk,
\S7.1, pp.~6--7.  Enumerate projective models with explicit
smooth, geometrically connected commutative group
structures, including identity, inverse and multiplication
on finite affine covers.  The defining equations,
smoothness, geometric connectedness, covers and group
identities admit finite algebraic verification.  The
verified models are abelian varieties, and every abelian
variety over \(k\) has one such finite presentation.
For the connectedness test one may compute
\(H^0(Y,\mathcal O_Y)\): for a nonempty smooth projective
\(Y\) it equals \(k\) precisely when \(Y\) is geometrically
connected.  This finite-dimensional algebra, smoothness
and equalities of the specified morphisms can be computed
from projective and affine equations.
Also enumerate \(\operatorname{Spec}\mathbb F_{q^s}\)
and finite disjoint unions thereof, all finite lists of
the verified models, and all integers \(n_\nu\).

For each such \(T\), enumerate all finite combinations of
subvarieties in the two codimensions (CC-XC.C11).
Dimension and containment are decidable by polynomial
ideal computation.  Apply the terminating verifier
CC-XC05 to every resulting candidate, scheduling the
enumerations and calculations so every candidate
eventually receives its finite required computation.
For example, dovetail all the finitely encoded tasks.
Accept and return the first nonzero determinant.
Only exact rational equalities and inequalities enter
the acceptance test.

If the search accepts, CC-XC06 supplies the claimed
splitting and membership, so there are no false positives.
If membership holds, CC-XC06 supplies a finite rational
certificate.  Its models and cycles eventually occur
in the enumeration, and its verifier terminates with
a nonzero determinant.  Thus the search eventually
halts.  Multiplying each of \(f,g\) by a common
nonzero integer scales \(a\) by a nonzero integer
and preserves nonvanishing, so even integral cycle
coefficient lists suffice for the search.

The twists cause no hidden infinite operation within
one candidate.  The displayed codimensions already
account for them: \(f_\nu\) shifts cohomological degree
by \(2n_\nu\), \(g_\nu\) shifts it back, and \(a=gf\)
has degree zero.  Alternatively, choose one integer
\(m\ge\max(0,n_\nu)\), tensor both source and target by
\(\mathbb Q(-m)\), and use the usual Lefschetz projector
on each of finitely many copies of \(\mathbb P^1\).
This realizes every remaining twist \(n_\nu-m\le0\)
and the source twist \(-m\) as summands of ordinary
effective motives by finitely many cycle operations.
For the unshifted \(h(X)\) considered here, one may
even restrict the search to \(n_\nu=0\): CC-XC01
supplies an effective finite abelian/Artin target
whenever CC-XC02 holds.  Hence this restriction does
not change which inputs cause the search to halt.
\(\square\)

\paragraph{General smooth projective schemes.}
The statements extend to a nonempty smooth projective
scheme by treating its finitely many connected components;
components of different dimensions have their own
correspondence codimensions.  A connected component
has a finite field of constants \(k'/k\) over which
it is geometrically integral.  The zeta computation
reduces to the preceding one using
\(Z_{X/k}(t)=Z_{X/k'}(t^{[k':k]})\), and degrees and
cycles are computed componentwise.  Equivalently,
the point-count reconstruction proof of Proposition~8.1
uses the same Betti bound without requiring geometric
integrality.  Projectors are still obtained from the
common degree-weight factors, and the trace proof
uses the full direct sum of geometric components.
The empty scheme is certified by the zero object,
with determinant of its zero-dimensional realization
equal to one.

\subsection{Exact literature scope and the remaining gate}

The cited algorithmic ingredients are unconditional
for a smooth projective variety over a finite field.
In contrast, Poonen--Testa--van Luijk,
Theorem~8.36(b), author-manuscript p.~19, assumes
their \(E^r(X,\ell)\) when it identifies halting with
the Tate conjecture in codimension \(r\).  We do not
remove that hypothesis from their theorem or use it
here.  Our acceptance test asks whether an explicitly
given cross-correspondence has a nonzero determinant,
which is determined by its finite trace data.

Thus the remaining arbitrary-variety problem is
whether the enumerated search always finds a certificate.
Neither its universal termination nor a cycle-degree
bound is proved.  A numerical isomorphism supplied
without homologous splitting cycles is also not,
by itself, asserted to satisfy this test.  The result
provides a finite positive verifier, rational coefficient
descent and prime independence for every accepted
variety.  It does not settle the full Tate conjecture
or the equivalent arbitrary-motive effectivity gate.

\subsection{Primary references}

\begin{itemize}
\item Deligne, P. (1974). La conjecture de Weil: I.
\emph{Publications Mathématiques de l'IHÉS, 43},
273--307. Theorem~1.6, p.~276.
\url{https://www.numdam.org/item/PMIHES_1974__43__273_0.pdf}.
\item Katz, N.~M., \& Messing, W. (1974).
Some consequences of the Riemann hypothesis for
varieties over finite fields.
\emph{Inventiones Mathematicae, 23}, 73--77.
\url{https://doi.org/10.1007/BF01405203}.
Theorem~2, p.~76, common algebraic Künneth projectors;
primary author-hosted scan:
\url{https://web.math.princeton.edu/~nmk/old/katzmessing.pdf}.
\item Poonen, B., Testa, D., \& van Luijk, R. (2015).
Computing Néron--Severi groups and cycle class groups.
\emph{Compositio Mathematica, 151}(4), 713--734.
\url{https://doi.org/10.1112/S0010437X14007878}.
Inspected author manuscript:
\url{https://math.mit.edu/~poonen/papers/compute_ns.pdf}.
Exact author-page locators: \S7.1, pp.~6--7;
Propositions~8.1--8.2, p.~10; Lemma~8.7,
p.~11; Remark~8.18, p.~14; and the differently
hypothesized Theorem~8.36, p.~19.
\item The Stacks Project Authors. (2026).
\emph{The Stacks Project}.
Lemma~45.9.7, Tag~0FHH,
\url{https://stacks.math.columbia.edu/tag/0FHH};
Lemma~45.10.3, Tag~0FHR,
\url{https://stacks.math.columbia.edu/tag/0FHR}.
The inspected Weil-cohomology chapter PDF identifies
version ed88ff78, compiled July~14, 2026; the displayed
lemma statements and proofs were directly inspected.
\end{itemize}

% END consequences_criteria_cycle_certificates.tex

% BEGIN consequences_criteria_fermat_certificates.tex
\section{Geometric constructions that force the finite certificate
search to
terminate}\label{geometric-constructions-that-force-the-finite-certificate-search-to-terminate}

\textbf{Scope and inputs.} Fix a finite field \(k=\mathbf F_q\) of
characteristic \(p\). Call a smooth projective \(k\)-scheme
\textbf{certified} if it satisfies the equivalent conditions of CC-XC06:
there are rational algebraic correspondences giving its cohomological
motive as a direct summand of a finite sum of Tate twists of abelian
varieties and finite étale schemes. The same correspondences can be
chosen to give the splitting for every \(\ell\ne p\), and also in the
rational numerical category. Empty schemes and finite disjoint unions
are allowed.

This section uses the abelian-sector fullness and numerical--homological
comparison already isolated in AQ.15 and CC-AG01, and the proved
rational certificate theorem CC-XC05--07. It does not assume the Tate
conjecture for an arbitrary variety. The geometric input for Fermat
varieties is the classical construction of Shioda and Katsura (1979,
Lemmas 1.1--1.3, pp.~98--100; Lemma 2.1, p.~103). The statements below
are applications and complete deductions from these inputs, with no
novelty claim. Literature cutoff: 20 November 2026; checks: 21 November
2026 UTC.

\subsection{CC-XC08 --- Closure operations, including descent of
membership}\label{cc-xc08-closure-operations-including-descent-of-membership}

The following assertions hold.

\begin{enumerate}
\def\labelenumi{\arabic{enumi}.}
\tightlist
\item
  Abelian varieties, finite étale schemes, and smooth projective curves
  are certified. Finite products and finite disjoint unions of certified
  schemes are certified.
\item
  If \(f:Y\to X\) is a proper, dominant, generically finite morphism
  between smooth projective geometrically integral schemes, and \(Y\) is
  certified, then \(X\) is certified. Separability is not needed for
  this assertion.
\item
  If \(P\) is smooth projective and \(Z\hookrightarrow P\) is a smooth
  closed subscheme of pure codimension two, then certification of \(P\)
  and \(Z\) implies certification of \(\operatorname{Bl}_ZP\).
\item
  For every finite extension \(k'/k\), a smooth projective \(X/k\) is
  certified if and only if \(X_{k'}\) is certified over \(k'\).
\end{enumerate}

Assertion 4 concerns membership in the whole abelian--Artin category
with a freely chosen target. It does not say that every specified
rational object over \(k'\), or over \(\bar k\), descends to a specified
object over \(k\). The orbit and Schur-index obstructions in CC-AG09--13
remain applicable to that different descent problem.

\textbf{Complete proof.}

\textbf{Finite sums and products.} For one auxiliary prime
\(\ell\ne p\), write \(\mathcal H_{\ell,E}(k)\) for the homological
category used in CC-XC01, with \(E=\overline{\mathbf Q}_\ell\). CC-XC06
identifies certification with membership in its fully embedded
abelian--Artin subcategory \(\mathcal A_{q,E}\). This subcategory is
closed under finite sums, tensor products, Tate twists, duals and direct
summands. Künneth identifies a product with the tensor product, so
products and disjoint unions preserve membership and hence
certification. Abelian varieties and finite étale schemes are
generators.

This use of tensor closure matters: a product of an abelian variety and
a nontrivial finite étale scheme need not itself be an abelian variety.
Membership in the tensor category followed by CC-XC06 supplies the
permitted finite certificate target.

\textbf{Generically finite images.} Put \(\delta=[k(Y):k(X)]>0\).
Pullback and proper pushforward are algebraic correspondences and
satisfy

\[
 f_*f^*(a)=f_*(1)\smile a=\delta a
 \qquad(a\in H^*(X_{\bar k},E)).
 \tag{CC-XC.F11}
\]

Indeed, \(f_*[Y]=\delta[X]\), and the displayed equality is the
projection formula. Thus \(f^*\) and \(\delta^{-1}f_*\) split
\(h_{\ell,E}(X)\) off \(h_{\ell,E}(Y)\). Since \(\mathcal A_{q,E}\) is
closed under summands, \(X\) belongs to it. Apply CC-XC06. The identity
depends on the full degree, including an inseparable degree if present;
division by this positive integer is allowed with rational coefficients.
This proves the required trace splitting directly from cycle pushforward
and the projection formula.

\textbf{Codimension-two blowups.} Write

\[
 B=\operatorname{Bl}_ZP,\quad
 \beta:B\to P,\quad
 j:E_Z\hookrightarrow B,\quad
 \rho:E_Z=\mathbf P(N_{Z/P})\to Z.
\]

Here \(\mathbf P(N_{Z/P})\) parametrizes normal lines. Use
\(\xi=c_1(\mathcal O_{E_Z}(1))\); the normal line of \(E_Z\) in \(B\) is
\(\mathcal O_{E_Z}(-1)\). Consider the algebraic cohomological maps

\[
 \begin{aligned}
 \Phi:\ H^i(P)\oplus H^{i-2}(Z)(-1)&\longrightarrow H^i(B),
 &\Phi(u,v)&=\beta^*u+j_*\rho^*v,\\
 \Psi:\ H^i(B)&\longrightarrow H^i(P)\oplus H^{i-2}(Z)(-1),
 &\Psi(w)&=(\beta_*w,-\rho_*j^*w).
 \end{aligned}
 \tag{CC-XC.F12}
\]

Every map here is induced by a graph, its transpose, or their
composition. The twists in the second summand are the usual étale Tate
twists.

We check \(\Psi\Phi=1\). First, \(\beta_*\beta^*=1\) by (CC-XC.F11) with
degree one. Next, \(\rho_*\rho^*=0\): the fibres of \(\rho\) have
dimension one, so \(\rho_*1\in H^{-2}(Z)(-1)=0\), and the projection
formula applies. If \(i_Z:Z\hookrightarrow P\), then
\(\beta j=i_Z\rho\). It follows that

\[
 \beta_*j_*\rho^*=(i_Z)_*\rho_*\rho^*=0,\qquad
 \rho_*j^*\beta^*=\rho_*\rho^*i_Z^*=0.
\]

Finally, self-intersection of the divisor \(E_Z\) gives
\(j^*j_*v=-\xi\smile v\), while integration over the projective-line
fibres gives \(\rho_*(\xi\smile\rho^*v)=v\). Thus the lower-right entry
of \(\Psi\Phi\) is also the identity, with the minus sign in \(\Psi\)
cancelling the normal-bundle sign.

The blowup cohomology formula, Shioda--Katsura, Lemma 2.1, p.~103,
equation (2.1), gives

\[
 \dim H^i(B)=\dim H^i(P)+\dim H^{i-2}(Z).
\]

Consequently the left inverse just verified is the inverse, in every
degree. This proves the algebraic isomorphism

\[
 h_{\ell,E}(B)\simeq h_{\ell,E}(P)\oplus h_{\ell,E}(Z)(-1).
 \tag{CC-XC.F13}
\]

It is Frobenius-equivariant because all the maps are defined over \(k\).
Hence \(B\) is certified whenever \(P\) and \(Z\) are. Shioda--Katsura
write their positive twist as tensoring with \(H^2(\mathbf P^1)\); their
twist \((1)\) in this formula is therefore our \((-1)\). Disconnected
smooth centres cause no change: both their cohomology and the formulas
split by components.

\textbf{Finite-field descent.} The forward implication in assertion 4
follows by base change of the rational splitting cycles. For the
converse, suppose \(X_{k'}\) is certified and put \(e=[k':k]\). By
CC-XC06 it has a splitting into a finite sum

\[
 T'=\bigoplus_\nu h(Y_\nu)(n_\nu),
 \tag{CC-XC.F14}
\]

with rational splitting correspondences, where each \(Y_\nu\) is an
abelian variety or a finite étale scheme over \(k'\).

There is an additive induction operation on these homological
correspondences: regard a smooth projective \(k'\)-scheme as a
\(k\)-scheme, and extend a correspondence on \(U\times_{k'}V\) by zero
along the open-and-closed immersion into \(U\times_kV\). It preserves
identities and composition. To verify this directly, extend scalars to
\(\bar k\). The resulting schemes are disjoint unions indexed by the
embeddings \(k'\hookrightarrow\bar k\). An induced correspondence is
block diagonal, with its conjugate in each matching block. Composition
and identities are therefore block diagonal with exactly the asserted
entries. This also proves that a zero cohomological action remains zero
after induction. Thus the given splitting of \(X_{k'}\) induces a
splitting of its underlying \(k\)-scheme into the induced summands of
(CC-XC.F14).

If \(Y_\nu\) is finite étale over \(k'\), its underlying \(k\)-scheme is
finite étale. If \(Y_\nu=A\) is an abelian variety over \(k'\), let

\[
 B=\operatorname{Res}_{k'/k}A.
\]

Because the extension is finite separable, \(B\) is an abelian variety
over \(k\), and

\[
 B_{k'}\simeq
 \prod_{\sigma\in\operatorname{Gal}(k'/k)} A^\sigma .
 \tag{CC-XC.F15}
\]

These are the exact restriction-of-scalars properties recalled and
justified by Milne (1972, §1, pp.~177--178). Projection onto the
identity factor, and insertion into that factor with zero in all other
factors, are \(k'\)-morphisms whose composite on \(A\) is the identity.
Their graph correspondences show, after viewing the schemes over \(k\),
that

\[
 h(A\text{ viewed over }k)
 \quad\text{is a summand of}\quad
 h(B_{k'}\text{ viewed over }k)
 =h(B)\otimes h(\operatorname{Spec}k').
 \tag{CC-XC.F16}
\]

The object on the right is explicitly abelian \textbf{times Artin}. It
belongs to the tensor-closed category \(\mathcal A_{q,E}\); we have not
identified it with the motive of a single abelian variety. Tate twists
preserve this conclusion. Therefore induction of every summand in
(CC-XC.F14), and then the underlying \(k\)-scheme of \(X_{k'}\), belong
to \(\mathcal A_{q,E}\).

Finally the finite étale projection

\[
 \tau:X_{k'}\text{ viewed over }k\longrightarrow X
\]

has degree \(e\), and \(\tau_*\tau^*=e\,1\). Pullback and \(e^{-1}\)
times trace split off \(h(X)\). Hence \(h_{\ell,E}(X)\) is in the
abelian--Artin category, and CC-XC06 supplies a rational certificate
with its specified allowed target over \(k\). This argument descends
membership, without asserting descent of any prescribed simple summand.

\textbf{Curves.} By this finite-extension assertion, it suffices first
to treat a smooth projective geometrically connected curve \(C/k\) with
a rational point \(P\): every such curve acquires a point over a finite
extension. Put \(g=\dim J(C)\). The point supplies algebraic projectors
onto \(H^0(C)\) and \(H^2(C)\), whose motives are \(\mathbf Q\) and
\(\mathbf Q(-1)\). If \(g=0\), these exhaust the cohomology, and the
assertion follows directly.

For \(g>0\), the Abel--Jacobi map \(a:C\to J=J(C)\) induces an
isomorphism

\[
 a^*:H^1(J_{\bar k},E)\xrightarrow{\sim}H^1(C_{\bar k},E).
 \tag{CC-XC.F17}
\]

Milne (1986, Proposition 9.1, p.~195) states the stronger abelianized
fundamental-group isomorphism over a separably closed field; the
prime-to-\(p\) part needed here is proved on p.~196 using Lemmas
9.2--9.3. Apply it after base change to \(\bar k\). Because \(a\) is
defined over \(k\), (CC-XC.F17) is Frobenius-equivariant.

By Poincaré duality, its transpose is the algebraic pushforward
isomorphism

\[
 a_*:H^1(C_{\bar k},E)
       \xrightarrow{\sim}H^{2g-1}(J_{\bar k},E)(g-1).
 \tag{CC-XC.F18}
\]

Use the algebraic degree projectors on \(J\) from Katz--Messing, Theorem
2, or CC-XC05. Together with the point projectors, the graph of \(a\)
therefore gives a morphism

\[
 h_{\ell,E}(C)\longrightarrow
 \mathbf Q\ \oplus\
 h^{2g-1}_{\ell,E}(J)(g-1)\
 \oplus\ \mathbf Q(-1)
 \tag{CC-XC.F19}
\]

whose realization is an isomorphism. The target belongs to
\(\mathcal A_{q,E}\). The algebraic-inverse construction in the proof of
CC-XC02 applies to this target as well: transfer the Poincaré duality
pairing to the target, lift that pairing using fullness inside the
abelian sector, and transpose the given map. It produces the inverse
without assuming Tate surjectivity for \(C\times J\). Thus \(C\) is
certified by CC-XC06.

For an arbitrary smooth projective curve, a finite extension separates
its finitely many geometric connected components and gives each a point.
The proved finite-union and finite-extension assertions then apply. This
completes all four assertions. \(\square\)

\subsection{CC-XC09 --- Fermat and diagonal hypersurfaces in every
dimension}\label{cc-xc09-fermat-and-diagonal-hypersurfaces-in-every-dimension}

Let \(m\ge2\), let \(p\nmid m\), and for \(r\ge0\) let

\[
 F_m^r:\quad x_0^m+\cdots+x_{r+1}^m=0
       \ \subset\ \mathbf P_k^{\,r+1}.
 \tag{CC-XC.F20}
\]

Then \(F_m^r\) is certified. More generally, every diagonal hypersurface

\[
 D:\quad a_0x_0^m+\cdots+a_{r+1}x_{r+1}^m=0,
 \qquad a_i\in k^\times,
 \tag{CC-XC.F21}
\]

is certified.

Consequently CC-XC07's finite search terminates for each of these
explicit inputs. It returns rational algebraic splitting correspondences
valid for every \(\ell\ne p\), and an actual object of the rational
numerical abelian--Artin category isomorphic to its numerical motive.
The Frobenius semisimplicity, Tate surjectivity, and
numerical--homological consequences of CC-XC06 apply to their products,
products with abelian varieties and finite étale schemes, and finite
extensions of the base field.

\textbf{Complete proof.} The partial derivatives of the equations in
(CC-XC.F20)--(CC-XC.F21) are nonzero scalar multiples of \(x_i^{m-1}\).
They cannot vanish simultaneously at a projective point, so these
hypersurfaces are smooth. In dimension zero they are finite étale. For
\(r\ge1\), the Fermat hypersurface is geometrically integral: two
distinct positive-degree hypersurface components in projective space of
dimension at least two would meet, and their intersection would be
singular in their union; a repeated component is likewise incompatible
with smoothness.

Choose a finite extension \(k'/k\) containing an element \(\epsilon\)
with \(\epsilon^m=-1\). This is possible because the polynomial has
roots in the algebraic closure, each in a finite field. It is enough to
prove certification over \(k'\), by CC-XC08(4). In characteristic two,
\(-1=1\) and \(\epsilon=1\) is already available; no primitive \(2m\)-th
root of unity is required.

Shioda--Katsura's construction can be used over this \(k'\). For
positive integers \(u,v\), its rational map is

\[
 \begin{split}
 F_m^u\times F_m^v&\dashrightarrow F_m^{u+v},\\
 (x,y)&\longmapsto
 (x_0y_{v+1}:\cdots:x_uy_{v+1}:
   \epsilon x_{u+1}y_0:\cdots:\epsilon x_{u+1}y_v).
 \end{split}
 \tag{CC-XC.F22}
\]

The displayed coordinates satisfy the Fermat equation: the first group
of \(m\)-th powers sums to \(-x_{u+1}^m y_{v+1}^m\), and the second to
\(+x_{u+1}^m y_{v+1}^m\). Its base scheme is

\[
 Z=F_m^{u-1}\times F_m^{v-1},
 \qquad x_{u+1}=y_{v+1}=0,
\]

a smooth codimension-two centre in \(F_m^u\times F_m^v\).
Shioda--Katsura, Lemma 1.1, p.~98, proves that the map has degree \(m\);
Lemma 1.2, p.~99, proves that the blowup of this centre resolves the
map:

\[
 \psi:\operatorname{Bl}_{Z}(F_m^u\times F_m^v)
       \longrightarrow F_m^{u+v}.
 \tag{CC-XC.F23}
\]

Lemma 1.3, p.~100, moreover gives a finite étale degree-\(m\) map on the
specified dense open. Only its dominance and degree are needed here.
These lemmas hold in arbitrary characteristic with \(p\nmid m\); all the
displayed formulas, the centre, the blowup and the resolved morphism are
defined over \(k'\).

Proceed by induction on \(r\). The case \(r=0\) is finite étale, and
\(r=1\) is a smooth projective curve, so CC-XC08 applies. For \(r\ge2\),
use \(u=r-1\) and \(v=1\). The source product

\[
 P=F_m^{r-1}\times F_m^1
 \quad\text{and its centre}\quad
 Z=F_m^{r-2}\times F_m^0
\]

are certified by induction and product closure. Hence their
codimension-two blowup is certified. Its degree-\(m\) morphism
(CC-XC.F23) makes \(F_m^r\) a cohomological summand, by the explicit
pullback/trace splitting (CC-XC.F11). Thus \(F_m^r\) is certified over
\(k'\), and finite-field descent proves the statement over \(k\).

For (CC-XC.F21), choose a finite extension containing \(c_i\) with
\(c_i^m=a_i\). The change of coordinates \(y_i=c_i x_i\) identifies
\(D\) over this extension with \(F_m^r\). Finite-field descent again
applies. This proves the theorem.

The assertions about search termination and simultaneous rational
splitting are precisely CC-XC06--07 applied to the membership just
proved. Products and further finite extensions remain certified by
CC-XC08, so the corresponding consequences of CC-XC06 apply there too.
No Chow-projector lifting or arbitrary-variety Tate assertion is being
added. \(\square\)

\textbf{Exact status improvement.} For a specified smooth diagonal
hypersurface with \(p\nmid m\), the existential cross-cycle gate in
CC-XC02 has been discharged by the stated geometric construction and the
existing abelian-sector inputs. The certificate search is therefore
guaranteed to halt on this family. Its runtime and the degrees of the
first certificate are not bounded here. The argument does not extend to
a general smooth hypersurface merely from equality of dimension, degree,
or zeta function.

\subsection{Primary source scope and
bibliography}\label{primary-source-scope-and-bibliography}

Shioda--Katsura work over an algebraically closed field in the
cohomological discussion. Their defining formulas and the resolution in
Lemmas 1.1--1.2 descend to a finite field containing \(\epsilon\), and
their cohomological lemma is applied after geometric base change; the
graph maps in this proof retain their finite-field Frobenius actions.
Remark 1.8, p.~102, imposes a primitive \(2m\)-th root for the larger
quotient diagram. We use only the map and blowup lemmas, and have not
imposed that stronger condition in characteristic two.

Fu--Vial (2020, §5.3) discuss Fermat varieties of abelian type, but that
section explicitly assumes the complex base field. It is not the source
for the finite-field assertion proved here. The
characteristic-independent geometric input used here is Shioda--Katsura.

\begin{itemize}
\tightlist
\item
  Katz, N. M., \& Messing, W. (1974). Some consequences of the Riemann
  hypothesis for varieties over finite fields. \emph{Inventiones
  Mathematicae, 23}, 73--77. \url{https://doi.org/10.1007/BF01405203} . Exact
  locator: Theorem 2, p.~76. Author scan:
  \url{https://web.math.princeton.edu/~nmk/old/katzmessing.pdf}
  .
\item
  Milne, J. S. (1972). On the arithmetic of abelian varieties.
  \emph{Inventiones Mathematicae, 17}, 177--190. Exact locator: §1,
  pp.~177--178, restriction-of-scalars definition, existence, separable
  product decomposition and abelian-variety conclusion.
  \url{https://www.jmilne.org/math/articles/1972a.pdf} .
\item
  Milne, J. S. (1986). Jacobian varieties. In G. Cornell \& J. H.
  Silverman (Eds.), \emph{Arithmetic geometry} (pp.~167--212). Springer.
  Exact locators: Proposition 9.1, p.~195; prime-to-characteristic part
  of the proof and Lemmas 9.2--9.3, p.~196.
  \url{https://www.jmilne.org/math/articles/1986c.pdf} . Author's errata page:
  \url{https://www.jmilne.org/math/articles/1986c.html} ; its listed
  corrections do not alter the cited statement.
\item
  Shioda, T., \& Katsura, T. (1979). On Fermat varieties. \emph{Tohoku
  Mathematical Journal, Second Series, 31}(1), 97--115.
  \url{https://doi.org/10.2748/tmj/1178229881} . Exact locators: Lemma 1.1,
  p.~98; Lemma 1.2, p.~99; Lemma 1.3, p.~100; Lemma 2.1 and equation
  (2.1), p.~103. Primary scan:
  \url{https://www.jstage.jst.go.jp/article/tmj1949/31/1/31_1_97/_pdf} .
\item
  Fu, L., \& Vial, C. (2020). Distinguished cycles on varieties with
  motive of abelian type and the section property. \emph{Journal of
  Algebraic Geometry, 29}, 53--107. Used only for the literature-scope
  comparison, not as a finite-field theorem: conventions, p.~60, and
  §5.3, p.~92. Author manuscript:
  \url{https://irma.math.unistra.fr/~lfu/articles/DistinguishedCycles.pdf}
  .
\end{itemize}

% END consequences_criteria_fermat_certificates.tex

% BEGIN consequences_criteria_separated_hypersurfaces.tex
\section{A larger explicit family: sums of binary
forms}\label{a-larger-explicit-family-sums-of-binary-forms}

\textbf{Scope.} This is a further application of the finite rational
certificate theorem CC-XC05--07 and the geometric closure theorem
CC-XC08, with the same explicitly stated abelian-sector inputs. The
source geometry is Shioda--Katsura (1979, Remark 1.10, p.~102), which
extends their construction from Fermat equations to sums of homogeneous
forms in disjoint variables. A direct verification of the part of that
construction needed here is included. No claim of novelty or of the Tate
conjecture for arbitrary hypersurfaces is made.

\subsection{CC-XC10 --- Separated-variable joins and guaranteed
termination}\label{cc-xc10-separated-variable-joins-and-guaranteed-termination}

Let \(k=\mathbf F_q\) have characteristic \(p\), and let \(m\ge2\) with
\(p\nmid m\).

\textbf{Join criterion.} Let \(r,s\ge1\), and let

\[
 f(x_0,\ldots,x_r),\qquad g(y_0,\ldots,y_s)
\]

be homogeneous polynomials of degree \(m\). Suppose the projective
hypersurfaces

\[
 F=\{f=0\}\subset\mathbf P^r,\qquad
 G=\{g=0\}\subset\mathbf P^s
\]

are smooth. Define

\[
 X_f=\{f(x)+u^m=0\},\quad
 X_g=\{g(y)+v^m=0\},\quad
 Z_{f,g}=\{f(x)+g(y)=0\}.
 \tag{CC-XC.S24}
\]

These are smooth projective hypersurfaces of dimensions \(r,s,r+s\). If
\(F,G,X_f,X_g\) are certified in the sense of CC-XC08, then \(Z_{f,g}\)
is certified.

\textbf{Explicit family.} For every \(t\ge1\), every collection of
nonzero squarefree homogeneous binary forms \(f_1,\ldots,f_t\) of degree
\(m\) over \(k\), and every \(a\in k^\times\), both hypersurfaces

\[
 \begin{aligned}
 B_t(f_1,\ldots,f_t)
   &:\ \sum_{j=1}^t f_j(x_{2j-2},x_{2j-1})=0
       &&\subset\mathbf P^{2t-1},\\
 A_t(f_1,\ldots,f_t;a)
   &:\ \sum_{j=1}^t f_j(x_{2j-2},x_{2j-1})+a z^m=0
       &&\subset\mathbf P^{2t}
 \end{aligned}
 \tag{CC-XC.S25}
\]

are smooth and certified. Here squarefree means no repeated linear
factor after extending to \(\bar k\); equivalently the corresponding
degree-\(m\) subscheme of \(\mathbf P^1\) is finite étale. In particular
the finite certificate search of CC-XC07 is guaranteed to halt on every
input in (CC-XC.S25). These include hypersurfaces whose binary blocks
have arbitrary distinct branch points, rather than only diagonal
monomials.

\textbf{Complete proof of the join criterion.}

First, the affine gradients of \(f\) and \(g\) vanish only at the
origin. Indeed, if all first partials of \(f\) vanish at a nonzero
vector \(x\), Euler's identity gives

\[
 m f(x)=\sum_i x_i\frac{\partial f}{\partial x_i}(x)=0.
\]

Since \(p\nmid m\), this would be a singular point of the projective
hypersurface \(F\). The same reasoning applies to \(g\). It follows at
once that every hypersurface in (CC-XC.S24) is smooth: vanishing of its
partial derivatives would force all its projective coordinates to
vanish. Their positive dimensions imply geometric integrality, by the
component-intersection argument in CC-XC09.

Choose a finite extension containing \(\epsilon\) with
\(\epsilon^m=-1\). It is enough to prove certification there, by
CC-XC08(4). For clarity use \(k\) for this enlarged field during the
construction. The map

\[
 \phi:X_f\times X_g\dashrightarrow Z_{f,g},\qquad
 ((x:u),(y:v))\longmapsto(x_0v:\cdots:x_rv:
                              \epsilon u y_0:\cdots:\epsilon u y_s)
 \tag{CC-XC.S26}
\]

lands in the stated target because

\[
 f(xv)+g(\epsilon uy)
 =v^m(-u^m)+(-1)u^m(-v^m)=0.
\]

At a point of \(X_f\), not all the \(x_i\) vanish, since that would also
force \(u=0\). Similarly, some \(y_j\) is nonzero on \(X_g\). On the
product chart \(x_i y_j\ne0\), the ideal generated by the coordinates in
(CC-XC.S26) is therefore exactly the ideal generated by \(u/x_i\) and
\(v/y_j\), up to the line-bundle trivialization. Thus the base scheme is

\[
 C=F\times G=\{u=v=0\}\subset X_f\times X_g.
\]

It is smooth of codimension two. Blowing up its ideal makes this
generating linear system invertible, and hence resolves the rational map
to projective space. The Fermat-type equation already verified on the
dense domain of \(\phi\) then holds on the blowup. We obtain a proper
morphism

\[
 \widetilde\phi:\operatorname{Bl}_{F\times G}(X_f\times X_g)
                     \longrightarrow Z_{f,g}.
 \tag{CC-XC.S27}
\]

This is exactly the portion of Shioda--Katsura, Remark 1.10, p.~102,
required here; the preceding ideal calculation verifies the resolution
directly.

We also verify its degree. The exceptional divisor maps into
\(\{f(z_x)=g(z_y)=0\}\): over \((x,y)\in F\times G\), the leading
coordinates of the linear system have their two blocks proportional to
\(x\) and \(y\). Points of the original domain with \(u=0\) or \(v=0\)
also map into this closed set. Take a geometric point represented as
\((z_x:z_y)\) in the dense open where \(f(z_x)\ne0\). Then
\(g(z_y)=-f(z_x)\), so both coordinate blocks are nonzero, and its
entire inverse image lies in the original domain with \(uv\ne0\). Each
root

\[
 u^m=-f(z_x)
\]

gives a preimage

\[
 ((z_x:u),(z_y/(\epsilon u):1)).
 \tag{CC-XC.S28}
\]

The first point lies on \(X_f\), and

\[
 g(z_y/(\epsilon u))
 =\frac{-f(z_x)}{(-1)(-f(z_x))}=-1,
\]

so the second lies on \(X_g\). Conversely, in every preimage of this
open point, both \(u\) and \(v\) are nonzero. Normalize \(v=1\), and
then use the scalar of the first projective factor to make its
\(x\)-block equal to \(z_x\). The equations force exactly the
presentation (CC-XC.S28). More precisely, on a target chart where one
coordinate in \(z_x\) is normalized to \(1\), these formulas and the
normalizations are regular inverse maps identifying the inverse image
with the cover defined by \(u^m+f(z_x)=0\). This is finite étale of
degree \(m\), because \(p\nmid m\) and \(f(z_x)\) is invertible on the
chosen open. Thus the morphism is dominant and generically finite of
degree \(m\).

Since \(X_f,X_g,F,G\) are certified, so are their two products.
CC-XC08(3) certifies the blowup in (CC-XC.S27), and CC-XC08(2) certifies
its target by the degree-\(m\) pullback/trace splitting. Finally
CC-XC08(4) descends this conclusion to the original field. This proves
the join criterion.

\textbf{Complete proof for the binary-form family.}

For each squarefree binary \(f_j\), its gradient has no nonzero common
zero by the same Euler argument: a nonzero common zero would be a
repeated point of its degree-\(m\) divisor on \(\mathbf P^1\). As the
blocks involve disjoint variables, the gradient of \(\sum f_j\) vanishes
only when every block is zero. Adding \(a z^m\) preserves this property,
since \(ma\ne0\). Thus both hypersurfaces in (CC-XC.S25) are smooth.

We prove simultaneously, by induction on \(t\), the assertions for all
choices of the binary forms and all \(a\ne0\). For \(t=1\), \(B_1\) is
finite étale and \(A_1\) is a smooth projective curve, so CC-XC08(1)
applies.

Suppose \(t\ge2\) and all assertions with \(t-1\) blocks have been
proved. To prove the assertion for \(B_t\), put

\[
 f=\sum_{j<t}f_j,\qquad g=f_t.
\]

The four varieties in the join criterion are respectively

\[
 F=B_{t-1},\qquad G=B_1(f_t),\qquad
 X_f=A_{t-1}(f_1,\ldots,f_{t-1};1),\qquad
 X_g=A_1(f_t;1).
 \tag{CC-XC.S29}
\]

All are certified by induction or the base case. Their dimensions fit
the join criterion with \(r=2t-3\ge1\) and \(s=1\), and its target is
\(B_t\). This proves the \(B_t\) assertion for every collection of \(t\)
squarefree binary forms.

Now prove the \(A_t\) assertion. Put

\[
 f=\sum_{j<t}f_j+a z^m,\qquad g=f_t.
\]

Here the four input varieties are

\[
 \begin{aligned}
 F&=A_{t-1}(f_1,\ldots,f_{t-1};a),&
 G&=B_1(f_t),\\
 X_f&=B_t(f_1,\ldots,f_{t-1},h),&
 X_g&=A_1(f_t;1),
 \end{aligned}
 \qquad h(z,u)=a z^m+u^m.
 \tag{CC-XC.S30}
\]

The binary form \(h\) is squarefree because \(p\nmid m\) and \(a\ne0\).
Thus \(X_f\) is certified by the \(B_t\) assertion just proved, not by
an unproved \(A_t\) induction hypothesis. The other three inputs are
certified by induction or the base case. The dimensions are
\(r=2t-2\ge2\) and \(s=1\), and the join target is \(A_t\). This
completes the simultaneous induction without a circular step.

Certification gives all the rational, prime-independent splitting and
numerical-category consequences in CC-XC06. The positive search of
CC-XC07 therefore terminates on every member of the stated family. This
is a specified geometric class with a finite inductive construction; it
does not give a termination theorem for a general homogeneous
polynomial. \(\square\)

\subsection{Exact primary reference}\label{exact-primary-reference}

Shioda, T., \& Katsura, T. (1979). On Fermat varieties. \emph{Tohoku
Mathematical Journal, Second Series, 31}(1), 97--115.
\url{https://doi.org/10.2748/tmj/1178229881} . \textbf{Exact locator:} Remark
1.10, p.~102, equation \((1.1')\), which states that the same
construction applies to smooth \(f=0\) and \(g=0\) in disjoint variable
sets and their three displayed augmented/sum hypersurfaces. The original
rational-map and resolution formulas are Lemmas 1.1--1.2, pp.~98--99.
Primary scan:
\url{https://www.jstage.jst.go.jp/article/tmj1949/31/1/31_1_97/_pdf} .

The local primary scan was directly inspected, including printed p.~102.
The full simultaneous induction and the degree/base-ideal verification
above are the deductions checked in this module; the general blowup
cohomology and trace closure inputs are those already checked in
CC-XC08.

% END consequences_criteria_separated_hypersurfaces.tex

% BEGIN consequences_criteria_rational_spectral_target.tex
\section{A rational abelian-sector target for arbitrary smooth
projective
cohomology}\label{a-rational-abelian-sector-target-for-arbitrary-smooth-projective-cohomology}

The remaining cross-cycle question in cote 15 can be put against one
\textbf{rational} target determined by the Frobenius polynomials of the
specified variety. The rational orbit and Schur-index obstructions in
CC-AG08 automatically vanish for the actual cohomology of a smooth
projective variety. The reason is crystalline Frobenius at the finite
places and the hard-Lefschetz pairing at the real places. This is a
deduction from those established results and the previously scoped
abelian-sector classification, with no historical-priority claim.

The arithmetic divisibility statement below is unconditional. Its
interpretation inside the rational abelian--Artin--Tate category uses
exactly the all-abelian inputs already recorded in CC-TATE01--02,
AQ-ABELIAN-RECOGNITION, and CC-AG08. No Tate assertion, numerical
realization, or Frobenius semisimplicity is assumed for the specified
variety.

\subsection{Primary inputs and
conventions}\label{primary-inputs-and-conventions}

Let \(k=\mathbf F_q\), \(q=p^a\), and let \(X/k\) be a nonempty smooth
projective scheme, pure of dimension \(d\). Use geometric Frobenius
\(F\) on étale cohomology. For \(0\leq i\leq 2d\), put \[
 P_i(t)=\det(t-F\mid H^i_{\mathrm{et}}
                  (X_{\bar k},\mathbf Q_\ell))\in\mathbf Z[t].
                                                        \tag{CC-XC.R11}
\] This is the monic characteristic polynomial, rather than the
\(\det(1-tF)\) convention. It is independent of \(\ell\ne p\); its roots
are nonzero integral \(q\)-Weil numbers of weight \(i\).

The precise primary results used are:

\begin{itemize}
\tightlist
\item
  Deligne (1974), Theorem 1.6, printed p.~276, gives the finite-field
  Weil theorem in the smooth projective case. Katz and Messing (1974),
  Theorem 1, printed p.~74, identify the characteristic polynomials in a
  reasonable Weil cohomology with the étale ones. Their immediately
  preceding paragraph explicitly includes crystalline cohomology. Their
  Corollary 1(2), p.~75, also transfers hard Lefschetz to those
  theories.
\item
  Put \(K_0=W(k)[1/p]\), and let \(\sigma\) be Witt Frobenius, of order
  \(a\). The rational crystalline cohomology
  \(C_i=H^i_{\mathrm{cris}}(X/W(k))[1/p]\) is a finite-dimensional
  \(K_0\)-space with an invertible \(\sigma\)-semilinear Frobenius
  \(\varphi\). Thus \(\varphi^a\) is \(K_0\)-linear. The Stacks Project,
  Remark 60.24.13, tag 07MY, gives perfectness, and Theorem 60.26.4, tag
  07N5, gives Frobenius invertibility after inverting \(p\). Their
  hypotheses hold for the structure crystal over \((W(k),(p),\gamma)\);
  Situation 60.26.1 and Definition 60.26.2, tags 07N2--07N3, specify the
  Frobenius lift and the nondegenerate structure.
\item
  Deligne (1980), Theorem 4.1.1, printed p.~217, gives hard Lefschetz
  for a smooth projective variety of pure dimension over an
  algebraically closed field and an ample line bundle. We use its
  ordinary Poincaré-duality consequence, with an ample line bundle
  already defined over \(k\).
\item
  The rational simple-object index formula is the previously scoped
  CC-AG08: Milne (1994), corrected author version of 2 November 2019,
  Proposition 2.8, p.~14, and Theorems 2.15--2.16, p.~18. Here
  \(\operatorname{End}\) means the endomorphism algebra of the
  cohomological motive, with the geometric-Frobenius convention
  specified on Milne's printed p.~17.
\end{itemize}

For clarity, Katz--Messing state their argument for an absolutely
irreducible smooth projective variety. The polynomial equality needed
here also follows for our \(X\). Partition its geometric connected
components into Frobenius orbits. For an orbit of length \(s\), one
component is defined over \(\mathbf F_{q^s}\) and is geometrically
connected and smooth, hence geometrically irreducible. In either
cohomology theory, after extending coefficients, the orbit contributes
the direct sum of the \(s\) component spaces, cyclically permuted by
\(F\), with \(F^s\) on a chosen component equal to its
\(\mathbf F_{q^s}\)-Frobenius. In the crystalline statement this means
the linear base change of \(\varphi^a\), whose \(s\)-th power agrees
with the \(\varphi^{as}\) for the enlarged constant field; it does not
mean a semilinear extension acting on the new coefficient field. A
cyclic block operator \(B\) satisfies \[
 \det(t-B)=\det(t^s-B^s\mid\text{one component});
                                                        \tag{CC-XC.R12}
\] one sees this by using \(B\) to identify successive component spaces,
which gives the usual block companion matrix. Katz--Messing on the
chosen component proves the equality on the orbit, and then on their
direct sum. The crystalline coefficient extension used here is
proper-smooth base change: Stacks, Remark 60.24.9, tag 07MU, modulo
\(p^n\), followed by perfect \(p\)-adic completion as in 07MY. We
therefore have, without assuming semisimplicity, \[
 \det(t-\varphi^a\mid C_i)=P_i(t).
                                                        \tag{CC-XC.R13}
\]

\subsection{CC-XC11: automatic Schur-index divisibility and its rational
target}\label{cc-xc11-automatic-schur-index-divisibility-and-its-rational-target}

Factor the monic rational polynomial as \[
 P_i(t)=\prod_f f(t)^{m_f},
                                                        \tag{CC-XC.R14}
\] where the \(f\) are distinct monic irreducible polynomials in
\(\mathbf Q[t]\) and \(m_f>0\). Choose a root \(\alpha_f\), put
\(K_f=\mathbf Q(\alpha_f)\), and normalize every \(p\)-adic valuation by
\(v_p(p)=1\). Define \[
 d_f=\operatorname{lcm}_v
 \operatorname{den}\begin{cases}
 i/2,&v\text{ real},\\[2pt]
 -[K_{f,v}:\mathbf Q_p]\,v_p(\alpha_f)/a,&v\mid p,\\[2pt]
 0,&\text{otherwise}.
 \end{cases}
                                                        \tag{CC-XC.R15}
\] The entries are taken in \(\mathbf Q/\mathbf Z\), and the denominator
of zero is one. These are exactly the local invariants and Schur index
of the rational abelian-sector simple object labelled by the Galois
orbit of \(\alpha_f\).

\textbf{Theorem.} For every factor in (CC-XC.R14), \[
                         d_f\mid m_f.                  \tag{CC-XC.R16}
\] Consequently there is a unique rational isomorphism class
\(M_i(X)\in\mathcal A_q\), pure of weight \(i\), with polynomial
\(P_i(t)\). If \(S_f\) is the rational simple with label \(f\), it is \[
 M_i(X)=\bigoplus_f S_f^{\oplus m_f/d_f},\qquad
 M(X)=\bigoplus_{i=0}^{2d}M_i(X).
                                                        \tag{CC-XC.R17}
\] For every \(\ell\ne p\), its realization is isomorphic to the
semisimplification of the graded Frobenius module \(H^*(X_{\bar k},
\mathbf Q_\ell)\). The target admits a rational presentation \[
 M(X)=(T,e),\qquad
 T=\bigoplus_{\nu=1}^{s}h(Y_\nu),
                                                        \tag{CC-XC.R18}
\] where every \(Y_\nu/k\) is an abelian variety or a finite étale
scheme and \(e\) is a rational numerical projector. By the established
numerical--homological equality on this sector, it is a projector in
every corresponding homological realization. No Tate twist is needed in
(CC-XC.R18).

The unique isomorphism class in (CC-XC.R17) is a spectral target. The
statement does not identify \(h(X)\) with it.

\subsubsection{Complete proof}

\textbf{The \(p\)-adic denominators.} Fix \(f\), and fix a place
\(v\mid p\) of \(K_f\). The place corresponds to an irreducible factor
\(f_v\in\mathbf Q_p[t]\) of \(f\), of degree \[
 n_v=\deg f_v=[K_{f,v}:\mathbf Q_p].
\] Set \(F_{\mathrm{cris}}=\varphi^a\) and choose an integer \(B\) at
least \(\dim_{K_0} C_i\). On \(C_i\) take the entire
\(\mathbf Q_p\)-primary block \[
 U_v=\ker\bigl(f_v(F_{\mathrm{cris}})^B\bigr).
                                                        \tag{CC-XC.R19}
\] Its dimension is \(m_f n_v\): after an algebraic closure of the
coefficient field it is precisely the sum of the generalized eigenspaces
for the \(n_v\) distinct roots of \(f_v\), each of algebraic
multiplicity \(m_f\).

It is essential here to take the whole \(f_v\)-block, not an individual
factor over \(K_0\). Since the coefficients of \(f_v\) lie in
\(\mathbf Q_p\), they are fixed by \(\sigma\). Also
\(\varphi F_{\mathrm{cris}}=F_{\mathrm{cris}}\varphi\). Consequently
(CC-XC.R19) is preserved by \(\varphi\) and by its inverse. In a
\(K_0\)-basis of \(U_v\), write \(\varphi=A\sigma\) with
\(A\in\mathrm{GL}_{m_fn_v}(K_0)\). Then \[
 \det(F_{\mathrm{cris}}\mid U_v)
   =\det A\cdot\sigma(\det A)\cdots\sigma^{a-1}(\det A)
   =N_{K_0/\mathbf Q_p}(\det A).
                                                        \tag{CC-XC.R20}
\] The extension \(K_0/\mathbf Q_p\) is unramified, so \(v_p(\det A)\)
is an integer and \(\sigma\) preserves the valuation. All roots of the
irreducible \(f_v\) have the same \(p\)-adic valuation as \(\alpha_f\)
at \(v\). Taking valuations of (CC-XC.R20) gives \[
 m_f n_v\,v_p(\alpha_f)=a\,v_p(\det A),\qquad
 m_f\,\frac{[K_{f,v}:\mathbf Q_p]v_p(\alpha_f)}a
        \in\mathbf Z.                                 \tag{CC-XC.R21}
\] Thus each finite-place denominator in (CC-XC.R15) divides \(m_f\).
This argument used the characteristic polynomial, with its full
algebraic multiplicities, and never discarded a possible Jordan block.

\textbf{A linear algebra fact for the real places.} Suppose \(V\) is a
finite-dimensional vector space over an algebraically closed
characteristic-zero field, \(B\) a nondegenerate alternating form, and
\(F\) an invertible operator satisfying \[
 B(Fx,Fy)=c\,B(x,y),\qquad c\ne0.
                                                        \tag{CC-XC.R22}
\] If \(V_\lambda\) denotes the generalized \(\lambda\)-eigenspace, then
\[
 B(V_\lambda,V_\mu)=0\quad\hbox{unless }\lambda\mu=c.
                                                        \tag{CC-XC.R23}
\] Indeed the adjoint of \(F\) is \(cF^{-1}\). If \(\lambda\mu\ne c\),
then \(cF^{-1}-\lambda\) is invertible on \(V_\mu\). For
\(x\in V_\lambda\), choose \(N\) with \((F-\lambda)^N x=0\). Every
\(y\in V_\mu\) equals \((cF^{-1}-\lambda)^N z\) for some \(z\in V_\mu\),
and the adjoint identity gives \[
 B(x,y)=B((F-\lambda)^N x,z)=0.
 \] The generalized eigenspaces sum to all of \(V\), so the
nondegeneracy of \(B\) implies that its pairing between \(V_\lambda\)
and \(V_{c/\lambda}\) is perfect. In particular, if \(\lambda^2=c\), the
restriction to \(V_\lambda\) is a nondegenerate alternating form, and \[
                         \dim V_\lambda\text{ is even}. \tag{CC-XC.R24}
\] This proof does not require \(F\) to be diagonalizable.

\textbf{Odd-degree cohomology.} First suppose \(i\le d\) is odd. Choose
an ample line bundle on \(X\) over \(k\), and view its first Chern class
as \(L\in H^2(X_{\bar k},\mathbf Q_\ell)\) after trivializing the Tate
twist. Then \(F(L)=qL\). Hard Lefschetz and Poincaré duality show that
\[
 B_i(x,y)=\operatorname{Tr}(x\smile y\smile L^{d-i})
                                                        \tag{CC-XC.R25}
\] is nondegenerate and alternating on \(H^i\). The trace means the sum
of the component traces when \(X_{\bar k}\) is disconnected; this
pairing is still perfect. Frobenius acts on a component's top cohomology
by \(q^d\) and permutes the components. Therefore \[
 \begin{aligned}
 q^{d-i}B_i(Fx,Fy)
 &=\operatorname{Tr}
       \bigl(F(x\smile y\smile L^{d-i})\bigr)\\
 &=q^d B_i(x,y),
 \end{aligned}
 \qquad B_i(Fx,Fy)=q^iB_i(x,y).                         \tag{CC-XC.R26}
\] Apply (CC-XC.R24) after extending coefficients to
\(\overline{\mathbf Q}_\ell\). Every occurrence of \(\lambda=q^{i/2}\)
or \(\lambda=-q^{i/2}\) has even algebraic multiplicity.

For odd \(i>d\), put \(j=2d-i<d\). Hard Lefschetz gives the isomorphism
\[
 L^{i-d}:H^j\xrightarrow{\sim}H^i,\qquad
 F_i L^{i-d}=q^{i-d}L^{i-d}F_j.                         \tag{CC-XC.R27}
\] It carries generalized eigenspaces for \(\beta\) to those for
\(q^{i-d}\beta\), preserving their dimensions. The real eigenvalues
\(\pm q^{i/2}\) therefore correspond to \(\pm q^{j/2}\), whose
multiplicities are even by the preceding case.

Now if \(K_f\) has a real embedding, a real conjugate of \(\alpha_f\)
must be one of \(\pm q^{i/2}\), by its Weil absolute value. Its
irreducible polynomial divides \(t^2-q^i\). When this splits over
\(\mathbf Q\), each relevant linear factor has even exponent in odd
degree. When it is irreducible, the centre is the real quadratic field
\(\mathbf Q(q^{i/2})\); each of its two roots has generalized eigenspace
dimension \(m_f\), so \(m_f\) is again even. For even \(i\), the real
invariant \(i/2\) is zero and imposes no condition. Thus every
real-place denominator in (CC-XC.R15) also divides \(m_f\).

The other invariants vanish. Taking their least common multiple proves
(CC-XC.R16).

\textbf{The rational motive and its effective presentation.} Each root
\(\alpha\) is a \(p\)-unit: both \(\alpha\) and its complex conjugate
\(q^i/\alpha\) are algebraic integers, so their valuations at a prime
away from \(p\) are nonnegative and sum to zero. Thus the roots belong
to the Weil-label group of CC-AG08. They have constant multiplicity
\(m_f\) on each rational Galois orbit. Equation (CC-XC.R16) is exactly
the second divisibility condition in CC-AG08. That theorem gives
(CC-XC.R17), including uniqueness of the rational isomorphism class. Its
realization is semisimple, with characteristic polynomial \(P_i\). Over
any characteristic-zero field, two semisimple operators with the same
characteristic polynomial are isomorphic as modules over the polynomial
ring: decompose them by irreducible factors, whose simple modules are
the corresponding residue fields. This proves the assertion about every
\(\ell\)-adic semisimplification.

For the presentation without twists, consider first a root \(\alpha\) in
positive degree \(i\). Choose \(\gamma\) with \(\gamma^i=\alpha\). It is
integral, and every complex conjugate has absolute value \(q^{1/2}\), so
it is a weight-one integral \(q\)-Weil number. Honda's Main Theorem
supplies an abelian variety \(B/k\) with \(\gamma\) on \(H^1\). The
Künneth summand \(H^1(B)^{\otimes i}\) of \(H^i(B^i)\) then has
eigenvalue \(\alpha\). In degree zero, every root is a root of unity,
and a root of order \(n\) occurs in
\(h(\operatorname{Spec}\mathbf F_{q^n})\).

Thus each simple \(S_f\) appearing in (CC-XC.R17) occurs in the rational
semisimple decomposition of the cohomology motive of one such abelian
variety or finite étale scheme. Indeed its realization contains a label
in the orbit \(f\), and the classification CC-AG08 identifies the unique
rational simple with that orbit. Take enough copies to accommodate the
desired integer number \(m_f/d_f\) of each rational simple.
Semisimplicity then supplies a rational summand projector on their
finite direct sum \(T\). On this sector numerical and homological
equivalence agree, so a rational cycle representative of the projector
gives (CC-XC.R18) in every realization. This completes the proof.
\(\square\)

The index \(d_f\) and the abstract rational isomorphism class (CC-XC.R17)
can be computed from the factorizations of \(P_i\), number-field local
degrees and valuations. This assertion does not say that those
polynomials alone print equations for the varieties \(Y_\nu\) or cycles
representing \(e\); a presentation in the already understood abelian
sector is additional concrete data to be chosen. Once chosen, it remains
fixed throughout the cross-cycle search below.

\subsection{CC-XC12: one fixed rational target and a finite cycle
certificate}\label{cc-xc12-one-fixed-rational-target-and-a-finite-cycle-certificate}

Fix a rational presentation (CC-XC.R18) of the target just proved to
exist. Write \(r_\nu=\dim Y_\nu\), treating components of a finite étale
\(Y_\nu\) as dimension zero. For rational cycles \[
 z_\nu\in\mathrm{CH}^d(X\times Y_\nu)_{\mathbf Q},\qquad
 w_\nu\in\mathrm{CH}^{r_\nu}(Y_\nu\times X)_{\mathbf Q},
                                                        \tag{CC-XC.R28}
\] let \[
 f=e\circ z:h(X)\longrightarrow M(X),\qquad
 g=w\circ e:M(X)\longrightarrow h(X).
                                                        \tag{CC-XC.R29}
\] These expressions refer first to correspondences; the numerical
projector \(e\) is used through its homological class on the abelian
sector. The codimensions in (CC-XC.R28) give degree-zero cohomological
maps in both directions.

\textbf{Theorem.} The following conditions are equivalent.

\begin{enumerate}
\def\labelenumi{\arabic{enumi}.}
\tightlist
\item
  For some \(\ell\ne p\), the object \(h_{\ell,E}(X)\), with
  \(E=\overline{\mathbf Q}_\ell\), belongs to the abelian sector of the
  homological category used in CC-XC02.
\item
  For some such \(\ell\), the determinant is nonzero on the image of the
  single projected cross-cycle map \[
  \bigoplus_\nu
      \mathrm{CH}^d(X\times Y_\nu)_{\mathbf Q}\otimes E
    \longrightarrow
      \operatorname{Hom}_{E[F]}
      \bigl(H^*(X_{\bar k},E),R_E(M(X))\bigr),\quad
        z\longmapsto R_E(ez).
                                                       \tag{CC-XC.R30}
  \]
\item
  There are rational cycles (CC-XC.R28) for which \[
      \det(gf\mid H^i(X_{\bar k},\mathbf Q_\ell))\ne0
         \quad\text{for every }0\le i\le2d.
                                                       \tag{CC-XC.R31}
  \] Here the determinants are rational numbers independent of
  \(\ell\ne p\), so the condition can be checked at no chosen
  cohomological prime.
\item
  Rational correspondences give mutually inverse isomorphisms between
  \(h(X)\) and this fixed \(M(X)\) in every \(\ell\)-adic homological
  realization, \(\ell\ne p\).
\end{enumerate}

The certificates in (3) are finite and their verification is the
intersection-number procedure of CC-XC05. Enumerating only the cycles
(CC-XC.R28), with \(T\) and \(e\) fixed, is a one-sided procedure that
halts exactly when these conditions hold. It does not decide
nonmembership by failing to halt.

\subsubsection{Complete proof}\label{complete-proof-1}

CC-XC11 shows that the fixed \(R_E(M(X))\) is precisely the semisimple
target required in CC-XC02, with the same dimension in each degree. Its
proof uses only a presentation by abelian varieties and finite étale
schemes, which (CC-XC.R18) provides. CC-XC02 therefore applies verbatim
to (CC-XC.R30): (1) and (2) are equivalent. If they hold, an invertible
projected map has an algebraic inverse in the \(E\)-coefficient
homological category, by the duality argument proved there.

Write such an inverse pair as finite sums of rational cycle
representatives with coefficients in \(E\): \[
 f_E=\sum_{u=1}^{b}c_u\,e z_u,\qquad
 g_E=\sum_{v=1}^{c}d_v\,w_v e.
 \] The projector is already rational. The product \[
 D(x,y)=\prod_{i=0}^{2d}
       \det\!\left(
       \left.\sum_{u,v}x_u y_v\,w_v e z_u
       \right|H^i(X_{\bar k},E)\right)
                                                        \tag{CC-XC.R32}
 \] is a polynomial in finitely many variables over \(E\), nonzero at
the coefficients of \(f_E,g_E\). A nonzero polynomial over an extension
field of \(\mathbf Q\) cannot vanish on every rational tuple: induct on
the number of variables, using the infinitely many rational values for
the last variable. Choose such a tuple. It gives rational \(f,g\)
satisfying (CC-XC.R31) at the selected prime.

For these rational correspondences, CC-XC05 supplies the
prime-independent rational characteristic polynomial of \(gf\) in every
degree. Explicitly, the rational Frobenius polynomials \(P_i\) are
pairwise coprime by their weights; the Chinese remainder theorem, using
the full characteristic polynomials \(P_i\) including their repeated
factors, gives rational polynomials in Frobenius acting as cohomological
projectors. Their construction needs no semisimplicity. The intersection
trace of the composition of such a projector with \((gf)^n\), including
its parity sign, gives the trace in the single degree \(i\). Newton's
identities then give that degree's characteristic polynomial. All these
intersection degrees are rational and independent of \(\ell\), as proved
in CC-XC05. Therefore the nonzero determinants just obtained are nonzero
in every realization. This proves (3).

Conversely suppose (3). On every degree, Cayley--Hamilton expresses the
inverse of \(gf\) as a polynomial in \(gf\) with rational coefficients:
the constant term of its characteristic polynomial is nonzero. Combine
these inverse polynomials using the rational cohomological projectors
just described. This gives a rational correspondence \(u\) acting as
\((gf)^{-1}\) in every degree and every \(\ell\ne p\). Set \(g'=ug\).
Then \(g'f=1\) on every realization of \(X\). Since the dimensions of
\(X\) and \(M(X)\) agree in each degree by CC-XC11, \(f\) is an
isomorphism there and \(fg'=1\) on \(M(X)\) as well. Identities of these
linear actions are identities in each homological category, by its
definition. Thus (4) holds. Condition (4) implies (1), completing the
equivalences.

Finally, rational algebraic cycles over the fixed finite field, with
rational coefficients, can be enumerated as in CC-XC07. The fixed
numerical projector is represented by a fixed rational matrix of cycles
on the products \(Y_\nu\times Y_\mu\). For each pair of lists
(CC-XC.R28), the finite compositions, intersections, traces, and Newton
identities are exactly the certificate verifier CC-XC05. A successful
certificate proves (3). If (1) holds, the finite rational specialization
above provides some certificate in the enumeration, so the procedure
eventually finds it. No step asserts termination when (1) fails.
\(\square\)

\subsection{What this removes and what
remains}\label{what-this-removes-and-what-remains}

For arbitrary smooth projective \(X/\mathbf F_q\), neither Galois-orbit
multiplicities nor local Schur indices are an additional rational-target
obstruction: (CC-XC.R16) proves their required divisibilities. The
spectral target is uniquely fixed before searching for
cross-correspondences. The remaining condition is a finite positive
certificate on \(X\) and that target, rather than a search for a target
isomorphism class.

The proof has not made Frobenius on \(H^*(X)\) semisimple:
semisimplicity follows only after the determinant certificate, because
then \(H^*(X)\) is isomorphic to the semisimple abelian-sector
realization. Nor has it supplied the cross-cycles for all \(X\). The
arbitrary-variety Tate/effectivity question therefore retains this
explicit cycle-existence gate. Projectivity is used in the real-place
argument through hard Lefschetz; the proof is not asserted for arbitrary
smooth nonproper schemes or without that pairing. The already proved
nonexistence of a \(\mathbf Q\)-valued fibre functor on the whole
rational abelian category also remains unchanged.

\subsection{APA bibliography and exact
locators}

Deligne, P. (1974). La conjecture de Weil: I. \emph{Publications
Mathématiques de l'IHÉS, 43}, 273--307. Theorem 1.6, p.~276.
\url{https://www.numdam.org/item/PMIHES_1974__43__273_0.pdf}

Deligne, P. (1980). La conjecture de Weil: II. \emph{Publications
Mathématiques de l'IHÉS, 52}, 137--252. Theorem 4.1.1, p.~217; the
paired odd-Betti-number consequence is Corollary 4.1.5, p.~218.
\url{https://publications.ias.edu/sites/default/files/Number40.pdf}

Honda, T. (1968). Isogeny classes of abelian varieties over finite
fields. \emph{Journal of the Mathematical Society of Japan, 20}(1--2),
83--95. Main Theorem, p.~83. \url{https://doi.org/10.2969/jmsj/02010083}

Katz, N. M., \& Messing, W. (1974). Some consequences of the Riemann
hypothesis for varieties over finite fields. \emph{Inventiones
Mathematicae, 23}, 73--77. Theorem 1, p.~74; Corollary 1, p.~75; Theorem
2, p.~76, for the cohomological-projector input already used in CC-XC05.
\url{https://web.math.princeton.edu/~nmk/old/katzmessing.pdf}

Milne, J. S. (1994). Motives over finite fields. In U. Jannsen, S. L.
Kleiman, \& J.-P. Serre (Eds.), \emph{Motives} (Proceedings of Symposia
in Pure Mathematics, Vol. 55, Part 1, pp.~401--459). American
Mathematical Society. Corrected author version, 2 November 2019:
Propositions 2.4, 2.6, 2.8, pp.~13--14; endomorphism convention, p.~17;
Theorems 2.15--2.16, p.~18. The full-Tate hypothesis in its Section 2 is
used here only through the already justified abelian-sector restriction
CC-AG08. \url{https://www.jmilne.org/math/xnotes/Seattle1.pdf}

The Stacks Project Authors. (2026). \emph{The Stacks Project}. Remark
60.24.9 (base change), tag 07MU; Remark 60.24.13 (perfectness), tag
07MY; Situation 60.26.1, tag 07N2; Definition 60.26.2, tag 07N3; Theorem
60.26.4 (Frobenius after inverting \(p\)), tag 07N5.
\url{https://stacks.math.columbia.edu/tag/07MU} ;
\url{https://stacks.math.columbia.edu/tag/07MY} ;
\url{https://stacks.math.columbia.edu/tag/07N2} ;
\url{https://stacks.math.columbia.edu/tag/07N3} ;
\url{https://stacks.math.columbia.edu/tag/07N5}

Proved internal inputs: CC-AG08 (rational abelian-sector
classification), CC-XC02 (one projected determinant criterion), and
CC-XC05--07 (finite rational-cycle certificates and their one-sided
enumeration). The external computability and intersection-theory inputs
to the last three results are retained with their exact locators in
research/consequences\_criteria\_cycle\_certificates.tex.

% END consequences_criteria_rational_spectral_target.tex

% BEGIN consequences_criteria_proper_parity.tex
\section{The proper-smooth parity question and the rational spectral
target}\label{the-proper-smooth-parity-question-and-the-rational-spectral-target}

\subsection{Exact known results answering the parity
question}\label{exact-known-results-answering-the-parity-question}

Deligne writes, immediately after Corollary 4.1.5:

\begin{quote}
Ce résultat devrait valoir sans hypothèse de projectivité, mais je ne
sais pas le prouver en caractéristique \(\ne0\).
\end{quote}

The source is \emph{Weil II}, printed p.~218.

The question following Deligne's Corollary 4.1.5 in \emph{Weil II},
printed p.~218, concerns removing projectivity from the odd Betti-number
parity statement in the \textbf{proper smooth} setting. That
interpretation is explicitly identified by Sun and Zheng (2016, p.~236).
The statement is now known. Two results must be distinguished:

\begin{longtable}[]{@{}
  >{\raggedright\arraybackslash}p{(\columnwidth - 4\tabcolsep) * \real{0.3333}}
  >{\raggedright\arraybackslash}p{(\columnwidth - 4\tabcolsep) * \real{0.3333}}
  >{\raggedright\arraybackslash}p{(\columnwidth - 4\tabcolsep) * \real{0.3333}}@{}}
\toprule\noalign{}
\begin{minipage}[b]{\linewidth}\raggedright
Statement
\end{minipage} & \begin{minipage}[b]{\linewidth}\raggedright
Exact primary result
\end{minipage} & \begin{minipage}[b]{\linewidth}\raggedright
Scope
\end{minipage} \\
\midrule\noalign{}
\endhead
\bottomrule\noalign{}
\endlastfoot
Odd-degree Betti numbers are even & Suh (2012), Corollary 2.2.3, p.~298
& Any proper smooth variety over any field; cohomology over a separably
closed extension; every odd degree \(r\ge1\) \\
Each of \(+q^{r/2}\) and \(-q^{r/2}\) has even algebraic multiplicity on
\(H^r\) for odd \(r\) & Suh (2012), Theorem 3.3.1, p.~301, applied using
proper-smooth purity; independently and more strongly, Sun--Zheng
(2016), Corollaries 1.2 and 1.4, pp.~236--237 & Proper smooth
\(X/\mathbf F_q\), without a projectivity or semisimplicity
assumption \\
\end{longtable}

Suh's Theorem 3.3.1 is stated for every separated scheme of finite type
over \(\mathbf F_q\): each of the two real numbers \(\pm q^{r/2}\),
\(r\ge1\) odd, has even multiplicity as a reciprocal zero or pole of the
zeta function. For a proper smooth scheme, purity of \(H^j\) isolates
weight \(r\) in degree \(r\); thus there is no cancellation with another
degree, and the theorem gives the second row. For a general scheme, its
zeta-function statement does not by itself give the same assertion in
each individual cohomological degree.

The later theorem supplies an actual equivariant form. Sun--Zheng,
Theorem 1.1, p.~236, states that for a proper equidimensional scheme
over an arbitrary field, its \(n\)-th intersection cohomology has a
Galois-equivariant perfect pairing with values in
\(\mathbf Q_\ell(-n)\), symmetric for even \(n\) and alternating for odd
\(n\). Their \textbf{Corollary 1.2 on the same page} gives the direct
ordinary-cohomology statement for every proper smooth scheme: \[
 H^n(X_{\bar k},\mathbf Q_\ell)\otimes
 H^n(X_{\bar k},\mathbf Q_\ell)
       \longrightarrow\mathbf Q_\ell(-n).
                                                        \tag{CC-XC.P33}
\] This is stronger than evenness of the total odd Betti number.

Over \(\mathbf F_q\), let \(\mu_{\lambda,e}\) count the Jordan blocks of
size \(e\) with eigenvalue \(\lambda\) on \(IH^n\). Sun--Zheng,
Corollary 1.4, p.~237, gives \[
 \mu_{\lambda,e}=\mu_{q^n/\lambda,e},\qquad
 \mu_{q^{n/2},e},\,\mu_{-q^{n/2},e}\in2\mathbf Z
                   \quad\hbox{when }n+e\text{ is even}.
                                                        \tag{CC-XC.P34}
\] In particular, for odd \(n\) each total algebraic multiplicity
\(\sum_e e\mu_{\pm q^{n/2},e}\) is even. The statement explicitly
accommodates nontrivial Jordan blocks. For smooth schemes \(IH^n=H^n\);
for singular proper schemes this theorem remains an intersection-
cohomology theorem, not an assertion of the same purity and parity for
ordinary cohomology.

The properness qualification cannot be discarded from raw odd Betti
parity. For example, \[
 H^1(\mathbf G_{m,\bar k},\mathbf Q_\ell)
      \simeq\mathbf Q_\ell(-1)
                                                        \tag{CC-XC.P35}
\] has dimension one. To see this directly, the Gysin sequence for
\(\mathbf G_m=\mathbf P^1\setminus\{0,\infty\}\) identifies this group
with the kernel of the sum map
\(\mathbf Q_\ell(-1)^2\to\mathbf Q_\ell(-1)\). Thus proper-smooth parity
is the relevant settled question, and it is distinct from statements
about arbitrary open varieties or arbitrary singular proper ordinary
cohomology.

\subsection{CC-XC13: rational spectral targets for proper smooth
schemes}\label{cc-xc13-rational-spectral-targets-for-proper-smooth-schemes}

\textbf{Theorem.} All arithmetic-divisibility and rational-target
conclusions of CC-XC11 hold when ``smooth projective'' is replaced by
``proper smooth''. More explicitly, let \(X/\mathbf F_q\), \(q=p^a\), be
a nonempty proper smooth scheme, pure of dimension \(d\), and write \[
 \det(t-F\mid H^i(X_{\overline{\mathbf F}_q},
                         \mathbf Q_\ell))
             =\prod_f f(t)^{m_f}.
                                                        \tag{CC-XC.P36}
\] For a root \(\alpha_f\) and \(K_f=\mathbf Q(\alpha_f)\), the
finite-field index \(d_f\) in (CC-XC.R15) divides \(m_f\). Hence the
unique rational abelian-sector object \[
 M_i(X)=\bigoplus_f S_f^{\oplus m_f/d_f}
                                                        \tag{CC-XC.P37}
\] has the semisimplification of \(H^i(X)\) as its \(\ell\)-adic
realization for every \(\ell\ne p\). Their sum has a presentation as a
rational summand of a finite direct sum of cohomology motives of abelian
varieties and finite étale schemes, with no Tate twist.

The arithmetic assertion is unconditional. The abelian-sector
interpretation uses the same explicitly scoped all-abelian
classification as CC-XC11; no arbitrary-variety Tate assertion is
introduced.

\subsubsection{Complete proof of the
extension}\label{complete-proof-of-the-extension}

The comparison input needed for the crystalline part also has the
precise proper-smooth extension. Suh, \textbf{Theorem 1.4, p.~296, with
proof on p.~297}, states that the conclusions of his Theorems 1.2 and
1.3 hold for every proper smooth variety over the finite field. Theorem
1.2 is the equality between crystalline and étale Frobenius
characteristic polynomials. Thus, with \(K_0=W(\mathbf F_q)[1/p]\) and
semilinear \(\varphi\) on \(C_i=H^i_{\mathrm{cris}}(X)[1/p]\), \[
 \det(t-\varphi^a\mid C_i)
     =\det(t-F\mid H^i_{\mathrm{et}}(X_{\bar k},
                                      \mathbf Q_\ell)).
                                                        \tag{CC-XC.P38}
\] The same result gives integral coefficients, and proper-smooth purity
gives weight \(i\). The perfectness and Frobenius-invertibility results
in Stacks tags 07MY and 07N5 require proper smoothness, so they apply
directly. Disconnected schemes are handled by the finite component
argument in CC-XC11, or by applying the proper-smooth statements to
their components.

Fix an irreducible \(\mathbf Q_p\)-factor \(f_v\) of one \(f\). The
whole primary block for \(f_v\) is \(\varphi\)-stable exactly as in
(CC-XC.R19), because its polynomial coefficients are fixed by Witt
Frobenius. It has dimension \(m_f[K_{f,v}:\mathbf Q_p]\). The
determinant norm calculation (CC-XC.R20)--(CC-XC.R21), which did not use
projectivity, therefore gives \[
 m_f[K_{f,v}:\mathbf Q_p]\,
          v_p(\alpha_f)/a\in\mathbf Z.
                                                        \tag{CC-XC.P39}
\] All finite-place index denominators divide \(m_f\).

For odd \(i\), use the perfect alternating form (CC-XC.P33) from
Sun--Zheng, Corollary 1.2. Since geometric Frobenius acts on
\(\mathbf Q_\ell(-i)\) by \(q^i\), it is a similitude of multiplier
\(q^i\) for this form. The complete generalized-eigenspace argument
(CC-XC.R22)--(CC-XC.R24) now applies in every odd degree directly: the
generalized eigenspaces at each of \(q^{i/2}\) and \(-q^{i/2}\) have
even dimension. Equivalently, one may use Corollary 1.4, whose more
detailed Jordan-block statement was recorded above.

If \(K_f\) has a real place, \(f\) divides \(t^2-q^i\). Thus for odd
\(i\) the multiplicity \(m_f\) is even, both for a rational real root
and for the real-quadratic factor with its two roots. For even \(i\),
the real-place invariant is zero. Combining this with (CC-XC.P39) proves
\(d_f\mid m_f\).

All remaining arguments of CC-XC11 use only integrality, weights, the
rational abelian-sector classification, and Honda's theorem. They
therefore apply without change: the eigenvalues are \(p\)-units, the
orbit multiplicities descend as in CC-AG08, positive weights are
realized by tensor powers of \(H^1\) after taking a weight-one root, and
degree zero is realized by finite étale schemes. This proves the stated
unique rational target and its untwisted presentation. No semisimplicity
conclusion for \(X\) was used or obtained. \(\square\)

This extension concerns the spectral target. It does not, on its own,
extend the projective intersection-computability and cycle-enumeration
algorithms in CC-XC12 to arbitrary proper nonprojective schemes. The
separate existence of cross-correspondences identifying \(X\) with its
target remains the required arbitrary-variety condition.

\subsection{APA bibliography and precise source
provenance}\label{apa-bibliography-and-precise-source-provenance}

Deligne, P. (1980). La conjecture de Weil: II. \emph{Publications
Mathématiques de l'IHÉS, 52}, 137--252. Corollary 4.1.5 and following
question, printed p.~218.
\url{https://publications.ias.edu/sites/default/files/Number40.pdf}

Suh, J. (2012). Symmetry and parity in Frobenius action on cohomology.
\emph{Compositio Mathematica, 148}(1), 295--303.
\url{https://doi.org/10.1112/S0010437X11007056} Theorem 1.4, p.~296, proof
p.~297; Theorem 2.2.2 and Corollary 2.2.3, p.~298; Theorem 3.3.1,
p.~301, proof completed on p.~302; Corollary 3.3.5, p.~302. The
publisher's journal PDF was inspected; the later arXiv upload 2410.01184
is not being treated as a new 2024 result.
\url{https://www.cambridge.org/core/services/aop-cambridge-core/content/view/CE57E10E819362A8F1DB4EC6C7DCBF26/S0010437X11007056a.pdf/symmetry-and-parity-in-frobenius-action-on-cohomology.pdf}

Sun, S., \& Zheng, W. (2016). Parity and symmetry in intersection and
ordinary cohomology. \emph{Algebra \& Number Theory, 10}(2), 235--307.
\url{https://doi.org/10.2140/ant.2016.10.235} Theorem 1.1 and Corollary 1.2,
p.~236; Corollaries 1.3--1.4 and Remark 1.5(2), p.~237. The published
primary PDF, including these rendered pages, was inspected directly.
\url{https://msp.org/ant/2016/10-2/ant-v10-n2-p02-s.pdf}

The Stacks Project Authors. (2026). \emph{The Stacks Project}. Remark
60.24.13, tag 07MY, and Theorem 60.26.4, tag 07N5.
\url{https://stacks.math.columbia.edu/tag/07MY} ;
\url{https://stacks.math.columbia.edu/tag/07N5}

Internal proved inputs: CC-XC11, including its complete primary-block
determinant and generalized-eigenspace proofs, and the scoped rational
classification CC-AG08.

% END consequences_criteria_proper_parity.tex

% BEGIN consequences_criteria_effective_target.tex
\section{Effective construction of the fixed rational spectral
target}\label{effective-construction-of-the-fixed-rational-spectral-target}

CC-XC11 identifies the rational isomorphism class of the target for a
smooth projective variety from its Frobenius polynomials. Here the
remaining choice of a concrete abelian-sector presentation is also made
effective. This construction always terminates on the stated spectral
input. The subsequent search for cross-cycles from the arbitrary variety
retains its one-sided termination condition.

The inputs from the established abelian sector are precisely Tate
surjectivity, numerical--homological equality, and Frobenius
semisimplicity on abelian varieties and their products, as recorded in
AQ-ABELIAN-RECOGNITION and CC-TATE01--02. The algorithm never assumes
any of these properties for the arbitrary input variety.

\subsection{CC-XC14: a terminating concrete-target
algorithm}\label{cc-xc14-a-terminating-concrete-target-algorithm}

Fix \(q=p^a\). Suppose a finite list \[
 P_0(t),\ldots,P_N(t)\in\mathbf Z[t]                       \tag{CC-XC.E1}
\] has the following properties. Each polynomial is monic, has nonzero
constant term, and all its roots are integral \(q\)-Weil numbers of
weight equal to its index. The polynomial \(1\), with no roots, is
permitted. For every rational irreducible factor \(f^{m_f}\) of \(P_i\),
the Schur index \(d_f\) from (CC-XC.R15) divides \(m_f\). These are
promised conditions on the input; they are satisfied by the Frobenius
polynomials of every smooth projective \(X/\mathbf F_q\), by CC-XC11,
and also by the polynomials of a proper smooth scheme supplied as input,
by CC-XC13.

If a smooth scheme has components of different dimensions, split it into
the finitely many open-and-closed unions of components of each fixed
dimension. These unions are defined over \(k\), and their degreewise
polynomials multiply to the polynomials of the whole scheme. The index
\(d_f\) depends only on \(f,q,i\), so its divisibility survives addition
of the factor multiplicities. The final cross-cycle procedure below is
applied separately to these pure-dimensional pieces; the whole object
belongs to the abelian sector exactly when each piece does, since that
sector is closed under finite direct sums and direct summands.

\textbf{Theorem.} There is a terminating algorithm that constructs
finitely many explicitly presented abelian varieties and finite étale
schemes \(Y_\nu/\mathbf F_q\), together with a rational matrix of
algebraic cycles representing a numerical idempotent \[
 e\in\operatorname{End}_{\mathcal A_q}(T),\qquad
 T=\bigoplus_\nu h(Y_\nu),                               \tag{CC-XC.E2}
\] such that, for every \(\ell\ne p\), \[
 \det(t-F\mid R_\ell(T,e)^i)=
 \begin{cases}
 P_i(t),&0\le i\le N,\\
 1,&i>N.
 \end{cases}                                           \tag{CC-XC.E3}
\] Here \(R_\ell(T,e)^i\) is the degree-\(i\) realization; the projector
is valid in every such realization. The object \((T,e)\) is the unique
rational spectral target specified by the list.

For a smooth projective input variety \(X/\mathbf F_q\) given by
equations, one may first compute its Frobenius polynomials and then run
this algorithm. After these two terminating steps, the fixed-target
cross-cycle search in CC-XC12 halts exactly when the cohomological
motive of \(X\) belongs to the abelian sector. This conclusion does not
assert termination of that last search for every \(X\), and gives no
complexity or cycle-degree bound.

\subsubsection{Complete proof and
algorithm}\label{complete-proof-and-algorithm}

\textbf{1. Computational operations and numerical ranks on abelian
products.} A smooth projective geometrically integral variety over a
finite field is called a ``nice variety'' in Poonen--Testa--van Luijk
(2015). Their Proposition 8.1, author p.~10, computes its zeta function.
Separating its factors by their weights then computes the degreewise
Frobenius polynomials, as in the proof of Proposition 8.2 on that page.
Every abelian variety, including every product used below, meets these
hypotheses.

Let \(Z/k\) be an abelian variety of dimension \(D\), and fix
\(0\le r\le D\). Write \(\operatorname{Num}^r(Z)_{\mathbf Q}\) for
rational codimension-\(r\) cycles over \(k\) modulo numerical
equivalence. Its dimension is the computable integer \[
 R_r(Z)=
 \operatorname{ord}_{t=q^r}
       \det(t-F\mid H^{2r}(Z_{\bar k},\mathbf Q_\ell)).
                                                        \tag{CC-XC.E4}
\] Indeed, semisimplicity identifies the right side with the dimension
of the Frobenius-invariant subspace of \(H^{2r}(Z)(r)\). Tate
surjectivity and numerical--homological equality, with the
scalar-injectivity statement in AQ.15, identify that space with
\(\operatorname{Num}^r(Z)_{\mathbf Q}\otimes\mathbf Q_\ell\). Only the
known abelian sector is used in this rank computation.

The complementary intersection pairing between
\(\operatorname{Num}^r(Z)_{\mathbf Q}\) and
\(\operatorname{Num}^{D-r}(Z)_{\mathbf Q}\) is perfect. By the
definition of numerical equivalence it is nondegenerate on both sides;
the spaces are finite-dimensional, so it is a perfect pairing and their
dimensions agree. Enumerate \(R=R_r(Z)\) cycles \(z_1,\ldots,z_R\) of
codimension \(r\) and \(R\) cycles \(w_1,\ldots,w_R\) of complementary
codimension until \[
             \det(z_u\cdot w_v)_{1\le u,v\le R}\ne0.
                                                        \tag{CC-XC.E5}
\] Such a pair of lists exists because cycles over \(k\) span both
numerical spaces. Each candidate is finitely checkable:
Poonen--Testa--van Luijk, Lemma 8.7, author p.~11, computes the
intersection numbers. Exhaustive enumeration therefore terminates. When
\(R=0\), use empty lists.

The classes of the \(z_u\) form a rational basis. For any further
rational cycle \(z\), the intersection numbers \(z\cdot w_v\), followed
by inversion of the rational matrix in (CC-XC.E5), give its coordinates
in that basis. In particular, numerical equality is now decidable for
cycles of this codimension on \(Z\). The argument is the
finite-dimensional intersection-matrix method appearing in
Poonen--Testa--van Luijk, Theorem 8.15(b), author p.~13, specialized
here to a known arithmetic numerical rank. It does not call their
conditional algorithm on an arbitrary variety.

\textbf{2. Effective numerical endomorphisms of an abelian sum.} Let \[
                   T_+=\bigoplus_{\nu=1}^{s}h(A_\nu)
                                                        \tag{CC-XC.E6}
\] be a finite sum of abelian varieties over \(k\). For every ordered
pair \((\nu,\mu)\), apply Step 1 to \[
 Z=A_\nu\times A_\mu,\qquad r=\dim A_\nu.
 \] This computes a rational cycle basis for the block
\(\operatorname{Hom}(h(A_\nu),h(A_\mu))\) in the numerical category,
with a finite coordinate test for every cycle in the block. Taking all
blocks gives a rational basis for the finite-dimensional algebra
\(\operatorname{End}(T_+)\).

Composition is also computable. Represent the two correspondences by
cycles, pull them to the triple product, intersect, and push forward to
the first and third factors. All these varieties are smooth projective.
The effective Chow-intersection operation is Poonen--Testa--van Luijk,
Remark 8.18, author p.~14; the pullbacks along projections and proper
pushforwards are computable from equations and generic degrees, as in
the complete verifier CC-XC05. Step 1 gives the numerical coordinates of
the resulting cycle. Thus for each rational element \(e\) of this
algebra the test \[
                              e^2=e                     \tag{CC-XC.E7}
\] is a finite exact rational calculation. No test of numerical
equivalence on a product involving the arbitrary \(X\) is being used.

\textbf{3. Testing the image polynomial of a projector.} Suppose \(e\)
passes (CC-XC.E7). Numerical--homological equality on the abelian
products makes its realization an idempotent in every \(\ell\)-adic
realization. All cycles are defined over \(k\); consequently this
idempotent commutes with Frobenius.

Compute the degreewise Frobenius polynomials of every \(A_\nu\), and the
rational cohomological projectors \(\pi_{i,\nu}\) by the Chinese
remainder construction in CC-XC05. Here \(F_\nu\) is the correspondence
acting as geometric Frobenius, represented by the transpose graph of the
relative \(q\)-power endomorphism in the cohomological correspondence
convention. For \(n\ge0\), the rational number \[
 s_{i,n}=(-1)^i\sum_\nu
   \operatorname{Tr}_{\mathrm{corr}}
       \bigl(\pi_{i,\nu}F_\nu^n e_{\nu\nu}\bigr)
                                                        \tag{CC-XC.E8}
\] is computable by intersection with the diagonal. The notation
\(\operatorname{Tr}_{\mathrm{corr}}\) is the ordinary intersection
trace, which is the alternating cohomological trace. Hence the factor
\((-1)^i\) in (CC-XC.E8) gives exactly \[
 s_{i,n}=\operatorname{Tr}
       (F^n\mid\operatorname{im}e\cap H^i(T_+)).
 \] For the direct sum the trace is the sum of the diagonal block
traces; off-diagonal blocks do not contribute. The formula holds at
every \(\ell\ne p\) and is independent of it.

In particular \(s_{i,0}\) is the integer rank of the image in degree
\(i\). First test that it equals the desired degree of \(P_i\), using
the value zero outside the specified degrees. For a matched rank
\(b_i\), the \(b_i\) further traces \(s_{i,1},\ldots,s_{i,b_i}\)
determine the characteristic polynomial of Frobenius on that image by
Newton's identities. Compare it with \(P_i\). Test every degree up to
the maximum of \(N\) and twice the largest dimension of the \(A_\nu\).
The realization is zero above the latter bound, so any nonconstant
prescribed polynomial in that range rejects the candidate. All larger
degrees vanish on both sides. This gives a finite exact test of every
required image polynomial. For a zero desired rank, the rank test
suffices.

\textbf{4. Explicit degree-zero target.} Every integral weight-zero Weil
number is a root of unity, by Kronecker's theorem, so factor \[
                       P_0(t)=\prod_n\Phi_n(t)^{b_n},
                                                        \tag{CC-XC.E9}
\] with finitely many nonzero \(b_n\). One can find the indices \(n\) by
enumerating cyclotomic polynomials and matching the irreducible factors;
the root-of-unity property guarantees termination. For each \(n\), the
Frobenius on \(h(\operatorname{Spec}\mathbf F_{q^n})\) is a regular
cyclic permutation with characteristic polynomial \(t^n-1\). Choose by
the rational Chinese remainder algorithm a polynomial \(u_n(t)\) that is
\(1\) modulo \(\Phi_n(t)\) and \(0\) modulo every other cyclotomic
factor of \(t^n-1\).

The explicit rational correspondence \(u_n(F)\), a rational linear
combination of Frobenius graphs, is an idempotent on this finite étale
scheme and its image has polynomial \(\Phi_n(t)\). The identity holds as
a correspondence here: \(F^n=1\) and the permutation action of
\(\mathbf Q[t]/(t^n-1)\) is faithful. Take \(b_n\) copies. This
constructs a concrete rational Artin projector \((T_0,e_0)\) with
polynomial \(P_0\). It is zero if \(P_0=1\).

\textbf{5. Exhaustive search and termination.} Use the effective
enumeration of abelian varieties and finite lists established in
CC-XC07: enumerate projective equations together with explicit identity,
inverse and group law on finite affine covers, and retain the smooth
geometrically connected commutative group schemes. The finite algebraic
verifications and their exhaustiveness are explained in that proof. For
each finite list \(T_+\), Steps 1--2 terminate and give a finite
rational basis of \(\operatorname{End}(T_+)\). Enumerate all rational
coefficient tuples in that basis, retain the idempotents, and apply Step
3. The desired positive-degree list is \(P_i\) for \(1\le i\le N\), with
polynomial \(1\) in degree zero and outside that range. Dovetail these
searches over all lists so that each finite stage in each search is
eventually executed.

The promised index divisibilities give the unique rational object
\(\bigoplus_{i\ge1,f}S_f^{m_f/d_f}\) by CC-AG08. The
effective-presentation argument in CC-XC11 proves that it is a rational
summand of a finite sum of abelian varieties: take a weight-one root for
each positive weight eigenvalue, apply Honda, and take tensor powers and
enough copies. This argument needs only the promised spectral
conditions, not the existence of an \(X\). Consequently some enumerated
list contains it, and some rational idempotent in the computed basis has
exactly the required image polynomials. The dovetailed search eventually
accepts this projector.

Return \[
                     (T,e)=(T_0\oplus T_+,e_0\oplus e).
                                                        \tag{CC-XC.E10}
\] The construction satisfies (CC-XC.E3), and uniqueness follows from
CC-AG08. If every positive-degree polynomial is \(1\), this step can
instead return \(T_+=0\) immediately. The completely zero list returns
the zero object.

Finally, if a smooth projective variety \(X\) is the original input, its
Frobenius polynomials are computable by CC-XC05 using Poonen--Testa--van
Luijk, Propositions 8.1--8.2. CC-XC11 verifies all the promises, so the
construction just proved terminates and supplies the concrete target
required by CC-XC12. That theorem then proves the exact one-sided
termination assertion for the remaining cross-cycle search. \(\square\)

\subsection{Scope of the algorithmic
conclusion}\label{scope-of-the-algorithmic-conclusion}

The guaranteed stage uses cycle searches only on products of abelian
varieties, where the relevant ranks and equivalence relations are
already known. It produces equations and rational projector cycles,
rather than only a list of Weil labels. The search on products involving
\(X\) begins only after that target has been fixed, and its universal
termination is not proved. In particular this theorem does not assert an
algorithm deciding arbitrary Tate classes or numerical equivalence on
arbitrary smooth projective varieties.

When proper nonprojective smooth cohomology polynomials are supplied as
spectral input, the target construction still applies by CC-XC13. The
final assertion about computing from projective equations and applying
the fixed-target cycle verifier retains its projective scope.

\subsection{APA bibliography and exact
locators}

Poonen, B., Testa, D., \& van Luijk, R. (2015). Computing Néron--Severi
groups and cycle class groups. \emph{Compositio Mathematica, 151}(4),
713--734. \url{https://doi.org/10.1112/S0010437X14007878} Inspected author
manuscript:
\url{https://math.mit.edu/~poonen/papers/compute_ns.pdf} Exact
author-page locators: §7.1, pp.~6--7; Propositions 8.1--8.2, p.~10;
Lemma 8.7, p.~11; Theorem 8.15(b), p.~13; Remark 8.18, p.~14. Theorem
8.36(b), p.~19, has its own conditional hypotheses and is not used on
arbitrary \(X\) in this construction.

Honda, T. (1968). Isogeny classes of abelian varieties over finite
fields. \emph{Journal of the Mathematical Society of Japan, 20}(1--2),
83--95. Main Theorem, p.~83. \url{https://doi.org/10.2969/jmsj/02010083}

Milne, J. S. (2020). \emph{Algebraic number theory} (Version 3.08, 19
July 2020). Corollary 5.6, printed p.~87, for the Kronecker
root-of-unity criterion. \url{https://www.jmilne.org/math/CourseNotes/ANT.pdf}

The exact all-abelian primary inputs and their scope are those in
CC-TATE01--02 and AQ.15. The complete rational classification, cycle
enumeration, and intersection-trace proofs are CC-AG08, CC-XC05--07, and
CC-XC11--13.

% END consequences_criteria_effective_target.tex

% BEGIN consequences_criteria_one_way_certificate.tex
\section{A fixed rational duality and a one-way cross-cycle
certificate}\label{a-fixed-rational-duality-and-a-one-way-cross-cycle-certificate}

\subsection{CC-XC15: only one family of cross-cycles remains to be
found}\label{cc-xc15-only-one-family-of-cross-cycles-remains-to-be-found}

Let \(X/\mathbf F_q\) be nonempty, smooth projective, and pure of
dimension \(d\). Fix the concrete rational target \[
       M=M(X)=(T,e),\qquad T=\bigoplus_\nu h(Y_\nu),
                                                        \tag{CC-XC.D1}
\] constructed by CC-XC14, where the \(Y_\nu\) are abelian varieties or
finite étale schemes. Put \(r_\nu=\dim Y_\nu\). All the known
abelian-sector inputs and the geometric-Frobenius conventions of
CC-XC11--14 remain in force. Correspondences involving \(X\) are
interpreted in the homological categories of CC-XC02, with rational
cycle representatives; numerical computations below are confined to the
abelian-sector targets.

\textbf{Theorem.} One can effectively construct, using only cycles
between the \(Y_\nu\), a rational isomorphism and its inverse \[
              \delta_M:M\xrightarrow{\sim}M^\vee(-d).
                                                        \tag{CC-XC.D2}
\] Choose it once. Let \(\delta_X:h(X)\xrightarrow{\sim}h(X)^\vee(-d)\)
be the canonical Poincaré-duality isomorphism, and, for rational cycles
\[
        z_\nu\in\mathrm{CH}^d(X\times Y_\nu)_{\mathbf Q},
        \qquad f=e z:h(X)\longrightarrow M,
                                                        \tag{CC-XC.D3}
\] form the rational adjoint correspondence \[
 f^\dagger=
   \delta_X^{-1}\circ f^\vee(-d)\circ\delta_M:
          M\longrightarrow h(X).                       \tag{CC-XC.D4}
\] The finite rational number \[
 \mathcal D_X(f)=
   \prod_{i=0}^{2d}
       \det\bigl(f^\dagger f\mid H^i(X_{\bar k},
                                             \mathbf Q_\ell)\bigr)
                                                        \tag{CC-XC.D5}
\] is computable by intersection numbers and is independent of
\(\ell\ne p\). It is nonzero if and only if \(f\) is an isomorphism on
\(\ell\)-adic cohomology, for one, and then every, \(\ell\ne p\).

Such an \(f\) exists if and only if the equivalent abelian-sector
membership conditions of CC-XC12 hold. Consequently, after the
terminating construction of \(M,\delta_M,\delta_M^{-1}\), a one-sided
procedure need enumerate only the forward cycles (CC-XC.D3). It accepts
precisely when \(\mathcal D_X(f)\ne0\). An accepted \(f\) determines a
rational inverse correspondence valid in every \(\ell\)-adic homological
realization. No separately chosen reverse cross-cycles on
\(Y_\nu\times X\) are required.

Only invertibility of \(\delta_M\) is used. The theorem does not require
it to have any symmetry or positivity property. For schemes with several
dimensions, apply this statement to the finitely many pure-dimensional
open-and-closed pieces, as in CC-XC14.

\subsubsection{Complete proof}

\textbf{Existence of the fixed rational duality.} Poincaré duality for
\(X\) gives, in every degree, a Frobenius-equivariant isomorphism \[
 H^i(X_{\bar k},\mathbf Q_\ell)
      \simeq
 H^{2d-i}(X_{\bar k},\mathbf Q_\ell)^\vee(-d).
                                                        \tag{CC-XC.D6}
\] This does not require Frobenius semisimplicity. On passing to
semisimplifications, (CC-XC.D6) shows that \(R_\ell(M)\) and
\(R_\ell(M^\vee(-d))\) have the same degreewise Frobenius polynomials.
Both are objects of the rational abelian sector. The rational
recognition theorem AQ.16, equivalently the uniqueness in CC-AG08,
therefore gives an actual rational isomorphism (CC-XC.D2).

Here is a terminating construction of that isomorphism. Poincaré duality
on \(Y_\nu\) identifies \[
 h(Y_\nu)^\vee(-d)=h(Y_\nu)(r_\nu-d).
 \] Thus, before applying the two projectors, the forward and reverse
Hom blocks used to find \(\delta_M\) and its inverse are represented
respectively by \[
 \begin{aligned}
 h(Y_\nu)&\longrightarrow h(Y_\mu)(r_\mu-d):
 &\mathrm{CH}^{r_\nu+r_\mu-d}
                       (Y_\nu\times Y_\mu)_{\mathbf Q},\\
 h(Y_\nu)(r_\nu-d)&\longrightarrow h(Y_\mu):
 &\mathrm{CH}^{d}
                       (Y_\nu\times Y_\mu)_{\mathbf Q}.
 \end{aligned}                                         \tag{CC-XC.D7}
\] As usual, a codimension outside \([0,\dim(Y_\nu\times Y_\mu)]\) gives
the zero group. These formulas retain the Tate twists; replacing the
first codimension by a uniform \(d\) would generally be incorrect.

The finite numerical-basis procedure of CC-XC14 computes every block and
its compositions. If a finite étale factor occurs, its product with an
abelian variety is a finite disjoint union of abelian varieties over
explicit finite extensions of \(k\), viewed as \(k\)-schemes. A product
of two finite étale schemes decomposes by the finite-field
tensor-product algorithm in the same way, with zero-dimensional abelian
varieties. Cycles on such a component are the cycles over that
component's constant field. Intersection degrees over \(k\) are their
degrees over the constant field multiplied by its extension degree. This
nonzero rational factor preserves numerical equivalence and
nondegeneracy of the pairing. Thus the same known-rank basis search,
over those explicit finite fields and then on their finite direct sum,
computes these blocks as well.

Let \(e'\) be the dual projector \(e^\vee(-d)\) on \(T^\vee(-d)\).
Enumerate rational projected maps \[
       u=e' u e:T\longrightarrow T^\vee(-d),\qquad
       v=e v e':T^\vee(-d)\longrightarrow T
 \] in the computed finite-dimensional numerical spaces, and test \[
                        vu=e,\qquad uv=e'.              \tag{CC-XC.D8}
\] These are decidable finite rational equalities inside the known
abelian sector. Enumerate the two coefficient tuples jointly, so that no
unsuccessful candidate monopolizes the search. The rational isomorphism
already proved to exist supplies a pair satisfying (CC-XC.D8); the
search therefore terminates. Rational cycle representatives of that pair
give \(\delta_M,\delta_M^{-1}\). Numerical--homological equality on this
sector makes them mutually inverse in every \(\ell\)-adic realization.

\textbf{The determinant test.} Duality of rational correspondences is
given by transposition with the indicated Tate twists. Thus (CC-XC.D4)
is computable from the forward cycles, the fixed \(\delta_M\), and the
diagonal duality of \(X\). It is a rational correspondence from \(M\) to
\(h(X)\).

Fix a prime \(\ell\ne p\), and choose bases of the total cohomology
spaces \(V=H^*(X)\) and \(W=R_\ell(M)\), with the corresponding dual
bases. Let \(A:V\to W\) be the matrix of \(f\), and let \(B_X,B_M\) be
the invertible matrices of the two dualities, after choosing a basis of
the one-dimensional Tate twist. Both spaces have the same dimension by
CC-XC11. The total matrix of \(f^\dagger f\) is \[
                         B_X^{-1}A^{\mathrm t}B_M A.
                                                        \tag{CC-XC.D9}
\] Consequently \[
       \det(f^\dagger f)
           =\det(B_X)^{-1}\det(B_M)\det(A)^2,
                                                        \tag{CC-XC.D10}
\] and it is nonzero exactly when \(A\) is invertible. This identity is
used on the total cohomology: duality pairs degree \(i\) with degree
\(2d-i\), so it does not assert a separate square identity for each
individual degree. Every correspondence in (CC-XC.D4) respects the
cohomological degree, so the determinant on the left is the product in
(CC-XC.D5).

The endomorphism \(f^\dagger f\) is a rational algebraic correspondence
on the smooth projective \(X\). CC-XC05 computes its characteristic
polynomial in every degree by rational Frobenius projectors,
intersection traces and Newton's identities. Therefore (CC-XC.D5) is a
rational, prime-independent computable number. If it is nonzero at one
prime, it is nonzero at all; (CC-XC.D9) then shows that \(f\) is an
isomorphism in every realization.

Cayley--Hamilton, degree by degree and combined by those rational
projectors, gives a rational correspondence \(u_X\) acting as
\((f^\dagger f)^{-1}\) in every realization. Hence \[
                             u_X f^\dagger              \tag{CC-XC.D11}
\] is a rational inverse of \(f\) in each homological category. This
proves all assertions about an individual candidate.

Finally, if the membership condition of CC-XC12 holds, that theorem
supplies a rational forward isomorphism \(f\) to this fixed target. The
fixed-duality calculation makes its (CC-XC.D5) nonzero. Conversely any
successful test gives an isomorphism to \(M\) by (CC-XC.D11), and hence
membership. Enumerating the rational cycles in (CC-XC.D3), and running
the finite test, therefore halts exactly in the asserted cases. No
conclusion follows from a search that has not halted. \(\square\)

The rational verifier is defined through intersection traces. Its
construction does not require a \(\mathbf Q\)-valued fibre functor on
the entire rational abelian category. The arbitrary-variety problem is
now the existence of one family of forward cross-cycles passing this
explicit quadratic-composition determinant test. All target, projector,
and duality data are obtained in the terminating abelian-sector stage.

\subsection{Exact dependencies and source
locators}\label{exact-dependencies-and-source-locators}

This is a complete deduction from AQ.16 and CC-AG08 (rational
recognition), CC-XC11 (the rational spectral target), CC-XC12 (rational
all-prime correspondence certificates), CC-XC14 (effective numerical Hom
bases on the abelian sector), and CC-XC05 (intersection-trace
verification). Their primary hypotheses are retained.

For the classical duality and composition conventions, the directly
inspected primary locators are Milne (1994), corrected 2019 author
version, §1, p.~4 for correspondence composition, and the proof of
Proposition 1.1, p.~6, for \(h(Y)^\vee=h(Y)(\dim Y)\), together with the
exact codimension convention in AQ.15. The computation uses
Poonen--Testa--van Luijk (2015), Propositions 8.1--8.2, Lemma 8.7,
Theorem 8.15(b), and Remark 8.18, author pp.~10--14, exactly as scoped
in CC-XC14. No new external classification theorem is required.

Milne, J. S. (1994). Motives over finite fields. In U. Jannsen, S. L.
Kleiman, \& J.-P. Serre (Eds.), \emph{Motives} (Proceedings of Symposia
in Pure Mathematics, Vol. 55, Part 1, pp.~401--459). American
Mathematical Society. Corrected author version:
\url{https://www.jmilne.org/math/xnotes/Seattle1.pdf}

Poonen, B., Testa, D., \& van Luijk, R. (2015). Computing Néron--Severi
groups and cycle class groups. \emph{Compositio Mathematica, 151}(4),
713--734. \url{https://doi.org/10.1112/S0010437X14007878} Author manuscript:
\url{https://math.mit.edu/~poonen/papers/compute_ns.pdf}

% END consequences_criteria_one_way_certificate.tex

% BEGIN consequences_criteria_chow_generation.tex
\section{Universal Chow generation supplies the missing forward
cycles}\label{universal-chow-generation-supplies-the-missing-forward-cycles}

\textbf{Scope.} The spectral target in CC-XC14 exists for every smooth
projective variety over a finite field; the remaining general gate is
the existence of algebraic maps from that variety to its target. Here a
geometric hypothesis supplies those maps. The input is Vial's published
decomposition theorem for rational Chow motives, not an additional Tate
conjecture for the variety under consideration. The finite-field descent
below is at the level of actual rational Chow correspondences. The final
conversion to the fixed target and the terminating search use the
abelian-sector inputs already recorded in AQ.15 and CC-XC11--15.

All Chow groups in this section have rational coefficients. A
\textbf{universal domain} \(\Omega\) is an algebraically closed field of
infinite transcendence degree over its prime field. In particular,
\(\overline{\mathbf F}_p\) is not such a domain. A Chow motive of
abelian type is an object of the full tensor subcategory closed under
finite sums, direct summands and duals generated by abelian varieties;
allowing finite étale schemes as additional generators gives its
abelian--Artin version. These are statements about rational equivalence,
before passage to homological or numerical equivalence.

\subsection{Primary input and
conventions}\label{primary-input-and-conventions}

Vial (2017, Theorem 3.11, pp.~213--214) works with covariant Chow
motives and writes \(\mathbf 1(1)\) for the Lefschetz motive. If a
morphism \(N\to h(X)\), with \(N\) finite-dimensional in the sense of
Kimura, surjects on algebraically trivial \(\mathrm{CH}_i\) over a
universal domain for \(i<l\), his self-dual conclusion splits \(h(X)\)
into a summand built from \(N,N^\vee\) and twisted curves, and a twisted
summand of a smooth projective variety of dimension at most
\(\dim X-2l\). This theorem has no characteristic-zero assumption. Its
special case with \(l=\lfloor\dim X/2\rfloor\) is Theorem 5,
pp.~197--198. The algebraically closed base-extension descent used below
is Vial's Theorem 1, p.~196, proved in §1.1, p.~199.

We use the contravariant convention of the preceding modules. Thus a
positive Lefschetz twist in Vial's convention is a negative étale Tate
twist here. The assertion that a motive belongs to the category
generated by twisted curves or abelian varieties is unchanged by this
translation. We retain all twists in the argument and do not assert an
untwisted Chow presentation from Vial's theorem.

\subsection{CC-XC16 --- A geometric sufficient condition, with actual
finite-field
descent}\label{cc-xc16-a-geometric-sufficient-condition-with-actual-finite-field-descent}

Let \(k=\mathbf F_q\), let \(X/k\) be smooth projective and
geometrically connected of pure dimension \(d\ge2\), and let
\(\Omega\supset\bar k\) be a universal domain. Put
\(l=\lfloor d/2\rfloor\). Suppose there is a finite direct sum \(N\) of
Tate twists of motives of products of smooth projective curves over
\(\Omega\), and a morphism of rational Chow motives

\[
                 u:N\longrightarrow h(X_\Omega),
 \tag{CC-XC.G1}
\]

written in Vial's covariant convention, such that

\[
 u_*:\mathrm{CH}_i(N)_{\mathrm{alg}}
       \longrightarrow \mathrm{CH}_i(X_\Omega)_{\mathrm{alg}}
 \quad\text{is onto for every }0\le i<l.
 \tag{CC-XC.G2}
\]

Here the finite direct sum and its finitely many correspondences are
part of the hypothesis. They may be defined over \(\Omega\), rather than
over \(k\). No finite-dimensionality assumption on \(h(X)\) is imposed.

Then \(h(X)\) belongs to the rational Chow category generated by abelian
varieties, finite étale \(k\)-schemes and Tate twists. In particular,
\(X\) satisfies the equivalent conditions CC-XC06 and CC-XC12. The
always-terminating construction CC-XC14 followed by the one-way search
CC-XC15 terminates for \(X\), and the accepted rational correspondences
give its fixed-target isomorphism for every \(\ell\ne p\).

\textbf{Complete deduction from the stated primary input.}

Every curve motive is finite-dimensional, and this property is preserved
by the finite sums, products and twists in \(N\); these are the
curve-source facts used in Vial's theorem and recalled on p.~195. Apply
Theorem 3.11 with algebraic equivalence and \(l=\lfloor d/2\rfloor\).
For negative \(i\), the target \(\mathrm{CH}_i(X_\Omega)\) is zero, so
(CC-XC.G2) gives the entire range \(i<l\) required there. Chow duality
gives the self-dual hypothesis for \(h(X_\Omega)\). The remaining
variety has dimension at most

\[
                    d-2\lfloor d/2\rfloor\in\{0,1\}.
 \tag{CC-XC.G3}
\]

Its motive, as well as every other summand produced by the theorem, is
generated by curves and Tate motives. Thus \(h(X_\Omega)\) is of abelian
type. This step uses an actual Chow splitting furnished by the cited
theorem, not an equality of Frobenius polynomials.

Apply Vial's Theorem 1 to the motive \(h(X_{\bar k})\) and the extension
\(\Omega/\bar k\). It gives a Chow splitting over \(\bar k\). Expanding
finite sums, tensor products and duals of abelian motives, we may write
it in the form

\[
 h(X_{\bar k})\ \xrightarrow{\ a\ }\
 T_{\bar k}=\bigoplus_\nu h(A_\nu)(n_\nu)
 \ \xrightarrow{\ b\ }\ h(X_{\bar k}),\qquad ba=1,
 \tag{CC-XC.G4}
\]

with abelian varieties \(A_\nu/\bar k\) and integers \(n_\nu\). Products
of abelian varieties are abelian, their dual motives are Tate twists of
their motives, and Tate motives can be included as summands of powers of
an elliptic curve, so this form loses no summands.

There is a finite extension \(k'/k\) over which all of (CC-XC.G4) is
defined \textbf{as a Chow splitting}. Indeed, the finitely many
varieties, their structure maps, the finitely many cycles representing
\(a,b\), and the finitely many rational-equivalence relations witnessing
\(ba=\Delta_X\) use finitely many elements of \(\bar k\). Each such
element lies in a finite extension of \(k\); their compositum is finite.
Enlarging it if necessary also defines the group structures and the
smooth projective models. Consequently

\[
 h(X_{k'})\xrightarrow{a'} T'
       \xrightarrow{b'}h(X_{k'}),\qquad b'a'=1
 \quad\text{in rational Chow motives over }k'.
 \tag{CC-XC.G5}
\]

For completeness, the descent from (CC-XC.G5) to membership over \(k\)
is as follows. Put \(e=[k':k]\). Regard a smooth projective
\(k'\)-scheme as a smooth projective \(k\)-scheme. A correspondence on
\(Y\times_{k'} Z\) extends by zero along the open-and-closed immersion
\(Y\times_{k'} Z\hookrightarrow Y\times_k Z\). This operation preserves
rational equivalence, identities and composition. To check composition,
extend to \(\bar k\): the schemes split into components indexed by the
embeddings \(k'\hookrightarrow\bar k\), and the induced correspondence
is block diagonal, with the same conjugate correspondence in each
matching block. The composition is the same block-diagonal composition.
The verification descends to rational Chow groups because base change on
rational Chow groups is injective (Vial, Lemma 1.2, p.~200). Denote this
additive operation by \(\operatorname{Ind}_{k'/k}\).

For an abelian variety \(A/k'\), let \(B=\operatorname{Res}_{k'/k}A\).
This is an abelian variety over \(k\), and

\[
 B_{k'}\simeq\prod_{\sigma\in\operatorname{Gal}(k'/k)}A^\sigma .
 \tag{CC-XC.G6}
\]

These precise restriction-of-scalars properties are given in Milne
(1972, §1, pp.~177--178). Projection onto the identity factor and
insertion into that factor, using the zero sections in the other
factors, have composite \(1_A\). Their graph correspondences therefore
exhibit

\[
 \operatorname{Ind}_{k'/k}h(A)
 \quad\text{as a Chow direct summand of}\quad
 h(B)\otimes h(\operatorname{Spec}k').
 \tag{CC-XC.G7}
\]

The right side is explicitly an abelian motive tensored with an Artin
motive. The statement does not replace this product by a single abelian
variety. Twists commute with the induction operation. Thus
\(\operatorname{Ind}_{k'/k}T'\), and its direct summand
\(h(X_{k'}\text{ viewed over }k)\), lie in the stated Chow subcategory.

Finally, the finite étale projection \(\tau:X_{k'}\to X\), with source
now viewed over \(k\), supplies pullback and trace correspondences
satisfying

\[
                         \tau_*\tau^*=e\,1_{h(X)} .
 \tag{CC-XC.G8}
\]

This is the graph-correspondence composition identity: the projection
from the corresponding fibre product onto the diagonal has degree \(e\).
It also follows from proper pushforward and the projection formula.
Division by the positive integer \(e\) is allowed over \(\mathbf Q\).
Hence \(\tau^*\) and \(e^{-1}\tau_*\) split \(h(X)\) off the induced
object, proving the claimed Chow inclusion over the specified field
\(k\).

Pass these actual correspondences to any homological realization.
CC-XC06 gives certification; CC-XC12 identifies the resulting
homological object with the fixed rational spectral target. CC-XC14
constructs a concrete untwisted numerical presentation of that target,
and CC-XC15 proves that its forward-cycle search terminates exactly when
a forward isomorphism exists. This proves the algorithmic conclusion. It
does not assert a Chow isomorphism to the particular numerical projector
output by CC-XC14. \(\square\)

\subsection{CC-XC17 --- Universal zero-cycles, including rationally
chain connected
threefolds}\label{cc-xc17-universal-zero-cycles-including-rationally-chain-connected-threefolds}

Let \(X/k\) be smooth projective and geometrically connected with
\(2\le d\le3\). If

\[
              \deg:\mathrm{CH}_0(X_\Omega)_{\mathbf Q}
                         \xrightarrow{\ \sim\ }\mathbf Q ,
 \tag{CC-XC.G9}
\]

then \(h(X)\) belongs to the same rational Chow abelian--Artin category
and its fixed-target forward-cycle search terminates. In fact, Vial's
theorem in this case expresses the Chow motive as a direct summand of a
finite sum of twisted curve motives. In particular the conclusion
applies when \(X\) is geometrically rationally chain connected, in every
characteristic.

\textbf{Complete proof.} Algebraically trivial zero-cycles have degree
zero. Thus (CC-XC.G9) implies
\(\mathrm{CH}_0(X_\Omega)_{\mathrm{alg}}=0\). Apply CC-XC16 with \(N=0\)
and \(l=1\). More precisely, one can apply Vial's Theorem 3.11 directly
to \(0\to h(X)\) over \(k\), since its Chow-group hypothesis is tested
over \(\Omega\): the self-dual remainder has dimension at most
\(d-2\le1\), and its other terms are twisted curves. This proves the
asserted stronger curve-source form with rational coefficients; no
numerical projector-lifting conjecture enters.

For the rational-chain-connected assertion, over an algebraically closed
extension such as \(\Omega\), any two points are joined by a connected
finite chain of images of \(\mathbf P^1\). On \(\mathbf P^1\), the
difference of two points is principal. Proper pushforward to
\(X_\Omega\) therefore makes the endpoints of each member of the chain
rationally equivalent; adding along the chain makes its two endpoints
rationally equivalent. Since closed points of \(X_\Omega\) are
\(\Omega\)-points, every zero-cycle is rationally equivalent to its
degree times one chosen point. This proves (CC-XC.G9), even integrally
over \(\Omega\). The standard rational-chain-connected implication and
its transfer formulation are also recorded in the published
Kahn--Sujatha Proposition 3.1.4(2), p.~400; Proposition 3.1.1,
pp.~398--399, gives the equivalent universal rational zero-cycle
formulations.

This argument uses smoothness and projectivity of \(X\), and uses
rational coefficients for the resulting motivic splitting. It does not
infer integral descent of a diagonal decomposition, or universal
integral triviality over nonclosed extensions of \(k\). The restriction
\(d\le3\) is essential to the application: if \(d\ge4\), vanishing of
only \(\mathrm{CH}_0(X_\Omega)_{\mathrm{alg}}\) leaves a possible
residual motive of dimension \(d-2\), which need not be generated by
curves. \(\square\)

\subsection{\texorpdfstring{Why the universal-domain condition cannot be
replaced by
\(\bar k\)}{Why the universal-domain condition cannot be replaced by \textbackslash bar k}}\label{why-the-universal-domain-condition-cannot-be-replaced-by-bar-k}

There is a direct test case. Let \(C/k\) be a smooth projective
geometrically connected curve of positive genus. Its Jacobian has only
torsion \(\bar k\)-points: each point is defined over a finite field,
whose group of rational points is finite. Consequently
\(\mathrm{CH}_0(C_{\bar k})_{\mathbf Q}^0=J(C)(\bar k)\otimes\mathbf Q=0\).
Nevertheless \(H^1(C)\ne0\). Moreover, after enlarging to a universal
domain \(\Omega\), the generic point of \(J(C)\) over \(\bar k\),
embedded in \(\Omega\), is not torsion: an equation \([n]x=0\) would
force the generic point to lie in the finite kernel of \([n]\). Thus
\(J(C)(\Omega)\otimes\mathbf Q\ne0\), and the same degree-zero Chow
group over \(\Omega\) is nonzero. The canonical degree-zero
Chow/Jacobian identification for a pointed curve gives these equalities;
a point exists over \(\bar k\).

A fixed-field vanishing therefore does not verify the universal
vanishing required by the theorem. This example distinguishes the
hypotheses; it is not a counterexample to abelian-type membership, since
a curve is itself of abelian type. Vial's zero-cycle discussion in §3.1,
pp.~207--208, explicitly explains why a test over a finite field cannot
replace the field-extension condition. In CC-XC16--17 the universal
Chow-generation assumption is geometric input; the earlier integral
comparison of Albanese relations in CC-ALB01--02 does not supply it
automatically.

\subsection{Resulting change to the original cote 15
gate}\label{resulting-change-to-the-original-cote-15-gate}

The forward-cycle existence problem is discharged for the stated
universal-Chow-generation class, including every smooth projective
geometrically rationally chain connected surface or threefold. The
conclusion is actual Chow membership, hence stronger than the required
numerical or homological inclusion. Neither the rational-chain-connected
test nor the more general Chow-generation condition is asserted for an
arbitrary input variety. For such an input, the forward-cycle existence
gate remains; the one-sided search still has no negative-membership
stopping rule.

\subsection{Primary bibliography and exact
locators}\label{primary-bibliography-and-exact-locators}

\begin{itemize}
\tightlist
\item
  Vial, C. (2017). Remarks on motives of abelian type. \emph{Tohoku
  Mathematical Journal, Second Series, 69}(2), 195--220.
  \url{https://doi.org/10.2748/tmj/1498269623} . Published source:
  \url{https://www.jstage.jst.go.jp/article/tmj/69/2/69_1498269623/_pdf/-char/en}
  . Exact locators: conventions and curve finite-dimensionality, p.~195;
  Theorem 1, p.~196, and its §1.1 proof, p.~199; Theorem 5,
  pp.~197--198; Lemma 1.2, p.~200; universal-domain zero-cycle
  discussion, §3.1, pp.~207--208; Theorem 3.11 and complete proof,
  pp.~213--214. The journal version was received 24 October 2014 and
  revised 21 July 2015. The website's later release date is not a new
  theorem date.
\item
  Kahn, B., \& Sujatha, R. (2016). Birational motives, I: Pure
  birational motives. \emph{Annals of K-Theory, 1}(4), 379--440.
  \url{https://doi.org/10.2140/akt.2016.1.379} . Published source:
  \url{https://msp.org/akt/2016/1-4/akt-v1-n4-p02-s.pdf} . Exact locators:
  Proposition 3.1.1, pp.~398--399; Proposition 3.1.4(2), p.~400,
  including its proof. The older preprint's Corollary 3.1.2 is not the
  published locator for rational chain connectedness.
\item
  Milne, J. S. (1972). On the arithmetic of abelian varieties.
  \emph{Inventiones Mathematicae, 17}, 177--190.
  \url{https://www.jmilne.org/math/articles/1972a.pdf} . Exact locator: §1,
  pp.~177--178, for restriction of scalars of an abelian variety and its
  product decomposition after finite separable base change. This is the
  same primary input already inspected for CC-XC08.
\end{itemize}

% END consequences_criteria_chow_generation.tex

% BEGIN consequences_criteria_quadric_fibrations.tex
\section{Fibrations with small fibre Chow groups: a geometric reduction
of the cycle
gate}\label{fibrations-with-small-fibre-chow-groups-a-geometric-reduction-of-the-cycle-gate}

\textbf{Scope.} This is a deduction from Vial's published motivic
fibration theorem, with rational coefficients. It constructs a Chow
decomposition of the original variety. In the cases below where the
remaining summand comes from a curve or a finite étale scheme, the
forward-cycle existence gate in CC-XC15 is therefore discharged. No
claim about an arbitrary fibration is made from its generic fibre alone.

Write \(\mathcal C_{\mathrm{ab,Art}}(k)\) for the full rational Chow
subcategory generated by abelian varieties, finite étale schemes and
Tate twists, closed under sums, tensor products, duals and summands.
Smooth projective curves belong to this category: over \(\bar k\) this
is the curve/abelian equivalence in Vial (2017, p.~195), and the actual
restriction--corestriction proof in CC-XC16 descends this membership to
\(k\). Write \(\mathcal A_{q,E}\) for the fully embedded abelian--Artin
subcategory of the \(E=\overline{\mathbf Q}_\ell\)-linear homological
category used in CC-XC01--15. Thus certification means membership in
\(\mathcal A_{q,E}\), with its already-proved rational, all-\(\ell\)
consequences.

\subsection{\texorpdfstring{CC-XC18 --- A residual motive of dimension
\(\dim X-2n\)}{CC-XC18 --- A residual motive of dimension \textbackslash dim X-2n}}\label{cc-xc18-a-residual-motive-of-dimension-dim-x-2n}

Let \(k=\mathbf F_q\). Let \(f:X\to B\) be a flat, projective,
surjective morphism of smooth projective geometrically integral
\(k\)-varieties. Put

\[
       D=\dim X,\qquad b=\dim B\ge1,\qquad r=D-b\ge1.
 \tag{CC-XC.Q1}
\]

Let \(\Omega\supset\bar k\) be a universal domain. Suppose \(n\) is an
integer with \(1\le 2n\le D\) and

\[
 \dim_{\mathbf Q}\mathrm{CH}_j(X_t)_{\mathbf Q}=1
 \quad\text{for every }t\in B(\Omega)\text{ and every }0\le j<n.
 \tag{CC-XC.Q2}
\]

The notation includes singular and nonreduced fibres; Chow groups use
their usual scheme-theoretic cycles modulo rational equivalence. Then
there are a smooth projective \(k\)-variety \(Z\) of dimension \(D-2n\),
an idempotent \(e\in\mathrm{CH}^{\dim Z}(Z\times Z)_{\mathbf Q}\), and
an isomorphism of rational Chow motives

\[
 h(X)\simeq
 \left(\bigoplus_{i=0}^{r} h(B)(-i)\right)
       \oplus (Z,e)(-n).
 \tag{CC-XC.Q3}
\]

Here and below twists are in the contravariant convention of this
project. Consequently:

\begin{enumerate}
\def\labelenumi{\arabic{enumi}.}
\tightlist
\item
  \(h(X)\in\mathcal C_{\mathrm{ab,Art}}(k)\) if and only if both
  \(h(B)\) and \((Z,e)\) belong to that category.
\item
  \(X\) is certified if and only if \(B\) is certified and the
  homological image of \((Z,e)\) belongs to \(\mathcal A_{q,E}\).
\item
  If \(D-2n\le1\), the extra condition on \((Z,e)\) is automatic. Thus
  in this case Chow membership of \(X\) is equivalent to Chow membership
  of \(B\), and certification of \(X\) is equivalent to certification of
  \(B\).
\end{enumerate}

\textbf{Complete deduction.} Formula (CC-XC.Q3) is exactly Vial (2013,
Theorem 4.2, p.~1534), translated from his covariant convention. His
positive Lefschetz twist becomes our negative Tate twist; the idempotent
is transposed under the corresponding opposite-category convention, and
we have denoted the resulting idempotent by \(e\). The cited theorem
assumes (CC-XC.Q2) for every \(\Omega\)-point, and does not impose
characteristic zero, smoothness of \(f\), or a smooth discriminant. The
projective-surjective convention is fixed at the beginning of the paper,
p.~1522. We impose \(2n\le D\) to keep the displayed residual dimension
nonnegative.

The isomorphism is at the level of Chow correspondences. Its proof on
p.~1534 uses the graph/hyperplane splitting of Theorem 1.4 and the
self-dual base projector of Proposition 1.5. The fibre-Chow hypothesis
makes both low-dimensional Chow groups of the residual projector and of
its transpose vanish, and Proposition 4.1 then gives the stated residual
variety. We invoke that established theorem; equality of zeta functions
is not being used to infer (CC-XC.Q3).

Both categories in assertions 1 and 2 are closed under finite sums,
summands and Tate twists. If the two indicated objects belong,
(CC-XC.Q3) therefore puts \(h(X)\) in the category. Conversely, \(h(B)\)
is the \(i=0\) direct summand of (CC-XC.Q3), and \((Z,e)\) is a twist of
another direct summand, so membership of \(h(X)\) forces both. Passing
the actual Chow isomorphism to homological motives is legitimate without
assuming numerical--homological comparison for \(Z\): the latter
category is defined using the kernel of the chosen cycle realization.

If \(\dim Z=0\), smoothness and projectivity over a finite field imply
that \(Z\) is finite étale. If \(\dim Z=1\), its components are smooth
projective curves, possibly with larger finite constant fields; the
finite-field descent argument in CC-XC16 includes those components and
their Artin factors. Thus \(h(Z)\), and hence its summand \((Z,e)\),
lies in \(\mathcal C_{\mathrm{ab,Art}}(k)\) whenever \(\dim Z\le1\).
This proves assertion 3.

When certification follows, CC-XC14 constructs the fixed rational target
and CC-XC15's forward-only search terminates. This is a termination
consequence of actual cycle existence, with no claimed degree or runtime
bound. The theorem does not assert an effective construction of the
auxiliary \(Z\) from the fibre-Chow hypothesis alone. \(\square\)

\subsection{Quadric-fibration consequences of
CC-XC18}\label{quadric-fibration-consequences-of-cc-xc18}

Suppose every geometric fibre of \(f\) is a quadric hypersurface of
dimension \(r\). Retain flatness, smooth projective \(X,B\), and
projectivity and surjectivity of \(f\); the fibres themselves need not
be smooth. The quadric case of Vial (2013, Corollary 4.4 and its proof,
p.~1535) supplies (CC-XC.Q2) with

\[
 n=\left\lceil\frac r2\right\rceil,\qquad
 \dim Z=
 \begin{cases}
 b,&r\text{ even},\\
 b-1,&r\text{ odd}.
 \end{cases}
 \tag{CC-XC.Q4}
\]

This primary result is stated over an arbitrary field, with rational
Chow groups. It therefore applies in characteristic two as well as in
odd characteristic. We use its quadric case and do not impose additional
discriminant or corank hypotheses.

The following are exact consequences.

\begin{longtable}[]{@{}
  >{\raggedright\arraybackslash}p{(\columnwidth - 4\tabcolsep) * \real{0.3333}}
  >{\raggedright\arraybackslash}p{(\columnwidth - 4\tabcolsep) * \real{0.3333}}
  >{\raggedright\arraybackslash}p{(\columnwidth - 4\tabcolsep) * \real{0.3333}}@{}}
\toprule\noalign{}
\begin{minipage}[b]{\linewidth}\raggedright
Base and relative dimension
\end{minipage} & \begin{minipage}[b]{\linewidth}\raggedright
Remaining Chow summand
\end{minipage} & \begin{minipage}[b]{\linewidth}\raggedright
Consequence for the original \(X\)
\end{minipage} \\
\midrule\noalign{}
\endhead
\bottomrule\noalign{}
\endlastfoot
\(b=1\), any \(r\ge1\) & A summand of a curve motive if \(r\) is even,
or of a finite étale motive if \(r\) is odd &
\(h(X)\in\mathcal C_{\mathrm{ab,Art}}(k)\); its forward-cycle search
terminates. \\
\(b=2\), \(r\) odd & A twisted summand of a curve motive & Chow
membership of \(X\) is equivalent to that of \(B\); certification of
\(X\) is equivalent to that of \(B\). \\
\(b=2\), \(r\) even & One specified twisted projector summand of a
surface motive & The exact remaining gate is membership of \(h(B)\) and
of \((Z,e)\), not necessarily of the whole surface \(Z\). \\
General \(b\ge1\) & Dimension at most \(b\), and \(b-1\) for odd \(r\) &
The high-dimensional forward-cycle problem reduces to the base and this
single residual projector. \\
\end{longtable}

\textbf{Verification of these consequences.} For an integer \(j\),
\(j<r/2\) is equivalent to \(0\le j<\lceil r/2\rceil\) once the lower
bound \(j\ge0\) is included. Substituting this \(n\) into
\(D-2n=b+r-2\lceil r/2\rceil\) gives (CC-XC.Q4). Because \(b\ge1\), the
residual dimension is nonnegative and the range imposed in CC-XC18 is
satisfied. The first two rows now follow from assertion 3 and the
membership of curves. The third and fourth rows are precisely assertions
1 and 2, retaining the projector \(e\) instead of replacing it by an
unwarranted statement about all of \(h(Z)\).

In the first row, \(h(B)\) and the residual motive are generated by
curves and Artin motives. Their twists in (CC-XC.Q3) therefore prove
actual Chow inclusion, not just Kimura finite-dimensionality. In the
second row, if \(B\) is only known to be certified, the conclusion is
certification of \(X\); a Chow-membership assertion for \(B\) is not
inserted. Examples with a Chow abelian--Artin base include an abelian
surface and the universal-zero-cycle surfaces of CC-XC17. \(\square\)

\subsection{Scope of the geometric
input}\label{scope-of-the-geometric-input}

The theorem uses all geometric fibres; we have not replaced that
condition by a condition on only the generic fibre. Vial explicitly
distinguishes these hypotheses in the introduction, pp.~1523--1524,
including the change of Chow groups in singular fibres of a conic
bundle. In the quadric case the required lower Chow groups remain
one-dimensional for the entire allowed fibre class, as encoded in
Corollary 4.4; this is why degenerations cause no extra hypothesis in
the conclusion above.

The residual surface in the even-dimensional quadric case is left
visible. A Chow--Künneth decomposition or Kimura finite-dimensionality
alone is not used as a substitute for abelian-sector membership. Nor do
the above statements claim that numerical and rational equivalence agree
on all products of \(X\). The conclusions needed for cote 15 are the
stated Chow inclusion or the exact lower-dimensional membership
reduction, followed by the existing finite forward-cycle search.

\subsection{Primary bibliography and
locators}\label{primary-bibliography-and-locators}

\begin{itemize}
\tightlist
\item
  Vial, C. (2013). Algebraic cycles and fibrations. \emph{Documenta
  Mathematica, 18}, 1521--1553. \url{https://doi.org/10.4171/DM/435} .
  Publisher PDF:
  \url{https://ems.press/content/serial-article-files/26216?nt=1} . Exact
  locators: standing hypotheses and rational coefficients, pp.~1522 and
  1525--1526; Theorem 1.4 and Proposition 1.5, pp.~1528--1529; Theorem
  4.2 and its proof, p.~1534; Corollary 4.4 and its proof, p.~1535. The
  publisher version was received 23 March 2012 and revised 19 July 2013.
  Some circulated files have obsolete pagination; the locators here
  refer to the publisher's 1521--1553 version.
\item
  Vial, C. (2017). Remarks on motives of abelian type. \emph{Tohoku
  Mathematical Journal, Second Series, 69}(2), 195--220.
  \url{https://doi.org/10.2748/tmj/1498269623} . Exact input: p.~195,
  generation by curves and abelian varieties, together with the
  finite-field Chow descent proved in CC-XC16. Published source:
  \url{https://www.jstage.jst.go.jp/article/tmj/69/2/69_1498269623/_pdf/-char/en}
  .
\end{itemize}

% END consequences_criteria_quadric_fibrations.tex

% BEGIN consequences_criteria_quadric_equivalence.tex
\section{The published equivalence comparison for quadric
fibrations}\label{the-published-equivalence-comparison-for-quadric-fibrations}

\subsection{CC-XC19 --- Rational and numerical equivalence on the total
space}\label{cc-xc19-rational-and-numerical-equivalence-on-the-total-space}

Let \(f:X\to B\) be a flat dominant morphism of smooth projective
varieties over a finite field \(k=\mathbf F_q\), with every geometric
fibre a quadric. Suppose \(\dim B\le2\). All cycle groups and
equivalence relations in this statement have rational coefficients.

\textbf{Published theorem.} Vial (2013, Corollary 4.8, p.~1536) proves
that rational and numerical equivalence agree on \(X\). Equivalently,
for every codimension \(j\), the quotient map

\[
  \mathrm{CH}^{j}(X)_{\mathbf Q}
       \longrightarrow
  \mathrm N^{j}(X)_{\mathbf Q}
\tag{CC-XC.Q5}
\]

is an isomorphism. Here \(\mathrm N^j\) denotes cycles modulo numerical
equivalence. The hypotheses include singular quadric fibres; smoothness
is required of the total space and base, not of the fibration. There is
no restriction to odd characteristic in the stated theorem.

This is an application of a published result, not a new proof of that
result. The theorem and its complete printed proof have been inspected
in the publisher's version. Its proof uses the Chow decomposition of
Corollary 4.4 and the vanishing of numerically trivial rational Chow
groups in the relevant degrees on varieties of dimension at most two
over a finite field.

\textbf{Scope corollary.} For every finite extension \(k'/k\), every
prime \(\ell\ne\operatorname{char}k\), and every \(j\), rational,
\(\ell\)-adic homological and numerical equivalence agree on the
rational cycles of \(X_{k'}\). In particular, the standard conjecture
\(D\), that homological and numerical equivalence agree, holds on
\(X_{k'}\), and the cycle-class map on
\(\mathrm{CH}^{j}(X_{k'})_{\mathbf Q}\) is injective.

\textbf{Complete deduction of the scope corollary.} Flatness, dominance,
projectivity and smoothness are preserved by the finite field extension,
and the geometric fibres remain quadrics. Apply Vial's Corollary 4.8 to
\(f_{k'}\). Rational equivalence implies homological equivalence.
Homological equivalence implies numerical equivalence, because the
cycle-class map respects intersection products and degrees. These two
implications, together with equality of the first and third equivalence
relations, force equality of all three. Equivalently, the homologically
trivial subgroup lies in the numerically trivial subgroup of the
rational Chow group, which is zero by the published theorem. This proves
the asserted injectivity and \(D\) for every indicated \(\ell\).
\(\square\)

The conclusion concerns cycles on the specified total space and its
finite field extensions. Corollary 4.8 does not assert this comparison
on every power \(X^m\), or on every product with another variety. Those
products do not automatically satisfy its hypothesis of being a quadric
fibration over a base of dimension at most two. The result also does not
assert surjectivity of the cycle-class map onto Frobenius-invariant
cohomology, which is the separate Tate existence problem.

For cote 15, this closes the equivalence-comparison question on each
such \(X\), while retaining the exact motive-generation boundary in
CC-XC18. In particular, over a surface base the present comparison does
not supply abelian-sector membership of the base or of the residual
surface projector in the even-dimensional quadric case. Actual
membership of the total space still has precisely the conditions proved
in CC-XC18.

\subsection{Primary reference}

Vial, C. (2013). Algebraic cycles and fibrations. \emph{Documenta
Mathematica, 18}, 1521--1553. \url{https://doi.org/10.4171/DM/435} . Publisher
PDF: \url{https://ems.press/content/serial-article-files/26216?nt=1} . Exact
locator: Corollary 4.8 and its complete proof, p.~1536;
rational-coefficient convention, p.~1522; preceding decomposition,
Corollary 4.4, p.~1535. Received 23 March 2012; revised 19 July 2013.
These locators use the published 1521--1553 pagination.

% END consequences_criteria_quadric_equivalence.tex

% BEGIN consequences_criteria_supersingular_cycles.tex
\section{Geometrically Tate Chow motives and supersingular K3 cycle
certificates}\label{geometrically-tate-chow-motives-and-supersingular-k3-cycle-certificates}

\textbf{Scope.} The geometric input below is an established theorem
about actual rational Chow motives. The finite-field deductions retain
the Galois action through Artin motives. They require neither a general
Tate conjecture nor the added all-abelian Tate input used elsewhere in
the programme.

\subsection{CC-XC20 --- Descent of a geometric Tate
decomposition}\label{cc-xc20-descent-of-a-geometric-tate-decomposition}

Let \(X/\mathbf F_q\) be smooth projective of pure dimension \(d\), and
suppose that in rational Chow motives

\[
 h(X_{\bar{\mathbf F}_q})
       \simeq\bigoplus_{i=0}^{d}\mathbf 1(-i)^{b_{2i}}.
\tag{CC-XC.T1}
\]

Then there are rational Artin motives \(A_i\), factoring through one
finite cyclic Galois group, such that

\[
                  h(X)\simeq\bigoplus_{i=0}^{d} A_i(-i).
\tag{CC-XC.T2}
\]

Their finite-image rational Galois representations are
\(\mathrm{CH}^i(X_{\bar{\mathbf F}_q})_{\mathbf Q}\); their dimensions
are \(b_{2i}\). For every \(\ell\ne p\), the realization of \(A_i\) is
\(H^{2i}(X_{\bar{\mathbf F}_q},\mathbf Q_\ell(i))\), and all odd
cohomology vanishes. The finite-order Frobenius representations
determine the \(A_i\) up to rational Artin-motive isomorphism.

Every finite product of varieties satisfying (CC-XC.T1), over every
finite extension of \(\mathbf F_q\), has the following properties:

\begin{enumerate}
\def\labelenumi{\arabic{enumi}.}
\tightlist
\item
  Its rational Chow motive is Artin--Tate and hence belongs to the
  abelian--Artin category of CC-XC16.
\item
  Rational, \(\ell\)-adic homological and numerical equivalence agree in
  every codimension.
\item
  The rational cycle-class map, after tensoring with \(\mathbf Q_\ell\),
  is an isomorphism onto the Frobenius-invariant part of the
  corresponding Tate-twisted even cohomology.
\item
  The fixed-target forward-cycle search in CC-XC15 terminates.
\end{enumerate}

\textbf{Complete proof.} Choose an isomorphism in (CC-XC.T1) and its
inverse. A cycle is a finite sum of integral subschemes, and a
rational-equivalence relation has finitely many subvarieties and
rational functions as witnesses. Thus both correspondences and the two
identities between their compositions descend to one finite extension
\(k'=\mathbf F_{q^m}\). We obtain an actual Chow isomorphism

\[
                     h(X_{k'})\simeq
                         \bigoplus_i\mathbf1_{k'}(-i)^{b_{2i}}.
\tag{CC-XC.T3}
\]

Use the induction of Chow correspondences proved in CC-XC16: regard a
\(k'\)-scheme as a \(k\)-scheme, and extend a correspondence by zero
from the open-and-closed component defined by the \(k'\)-fibre product.
Its composition law is block diagonal after a separable splitting
extension, and rational restriction--corestriction makes this an
identity over \(k\). Applying induction to (CC-XC.T3) gives

\[
 h(X_{k'}\text{ as a }k\text{-scheme})
      \simeq\bigoplus_i h(\operatorname{Spec}k')(-i)^{b_{2i}}.
\tag{CC-XC.T4}
\]

The projection \(\tau:X_{k'}\to X\) satisfies
\(\tau_*\tau^*=m\,\mathrm{id}\) as a Chow correspondence. Dividing the
transfer by \(m\) exhibits \(h(X)\) as a direct summand of the
right-hand side of (CC-XC.T4). No numerical or homological replacement
of the Chow identity is used.

For finite étale \(k\)-schemes \(U,V\), the groups of Chow
correspondences between \(h(U)(-i)\) and \(h(V)(-j)\) vanish when
\(i\ne j\), since the product has dimension zero and has nonzero Chow
groups only in codimension zero. Consequently the idempotent cutting out
\(h(X)\) is block diagonal in the twist index \(i\), and its \(i\)-th
block defines the Artin motive \(A_i\). This proves (CC-XC.T2).

For completeness, the category of rational Artin motives split by \(k'\)
identifies with finite-dimensional rational representations of the
cyclic group \(G=\operatorname{Gal}(k'/k)\). The motive of a finite
étale scheme gives its permutation representation. Correspondences are
rational matrices constant on \(G\)-orbits of pairs of geometric points,
exactly the equivariant maps between these permutation representations.
Taking direct sums and idempotent images gives all finite-dimensional
representations: every such representation is a summand of a finite sum
of regular representations by Maschke's averaging argument. This also
proves faithfulness of the Artin realization. Finite-order rational
operators are semisimple, and their cyclotomic characteristic-polynomial
multiplicities determine their rational representations.

Over \(\bar k\), maps from a Tate summand to a different Tate summand
vanish, and its degree-\(2i\) cycle class is an isomorphism. Applying
this to (CC-XC.T1) proves

\[
 \mathrm{CH}^i(X_{\bar k})_{\mathbf Q}\otimes\mathbf Q_\ell
       \simeq H^{2i}(X_{\bar k},\mathbf Q_\ell(i)).
\tag{CC-XC.T5}
\]

The identification is Galois equivariant, so it identifies \(A_i\) as
asserted. In particular its action has finite image: a basis of the
finite-dimensional rational Chow group is defined over a common finite
extension.

The condition (CC-XC.T1) is preserved by products and finite extensions.
It remains to justify the cycle assertions for any one such product
\(Y/k\). Rational restriction to \(\bar k\) is injective. Indeed, an
equality of cycles modulo rational equivalence over \(\bar k\) descends
to a finite extension, where restriction followed by corestriction is
multiplication by its degree.

Moreover,

\[
 \mathrm{CH}^j(Y)_{\mathbf Q}
    =\mathrm{CH}^j(Y_{\bar k})_{\mathbf Q}^{\operatorname{Gal}(\bar k/k)}.
\tag{CC-XC.T6}
\]

To check surjectivity, represent an invariant geometric class over a
finite Galois extension \(k''/k\). Its invariance holds already in
\(\mathrm{CH}^j(Y_{k''})_{\mathbf Q}\), by the injectivity just proved.
The corestriction divided by \([k'':k]\) restricts to that class. The
actions on the finite-dimensional geometric Chow groups factor through a
finite group, so taking invariants commutes with extension of scalars to
\(\mathbf Q_\ell\). Formula (CC-XC.T5), applied to \(Y\), now proves the
asserted cycle-class isomorphism onto invariants.

The intersection pairing on the geometric rational Chow groups in
complementary codimensions is perfect: (CC-XC.T5) identifies its scalar
extension with Poincaré duality, and its matrix has rational entries. It
remains perfect on invariants. In fact, if a nonzero invariant vector
pairs nontrivially with some vector, averaging the latter over the
finite acting group preserves that nonzero pairing and produces an
invariant vector. By (CC-XC.T6), this is the numerical pairing over
\(k\). Its kernel is zero, proving rational equals numerical; the usual
rational--homological--numerical implications then give the middle
equality as well.

Finally Artin--Tate motives are contained in the abelian--Artin
category, and all the isomorphisms used above are given by rational
cycles. Thus the forward-cycle existence condition of CC-XC15 is
satisfied for \(Y\), and that search terminates. No a priori bound for
its termination time is asserted. \(\square\)

\subsection{CC-XC21 --- Supersingular K3 surfaces, their Hilbert schemes
and residual
projectors}\label{cc-xc21-supersingular-k3-surfaces-their-hilbert-schemes-and-residual-projectors}

Let \(S/\mathbf F_q\) be a smooth projective geometrically connected K3
surface which is geometrically supersingular. One can use the definition
\(\rho(S_{\bar{\mathbf F}_q})=22\); equivalence with the formal-Brauer
or slope-one definition is the established supersingular Tate theorem,
recorded in Fu--Li (2021, pp.~148 and 169).

The published motivic input is Fu--Li (2021, Theorem 1.3(iii),
pp.~152--153; proof in §4.5, p.~179). It implies

\[
 h(S_{\bar k})\simeq
       \mathbf1\oplus\mathbf1(-1)^{22}\oplus\mathbf1(-2).
\tag{CC-XC.T7}
\]

The same theorem gives a Tate Chow motive for a moduli space
\(M_H(S_{\bar k},v)\), where \(v=(r,c,s)\) is a \textbf{primitive} Mukai
vector, \(r>0\), \(v^2\ge0\), \(H\) is general for \(v\), and \[
             p\nmid r\quad\text{or}\quad
             p\nmid s\quad\text{or}\quad
             p\nmid(c\cdot\mathrm{NS}(S_{\bar k})).
\tag{CC-XC.T8}
\] Here ``general'' means that every semistable sheaf in question is
stable; the last expression means that at least one intersection
\(c\cdot D\) is not divisible by \(p\). These are the source's
Definitions 4.3 and 4.7, pp.~171 and 173. We include primitivity
explicitly from its standing moduli-space convention. A finite-field
form of such a smooth projective moduli space is allowed.

In particular, for every \(n\ge1\), the Hilbert scheme \(S^{[n]}\) has
geometric Tate Chow motive. Indeed it is the rank-one ideal-sheaf case
\(v=(1,0,1-n)\): this vector is primitive, satisfies (CC-XC.T8), has
\(v^2=2n-2\), and its rank-one ideal sheaves are stable for every ample
polarization. This agrees with the Hilbert-scheme step of the published
proof, p.~179.

\textbf{Complete finite-field deduction.} Apply CC-XC20 to each of these
varieties. Thus their actual rational Chow motives over \(\mathbf F_q\)
are Artin--Tate, all of their finite products satisfy the cycle
comparisons and Tate-surjectivity statements of CC-XC20, and their
fixed-target forward-cycle searches terminate. This applies in every
positive characteristic; no assumption of unirationality is made.

For \(S\) itself, the degree-zero and top-degree Artin factors are
trivial by geometric connectedness and the degree map. The middle factor
is exactly the rational Néron--Severi representation. If
\(A_{\mathrm{NS}}\) denotes its rational Artin motive, then

\[
                h(S)\simeq
                  \mathbf1\oplus A_{\mathrm{NS}}(-1)
                              \oplus\mathbf1(-2).
\tag{CC-XC.T9}
\]

After a finite extension defining all divisor classes, this becomes the
Tate decomposition (CC-XC.T7) over that finite field. This conclusion is
a Chow-motive isomorphism, not merely equality of characteristic
polynomials.

For the even-dimensional quadric fibration over a surface in CC-XC18,
suppose that its supplied residual surface \(Z\) is a supersingular K3
surface, or that its specified projector summand \((Z,e)\) is known by
Chow correspondences to be a summand of a finite sum of twisted motives
of the varieties just considered. Then \((Z,e)\) is Artin--Tate. The
residual-projector gate in CC-XC18 is discharged, and the total space
has Chow abelian--Artin membership exactly when its base does; the
analogous statement holds for homological certification. If the base is
itself geometrically Tate, then the total space is Artin--Tate and the
stronger all-products assertions of CC-XC20 apply to it. This
conditional identification of a residual projector does not assert that
the auxiliary surface of every even-dimensional quadric fibration is K3.
\(\square\)

\subsection{Primary-source scope and dependency
record}\label{primary-source-scope-and-dependency-record}

The K3 input is a theorem about rational Chow motives, not just the
algebraicity of \(H^2\). The proof in Fu--Li, §4.5, uses twisted derived
equivalences to reduce Artin invariant and therefore does not rely on
general supersingular unirationality. The corresponding current
Bragg--Lieblich locator is Proposition 5.2.5, p.~95, in
arXiv:1804.07282v6 (9 February 2019), in its \(p\ge3\) range. Fu--Li's
printed proof uses an older locator ``Proposition 5.15''; we do not copy
that locator as a reference to v6. The characteristic-two case is the
separate known unirational case discussed on their p.~148. Fu--Vial
(2021, Theorem 1, p.~3) supplies actual rational Chow isomorphisms for
twisted-derived-equivalent K3 surfaces over an arbitrary field. These
are inspected published or preprint inputs; their full foundational
proofs are not claimed to have been reconstructed here.

No conclusion is drawn for a general finite-height K3 surface from its
Tate theorem for divisors or from a theorem about its self-product
alone. No unirationality conclusion for Artin invariant ten is asserted.
The new finite-field deductions above use the precise geometric Chow
input that is stated.

\subsection{Primary bibliography}\label{primary-bibliography}

\begin{itemize}
\tightlist
\item
  Fu, L., \& Li, Z. (2021). Supersingular irreducible symplectic
  varieties. In G. Farkas, G. van der Geer, M. Shen, \& L. Taelman
  (Eds.), \emph{Rationality of varieties} (Progress in Mathematics, Vol.
  342, pp.~147--200). Birkhäuser.
  \url{https://doi.org/10.1007/978-3-030-75421-1_7} . Author-hosted published
  PDF:
  \url{https://irma.math.unistra.fr/~lfu/articles/Supersingular%20ISV.pdf}
  . Exact locators: Theorem 1.3(iii), pp.~152--153; Definitions 4.3 and
  4.7, pp.~171 and 173; §4.5, full concluding proof, p.~179;
  supersingularity conventions, pp.~148 and 169.
\item
  Fu, L., \& Vial, C. (2021). A motivic global Torelli theorem for
  isogenous K3 surfaces. \emph{Advances in Mathematics, 383}, 107674.
  \url{https://doi.org/10.1016/j.aim.2021.107674} . Published PDF:
  \url{https://irma.math.unistra.fr/~lfu/articles/DerivedEqK3.pdf}
  . Exact locator: Theorem 1, p.~3, over any field.
\item
  Bragg, D., \& Lieblich, M. (2019). \emph{Twistor spaces for
  supersingular K3 surfaces} (arXiv:1804.07282v6).
  \url{https://arxiv.org/abs/1804.07282v6} . Exact locator: Proposition 5.2.5,
  p.~95; standing \(p\ge3\) range in §1.1 and §3. This is the 9 February
  2019 version, whose notice removes the earlier material claiming
  general unirationality.
\item
  Vial, C. (2017). Remarks on motives of abelian type. \emph{Tohoku
  Mathematical Journal, Second Series, 69}(2), 195--220.
  \url{https://doi.org/10.2748/tmj/1498269623} . Lemma 1.2, p.~200, and the
  finite-field induction proof of CC-XC16 justify the rational Chow
  descent used above.
\end{itemize}

% END consequences_criteria_supersingular_cycles.tex

% BEGIN consequences_criteria_supersingular_liftings.tex
\section{Supersingular abelian generators and the precisely scoped 2026
lifting
theorem}\label{supersingular-abelian-generators-and-the-precisely-scoped-2026-lifting-theorem}

\subsection{CC-XC22 --- Finite-field consequences of geometric
supersingular abelian Chow
type}\label{cc-xc22-finite-field-consequences-of-geometric-supersingular-abelian-chow-type}

Let \(X/\mathbf F_q\) be smooth projective of pure dimension. Suppose
that \(h(X_{\bar{\mathbf F}_q})\) belongs to the rational Chow category
generated by supersingular abelian varieties and Tate motives, closed
under sums, tensor products, duals and direct summands. Then:

\begin{enumerate}
\def\labelenumi{\arabic{enumi}.}
\tightlist
\item
  \(h(X)\) belongs to the rational Chow abelian--Artin category over
  \(\mathbf F_q\), so its fixed-target forward-cycle search terminates.
\item
  For \(X\), every finite product of varieties satisfying this same
  hypothesis, and every finite extension of the base field, rational,
  \(\ell\)-adic homological and numerical equivalence agree in every
  codimension.
\item
  The rational cycle-class map, tensored with \(\mathbf Q_\ell\), is an
  isomorphism onto Frobenius-invariant even cohomology with the
  corresponding Tate twist, for every \(\ell\ne p\).
\item
  Frobenius acts semisimply in every cohomological degree.
\end{enumerate}

These conclusions do not require the general or all-abelian Tate
conjecture. If odd cohomology is nonzero, the conclusion is not that
\(h(X)\) is Artin--Tate.

\textbf{Complete deduction.} The geometric Chow hypothesis supplies a
finite presentation as a summand of a finite sum of twisted
abelian-variety motives: expand tensor products as products of abelian
varieties, and use their usual Chow duality for the dual factors. The
projection, inclusion and their rational-equivalence witnesses descend
to a common finite field extension. The finite-field induction and
trace-retraction construction of CC-XC16 then puts \(h(X)\) in the
actual Chow abelian--Artin category over \(\mathbf F_q\). This proves
the asserted cycle existence and search termination.

For the stronger comparison assertions, use the more precise published
result of Fu--Li (2021, Corollary 2.12(i), (ii), (iv), p.~160; proof
pp.~160--161). Over the algebraic closure, for a variety \(Y\)
satisfying the hypothesis, it gives

\[
 h(Y_{\bar k})\simeq
 \bigoplus_i\mathbf1(-i)^{b_{2i}(Y)}
 \oplus
 \bigoplus_i h^1(E)(-i)^{b_{2i+1}(Y)/2},
\tag{CC-XC.U1}
\]

where \(E/\bar k\) is a supersingular elliptic curve. This is a
Chow-motive decomposition, and the cited corollary identifies the
algebraically trivial part of each rational Chow group with a finite sum
of copies of \(E(\bar k)\otimes_{\mathbf Z}\mathbf Q\).

An elliptic curve over \(\bar{\mathbf F}_q\) is defined over a finite
extension of \(\mathbf F_q\), and each of its geometric points is
rational over a further finite extension. Every such point therefore
lies in a finite group. Hence

\[
               E(\bar{\mathbf F}_q)\otimes_{\mathbf Z}\mathbf Q=0.
\tag{CC-XC.U2}
\]

It follows from (CC-XC.U1) and the cited corollary that

\[
 \mathrm{CH}^j(Y_{\bar k})_{\mathbf Q}\otimes\mathbf Q_\ell
          \simeq H^{2j}(Y_{\bar k},\mathbf Q_\ell(j)).
\tag{CC-XC.U3}
\]

In particular, the geometric rational Chow group is finite-dimensional
and has a finite-image Galois action: finitely many cycles forming a
basis are defined over one finite extension.

The rational restriction--corestriction argument in CC-XC20 gives \[
 \mathrm{CH}^j(Y)_{\mathbf Q}
     \simeq \mathrm{CH}^j(Y_{\bar k})_{\mathbf Q}^{\operatorname{Gal}(\bar k/k)}.
\tag{CC-XC.U4}
\] Taking invariants commutes with tensoring by \(\mathbf Q_\ell\) for
this finite group action. Applying it to (CC-XC.U3) proves assertion 3.
The even-degree intersection pairing is perfect on geometric rational
Chow groups by (CC-XC.U3) and Poincaré duality. Averaging a
complementary class over the finite Galois group proves that it remains
perfect on invariants. This makes the numerical kernel of the rational
Chow group zero; the rational--homological--numerical implications give
assertion 2. Tensor products and finite base extensions preserve the
geometric supersingular abelian Chow hypothesis, so this argument
applies to all the products specified in the statement.

Finally the Chow presentation over \(\mathbf F_q\) from the first
paragraph realizes \(H^*(X)\) as a direct summand of a finite sum of
tensor products of abelian-variety cohomology, Artin representations and
Tate twists. Frobenius is semisimple on \(H^1\) of an abelian variety:
Tate (1966, Theorem 2(a), p.~140) identifies \(\mathbf Q[\pi]\) with the
centre of its semisimple rational endomorphism algebra, so the minimal
polynomial of \(\pi\) is squarefree. Abelian cohomology is the exterior
algebra on \(H^1\), and exterior products of a diagonalizable operator
are diagonalizable. Finite-order Artin actions, Tate twists, tensor
products, sums and direct summands preserve semisimplicity. This proves
assertion 4.

The fixed-target algorithm can also be run entirely in this
supersingular sector. In degree \(i>0\), every Frobenius eigenvalue has
all normalized crystalline slopes \(i/2\), by the geometric Chow
presentation. Each weight-one Honda root used in CC-XC11/14 therefore
has all slopes \(1/2\); its Honda abelian variety is supersingular.
Degree zero uses only Artin factors. The target construction, its
duality, and the required numerical Hom calculations thus use only
products of supersingular abelian varieties and finite étale schemes.
Assertions 2 and 3 just proved supply the required Tate surjectivity and
numerical--homological equality on these products. This verifies search
termination without invoking the all-abelian input. \(\square\)

\textbf{Source-consistency note.} The explicit multiplicities printed in
Fu--Li's Lemma 2.8, p.~158, contain an error: its claimed rank ten for
\(\operatorname{Sym}^4 h^1(E)\) cannot be a direct summand of the stated
rank-nine \((\operatorname{Sym}^2 h^1(E))^{\otimes2}\). In the source's
convention, in which \(\operatorname{Sym}^2 h^1(E)\) has rank three, the
correct even and odd multiplicities are \(2j+1\) and \(j+1\),
respectively. Its qualitative argument still gives a sum of Tate motives
in even degree and a sum of twists of \(h^1(E)\) in odd degree. Theorem
2.9 and Corollary 2.12 then obtain the Betti-number multiplicities used
in (CC-XC.U1) directly from realization; the incorrect Lemma 2.8 counts
are not used here.

\subsection{CC-XC23 --- Actual cycle generators from the 2026 lifting
theorem}\label{cc-xc23-actual-cycle-generators-from-the-2026-lifting-theorem}

\textbf{External input and version.} Fu--Huang--Li's \emph{Supersingular
Tate conjecture for irreducible symplectic varieties of known types: I},
arXiv:2609.38801v2, was submitted on \textbf{20 October 2026 at 07:58:38
UTC}, after v1 of 30 September 2026. The theorem cited here is
\textbf{Theorem 1.9, p.~5} of v2. The downloaded v2 text agrees with the
author-hosted 59-page version after removing only the arXiv header and
whitespace.

Let \(X/\mathbf F_q\) be smooth projective of dimension \(d\). Apply the
following hypotheses to \(X_{\bar k}\):

\begin{itemize}
\tightlist
\item
  It is irreducible symplectic in the source's sense: its geometric
  étale fundamental group is trivial and its global two-forms are
  spanned by a closed nowhere-degenerate two-form (Definition 2.1,
  p.~7).
\item
  It is \textbf{excellent}: there is a smooth projective lifting over a
  finite extension DVR of \(W(\bar k)\), with irreducible-symplectic
  geometric generic fibre, all Hodge numbers unchanged, and a relatively
  ample line bundle whose restriction is a primitive polarization of
  \(X_{\bar k}\). This is Definitions 1.6, p.~4, and 2.2/2.6, pp.~8--9.
\item
  It is Artin-supersingular, meaning that its formal Brauer group is the
  formal additive group. In this excellent setting the condition is
  equivalent to slope one for its second crystalline cohomology
  (Definition 2.14 and Proposition 2.15, p.~11).
\item
  \(p>d\).
\end{itemize}

Under these hypotheses the inspected preprint gives the following
geometric Chow generators:

\begin{longtable}[]{@{}
  >{\raggedright\arraybackslash}p{(\columnwidth - 6\tabcolsep) * \real{0.2500}}
  >{\raggedright\arraybackslash}p{(\columnwidth - 6\tabcolsep) * \real{0.2500}}
  >{\raggedright\arraybackslash}p{(\columnwidth - 6\tabcolsep) * \real{0.2500}}
  >{\raggedright\arraybackslash}p{(\columnwidth - 6\tabcolsep) * \real{0.2500}}@{}}
\toprule\noalign{}
\begin{minipage}[b]{\linewidth}\raggedright
Deformation type of the excellent lift
\end{minipage} & \begin{minipage}[b]{\linewidth}\raggedright
Additional condition
\end{minipage} & \begin{minipage}[b]{\linewidth}\raggedright
Conclusion of Theorem 1.9 over \(\bar k\)
\end{minipage} & \begin{minipage}[b]{\linewidth}\raggedright
Finite-field consequence
\end{minipage} \\
\midrule\noalign{}
\endhead
\bottomrule\noalign{}
\endlastfoot
\(K3^{[n]}\) & None beyond the displayed hypotheses & Tate Chow motive &
CC-XC20: actual Artin--Tate motive over \(\mathbf F_q\), with
all-products cycle comparison and Tate surjectivity. \\
\(\mathrm{OG10}\) & \(\sigma_0\ne12\) & Tate Chow motive & The same
conclusion from CC-XC20. \\
\(\mathrm{OG6}\) & \(\sigma_0\ne4\) & Tate Chow motive & The same
conclusion from CC-XC20. \\
\(\mathrm{Kum}^{n}\) & \(\sigma_0\ne3\) & Supersingular abelian Chow
motive & CC-XC22: actual abelian--Artin Chow membership, all-products
cycle comparison, Tate surjectivity and semisimple Frobenius. Odd
cohomology need not vanish. \\
\end{longtable}

Here \(\sigma_0\) is the source's Artin invariant. It is specified by
the \(p\)-primary discriminant group
\((\mathbf Z/p\mathbf Z)^{2\sigma_0}\) of the crystalline Tate lattice
with its Beauville--Bogomolov form; see §5.1, Proposition 5.1, p.~26. It
is not being replaced by the Artin invariant of an unspecified auxiliary
surface.

\textbf{Complete deduction of the finite-field column.} In the first
three rows, the preprint supplies exactly the geometric Tate Chow input
of CC-XC20. In the last row it supplies exactly the geometric
supersingular abelian Chow input of CC-XC22. Apply those theorems. Their
proofs descend finite cycle data and actual splitting identities;
consequently they construct the required motive inclusion over
\(\mathbf F_q\) even when the lifting or the geometric generating
abelian varieties were initially given only over \(\bar k\). All four
rows therefore discharge the forward-cycle existence condition, and
CC-XC15's search terminates. Products and finite base extensions are
covered by the same preceding theorems. \(\square\)

The new geometric content invoked here is the preprint's generator
theorem; the finite-field deductions do not reprove it. Its final
motivic proofs have been inspected: Theorem 7.13, p.~44, for
\(K3^{[n]}\), \(\mathrm{Kum}^n\) and \(\mathrm{OG10}\); and Theorem 8.9,
p.~49, for \(\mathrm{OG6}\). Their route uses a modular lift,
supersingular good reduction of the underlying K3 or abelian surface,
and specialization of an actual split Chow-motive inclusion. The
specialization statement is Lemma 6.12, pp.~39--40, which explicitly
allows smooth proper algebraic-space models with smooth projective
generic and special fibres. We inspected these statements and final
proofs and the stated dependency chain; an independent reconstruction of
all 59 pages is not claimed.

The lifting conditions and excluded Artin invariants remain genuine
boundaries of this application. Mere equality of the second crystalline
Newton slopes does not supply an excellent lifting. Likewise a variety
of one of the four named deformation types without the displayed lifting
and characteristic hypotheses is not covered by this theorem. The three
excluded maximal invariants are stated limitations of the source's
method, not asserted counterexamples (Remark 1.13, p.~6).

This gives actual cycle generators for the original finite-field variety
in the listed cases, including varieties not initially presented as
moduli spaces. It neither assumes nor proves that every smooth
projective variety belongs to the abelian--Artin sector.

\subsection{Primary bibliography and inspected
locators}\label{primary-bibliography-and-inspected-locators}

\begin{itemize}
\tightlist
\item
  Fu, L., Huang, X., \& Li, Z. (2026). \emph{Supersingular Tate
  conjecture for irreducible symplectic varieties of known types: I}
  (arXiv:2609.38801v2). \url{https://arxiv.org/abs/2609.38801v2} . Author PDF:
  \url{https://irma.math.unistra.fr/~lfu/articles/Supersingular_Tate_conjecture.pdf}
  . Exact locators: Definitions 1.6/2.2/2.6, pp.~4, 8--9; Theorem 1.9,
  p.~5; Definition 2.14 and Proposition 2.15, p.~11; Proposition 5.1,
  p.~26; Lemma 6.12, pp.~39--40; Propositions 7.10--7.12 and Theorem
  7.13, pp.~43--44; Propositions 8.6--8.8 and Theorem 8.9, pp.~47--49.
  This is an inspected October 2026 preprint, not a reported journal
  publication.
\item
  Fu, L., \& Li, Z. (2021). Supersingular irreducible symplectic
  varieties. In G. Farkas, G. van der Geer, M. Shen, \& L. Taelman
  (Eds.), \emph{Rationality of varieties} (Progress in Mathematics, Vol.
  342, pp.~147--200). Birkhäuser.
  \url{https://doi.org/10.1007/978-3-030-75421-1_7} . Published PDF:
  \url{https://irma.math.unistra.fr/~lfu/articles/Supersingular%20ISV.pdf}
  . Exact input: Definition 2.10, p.~159; Corollary 2.12 and its proof,
  pp.~160--161.
\item
  Tate, J. (1966). Endomorphisms of abelian varieties over finite
  fields. \emph{Inventiones Mathematicae, 2}, 134--144.
  \url{https://doi.org/10.1007/BF01404549} . Exact input: Theorem 2(a),
  p.~140, for the centre \(\mathbf Q[\pi]\); primary scan:
  \url{https://pazuki.perso.math.cnrs.fr/index_fichiers/Tate66.pdf} .
\end{itemize}

% END consequences_criteria_supersingular_liftings.tex

% BEGIN root_quadric_discriminant.tex
\section{The actual discriminant cover in the even-quadric cycle gate}

\subsection{BR-15-ORBIT-RETRACT: recovering finite Chow retractions from an orbit isomorphism}
Let $\mathcal C$ be an additive category with an invertible additive
autoequivalence $M\mapsto M(1)$. Its orbit category has
\[
 \operatorname{Hom}_{\mathcal C/(1)}(M,N)
 =\bigoplus_{j\in\mathbf Z}\operatorname{Hom}_{\mathcal C}(M,N(j)),
\]
with the usual shifted composition. If $M$ and $N$ are isomorphic in
this orbit category, then $M$ is a retract in $\mathcal C$ of a finite
direct sum of twists of $N$, and conversely $N$ is a retract of a finite
direct sum of twists of $M$. Consequently $M$ and $N$ belong to the same
full replete additive subcategories closed under retracts and all twists.

\textbf{Proof.} Write an orbit isomorphism and its inverse as finite
families $u=(u_j)$ and $v=(v_k)$, with
$u_j:M\to N(j)$ and $v_k:N\to M(k)$. The component of twist zero in
$vu=1_M$ is the identity of actual morphisms in $\mathcal C$
\[
 \sum_j v_{-j}(j)u_j=1_M. \tag{BR-OD.1}
\]
Choose the finite set $J$ of indices with $u_j\ne0$, and put
$P=\bigoplus_{j\in J}N(j)$. The maps
\[
 a=(u_j)_{j\in J}:M\longrightarrow P,\qquad
 b=(v_{-j}(j))_{j\in J}:P\longrightarrow M
\]
satisfy $ba=1_M$ by (BR-OD.1). Thus $M$ is a retract of $P$. Interchange
$u$ and $v$ for the other retraction. Closure under finite sums, twists
and retracts proves the final assertion. All sums here are finite
because the orbit Hom is a direct sum. No conservativity of a realization,
nilpotence conjecture, or lifting of a numerical idempotent is used.
$\square$

\subsection{BR-15-QUADRIC-DISCRIMINANT: replacing the unspecified surface projector}
Let $k$ be a field of characteristic different from two. Let $B$ be a
smooth projective $k$-scheme, and let $f:Q\to B$ be the flat, generically
smooth quadric
fibration associated with a line-bundle-valued quadratic form on a
vector bundle of even rank $r+2$, where $r\ge2$ is even. Assume $Q$ is
smooth, every geometric fibre has corank at most one, and the
discriminant divisor $D\subset B$ is smooth. We allow an empty
discriminant divisor. Define the discriminant double cover
\[
 T=\operatorname{Spec}_B Z(\operatorname{Cl}_0(f)). \tag{BR-OD.2}
\]
This is a smooth projective $k$-scheme; it need not be geometrically
connected. For every full replete additive subcategory $\mathcal A$ of
rational Chow motives over $k$ that is closed under direct summands and
all Tate twists,
\[
 h(Q)\in\mathcal A
 \quad\Longleftrightarrow\quad
 h(B)\in\mathcal A\ \hbox{and}\ h(T)\in\mathcal A. \tag{BR-OD.3}
\]
In particular this holds for the Chow category generated by abelian
varieties and Artin motives. After algebraic closure, it also holds for
the subcategory of Tate motives and for the subcategory generated by
supersingular abelian varieties. This identifies a concrete geometric
test in the even-relative-dimensional case of CC-XC18.
The same equivalence holds for membership in a full additive,
summand- and twist-closed category of rational $\ell$-adic homological
motives, for any invertible prime $\ell$. In particular it identifies
the original homological abelian--Artin forward-cycle gate with the
two gates for $B$ and the specified cover $T$.

\textbf{Exact published inputs.} We use rational coefficients throughout.
Bernardara--Marcolli--Tabuada (2018, Proposition 4.1, pp.~306--307,
diagram (4.2)) give a fully faithful additive functor
\[
 R:\operatorname{Chow}(k)_{\mathbf Q}/(-\otimes\mathbf Q(1))
   \longrightarrow\operatorname{NChow}(k)_{\mathbf Q}.
 \tag{BR-OD.4}
\]
It takes the orbit of $h(V)$ to the noncommutative Chow motive
$U(\operatorname{perf}_{\rm dg}V)$ of a smooth projective scheme $V$;
see their p.~307, proof of Theorem 1.1. Their equation (6.1) and the
proof of Theorem 1.2(ii), p.~308, give in this projective setting
\[
 U(Q)\simeq U(T)\oplus U(B)^{\oplus r}.
 \tag{BR-OD.5}
\]
Here the first decomposition comes from the Clifford component of the
quadric fibration; the even Clifford algebra becomes Azumaya on $T$;
and its rational noncommutative motive agrees with that of $T$.
The Azumaya step is Tabuada--Van den Bergh (2013, Theorem 2.1,
arXiv:1307.7946v2, p.~3): the required ranks are invertible in
$\mathbf Q$.

The field hypothesis deserves an explicit source qualification.
Kuznetsov's original 2008 paper assumes an algebraically closed field
of characteristic zero in its \S2, p.~1343. The general-field
quadric semiorthogonal decomposition used here is stated and justified
by Auel--Bernardara--Bolognesi (2014, Theorem 2.2.1 and its proof,
p.~270), which expressly removes that restriction. Tabuada (2020,
Proposition 15(i), p.~2346, complete proof on that page) gives the same
decomposition after inverting two and describes the smooth discriminant
cover, under the explicit simple-degeneration hypotheses above. We do
not transfer a characteristic-zero-only theorem to finite fields by
analogy.

\textbf{Complete deduction.} Put $N=h(T)\oplus h(B)^{\oplus r}$.
Equation (BR-OD.5) and full faithfulness of (BR-OD.4) lift to an
isomorphism $h(Q)\simeq N$ in the Tate orbit category: lift each
direction by fullness, and use faithfulness to lift the two inverse
identities. BR-15-ORBIT-RETRACT now gives actual Chow-motive
retractions
\[
 h(Q)\ \mid\ \bigoplus_{j\in J}N(j),\qquad
 N\ \mid\ \bigoplus_{k\in K}h(Q)(k) \tag{BR-OD.6}
\]
for finite sets $J,K$, where $M\mid P$ denotes a specified split
inclusion. If $h(B)$ and $h(T)$ belong to $\mathcal A$, the first
retraction puts $h(Q)$ in $\mathcal A$. Conversely the second puts
$N$ in $\mathcal A$, and projection onto its displayed summands gives
both $h(B),h(T)\in\mathcal A$; here $r\ge2$ ensures that a copy of
$h(B)$ is present. This proves (BR-OD.3). These are actual rational
Chow correspondences with identities modulo rational equivalence.
The original degree-zero-component argument does not specify a
particular graded direct-sum formula for $h(Q)$, and none is needed.

All the hypotheses, the Clifford centre, and the quoted decompositions
commute with field extension. Apply the same proof over $\bar k$ to
obtain the last two assertions. The subcategories in question are
closed under finite sums, summands, and twists by their definitions.
For the homological assertion, let $\mathcal A$ be the full preimage
of the specified homological subcategory under the rational Chow
realization functor. That preimage has all the closure properties used
in (BR-OD.3), since a functor preserves a displayed retraction. Thus
(BR-OD.3) applies to it. This argument assumes no equality between
numerical and homological equivalence for $B$, $T$, or $Q$.
$\square$

\subsection{BR-15-QUADRIC-TATE: finite-field cycle consequences and a concrete K3 case}
Suppose now that $k=\mathbf F_q$ has odd characteristic and the
preceding hypotheses hold. If $B_{\bar k}$ and $T_{\bar k}$ have Tate
Chow motives, then $h(Q)$ is Artin--Tate over $k$. On all finite
products of such total spaces, and after every finite base extension,
rational, homological and numerical equivalence coincide; the
rational cycle-class map tensored with $\mathbf Q_\ell$ is an
isomorphism onto the Frobenius-invariant even cohomology; and odd
cohomology vanishes. If the two geometric generators are instead of
supersingular abelian Chow type, the corresponding conclusions of
CC-XC22 hold, including semisimple Frobenius, without an assertion that
odd cohomology vanishes. In either case the fixed-target forward-cycle
search terminates.

\textbf{Proof.} BR-15-QUADRIC-DISCRIMINANT puts the geometric motive of
$Q$ in the same stated subcategory as its two generators. Apply
CC-XC20 or CC-XC22, respectively, whose proofs give actual finite-field
Chow splittings and all-product comparisons. This proves every
assertion, without the all-abelian 2026 Tate input in the supersingular
cases. $\square$

For example, take $B=\mathbf P^2_{\mathbf F_q}$ and suppose the
actual smooth discriminant double cover $T$ is a geometrically
supersingular K3 surface. The projective bundle decomposition gives
$h(B)=\mathbf1\oplus\mathbf1(-1)\oplus\mathbf1(-2)$, and the
published supersingular K3 result used in CC-XC21 makes $h(T_{\bar k})$
Tate. Therefore all the preceding Tate conclusions hold for $Q$.
This condition refers to the explicitly defined cover (BR-OD.2),
rather than to an unidentified projector of some auxiliary surface.

There is also a precise birational corollary. Let $Y/\mathbf F_q$ be a
smooth projective variety for which a smooth projective $Q$ satisfying
this Tate criterion admits a proper birational morphism $\beta:Q\to Y$.
Then the same conclusions hold for $Y$ and its finite products.
Indeed, the graph and transpose graph give actual Chow correspondences
$\beta^*,\beta_*$, and $\beta_*\beta^*=1_{h(Y)}$: the graph
composition is supported on the diagonal, and its generic multiplicity
is $\deg\beta=1$. Equivalently, on every product with a smooth
projective test scheme this is the projection formula and the
pushforward of the fundamental cycle. Thus $h(Y)$ is a Chow summand of
$h(Q)$, also after base change, and CC-XC20 applies.

In particular the conclusion applies to a smooth cubic fourfold
containing a plane whenever its blowup along that plane has the
displayed flat quadric-surface fibration over $\mathbf P^2$, smooth
simple discriminant, and geometrically supersingular K3 discriminant
cover. The blowup is smooth and its morphism to the fourfold is
birational. No assertion that every cubic with a plane has this
discriminant geometry is made.

\subsection{Exact references and scope}
\begin{itemize}
\item Auel, A., Bernardara, M., \& Bolognesi, M. (2014).
Fibrations in complete intersections of quadrics, Clifford algebras,
derived categories, and rationality problems.
\emph{Journal de Math\'ematiques Pures et Appliqu\'ees, 102}(1),
249--291. \url{https://doi.org/10.1016/j.matpur.2013.11.009}.
Theorem 2.2.1 and general-field proof, p.~270; Proposition 1.2.5,
p.~257, for smooth discriminant and simple degeneration.
\item Bernardara, M., Marcolli, M., \& Tabuada, G. (2018).
Some remarks concerning Voevodsky's nilpotence conjecture.
\emph{Journal f\"ur die reine und angewandte Mathematik, 738},
299--312. \url{https://doi.org/10.1515/crelle-2015-0068}.
Orbit category definition, p.~305; Proposition 4.1, pp.~306--307;
equation (6.1) and proof of Theorem 1.2(ii), p.~308.
\item Tabuada, G. (2020). Schur-finiteness (and Bass-finiteness)
conjecture for quadric fibrations and families of sextic du Val del
Pezzo surfaces. \emph{Documenta Mathematica, 25}, 2339--2354.
\url{https://doi.org/10.4171/DM/800}. Theorem 1, p.~2340;
Proposition 15(i), complete proof and discriminant-cover definition,
p.~2346.
\item Tabuada, G., \& Van den Bergh, M. (2013).
\emph{Noncommutative motives of Azumaya algebras}
(arXiv:1307.7946v2). \url{https://arxiv.org/pdf/1307.7946v2}.
Theorem 2.1, p.~3; proof \S8, pp.~14--18. The numbered statement is
used as a cited input; no new proof of Azumaya invariance is claimed.
\end{itemize}

The known decompositions and fully faithful functor are cited inputs.
The finite orbit-retraction argument and its consequences are proved
above. The corank, smoothness, line-bundle quadratic presentation and
odd-characteristic conditions have not been dropped. Arbitrary even
quadric fibrations remain under the broader, less explicit CC-XC18
reduction when these additional hypotheses are unavailable.

% END root_quadric_discriminant.tex

% BEGIN root_springer_archive.tex
\section{Cote 123: the tangent-rank prediction, Springer dimensions, and positivity}
\label{sec:cote123}

\subsection{The original letter and a corrected locator}
The original is Grothendieck's letter to Kostant dated 22 October 1969,
cote 123, physical PDF pages 22--23, manuscript and archive footers
21--22, in \url{https://grothendieck.umontpellier.fr/123.pdf}.
The official PDF has a cover. Thus the intake's locator ``123/023''
for the dimension question is a physical-PDF locator, not footer 23:
footer 23, physical page 24, is the nearly blank verso.
The entire two-page letter was inspected visually in this run.

On its second page Grothendieck asks how the tangent rank of the
adjoint quotient \(\chi:\mathfrak g\to\mathfrak h/W\), and of
the analogous group quotient, varies. For an orbit of dimension
\(\dim G-r-2\), he predicts rank \(r-1\) when all roots have
the same length, and strictly smaller rank otherwise. He also asks
for the dimension of the variety of Borel subgroups containing an
element, or whose Lie algebra contains a given Lie-algebra element.
Writing \(2N=\dim G-r\), an orbit of dimension \(2N-2i\)
is predicted to have such a variety of dimension \(i\).
The initial setting of the letter is a complex simple adjoint group.
These are mathematical questions in the letter itself, distinct
from the later intake's finite-field positivity interpretation.

\paragraph{Known answer to the dimension question.}
In the complex setting the asserted dimension is indeed
\[
 \dim\mathcal B_x=\frac{\dim Z_G(x)-r}{2}.
 \tag{123.1}
\]
Steinberg (1976), \textbf{Theorem 4.6, p.~217}, proves it for
unipotent group elements; \textbf{Theorem 6.5, p.~223}, proves
it for arbitrary group elements. His sufficient-characteristic
condition is automatic over \(\mathbf C\); it must not be
silently discarded when applying that paper in positive
characteristic. The Lie-algebra reduction is Springer's (1984)
\textbf{Lemma 1, p.~273}: split off the semisimple part and
pass to its connected reductive centralizer. In characteristic
zero a nilpotent element and its exponential have the same
centralizer and the same containing Borel subgroups. Applying
Steinberg's unipotent formula in that centralizer gives (123.1)
for the Lie-algebra question as well. This is an application
of known results, not a new proof of their geometry.

\subsection{BR-12301. The tangent rank in type A}
Over \(\mathbf C\), for \(X\in\mathfrak{gl}_n\), the tangent
rank of its adjoint quotient is the degree of its minimal
polynomial. For \(X\in\mathfrak{sl}_n\), it is that degree
minus one. In particular, a nilpotent matrix with largest
Jordan block of size \(b\) has ranks \(b\) and \(b-1\),
respectively. The same \(\mathfrak{sl}_n\) formula applies
to the Lie algebra of the adjoint group \(\mathrm{PGL}_n\).

\begin{proof}
The coefficients of the characteristic polynomial generate
the conjugation-invariant ring on \(\mathfrak{gl}_n\).
For clarity, one can check this assertion without choosing
generators from an abstract restriction theorem. Matrices
with distinct eigenvalues form a dense open subset; each is
conjugate to a diagonal matrix. The restriction of an
invariant polynomial is symmetric in the diagonal entries.
The elementary symmetric polynomials generate these symmetric
polynomials and lift to the characteristic coefficients.
Density proves that restriction is injective, which proves
the assertion. On trace-zero matrices the same argument uses
the hyperplane \(\sum_i\lambda_i=0\), and drops the first
characteristic coefficient. A polynomial on that hyperplane
invariant under the finite permutation group can be extended
invariantly to the whole diagonal space by averaging any
extension; hence its invariant ring is precisely the
symmetric-polynomial ring modulo the first coefficient.

In characteristic zero Newton identities replace the
characteristic coefficients, by a polynomial change of
coordinates with everywhere invertible triangular Jacobian,
with \(p_j(X)=\operatorname{Tr}(X^j)\), \(1\leq j\leq n\).
The differential is
\[
 dp_j|_X(Y)=j\operatorname{Tr}(X^{j-1}Y).
\]
The trace pairing on all matrices is nondegenerate, so their
span has dimension
\[
 \dim\langle 1,X,\ldots,X^{n-1}\rangle=\deg m_X,
\]
where \(m_X\) is the minimal polynomial. The last equality
is its defining minimality, together with division of
polynomials by \(m_X\).

The annihilator of \(\mathfrak{sl}_n\) for the trace pairing
is the scalar line: off-diagonal matrix units first force
a matrix in the annihilator to be diagonal, and differences
of diagonal units force all its entries to agree. The
quotient on \(\mathfrak{sl}_n\) uses \(p_2,\ldots,p_n\).
Their differentials therefore have rank equal to the
dimension of the image of \(\langle X,\ldots,X^{n-1}\rangle\)
modulo the scalar line. Together with that line these
powers span \(\mathbf C[X]\), so this rank is
\(\deg m_X-1\). For a nilpotent matrix \(m_X=t^b\).
Passing from \(\mathrm{SL}_n\) to its adjoint quotient
does not change the Lie algebra or its conjugation orbits.
This proves all assertions.
\end{proof}

\subsection{BR-12302. The strict nonsimply-laced rank prediction is false}
The strict inequality in the letter is refuted by the
complex simple adjoint group \(G=\mathrm{PSp}_4\), of type
\(C_2\). Put
\[
 J=\begin{pmatrix}0&I_2\\-I_2&0\end{pmatrix},\qquad
 N=\begin{pmatrix}0&I_2\\0&0\end{pmatrix}
       \in\mathfrak{sp}_4.
\]
Its orbit has dimension \(6=\dim G-r-2\), but
\(\operatorname{rank}(d\chi_N)=1=r-1\), despite the two
root lengths. The unipotent element represented by \(U=I+N\)
gives the same counterexample for the group quotient.
This refutes the stated strict inequality;
it does not assert a complete rank classification in every type.

\begin{proof}
Every element of \(\mathfrak{sp}_4\) has the form
\[
 Y=\begin{pmatrix}A&B\\C&-A^t\end{pmatrix},
 \qquad B=B^t,\quad C=C^t.
\]
Thus its dimension is \(4+3+3=10\) and its rank is two.
Direct multiplication gives
\[
 [Y,N]=\begin{pmatrix}-C&A+A^t\\0&C\end{pmatrix}.
\]
The Lie centralizer has \(C=0\), arbitrary symmetric
\(B\), and skew-symmetric \(A\), so dimension \(3+1=4\).
In characteristic zero an algebraic group stabilizer is
smooth; equivalently the conjugation orbit map has its
usual Lie-algebra tangent quotient. Therefore the orbit
has dimension \(10-4=6\). Passing through the finite
central isogeny to \(\mathrm{PSp}_4\) does not change it.

For completeness its adjoint quotient is explicitly
\(\chi(Y)=(c_2(Y),c_4(Y))\), where
\[
 \det(tI-Y)=t^4+c_2(Y)t^2+c_4(Y).
\]
Indeed, a regular semisimple symplectic matrix in the Lie
algebra has eigenvalues \(\lambda_1,\lambda_2,-\lambda_1,
-\lambda_2\). The symplectic form pairs opposite eigenspaces,
so an eigenbasis can be normalized to a symplectic basis.
Such matrices form a dense open subset. The normalizer
acts on the two \(\lambda_i\) by independent sign changes
and permutation. Its polynomial invariants are the
symmetric polynomials in \(\lambda_1^2,\lambda_2^2\),
generated by their sum and product. These lift to
\(-c_2\) and \(c_4\). Restriction and density prove
that these two coefficients generate the full invariant
ring, just as in the preceding proof.

Newton's identity gives \(c_2(Y)=-\tfrac12\operatorname{Tr}(Y^2)\).
Consequently \(dc_2|_N(Y)=-\operatorname{Tr}(NY)=-\operatorname{Tr}C\),
which is nonzero, for example when \(C=\operatorname{diag}(1,0)\).
Since \(N\) has rank two, all its three-by-three minors
vanish; the differential of \(c_4=\det\), whose entries
are those cofactors, is zero at \(N\). The quotient
differential thus has rank exactly one, as claimed.

Here is the group version, so that the letter's use of
unipotent orbits causes no ambiguity. Since \(N^2=0\) and
\(N^tJ=-JN\), direct expansion gives \(U^tJU=J\).
The centralizer of \(U=I+N\) is the centralizer of \(N\),
so its orbit also has dimension six. On a maximal torus
of \(\mathrm{Sp}_4\) the eigenvalues are
\(z_1,z_2,z_1^{-1},z_2^{-1}\). Put \(w_j=z_j+z_j^{-1}\).
Inversion-invariant Laurent polynomials in one variable
are polynomials in \(z+z^{-1}\), by the recurrence for
\(z^m+z^{-m}\). Independent inversions and permutation
therefore give invariant generators
\(t=w_1+w_2=\operatorname{Tr}g\) and
\(e=w_1w_2=c_2(g)-2\). Restriction to the dense regular
semisimple locus shows these generate the group invariant
ring, exactly as above. The central element \(-I\) sends
\((t,e)\) to \((-t,e)\); hence the adjoint group quotient
has coordinates \((t^2,e)\). At \(U\), \(t=4\), so
\((t,e)\mapsto(t^2,e)\) has invertible differential.

A tangent vector at \(U\) is \(UY\) with
\(Y\in\mathfrak{sp}_4\). Using
\(c_2(g)=((\operatorname{Tr}g)^2-\operatorname{Tr}(g^2))/2\),
\(U^2=I+2N\), and \(\operatorname{Tr}Y=0\), we obtain
\[
 dt_U(UY)=\operatorname{Tr}(NY),\qquad
 de_U(UY)=4\operatorname{Tr}(NY)-\operatorname{Tr}(U^2Y)
          =2\operatorname{Tr}(NY).
\]
The first functional is nonzero by the same choice of
\(C\), so the group quotient also has rank one, including
after descent to \(\mathrm{PSp}_4\).
\end{proof}

\subsection{BR-12303. The exact positive point-count consequence}
Let \(X_x\) be the Springer fibre of a nilpotent Lie-algebra
element \(x\) over \(\mathbf F_q\), with \(G\) connected reductive and residue
characteristic good for \(G\). Springer's (1984)
\textbf{Theorem 1, pp.~272--273}, says that \(H^j(X_x,\mathbf Q_\ell)\)
is pure of weight \(j\). His \textbf{Corollary 1 to Theorem 2,
p.~280}, gives odd vanishing, and \textbf{Lemma 4 on that page}
says that the eigenvalues on \(H^{2i}\) are \(q^i\) times
roots of unity. Here \(\ell\ne\operatorname{char}\mathbf F_q\).
These are the relevant finite-field inputs; purity alone
would not give the root-of-unity assertion.

There exists \(m\geq1\) such that, for all \(r\geq1\),
\[
 \#X_x(\mathbf F_{q^{mr}})
       =\sum_i b_{2i}q^{mri},\qquad
 Z(X_x/\mathbf F_{q^m},t)
       =\prod_i(1-q^{mi}t)^{-b_{2i}},
 \tag{123.2}
\]
where \(b_{2i}=\dim H^{2i}(X_x,\mathbf Q_\ell)\).
If all the root-of-unity factors are already one, one
can take \(m=1\). No Weyl-group trace is inserted in (123.2).

\begin{proof}
There are finitely many eigenvalues. Choose \(m\) divisible
by the orders of all their root-of-unity factors. The
eigenvalues of \(F^{mr}\) on \(H^{2i}\) are then all
\(q^{mri}\); this statement about traces needs no
semisimplicity assumption. The proper trace formula,
together with odd vanishing, proves the count. Substituting
it in \(Z(t)=\exp(\sum_{r\geq1}\#X_x(\mathbf F_{q^{mr}})t^r/r)\)
and using \(-\log(1-z)=\sum_{r\geq1}z^r/r\) proves the
displayed product as a formal power series. Every exponent
\(b_{2i}\) is nonnegative. This proves the stated consequence.
\end{proof}

\subsection{BR-12304. A Green function with a negative coefficient}
The intake SA20, \S4, D.018, goes from positive ordinary
counts to nonnegative coefficients of Green polynomials.
That implication is false. Use the Springer action and the
normalization of Springer (1984), \S5, p.~279:
\[
 Q_w(x)=\sum_i(-1)^i\operatorname{Tr}
       (Fw\mid H^i(X_x,\mathbf Q_\ell)).
 \tag{123.3}
\]
For \(G=\mathrm{SL}_2\), \(x=0\), odd \(q\), and the
nonidentity element \(s\in W=S_2\),
\[
 X_0=\mathbf P^1,\qquad Q_1(0)=1+q,
       \qquad Q_s(0)=1-q.
 \tag{123.4}
\]
The last polynomial has a negative coefficient. The
trace calculation holds for every odd \(q\). To invoke
the Green-function identification exactly as stated on
Springer's p.~279, take odd \(q\) sufficiently large;
the same polynomial is then a Green polynomial in that
source's terminology.

\begin{proof}
The regular semisimple orbit with eigenvalues \(1,-1\)
is the variety
\(U=(\mathbf P^1\times\mathbf P^1)\setminus\Delta\)
of ordered distinct eigenlines. The nonidentity Weyl element
interchanges the two lines. For \(x=0\), the simultaneous
resolution identifies its cohomology with that of \(X_0\);
this is precisely \textbf{Proposition 3 and the definition
of the Weyl action in \S5, p.~279}, of Springer (1984).
Thus it suffices to compute the swap action on \(H^*(U)\).

Projection \(U\to\mathbf P^1\) is a Zariski locally trivial
affine-line bundle: over a fixed line, complements to that
line form an affine line, and choosing a local complement
trivializes this torsor. Homotopy invariance gives cohomology
only in degrees zero and two, both of dimension one.
Let \(h_1,h_2\) be the two divisor classes on the product.
The divisor \(\Delta\) has class \(h_1+h_2\), as is seen
by intersecting it with each ruling. The localization
sequence for this smooth divisor yields
\[
 H^2(U,\mathbf Q_\ell)
   =\bigl(\mathbf Q_\ell h_1\oplus\mathbf Q_\ell h_2\bigr)
       /\mathbf Q_\ell(h_1+h_2).
\]
Here the adjacent \(H^1(\Delta)(-1)\) is zero and the
Gysin image is exactly the displayed divisor class.
The swap sends \(h_1\) to \(h_2=-h_1\), so it acts
by minus one on this one-dimensional group and by one
on \(H^0\). Both divisor classes are rational over
\(\mathbf F_q\), hence geometric Frobenius acts by
\(q\) on \(H^2\), and by one on \(H^0\).
Equation (123.3) now gives (123.4).
\end{proof}

The dimension formula and the ordinary positive count remain
valid. The failure occurs at the added Weyl trace, whose
character values need not be nonnegative. This is a concrete
boundary for any proposed passage from purity to positivity.
No RH or Dirichlet-GRH input is used in this section.

\subsection{Primary bibliography and inspection scope}
Grothendieck, A. (1969, October 22). \emph{Letter to B. Kostant}.
Montpellier archives, cote 123, physical pages 22--23,
footers 21--22. \url{https://grothendieck.umontpellier.fr/123.pdf}.
The letter also raises questions about opposed tori and rational
singularities; these are separate from the four statements proved
or cited above and are retained in the question register.

Springer, T. A. (1984). A purity result for fixed point varieties
in flag manifolds. \emph{Journal of the Faculty of Science,
University of Tokyo, Section IA, Mathematics, 31}(2), 271--282.
\url{https://doi.org/10.15083/00039529}.
Primary scan: \url{https://repository.dl.itc.u-tokyo.ac.jp/record/39538/files/jfs310202.pdf}.
Theorem 1, pp.~272--273; Lemma 1, p.~273; Proposition 3 and
the action convention, p.~279; Corollary 1 and Lemma 4, p.~280.
The stated results and the indicated action construction were read;
the full purity proof and its preceding geometric sources are
cited inputs, not independently reproved here.

Steinberg, R. (1976). On the desingularization of the unipotent
variety. \emph{Inventiones Mathematicae, 36}, 209--224.
\url{https://doi.org/10.1007/BF01390010}.
Theorem 4.6 and its characteristic convention, p.~217;
Remark 6.3 and Theorem 6.5, p.~223. Primary digitization:
\url{https://gdz.sub.uni-goettingen.de/download/pdf/PPN356556735_0036/LOG_0018.pdf}.
The complete short proofs of the two numbered statements were
read with their earlier classification and dimension results
retained as cited inputs.

% END root_springer_archive.tex

% BEGIN source_cote123_rational.tex
\section{Cote 123: the rational singularities of adjoint-quotient fibres}
\label{sec:SM-cote123-rational}

The letter to Kostant dated 22 October 1969, physical pages 22--23
of cote 123 (handwritten pages 21--22), is explicitly over
$\mathbb C$.  It asks whether the fibres of
$\mathfrak g\to\mathfrak h/W$, and of $G\to T/W$ when $G$ is
simply connected, have rational singularities.  Its stated test is
$R^if_*\mathcal O=0$ for $i>0$ on a desingularization.
This is a known positive result, independent of any adopted
analytic-number-theory input.

\paragraph{Finding F-SM-23 (known affirmative answer, with the reduction written out).}
Let $\mathfrak g$ be a complex semisimple Lie algebra.
Every scheme-theoretic fibre of its adjoint quotient
$\mathfrak g\to\mathfrak h/W$ is reduced and normal and has
rational singularities.  Let $G$ be a connected simply connected
complex semisimple algebraic group, with maximal torus $T$ and Weyl
group $W$.  The same assertions hold for every scheme-theoretic
fibre of $G\to T/W$.
Thus for each such fibre $F$ and every proper resolution
$f:Y\to F$,
\[
 f_*\mathcal O_Y=\mathcal O_F,\qquad
 R^if_*\mathcal O_Y=0\quad(i>0).
\]
These assertions use classical results; the following argument is a
complete deduction from the named inputs, not a new proof of those
inputs.

\begin{proof}
\emph{The Lie-algebra statement.}
Kraft and Procesi (1982, \S0.3, p.~540) state reducedness and
normality for every fibre.  Their immediately following \S0.4 on
the same page explicitly states that these fibres have rational
singularities, citing Hesselink.  The scope of \S0.2 is a reductive
complex group; these introductory assertions are not restricted
to the classical orbit-closure classification treated later in that
paper.  This supplies the first statement directly.

\emph{The group fibre as a scheme.}
Put $r=\operatorname{rank}G$.  Steinberg's fundamental-character
coordinates identify $T/W$ with $\mathbb A^r$ (1965,
Theorem~6.1(b), p.~62, and \S6.16, p.~65).
For a fibre $F$, Theorem~6.11, p.~64, gives an irreducible
underlying variety of codimension $r$, with a dense regular class
whose complement has codimension at least two.
The fibre ideal is generated by $r$ functions in the smooth
variety $G$.  Its height is $r$, so these functions form a regular
sequence locally.  Thus $F$ is Cohen--Macaulay.  Theorem~8.1,
p.~70, says that their differentials are independent on the
regular class.  Therefore $F$ is generically reduced; the absence
of embedded associated primes for a Cohen--Macaulay scheme implies
that $F$ is reduced.  Its complement to the smooth regular locus
has codimension at least two.  Serre's $R_1$ and $S_2$ criterion
then gives normality.  Normality is a deduction here; it is not
misattributed to the literal statement of Theorem~6.11.

\emph{Reduction to a unipotent variety.}
Let $s$ represent the unique closed semisimple class in $F$ and
put $L=Z_G(s)$.  This is a connected reductive group: reductivity
and connectedness are the semisimple-centralizer facts recorded
in Steinberg (1965, \S2.7 and Remark~2.10, p.~53).  The latter
remark states the connectedness result with proof deferred; we
use the established result rather than claim to reproduce that
proof.

Since $F$ is the closure of its regular class, the algebraic
semisimple-part map of Kraft--Procesi (\S0.2, pp.~539--540) gives
a $G$-equivariant morphism
\[
 j:F\longrightarrow G\cdot s=G/L.
\]
Its fibre over $s$ is $s\mathcal U(L)$, where $\mathcal U(L)$ is
the unipotent variety.  Indeed, Jordan decomposition gives that
description for every element with semisimple part $s$; conversely
every $su$, $u\in\mathcal U(L)$, has the same invariant-character
values as $s$ and belongs to $F$.  The standard homogeneous-space
descent, applied to $G\to G/L$, therefore gives
\[
 F \simeq G\times^L \mathcal U(L),\qquad
 [g,u]\longmapsto gsu g^{-1}.
 \tag{SM23.1}
\]
Here the right-hand fibre has its reduced structure.  This also
agrees with the scheme fibre: after the smooth faithfully flat
base change $G\to G/L$, the left side is
$G\times j^{-1}(s)$ and is reduced, since $F$ is reduced.
Faithfully flat descent of reducedness over $\mathbb C$ makes
$j^{-1}(s)$ reduced.  The preceding identification of reduced
closed subvarieties is consequently scheme-theoretic.

Hesselink (1976, Theorem~A and \S2, p.~249) supplies the
equivariant Springer resolution
\[
 \tau:Y_L\longrightarrow\mathcal U(L),\qquad
 \tau_*\mathcal O_{Y_L}=\mathcal O_{\mathcal U(L)},\qquad
 R^i\tau_*\mathcal O_{Y_L}=0\quad(i>0).
 \tag{SM23.2}
\]
The theorem is stated for a split reductive group in
characteristic zero, and hence applies to $L$.
The construction is the Springer resolution of the nilpotent
cone; the equivariant logarithm identifies that cone with
$\mathcal U(L)$ in characteristic zero.  Its equivariance is
essential for the next step.

Form the associated morphism
\[
 f:G\times^L Y_L\longrightarrow G\times^L\mathcal U(L)\simeq F.
\]
The smooth $L$-torsor $G\to G/L$ has sections \'etale locally.
For such a cover $V\to G/L$, the base change of $f$ is
$\operatorname{id}_V\times\tau$.
It is proper and birational, with smooth source, so these
properties descend and $f$ is a resolution.
Flat base change for $\tau$, followed by (SM23.2), gives
\[
 R^i(\operatorname{id}_V\times\tau)_*
       \mathcal O_{V\times Y_L}
 =
 \begin{cases}
 \mathcal O_{V\times\mathcal U(L)},&i=0,\\
 0,&i>0.
 \end{cases}
\]
The corresponding assertions for $Rf_*\mathcal O$ descend along
the resulting faithfully flat \'etale cover of $F$.
This proves rational singularities of $F$.

Finally, the independence of the chosen resolution in
characteristic zero is part of the standard rational-singularity
criterion: the definition with \emph{some} resolution is equivalent
to the definition with \emph{every} resolution.  A directly checked
primary research reference stating this equivalence is
Glazer (2019, \S2A, p.~488), which cites
Kempf--Knudsen--Mumford--Saint-Donat, I.3, pp.~50--51.
It yields the displayed conclusion for every resolution.
\end{proof}

\paragraph{Relation to the particular resolution in the letter.}
Steinberg (1976, Corollary~6.4, p.~223), explicitly attributed
there to Grothendieck, states that for the closure $F$ of a regular
conjugacy class and $t\in T\cap F$ the morphism
\[
 \{(x,gB):g^{-1}xg\in tU\}\longrightarrow F
\]
is a desingularization.  This applies under the present complex
simply connected hypothesis.  F-SM-23 and independence of the
resolution give the requested higher-direct-image vanishing for
this resolution as well.  Corollary~6.4 alone supplies a
resolution, not the higher-direct-image vanishing; the latter
comes from rational singularities as explained above.

\paragraph{Scope and status.}
This closes the rational-singularity question in the original
complex setting.  It is not presented as a new theorem of this
project, as a characteristic-positive analogue, or as a statement
about arbitrary individual orbit closures.  No general
non-simply-connected group assertion is needed here.

\subsection*{References for F-SM-23}
\begin{itemize}
\item Glazer, I. (2019). On rational singularities and counting points
of schemes over finite rings. \emph{Algebra \& Number Theory, 13}(2),
485--500. \url{https://doi.org/10.2140/ant.2019.13.485}.
Exact input: \S2A, p.~488.
\item Hesselink, W. H. (1976). Cohomology and the resolution of the
nilpotent variety. \emph{Mathematische Annalen, 223}(3), 249--252.
\url{https://doi.org/10.1007/BF01360956}.
Exact inputs: Theorem~A and \S2, p.~249.
\item Kraft, H., \& Procesi, C. (1982). On the geometry of conjugacy
classes in classical groups. \emph{Commentarii Mathematici
Helvetici, 57}, 539--602.
\url{https://doi.org/10.1007/BF02565876}.
Exact inputs: \S0.2, pp.~539--540; \S\S0.3--0.4, p.~540.
\item Steinberg, R. (1965). Regular elements of semi-simple algebraic
groups. \emph{Publications Math\'ematiques de l'IH\'ES, 25}, 49--80.
\url{https://doi.org/10.1007/BF02684397}.
Exact inputs: \S2.7 and Remark~2.10, p.~53;
Theorem~6.1(b), p.~62; Theorem~6.11, p.~64;
\S6.16, p.~65; Theorem~8.1, p.~70.
\item Steinberg, R. (1976). On the desingularization of the unipotent
variety. \emph{Inventiones Mathematicae, 36}, 209--224.
\url{https://doi.org/10.1007/BF01390010}.
Exact input: Corollary~6.4, p.~223.
\end{itemize}

% END source_cote123_rational.tex

% BEGIN consequences_criteria_apposition.tex
\section{Cote 123: apposition, opposition, and rigid quadruples in type
A}\label{cote-123-apposition-opposition-and-rigid-quadruples-in-type-a}

\textbf{Scope and status.} These results answer the relevant original
questions for the complex simple adjoint groups
\(G=\operatorname{PGL}_n(\mathbf C)\), \(n\geq2\). All subgroup
automorphisms are complex algebraic automorphisms. No RH or
Dirichlet-GRH input is used. The arithmetic assertion at the end
concerns subgroup schemes individually defined over a specified subfield
\(k\subset\mathbf C\). Literature cutoff: 20 November 2026; source
inspection and proofs: 21 November 2026 UTC. No historical novelty is
asserted.

Grothendieck's letter of 22 October 1969, cote 123, physical PDF page
22, printed footer 21, asks for the classification of quadruples
\((T,B,T',B')\) with the maximal tori in Kostant apposition, whether the
two Borels are opposite for some or all choices, whether a distinguished
conjugacy class can be invariant under outer automorphisms, and what
automorphism groups these quadruples have. The original page was
visually inspected. The distinctions between these questions are
preserved below.

\begin{longtable}[]{@{}
  >{\raggedright\arraybackslash}p{(\columnwidth - 2\tabcolsep) * \real{0.5000}}
  >{\raggedright\arraybackslash}p{(\columnwidth - 2\tabcolsep) * \real{0.5000}}@{}}
\toprule\noalign{}
\begin{minipage}[b]{\linewidth}\raggedright
Result
\end{minipage} & \begin{minipage}[b]{\linewidth}\raggedright
Precise conclusion in type \(A_{n-1}\)
\end{minipage} \\
\midrule\noalign{}
\endhead
\bottomrule\noalign{}
\endlastfoot
CC-AP01 & Every ordered apposition pair is conjugate to the diagonal and
Fourier pair. \\
CC-AP02 & Opposition is exactly a specified list of nonvanishing Fourier
minors. \\
CC-AP03 & Every pair admits opposite Borel choices. All choices are
opposite if and only if \(n\) is prime. For composite \(n\), an explicit
common unipotent subgroup disproves opposition. \\
CC-AP04 & The inner-conjugacy classes are the explicit free quotient
\(K_n\backslash(S_n\times S_n)\), where
\(K_n=(\mathbf Z/n)^2\rtimes(\mathbf Z/n)^\times\). Their number is
\((n!)^2/(n^2\varphi(n))\). \\
CC-AP05 & For every odd \(n\geq3\), no such class is invariant under the
nontrivial outer automorphism; all quadruples then have trivial full
automorphism group. For even \(n\geq4\), outer-invariant examples exist,
with automorphism group \(C_2\). An exact finite criterion covers every
class. \\
CC-AP06 & A quadruple whose four subgroup schemes are defined over
\(k\subset\mathbf C\) exists in split \(\operatorname{PGL}_{n,k}\) if
and only if \(k\) contains a primitive \(n\)-th root of unity. In
particular there is no such quadruple over \(\mathbf Q\) for
\(n\geq3\). \\
CC-AP07 & Adding a generating principal element of \(T\cap N_G(T')\)
gives the free quotient \((\mathbf Z/n)^2\backslash(S_n\times S_n)\),
with \((n!)^2/n^2\) classes. \\
\end{longtable}

\subsection{Inspected primary
convention}\label{inspected-primary-convention}

Use Kostant's convention of a Cartan \(\mathfrak t'\) in apposition to
\(\mathfrak t\) with respect to a principal element \(P\). The intrinsic
characterization needed here is

\[
 \mathfrak t=\mathfrak g^{\operatorname{Ad}P},
 \qquad
 \operatorname{Ad}P(\mathfrak t')=\mathfrak t',
 \qquad
 \operatorname{Ad}P|_{\mathfrak t'}\text{ is a Coxeter element}.
 \tag{CC-AP.1}
\]

A lift of a Weyl-group Coxeter element produces this configuration, and
apposition has this Coxeter action. These are Kostant (1959),
\textbf{Theorem 8.6 and the immediately following Corollary 8.6,
p.~1026}. Their statements were read in the OCR of the author's
collected-paper reprint. Kostant's own later account provides an
independently accessible primary check: \textbf{Theorem 0.1 and the
following apposition sentence, arXiv:1003.0046v1 p.~2; Theorems 1.2--1.3
p.~6; Proposition 1.8 and Theorem 1.9 p.~8}. His principal element is
\(\exp(2\pi i\rho^\vee/h)\); regularity alone is not its definition.

The original definition in 1959, §7.3, p.~1018, was located through
later references but was \textbf{not directly inspected} in the original
scan. No claim of such inspection is made. The proof below uses the
inspected equivalent characterization (CC-AP.1), together with the
cyclic-element construction explicitly identified as apposition by
Kostant himself in 2010. The short proof of his Theorem 1.9 was read;
the full 1959 structure theorem is a cited input.

When using (CC-AP.1) at group level, take the connected centralizer
determined by its Lie algebra. In an adjoint group the full centralizer
of a regular semisimple group element can be disconnected.

\subsection{CC-AP01 --- The diagonal and Fourier normal
form}\label{cc-ap01-the-diagonal-and-fourier-normal-form}

Put \(I=\mathbf Z/n\mathbf Z\), fix \(\zeta=e^{2\pi i/n}\), and use the
basis \(e_0,\ldots,e_{n-1}\) of \(V=\mathbf C^n\). Define

\[
 C e_j=e_{j+1},\qquad
 Qe_j=\zeta^j e_j,\qquad
 v_a=\sum_{j\in I}\zeta^{ja}e_j,
 \qquad F=(\zeta^{ja})_{j,a\in I}.
\]

Let \(T_0\) be the diagonal maximal torus of \(\operatorname{PGL}(V)\),
and set \(T_1=FT_0F^{-1}\).

\textbf{Proposition.} The ordered pair \((T_0,T_1)\) is in Kostant
apposition, and every ordered apposition pair in
\(\operatorname{PGL}_n(\mathbf C)\) is conjugate to it.

\textbf{Proof.} The vectors \(v_a\) form a basis because the usual
Vandermonde determinant of the distinct numbers \(\zeta^a\) is nonzero.
Direct calculation gives

\[
 Cv_a=\zeta^{-a}v_a,\qquad Qv_a=v_{a+1},
 \qquad QCQ^{-1}=\zeta C.
\]

Thus \(T_1\) is the connected centralizer of \(C\). The adjoint action
of \(Q\) fixes exactly the diagonal Cartan and acts on the \(n\)
eigenlines of \(T_1\) by an \(n\)-cycle. Such a cycle is a Coxeter
element of \(S_n\). This verifies apposition by (CC-AP.1). Equivalently,
\(C=\sum_{j=0}^{n-2}E_{j+1,j}+E_{0,n-1}\) is precisely the sum of the
simple-root vectors and the negative-highest-root vector for the reverse
standard positive system. All coefficients are nonzero, so Kostant's
cyclic-element construction applies. Projectively
\(Q=\exp(2\pi i\rho^\vee/n)\) for this positive system.

Conversely, take a pair \((T,T')\) and a principal \(P\) as in
(CC-AP.1). Conjugate \(T'\) to \(T_0\). The action of \(P\) on its
eigenlines is a Coxeter element, hence an \(n\)-cycle. After a
permutation of the eigenlines, a matrix representative of \(P\) has the
form

\[
 A e_j=c_j e_{j+1},\qquad c_j\ne0.
\]

Choose \(c\) with \(c^n=\prod_jc_j\) and choose nonzero diagonal entries
\(d_j\) recursively by \(d_{j+1}=c\,d_j/c_j\). The last equation agrees
with the first by the displayed product relation. Conjugation by
\(\operatorname{diag}(d_j)\) changes \(A\) to \(cC\), which represents
\(C\) in \(\operatorname{PGL}(V)\). Therefore the ordered pair becomes
\((T_1,T_0)\), since the fixed Lie algebra of \(\operatorname{Ad}C\) is
the Cartan centralizing \(C\).

Finally \(F^2=nR\), where \(Re_j=e_{-j}\). The permutation \(R\)
normalizes \(T_0\). It follows that \(F^{-1}T_0F=T_1\), while
\(F^{-1}T_1F=T_0\). Conjugating by \(F^{-1}\) therefore exchanges the
two tori and gives the asserted ordered normal form. \(\square\)

\subsection{CC-AP02 --- The exact opposition
test}\label{cc-ap02-the-exact-opposition-test}

Write permutations \(\pi,\sigma\in S_I\) as ordered lists. Every Borel
containing \(T_0\) stabilizes exactly one of the flags

\[
 {\cal E}^{\pi}_k=\langle e_{\pi(0)},\ldots,e_{\pi(k-1)}\rangle ,
\]

and every Borel containing \(T_1\) is obtained from

\[
 {\cal V}^{\sigma}_r=\langle v_{\sigma(0)},\ldots,v_{\sigma(r-1)}\rangle .
\]

Here \(0\leq k,r\leq n\). This is the full-flag description of Borels in
\(\operatorname{PGL}(V)\): the invariant subspaces for a maximal
diagonal torus are sums of its distinct one-dimensional weight spaces.

\textbf{Proposition.} The Borels
\(B_\pi=\operatorname{Stab}({\cal E}^\pi)\) and
\(B'_\sigma=\operatorname{Stab}({\cal V}^\sigma)\) are opposite if and
only if, for every \(1\leq k\leq n-1\),

\[
 \det\bigl(\zeta^{ja}\bigr)_{
 j\in\{\pi(k),\ldots,\pi(n-1)\},
 \ a\in\{\sigma(0),\ldots,\sigma(n-k-1)\}}\ne0.
 \tag{CC-AP.2}
\]

The order chosen on either displayed set affects only the sign.

\textbf{Proof.} First recall directly why opposite flags are
characterized by

\[
 {\cal E}_k\cap{\cal V}_{n-k}=0
 \quad(1\leq k\leq n-1).
 \tag{CC-AP.3}
\]

Two full flags admit a simultaneous adapted basis. One construction uses
the dimensions \(r_{ij}=\dim({\cal E}_i\cap{\cal V}_j)\): the
nonnegative integers

\[
 r_{ij}-r_{i-1,j}-r_{i,j-1}+r_{i-1,j-1}
\]

are the dimensions of the corresponding successive quotients of the
intersections. Their row and column sums are one, so the nonzero
positions form a permutation matrix. Choosing a vector at each of these
positions gives the simultaneous basis, by induction on the
intersections.

Order that basis according to the first flag, and let \(\eta(i)\) be the
position of its \(i\)-th vector in the second flag. If an increasing
pair \(i<j\), \(\eta(i)<\eta(j)\), occurs, the nontrivial root subgroup
sending the \(j\)-th basis vector to itself plus a multiple of the
\(i\)-th preserves both flags. Hence the common stabilizer contains a
nontrivial unipotent subgroup and is not a torus. If there is no
increasing pair, \(\eta\) is the reverse ordering. The common stabilizer
is then exactly the diagonal torus in this basis. The reverse ordering
is also equivalent to (CC-AP.3). This proves the equivalence with
opposition in the sense used in the letter.

For the flags in the statement, the quotient map
\(V\to V/{\cal E}^{\pi}_k\), restricted to \({\cal V}^{\sigma}_{n-k}\),
is represented by the square submatrix in (CC-AP.2). Its domain and
codomain both have dimension \(n-k\). Its kernel is the intersection in
(CC-AP.3). Thus (CC-AP.2) and (CC-AP.3) are equivalent. \(\square\)

\subsection{CC-AP03 --- ``Some'' always; ``all'' exactly in prime
dimension}\label{cc-ap03-some-always-all-exactly-in-prime-dimension}

\textbf{Theorem.} For \(G=\operatorname{PGL}_n(\mathbf C)\):

\begin{enumerate}
\def\labelenumi{\arabic{enumi}.}
\tightlist
\item
  Every apposition pair admits opposite containing Borels. More
  strongly, after the normal form of CC-AP01, the Fourier flag in the
  ordering \(0,1,\ldots,n-1\) is opposite to every coordinate flag, and
  the coordinate flag in that ordering is opposite to every Fourier
  flag.
\item
  Every quadruple with tori in apposition has opposite Borels if and
  only if \(n\) is prime.
\item
  If \(n\) is composite, the failure can be witnessed by a rank-one
  nilpotent matrix whose associated one-parameter unipotent group lies
  in both Borels.
\end{enumerate}

\textbf{Proof.} For the first assertion, fix the Fourier ordering
\(\sigma(a)=a\). For any coordinate ordering and any \(k\), the minor in
(CC-AP.2) has columns with exponents \(0,\ldots,n-k-1\) and distinct row
parameters \(\zeta^j\). It is a Vandermonde determinant and is nonzero.
Conversely, for the standard coordinate ordering its rows have
consecutive exponents \(k,\ldots,n-1\). Factoring \(\zeta^{ka}\) from
the column indexed by \(a\) again leaves a Vandermonde determinant in
the distinct column parameters \(\zeta^a\). This proves both
strengthened assertions.

If \(n=p\) is prime, every square Fourier minor is nonzero by the
Chebotarev minor theorem, in the exact primary formulation \textbf{Tao
(2005), Lemma 1.3, p.~123, proof pp.~124--125}; equivalently
arXiv:math/0308286v6, Lemma 1.3, pp.~3--4. Applying that theorem to
(CC-AP.2) proves opposition for every pair of containing Borels. CC-AP01
then covers every apposition pair. The cited minor theorem is a known
input; this application does not claim a new proof of it.

For the converse write \(n=ab\), where \(a,b>1\). The submatrix with
rows \(\{0,a\}\) and columns \(\{0,b\}\) is

\[
 \begin{pmatrix}1&1\\1&\zeta^{ab}\end{pmatrix}
 =\begin{pmatrix}1&1\\1&1\end{pmatrix}.
\]

Choose the coordinate ordering so that its last two indices are \(0,a\),
and choose the Fourier ordering so that its first two indices are
\(0,b\). Equation (CC-AP.2) fails at \(k=n-2\).

For a direct group-theoretic witness put

\[
 x=v_0-v_b,\qquad
 \lambda=e_0^*-e_a^*,\qquad
 N=x\otimes\lambda.
 \tag{CC-AP.4}
\]

The vector \(x\ne0\) belongs both to \({\cal E}_{n-2}^{\pi}\) and to
\({\cal V}_2^\sigma\), because its coordinates at \(0,a\) vanish. The
functional \(\lambda\) vanishes on both these subspaces: on the
coordinate subspace this is immediate, and on \(v_0,v_b\) it follows
from \(\zeta^{ab}=1\). Hence \(N\ne0\), \(N^2=0\), and \(I+tN\) is a
nontrivial one-parameter unipotent subgroup.

Every earlier step of the coordinate flag is contained in
\(\ker\lambda\), and every later step contains \(x\). The same holds for
the Fourier flag, with its dividing step at dimension two. Therefore
\(I+tN\) preserves both complete flags. Their common stabilizer contains
this nontrivial unipotent subgroup and cannot be a torus.

In the smallest example \(n=4\), take \(\zeta=i\), coordinate ordering
\((1,3,0,2)\), and Fourier ordering \((0,2,1,3)\). Rescaling (CC-AP.4)
gives the particularly simple witness

\[
 N=(e_1+e_3)(e_0^*-e_2^*).
\]

This is a counterexample inside the original complex simple adjoint
group of type \(A_3\), with apposition established before the flag
calculation. \(\square\)

\subsection{CC-AP04 --- Explicit classification and inner
rigidity}\label{cc-ap04-explicit-classification-and-inner-rigidity}

Let

\[
 K_n=N_G(T_0)\cap N_G(T_1),\qquad
 \Omega_n=S_I\times S_I.
\]

For \(a\in I^\times\) and \(b,d\in I\), define the projective class of
the matrix

\[
 g_{a,b,d}e_j=\zeta^{ad j}e_{aj+b}.
 \tag{CC-AP.5}
\]

\textbf{Theorem.} These are exactly the elements of \(K_n\). On the
coordinate and Fourier eigenline labels their action is

\[
 j\longmapsto aj+b,\qquad
 c\longmapsto a^{-1}c+d.
 \tag{CC-AP.6}
\]

Consequently

\[
 K_n\simeq(I\times I)\rtimes I^\times,\qquad
 a\cdot(b,d)=(ab,a^{-1}d).
\]

The inner-conjugacy classes of ordered quadruples are exactly

\[
 K_n\backslash\Omega_n,\qquad
 (a,b,d)\cdot(\pi,\sigma)
 =(a\pi+b,\ a^{-1}\sigma+d).
 \tag{CC-AP.7}
\]

This action is free. Thus every quadruple has trivial inner stabilizer,
and the number of classes is

\[
 \frac{(n!)^2}{n^2\varphi(n)}.
 \tag{CC-AP.8}
\]

\textbf{Proof.} An element normalizing \(T_0\) is monomial in the
coordinate basis, and an element normalizing \(T_1\) is monomial in the
Fourier basis. Thus a representative \(g\) of an element of \(K_n\)
satisfies, for permutations \(\pi_0,\sigma_0\) of \(I\) and nonzero
scalars \(u_j,w_c\),

\[
 ge_j=u_j e_{\pi_0(j)},\qquad
 gv_c=w_c v_{\sigma_0(c)}.
\]

Comparison of the coordinate indexed by \(\pi_0(j)\) yields

\[
 u_j\zeta^{jc}=w_c\zeta^{\pi_0(j)\sigma_0(c)}.
\]

Divide this equation by its instance at \(c=0\), and then use its
instance at \(j=0\) to remove the remaining scalar. The result is the
congruence

\[
 jc=(\pi_0(j)-\pi_0(0))(\sigma_0(c)-\sigma_0(0))
 \pmod n.
 \tag{CC-AP.9}
\]

At \(j=c=1\), the two factors on the right are units and are inverses.
Put \(a=\pi_0(1)-\pi_0(0)\), \(b=\pi_0(0)\), and \(d=\sigma_0(0)\).
Taking \(c=1\) and then \(j=1\) in (CC-AP.9) gives

\[
 \pi_0(j)=aj+b,\qquad \sigma_0(c)=a^{-1}c+d.
\]

The equation at \(c=0\) then gives \(u_j=u\,\zeta^{ad j}\), with a
scalar \(u\) independent of \(j\). This proves that the projective
matrix is (CC-AP.5).

Conversely (CC-AP.5) gives

\[
 g_{a,b,d}v_c
 =\zeta^{-b(a^{-1}c+d)}v_{a^{-1}c+d},
\]

so all the displayed matrices normalize both tori. Their action on the
two sets of eigenlines is faithful. Composition of the affine
permutations gives

\[
 (a,b,d)(a',b',d')
 =(aa',\,ab'+b,\,a^{-1}d'+d),
\]

which proves the claimed semidirect-product structure and its order
\(n^2\varphi(n)\).

By CC-AP01 every quadruple has a representative with torus pair
\((T_0,T_1)\), and precisely \(K_n\) remains available for conjugating
such representatives. This proves the orbit classification (CC-AP.7). If
a \(K_n\)-element preserves a full flag made from eigenlines, its
permutation of those eigenlines preserves every initial segment and is
therefore the identity. Preserving both flags makes both affine
permutations in (CC-AP.6) the identity. It follows that \(a=1,b=d=0\),
and the projective transformation is the identity. The action is free,
so counting proves (CC-AP.8). This also proves inner rigidity for every
original quadruple, not merely for an opposite one. \(\square\)

For example, there is one class in type \(A_1\), two in type \(A_2\),
and eighteen in type \(A_3\). Formula (CC-AP.7) is an explicit finite
classification in type A; it is not a classification in the other simple
types.

\subsection{CC-AP05 --- Outer action, distinguished classes, and full
automorphisms}\label{cc-ap05-outer-action-distinguished-classes-and-full-automorphisms}

Let \(\theta(g)=(g^{-1})^{\mathsf T}\). For \(n\geq3\) it represents the
nontrivial element of the outer automorphism group of \(G\).

\textbf{Theorem.} On the quotient (CC-AP.7), \(\theta\) acts by

\[
 (\pi,\sigma)\longmapsto
 \bigl(\pi^{\mathrm{rev}},-\sigma^{\mathrm{rev}}\bigr),
 \quad
 \pi^{\mathrm{rev}}(i)=\pi(n-1-i).
 \tag{CC-AP.10}
\]

A class is fixed if and only if some \(a\in I^\times\), \(b,d\in I\)
satisfy, for every \(0\leq i<n\),

\[
 a\pi(i)+b=\pi(n-1-i),\qquad
 a^{-1}\sigma(i)+d=-\sigma(n-1-i).
 \tag{CC-AP.11}
\]

For \(n\geq3\), the full automorphism group of an individual quadruple
is trivial if this condition fails and is \(C_2\) if it holds. If \(n\)
is odd, the condition always fails. If \(n\geq4\) is even, it holds for
explicit examples.

\textbf{Proof.} We first justify the outer-automorphism statement in the
precise algebraic category being used. Any algebraic automorphism of
\(G\), after an inner adjustment, preserves a chosen Borel and its
maximal torus. Its action on the roots then preserves the positive
system and induces an automorphism of the chain \(A_{n-1}\). There are
only identity and reversal, with only identity when \(n=2\). The
composition of \(\theta\) with conjugation by the reverse permutation
matrix preserves this Borel and torus and induces reversal.

If the induced root action is the identity, the automorphism is the
identity on the torus, since its character lattice for the adjoint group
is generated by the roots. On the Lie algebra root spaces it has the
form \(E_{ij}\mapsto c_{ij}E_{ij}\). Preservation of brackets gives
\(c_{ij}c_{ji}=1\), and, for distinct \(i,j,k\),
\(c_{ik}=c_{ij}c_{jk}\). Fixing one index shows \(c_{ij}=s_i/s_j\) for
suitable \(s_i\ne0\); the two-index case follows directly from
\(c_{ij}c_{ji}=1\). Thus conjugation by \(\operatorname{diag}(s_i)\)
gives the same differential. Automorphisms of the connected complex
group are determined by their differential, so the automorphism is
inner. Finally, for \(n\geq3\), reversal of the simple roots cannot
arise from an inner automorphism preserving this Borel and torus: its
conjugating element would belong to \(B\cap N_G(T)=T\), whose action on
the roots is the identity. This proves \(\operatorname{Out}(G)=C_2\)
with representative \(\theta\) when \(n\geq3\), and
\(\operatorname{Out}(G)=1\) when \(n=2\).

Both \(T_0\) and \(T_1\) are preserved by \(\theta\): inverse transpose
preserves diagonal matrices and sends the algebra of circulant matrices
to itself. On a full flag \(\theta\) sends the stabilizer to the
stabilizer of the flag of annihilators, in reverse order, using the
standard bilinear form to identify \(V^*\) with \(V\). The dual
coordinate basis is the same coordinate basis. For the Fourier basis one
has

\[
 v_a^{\mathsf T}v_b=n\,\mathbf 1_{a+b=0},
\]

so its dual basis has labels \(-a\). This proves (CC-AP.10). Combining
it with (CC-AP.7) proves the equivalence with (CC-AP.11).

By CC-AP04 the inner stabilizer is trivial. Consequently the
automorphism group of a quadruple injects into
\(\operatorname{Out}(G)\), and its image is exactly the stabilizer of
its inner-conjugacy class. In particular, if (CC-AP.11) holds, an inner
adjustment of \(\theta\) fixes the quadruple. Its square is inner and
fixes the quadruple, hence is the identity. This proves the asserted
automorphism-group classification.

Now suppose \(n\geq3\) is odd and (CC-AP.11) holds. The affine
permutation

\[
 R(j)=aj+b
\]

reverses an ordering of \(n\) elements. It is an involution with exactly
one fixed point. Therefore \(a^2=1\pmod n\), and the equation
\((a-1)j=-b\) has exactly one solution. The latter fact implies
\(\gcd(a-1,n)=1\). Multiplying the identity \((a-1)(a+1)=0\) by the
inverse of \(a-1\) gives \(a=-1\pmod n\).

The second equation of (CC-AP.11) says that

\[
 S(c)=-a^{-1}c-d
\]

also reverses an ordering of \(n\) elements. With \(a=-1\), this is the
translation \(S(c)=c-d\). Since it is an involution and \(n\) is odd,
\(2d=0\) implies \(d=0\). Then \(S\) is the identity, which cannot
reverse an ordering of at least three distinct elements. This
contradiction proves the assertion for every odd \(n\geq3\), including
composite \(n\).

For an even example write \(n=2m\), with \(m\geq2\), and take

\[
 \pi=(0,1,\ldots,m-1,\ 2m-1,2m-2,\ldots,m),\qquad
 \sigma=(0,1,\ldots,2m-1).
\]

Equation (CC-AP.11) holds with \(a=1,b=m,d=1\). The corresponding class
is outer-invariant, and its automorphism group is \(C_2\), as already
proved. \(\square\)

\textbf{Direct answer to the original distinguished-class question.} In
every type \(A_{n-1}\) with odd \(n\geq3\), there is no inner-conjugacy
class invariant under all outer automorphisms. Thus a distinguished
class satisfying Grothendieck's stated necessary invariance condition
cannot exist there. In type \(A_2\), the two classes of CC-AP04 are
exchanged. Every quadruple in these odd types has trivial full algebraic
automorphism group. In type \(A_1\) the sole class is invariant and its
automorphism group is trivial. For even \(n\geq4\) the theorem supplies
invariant classes and the exact automorphism test, but does not assert a
unique canonical choice among them.

For a completely explicit description in type \(A_2\), the two classes
are distinguished by
\(\operatorname{sgn}(\pi)\operatorname{sgn}(\sigma)\). Translations of
\(\mathbf Z/3\) are even permutations, while multiplication by \(-1\) is
odd on both sets of labels, so this product is \(K_3\)-invariant. The
count of two classes proves it is complete. Representatives are
\((\pi,\sigma)=((0,1,2),(0,1,2))\) and \(((0,1,2),(0,2,1))\). Reversal
changes the row sign, whereas negation followed by reversal leaves the
Fourier sign unchanged; thus \(\theta\) exchanges the two classes
directly. Both representatives have opposite Borels by CC-AP03.

\subsection{CC-AP06 --- The exact split-field
obstruction}\label{cc-ap06-the-exact-split-field-obstruction}

Let \(k\subset\mathbf C\) be a field, and retain \(n\geq2\). Here a
\(k\)-defined quadruple means that each of \(T,B,T',B'\) is a subgroup
scheme of the split group \(\operatorname{PGL}_{n,k}\) defined over
\(k\), and that after extending to \(\mathbf C\) it has the apposition
property above.

\textbf{Theorem.} Such a quadruple exists if and only if
\(\zeta_n\in k\) for a primitive \(n\)-th root of unity \(\zeta_n\).

\textbf{Proof.} A \(k\)-defined Borel in split \(\operatorname{PGL}_n\)
stabilizes a full flag over \(k\). A \(k\)-defined maximal torus inside
it is split. To see the latter assertion concretely, lift the torus and
Borel to \(\operatorname{GL}_n\) and choose a \(k\)-basis adapted to the
flag. The lifted torus has dimension \(n\) and is a torus after base
change. Projection of the upper triangular group onto its diagonal
quotient restricts to an injective map on this torus: its kernel would
be both a torus subgroup and a unipotent subgroup, and hence trivial in
characteristic zero. By equal dimension this map is an isomorphism onto
\((\mathbf G_m)^n\). Thus the torus is split, and its one-dimensional
weight spaces admit a basis over \(k\).

Apply this argument to both \(T\subset B\) and \(T'\subset B'\). Choose
a \(k\)-basis diagonalizing the first torus. There is an
\(A\in\operatorname{GL}_n(k)\) such that the second torus is
\(AT_0A^{-1}\). CC-AP01 gives a complex transformation preserving the
first torus and carrying the second to \(FT_0F^{-1}\). Transformations
preserving \(T_0\) are monomial, so this says exactly that \(A\) is
obtained from \(F\) by row and column permutations and nonzero row and
column scalings.

All entries of \(A\) are therefore nonzero. Its cross-ratios

\[
 \frac{A_{ic}A_{jd}}{A_{id}A_{jc}}\qquad(i\ne j,\ c\ne d)
 \tag{CC-AP.12}
\]

belong to \(k\). Row and column scalings leave these ratios unchanged,
and row and column permutations only relabel them. For the Fourier
matrix, taking rows \(0,1\) and columns \(0,1\) gives the ratio
\(\zeta\). Hence \(\zeta\in k\).

Conversely, if a primitive \(\zeta\) belongs to \(k\), the Fourier
matrix, both tori, and the coordinate and Fourier flags are all defined
over \(k\). They give the required quadruple by CC-AP01. \(\square\)

For \(n=3\) and \(k=\mathbf Q\), this is an independently proved
arithmetic nonsplitness obstruction: a full rational apposition
quadruple does not exist in split \(\operatorname{PGL}_3\), even though
the diagonal torus and the connected centralizer of the rational
cyclic-shift matrix give a rationally defined pair of tori. The second
torus in that pair is not split over \(\mathbf Q\).

The supplied SA20 intake states an arithmetic nonsplitness result
labelled D.016 for ``RigKos in \(A_2/\mathbf Q\)'' but explicitly says
that Part D was read only at statement level. The full definition and
proof of that object were not supplied here. CC-AP06 therefore
establishes the concrete subgroup-scheme assertion above and does not
identify it with an unseen moduli functor, a coarse moduli point, or
every possible twisted form.

\subsection{CC-AP07 --- Classification with a specified generating
principal
element}\label{cc-ap07-classification-with-a-specified-generating-principal-element}

The original letter also asks for the classification when a generating
principal element of the finite group \(T\cap N_G(T')\) is included in
the data.

\textbf{Corollary.} In type \(A_{n-1}\), this group is cyclic of order
\(n\). The inner-conjugacy classes of the ordered data
\((T,B,T',B',P)\), where \(P\) generates that group, are

\[
 K_n\backslash(\Omega_n\times I^\times),
 \quad
 (a,b,d)\cdot(\pi,\sigma,u)
 =(a\pi+b,\ a^{-1}\sigma+d,\ a^{-1}u).
 \tag{CC-AP.13}
\]

Equivalently, after normalizing \(P=Q\), they are

\[
 (I\times I)\backslash\Omega_n,\qquad
 (b,d)\cdot(\pi,\sigma)=(\pi+b,\sigma+d).
 \tag{CC-AP.14}
\]

Both actions are free and their common number of orbits is
\((n!)^2/n^2\).

\textbf{Proof.} In the normal form, an element of \(T_0\cap N_G(T_1)\)
belongs to \(K_n\) and has trivial permutation of the coordinate
eigenlines. By (CC-AP.6), this forces \(a=1,b=0\), so the element is
\(g_{1,0,d}=Q^d\). Conversely every such power normalizes \(T_1\). Hence

\[
 T_0\cap N_G(T_1)=\langle Q\rangle\simeq C_n.
\]

Its generators are precisely \(Q^u\), \(u\in I^\times\). They are
principal: each is conjugate to \(Q\) by permuting its \(n\) distinct
eigenvalues. Direct conjugation, or (CC-AP.5), gives in the projective
group

\[
 g_{a,b,d}Q^u g_{a,b,d}^{-1}=Q^{a^{-1}u}.
\]

This proves (CC-AP.13). Its action is free because its action on the
first two coordinates is already free by CC-AP04. Therefore its orbit
count is

\[
 \frac{(n!)^2\varphi(n)}{n^2\varphi(n)}
 =\frac{(n!)^2}{n^2}.
\]

Given any generator exponent \(u\), choose \(a=u\) to change it to
\(1\). The remaining transformations fixing \(u=1\) have \(a=1\) and are
exactly the two independent translations in (CC-AP.14). This proves the
alternative classification, its freeness, and its count. \(\square\)

These two counts agree with the letter's \(|W|^2/(hz\varphi(h))\) and
\(|W|^2/(hz)\): in type \(A_{n-1}\), \(W=S_n\), \(h=n\), and \(z=n\) is
the order of the center of the simply connected cover
\(\operatorname{SL}_n\).

\subsection{Primary bibliography and inspection
record}\label{primary-bibliography-and-inspection-record}

Grothendieck, A. (1969, October 22). \emph{Letter to B. Kostant}.
Montpellier archives, cote 123, physical PDF page 22, footer 21.
\url{https://grothendieck.umontpellier.fr/123.pdf}. The original page was
visually read, including the distinct classification, opposition,
distinguished-class, and automorphism questions.

Kostant, B. (1959). The principal three-dimensional subgroup and the
Betti numbers of a complex simple Lie group. \emph{American Journal of
Mathematics, 81}(4), 973--1032. \url{https://doi.org/10.2307/2372999}. Theorem
8.6 and the following Corollary 8.6, p.~1026, were read in the OCR of
the author's reprint in \emph{Collected Papers}, Vol. I, pp.~130--189
(2009), \url{https://doi.org/10.1007/b94535_11}. The source scan itself was
not obtained; original §7.3, p.~1018, is a located but uninspected
reference. OCR inspected at
\url{https://dokumen.pub/collected-papers-volume-i-1955-1966-2009nbsped-9780387095820-0387095829.html}.

Kostant, B. (2010). Experimental evidence for the occurrence of \(E_8\)
in nature and the radii of the Gosset circles. \emph{Selecta
Mathematica, 16}, 419--438. \url{https://doi.org/10.1007/s00029-010-0023-1}.
Inspected author preprint: \url{https://arxiv.org/pdf/1003.0046v1} (1 March
2010). Exact locators in this version: Theorem 0.1 and apposition
sentence p.~2; Proposition 1.1 p.~5; Theorems 1.2--1.3 p.~6; Theorem 1.7
and principal-element discussion p.~7; Proposition 1.8 and Theorem 1.9,
including its complete short proof, p.~8.

Tao, T. (2005). An uncertainty principle for cyclic groups of prime
order. \emph{Mathematical Research Letters, 12}(1), 121--127.
\url{https://doi.org/10.4310/MRL.2005.v12.n1.a11}. Lemma 1.3 p.~123; proof
pp.~124--125; preliminary Lemma 1.2 p.~123. The journal PDF was directly
downloaded and the statements and complete indicated proofs read:
\url{https://www.intlpress.com/site/pub/files/_fulltext/journals/mrl/2005/0012/0001/MRL-2005-0012-0001-a011.pdf}.
Also inspected: \url{https://arxiv.org/pdf/math/0308286v6} (22 July 2004),
Lemmas 1.2--1.3 pp.~3--4. The prime Fourier-minor statement is credited
there to Chebotarev; no original Chebotarev theorem number is invented.

\textbf{Proof provenance.} CC-AP01--07 give complete proofs of the
stated type-A deductions, with Kostant's apposition/Coxeter
characterization and the prime Fourier-minor theorem explicitly retained
as known inputs. The finite normalizer and outer-action calculations,
the composite unipotent witness, the field-of-definition obstruction,
and the classification with a specified generator are proved here. This
is neither a proof of RH nor a classification of apposition quadruples
in every simple type.

% END consequences_criteria_apposition.tex

% BEGIN consequences_criteria_apposition_general.tex
\section{Cote 123: ordered apposition quadruples in every simple
type}\label{cote-123-ordered-apposition-quadruples-in-every-simple-type}

\textbf{Scope.} Let \(G\) be a connected complex simple \textbf{adjoint}
algebraic group. Write \(r=\operatorname{rank}G\), \(h\) for its Coxeter
number, \(W\) for its Weyl group, and \(z=|Z(G_{\mathrm{sc}})|\).
Automorphisms are complex algebraic automorphisms. An ordered apposition
pair \((T,T')\) means that some principal element \(c\in T\) has order
\(h\), satisfies \(Z_G(c)^\circ=T\), normalizes \(T'\), and induces a
Coxeter element on \(\mathfrak t'\). This is the inspected Kostant
convention used in CC-AP01; the order matters. The full centralizer of
\(c\) need not be connected.

This module answers the classification, inner-rigidity, and
existence-of-an-outer-invariant-class questions on cote 123, physical
PDF page 22, footer 21, for every simple type. It also answers the
``some opposite Borels'' question for arbitrary two maximal tori. It
does not claim that the outer-invariant class is unique or canonical.
The stronger ``all containing Borels are opposite'' question is settled
throughout the classical types, by CC-AP03 and CC-APG06--08 below; no
classification of that property in the exceptional types is inferred. No
RH or GRH input is used; no historical novelty is asserted. Literature
cutoff: 20 November 2026. Proof and source inspection: 21 November 2026
UTC.

\begin{longtable}[]{@{}
  >{\raggedright\arraybackslash}p{(\columnwidth - 2\tabcolsep) * \real{0.5000}}
  >{\raggedright\arraybackslash}p{(\columnwidth - 2\tabcolsep) * \real{0.5000}}@{}}
\toprule\noalign{}
\begin{minipage}[b]{\linewidth}\raggedright
ID
\end{minipage} & \begin{minipage}[b]{\linewidth}\raggedright
Conclusion
\end{minipage} \\
\midrule\noalign{}
\endhead
\bottomrule\noalign{}
\endlastfoot
CC-APG01 & All ordered apposition pairs form one \(G\)-orbit. Their tori
intersect trivially. Every quadruple has trivial inner stabilizer, and
its full automorphism group is exactly the stabilizer of its
inner-conjugacy class in \(\operatorname{Out}(G)\). \\
CC-APG02 & An explicit finite root-datum group
\(K=A_h\times_{(\mathbf Z/h)^\times}N_W(\langle C\rangle)\) classifies
the quadruples by its free action on \(W\times W\). Its order is
\(hz\varphi(h)\). The two counts in the letter are
\(|W|^2/(hz\varphi(h))\) and, with a generating principal element,
\(|W|^2/(hz)\). \\
CC-APG03 & Given any Borel subgroup \(B\) and any maximal torus \(S\),
some Borel containing \(S\) is opposite to \(B\). \\
CC-APG04 & An inner-conjugacy class of apposition quadruples invariant
under all outer automorphisms exists if and only if the simple type is
not \(A_{2m}\), \(m\geq1\). In every allowed type there is a quadruple
whose full automorphism group is \(\operatorname{Out}(G)\). \\
CC-APG05 & For every quadruple, explicit membership tests in \(A_h\) and
\(N_h\), with equality of their multipliers, determine its full
automorphism subgroup of \(\operatorname{Out}(G)\). In \(A_{2m}\) it is
always trivial. \\
CC-APG06 & In type \(C_r\), every pair of containing Borels is opposite
if and only if \(r\leq2\). For every \(r\geq3\), an explicit pair of
selfdual eigenline flags witnesses failure. In particular the answer is
affirmative in \(B_2=C_2\). \\
CC-APG07 & In type \(B_r\), every pair of containing Borels is opposite
if and only if \(r\leq2\). An explicit scaled Fourier minor gives
nonopposite choices for every \(r\geq3\). \\
CC-APG08 & In every simple type \(D_r\), \(r\geq4\), there are
nonopposite containing Borels. An isotropic eigenline in the cyclic
element's two-dimensional kernel gives a direct witness. \\
\end{longtable}

\subsection{Known structure inputs, with inspected
locators}\label{known-structure-inputs-with-inspected-locators}

Kostant (2010), arXiv:1003.0046v1, gives the following primary inputs.
For a maximal torus \(T\), simple roots \(\alpha_1,\ldots,\alpha_r\),
highest root \(\theta=\sum_i n_i\alpha_i\), and
\(\alpha_i(\rho^\vee)=1\), put

\[
 \zeta=e^{2\pi i/h},\qquad
 c=\exp(2\pi i\rho^\vee/h),\qquad
 x=\sum_{i=1}^r b_i e_{\alpha_i}+b_0e_{-\theta},
 \qquad b_i\ne0.
 \tag{CC-APG.1}
\]

Here \(h=1+\sum_i n_i\). \textbf{Proposition 1.1, p.~5}, identifies the
\(\zeta\)-eigenspace of \(\operatorname{Ad}c\) with the span of the
displayed \(r+1\) root vectors. \textbf{Theorem 1.2, p.~6}, says that an
element in this eigenspace is regular semisimple exactly when all
\(r+1\) coefficients are nonzero. \textbf{Theorem 1.3, p.~6}, says that
\(T'=Z_G(x)\) is a maximal torus, that
\(C=\operatorname{Ad}c|_{\mathfrak t'}\) is a Coxeter element, and that
its \(\zeta\)-eigenspace is exactly \(\mathbf Cx\). The cyclic-element
characterization immediately before \textbf{Theorem 1.4, p.~6}, together
with that theorem, supplies a homogeneous invariant \(J_h\) of degree
\(h\) for which \(J_h(x)\ne0\).

The \textbf{principal-element paragraph at the end of p.~7}, immediately
before Proposition 1.8, states that the adjoint group has one conjugacy
class of regular elements of order \(h\), and that every representative
in \(G\) of a Weyl Coxeter element belongs to that class.
\textbf{Proposition 1.8 and Theorem 1.9, p.~8}, identify all Cartans
Coxeter for \(c\) with the cyclic centralizers above and assert that
\(T\) acts transitively on them. The short proofs of Theorems 1.7 and
1.9 were read in full. These are known inputs, not reproved claims of
this module.

Kostant (1959), \textbf{Theorem 8.6 and the immediately following
Corollary 8.6, p.~1026}, give the original Coxeter/apposition
statements; their author-reprint OCR was inspected. The original
definition in §7.3, p.~1018, remains a located but uninspected scan
reference. The equivalent characterization above is directly supported
by the inspected primary statements.

For Bruhat decomposition we use Steinberg (1965), \textbf{§2.4(a)--(b),
p.~52}. For the orthogonal diagram-orbit argument we use the
\textbf{proof of Theorem 9.4, pp.~72--73}, which explicitly excludes
precisely \(A_n\) with \(n\) even. Our application is written out below;
Steinberg's theorem itself concerns rational cross-sections, not the
quadruple assertion. Standard torus and pinning facts are recorded in
Milne (2017), \textbf{Corollary 17.36, p.~363}, and \textbf{Propositions
23.44--23.45 and Corollaries 23.46--23.49, pp.~497--499}; the
author-hosted text of these statements and their short proofs was
directly inspected.

\subsection{CC-APG01 --- One pair orbit and rigidity of every
quadruple}\label{cc-apg01-one-pair-orbit-and-rigidity-of-every-quadruple}

\textbf{Proposition.} All ordered apposition pairs are conjugate to the
pair \((T,T')\) in (CC-APG.1). For that pair,

\[
 T\cap N_G(T')=\langle c\rangle,\qquad
 T\cap T'=\{1\}.
 \tag{CC-APG.2}
\]

Every ordered quadruple \(\mathcal Q=(T,B,T',B')\), with \(T\subset B\),
\(T'\subset B'\), and \((T,T')\) in apposition, has trivial stabilizer
in \(G=\operatorname{Inn}(G)\). If \([\mathcal Q]\) denotes its
inner-conjugacy class, the quotient map induces an isomorphism

\[
 \operatorname{Aut}(G;\mathcal Q)
 \;\xrightarrow{\ \sim\ }\;
 \operatorname{Stab}_{\operatorname{Out}(G)}([\mathcal Q]).
 \tag{CC-APG.3}
\]

In particular its automorphism group is finite and embeds in
\(\operatorname{Out}(G)\).

\textbf{Proof.} Given an ordered apposition pair, choose its principal
element. Conjugacy of principal elements carries it to \(c\), and
carries the first torus to \(Z_G(c)^\circ=T\). Kostant's Theorem 1.9
then supplies a conjugation by \(T\) carrying its second torus to the
chosen \(T'\). This proves the single-orbit assertion.

Let \(t\in T\cap N_G(T')\). As \(t\) commutes with \(c\), its action on
\(\mathfrak t'\) preserves the one-dimensional \(\zeta\)-eigenspace
\(\mathbf Cx\). Hence \(\operatorname{Ad}(t)x=\lambda x\) for some
\(\lambda\in\mathbf C^\times\). Comparison of the nonzero root-vector
coefficients in (CC-APG.1) gives

\[
 \alpha_i(t)=\lambda\quad(1\leq i\leq r),
 \qquad (-\theta)(t)=\lambda.
\]

The highest-root relation implies \(\lambda^{-(h-1)}=\lambda\), so
\(\lambda^h=1\). Because \(G\) is adjoint, the simple roots are a basis
of \(X^*(T)\). They determine \(t\) uniquely. If \(\lambda=\zeta^j\),
their values agree with those of \(c^j\), so \(t=c^j\). Conversely \(c\)
normalizes \(T'\), proving the first equality in (CC-APG.2).

If also \(t\in T'\), then \(t\) centralizes \(x\), whence \(\lambda=1\)
and \(t=1\). This proves the second equality.

An element of \(G\) preserving both \(T\) and a containing Borel \(B\)
lies in \(N_G(T)\cap B=T\): the Weyl group acts freely on the chambers,
or equivalently on the Borels containing \(T\). If it preserves the
whole quadruple, it also lies in \(T'\), and is therefore the identity.
This proves inner rigidity.

The kernel of the map in (CC-APG.3) is this inner stabilizer. Its image
lies in the stated class stabilizer. Conversely, if an outer class
preserves \([\mathcal Q]\), choose a representative algebraic
automorphism \(a\). There is \(g\in G\) with
\(\operatorname{Int}(g)a(\mathcal Q)=\mathcal Q\); thus the outer class
lies in the image. Injectivity and surjectivity prove (CC-APG.3),
including its group structure. \(\square\)

The two torus positions are genuinely distinct in this convention. The
reverse normalizer intersection will have order \(z\), whereas
\(T\cap N_G(T')\) has order \(h\); no general symmetry of ordered
apposition is being assumed.

\subsection{CC-APG02 --- Exact finite classification and the two
counts}\label{cc-apg02-exact-finite-classification-and-the-two-counts}

Put

\[
 K=N_G(T)\cap N_G(T'),\quad
 W_0=N_G(T)/T,\quad W_1=N_G(T')/T',
 \quad U_h=(\mathbf Z/h\mathbf Z)^\times.
\]

Retain \(C\in W_1\) induced by \(c\). Define the following two finite
groups and homomorphisms:

\[
\begin{aligned}
 A_h&=\{w\in W_0:\ w(c)=c^a\text{ for some }a\in U_h\},
 &\mu_0(w)&=a,\\
 N_h&=N_{W_1}(\langle C\rangle),
 &w C w^{-1}&=C^{\mu_1(w)}.
\end{aligned}
 \tag{CC-APG.4}
\]

The exponents are unique because \(c\) and \(C\) have order \(h\). The
first group has a purely integral root-datum description:

\[
 A_h=\{w\in W_0:\ w\rho^\vee-a\rho^\vee\in hP^\vee
                   \text{ for some }a\in U_h\},
 \tag{CC-APG.5}
\]

where \(P^\vee=X_*(T)\) is the coweight lattice for adjoint \(G\).

\textbf{Theorem.} Both maps in (CC-APG.4) are surjective. The natural
map to the two Weyl groups identifies

\[
 K \simeq A_h\times_{U_h}N_h.
 \tag{CC-APG.6}
\]

If \(F=(T')^c\) is the finite fixed subgroup of \(T'\), there is an
exact sequence

\[
 1\longrightarrow \langle c\rangle\times F
 \longrightarrow K\xrightarrow{\ \mu\ }U_h
 \longrightarrow1,
 \qquad F\simeq Z(G_{\mathrm{sc}}).
 \tag{CC-APG.7}
\]

The last isomorphism is as finite complex diagonalizable groups and need
not be canonical. In particular \(|K|=hz\varphi(h)\). No splitting of
the extension by \(U_h\) is asserted.

After choosing one base Borel for each torus, the inner-conjugacy
classes of ordered apposition quadruples are exactly

\[
 K\backslash(W_0\times W_1),
 \qquad
 k\cdot(w_0,w_1)=(p_0(k)w_0,p_1(k)w_1).
 \tag{CC-APG.8}
\]

This action is free. Consequently their number is

\[
 \frac{|W|^2}{hz\varphi(h)}.
 \tag{CC-APG.9}
\]

If the data also include a generating principal element \(P\) of
\(T\cap N_G(T')\), the classification is the free quotient

\[
 (\langle c\rangle\times F)\backslash(W_0\times W_1),
 \qquad\text{with number }\frac{|W|^2}{hz}.
 \tag{CC-APG.10}
\]

Equations (CC-APG.4)--(CC-APG.8) are a finite root-datum classification
in every type, rather than just a formula for the cardinality.

\textbf{Proof.} We give the eigenvector, normalizer, and lattice steps
explicitly.

First, for every \(a\in U_h\), the Coxeter element \(C\) has a regular
eigenvector with eigenvalue \(\zeta^a\). Indeed, \(C\) is defined over
the rational cocharacter space of \(T'\). Its one-dimensional
\(\zeta\)-eigenspace has a nonzero vector \(v\) defined over
\(\mathbf Q(\zeta)\). This vector is a scalar multiple of \(x\), so
every root is nonzero on \(v\). Applying the automorphism
\(\zeta\mapsto\zeta^a\) to the coordinates gives a
\(\zeta^a\)-eigenvector; roots have rational coefficients, so their
nonzero values stay nonzero. In particular every primitive power \(C^b\)
has a regular \(\zeta\)-eigenvector.

We will also need

\[
 Z_{W_1}(C)=\langle C\rangle.
 \tag{CC-APG.11}
\]

To prove it, let \(w\) commute with \(C\). Then \(wv=\lambda v\),
because the \(\zeta\)-eigenspace is one-dimensional. Since
\(J_h(v)\ne0\) and \(J_h\) is \(W_1\)-invariant, homogeneity implies
\(\lambda^h=1\). The resulting homomorphism \(Z_{W_1}(C)\to\mu_h\) is
injective: a Weyl element fixing a regular complex vector is the
identity. For completeness, write such a vector \(v=a+ib\) in the real
Cartan space. Carry \(a\) into a closed Weyl chamber. Its stabilizer is
the reflection group generated by the walls through \(a\); this follows
by joining adjacent chambers whose closures contain \(a\). For every
root of this stabilizer's root system, its value on \(b\) is nonzero,
since its value on \(a\) is zero and its value on \(v\) is nonzero. Thus
\(b\) lies in an open chamber for that reflection group and has trivial
stabilizer. This proves the asserted injectivity. Its image contains the
scalar actions of \(1,C,\ldots,C^{h-1}\), namely all of \(\mu_h\),
proving (CC-APG.11).

By CC-APG01, \(S=T\cap N_G(T')=\langle c\rangle\) is intrinsic to the
ordered pair. Therefore every \(k\in K\) normalizes \(S\), defining
\(\mu(k)\in U_h\) by

\[
 kck^{-1}=c^{\mu(k)}.
\]

This map is surjective. Fix \(a\in U_h\). The element \(c^a\) is regular
of order \(h\), since its root values are the primitive-power transforms
of those of \(c\). The principal conjugacy statement gives \(n\in G\)
with \(ncn^{-1}=c^a\). Both centralizers have identity component \(T\),
so \(n\) normalizes \(T\). The Cartan
\(\operatorname{Ad}(n)\mathfrak t'\) is stable under \(c\); on it the
action is conjugate to \(C^{a^{-1}}\). The preceding eigenvector
argument gives a regular \(\zeta\)-eigenvector there. That vector
belongs to \(\mathfrak g(\zeta)\), so Kostant's Theorems 1.2--1.3
identify its Cartan as a cyclic centralizer Coxeter for \(c\). By
Theorem 1.9, some \(t\in T\) carries \(nT'n^{-1}\) to \(T'\). Then
\(k=tn\) belongs to \(K\) and \(\mu(k)=a\).

Now \(F=(T')^c\) centralizes \(c\), so it normalizes \(Z_G(c)^\circ=T\)
and belongs to \(K\). If \(k\in\ker\mu\), its image in \(W_1\)
centralizes \(C\), and hence equals \(C^j\) by (CC-APG.11). It follows
that \(c^{-j}k\in T'\) and centralizes \(c\), so \(c^{-j}k\in F\).
Conversely \(\langle c\rangle\) and \(F\) lie in \(\ker\mu\). They
commute and intersect trivially by \(T\cap T'=1\). Thus

\[
 \ker\mu=\langle c\rangle\times F.
 \tag{CC-APG.12}
\]

We identify \(F\) and verify finiteness integrally. Write \(Q\subset P\)
for the root and weight lattices of \(T'\). Its character lattice is
\(Q\), because \(G\) is adjoint. The induced action \(D\) of \(C\) on
characters is again a Coxeter element (using the inverse convention
makes no difference). Choose simple reflections so that
\(D=s_1\cdots s_r\), and let \(\omega_i\) be the fundamental weights.
The reflection formula gives

\[
 (1-D)\omega_i=s_1\cdots s_{i-1}\alpha_i.
 \tag{CC-APG.13}
\]

The vectors on the right form a \(\mathbf Z\)-basis of \(Q\): the
\(i\)-th equals \(\alpha_i\) plus an integral combination of
\(\alpha_1,\ldots,\alpha_{i-1}\). Hence \(1-D\) is invertible over
\(\mathbf Q\) and \((1-D)P=Q\). Its induced map yields

\[
 P/Q \ \simeq\ Q/(1-D)Q.
 \tag{CC-APG.14}
\]

The right side is \(X^*(F)\). The left side is
\(X^*(Z(G_{\mathrm{sc}}))\). This proves \(F\simeq Z(G_{\mathrm{sc}})\),
\(|F|=z\), and (CC-APG.7). It also proves \(T'\cap N_G(T)=F\): for
\(t'\) in that intersection, write \(t'ct'^{-1}=c^a\). Projecting to
\(W_1\), where \(t'\) is trivial, gives \(C=C^a\), hence \(a=1\) and
\(t'\in F\).

Consider the map \(K\to W_0\times W_1\). Its kernel is \(T\cap T'=1\);
its two multiplier values agree, so its image lies in the fiber product
of (CC-APG.6). Its projection onto \(W_0\) is exactly \(A_h\). To see
the nontrivial inclusion, represent \(w_0\in A_h\) by \(n\in N_G(T)\).
Then \(ncn^{-1}=c^a\) for its specified \(a\), and the same
Cartan-transitivity argument used above adjusts \(n\) by an element of
\(T\) to lie in \(K\), without changing \(w_0\).

Its projection onto \(W_1\) is all of \(N_h\). Given \(w_1\in N_h\),
choose \(k\in K\) with \(\mu(k)=\mu_1(w_1)\), using the proved
surjectivity. Then \(w_1p_1(k)^{-1}\) centralizes \(C\), so it equals
\(C^j\). Multiplication by \(c^j\) adjusts \(p_1(k)\) to \(w_1\). This
also proves surjectivity of \(\mu_1\); surjectivity of \(\mu_0\) follows
from that of \(\mu\).

Finally, for any pair \((w_0,w_1)\) with matching multipliers, first
lift \(w_0\) to \(k\in K\). The discrepancy \(w_1p_1(k)^{-1}\) is a
power \(C^j\), and multiplication by \(c^j\) corrects the second
projection while leaving the first unchanged, since \(c\in T\). This
proves surjectivity onto the fiber product. The exponential lattice
\(X_*(T)=P^\vee\) gives (CC-APG.5) directly.

All pair representatives are conjugate by CC-APG01. Once the pair is
fixed, precisely \(K\) remains for conjugation. Borels containing either
torus form a free transitive set for the corresponding Weyl group, so
(CC-APG.8) follows. Inner rigidity makes this action free, giving
(CC-APG.9).

For the enriched data, every generator \(c^u\), \(u\in U_h\), is
principal and makes \(T'\) Coxeter. Indeed \(C^u\) has a regular
\(\zeta\)-eigenvector, and conjugating \(c^u\) to \(c\) lets one apply
Kostant's Theorems 1.2--1.3 to that vector, just as above. The group
\(K\) acts on \(u\) by multiplication by \(\mu(k)\), and this action is
transitive. Normalize the generator to \(c\); its remaining stabilizer
is \(\ker\mu=\langle c\rangle\times F\). The action on the Borel pair
stays free. This proves (CC-APG.10). \(\square\)

The values of \(h,z\) in the formula are the usual root-datum values:

\begin{longtable}[]{@{}lrr@{}}
\toprule\noalign{}
Type & \(h\) & \(z\) \\
\midrule\noalign{}
\endhead
\bottomrule\noalign{}
\endlastfoot
\(A_r\) & \(r+1\) & \(r+1\) \\
\(B_r,C_r\) & \(2r\) & \(2\) \\
\(D_r\) & \(2r-2\) & \(4\) \\
\(E_6\) & \(12\) & \(3\) \\
\(E_7\) & \(18\) & \(2\) \\
\(E_8\) & \(30\) & \(1\) \\
\(F_4\) & \(12\) & \(1\) \\
\(G_2\) & \(6\) & \(1\) \\
\end{longtable}

For example the counts are four in \(B_2=C_2\), twelve in \(G_2\), and
768 in \(D_4\). In type \(A_{n-1}\), (CC-APG.9) reduces to the already
independently calculated \((n!)^2/(n^2\varphi(n))\). These formulas
agree exactly with the two expressions on the inspected original page.

\subsection{CC-APG03 --- Opposite Borels exist for any two maximal
tori}\label{cc-apg03-opposite-borels-exist-for-any-two-maximal-tori}

\textbf{Proposition.} Let \(G\) be a connected complex reductive group.
Given a Borel subgroup \(B_0\) and a maximal torus \(S\), there exists a
Borel \(B_1\supset S\) opposite to \(B_0\). In particular, for any two
maximal tori and any containing Borel for the first, some containing
Borel for the second is opposite to it.

\textbf{Proof.} Choose opposite Borels \(B,B^{-}\) containing \(S\), and
let \(U,U^{-}\) be their unipotent radicals. Write \(B_0=gBg^{-1}\). By
the precise Bruhat decomposition in Steinberg §2.4,

\[
 g=u\dot w b,\qquad
 u\in U\cap\dot wU^{-}\dot w^{-1},\quad b\in B,
\]

where \(\dot w\in N_G(S)\). Put \(D=\dot wB^{-}\dot w^{-1}\). This is a
Borel containing \(S\), and \(u\in D\). The Borels
\(\dot wB\dot w^{-1}\) and \(D\) are opposite. Conjugate both by \(u\);
the first becomes \(B_0\), while the second remains \(D\), since
\(u\in D\). Thus \(B_1=D\) is the required Borel. \(\square\)

This proof does not use apposition. It does not imply opposition for
every containing Borel, as the explicit composite type-A counterexamples
show.

\subsection{CC-APG04 --- Exact existence criterion for an
outer-invariant
class}\label{cc-apg04-exact-existence-criterion-for-an-outer-invariant-class}

Let \(\Gamma=\operatorname{Out}(G)\), identified with the automorphism
group of its Dynkin diagram. The pinning realizes \(\Gamma\) as
algebraic automorphisms preserving a fixed Borel \(B_1\), maximal torus
\(T_1\), and the collection of simple root vectors. This identification
and section are the inspected Milne Propositions 23.44--23.45 and
Corollaries 23.46--23.47.

\textbf{Theorem.} There exists an inner-conjugacy class of ordered
apposition quadruples invariant under every element of \(\Gamma\) if and
only if \(G\) is not of type \(A_{2m}\), \(m\geq1\). In every allowed
type one may choose a representative fixed by the entire pinned group
\(\Gamma\). It then has full automorphism group exactly \(\Gamma\). One
may also include a generating principal element fixed by \(\Gamma\).

\textbf{Proof.} We first construct the representative whenever the
diagram is not \(A_{2m}\). Each \(\Gamma\)-orbit of simple roots then
consists of pairwise orthogonal roots. This can be read directly from
the diagrams: the two arms exchanged in \(D_r\), all three outer
vertices in \(D_4\), the exchanged pairs in \(E_6\), and the reflected
pairs in \(A_{2m-1}\) contain no adjacent vertices. In \(A_{2m}\) the
middle two vertices are adjacent, which is the unique exception. When
\(\Gamma=1\), the orbits are singletons. The same precise fact is used
in the proof of Steinberg's Theorem 9.4, pp.~72--73.

For each simple root \(\alpha_i\), the pinning gives a root homomorphism
\(\phi_i:\operatorname{SL}_2\to G\), compatible with the action of
\(\Gamma\). Orthogonal simple roots generate commuting rank-one
subgroups: their roots have no sum or difference in the root system, and
the Chevalley root commutator relations give commutation of all their
positive and negative root subgroups. Therefore, for each orbit \(O\),
the map

\[
 \phi_O:\operatorname{SL}_2\longrightarrow G,\qquad
 A\longmapsto\prod_{i\in O}\phi_i(A)
\]

is a group homomorphism. Its product is independent of ordering, and its
image is fixed pointwise by \(\Gamma\). Being the image of a connected
group containing the identity, it lies in \(H=(G^\Gamma)^\circ\).

Choose the same standard Weyl representative

\[
 s=\begin{pmatrix}0&1\\-1&0\end{pmatrix}
\]

in every \(\operatorname{SL}_2\), order the finitely many
\(\Gamma\)-orbits in any way, and put

\[
 n=\prod_O\phi_O(s)\in H.
 \tag{CC-APG.15}
\]

Its Weyl image relative to \(T_1\) is a product of every simple
reflection exactly once. It is consequently a Coxeter element. By
Kostant's principal-lift statement, \(n\) is regular semisimple of order
\(h\) in adjoint \(G\). Thus

\[
 T_0=Z_G(n)^\circ
\]

is a maximal torus, and the ordered pair \((T_0,T_1)\) is in apposition
with principal element \(n\). Both tori are \(\Gamma\)-stable, and \(n\)
is fixed by \(\Gamma\).

We must also produce a \(\Gamma\)-stable Borel containing \(T_0\);
stability of the torus alone would not imply this. Since \(H\) is a
connected linear algebraic group and \(n\in H\) is semisimple, \(n\)
lies in a torus \(S\subset H\), by Milne Corollary 17.36. The torus
\(S\) centralizes \(n\), so \(S\subset T_0\). No root \(\alpha\) of
\((G,T_0)\) can restrict trivially to \(S\), since that would imply
\(\alpha(n)=1\), contradicting regularity of \(n\).

Consequently one may choose

\[
 v\in X_*(S)\otimes_{\mathbf Z}\mathbf R
\]

outside the finitely many hyperplanes defined by the restrictions of
these roots. Then \(\alpha(v)\ne0\) for every root. As
\(S\subset G^\Gamma\), the vector \(v\) is fixed by \(\Gamma\). The
positive system \(\{\alpha:\alpha(v)>0\}\) is therefore
\(\Gamma\)-stable, and its Borel \(B_0\supset T_0\) is
\(\Gamma\)-stable. Together with the original pinned \(B_1\), this gives
a quadruple

\[
 (T_0,B_0,T_1,B_1)
\]

fixed by all of \(\Gamma\), and the generator \(n\) is fixed as well.
CC-APG01 says that the full automorphism group injects into \(\Gamma\).
Since it contains the pinned copy just constructed, it equals
\(\Gamma\).

For the converse, suppose the type is \(A_{2m}\), so
\(G=\operatorname{PGL}_n\) with odd \(n=2m+1\geq3\). We recall the
complete obstruction proved in CC-AP04--05. The diagonal/Fourier normal
form identifies the ordered Borel choices with two permutations
\(\pi,\sigma\) of \(\mathbf Z/n\). Every simultaneous inner normalizer
acts on their labels by

\[
 j\longmapsto aj+b,\qquad
 c\longmapsto a^{-1}c+d,\qquad a\in(\mathbf Z/n)^\times.
\]

Inverse transpose, representing the nontrivial outer class, sends their
orderings to \(\pi^{\mathrm{rev}}\) and \(-\sigma^{\mathrm{rev}}\). If
the class were invariant, the affine permutation \(j\mapsto aj+b\) would
reverse an ordering of \(n\) distinct entries. It would be an involution
with exactly one fixed point. Hence \(a^2=1\), and its fixed-point
equation has exactly one solution, forcing \(\gcd(a-1,n)=1\). Therefore
\(a=-1\pmod n\).

The Fourier condition would then say that the translation
\(c\mapsto c-d\) also reverses an ordering of \(n\) distinct entries. It
is an involution, so \(2d=0\). Since \(n\) is odd, \(d=0\), making it
the identity, which cannot reverse such an ordering. This contradiction
excludes every outer-invariant class in type \(A_{2m}\), and a fortiori
excludes an invariant enriched class. \(\square\)

Thus the letter's necessary outer-invariance condition for a
distinguished class fails exactly in these even-rank \(A\)-types. In all
the other types an invariant class exists; uniqueness, a preferred
choice among invariant classes, and a count of all invariant classes are
separate questions not asserted here.

\subsection{CC-APG05 --- A finite test for the full automorphism group
of every
class}\label{cc-apg05-a-finite-test-for-the-full-automorphism-group-of-every-class}

\textbf{Corollary.} In a type other than \(A_{2m}\), choose the
\(\Gamma\)-fixed base quadruple and principal element \(c\) constructed
in CC-APG04. Use its two Borels as origins for the \(W_0\)- and
\(W_1\)-torsors. The group \(\Gamma\) acts on each Weyl group,
preserving \(K\subset W_0\times W_1\). For a representative
\((w_0,w_1)\), put

\[
 d_i(\gamma)=\gamma(w_i)w_i^{-1}\quad(i=0,1).
\]

Its full automorphism group, embedded in \(\Gamma\), is exactly

\[
\begin{split}
 \Gamma_{w_0,w_1}=\{\,\gamma\in\Gamma:\;&
 d_0(\gamma)\in A_h,\quad d_1(\gamma)\in N_h,\\
 &\mu_0(d_0(\gamma))=\mu_1(d_1(\gamma))\,\}.
\end{split}
 \tag{CC-APG.16}
\]

Here \(A_h,N_h,\mu_0,\mu_1\) are defined as in (CC-APG.4), now for this
base pair and \(c\). In type \(A_{2m}\) the full automorphism group is
trivial for every quadruple.

\textbf{Proof.} An automorphism \(\gamma\in\Gamma\) fixes both base tori
and both base Borels. Its action on the two Borel torsors is therefore
\((w_0,w_1)\mapsto(\gamma(w_0),\gamma(w_1))\); there is no affine
translation of either origin. It fixes \(c\), so it also preserves the
intrinsic simultaneous normalizer \(K\), its two Weyl projections, and
its multiplier map.

The orbit \(K\cdot(w_0,w_1)\) is \(\gamma\)-stable if and only if there
is \(k\in K\) with \[
 p_i(k)w_i=\gamma(w_i)\quad(i=0,1).
\] Equivalently \((d_0(\gamma),d_1(\gamma))\in K\). The fiber-product
identification (CC-APG.6) gives precisely the three conditions in
(CC-APG.16). By (CC-APG.3), this orbit stabilizer is the full
automorphism group of the quadruple. For type \(A_{2m}\), the only
nontrivial outer element fixes no class, by CC-APG04, so (CC-APG.3)
gives triviality. \(\square\)

Every part of this test is finite. If torus coordinates are taken
relative to the \(\Gamma\)-stable base Borel instead of the cyclic
coordinates of (CC-APG.1), write \(c=\exp(2\pi i\eta/h)\), with
\(\eta\in X_*(T_0)\) specified modulo \(hX_*(T_0)\). Then the first
condition is computed by \[
 w\eta-a\eta\in hX_*(T_0).
\] This is the same root-lattice test as (CC-APG.5) after a change of
the chosen positive system; it does not require the \(\Gamma\)-stable
Borel to be the cyclic construction's original Borel.

\subsection{CC-APG06 --- The complete ``all opposite'' answer in type
C}\label{cc-apg06-the-complete-all-opposite-answer-in-type-c}

\textbf{Theorem.} Let \(G=\operatorname{PSp}_{2r}(\mathbf C)\), with
\(r\geq1\). Every ordered apposition pair admits only opposite choices
of its two containing Borels if and only if \(r\leq2\). For each
\(r\geq3\) the construction below gives a concrete nonopposite pair of
containing Borels. Here \(C_1=A_1\) and \(C_2=B_2\) are included with
their usual low-rank identifications.

\textbf{Proof.} Put \(n=2r\), \(\zeta=e^{2\pi i/n}\), and equip
\(V=\mathbf C^n\), with basis \(e_0,\ldots,e_{n-1}\), with the
alternating form

\[
 \langle e_i,e_j\rangle=(-1)^i\,\mathbf 1_{i+j=n-1}.
 \tag{CC-APG.17}
\]

It is nondegenerate, and alternating because \(n-1\) is odd. Define the
cyclic shift \(X e_j=e_{j+1}\), with indices modulo \(n\), and set

\[
 v_a=\sum_{j=0}^{n-1}\zeta^{ja}e_j,\qquad
 Qe_j=\zeta^j e_j.
\]

Direct substitution in (CC-APG.17) gives
\(\langle Xu,v\rangle+\langle u,Xv\rangle=0\), including the two
wrap-around terms. Thus \(X\in\mathfrak{sp}(V)\). Moreover,

\[
 Xv_a=\zeta^{-a}v_a,\quad
 QXQ^{-1}=\zeta X,\quad
 \langle v_a,v_b\rangle
   =n\zeta^{-b}\,\mathbf 1_{b=a+r\;(\mathrm{mod}\ n)}.
 \tag{CC-APG.18}
\]

The Fourier vectors form a basis. Let \(T_0\) be the diagonal maximal
torus of \(\operatorname{PSp}(V)\), and \(T_1\) the maximal torus
diagonal in this Fourier basis and preserving the form. It is the group
centralizer of \(X\), interpreted in the adjoint quotient.

We verify Kostant apposition rather than inferring it from the Fourier
appearance alone. Label the weights of \(e_0,\ldots,e_{r-1}\) as
\(\varepsilon_0,\ldots,\varepsilon_{r-1}\); the weight of \(e_{n-1-i}\)
is \(-\varepsilon_i\). Take the simple roots

\[
 \alpha_i=\varepsilon_i-\varepsilon_{i-1}\quad(1\leq i<r),
 \qquad \alpha_r=-2\varepsilon_{r-1}.
\]

Their highest root is \(\theta=-2\varepsilon_0
=2\alpha_1+\cdots+2\alpha_{r-1}+\alpha_r\). The matrix \(X\) is the sum,
with all coefficients one, of the simple-root vectors

\[
 E_{i,i-1}+E_{n-i,n-i-1}\quad(1\leq i<r),
 \qquad E_{r,r-1},
\]

and the negative-highest-root vector \(E_{0,n-1}\). Each displayed
vector satisfies the symplectic infinitesimal condition for (CC-APG.17),
and the sum is exactly \(X\). This is Kostant's cyclic construction,
with \(h=2r=n\), so \((T_0,T_1)\) is an ordered apposition pair by the
inspected Theorems 1.2--1.3.

One may also see the principal element explicitly. The matrix \(Q\)
multiplies the form by \(\zeta^{-1}\). Choose \(\lambda^2=\zeta\); then
\(\lambda Q\in\operatorname{Sp}(V)\), and its adjoint class
\(c=[\lambda Q]\) has all simple-root values equal to \(\zeta\). It is
the principal element of (CC-APG.1). Its action sends \(v_a\) to a
multiple of \(v_{a+1}\), compatibly with (CC-APG.18). By CC-APG01 it
suffices to analyze Borel choices for this one pair.

We spell out the Borel and opposition conditions. Borels of
\(\operatorname{Sp}(V)\), and hence of its adjoint quotient, stabilize
complete selfdual flags \(F_\bullet\), where \[
 F_{n-k}=F_k^\perp.
\] Indeed a basis for an isotropic flag can be completed to a symplectic
basis, and its stabilizer is the standard triangular Borel in that
basis. The group is transitive on these flags, and all Borels arise in
this way. If the Borel contains one of our tori, its flag is a flag of
sums of the distinct torus eigenlines.

Thus the choices for \(T_0\) are orderings \(\pi\) of the coordinate
labels satisfying \[
 \pi(n-1-j)=n-1-\pi(j);
\] the choices for \(T_1\) are orderings \(\sigma\) of the Fourier
labels satisfying \[
 \sigma(n-1-j)=\sigma(j)+r\pmod n.
 \tag{CC-APG.19}
\] Equivalently the first \(r\) labels choose one member of each dual
pair, in any order, and the remaining labels are the reversed dual
partners.

Two such symplectic Borels are opposite exactly when their complete
flags are opposite as ordinary linear flags. In one direction an
opposite symplectic Borel pair can be conjugated to upper and lower
triangular symplectic Borels, whose linear flags are opposite. This also
follows directly from the root decomposition: after fixing their common
maximal torus, no common root subgroup means that the two positive root
systems are negatives. Conversely, suppose the two full flags
\(F_\bullet,G_\bullet\) are opposite. They give lines \[
 L_i=F_i\cap G_{n+1-i},\qquad
 V=\bigoplus_{i=1}^n L_i.
\] Selfduality shows \(\langle L_i,L_j\rangle=0\) unless \(i+j=n+1\):
use the \(F\)-flags when \(i+j\leq n\), and the \(G\)-flags when
\(i+j\geq n+2\). Nondegeneracy then makes the remaining pairings
nonzero. Their common symplectic stabilizer is the diagonal torus in
these lines, with paired diagonal entries inverse. Hence it is a maximal
torus, proving symplectic opposition as defined in the letter.

The full-flag Fourier test from CC-AP02 therefore applies with the
restrictions (CC-APG.19): for every \(1\leq k<n\), the square minor with
rows \(\{\pi(k),\ldots,\pi(n-1)\}\) and columns
\(\{\sigma(0),\ldots,\sigma(n-k-1)\}\) must be nonzero.

For \(r=1\), every relevant minor is an entry of the \(2\times2\)
Fourier matrix and is nonzero. For \(r=2\), the \(1\times1\) minors are
nonzero, and all \(3\times3\) minors are nonzero because the inverse
Fourier matrix has no zero entry and cofactors are its entries times the
nonzero determinant. For a relevant \(2\times2\) minor, the first two
Fourier labels are not dual partners; their difference is therefore odd
modulo four. If its distinct rows are \(u,v\) and columns are \(a,b\),
its determinant is a nonzero scalar times \[
 1-\zeta^{(v-u)(b-a)}.
\] Here \(v-u\not\equiv0\pmod4\), and \(b-a\) is a unit modulo four. The
displayed value is nonzero. All required minors are consequently
nonzero, for every pair of selfdual flags.

Now let \(r\geq3\). Choose the coordinate flag's first two labels to be
\(n-1,r-1\), and the Fourier flag's first two labels to be \(0,2\). Each
pair consists of two distinct nondual labels, so each extends to a
selfdual ordering of the form (CC-APG.19), by choosing one member of
each remaining dual pair and then appending the reversed partners. The
last two coordinate labels are \(r,0\). The relevant minor at \(k=n-2\)
has row set \(\{0,r\}\) and column set \(\{0,2\}\); it is \[
 \begin{pmatrix}1&1\\1&\zeta^{2r}\end{pmatrix}
 =\begin{pmatrix}1&1\\1&1\end{pmatrix}.
\] It is singular. Equivalently, \(v_0-v_2\ne0\) has zero coordinates in
rows \(0,r\), so it belongs both to the coordinate \(F_{n-2}\) and the
Fourier \(G_2\). The full flags fail the opposition condition. By the
proved symplectic equivalence, their Borels are not opposite. This
supplies the stated counterexample in every rank \(r\geq3\). \(\square\)

In particular \(C_2\) has four inner-conjugacy classes by CC-APG02, and
all four have opposite Borels. This differs from the type \(A_3\)
behavior at the same matrix dimension: the selfduality restrictions on
allowed Fourier flags are essential.

\subsection{CC-APG07 --- The complete ``all opposite'' answer in type
B}\label{cc-apg07-the-complete-all-opposite-answer-in-type-b}

\textbf{Theorem.} For the adjoint group
\(\operatorname{SO}_{2r+1}(\mathbf C)\) of type \(B_r\), every pair of
containing Borels for an apposition pair is opposite if and only if
\(r\leq2\). For \(r\geq3\), explicit nonopposite choices are given
below.

\textbf{Proof.} The cases \(B_1=A_1\) and \(B_2=C_2\) follow from the
proved low-rank results; apposition and opposition are invariant under
these group isomorphisms. Let \(r\geq3\), put \(m=2r\), and use the
basis \(e_0,\ldots,e_m\) of \(V=\mathbf C^{m+1}\), with symmetric form

\[
 (e_i,e_j)=(-1)^i\mathbf 1_{i+j=m}.
 \tag{CC-APG.20}
\]

It is symmetric because \(m\) is even. Put

\[
 X=\sum_{i=0}^{m-1}E_{i+1,i}
       +E_{0,m-1}+E_{1,m}.
 \tag{CC-APG.21}
\]

This lies in \(\mathfrak{so}(V)\), as follows either by direct
substitution in (CC-APG.20) or from the root vectors just below. Let
\(T_0\) be the diagonal torus, with weights \(\varepsilon_i\) on
\(e_i\), \(0\leq i<r\), zero on \(e_r\), and \(-\varepsilon_i\) on
\(e_{m-i}\). Choose the simple roots

\[
 \alpha_i=\varepsilon_i-\varepsilon_{i-1}\quad(1\leq i<r),
 \qquad \alpha_r=-\varepsilon_{r-1}.
\]

The simple-root vectors are \[
 E_{i,i-1}+E_{m+1-i,m-i}\quad(1\leq i<r),
 \qquad E_{r,r-1}+E_{r+1,r}.
\] Their sum is the subdiagonal sum in (CC-APG.21). The highest root is
\(\theta=-\varepsilon_0-\varepsilon_1
=\alpha_1+2\alpha_2+\cdots+2\alpha_r\); its negative-root vector is
\(E_{0,m-1}+E_{1,m}\). Thus \(X\) is the cyclic element with all simple
and negative-highest coefficients nonzero. Kostant's theorem gives the
ordered apposition pair \((T_0,T_1)\), where \(T_1=Z_G(X)\).

Here is an eigenbasis making the counterexample explicit. Choose
\(a\in\mathbf C^\times\) with \(a^m=2\), put \(\zeta=e^{2\pi i/m}\), and
set \(\lambda_j=a\zeta^j\), \(j\in\mathbf Z/m\). Define \[
 v_j=e_0+\sum_{i=1}^{m-1}\lambda_j^{m-i}e_i+e_m,
 \qquad w=e_0-e_m.
 \tag{CC-APG.22}
\] The row equations in (CC-APG.21) give \[
 Xv_j=\lambda_jv_j,\qquad Xw=0.
\] These \(m+1\) eigenvectors have distinct eigenvalues and form a
basis. Moreover, \[
 (v_j,v_b)=2m\,\mathbf 1_{b=j+r},\qquad
 (w,v_j)=0,\qquad(w,w)=-2.
 \tag{CC-APG.23}
\] For the first identity, with \(\lambda=\lambda_j\),
\(\mu=\lambda_b\), expand the form as \[
 2+\sum_{i=1}^{m-1}(-1)^i\lambda^{m-i}\mu^i
   =2\sum_{i=0}^{m-1}(-\mu/\lambda)^i.
\] The ratio is an \(m\)-th root of unity, so the sum has exactly the
claimed values. Thus \(T_1\) has the dual eigenline pairs
\(v_j,v_{j+r}\), while \(w\) is its nonisotropic zero-weight line.

A torus-containing Borel of the odd orthogonal group is specified by a
complete isotropic flag: choose one line from each dual pair, order
those \(r\) lines, put the zero-weight line in the middle of the
associated complete flag, and append the reversed dual partners. The
argument of CC-APG06 for symplectic flags also proves the analogous
opposition assertion here. In particular opposite orthogonal Borels have
opposite complete linear flags; the middle self-paired line contributes
diagonal entry \(1\) in the common special-orthogonal torus.

Choose the first two coordinate lines to be \(e_{m-1},e_{r-1}\), and the
first two Fourier lines to be \(v_0,v_2\). In each case these are
distinct nondual lines when \(r\geq3\), so they extend to
torus-containing Borel flags. The coordinate \(F_{m-1}=F_2^\perp\) omits
exactly rows \(1,r+1\). The matrix for the map \[
 \langle v_0,v_2\rangle\longrightarrow V/F_{m-1}
\] in those two rows is \[
 \begin{pmatrix}
 a^{m-1}&a^{m-1}\zeta^{-2}\\
 a^{r-1}&a^{r-1}\zeta^{-2}
 \end{pmatrix}.
\] It is singular. The nonzero vector \(v_2-\zeta^{-2}v_0\) belongs to
\(F_{m-1}\cap G_2\). Since \(\dim V=m+1\), this is exactly a failure of
the full-flag opposition condition. The containing orthogonal Borels are
therefore not opposite. By the single pair orbit, these are
counterexamples for the original apposition question in every type
\(B_r\), \(r\geq3\). \(\square\)

\subsection{CC-APG08 --- Nonopposition in every simple type
D}\label{cc-apg08-nonopposition-in-every-simple-type-d}

\textbf{Theorem.} In the adjoint group of type \(D_r\), \(r\geq4\), an
apposition pair always has some choices of containing Borels which are
not opposite.

\textbf{Proof.} Work first in \(\operatorname{SO}_{2r}(\mathbf C)\);
quotienting its center preserves tori, Borels, their root systems, and
the apposition and opposition properties. Put \(n=2r\). On the basis
\(e_0,\ldots,e_{n-1}\), choose the symmetric form \[
 (e_i,e_{n-1-i})=s_i,\qquad
 s_i=(-1)^i\ (0\leq i<r),\quad
 s_{n-1-i}=s_i,
 \tag{CC-APG.24}
\] with other entries zero. For the diagonal torus \(T_0\), write
\(\varepsilon_i\) for the weight on \(e_i\), \(0\leq i<r\), so the
weight on \(e_{n-1-i}\) is \(-\varepsilon_i\).

Take the simple roots \[
 \alpha_i=\varepsilon_i-\varepsilon_{i-1}\quad(1\leq i<r),
 \qquad \alpha_r=-\varepsilon_{r-2}-\varepsilon_{r-1}.
\] The highest root is \[
 \theta=-\varepsilon_0-\varepsilon_1
 =\alpha_1+2\alpha_2+\cdots+2\alpha_{r-2}
    +\alpha_{r-1}+\alpha_r.
\] The following skew-adjoint matrices are root vectors for these simple
roots and the negative highest root: \[
\begin{aligned}
 u_i&=E_{i,i-1}+E_{n-i,n-i-1} &&(1\leq i<r),\\
 u_r&=E_{r,r-2}+E_{r+1,r-1},&
 u_0&=E_{0,n-2}+E_{1,n-1}.
\end{aligned}
\] Their form compatibility follows directly from (CC-APG.24); the two
entries of each \(u_i\) cancel under the adjoint condition. Put \[
 X=u_0+u_1+\cdots+u_r.
 \tag{CC-APG.25}
\] Kostant's cyclic theorem makes \(X\) regular semisimple and gives the
apposition pair \((T_0,T_1)\), with \(T_1=Z_{\operatorname{SO}(V)}(X)\).
We use its image in the adjoint group at the end.

The row equations of (CC-APG.25) show that \[
 \ker X=\langle u,v\rangle,\qquad
 u=e_0-e_{n-1},\quad v=e_{r-1}-e_r.
 \tag{CC-APG.26}
\] Indeed all coordinates except \(0,n-1,r-1,r\) vanish in the kernel,
and the remaining equations are \(x_0+x_{n-1}=0\) and \(x_{r-1}+x_r=0\).
The restricted form satisfies \[
 (u,u)=-2,\qquad
 (v,v)=-2(-1)^{r-1},\qquad(u,v)=0.
\] Choose \(\lambda^2=(-1)^r\). Then \[
 w_+=u+\lambda v,\qquad w_-=u-\lambda v
\] span the two isotropic lines in this nondegenerate plane, and
\((w_+,w_-)=-4\).

These are \(T_1\)-eigenlines. To justify this directly, \(T_1\)
preserves \(\ker X\) and its form. It is connected, so it preserves each
of the two isotropic lines. Moreover the special-orthogonal torus of
this plane, acting trivially on its orthogonal complement, centralizes
\(X\): the complement is \(X\)-stable by skew-adjointness and
nondegeneracy of the kernel. Thus it lies in \(T_1\), acts nontrivially
on both lines, and makes them a genuine pair of distinct torus weight
lines.

The coordinate line \(\mathbf Ce_1\) is an isotropic \(T_0\)-eigenline
and \[
 (e_1,w_+)=0.
 \tag{CC-APG.27}
\] This follows because its only paired coordinate is \(n-2\), which is
absent from the support \(0,n-1,r-1,r\) of \(w_+\) for \(r\geq4\).

Choose a Borel containing \(T_0\) whose first isotropic line is
\(\mathbf Ce_1\), and one containing \(T_1\) whose first isotropic line
is \(\mathbf Cw_+\). Such choices exist: in either torus eigenbasis, put
the desired line first and extend by one line from each remaining dual
pair to a maximal isotropic flag. Its stabilizer is a Borel. In the even
orthogonal case the two choices of final maximal-isotropic extension can
have the same Borel; this does not change the first line or the
argument.

If these two Borels were opposite, their first isotropic lines would
pair nontrivially. Indeed conjugate an opposite pair to standard upper
and lower triangular orthogonal Borels. Their first lines are the paired
endpoint lines \(e_0,e_{n-1}\); the common invariant form has a nonzero
pairing there, which is preserved by conjugation. The first line of a
Borel's isotropic flag is unique, so the implication applies to our
chosen lines independently of the possible final-extension ambiguity.
Equation (CC-APG.27) contradicts it. The two Borels are not opposite.

Finally, under the finite central quotient
\(\operatorname{SO}_{2r}\to G_{\mathrm{ad}}\), the cyclic Cartans still
give apposition by the adjoint principal element, and opposite Borels
lift to opposite Borels. The same pair therefore refutes the ``all''
assertion for the required adjoint group. \(\square\)

The complete classical answer to the universal-choice question is
consequently:

\begin{longtable}[]{@{}ll@{}}
\toprule\noalign{}
Simple type & Are all containing Borel choices opposite? \\
\midrule\noalign{}
\endhead
\bottomrule\noalign{}
\endlastfoot
\(A_r\) & Yes exactly when \(r+1\) is prime. \\
\(B_r\), \(r\geq2\) & Yes exactly when \(r=2\). \\
\(C_r\), \(r\geq2\) & Yes exactly when \(r=2\). \\
\(D_r\), \(r\geq4\) & No. \\
\end{longtable}

In every row, some opposite choices still exist by CC-APG03.

\subsection{Primary bibliography and proof
provenance}\label{primary-bibliography-and-proof-provenance}

Grothendieck, A. (1969, October 22). \emph{Letter to B. Kostant}.
Montpellier archives, cote 123, physical PDF page 22, footer 21.
\url{https://grothendieck.umontpellier.fr/123.pdf}. The original page and its
two counts were visually inspected.

Kostant, B. (1959). The principal three-dimensional subgroup and the
Betti numbers of a complex simple Lie group. \emph{American Journal of
Mathematics, 81}(4), 973--1032. \url{https://doi.org/10.2307/2372999}. Theorem
8.6 and the immediately following Corollary 8.6, p.~1026; author-reprint
OCR inspected at
\url{https://dokumen.pub/collected-papers-volume-i-1955-1966-2009nbsped-9780387095820-0387095829.html}.
Author reprint in \emph{Collected Papers}, Vol. I, pp.~130--189 (2009),
\url{https://doi.org/10.1007/b94535_11}. Original definition §7.3, p.~1018,
is not represented as a directly inspected scan.

Kostant, B. (2010). Experimental evidence for the occurrence of \(E_8\)
in nature and the radii of the Gosset circles. \emph{Selecta
Mathematica, 16}, 419--438. \url{https://doi.org/10.1007/s00029-010-0023-1}.
Inspected version: \url{https://arxiv.org/pdf/1003.0046v1} (1 March 2010).
Proposition 1.1 p.~5; Theorems 1.2--1.4 p.~6; Theorem 1.7 and
principal-element paragraph p.~7; Proposition 1.8 and Theorem 1.9 p.~8.
The two short conjugacy proofs were read in full; the underlying 1959
theorems remain known inputs.

Milne, J. S. (2017). \emph{Algebraic groups: The theory of group schemes
of finite type over a field} (Cambridge Studies in Advanced Mathematics,
Vol. 170). Cambridge University Press. Author-hosted text:
\url{https://www.jmilne.org/math/Books/iAG2017.pdf}. Corollary 17.36 p.~363;
Propositions 23.44--23.45 and Corollaries 23.46--23.49 pp.~497--499. The
cited statements and short proofs were read in the extracted
author-hosted text; no whole-book verification is claimed.

Steinberg, R. (1965). Regular elements of semi-simple algebraic groups.
\emph{Publications Mathématiques de l'IHÉS, 25}, 49--80.
\url{https://doi.org/10.1007/BF02684397}. Primary scan:
\url{https://www.numdam.org/item/PMIHES_1965__25__49_0/}. §2.4(a)--(b)
p.~52; Theorem 9.4 and its complete proof pp.~72--73. Its
orbit-of-orthogonal-roots argument is used explicitly, with no
attribution of our quadruple statement to that theorem.

\textbf{Proof provenance.} CC-APG01--08 are complete deductions from the
specifically listed structure inputs and the already fully proved type-A
normalizer and outer-action calculations in CC-AP04--05. The
centralizer-of-Coxeter argument, integral computation of the finite
fixed torus, exact fiber product, free orbit counts, Bruhat opposition
argument, invariant-quadruple construction, full finite automorphism
test, and classical matrix and flag calculations are all supplied.
Nothing here independently reproves the whole Kostant theory, the
algebraic-group structure theorem, RH, or Dirichlet GRH.

% END consequences_criteria_apposition_general.tex

% BEGIN root_exceptional_opposition.tex
\section{Exceptional nonopposition from an exterior-square pairing}

\subsection{BR-123-WEDGE: a root-pair obstruction}
We continue the original question in cote 123, physical page 22,
footer 21: whether all choices of containing Borels for tori in
Kostant apposition are opposite.  We give complete counterexamples
in types $F_4$ and $E_8$.  The argument uses the adjoint representation
alone, and does not assume a table of exceptional representation
dimensions or a numerical diagonalization.

The known structural inputs are Kostant (2010), Theorems 1.2--1.3,
p.~6, and the principal-element paragraph on p.~7; and Milne (2017),
Proposition 22.17 and Theorems 22.18--22.19, pp.~467--468,
and Proposition 22.41, p.~477.  The latter give the highest-weight
line and complete reducibility in characteristic zero.  The explicit
Euclidean root sets are in Milne, \S24(l), pp.~539--542.  The finite
root calculations actually needed below are displayed, rather than
replaced by a claim that these sources contain the counterexamples.

Let $\mathfrak g=\operatorname{Lie}G$, and let $K$ denote its Killing
form in this section only.  It induces the invariant bilinear form
\[
 K_2(a\wedge b,u\wedge v)
       =K(a,u)K(b,v)-K(a,v)K(b,u)
       \quad\hbox{on }\bigwedge^2\mathfrak g.
 \tag{EX.1}
\]
For a maximal torus with chosen simple roots, suppose that distinct
roots $\theta,\delta$ satisfy the following explicit conditions:
$\theta+\alpha_i$ is never a root; $\delta+\alpha_i$ is a root
only for $i=j$, when it equals $\theta$; and $-1$ belongs to the
Weyl group.  Choose nonzero root vectors and put
\[
                         z=e_\theta\wedge e_\delta.
\]
Every simple positive root operator kills $z$: it kills the first
factor, and its only possible nonzero action on the second produces
a multiple of the first factor.  The simple root operators generate
the positive nilpotent Lie algebra, so the whole unipotent radical
fixes $z$.  It is also a torus eigenvector.  It therefore spans the
highest line of its cyclic subrepresentation $V=\langle Gz\rangle$.
By complete reducibility and Proposition 22.17 this subrepresentation
is irreducible: in a direct-sum decomposition, every summand receiving
the cyclic vector has that same highest weight, and the cyclic
diagonal copy has one-dimensional highest-weight space.

A representative of $-1$ sends $z$ to a nonzero scalar multiple of
$e_{-\theta}\wedge e_{-\delta}$.  This vector belongs to $V$ and
has nonzero $K_2$-pairing with $z$, because opposite root spaces pair
nondegenerately and the two cross-pairings vanish.  Hence $K_2|_V$
is nonzero; its radical is invariant, so irreducibility makes the
restriction nondegenerate.  It follows also directly that highest
lines for opposite Borels in $V$ pair nontrivially: conjugate the
opposite Borels simultaneously to the chosen positive and negative
Borels, whose highest lines are the two just displayed, and use
invariance and uniqueness of the highest line.

Now let $(T_0,T_1)$ be an apposition pair, with principal element
$c\in T_0\cap N_G(T_1)$ of order $h$, acting on the second root
system as a Coxeter element $C$.  Choose the positive root system
of $T_0$ so that
\[
 \alpha_i(c)=\zeta=e^{2\pi i/h},\qquad
 \alpha(c)=\zeta^{\operatorname{ht}(\alpha)}.
 \tag{EX.2}
\]
Suppose that the unordered pairs $(\alpha,\beta)$ in the first
root system and $(\mu,C^d\mu)$ in the second are Weyl translates
of $(\theta,\delta)$, and that
\[
             h\mid d\bigl(\operatorname{ht}(\alpha)
                            -\operatorname{ht}(\beta)\bigr).
 \tag{EX.3}
\]
The vector $z_0=e_\alpha\wedge e_\beta$ is then a highest-line
vector for some Borel $B_0$ containing $T_0$.  If
$u\in\mathfrak g_\mu(T_1)$ is nonzero, then
$z_1=u\wedge\operatorname{Ad}(c^d)u$ is nonzero and is a
highest-line vector for some Borel $B_1$ containing $T_1$.
Both lie in the same subrepresentation $V$: transport the highest
wedge to the other torus by conjugation and then use the stated
Weyl translates.

Write $a=K(e_\alpha,u)$ and $b=K(e_\beta,u)$.  Invariance gives
\[
 K(e_\alpha,\operatorname{Ad}(c^d)u)
       =\zeta^{-d\operatorname{ht}(\alpha)}a,
\]
and similarly for $\beta$.  Consequently
\[
 K_2(z_0,z_1)=ab\bigl(
       \zeta^{-d\operatorname{ht}(\beta)}
       -\zeta^{-d\operatorname{ht}(\alpha)}\bigr)=0.
 \tag{EX.4}
\]
This conclusion does not require either $a$ or $b$ to be nonzero.
The highest-line test proved above shows that $B_0,B_1$ cannot be
opposite.  This proves the obstruction.

For its applications, one can first choose the explicitly ordered
Coxeter element on $T_1$ and a lift $c$.  Kostant's p.~7 statement
makes it principal and regular; its centralizer identity component
is $T_0$.  Conjugating this principal element gives (EX.2), while
transporting the chosen root coordinates on $T_1$.  Thus prescribing
the Coxeter word on the second root system and the principal
coordinates on the first is legitimate.  CC-APG01 then transports
one counterexample to every ordered apposition pair.

\subsection{BR-123-OPPF4: a counterexample in type \texorpdfstring{$F_4$}{F4}}
Use the Euclidean root system
\[
 \{\pm e_i,\ \pm e_i\pm e_j,\
             \tfrac12(\pm e_1\pm e_2\pm e_3\pm e_4)\}
\]
with the simple roots
\[
 a_1=e_2-e_3,\quad a_2=e_3-e_4,\quad a_3=e_4,\quad
 a_4=\tfrac12(e_1-e_2-e_3-e_4).
 \tag{EX.5}
\]
The root-coroot Cartan matrix, with entry
$A_{ij}=\langle a_i,a_j^\vee\rangle$, is
\[
 A=\begin{pmatrix}
 2&-1&0&0\\-1&2&-2&0\\0&-1&2&-1\\0&0&-1&2
 \end{pmatrix}.
\]
Put $\theta=e_1+e_2=2a_1+3a_2+4a_3+2a_4$ and
$\delta=e_1+e_3=\theta-a_1$.  From the displayed root set,
$\theta+a_i$ is never a root, and $\delta+a_i$ is a root only
for $i=1$, giving $\theta$.  The element $-1$ is a Weyl element:
multiply the four reflections in the mutually orthogonal short
roots $e_i$.  Thus the highest-wedge hypotheses hold.

On the first torus take
\[
                 \alpha=e_2+e_3,
                 \qquad\beta=\theta.
\]
Starting with $(\theta,\delta)$, apply the simple reflections
$s_2,s_3,s_4$ in that order.  They fix $\theta$ and send
\[
 e_1+e_3\longmapsto e_1+e_4
       \longmapsto e_1-e_4\longmapsto e_2+e_3.
\]
This is the required unordered-pair Weyl certificate.  Its heights
are $\operatorname{ht}(\alpha)=5$ and
$\operatorname{ht}(\beta)=11$; equivalently the height functional
on Euclidean coordinates is $(8,3,2,1)$.

On the second root system choose $C=s_1s_2s_3s_4$, with the rightmost
reflection applied first, and take $\mu=\theta$.  Direct reflection
gives
\[
 C(e_1+e_2)=e_1+e_3,\qquad C^2(e_1+e_2)=e_1+e_4.
\]
The pair $(\mu,C^2\mu)$ is the translate of
$(\theta,\delta)$ by $s_2$.  Since $h=12$ and
$2(11-5)=12$, (EX.3) holds with $d=2$.  Formula (EX.4) supplies
nonopposite containing Borels.  This proves the claimed failure
of the all-choice assertion in type $F_4$.  $\square$

\subsection{BR-123-OPPE8: a counterexample in type \texorpdfstring{$E_8$}{E8}}
Use the roots $\pm e_i\pm e_j$ and the half-vectors
$\frac12\sum_i\varepsilon_i e_i$ with $\varepsilon_i=\pm1$
and an even number of minus signs.  Number the simple roots by
\[
\begin{split}
 a_1&=\tfrac12(e_1-e_2-e_3-e_4-e_5-e_6-e_7+e_8),\\
 a_2&=e_1+e_2,\qquad a_3=e_2-e_1,\\
 a_i&=e_{i-1}-e_{i-2}\quad(4\leq i\leq8).
\end{split}
 \tag{EX.6}
\]
Every root has squared length two.  The highest root and its
adjacent root are
\[
 \begin{split}
 \theta&=e_7+e_8
       =2a_1+3a_2+4a_3+6a_4+5a_5+4a_6+3a_7+2a_8,\\
 \delta&=e_6+e_8=\theta-a_8.
 \end{split}
\]
Inspection of the explicit root set shows that $\theta+a_i$
is never a root and that the sole root among $\delta+a_i$ is
$\delta+a_8=\theta$.  In the half-integral case the putative
sum has an entry of absolute value $3/2$; in the integral cases
it has too many nonzero entries or an entry of absolute value two.
The product of reflections in
$e_1-e_2,e_1+e_2,\ldots,e_7-e_8,e_7+e_8$ is $-1$.
The highest-wedge hypotheses therefore hold again.

For the first torus take
$\alpha=e_8-e_1$ and $\beta=\theta$.
Apply the reflections $s_7,s_6,s_5,s_4,s_2$ in that order to
$(\theta,\delta)$.  They all fix $\theta$ and send its other
member successively through
\[
 e_6+e_8,\ e_5+e_8,\ e_4+e_8,\ e_3+e_8,
                    \ e_2+e_8,\ e_8-e_1.
\]
This proves the required pair conjugacy.  The height functional
is $(0,1,2,3,4,5,6,23)$, because it takes value one on every
simple root in (EX.6).  Thus the two heights are $23$ and $29$.

For the second torus put $C=s_1s_2\cdots s_8$ and
$\mu=e_2+e_8$.  The reflection formula
$s_i(v)=v-(v,a_i)a_i$ gives
\[
 C^j\mu=e_{j+2}+e_8\quad(0\leq j\leq5).
\]
The pair $(\mu,C^5\mu)=(e_2+e_8,\theta)$ is obtained from
$(\theta,\delta)$ by $s_7,s_6,s_5,s_4$ in application order.
Here $h=30$, and $5(29-23)=30$.  The obstruction applies with
$d=5$, giving nonopposite containing Borels in type $E_8$.
This proves the claim.  $\square$

\paragraph{Exact supplementary certificate.}
The script \nolinkurl{research/root_exceptional_wedge_check.py}
uses only integers and rational numbers.  It checks weighted Gram
symmetry of the listed Cartan matrices, root counts $48,240$, the raising tests for the
highest wedges, the explicit Weyl words and Coxeter powers, and
the height divisibility.  It also enumerates the highest-pair
orbits, of sizes $96$ and $6720$.  The proofs above need only
the displayed reflection calculations; the enumeration is an
additional finite check, not a numerical substitute for the
highest-line and invariant-pairing argument.

\paragraph{Reference.}
Milne, J. S. (2017). \emph{Algebraic groups: The theory of group
schemes of finite type over a field}. Cambridge University Press.
\url{https://www.jmilne.org/math/Books/iAG2017.pdf}.
Proposition 22.17 and Theorems 22.18--22.19, pp.~467--468;
Proposition 22.41, p.~477; exceptional root systems,
\S24(l), pp.~539--542.  Kostant's exact version and locators
are the same as those already given in CC-APG01--08.

% END root_exceptional_opposition.tex

% BEGIN root_g2_opposition.tex
\section{BR-123-OPPG2: every containing Borel pair is opposite in type \texorpdfstring{$G_2$}{G2}}

For the original question in cote 123, physical page 22/footer 21,
the answer to ``all choices'' is affirmative in type $G_2$.
We prove this using both fundamental representations, keeping the
restriction from the ambient orthogonal group explicit.

\subsection{The seven-dimensional cyclic model}
Let $\alpha$ be the short simple root and $\beta$ the long one,
normalized by
\[
 (\alpha,\alpha)=2,\qquad(\beta,\beta)=6,
                  \qquad(\alpha,\beta)=-3.
\]
The fundamental weights are $\omega_1=2\alpha+\beta$ and
$\omega_2=3\alpha+2\beta=\theta$.  Let $V=L(\omega_1)$.
The Weyl dimension formula, Etingof (n.d.), Proposition 28.1,
p.~125, gives
\[
 \dim V=\frac21\frac33\frac54\frac75\frac96\frac{12}9=7.
 \tag{G2.1}
\]
Here the positive roots in order are
$\alpha,\beta,\alpha+\beta,2\alpha+\beta,3\alpha+\beta,
3\alpha+2\beta$, and $\rho=5\alpha+3\beta$.
The Weyl orbit of $\omega_1$ consists of the six short roots,
each of multiplicity one.  The single remaining weight must be
Weyl-fixed, hence zero.  Thus all seven weights have multiplicity
one.  Since $w_0=-1$, the module is self-dual by Etingof,
Proposition 33.5, p.~148.  Its nondegenerate invariant bilinear
form is symmetric: its transpose differs by a scalar by Schur's
lemma, so it is symmetric or alternating, and a nondegenerate
alternating form cannot have odd dimension.  Denote this form by $B$.

Choose a principal $\mathfrak{sl}_2$ triple with raising element
$N=e_\alpha+e_\beta$ and semisimple element $2\rho^\vee$;
the construction is Etingof, \S33.3, p.~149.  The weights of
$2\rho^\vee$ on $V$ are $-6,-4,-2,0,2,4,6$.
The highest-weight-six $\mathfrak{sl}_2$ summand already has
dimension seven, so the restriction is that single irreducible
module.  Consequently a weight basis $e_0,\ldots,e_6$ can be
chosen with
\[
 N=\sum_{i=0}^{5}E_{i+1,i},\qquad
 B(e_i,e_j)=(-1)^i\delta_{i+j,6}.
 \tag{G2.2}
\]
Indeed the first normalization uses the nonzero successive raising
coefficients, and invariance of $B$ then gives alternating successive
anti-diagonal entries; scaling $B$ sets the first entry to one.
The weight list, in this basis order, is
\[
 -(2\alpha+\beta),\quad-(\alpha+\beta),\quad-\alpha,
 \quad0,\quad\alpha,\quad\alpha+\beta,\quad2\alpha+\beta.
 \tag{G2.3}
\]

The representation embeds $\mathfrak g_2$ into
$\mathfrak{so}(V,B)$, since its kernel is an ideal and the
representation is nontrivial.  The grade $-5$ part of the latter
Lie algebra for $\rho^\vee$ is one-dimensional, spanned by
\[
                         Y=E_{0,5}+E_{1,6}.
 \tag{G2.4}
\]
To check this directly, a matrix entry of grade $-5$ has row
index five less than its column index, so only $(0,5)$ and
$(1,6)$ occur; the skew-adjointness equation for (G2.2) equates
their coefficients.  The nonzero negative-highest-root space
of $\mathfrak g_2$ has precisely this grade and therefore maps
onto this line.  In particular $X=N+Y$ belongs to
$\mathfrak g_2$ and has all simple and negative-highest-root
coefficients nonzero.  Kostant (2010), Theorems 1.2--1.3,
p.~6, identifies its centralizing torus $T_1$ as a torus in
apposition to the original weight torus $T_0$.

Put $\zeta=e^{2\pi i/6}$, choose $a$ with $a^6=2$, and set
$\lambda_j=a\zeta^j$ for $j\in\mathbf Z/6$.  The seven
eigenvectors are
\[
 v_j=e_0+\sum_{i=1}^{5}\lambda_j^{6-i}e_i+e_6,
       \quad Xv_j=\lambda_jv_j,
       \qquad w=e_0-e_6,\quad Xw=0.
 \tag{G2.5}
\]
These identities follow by multiplication using (G2.2)--(G2.4):
the row-one equation is $2=\lambda_j^6$, the interior equations
shift a power, and the boundary equations are immediate.
The distinct eigenvalues prove that these vectors form a basis.
The lines $\mathbf C v_j$ are exactly the six nonzero weight
lines of $T_1$, and $\mathbf Cw$ is its zero-weight line.

The principal element $c=\rho^\vee(\zeta)$ satisfies
$cXc^{-1}=\zeta X$, so $c$ permutes these six lines cyclically
by $j\mapsto j-1$.  On $T_1$ it acts as a Coxeter element.
In the rank-two $G_2$ root system a Coxeter element rotates the
short-root hexagon by $60$ degrees in one orientation: the
simple-reflection product has order six and trace one on the
root plane, as follows immediately from the above inner products.
Thus consecutive indices $j,j+1$ correspond to adjacent short
roots in the hexagon.  Their sum is a long root.

\subsection{Both fundamental highest-line pairings}
Under the canonical equivariant identification
$\mathfrak{so}(V,B)\simeq\bigwedge^2V$, the adjoint
representation of $G_2$ is an irreducible subrepresentation
of highest weight $\omega_2=\theta$.  Each long-root weight
occurs just once in $\bigwedge^2V$, from the unique unordered
pair of short roots whose sum is that long root.  In the
$T_0$ weight basis these six lines are
\[
 \mathbf C(e_i\wedge e_k),\qquad
 (i,k)\in\{(0,1),(0,2),(1,4),(2,5),(4,6),(5,6)\}.
 \tag{G2.6}
\]
This list is verified by adding the weights in (G2.3).
In the $T_1$ weight basis they are exactly
\[
                         \mathbf C(v_j\wedge v_{j+1}),
                         \qquad j\in\mathbf Z/6.
 \tag{G2.7}
\]
Indeed the six adjacent sums are the six distinct long roots,
and each has multiplicity one, so its wedge line is the
corresponding line in the adjoint subrepresentation.

The form $B$ induces
\[
 B_2(u\wedge u',v\wedge v')
        =B(u,v)B(u',v')-B(u,v')B(u',v)
 \tag{G2.8}
\]
on $\bigwedge^2V$.  Its restriction to the adjoint
subrepresentation is nondegenerate: it pairs the lines
$e_5\wedge e_6$ and $e_0\wedge e_1$ nontrivially,
so it is a nonzero invariant form on an irreducible module.

First, all entries $B(e_i,v_j)$ with
$i\in\{0,1,2,4,5,6\}$ are nonzero by (G2.2) and (G2.5).
Consequently every pair of highest lines for the two tori in
the $\omega_1$ representation has nonzero pairing.

For the $\omega_2$ representation, write
\[
 (r_0,r_1,r_2,r_3,r_4,r_5,r_6)=(0,5,4,3,2,1,0).
\]
The coefficient of $e_i$ in $v_j$ is $\lambda_j^{r_i}$.
For $(i,k)$ in (G2.6), (G2.8) is a nonzero sign times
\[
 \det\begin{pmatrix}
  \lambda_j^{r_{6-i}}&\lambda_{j+1}^{r_{6-i}}\\
  \lambda_j^{r_{6-k}}&\lambda_{j+1}^{r_{6-k}}
 \end{pmatrix}
  =\lambda_j^{r_{6-i}+r_{6-k}}
        (\zeta^{r_{6-k}}-\zeta^{r_{6-i}}).
 \tag{G2.9}
\]
The residues $r_s$ are distinct except for $s=0,6$,
and none of the pairs in (G2.6) is $(0,6)$.
Thus (G2.9) never vanishes.  Every pair of highest lines
in the second fundamental representation also has nonzero
pairing.

We finish by proving why these two tests suffice.  Fix a
Borel $B_+$ and highest vectors $z_1,z_2$ in its two
fundamental representations.  Their invariant forms pair
opposite weights nondegenerately.  For another Borel
$gB_+g^{-1}$ write $g=b_1n_ub_2$ by the Bruhat
decomposition, with $b_1,b_2\in B_+$ and $u\in W$.
Since $B_+$ fixes each highest line, invariance gives
\[
             (z_i,gz_i)=\text{nonzero scalar}\,(z_i,n_uz_i).
\]
The last pairing is nonzero exactly when
$u\omega_i=-\omega_i$.  Therefore nonzero pairings for
both $i=1,2$ force $u=-1=w_0$, because the fundamental
weights form a basis of the root plane.  This is exactly
the opposite-Borel Bruhat cell.  Conversely that cell
has the stated nonzero pairings.  This argument uses
Milne (2017), Theorem 21.80, p.~450, for Bruhat
decomposition and Theorem 21.84, p.~451, for opposite Borels.

For any $B_0\supset T_0$ and $B_1\supset T_1$, their
fundamental highest lines are among the lines just tested.
Thus $B_0,B_1$ are opposite.  CC-APG01 transfers this result
from the constructed ordered pair to every apposition pair
in type $G_2$.  This proves the assertion. $\square$

\paragraph{Primary representation-theoretic reference.}
Etingof, P. (n.d.). \emph{18.745: Lie groups and Lie algebras, I}
[Course notes, 223-page author PDF]. Massachusetts Institute
of Technology. \url{https://math.mit.edu/~etingof/lnlg.pdf}.
Proposition 28.1, p.~125; Proposition 33.5, p.~148;
\S33.3, p.~149.  The finite dimension product and the
matrix and minor calculations in this section are carried
out here; no result about apposition in this source is assumed.

% END root_g2_opposition.tex

% BEGIN consequences_criteria_exceptional_minuscule.tex
\section{CC-APG09--11: short extremal-weight orbits obstruct opposition in types \(E_6,E_7\)}

All groups in this section are over \(\mathbf C\). An ordered pair
\((T_0,T_1)\) is in principal apposition in the sense fixed in
CC-APG01: a principal element \(c\in T_0\) of order the Coxeter
number \(h\) normalizes \(T_1\) and acts on it as a Coxeter element.
The group \(G\) is adjoint. Representations below are taken on its
simply connected cover \(\widetilde G\); highest lines and Borel
subgroups descend through the central isogeny. No assertion that
the representations descend linearly to \(G\) is used.

\subsection{CC-APG09: the short-orbit criterion}

\paragraph{Statement.}
Choose a positive system for \(T_0\) in which
\[
        c=\rho^\vee(\zeta),\qquad \zeta=e^{2\pi i/h};
        \qquad \alpha_j(c)=\zeta\quad\hbox{for every simple root}.
\]
Let \(V=L(\lambda)\) be the irreducible representation of
\(\widetilde G\) with a fundamental highest weight
\(\lambda=\omega_i\). Transport the root labels to \(T_1\) by an
inner conjugacy. Suppose the Coxeter action \(C\) on the extremal
weight orbit \(W\lambda\) has an element \(\mu\) with
\[
                         C^d\mu=\mu,\qquad 0<d<h.
\]
Then some \(B_0\supset T_0\) and \(B_1\supset T_1\) are not opposite.
The hypothesis concerns only the extremal orbit; it makes no
assumption about other weights or the dimension of \(V\).

\paragraph{Proof.}
Existence of \(V\), uniqueness of its highest line, and Weyl
invariance of weight multiplicities are the standard highest-weight
theorems: Milne (2017), Lemma 22.15, pp.~466--467,
Proposition 22.17, p.~467, and Theorems 22.18--22.20,
pp.~467--468. In particular, every extremal
weight space \(V_\nu\), \(\nu\in W\lambda\), is one-dimensional and
is the highest line for a Borel containing the corresponding torus.

Choose any lift \(\widetilde c\) of \(c\) to
\(\widetilde T_0\). It normalizes \(\widetilde T_1\).
The line \(V_\mu\) for \(T_1\) is preserved by
\(\widetilde c^d\), because \(C^d\mu=\mu\).
Thus there is \(b\in\mathbf C^\times\) such that
\(\widetilde c^d u=bu\) for \(0\ne u\in V_\mu\).
Choose \(B_1\supset T_1\) for which this is the highest line.

For \(T_0\), the two extremal weights
\[
               \nu_0=\lambda,\qquad
               \nu_1=s_i\lambda=\lambda-\alpha_i
\]
give eigenvalues of \(\widetilde c^d\) whose ratio is
\(\alpha_i(c)^d=\zeta^d\ne1\).
Hence at least one of them, say \(\nu\), has eigenvalue
\(a=\nu(\widetilde c^d)\ne b\).
Let \(f\) be a nonzero element of the one-dimensional dual weight
space \(V^*_{-\nu}\). Equivalently, \(f\) is supported on \(V_\nu\)
and annihilates all other \(T_0\)-weight spaces. Consequently
\[
       f(\widetilde c^d u)=a f(u)=b f(u),
                  \qquad\hbox{so}\qquad f(u)=0.
\]
The weights \(-W\lambda\) form the extremal orbit of \(V^*\).
Therefore the line of \(f\) is the highest line of \(V^*\) for
some \(B_0\supset T_0\).

If \(B_0,B_1\) were opposite, their highest lines in \(V^*,V\)
would have nonzero natural evaluation. Indeed, conjugate the
pair to a standard pair \(B_+,B_-\). If the highest weight of
\(V\) for \(B_+\) is \(\eta\), its highest line for \(B_-\)
is \(V_{w_0\eta}\), and the highest line of \(V^*\) for \(B_+\)
is \(V^*_{-w_0\eta}\). These one-dimensional dual spaces pair
nondegenerately. This contradicts \(f(u)=0\).
Thus \(B_0,B_1\) are not opposite. \(\square\)

\paragraph{Construction with the specified coordinates.}
For each application, start with \(T_1\) and its explicitly labeled
base, and choose a representative \(c\in N_G(T_1)\) of the
specified product \(C=s_1\cdots s_r\). The converse principal-class
statement in Kostant (2010), p.~7, says that this representative
is regular of order \(h\) and belongs to the principal conjugacy
class. Put \(T_0=Z_G(c)^\circ\), and choose a positive base for
\(T_0\) in which all simple roots take the value \(\zeta\) on \(c\).
Choose the representation with the indicated fundamental highest
weight for the original based torus \(T_1\). Inner conjugacy of
based torus pairs transports this to a fundamental highest weight
for the chosen based torus \(T_0\); its index there is immaterial
to CC-APG09. Thus the displayed Coxeter word and the principal
coordinates refer to the same representation and apposition pair.
No conjugacy theorem for Coxeter elements is needed in this
construction. A change of lift of \(c\) only multiplies all relevant
eigenvalues by a common central scalar. Finally CC-APG01 transfers
the resulting counterexample to every ordered apposition pair.

\subsection{Exact finite certificates}

Use fundamental-weight coordinates
\[
       v=(\langle v,\alpha_1^\vee\rangle,\ldots,
                          \langle v,\alpha_r^\vee\rangle).
\]
The simply laced Cartan matrix has diagonal entries \(2\), entry
\(-1\) on the edges
\[
     (1,3),(2,4),(3,4),(4,5),(5,6),
          \quad\hbox{and additionally }(6,7)\hbox{ for }E_7,
\]
and zero elsewhere. This specifies the numbering completely;
it agrees with Milne (2017), Theorem C.56 and the following
Dynkin diagrams, pp.~625--626.
For this symmetric matrix \(A\), reflections are
\[
               s_j(v)_k=v_k-v_j A_{jk}.                 \tag{E.1}
\]
Throughout \(C=s_1\cdots s_r\), so \(s_r\) is applied first.
Listed reflection words are instead explicitly in application
order. Each equality below can therefore be checked by integer
substitution into (E.1), without a representation matrix.

\subsection{CC-APG10: \(E_6\)}

Take \(\lambda=\omega_1=(1,0,0,0,0,0)\) and \(h=12\).
The reflection word in application order
\[
                           1,3,4,2,5,4
\]
sends \(\lambda\) to
\(\mu=(0,0,1,-1,0,1)\).
The complete short cycle is
\[
\begin{split}
 (0,0,1,-1,0,1)
     &\xmapsto{C}(-1,0,0,1,-1,0)\\
     &\xmapsto{C}(1,0,-1,0,1,-1)
       \xmapsto{C}(0,0,1,-1,0,1).
\end{split}                                                   \tag{E.2}
\]
Thus \(\mu\in W\lambda\) and \(C^3\mu=\mu\), with \(0<3<12\).
CC-APG09 gives a nonopposite containing Borel pair.
The single ordered-apposition-pair orbit proved in CC-APG01
transfers this conclusion to every pair in apposition in type
\(E_6\). In particular, the answer to the original
\emph{all choices} question is negative. \(\square\)

\subsection{CC-APG11: \(E_7\)}

Take \(\lambda=\omega_7=(0,0,0,0,0,0,1)\) and \(h=18\).
The reflection word in application order
\[
                         7,6,5,4,2,3,1,4,3,5,4,6
\]
sends \(\lambda\) to
\(\mu=(0,1,0,-1,1,-1,1)\).
The short cycle is
\[
  (0,1,0,-1,1,-1,1)
     \xmapsto{C}(0,-1,0,1,-1,1,-1)
     \xmapsto{C}(0,1,0,-1,1,-1,1).                         \tag{E.3}
\]
Thus \(C^2\mu=\mu\), with \(0<2<18\).
CC-APG09 again gives a nonopposite containing Borel pair,
and CC-APG01 transfers it to every apposition pair.
The original \emph{all choices} question is negative in type
\(E_7\). \(\square\)

\paragraph{Scope of the computation.}
The script \nolinkurl{research/cc_minuscule_orbits.py} uses only
the integer reflection formula (E.1). It checks the full Coxeter
matrix orders \(12,18\) on a basis, and independently enumerates
the extremal Weyl orbits and obtains cycle lengths
\[
                  E_6:\ 3+12+12,\qquad
                  E_7:\ 2+18+18+18.
\]
The proofs above need only the displayed words and short cycles.
They do not use the full enumerations, the names or dimensions
of minuscule representations, a numerical approximation, or a
claim that the computational output is a literature theorem.
The usual root-system Coxeter numbers are also directly compatible
with the highest-root coefficient lists
\((1,2,2,3,2,1)\) and \((2,2,3,4,3,2,1)\), respectively:
their sums are \(11\) and \(17\).

\paragraph{Primary inputs and bibliography.}
The apposition and principal-element inputs are Kostant (2010),
Proposition 1.1, p.~5; Theorems 1.2--1.4, p.~6; the principal
order-\(h\) paragraph and Theorem 1.7, p.~7; and
Proposition 1.8 and Theorem 1.9, p.~8.
The representation-theoretic inputs are the numbered Milne
statements cited above; the exceptional root systems are described
in his \S24(l), pp.~541--542, and the labeled Dynkin diagrams in
Theorem C.56 and the following figure, pp.~625--626. These are retained as known
inputs, not re-proved literature claims.

Kostant, B. (2010). Experimental evidence for the occurrence of
\(E_8\) in nature and the radii of the Gosset circles.
\emph{Selecta Mathematica, 16}, 419--438.
\url{https://doi.org/10.1007/s00029-010-0023-1}.
Author version:
\url{https://arxiv.org/abs/1003.0046v1}.

Milne, J. S. (2017). \emph{Algebraic groups: The theory of group
schemes of finite type over a field}. Cambridge University Press.
Author PDF: \url{https://www.jmilne.org/math/Books/iAG2017.pdf}.

% END consequences_criteria_exceptional_minuscule.tex

% BEGIN root_effectivity.tex
\section{Spectral effectivity and the full Tate conjecture: the precise numerical-motive equivalence}

The SA06 intake, row D.4f and its successor-work instruction at
lines 501--503, asks whether effectivity detected by Frobenius is
equivalent to the Tate conjecture for all motives.  Its earlier proof
file is absent, but this question has the following precise answer in
the category of rational numerical motives.  This is a deduction from
classical results, not a new solution of the Tate conjecture.

Fix $k=\mathbb F_q$.  Let $\mathcal N(k)$ be the pseudo-abelian
category of pure numerical motives with rational coefficients and all
integer Tate twists.  Call an object \emph{effective} if it is a
direct summand of a finite direct sum of motives $h(X)$ with $X/k$
smooth projective, without a positive Tate twist.  The Frobenius
endomorphism $\pi$ is central.  On a simple object $S$, the algebra
$\mathbb Q[\pi_S]$ is a number field, so $\pi_S$ defines an algebraic
number up to conjugacy.  Call $M$ \emph{spectrally integral} if all
these numbers for its simple constituents are algebraic integers.
This definition does not presume that an $\ell$-adic realization
factors through numerical equivalence.

The structural inputs are Jannsen's semisimplicity theorem and
Frobenius centrality, in the exact formulations of Geisser (1998),
Theorem 2.1, p.~111, and Propositions 2.3--2.4, p.~112.
The decisive known result is Geisser, Theorem 1.1, p.~109,
equivalently Theorem 2.7, p.~114: the full Tate conjecture over $k$
is equivalent to the assertion that a simple numerical motive whose
Frobenius is $1$ is the unit.  We use this theorem as a cited input.
Here \emph{full Tate} means equality of the pole order at $s=r$
and the dimension of codimension-$r$ numerical cycle classes, for
all smooth projective $X/k$ and all $r$; it is not merely the
surjectivity assertion for one variety and one degree.

\subsection{BR-EFF01. The equivalence and its weight-zero test}
The following assertions are equivalent:
\begin{enumerate}
\item The full Tate conjecture holds over $k$.
\item Every spectrally integral object of $\mathcal N(k)$ is effective.
\item Every object $M$ of $\mathcal N(k)$ with $\pi_M=1_M$ is effective.
\item Every simple object $S$ with $\pi_S=1_S$ is effective.
\end{enumerate}
For every effective object, spectral integrality holds without
any Tate-conjecture assumption.  Thus item 2 can equivalently be
written ``effective if and only if spectrally integral.''

\paragraph{Complete deduction.}
For smooth projective $X$, its cohomological Frobenius characteristic
polynomials are monic integral polynomials, and their product annihilates Frobenius
in the category of homological motives.  They therefore annihilate
it numerically.  The minimal polynomial on a numerical summand
divides such a polynomial, so all its roots are algebraic integers.
This proves the unconditional necessity.

Assume item 4, and let $S$ be simple with $\pi_S=1$.  Since $S$
is effective, it is a summand of some finite sum of $h(X)$.
The algebraic K\"unneth projectors over $\mathbb F_q$ give
$h(X)=\bigoplus_{i\geq0}h^i(X)$; the precise cited input is
Katz--Messing (1974), Theorem 2(1), pp.~76--77.
For $i>0$, a polynomial annihilating $\pi_{h^i(X)}$ has all roots
of absolute value $q^{i/2}>1$.  In particular $\pi_{h^i(X)}-1$
is invertible: a B\'ezout identity between its annihilating
polynomial and $T-1$ constructs the inverse.  Centrality of
Frobenius now forces every map $S\to h^i(X)$ and $h^i(X)\to S$
to be zero.  Hence $S$ is a summand of a finite sum of $h^0(X)$.

These degree-zero objects are Artin motives: $h^0(X)$ is the
motive of the finite \'etale scheme of geometric connected
components.  For a geometrically connected component the
degree-zero projector is obtained from a degree-one rational
zero-cycle, obtained by dividing any closed point by its degree.
For a component with constant field $\mathbb F_{q^a}$ this gives
$h(\operatorname{Spec}\mathbb F_{q^a})$.  In its permutation
representation Frobenius is a cyclic permutation, and the
$1$-eigenspace is the constant line.  The rational projector is
$a^{-1}\sum_{j=0}^{a-1}\pi^j$.  Thus the part with Frobenius $1$
in any finite sum of such Artin objects is a sum of units.  The
simple object $S$ is consequently the unit.  Geisser's cited
theorem proves item 1.

Suppose now that item 1 holds.  The same cited theorem says that
simple numerical motives are determined by their Frobenius number.
Let $S$ be simple and spectrally integral, with Frobenius number
$\alpha$.  Its weight is an integer $m$, and all conjugates have
absolute value $q^{m/2}$.  Because $\alpha$ is a nonzero algebraic
integer, the norm shows $m\geq0$.

If $m=0$, Kronecker's theorem says that $\alpha$ is a root of
unity.  The rational Artin representation for its cyclotomic
polynomial is a summand of the regular representation of a finite
cyclic quotient of $\operatorname{Gal}(\overline k/k)$, hence is
an effective Artin motive.  Its simple constituent with Frobenius
$\alpha$ is isomorphic to $S$ by Geisser's characterization.

If $m>0$, choose an algebraic root $\beta$ with $\beta^m=\alpha$.
It is integral and each conjugate has absolute value $q^{1/2}$.
Honda--Tate supplies an abelian variety $B/k$ having $\beta$ among
its $H^1$ Frobenius eigenvalues.  The tensor power
$h^1(B)^{\otimes m}$ is an effective summand of $h(B^m)$, and its
realization has $\alpha$ among its eigenvalues.  Under full Tate,
numerical and homological equivalence agree and realization is
faithful.  Its decomposition into simple numerical summands
therefore contains a summand $S_\alpha$ with Frobenius number
$\alpha$.  Geisser's characterization gives $S\simeq S_\alpha$.
Thus $S$ is effective.  A finite direct sum of effective motives
is effective, proving item 2.  The implications $2\Rightarrow3
\Rightarrow4$ are immediate.  This proves all assertions.
\(\square\)

This proof never assumes a rational neutral fibre functor.
The quaternion obstruction in AQ-RATIONAL-FIBRE is compatible
with every step.  It also does not identify numerical and Chow
effectivity: a lift of a numerical projector to a Chow projector
requires additional information.

\subsection{BR-EFF02. The effective cone in the now established abelian sector}
Let $\mathcal A_q$ be the rational numerical category generated by
abelian, Artin, and Tate motives, with the same conventions as
CC-AG01.  Using the cited all-abelian Tate result in CC-TATE01--04,
an object of $\mathcal A_q$ is effective if and only if it is
spectrally integral.  In fact it then has an effective presentation
using only abelian varieties and finite \'etale schemes.  Over
$E=\overline{\mathbb Q}_\ell$, the equivalence of CC-AG01 identifies
this full subcategory with the finite $W(q)$-graded vector spaces
whose nonzero components are indexed by
\[
 W^+(q)=\{\alpha\in W(q):\alpha\text{ is an algebraic integer}\}.
 \tag{EFF.1}
\]
This is a tensor monoid of labels, not a group of labels.

\paragraph{Complete deduction.}
Necessity is the unconditional first paragraph of BR-EFF01.
For sufficiency, apply its construction to every simple constituent.
The required comparisons occur entirely inside $\mathcal A_q$.
Their validity follows from the all-abelian Tate theorem and its
full-faithfulness consequence AQ-ABELIAN-RECOGNITION: for simple
objects in this sector, an intersecting Frobenius spectrum gives a
nonzero morphism after scalar extension, hence a nonzero rational
morphism, and such a morphism is an isomorphism by semisimplicity.
In weight zero use the finite \'etale construction; in positive
weight use $h^1(B)^{\otimes m}$ as above.  These constructions
stay effective before any coefficient extension.

After extension to $E$, all rank-one $W(q)$-labels split by
CC-AG01.  Each label in $W^+(q)$ is effective by the same
construction, and each effective label is integral by necessity.
Finite sums and summands give (EFF.1).  Multiplication preserves
integrality, so $W^+(q)$ is a monoid.  For example $q\in W^+(q)$
whereas $q^{-1}\notin W^+(q)$, so it is not a group.
\(\square\)

\subsection{BR-EFF03. The slope test and the essential Artin correction}
For a Weil $q$-number $\alpha$ of weight $m$, spectral integrality
is equivalent to
\[
 \frac{v(\alpha)}{v(q)}\geq0
 \quad\text{for every finite place }v\mid p
 \text{ of a number field containing }\alpha.
 \tag{EFF.2}
\]
Consequently (EFF.2) is an exact effectivity test in the abelian
sector of BR-EFF02.  For all numerical motives the corresponding
universal assertion is equivalent to full Tate by BR-EFF01.

\paragraph{Complete proof.}
The definition of a Weil number gives $q^N\alpha$ integral for
some $N$.  At places away from $p$, this implies $v(\alpha)\geq0$.
Choose a normal field containing $\alpha$ and its conjugates.
For a complex-conjugation automorphism $c$ one has
$\alpha c(\alpha)=q^m$.  The same nonnegativity holds for
$c(\alpha)$ away from $p$, and their valuations sum to zero
there.  Both valuations are therefore zero.  Hence the only
possible denominators of $\alpha$ occur over $p$, and (EFF.2)
is precisely their absence.  This proves the equivalence.
\(\square\)

There is a small but necessary source correction in using an
effective realization theorem.  Milne--Ramachandran (2009),
Proposition 1.7 in the author final version of May 2, 2008,
p.~4, is printed for every effective semisimple Tate structure
with an abelian variety as its sole source.  Its unrestricted
weight-zero wording has the following counterexample.
The one-dimensional space with Frobenius $-1$ is effective,
semisimple, and of weight zero.  For any abelian variety $B/k$,
$H^0(B,\mathbb Q_\ell)$ has Frobenius $1$, whereas every
positive-degree eigenvalue has absolute value greater than one.
Thus $-1$ never occurs in $H^*(B,\mathbb Q_\ell)$.
The space does occur in
$H^0(\operatorname{Spec}\mathbb F_{q^2},\mathbb Q_\ell)$, as
the anti-invariant summand.  This proves the counterexample and
the required repair in that case.  The general weight-zero repair
is the finite cyclic Artin construction in BR-EFF01; the positive
weight construction there is complete.  The effective targets of
CC-XC01 already include these Artin summands.  No refutation of
the generalized Tate conjecture is inferred from this wording
correction.

\paragraph{Exact sources.}
\begin{itemize}
\item Geisser, T. (1998). Tate's conjecture, algebraic cycles and
rational $K$-theory in characteristic $p$. \emph{$K$-Theory,
13}(2), 109--122.
\url{https://www2.rikkyo.ac.jp/web/geisser/Tate_KtheoryJournal.pdf}.
Theorem 1.1, p.~109; Theorem 2.1 and Corollary 2.2, p.~111;
Propositions 2.3--2.4, p.~112; Theorem 2.7, p.~114.
\item Milne, J. S., \& Ramachandran, N. (2009). Motivic complexes
over finite fields and the ring of correspondences at the generic
point. \emph{Pure and Applied Mathematics Quarterly, 5}(4),
1219--1252. Operative author final version, May 2, 2008:
\url{https://www.jmilne.org/math/articles/2009a.pdf}.
Proposition 1.7, p.~4; Remark 1.12, p.~6; Lemma 2.1 and
Proposition 2.2, p.~8. The weight-zero correction above is explicit.
\item Milne, J. S. (2007, September 19). \emph{The Tate conjecture
over finite fields (AIM talk)}.
\url{https://www.jmilne.org/math/xnotes/TateAimv1.pdf}.
Section 1, pp.~1--4, particularly footnote 4, p.~4, on motivic
effectivity and generalized Tate. This contextual statement is
not substituted for the precise numerical-motive theorem proved here.
\end{itemize}

% END root_effectivity.tex

% BEGIN root_group_rank.tex
\section{The group quotient in cote 123: central stabilizers and an exact type-A rank formula}

The rank question in the letter, cote 123 physical page 23, footer 22,
mentions both the Lie-algebra quotient and $G\to T/W$ for the original
complex adjoint group.  The simply connected qualification occurs later,
for rational singularities.  These scopes matter: a regular element of
an adjoint group can have quotient differential of rank zero.  The
following complete deductions give a finite formula in every semisimple
type and a closed matrix formula in type A.

\subsection{BR-12305. The differential of a finite linear quotient}
Let a finite group $Z$ act linearly on a finite-dimensional complex
vector space $V$, let $\delta:V\to V/Z$ be the affine quotient, and
let $H=Z_x$ be the stabilizer of $x\in V$.  Identify $T_xV$ with $V$.
Then
\[
 \operatorname{im}(d\delta_x)^*=(V^*)^H.
 \tag{GR.1}
\]
Consequently, for any linear map $L:U\to V$, the rank of
$d\delta_x\circ L$ is the rank of the restrictions to $L(U)$ of
the $H$-invariant linear forms on $V$.

\paragraph{Complete proof.}
Every invariant polynomial has an $H$-invariant differential at $x$,
since $H$ fixes $x$.  This proves one inclusion in (GR.1).
Conversely, let $\lambda\in(V^*)^H$.  The maximal ideals of the
finitely many distinct points of $Zx$ are pairwise comaximal, as are
their squares.  The Chinese remainder theorem supplies a polynomial
$e$ with $e\equiv1\pmod{\mathfrak m_x^2}$ and
$e\equiv0\pmod{\mathfrak m_y^2}$ for every $y\in Zx$, $y\ne x$.
Put
\[
 f(v)=\frac1{|H|}\sum_{z\in Z}
       \lambda(z^{-1}v-x)e(z^{-1}v).
\]
Reindexing the sum shows that $f$ is $Z$-invariant.  At $v=x$,
the terms with $z\notin H$ have zero first jet.  A term with
$z\in H$ has differential $\lambda\circ z^{-1}=\lambda$;
its factor $\lambda(z^{-1}x-x)$ is zero, so differentiating $e$
adds nothing.  Therefore $df_x=\lambda$.  This proves equality.
The image of the dual of a composite is the restriction of this
image along $L$, so its dimension gives the final assertion.
\(\square\)

No smoothness of $V/Z$ at $\delta(x)$ was assumed.  Its Zariski
tangent space can have dimension larger than $\dim V$; (GR.1)
still computes the actual differential rank.

\subsection{BR-12306. A finite character formula for every central quotient}
Let $\widetilde G$ be a connected simply connected complex
semisimple group of rank $r$, let $Z$ be a subgroup of its center,
and put $G=\widetilde G/Z$.  For the fundamental representations
$\rho_i$ of $\widetilde G$, write
$\chi_i=\operatorname{Tr}\rho_i$ and let $\xi_i:Z\to\mathbb C^*$
be their central characters.  Steinberg (1965), Theorem 6.1(b),
p.~62, and the quotient discussion (6.9)--(6.10), pp.~63--64,
give the polynomial coordinates
\[
 \widetilde\chi=(\chi_1,\ldots,\chi_r):
 \widetilde G\longrightarrow\mathbb A^r
       =\widetilde T/W.
\]
For $g\in G$, choose a lift $\widetilde g$, put
$x_i=\chi_i(\widetilde g)$, and define
\[
 H_x=\{z\in Z:\xi_i(z)x_i=x_i\text{ for every }i\},
 \qquad J_x=\{i:\xi_i|_{H_x}=1\}.
 \tag{GR.2}
\]
If $\chi_G:G\to T/W$ is the group quotient, then
\[
 \operatorname{rank}d\chi_G|_g
 =\dim_{\mathbb C}\operatorname{span}
 \left\{Y\longmapsto
 \operatorname{Tr}\bigl(\rho_i(\widetilde g)d\rho_i(Y)\bigr)
 :i\in J_x\right\}\subseteq\mathfrak g^*.
 \tag{GR.3}
\]
This is an explicit finite formula: evaluate the fundamental
characters, determine the finite stabilizer (GR.2), and compute
the rank of the displayed $|J_x|$ linear forms on a Lie-algebra
basis.  It applies to arbitrary elements, including those with
nontrivial unipotent part.

\paragraph{Complete proof.}
The central isogeny $\widetilde G\to G$ is \'etale in characteristic
zero, and its differential is an isomorphism.  Its action on the
fundamental-character coordinates is diagonal with weights $\xi_i$,
because a central element acts as the scalar $\xi_i(z)$ in an
irreducible representation.  Taking invariants in either order gives
\[
 \mathbb C[G]^G
 =\bigl(\mathbb C[\widetilde G]^{\widetilde G}\bigr)^Z
 =\mathbb C[x_1,\ldots,x_r]^Z.
\]
Thus the lifted quotient map is the composite of
$\widetilde\chi$ with $\delta:\mathbb A^r\to\mathbb A^r/Z$.
The $H_x$-invariant linear forms are precisely the span of $dx_i$
with $i\in J_x$.  BR-12305 gives the rank as that of these
coordinate differentials composed with $d\widetilde\chi$.

Identify $T_{\widetilde g}\widetilde G$ with $\mathfrak g$ by
$Y\mapsto\widetilde gY$.  Differentiating
$\chi_i(\widetilde g\exp(tY))$ at zero gives exactly the trace
in (GR.3).  Replacing the lift by $z\widetilde g$ leaves $H_x$
and $J_x$ unchanged and multiplies each displayed form by the
nonzero scalar $\xi_i(z)$.  The rank is therefore independent
of the lift.  This completes the proof. \(\square\)

For a regular lift, Steinberg's Theorem 8.1, p.~70, makes all
$d\chi_i$ independent.  In that case (GR.3) simplifies to
$\operatorname{rank}d\chi_G|_g=|J_x|$.  In particular the
simply connected case has full rank $r$ at every regular element,
whereas the adjoint quotient can have a smaller rank.

\subsection{BR-12307. Closed formula for all type-A central quotients}
Let $c$ divide $n$, let $G=\operatorname{SL}_n/\mu_c$, and let
$A\in\operatorname{SL}_n(\mathbb C)$ lift $g\in G$.
Write its characteristic polynomial as
\[
 \det(tI-A)=\sum_{i=0}^{n}(-1)^i c_i(A)t^{n-i},
 \qquad c_0=c_n=1.
\]
Define
\[
 d=\gcd\bigl(c,\{i:1\leq i<n,\ c_i(A)\ne0\}\bigr),
 \tag{GR.4}
\]
with the convention that the gcd of the empty index set and $c$
is $c$.  If $m_{A^d}$ is the minimal polynomial of $A^d$, then
\[
 \boxed{\operatorname{rank}d\chi_G|_g
        =\deg m_{A^d}-1.}
 \tag{GR.5}
\]
The formula covers $c=1$ (the simply connected group) and $c=n$
(the adjoint group $\operatorname{PGL}_n$).

\paragraph{Complete proof.}
For $\operatorname{SL}_n$, the fundamental characters are
$c_i=\operatorname{Tr}(\bigwedge^i A)$, $1\leq i<n$.
The scalar $\zeta I\in\mu_c$ acts with central character
$\zeta^i$.  Hence (GR.2) gives $H_x=\mu_d$ and
\[
 J_x=\{d,2d,\ldots,n-d\}.
\]
All nonzero $c_i$ have index divisible by $d$, and $d\mid n$,
so $\det(tI-A)=f(t^d)$ for a monic polynomial $f$ of degree
$h=n/d$.  The Cayley--Hamilton identity gives $f(A^d)=0$.
In particular $\deg m_{A^d}\leq h$.

To compute the differentials without an eigenvalue-genericity
assumption, differentiate the identity
$\det(I+tA)=\sum_i c_i t^i$ in the tangent direction $AY$.
The determinant derivative and the formal geometric series give
\[
 dc_i|_A(AY)=\sum_{j=1}^{i}(-1)^{j-1}c_{i-j}(A)
                                  \operatorname{Tr}(A^jY).
 \tag{GR.6}
\]
For $i=md$, a nonzero coefficient on the right forces $d\mid j$.
Writing $B=A^d$, the forms $dc_{md}$, $1\leq m<h$, are
therefore related to
$Y\mapsto\operatorname{Tr}(B^jY)$, $1\leq j<h$, by a
triangular matrix whose diagonal entries are
$(-1)^{md-1}$, and are nonzero.

The annihilator of $\mathfrak{sl}_n$ in
$\mathfrak{gl}_n$ under the trace pairing is the scalar line.
The dimension of these forms restricted to $\mathfrak{sl}_n$
is consequently
\[
 \dim\operatorname{span}\{I,B,\ldots,B^{h-1}\}-1
 =\deg m_B-1.
\]
The last equality uses $\deg m_B\leq h$ and the defining
minimality of the minimal polynomial.  Equation (GR.3) now
proves (GR.5).  No Jordan block was discarded and no
semisimplicity of $A$ was used. \(\square\)

Equivalently, $\deg m_{A^d}$ is the sum, over distinct nonzero
eigenvalues $\eta$ of $A^d$, of the largest Jordan-block size
among eigenvalues $\lambda$ of $A$ with $\lambda^d=\eta$.
Indeed the polynomial $t^d$ has nonzero derivative at each
nonzero $\lambda$, so it preserves the sizes of the local
Jordan blocks, while it can merge their eigenvalue labels.
This gives a direct formula from the Jordan data of $A$.

\subsection{BR-12308. Regular adjoint elements with rank-zero quotient differential}
For every $n\geq2$, the group $\operatorname{PGL}_n(\mathbb C)$
has a regular semisimple element $g$ with
$\operatorname{rank}d\chi_G|_g=0$.
In the simply connected group $\operatorname{SL}_n$, the
corresponding regular lift has rank $n-1$.

\paragraph{Complete counterexample.}
Fix a primitive $n$th root of unity $\zeta$, set
$Q=\operatorname{diag}(1,\zeta,\ldots,\zeta^{n-1})$, and
take $g=[Q]$.  Its eigenvalues are distinct, so its connected
centralizer is a maximal torus of dimension $n-1$.  It is
therefore regular semisimple.  Choose a scalar $a\ne0$ with
$A=aQ\in\operatorname{SL}_n$.  Its characteristic polynomial
is $t^n-a^n$, so $c_i(A)=0$ for every $1\leq i<n$.
For $G=\operatorname{PGL}_n$ one has $c=n$, hence $d=n$ in
(GR.4), and $A^n=a^n I$ has minimal polynomial of degree one.
Equation (GR.5) gives rank zero.  For the simply connected
group one instead has $c=d=1$, and the distinct eigenvalues
give $\deg m_A=n$, hence rank $n-1$.
\(\square\)

This example concerns the finite central quotient of the
\emph{group} invariant space.  It does not contradict Kostant's
regular-point criterion for the Lie-algebra quotient or Steinberg's
simply connected group theorem.  The original rank questions
must therefore keep their two quotient maps separate.

\paragraph{Primary input.}
Steinberg, R. (1965). Regular elements of semi-simple algebraic
groups. \emph{Publications Math\'ematiques de l'IH\'ES, 25},
49--80. \url{https://doi.org/10.1007/BF02684397}.
Theorem 6.1(b), p.~62; (6.9)--(6.10), pp.~63--64;
Theorem 8.1, p.~70. The finite-quotient differential lemma and
the formulas (GR.3)--(GR.6) are completely proved above.

% END root_group_rank.tex

% BEGIN consequences_criteria_adjoint_rank.tex
\section{Cote 123: exact Lie-algebra quotient ranks and the subregular question}\label{cote-123-exact-lie-algebra-quotient-ranks-and-the-subregular-question}

\textbf{Scope.} Complex semisimple Lie algebras and complex simple adjoint algebraic groups. These are unconditional characteristic-zero results; neither RH nor Dirichlet GRH is used. The original question is Grothendieck's letter to Kostant of 22 October 1969, cote 123, physical PDF page 23, handwritten/printed footer 22, visually inspected. It asks about both the Lie-algebra and group quotient maps, proposes a strict rank drop below \(r-1\) for subregular elements in nonsimply-laced types, and suggests uniqueness of the subregular orbit might characterize ADE. The two proposed distinctions are false. The group quotient away from the unipotent locus must be treated separately, as in BR-12305--08.

\textbf{Literature and proof status.} Richardson (1987) is a primary announcement with proof outlines, as he expressly says on p.~228. Its stated results are used as known inputs, not represented as independently re-proved here. Its exceptional table has three demonstrable printing/omission errors and one explicitly unresolved entry. The errors are corrected below using its own theorems. The formerly unresolved entry is explicitly settled in Dinar (2020), Corollary 5.5; a separate short deduction from Johnson--Sommers' theorem is proved below. Johnson--Sommers was inspected in arXiv:1706.04820v2, 5 February 2018, and Dinar in arXiv:1905.01829v4, 7 April 2020; both journal publications are 2020. Literature cutoff: 20 November 2026. Actual inspection date: 21 November 2026 UTC. No novelty is claimed.

\subsection{CC-RK01 --- Reduction to nilpotent elements and precise classical locators}\label{cc-rk01-reduction-to-nilpotent-elements-and-precise-classical-locators}

Let \(\chi_{\mathfrak g}:\mathfrak g\to\mathfrak t/W\) be the adjoint quotient and write \(x=s+e\) for the Jordan decomposition. Set

\[
 \mathfrak l=\mathfrak g^s,
 \qquad \mathfrak l_{\mathrm{der}}=[\mathfrak l,\mathfrak l]
       =\bigoplus_j\mathfrak l_j,
 \qquad e=(e_j)_j.
\]

Richardson's \textbf{Proposition 1.2, p.~229}, gives

\[
 \operatorname{rank}(d\chi_{\mathfrak g})_x
   =\dim\mathfrak z(\mathfrak l)
       +\sum_j\operatorname{rank}(d\chi_{\mathfrak l_j})_{e_j}.
 \tag{CC-RK.1}
\]

His notation \(\mathfrak g(s)\) in that proposition is the fixed subspace of the group centralizer; for semisimple \(s\), the centralizer is a connected Levi subgroup and that subspace is \(\mathfrak z(\mathfrak l)\). The simple-factor sum is the product decomposition of the quotient. Thus the following nilpotent data, together with the exceptional table below, determine the rank at every element.

For a nilpotent matrix let \(\lambda=(\lambda_1\geq\cdots\geq\lambda_k)\) denote its Jordan partition.

\begin{longtable}{@{}p{.23\linewidth}p{.38\linewidth}p{.29\linewidth}@{}}
\toprule\noalign{}
\begin{minipage}[b]{\linewidth}\raggedright
Type
\end{minipage} & \begin{minipage}[b]{\linewidth}\raggedleft
Exact nilpotent rank
\end{minipage} & \begin{minipage}[b]{\linewidth}\raggedright
Primary locator
\end{minipage} \\
\midrule\noalign{}
\endhead
\bottomrule\noalign{}
\endlastfoot
\(A_r=\mathfrak{sl}_{r+1}\) & \(\lambda_1-1\) & Richardson, Theorem 4.2.1, p.~233 \\
\(B_r=\mathfrak{so}_{2r+1}\), \(C_r=\mathfrak{sp}_{2r}\) & \(\lfloor\lambda_1/2\rfloor\) & Richardson, Theorem 4.3.3, p.~233 \\
\(D_r=\mathfrak{so}_{2r}\) & \(\lfloor\lambda_1/2\rfloor+\mathbf1_{\{k=2,\ \lambda_1\text{ odd}\}}\) & Richardson, Lemma 4.4.1 and Theorem 4.4.2(a)--(c), p.~234 \\
\end{longtable}

The partition conditions are in \textbf{§4.1, p.~232}: even parts have even multiplicity in orthogonal types; odd parts have even multiplicity in symplectic type. In \(D_r\), a very even partition labels two orbits, and both have the displayed rank. The compact \(D_r\) formula is precisely a restatement of the three cases in Theorem 4.4.2: with two parts, either both are odd or the partition is \((r,r)\) with \(r\) even. This is a citation of a known computation, not a novelty claim for the all-element matrix formulas separately proved in the root module.

\subsection{CC-RK02 --- Completed exceptional nilpotent rank table}\label{cc-rk02-completed-exceptional-nilpotent-rank-table}

The following table groups the \textbf{Bala--Carter labels in Richardson's Carter convention} by the rank of the quotient differential. Every unlisted \textbf{nonzero} nilpotent orbit has rank \(1\); the zero orbit always has rank \(0\). This expressly excludes zero from the source's abbreviated ``all other'' rows.

\begin{longtable}{@{}p{.08\linewidth}p{.08\linewidth}p{.74\linewidth}@{}}
\toprule\noalign{}
\begin{minipage}[b]{\linewidth}\raggedright
Type
\end{minipage} & \begin{minipage}[b]{\linewidth}\raggedleft
Rank
\end{minipage} & \begin{minipage}[b]{\linewidth}\raggedright
Orbits
\end{minipage} \\
\midrule\noalign{}
\endhead
\bottomrule\noalign{}
\endlastfoot
\(E_6\) & 6 & \(E_6\) \\
& 5 & \(E_6(a_1)\) \\
& 4 & \(D_5\) \\
& 3 & \(E_6(a_3)\), \(D_5(a_1)\), \(A_5\) \\
& 2 & \(A_4+A_1\), \(D_4\), \(A_4\) \\
\(E_7\) & 7 & \(E_7\) \\
& 6 & \(E_7(a_1)\) \\
& 5 & \(E_7(a_2)\) \\
& 4 & \(E_7(a_3)\), \(E_6\), \(D_6\) \\
& 3 & \(E_7(a_4)\), \(D_6(a_1)\), \(E_6(a_1)\), \(D_5+A_1\), \(D_5\) \\
& 2 & \(A_6\), \(E_7(a_5)\), \(E_6(a_3)\), \(D_6(a_2)\), \(D_5(a_1)+A_1\), \(A_5+A_1\), \((A_5)'\), \(D_5(a_1)\), \(D_4+A_1\), \((A_5)''\), \(D_4\) \\
\(E_8\) & 8 & \(E_8\) \\
& 7 & \(E_8(a_1)\) \\
& 6 & \(E_8(a_2)\) \\
& 5 & \(E_8(a_3)\), \(E_7\) \\
& 4 & \(E_8(a_4)\), \(E_8(b_4)\), \(E_7(a_1)\) \\
& 3 & \(E_8(a_5)\), \(E_8(b_5)\), \(D_7\), \(E_7(a_2)\), \(E_6+A_1\), \(E_6\) \\
& 2 & \(E_8(a_6)\), \(D_7(a_1)\), \(E_8(b_6)\), \(E_7(a_3)\), \(E_6(a_1)+A_1\), \(A_7\), \(D_7(a_2)\), \(D_6\), \(D_5+A_2\), \(E_6(a_1)\), \(E_7(a_4)\), \(D_6(a_1)\), \(D_5+A_1\), \(D_5\) \\
\(F_4\) & 4 & \(F_4\) \\
& 3 & \(F_4(a_1)\) \\
& 2 & \(F_4(a_2)\), \(B_3\), \(C_3\) \\
\(G_2\) & 2 & \(G_2\) \\
\end{longtable}

\textbf{Derivation and source corrections.} Richardson's \textbf{Appendix, pp.~239--240}, lists the nonzero exponents at each orbit; \textbf{Lemma 2.1, Proposition 2.2 and Corollary 2.3, p.~230}, identify their number with the differential rank. The table above counts those lists, with the following explicit changes. Both appendix pages were visually inspected, so these are defects in the printed source, not OCR repairs.

\begin{enumerate}
\def\labelenumi{\arabic{enumi}.}
\item
  The regular \(E_8\) exponent list ends with printed \(27\); the correct last exponent is \(29\). The ordinary exponent list is \(1,7,11,13,17,19,23,29\), also directly inspected in Dinar's \textbf{Table 2, arXiv v4 p.~11}. The rank is \(8\) either way, but the exponent itself matters in CC-RK03.
\item
  The \(E_8(a_1)\) row includes a spurious \(9\), giving eight printed entries. Richardson's \textbf{Lemma 5.1, p.~234}, says the subregular exponents are precisely the first \(r-1\) ordinary exponents. Thus this row is \(1,7,11,13,17,19,23\), and its rank is \(7\).
\item
  The class \(E_6\subset E_8\) is omitted from the rank-\(3\) rows. A regular nilpotent in the standard \(E_6\) Levi has exactly the common ordinary exponents \(1,7,11\), by Richardson's \textbf{Propositions 3.1 and 6.1, pp.~231, 234}, or the multiplicity formulation \textbf{§8.1, p.~238}. Its rank is therefore \(3\). This is a theorem-based insertion, not a conjectural correction.
\item
  The \(E_8(a_2)\) entry has the source's explicit question mark at \(19\). The complete deduction below removes that question mark using Johnson--Sommers.
\end{enumerate}

No independent reconstruction of every exceptional-orbit computation is claimed. The table is the application of those known primary results with the four changes explicitly justified here.

\subsection{\texorpdfstring{CC-RK03 --- The formerly unresolved \(E_8(a_2)\) entry has rank six}{CC-RK03 --- The formerly unresolved E\_8(a\_2) entry has rank six}}\label{cc-rk03-the-formerly-unresolved-e_8a_2-entry-has-rank-six}

\textbf{Proposition.} For \(e\in\mathfrak e_8\) in the orbit \(E_8(a_2)\),

\[
 \operatorname{rank}(d\chi_{\mathfrak e_8})_e=6,
 \qquad
 \operatorname{Exp}(\mathfrak e_8,e)=\{1,7,11,13,17,19\}.
 \tag{CC-RK.2}
\]

\textbf{Direct known theorem.} Dinar, arXiv:1905.01829v4, \textbf{Corollary 5.5, p.~14, proof p.~15}, states that the quotient differential has rank \(s\) at the distinguished nilpotent orbit \(Z_r(a_{r-s})\) in the list fixed at the end of \textbf{§4, p.~9}. That list includes \(E_8(a_2)\), hence \(r=8,s=6\). The opening of \textbf{§5, p.~12}, explicitly identifies this as the case left undecided by Richardson. \textbf{Table 1, p.~10}, and \textbf{Table 2, p.~11}, give the relevant orbit and ordinary exponent data. The statement, its short proof, and the preceding local argument in Lemma 5.1, Proposition 5.3 and Lemma 5.4, pp.~12--14, were inspected. This is not a claim to have independently reconstructed the whole paper or its separate integrability theorem.

\textbf{Inputs for a separate complete deduction.} Johnson--Sommers, arXiv:1706.04820v2, \textbf{§1.3, equations (1.1)--(1.2), pp.~2--3; Definition 1.2 and Theorem 1.3(1)--(2), p.~3; Figure 1(b), p.~4; §5.9, p.~15}. In their notation this orbit is the first-family case \(s=2\), induced from a Levi with two orthogonal simple roots. Their parameter \(s\) counts the simple roots of the inducing Levi semisimple part and is different from Dinar's parameter \(s=6\). In \(E_8\), the short-root representation \(V_\phi\) is the adjoint representation. Their theorem identifies the lowest-degree adjoint copy in the ideal of the orbit closure inside the nilpotent cone at degree \(m_{8-2+1}=m_7=23\), and an additional copy at degree \(m_8=29\). The statement of \textbf{Corollary 3.2, p.~9}, also gives normality of this closure, but normality is not needed separately in the following proof.

\textbf{Complete proof of the deduction.} Let \(A=\mathbf C[\mathcal N]\), and let \(I\subset A\) be the homogeneous ideal defining \(\overline{G e}\). The adjoint isotypic multiplicity in \(A\), degree by degree, is one in each degree

\[
 1,7,11,13,17,19,23,29
\]

and zero in all other degrees. This is Kostant's result recalled by Johnson--Sommers in §1.3: these copies are the derivatives of the eight fundamental invariants. Their Theorem 1.3(1) says that \(I\) contains no adjoint copy of degree less than \(23\), and that it contains the copy in degree \(23\). Part (2) says it contains the copy in degree \(29\). Each homogeneous graded piece is finite dimensional and is a semisimple rational \(G\)-module over \(\mathbf C\). Applying \(\operatorname{Hom}_G(\mathfrak e_8,-)\) to

\[
 0\longrightarrow I_j\longrightarrow A_j
     \longrightarrow (A/I)_j\longrightarrow0
\]

is therefore exact. Consequently the adjoint multiplicity in \(A/I=\mathbf C[\overline{G e}]\) is one in the six degrees \(1,7,11,13,17,19\) and zero in the other degrees. Richardson's \textbf{Proposition 2.5 and Corollary 2.6, p.~231}, identify the total multiplicity with the quotient-differential rank and the graded multiplicities with the orbit exponents. This gives (CC-RK.2). \(\square\)

Equivalently, the degree-\(19\) derivative copy of a degree-\(20\) fundamental invariant cannot vanish on the entire orbit. If it vanished at \(e\), equivariance would make it vanish on \(Ge\), hence on its closure, contradicting the minimum degree \(23\) of an adjoint copy in \(I\). This explains directly why the missing exponent is present.

\textbf{Additional source precision.} The arXiv v2 line in §5.9 describing the \emph{subregular} \(E_8(a_1)\) ideal prints \(V_{17}\); Theorem 1.3 with \(s=1\) instead forces \(V_{29}\). Our argument uses the numbered theorem and the separate, consistent \(E_8(a_2)\) entry \(V_{23}\), not that erroneous subregular line. No unsupported claim about whether a later journal rendering corrects this typo is made.

\subsection{CC-RK04 --- Subregular uniqueness and rank do not characterize ADE}\label{cc-rk04-subregular-uniqueness-and-rank-do-not-characterize-ade}

\textbf{Known Lie-algebra statement.} For every complex simple Lie algebra \(\mathfrak g\), of dimension \(N\) and rank \(r\), there is exactly one nilpotent orbit of dimension \(N-r-2\). For \(x\) in that orbit,

\[
 \operatorname{rank}(d\chi_{\mathfrak g})_x=r-1.
 \tag{CC-RK.3}
\]

The uniqueness statement is Richardson's \textbf{§5 opening paragraph, p.~234}, and the rank statement is his \textbf{Lemma 5.1 on the same page}, for all simple types. Johnson--Sommers independently restate the uniqueness in \textbf{§1.2, p.~2}. Richardson attributes the rank theorem to Slodowy; the exact inspected numbered locator used here is Richardson's Lemma 5.1, not an invented Slodowy theorem number. In type \(A_1\), the subregular orbit is the zero orbit, and the same formula gives rank zero.

\textbf{Unipotent consequence, with complete proof.} The original adjoint group \(G\) also has exactly one unipotent orbit of dimension \(N-r-2\). Embed \(G\) faithfully in \(\operatorname{GL}(\mathfrak g)\) by its adjoint representation. For a nilpotent \(X\) in its Lie algebra, the finite matrix sum \(\exp(X)\) is unipotent and lies in \(G\), since it is the usual complex Lie-group exponential. Conversely, if \(u\in G\) is unipotent, put \(M=u-I\) and consider the polynomial curve

\[
 u^t=\sum_{j\geq0}\binom tj M^j.
\]

The sum is finite. At every nonnegative integer \(t\) its value is a power of \(u\), hence lies in the Zariski-closed subgroup \(G\). The inverse image of \(G\) under this curve is closed in \(\mathbf A^1\) and contains infinitely many points, so it is all of \(\mathbf A^1\). Differentiating at \(t=0\) gives \(\log u\in\operatorname{Lie}(G)\), and this matrix is nilpotent. The finite exponential and logarithm polynomials are inverse on the nilpotent and unipotent varieties, respectively, by their formal power-series identities modulo a sufficiently high power of a nilpotent matrix. Both commute with conjugation. They therefore give a \(G\)-equivariant isomorphism of these varieties, identifying their individual orbits and orbit dimensions. The known unique nilpotent orbit transfers to the claimed unique unipotent orbit. \(\square\)

Thus the uniqueness speculation is refuted in all simple types, as is the proposed strict inequality below \(r-1\) for the Lie-algebra quotient in nonsimply-laced types. The next proposition transfers the rank statement specifically at unipotent elements and therefore settles the original subregular rank assertion for both maps. The distinction between general Lie and group ranks remains essential away from this locus, as in BR-12305--08.

\subsection{CC-RK05 --- All nilpotent ranks transfer to the unipotent group quotient}\label{cc-rk05-all-nilpotent-ranks-transfer-to-the-unipotent-group-quotient}

Let \(G\) be any connected complex semisimple algebraic group, with maximal torus \(T\), Lie algebras \(\mathfrak g,\mathfrak t\), and Weyl group \(W\). The group need not be simply connected. Denote the two quotient maps by

\[
 \chi_{\mathfrak g}:\mathfrak g\longrightarrow\mathfrak t/W,
 \qquad
 \chi_G:G\longrightarrow T/W.
\]

\textbf{Proposition.} If \(e\in\mathfrak g\) is nilpotent and \(u=\exp(e)\), then

\[
 \operatorname{rank}(d\chi_G)_u
       =\operatorname{rank}(d\chi_{\mathfrak g})_e.
 \tag{CC-RK.4}
\]

Consequently the whole nilpotent table in CC-RK01--03 also gives the ranks at the corresponding unipotent elements in every central form of \(G\). In particular, the unique subregular unipotent orbit has quotient-differential rank \(r-1\) in \textbf{every} simple type. The ranks in \(B_r,C_r,F_4,G_2\) are \(r-1,r-1,3,1\), respectively, for both original maps.

\textbf{Primary structural inputs.} The Lie-algebra restriction isomorphism is recalled by Richardson in \textbf{§1, p.~229}. The group restriction isomorphism
\(\mathbf C[G]^G\simeq\mathbf C[T]^W\), without a simply-connected hypothesis, is Steinberg's \textbf{Corollary 6.4, p.~63}, proved from \textbf{Lemmas 6.2--6.3, p.~62}. These exact printed pages and statements were inspected. The following is a complete local analytic deduction; no global equivalence of the two ambient quotient maps is asserted.

\textbf{Proof.} The torus exponential has derivative the identity at \(0\). Choose a sufficiently small \(W\)-invariant open neighborhood \(V\subset\mathfrak t\) of \(0\) on which it is a biholomorphism onto its image \(V_T\subset T\), a \(W\)-invariant neighborhood of \(1\). Since the exponential is \(W\)-equivariant, it induces a biholomorphism

\[
 \psi:V/W\xrightarrow{\sim}V_T/W.
 \tag{CC-RK.5}
\]

Here the two finite-group quotients are open analytic neighborhoods in the analytifications of \(\mathfrak t/W\) and \(T/W\). Indeed a \(W\)-invariant open set is saturated for the finite orbit quotient, and its holomorphic functions on the quotient are the invariant holomorphic functions; applying this to the exponential and its equivariant inverse proves the assertion.

The open set \(U=\chi_{\mathfrak g}^{-1}(V/W)\) contains the nilpotent cone. A regular semisimple \(x\in U\) is conjugate to a \(t\in V\): its Cartan conjugacy class is a \(W\)-orbit, and its quotient value lies in \(V/W\). Equivariance of the exponential and the two restriction isomorphisms give

\[
 \chi_G(\exp x)=\psi(\chi_{\mathfrak g}(x))
 \tag{CC-RK.6}
\]

for all such \(x\). Regular semisimple elements form a Zariski-open dense subset of \(\mathfrak g\), hence are analytically dense in \(U\). Both sides of (CC-RK.6) are holomorphic on \(U\); after embedding the affine target \(T/W\) into an affine space, equality of each coordinate extends by continuity to all \(U\).

For completeness, the derivative of the exponential in a faithful matrix representation satisfies

\[
 (L_{u^{-1}})_*(d\exp)_e(Y)
   =\int_0^1 \exp(-t\,\operatorname{ad}e)Y\,dt
   =\left(\sum_{j\geq0}
       \frac{(-1)^j(\operatorname{ad}e)^j}{(j+1)!}\right)Y.
 \tag{CC-RK.7}
\]

The integral identity follows by differentiating the exponential power series term by term, or by integrating
\(\exp((1-t)e)Y\exp(te)\) before left translation. Since \(\operatorname{ad}e\) is nilpotent, the displayed sum is finite and has constant term \(1\). It is therefore an invertible endomorphism of \(\mathfrak g\), so \((d\exp)_e\) is invertible. Differentiating (CC-RK.6) at \(e\) gives

\[
 (d\chi_G)_u\circ(d\exp)_e
       =(d\psi)_0\circ(d\chi_{\mathfrak g})_e.
\]

Both \((d\exp)_e\) and \((d\psi)_0\) are isomorphisms. Taking ranks proves (CC-RK.4). Analytic and algebraic tangent differentials of these algebraic quotient maps agree, so this is the rank asked for in the original letter. \(\square\)

\subsection{Exact bibliography and available primary files}\label{cc-rk-bibliography}

Grothendieck, A. (1969, October 22). \emph{Letter to B. Kostant} {[}Manuscript, cote 123, physical PDF p.~23, footer 22{]}. Université de Montpellier. \url{https://grothendieck.umontpellier.fr/123.pdf}

Richardson, R. W. (1987). Derivatives of invariant polynomials on a semisimple Lie algebra. In M. Cowling, C. Meaney, \& W. Moran (Eds.), \emph{Miniconference on harmonic analysis and operator algebras} (Proceedings of the Centre for Mathematical Analysis, Australian National University, Vol. 15, pp.~228--241). Australian National University. \url{https://maths.anu.edu.au/files/CMAProcVol15-Richardson_1.pdf}

Johnson, B., \& Sommers, E. (2020). Equations for some nilpotent varieties. \emph{International Mathematics Research Notices, 2020}(14), 4433--4464. \url{https://doi.org/10.1093/imrn/rny132} . Inspected author version: arXiv:1706.04820v2, 5 February 2018, \url{https://arxiv.org/pdf/1706.04820v2} . All page locators in this module refer to that version.

Dinar, Y. (2020). On integrability of transverse Lie--Poisson structures at nilpotent elements. \emph{Journal of Geometry and Physics, 155}, Article 103690. \url{https://doi.org/10.1016/j.geomphys.2020.103690} . Inspected author version: arXiv:1905.01829v4, 7 April 2020, \url{https://arxiv.org/pdf/1905.01829v4} . All Dinar page locators in this module refer to that version.

Steinberg, R. (1965). Regular elements of semi-simple algebraic groups. \emph{Publications Mathématiques de l'IHÉS, 25}, 49--80. \url{https://doi.org/10.1007/BF02684397} . The general restriction isomorphism used here is Corollary 6.4, p.~63, with Lemmas 6.2--6.3, p.~62.

Local primary snapshots: \nolinkurl{research/cc_sources/Richardson1987.pdf}, its extracted text, and visually inspected table images \nolinkurl{Richardson1987_table-12.png}/\nolinkurl{Richardson1987_table-13.png}; \nolinkurl{research/cc_sources/JohnsonSommers2018.pdf} and its extracted text; \nolinkurl{research/cc_sources/Dinar2019v4.pdf} and its extracted text. The source and check ledgers accompanying this module record the exact retrieval references, inspection scopes, and actual proof-reading events.

% END consequences_criteria_adjoint_rank.tex

% BEGIN root_henselian.tex
\section{Cote 66, subject 21: the exact affine henselian comparison}

\subsection{Original statement and coefficient scope}
The official cote 66 PDF, physical page 4, archive footer 3, subject 21,
asks when
\[
 H^*(X,F)\longrightarrow H^*(X_0,F_0)
\]
is an isomorphism for a closed subscheme $X_0\subset X$, and asks
\emph{``Suffit-il que ce soit vrai en dim 0?''}
The margin identifies \emph{cohomologie \'etale}. The specified cases
are $X=\operatorname{Spec}A$, $X_0=\operatorname{Spec}(A/I)$ with
$A$ separated and complete for the $I$-adic topology, and \'etale
localization along a closed subscheme of an affine scheme.
The sheet does not explicitly write the quantifier on $F$ or restrict
it to torsion. We state those choices: the all-torsion-sheaf affine
question has the following exact answer, while a test on one fixed
sheaf does not suffice.

\subsection{BR-66-01. Degree zero for all torsion coefficients is an exact affine test}
Let $A$ be a commutative ring and $I\subset A$ any ideal, and set
$X=\operatorname{Spec}A$, $Z=\operatorname{Spec}(A/I)$. The following
conditions are equivalent:
\begin{enumerate}
\item $(A,I)$ is a henselian pair.
\item For every torsion abelian sheaf $F$ on $X_{\text{ét}}$,
      the restriction $H^0(X,F)\to H^0(Z,F|_Z)$ is an isomorphism.
\item For every such $F$ and every $q\geq0$,
      $H^q(X,F)\to H^q(Z,F|_Z)$ is an isomorphism.
\item For every finite morphism $f:Y\to X$, restriction gives a
      bijection from the clopen subsets of $Y$ to those of
      $Y\times_X Z$.
\end{enumerate}
In particular, the degree-zero condition can be tested just on the
sheaves $f_*\underline{\mathbf F}_2$ for all finite $f$.

\paragraph{Exact known inputs.}
Gabber's affine analogue of proper base change gives (1)$\Rightarrow$(3)
for arbitrary torsion sheaves, with no noetherian assumption and no
invertibility assumption on their torsion orders. The exact statement
used is The Stacks Project, Theorem 59.82.7, Tag \nolinkurl{09ZI}; the
original article is Gabber (1994), pp.~325--335. Its complete published
proof is not reproduced here. The equivalence (1)$\Leftrightarrow$(4)
is the idempotent criterion in The Stacks Project, Lemma 15.11.6(1),(3),
Tag \nolinkurl{09XI}. We prove the short implication from the stated
sheaf test to that criterion, to settle the quantifier precisely.

\paragraph{Complete proof of the additional implication and equivalence.}
The implication (3)$\Rightarrow$(2) is immediate. Assume (2), and
fix a finite morphism $f:Y\to X$. The sheaf
$F=f_*\underline{\mathbf F}_2$ is annihilated by $2$, so (2) applies.
Put $Y_Z=Y\times_X Z$ and $f_Z:Y_Z\to Z$. Proper base change in
degree zero for the finite morphism gives
\[
 (f_*\underline{\mathbf F}_2)|_Z
       \simeq (f_Z)_*\underline{\mathbf F}_2.
\]
An exact inspected locator for this known input is The Stacks
Project, Theorem 59.91.11, Tag \nolinkurl{095T}; only its degree-zero
finite-morphism case is used. By the definition of direct image,
condition (2) therefore gives
\[
 \Gamma(Y,\underline{\mathbf F}_2)
       \xrightarrow{\sim}\Gamma(Y_Z,\underline{\mathbf F}_2).
 \tag{HEN.1}
\]
A section of the constant two-element \'etale sheaf is the same as
a decomposition of the scheme into two open and closed subsets:
it is a section of the disjoint union of two copies of the scheme,
and the inverse images of its two components are the two subsets.
Thus (HEN.1) is precisely condition (4). Since $Y$ is affine,
these decompositions are also precisely the idempotents of its ring,
and restriction is reduction modulo $I$. The cited idempotent
criterion implies (1). The same proof works if (2) is assumed only
for this displayed family of finite direct images. Combining it with
Gabber's theorem proves all stated equivalences. \(\square\)

\subsection{BR-66-02. Both affine cases named on the sheet are affirmative}
If $A$ is $I$-adically separated and complete, the desired comparison
holds for every torsion abelian \'etale sheaf and every degree.
The same holds for the henselization $(A^h,IA^h)$ of any affine pair,
with closed subscheme $\operatorname{Spec}(A/I)$.
These are direct applications of known results: The Stacks Project,
Lemma 15.11.4, Tag \nolinkurl{0ALJ}, says an $I$-adically complete pair is
henselian; Lemma 15.12.1, Tag \nolinkurl{0A02}, constructs the henselization
as the filtered colimit of \'etale neighborhoods inducing an isomorphism
on the closed subscheme, proves it henselian, and identifies
$A^h/IA^h=A/I$. Apply BR-66-01. This identifies the affine \'etale
localization intended by the second case without asserting that an
arbitrary unrelated pair is henselian.

\subsection{BR-66-03. Testing only one coefficient sheaf in degree zero fails}
Even for smooth affine schemes over $\mathbf C$, comparison in degree
zero for one fixed finite constant sheaf does not imply comparison
in degree one for that same sheaf.

\paragraph{Complete counterexample.}
Take $X=\operatorname{Spec}\mathbf C[t,t^{-1}]$,
$Z=\{t=1\}=\operatorname{Spec}\mathbf C$, and
$F=\underline{\mathbf Z/2}$. Both schemes are connected, so
restriction on $H^0$ is the identity on $\mathbf Z/2$.
The morphism
\[
 \operatorname{Spec}\mathbf C[u,u^{-1}]\longrightarrow X,
 \qquad t=u^2,
\]
is a torsor for $\{1,-1\}\simeq\mathbf Z/2$. It is finite \'etale:
its ring is free of rank two over $\mathbf C[t,t^{-1}]$, and the
derivative $2u$ is a unit. The sign action is free and acts transitively
on every geometric fibre. It is not a trivial torsor, because a section
would give a unit $a(t)\in\mathbf C[t,t^{-1}]$ with $a(t)^2=t$.
Every Laurent-polynomial unit is $ct^m$: the least and greatest
exponents of a product are sums of the corresponding exponents,
so a product equal to $1$ forces the difference between the greatest
and least exponent of each factor to be zero. Then $a(t)^2=t$
would require $2m=1$, which is impossible.

The classification of abelian-sheaf torsors by $H^1$ (The Stacks
Project, Lemma 21.4.3, Tag \nolinkurl{03AJ}, with the section criterion
in Lemma 21.4.2, Tag \nolinkurl{03AI}) therefore gives
a nonzero class in $H^1(X,F)$. Its restriction to $Z$ is trivial:
the fibre consists of the two points $u=1,-1$ and has a section.
Consequently $H^1(X,F)\to H^1(Z,F|_Z)$ is not injective.
This proves the assertion. \(\square\)

The result answers the original affine torsion formulation, including
its degree-zero test with the universal coefficient quantifier.
It does not assert a general comparison theorem for arbitrary
quasicompact schemes or for nontorsion coefficients. The unseen
SA22 run's specific nonseparated counterexample and its residual
\texttt{U:henselian-comparison} are not identified with this elementary
fixed-sheaf counterexample.

\subsection*{Exact references and source inspection}
Grothendieck, A. (n.d.). \emph{List of subjects}, cote 66,
physical PDF p.~4, archive footer 3, subject 21.
\url{https://grothendieck.umontpellier.fr/66.pdf}.
The complete sheet and an enlarged subject-21 crop were visually read
in this run; ``66/004'' in the intake denotes the physical PDF page.

Gabber, O. (1994). Affine analog of the proper base change theorem.
\emph{Israel Journal of Mathematics, 87}, 325--335.
\url{https://doi.org/10.1007/BF02773001}.
The publisher abstract and bibliographic record were inspected;
the accessible complete theorem and proof used are the following
Stacks Project version, so no unretrieved theorem number from the
original article is asserted as newly checked.

The Stacks Project Authors. (2026). \emph{The Stacks Project}.
Theorem 59.82.7 (Gabber), \url{https://stacks.math.columbia.edu/tag/09ZI};
Lemma 15.11.6, \url{https://stacks.math.columbia.edu/tag/09XI};
Theorem 59.91.11, \url{https://stacks.math.columbia.edu/tag/095T};
Lemma 15.11.4, \url{https://stacks.math.columbia.edu/tag/0ALJ};
Lemma 15.12.1, \url{https://stacks.math.columbia.edu/tag/0A02}.
Lemmas 21.4.2--21.4.3 in \url{https://stacks.math.columbia.edu/tag/03AG}.
All displayed statements were directly inspected on 21 November 2026.

% END root_henselian.tex

% BEGIN root_comparison.tex
\section{The comparison-cohomology problems recovered in cote 18}

\subsection{Original question and precise modern scope}
The original source is Grothendieck's Montpellier archive, cote 18,
physical PDF page 43, manuscript and printed footer 42, headed
``Pbs.''  Its four numbered headings read ``Betti, De-Rham,''
``Betti, $\ell$-adique,'' ``DeRham, $p$-adique,'' and
``D\'efinition de cohomologie $p$-adique en car $p>0$?''
The first heading also branches into comparison problems involving
different embeddings and additional structures.  We separate those
branches from the four precise cohomological assertions below.
The original sheet was visually inspected in this run.

These are answers by known comparison theorems, with their hypotheses
retained.  They are not findings newly proved from RH.  The IDs
\nolinkurl{KN-COMP01--04} identify source-checked answers, not claims of
new comparison theorems.  The unrelated uppercase $K$ on physical
page 45/footer 44 is still not defined by the recovered adjacent
sheets; no meaning is assigned to it here.

\paragraph{KN-COMP01: Betti--de Rham comparison.}
Let $X$ be a smooth separated finite-type scheme over $\mathbf C$.
There is a canonical comparison
\[
 \mathbb H^i(X,\Omega^\bullet_{X/\mathbf C})
       \simeq H^i(X^{\rm an},\mathbf C).
 \tag{COMP.1}
\]
This is Grothendieck (1966), Theorem $1'$,
\emph{Publications Math\'ematiques de l'IH\'ES} 29, p.~96;
Theorem 1, p.~95, gives the smooth affine form.  The theorem is
cited, not reproved.  If a variety is defined over a subfield and
an embedding into $\mathbf C$ is specified, (COMP.1) is applied to
that complex base change.  This answers the cohomological comparison
for each specified embedding.  A preferred rational identification
between the Betti fibre functors of two different embeddings is
additional data; the complex comparison alone does not specify it.

\paragraph{KN-COMP02: Betti--\texorpdfstring{$\ell$}{ell}-adic comparison.}
For a separated finite-type complex scheme $X$ and every finite
constant abelian coefficient group $A$, the analytification comparison gives
\[
 H^i_{\rm et}(X,A)\simeq H^i(X^{\rm an},A).
 \tag{COMP.2}
\]
The precise general input is Artin, SGA~4, Expos\'e XVI,
Theorem 4.1; in the inspected re-edition its statement occupies
pp.~13--14.  For $f:X\to S$, it requires $f$ proper or the
coefficient sheaf constructible, and identifies the corresponding
higher direct images.  Constant finite $A$ is constructible, so the
absolute comparison above is included.  The familiar $\mathbf Z_\ell$
and $\mathbf Q_\ell$ comparisons use the compatible finite-level maps.
For example, for smooth proper $X$ the finite groups at each level
satisfy the Mittag--Leffler condition: for fixed level the descending
images stabilize.  Taking the inverse limit in (COMP.2), then
inverting $\ell$, gives the stated adic comparisons.  The comparison
is over the chosen complex base change; it does not supply an
unspecified rational period-torsor point.

\paragraph{KN-COMP03: crystalline cohomology in characteristic
\texorpdfstring{$p$}{p}.}
For an $\mathbf F_p$-scheme $S$, the crystalline site relative to
$\mathbf Z_p$ uses nilpotent divided-power thickenings of opens of
$S$, compatible with the canonical divided powers on $(p)\subset\mathbf Z_p$; its structure sheaf assigns the functions on the thickening.
Its derived global sections define $R\Gamma_{\rm crys}(S/\mathbf Z_p)$.
A crystal has the prescribed base-change isomorphisms along these
thickenings.  This is the directly inspected definition in
Bhatt--de Jong (2011), \S1, pp.~1--2.
Their Theorem 3.6, p.~8, computes crystalline cohomology by the
de Rham complex on a completed divided-power envelope.  Their
Corollary 3.8, p.~9, gives, for $X$ proper smooth over $\mathbf Z_p$
and $S=X\otimes\mathbf F_p$,
\[
 R\Gamma_{\rm crys}(S/\mathbf Z_p)
        \simeq R\Gamma(X,\Omega^\bullet_{X/\mathbf Z_p}).
 \tag{COMP.3}
\]
The left side is intrinsic to $S$, so the comparison also explains
the lift-independent Frobenius on this de Rham cohomology.
These are known results, with no assumption that the special fibre
has a chosen lift in the definition itself.

\paragraph{KN-COMP04: de Rham--\texorpdfstring{$p$}{p}-adic and
integral mixed-characteristic comparisons.}
Let $K/\mathbf Q_p$ be complete discretely valued with perfect
residue field, $C$ a complete algebraically closed extension, and
$X_0/K$ proper smooth.  Bhatt--Morrow--Scholze (2018), Theorem 1.7,
author version p.~5, gives the canonical comparison
\[
 H^i_{\rm dR}(X_0/K)\otimes_K B_{\rm dR}
    \simeq H^i_{\rm et}(X_0\widehat\otimes_K C,\mathbf Z_p)
                    \otimes_{\mathbf Z_p}B_{\rm dR}.
 \tag{COMP.4}
\]
For a proper smooth formal scheme $\mathfrak X/\mathcal O_C$, their
Theorem 1.8, pp.~5--6, supplies the perfect complex
$A=R\Gamma_{A_{\rm inf}}(\mathfrak X)$, where
$A_{\rm inf}=W(\mathcal O_C^\flat)$, with Frobenius and comparisons
\[
\begin{aligned}
 A\otimes^{L}_{A_{\rm inf}}W(\kappa_C)
   &\simeq R\Gamma_{\rm crys}(\mathfrak X_{\kappa_C}/W(\kappa_C)),\\
 A\otimes^{L}_{A_{\rm inf}}\mathcal O_C
   &\simeq R\Gamma_{\rm dR}(\mathfrak X),\\
 A[1/\mu]&\simeq
 R\Gamma_{\rm et}(\mathfrak X_C,\mathbf Z_p)
       \otimes^L_{\mathbf Z_p}A_{\rm inf}[1/\mu].
\end{aligned}
 \tag{COMP.5}
\]
Here $\kappa_C$ is the residue field and
$\mu=[\epsilon]-1$ comes from compatible $p$-power roots of unity.
These derived comparisons retain torsion information; replacing them
by unqualified integral cohomology-group isomorphisms would lose
their hypotheses.  The cited theorem also treats the crystalline
$A_{\rm crys}$ specialization.

\subsection{Consequences for the requested bridge}
The four numbered headings therefore have precise affirmative
cohomological answers in the scopes displayed above.  They explain
how Betti, de Rham, crystalline and \emph{generic-fibre} adic
realizations can belong to one arithmetic-geometric comparison
system.  The period rings and the divided-power or integral
comparison complexes are part of that system's mathematical data.
The adopted RH and Dirichlet-GRH theorems determine the zero locations
of the abelian Artin--Tate factors already isolated in BR-AT-01.
The comparison results identify realizations; they do not by
themselves construct the algebraic cross-correspondences still
required in CC-XC01--04, or choose a rational point on the period
torsors in the remaining cote 18 questions.

\paragraph{Primary references and inspected locators.}
Grothendieck, A. (1966). On the de Rham cohomology of algebraic
varieties. \emph{Publications Math\'ematiques de l'IH\'ES, 29},
95--103. \url{https://doi.org/10.1007/BF02684807}.
Theorems 1 and $1'$, pp.~95--96; primary scan:
\url{https://www.numdam.org/article/PMIHES_1966__29__95_0.pdf}.

Artin, M. (1973). Th\'eor\`eme de changement de base par un
morphisme lisse, et applications. In M. Artin, A. Grothendieck,
\& J.-L. Verdier (Eds.), \emph{Th\'eorie des topos et cohomologie
\'etale des sch\'emas (SGA 4)}, Tome 3, Expos\'e XVI,
Theorem 4.1. Lecture Notes in Mathematics 305. Springer.
Inspected re-edition, version 71766d9 (30 July 2024), pp.~13--14:
\url{https://www.normalesup.org/~forgogozo/SGA4/16/16.pdf}.
The version-specific pages are not asserted as original pagination.

Bhatt, B., \& de Jong, A. J. (2011).
\emph{Crystalline cohomology and de Rham cohomology}
[Preprint, arXiv:1110.5001]. \S1 pp.~1--2; Theorem 3.6 p.~8;
Corollary 3.8 p.~9.
\url{https://www.math.columbia.edu/~dejong/papers/crystalline-comparison.pdf}.

Bhatt, B., Morrow, M., \& Scholze, P. (2018).
Integral $p$-adic Hodge theory.
\emph{Publications Math\'ematiques de l'IH\'ES, 128}, 219--397.
\url{https://doi.org/10.1007/s10240-019-00102-z}.
Inspected author version dated 8 January 2019, Theorems 1.7--1.8,
pp.~5--6:
\url{https://www.math.uni-bonn.de/people/scholze/integralpadicHodge.pdf}.

% END root_comparison.tex

% BEGIN root_adelic_genericity.tex
\section{Cote 18, physical 49/footer 48: simultaneous local genericity}

\subsection{The original assertion and its two hypotheses}
The bottom of the original sheet, cote 18, physical PDF page 49,
manuscript/footer 48, considers a connected affine algebraic group
\(G/\mathbf Q\), a compact open subgroup
\(U\subset G(\mathbf A_f)\), and a coset \(P=Ua\).
It asserts the existence in \(P\) of an element whose component at
every finite prime is generic over \(\mathbf Q\).  The preceding
lines apply this to tuples of comparison matrices; physical page 51,
footer 50, continues the assertion in terms of the local points of
the variety of tuples.  Here \(\mathbf A_f\) denotes the finite
adeles, written with the archive's hatted rational notation there.

The group is explicitly connected on this sheet.  The proposed
motivic application separately assumes a uniform form of the Tate
conjecture giving a finite-index adelic Galois image.  The local
genericity assertion below is unconditional.  Its application to a
specified Galois representation retains the compact-open-image
hypothesis; the adopted RH and Dirichlet GRH are not used to supply it.

\subsection{BR-18-AD01: the adelic assertion, including finite tuples}
Let \(G\) be a connected affine algebraic group over \(\mathbf Q\),
let \(r\geq1\), and let
\(\mathcal O\subset G(\mathbf A_f)^r\) be a nonempty open set.
There exists \(g=(g_1,\ldots,g_r)\in\mathcal O\) such that,
for every rational prime \(\ell\), evaluation is injective:
\[
 \mathbf Q[G^r]\longrightarrow\mathbf Q_\ell,
 \qquad f\longmapsto f(g_{1,\ell},\ldots,g_{r,\ell}).
 \tag{18.AD1}
\]
The set of such tuples is dense in \(G(\mathbf A_f)^r\).
In particular it meets every coset of a compact open subgroup.
Genericity in (18.AD1) is relative to \(\mathbf Q\), exactly as in
the source; it does not mean that a \(\mathbf Q_\ell\)-valued point
is generic after allowing arbitrary \(\mathbf Q_\ell\)-coefficients.

\paragraph{Known local input.}
Poonen (2017), Proposition 3.5.75 and Remark 3.5.76, p.~87,
give the following stronger local form of Zariski density: in an
irreducible variety over a local field, every open neighborhood of
a smooth rational point contains a point outside any prescribed
countable union of proper closed subvarieties.  This is the known
local result; we give only the additional adelic deduction.
For the group hypotheses, the same source, Proposition 5.2.4,
pp.~121--122, makes the identity component geometrically
irreducible, and Corollary 5.2.18, p.~125, gives smoothness in
characteristic zero.

\paragraph{Complete proof of the deduction.}
Put \(X=G^r\).  The cited group facts imply that \(X\) and every
\(X_{\mathbf Q_\ell}\) are smooth and geometrically integral.
The coordinate ring \(A=\mathbf Q[X]\) is countable: it is a
finitely generated algebra over the countable field \(\mathbf Q\).
For each nonzero \(f\in A\), its image in
\(A\otimes_{\mathbf Q}\mathbf Q_\ell\) remains nonzero, by
faithful flatness of field extension.  Its zero set
\(Z_{f,\ell}\) is consequently a proper closed subvariety of
\(X_{\mathbf Q_\ell}\).  The local result applies to this
countable family.  Thus every nonempty analytic open subset of
\(X(\mathbf Q_\ell)\) contains a point \(x_\ell\) with
\(f(x_\ell)\ne0\) for all \(0\ne f\in A\).

Choose a point of \(\mathcal O\) and then a basic restricted-product
neighborhood of it contained in \(\mathcal O\).  Explicitly, after
choosing an affine group model of \(G\) outside finitely many
primes, this neighborhood has the form
\[
       \prod_{\ell\in S}V_\ell
          \times\prod_{\ell\notin S}\mathcal G(\mathbf Z_\ell)^r,
 \tag{18.AD2}
\]
where \(S\) is finite and every \(V_\ell\) is a nonempty
analytic open subset of \(X(\mathbf Q_\ell)\).
Each \(\mathcal G(\mathbf Z_\ell)^r\) appearing here is also
nonempty and analytically open: it contains the identity, and in
an affine embedding its integrality conditions are open conditions
in \(\mathbf Q_\ell\).  These are precisely the basic sets in
the restricted-product topology.

For each \(\ell\), use the preceding local conclusion to choose
\(x_\ell\) in the corresponding factor of (18.AD2).  The tuple
\((x_\ell)_\ell\) is an adelic point because it is integral
outside \(S\); it lies in \(\mathcal O\).  Every evaluation map
in (18.AD1) has zero kernel by construction.  Since
\(\mathcal O\) was arbitrary, the set of such points is dense.
Applying this with \(r=1\) and \(\mathcal O=Ua\) proves the
original assertion. \(\square\)

\paragraph{The precise Galois-image application.}
Suppose a group homomorphism
\(\rho:\Pi\to G(\mathbf A_f)\) has compact open image \(U\).
For prescribed \(a_1,\ldots,a_r\in G(\mathbf A_f)\), there
exist \(\sigma_1,\ldots,\sigma_r\in\Pi\) such that the tuple
\[
       (\rho(\sigma_1)a_1,\ldots,\rho(\sigma_r)a_r)
 \tag{18.AD3}
\]
satisfies (18.AD1) at every \(\ell\).  Indeed,
\(Ua_1\times\cdots\times Ua_r\) is nonempty and open; apply
the theorem and lift each of its \(U\)-coordinates through
the surjection \(\rho:\Pi\twoheadrightarrow U\).
This proves the finite-tuple step of the archive argument under
its stated image hypothesis.

It matters that this is an existence statement for arbitrary
elements of a compact open image.  Even a dense distinguished
subset need not contain a generic element.  For example,
\(\mathbf Z\subset\widehat{\mathbf Z}\subset\mathbf A_f\)
is dense in the compact open additive subgroup
\(\widehat{\mathbf Z}\): any finite list of residue conditions
has an integer solution by the Chinese remainder theorem.
Nevertheless each integer \(n\) is killed under evaluation by
the nonzero polynomial \(t-n\in\mathbf Q[t]\) at every prime.
Thus density of a selected collection of arithmetic elements
does not by itself provide the element of (18.AD1).

\paragraph{Bibliographic locator.}
Poonen, B. (2017). \emph{Rational points on varieties}.
Graduate Studies in Mathematics, 186. American Mathematical Society.
Proposition 3.5.75 and Remark 3.5.76, p.~87;
Proposition 5.2.4, pp.~121--122; Corollary 5.2.18, p.~125.
Author-hosted text:
\url{https://math.mit.edu/~poonen/papers/Qpoints.pdf}.


% END root_adelic_genericity.tex

% BEGIN root_deligne.tex
\section{Deligne's explicitly stated problem and the geometric weight boundary}
\label{sec:deligne-problem}

\subsection{D-WII-1.8.15: good degenerations of mixed Hodge structures}

There is an actual question next to the cited local-monodromy theorem.
Deligne (1980), \emph{La conjecture de Weil II}, 1.8.14--1.8.15,
printed pp.~178--179 (PDF pp.~42--43), asks to
``D\'egager une classe de `bonnes' variations de structures de Hodge mixtes''
on the punctured disc. In precise mathematical terms, for a good
unipotent variation $(V,W,F)$ with $N=\log T$, he asks for a finite
increasing filtration $M$ such that
\[
 NM_a\subseteq M_{a-2},\qquad
 N^b:\operatorname{gr}^{M}_{a+b}\operatorname{gr}^{W}_a V
       \xrightarrow{\sim}
       \operatorname{gr}^{M}_{a-b}\operatorname{gr}^{W}_a V
 \quad(b\geq0).
\]
He also asks for an asymptotic nilpotent orbit and for the filtration
$M$ to supply a mixed Hodge structure for the nilpotent orbit.
This is a faithful formal restatement of the displayed problem, not
an attribution of a question about the Riemann zeta function.

\paragraph{Answer: known geometric theorem, independent of the adopted RH.}
Admissible graded-polarizable variations provide the requested class.
For the precise definition use Kashiwara (1986), \S1.8,
conditions (1.8.1)--(1.8.4), pp.~993--994, and \S1.9: quasi-unipotence,
relative monodromy, and locally free extension of the Hodge and weight
graded bundles. Proposition 4.5.1 and Proposition 5.2.1,
pp.~1008--1010, identify the limiting data and the mixed Hodge
structure for $M$. His Theorem 4.4.1 and Theorem 4.5.2,
pp.~1007--1008, establish the compatibility of relative filtrations
and the normal-crossing extension criterion. These are theorems about
the defined class; the assertion is not that every variation satisfying
only Deligne's initial three axioms is automatically admissible.

For an explicit asymptotic estimate, Pearlstein (2000),
arXiv:math/0002030v3, Theorem 5.1(i)--(iii), p.~30,
gives a limiting filtration $F_\infty$ and the horizontal orbit
$F_{\rm nilp}(z)=\exp(zN)F_\infty$, with
\[
 d\bigl(F(z),F_{\rm nilp}(z)\bigr)
 \leq K(\operatorname{Im}z)^\beta
                  e^{-2\pi\operatorname{Im}z}
 \quad(\operatorname{Im}z\gg0)
\]
in the mixed Hodge metric. This supplies the asymptotic part in a
stronger quantitative form. The relative-monodromy mixed Hodge
structure and this estimate together answer the stated existence
problem. The question is settled by the cited Hodge-theoretic work;
no new proof or historical priority claim is made here.

\subsection{D-WII-1.9.5: the arithmetic base is a different gate}

At the end of 1.9.5, printed p.~181 (PDF p.~45), Deligne contrasts
normal-crossing divisors in a regular scheme of finite type over
$\mathbb Z$ with the preceding finite-field theory, already mentioning
$\operatorname{Spec}\mathbb Z[1/\ell]$. This observation must be
separated from the solved complex-analytic problem 1.8.15.
For the proper smooth local-field case a precise modern statement is
Scholze (2012), \emph{Perfectoid spaces}, Conjecture 9.3, p.~48
of arXiv:1111.4914:
for $H^i(X_{\overline K},\mathbb Q_\ell)$ the $j$th monodromy-graded
piece has Frobenius weight $i+j$. His Theorem 9.6, p.~49, proves
this in its stated set-theoretic complete-intersection situation
inside a projective smooth toric variety. We retain all of those
hypotheses; it is not a theorem for arbitrary proper smooth $X/K$.

BR-WM-01 identifies a concrete missing datum: a Frobenius
characteristic polynomial and the covariance equation do not determine
the monodromy filtration. Hence the adopted global zero-location
statements do not supply this datum through a reconstruction from
Euler polynomials. This is a proved failure of that proposed route,
not a metamathematical independence claim. The remaining gate for a
specified degeneration is the purity of each actual monodromy-graded
piece, with its actual $N$ and geometric realization.

\subsection{The universal stable curve: purity of a complex and weights of a fibre}

The RESIDUE intake, \S4(1), mixes two valid constructions with an
invalid identification. For a proper map from a smooth scheme over
a finite field, the constant complex is pure, and proper pushforward
preserves purity in the sense of mixed complexes. The exact source is
Beilinson, Bernstein, and Deligne (1982), Stabilit\'es 5.1.14 and its dual
form, Corollary 5.4.4, Theorem 5.4.5, and Corollary 5.4.6,
pp.~128--129 and 142. The splitting statements there hold after
base change to the algebraic closure: arithmetic Frobenius
semisimplicity is not silently added.

The universal stable-curve application uses Sun's stack decomposition
theorem, with its finite-diagonal and affine-stabilizer hypotheses,
as proved in CC-ST01 below. A splitting does not in general descend
merely from its existence on scheme charts. It does not say that every ordinary
cohomology group of a singular geometric fibre is pure of its
cohomological degree. BR-NODE-01 computes the actual group: its
weight-zero graph subspace and weight-one normalization part need
not have a single weight. Nearby smooth cohomology is a different
object. In the one-node genus-one smoothing its monodromy filtration
has graded pieces of weights zero and two, while the nodal special
fibre has only a one-dimensional weight-zero $H^1$.

For that last nearby-cycle statement the exact classical input is
SGA~7~II, Expos\'e XV, Th\'eor\`eme 3.4(ii)--(iii), printed
pp.~195--196: specialization, vanishing cycles, and the
Picard--Lefschetz inertia formula. The determinant
$\det H^1=\mathbb Q_\ell(-1)$ identifies the quotient twist once
the trivial specialization line has been identified. This explicit
case needs no general mixed-characteristic weight--monodromy theorem.
The entrance-path claims RESIDUE H.3, H.5, H.6 have their own
categorical proofs, unavailable in the supplied packet. Their
reported status is carried as an intake report, not promoted to a
newly checked theorem by this geometric calculation.

% END root_deligne.tex

% BEGIN root_core.tex
% Root mathematical contribution. Finding IDs identify deductions checked in this run.
\section{The exact abelian Artin--Tate consequence}
\label{sec:root-artin-tate}

We use the cohomological convention $\mathbb Q(1)$ of weight $-2$,
so that a geometric Frobenius at a good prime acts on its $\ell$-adic
realization by $p^{-1}$.  Coefficients may be extended to a number field
containing the values of the finite characters involved.  A rational
statement is recovered by grouping coefficient-conjugate characters.

\paragraph{Finding BR-AT-01 (the full shifted-zero consequence).}
Let $M$ be a finite direct sum of rank-one pure motives
\[
 M=\bigoplus_{\chi,n} (\chi\otimes\mathbb Q(n))^{\oplus m_{\chi,n}},
 \qquad m_{\chi,n}\in\mathbb Z_{\geq0},\quad n\in\mathbb Z,
\]
where $\chi$ is a finite-order character of
$\operatorname{Gal}(\overline{\mathbb Q}/\mathbb Q)$ with abelian image.
Use the primitive Dirichlet character associated with each $\chi$ for
its full Artin $L$-function.  Then
\[
 L(M,s)=\prod_{\chi,n}L(\chi,s+n)^{m_{\chi,n}}.
\]
Under the adopted Alp\"oge theorem, the nontrivial zeros contributed by
the summand of twist $n$ have real part $\tfrac12-n$.
These zero multiplicities add in the displayed product: none is cancelled
by a pole of another factor.  For a pure summand of weight $w=-2n$, this
is the expected line $\Re s=(w+1)/2$.

For a mixed object $V$ whose weight-graded pieces have this form, and a
finite set $S$ containing every ramified prime, the precise assertion is
\[
 L^S(V,s)=L^S(\operatorname{gr}^W V,s)
 =\prod_{\chi,n}L^S(\chi,s+n)^{m_{\chi,n}}.
\]
If Euler factors at primes in $S$ are deleted, the resulting extra zeros
lie on the lines $\Re s=-n$; they must not be called violations of the
nontrivial-zero assertion.  At ramified primes, no general multiplicativity
of inertia invariants for mixed extensions is being assumed.

\begin{proof}
At an unramified prime $p$, the eigenvalue for $\chi\otimes\mathbb Q(n)$
is $\chi(p)p^{-n}$, giving the Euler factor
$(1-\chi(p)p^{-s-n})^{-1}$.  The identity of the convergent Euler
products follows in a right half-plane and extends to their meromorphic
continuations.  This also proves the direct-sum formula.  On replacing
$z=s+n$, the strip $0<\Re z<1$ becomes
$-n<\Re s<1-n$; the adopted Dirichlet theorem places every zero there
on $\Re s=1/2-n$.  A primitive nonprincipal Dirichlet $L$-function is
entire; a principal primitive factor is $\zeta(s+n)$ and has its only
pole at $s=1-n$.  Such poles have integral real part, while the indicated
nontrivial zeros have half-integral real part.  Hence they cannot cancel.

For a filtration preserved by Frobenius, choose a basis respecting the
filtration.  Its matrix is block upper triangular; the determinant of
$1-p^{-s}\operatorname{Frob}_p$ is the product of the determinants of
the diagonal blocks.  At every $p\notin S$ this gives the mixed formula.
Deleting a nonzero local rank-one factor multiplies the full $L$-function
by $1-\alpha p^{-s-n}$ with $|\alpha|=1$.  Its zeros satisfy
$p^{-\Re s-n}=1$, hence $\Re s=-n$.  A ramified one-dimensional
character with zero inertia-invariant space contributes the factor $1$,
so there is nothing to delete in that case.  This proves all the assertions.
\end{proof}

The elementary Euler-factor argument is included because it specifies the
mixed, twisted, and incomplete forms not stated by the Lean declarations.
It does not assert that every general mixed-motivic $L$-function is
determined at every bad prime by its semisimplification.  In particular,
calling an object ``mixed Artin--Tate'' does not make its extension class
part of a Dirichlet character.

\paragraph{Archimedean convention.}
For $\chi(-1)=(-1)^a$, $a\in\{0,1\}$, the factor of
$\chi\otimes\mathbb Q(n)$ is $\Gamma_{\mathbb R}(s+n+a)$, where
$\Gamma_{\mathbb R}(s)=\pi^{-s/2}\Gamma(s/2)$.  Its real Frobenius
acts as $\chi(-1)(-1)^n$.  Both the Hodge type $(-n,-n)$ and that
sign are needed to recover the archimedean factor.  The sign alone does
not recover the weight: $\mathbb Q(0)$ and $\mathbb Q(2)$ both have
trivial real Frobenius.  The corrected real-Frobenius classification at
cote 18/038 is treated with the original text below.

\subsection{Mixed coefficients: a rigorous obstruction to reconstructing extensions}

Deligne--Goncharov (2005), \S1.1, equations (1.1.2), (1.1.4), and the
proof of Proposition 1.7, construct the Kummer extension $K(x)$ and
identify
\[
 \operatorname{Ext}^1_{\mathrm{MT}(\mathbb Q)}
   (\mathbb Q(0),\mathbb Q(1))=\mathbb Q^\times\otimes_{\mathbb Z}\mathbb Q.
\]
These exact locators are in the supplied file \nolinkurl{Sec1ccc.tex}; this
identification is a cited theorem, not a new result here.

\paragraph{Finding BR-MIX-01 (the Kummer extension survives identical
good Euler factors).}
There are nonisomorphic mixed Tate motives with the same weight-graded
object and the same Euler polynomials at every prime of simultaneous
good reduction.  Explicitly,
\[
 0\longrightarrow\mathbb Q(1)\longrightarrow K(2)
 \longrightarrow\mathbb Q(0)\longrightarrow0
\]
is nonsplit, whereas $\mathbb Q(0)\oplus\mathbb Q(1)$ is split.  For
$S$ containing $2$ and the realization prime $\ell$, both have
\[
 L^S(s)=\zeta^S(s)\zeta^S(s+1).
\]
Thus even complete knowledge of these good Euler factors, their
meromorphic continuation, and all their zeros does not reconstruct
the mixed extension.  This refutes that particular reconstruction
rule; it is not a logical independence theorem about RH.

\begin{proof}
The $2$-adic valuation gives a homomorphism
$v_2:\mathbb Q^\times\to\mathbb Z$ and hence a linear map
$\mathbb Q^\times\otimes\mathbb Q\to\mathbb Q$ taking $2\otimes1$
to $1$.  Consequently the class of $2$ is nonzero, and the cited
Ext-identification makes its extension nonsplit.  An isomorphism to
the split object would preserve the unique weight filtration and
give a splitting, a contradiction.  By Deligne--Goncharov,
Proposition 1.7 and its proof, this Kummer object is unramified away
from $2$ (and the $\ell$-adic realization is used away from $\ell$).
For $p\notin S$ its Frobenius matrix in a weight-adapted basis is
upper triangular with diagonal entries $p^{-1},1$.  Its Euler
polynomial is therefore $(1-T)(1-p^{-1}T)$, the polynomial of the
split object.  Taking their Euler products proves the identity.
\end{proof}

The same proof works for every $x\in\mathbb Q^\times$ with a nonzero
valuation somewhere, after deleting the primes dividing $x$.
The unipotent radical of the mixed Tate motivic Galois group, its
extension classes, and its period matrix are additional information.
This is a concrete reason why a theorem about the zero divisor of
these $L$-functions does not by itself compute the groups in step 17
of \emph{Remark}.  In particular the tentative footnote there gives
neither the degrees nor the primary parts required for a $K$-group
vanishing statement.  The assertion that \emph{all} positive-degree
$K$-groups of $\mathbb Z$ vanish is false already for
$K_1(\mathbb Z)=\{\pm1\}$; it is not an assertion made by that footnote.

\section{Cote 1: what the Hilbert-space trace theorem settles}
\label{sec:root-trace}

The supplied scans 1/079--080 contain the Fredholm determinant estimates
and the trace/eigenvalue formula; 1/111--112 contain the genus-zero
argument.  They do not mention $\xi$, $\Xi$, a Riemann-zero operator,
or RH.  Those are applications in the supplied SA23 intake, \S4.1,
and in this work.  An exact attribution therefore distinguishes the
original analytic theorem from its later arithmetic use.

Write
\[
 \xi(s)=\tfrac12s(s-1)\pi^{-s/2}\Gamma(s/2)\zeta(s),\qquad
 \Xi(t)=\xi(\tfrac12+it),\qquad F(w)=\frac{\Xi(\sqrt w)}{\Xi(0)}.
\]
Evenness of $\Xi$ makes $F$ single-valued and entire.  The classical
Hadamard product and zero-counting facts used here are Titchmarsh--
Heath-Brown (1986), \S2.12 (the $\xi$ product), Theorems 9.3--9.4,
and equation (9.4.3), p.~214 (the Riemann--von Mangoldt count).
We use their statements, not a
new proof of those analytic inputs.

\paragraph{Finding BR-TR-01 (a positive trace-class realization, and its limit).}
There is a diagonal trace-class operator $B$ with
\[
 F(w)=\det(I-wB).
\]
It can be constructed before RH by taking one zero $t_j$ from each
pair $\{t,-t\}$ of $\Xi$, counted with multiplicity, and putting
$B=\operatorname{diag}(t_j^{-2})$ on $\ell^2$.
For this normal realization, RH is equivalent to $B\geq0$.
Consequently the adopted RH gives such a positive realization.
Every trace-class operator $T$ with $\det(I-wT)=F(w)$ has the same
nonzero eigenvalues with algebraic multiplicity and satisfies
\[
 \operatorname{Tr}(T^k)=\sum_jt_j^{-2k}>0\quad(k\geq1)
\]
under RH.  It is false that every such $T$ is positive or self-adjoint.
Moreover, no trace-class operator has $\Xi(t)$ itself as an ordinary
Fredholm determinant in an affine variable $t$.

\begin{proof}
The nontrivial zeta zeros have counting function $O(T\log(T+2))$.
Partial summation, or a dyadic decomposition, gives
$\sum_j|t_j|^{-2}<\infty$; finitely many bounded zeros cause no issue
because $\Xi(0)\ne0$.  Thus the displayed diagonal operator is trace
class.  The entire even function $\Xi$ has order one, so $F$ has
order at most $1/2$.  Hadamard factorization in order less than one,
together with $F(0)=1$, gives the genus-zero product
$F(w)=\prod_j(1-wt_j^{-2})$, which is precisely its Fredholm determinant.

If RH holds, all $t_j$ are real and nonzero; thus every diagonal
entry of $B$ is positive.  Conversely $B\geq0$ implies
$t_j^{-2}>0$, so $t_j$ is real.  Indeed, for a nonzero complex number
$t$, $t^2$ positive real forces $t$ real; the purely imaginary
alternative has negative square.  Hence every zero of $\xi(s)$
has real part $1/2$.

For a general trace-class $T$, the Fredholm zero theorem identifies
its nonzero eigenvalues and their algebraic multiplicities with the
reciprocals of the zeros of its determinant.  The trace theorem then
gives the power-trace formula; the sum converges absolutely.  To see
why positivity of $T$ does not follow, start with $B$ under RH and
replace its first two diagonal entries $b_1,b_2$ by the block
$\left(\begin{smallmatrix}b_1&1\\0&b_2\end{smallmatrix}\right)$.
This is a finite-rank perturbation, remains trace class, has the
same determinant, and is not self-adjoint in the given Hilbert
inner product.  It therefore is not a positive operator there.

Finally suppose $\Xi(t)=c\det(I+(t-a)T)$ for trace-class $T$ and
$c\ne0$.  Its zeros would be $a-\lambda_j^{-1}$, with the nonzero
eigenvalues $\lambda_j$ of $T$ counted algebraically.  The trace-class
eigenvalue inequality forces $\sum_j|\lambda_j|<\infty$.
But the Riemann--von Mangoldt asymptotic gives divergent sum of the
reciprocals of the distances of the $\Xi$ zeros from any fixed $a$.
For example, in dyadic height intervals the count is of order
$2^k k$, whereas every such distance is at most $C2^k$ for large
$k$ (the imaginary parts of the $t$-zeros are bounded by $1/2$).
The sum over that interval is at least a positive constant times
$k$.  This contradicts the required summability.  A zero at $a$
would already contradict the nonzero determinant value at $t=a$.
\end{proof}

The construction uses the zero set; it does not solve the stronger
geometric Hilbert--P\'olya problem of producing the operator from
primes, correspondences, or a cohomology theory independently of its
prescribed spectrum.  The result settles the exact positive normal
realization and its trace inequalities, not that additional
construction problem.

\paragraph{Finding BR-TR-02 (a strictly positive moment hierarchy).}
Under RH put $b_j=t_j^{-2}>0$ and
$m_r=\sum_j b_j^{r+1}$ for $r\geq0$.  For every integer $d\geq0$,
both Hankel matrices $(m_{a+b})_{0\leq a,b\leq d}$ and
$(m_{a+b+1})_{0\leq a,b\leq d}$ are positive definite.
For $x\geq0$,
\[
 H(x)=-\frac{F'(-x)}{F(-x)}
      =\operatorname{Tr}\bigl(B(I+xB)^{-1}\bigr)
      =\sum_j\frac{b_j}{1+xb_j},
 \qquad (-1)^r H^{(r)}(x)>0\quad(r\geq0).
\]
These are quantitative inequalities in the cote 1 trace framework,
not assertions that its unspecified geometric operators are positive.

\begin{proof}
The positive measure $\mu=\sum_jb_j\delta_{b_j}$ is finite because
$B$ is trace class, and it has infinitely many distinct support points:
$\Xi$ has infinitely many zeros and finite multiplicity at each point.
For a nonzero polynomial $P(u)=\sum_{a=0}^d c_au^a$,
\[
 \sum_{a,b=0}^d\overline{c_a}c_bm_{a+b}
   =\int |P(u)|^2\,d\mu(u)>0.
\]
A nonzero polynomial cannot vanish on infinitely many distinct points.
Inserting an additional factor $u>0$ proves the second positive
definiteness assertion.  The logarithmic derivative follows from
the normally convergent product for $F$; for $x\geq0$ all its
denominators are positive.  On compact subintervals of $[0,\infty)$,
termwise differentiation is dominated by
$r!\|B\|^r\sum_jb_j$.  Thus
\[
 (-1)^rH^{(r)}(x)
 =r!\sum_j\frac{b_j^{r+1}}{(1+xb_j)^{r+1}}>0.
\]
This proves convergence, the identities, and strictness.
\end{proof}

\section{Weights, monodromy, and the additional geometric input}
\label{sec:root-weights}

Deligne (1980), \emph{La conjecture de Weil II}, Theorem 1.8.4,
states that for a punctually pure lisse sheaf of weight $w$ on the
punctured smooth curve over a finite field, the $j$th graded piece
of its local monodromy filtration is pure of weight $w+j$.
This is a theorem with a geometric-purity hypothesis, not an
implication from a list of zeros of $\zeta(s)$.
The mixed-characteristic question for $H^i$ of a smooth proper
variety over a local field is the corresponding assertion that
$\operatorname{gr}_j^M H^i$ has weight $i+j$.
An exact modern statement attributing it to Deligne is Scholze (2012),
\emph{Perfectoid Spaces}, Conjecture 9.3; Theorem 9.6 proves the
set-theoretic complete-intersection case in smooth projective toric
varieties.  No general mixed-characteristic proof is contained in
the supplied RH declarations.

\paragraph{Finding BR-WM-01 (a determinant and the Tate covariance do not
determine monodromy purity).}
For every prime power $q>1$ there are two Weil--Deligne linear
algebra data with the same semisimple Frobenius and the same
Frobenius determinant, both satisfying $FNF^{-1}=q^{-1}N$, for
which the monodromy-purity statement centered at weight $1$ has
different truth values.

\begin{proof}
Take a two-dimensional vector space with basis $e_0,e_2$ and
$Fe_0=e_0$, $Fe_2=qe_2$.  In the first datum put
$Ne_2=e_0$, $Ne_0=0$; in the second put $N=0$.
The covariance is checked on the basis and holds in both cases.
Their common Frobenius determinant is $(1-T)(1-qT)$.
For the first datum the monodromy filtration centered at zero is
$M_{-2}=0$, $M_{-1}=M_0=\langle e_0\rangle$, $M_1=V$.
Its graded pieces in degrees $-1$ and $1$ have eigenvalues $1$
and $q$, hence weights $0$ and $2$, which are $1-1$ and $1+1$.
For the second datum the monodromy filtration is concentrated in
degree zero, so its only graded piece is $V$, with weights $0,2$;
it is not pure of weight $1$.  This proves the claim.
\end{proof}

These data refute the proposed \emph{linear-algebra deduction} from
Frobenius roots and covariance.  They are not a counterexample to
Deligne's geometric conjecture: no realization as the cohomology
of a smooth proper variety of the required weight is asserted for
the second datum.  Also, the local Euler factor defined on
$\ker N$ differs between the two examples.  It would be incorrect
to call these examples indistinguishable by every ramified local
$L$-factor.  What fails is precisely the attempted passage from
a determinant on $V$ to its monodromy-graded purity.

\subsection{The tame triangle and the nodal universal curve}

The intake SA20, row C.008, attributes its triangle to cote 122;
the original scan for that row is not among the supplied files.
Its finite $\mathbb Z/\ell^a$ coefficients and twists type the
duality maps.  Numerical Frobenius weights require a rational
$\ell$-adic realization and a purity theorem.  After that passage,
$H^1(I_t,\mathbb Q_\ell)=\mathbb Q_\ell(-1)$ has weight $2$;
the monodromy logarithm lowers weights by $2$.  These facts explain
the appearance of adjacent twists, but do not prove the
mixed-characteristic monodromy-purity assertion.

There is a further distinction in the RESIDUE intake, H.3--H.6:
the ordinary cohomology of a singular nodal fibre is not the same
filtered object as the nearby smooth cohomology with its full
monodromy filtration.  The following elementary calculation makes
the distinction exact.

\paragraph{Finding BR-NODE-01 (the weight-zero part belongs to the dual graph).}
Let $C/\mathbb F_q$ be a geometrically connected proper nodal curve,
with $\ell\ne\operatorname{char}\mathbb F_q$, whose geometric
components, nodes, and two branches at every node are all defined over
$\mathbb F_q$.  Write $\widetilde C=\coprod_v C_v$ for the
normalization and $\Gamma$ for the dual graph.  There is a
Frobenius-equivariant exact sequence
\[
 0\longrightarrow H^1(\Gamma,\mathbb Q_\ell)
 \longrightarrow H^1(C_{\overline{\mathbb F}_q},\mathbb Q_\ell)
 \longrightarrow\bigoplus_vH^1(C_v,\mathbb Q_\ell)
 \longrightarrow0.
\]
The left term has weight zero and the right term weight one.
In particular a split rational curve with one ordinary node has
$H^1\simeq\mathbb Q_\ell$ of weight zero.  Its smooth genus-one
nearby fibre, if chosen in a semistable degeneration, also has the
weight-two monodromy quotient; that quotient is not an additional
piece of the singular fibre's ordinary $H^1$.

\begin{proof}
For the finite normalization map $\nu$, the exact sequence of
\'etale sheaves is
\[
 0\longrightarrow\mathbb Q_\ell
 \longrightarrow\nu_*\mathbb Q_\ell
 \longrightarrow\bigoplus_{e\in E(\Gamma)}\mathbb Q_{\ell,e}
 \longrightarrow0.
\]
Away from the nodes the first map is an isomorphism.  At a node
its stalk is the diagonal inclusion $\mathbb Q_\ell\to
\mathbb Q_\ell^2$, whose cokernel is one-dimensional; choosing
an orientation identifies the quotient with the difference of
branch values.  Finite pushforward and cohomology give the usual
normalization cohomology.  The first part of the long exact
sequence is consequently
\[
 0\to\mathbb Q_\ell\to\mathbb Q_\ell^{V(\Gamma)}
 \xrightarrow{\partial}\mathbb Q_\ell^{E(\Gamma)}
 \to H^1(C)\to\bigoplus_vH^1(C_v)\to0.
\]
The cokernel of the incidence map $\partial$ is graph cohomology.
All vertices and branches are rational, so Frobenius acts trivially
on this graph complex, giving weight zero.  Each $C_v$ is a smooth
proper curve; its $H^1$ has weight one by the finite-field Weil
theorem.  If $C$ is rational with one node, its graph has one
loop and its normalization has no $H^1$, proving the special case.
For a semistable genus-one smoothing, use the cited Tate-uniformization
calculation (Ochiai, 2010, \S4, p.~4, displayed Tate-module and
$H^1$ exact sequences; Berger, 2004, \S II.4.1--II.4.3 gives the
period calculation): the nonzero rank-one monodromy pairs the
weight-zero subspace with a weight-two quotient, exactly as in the
first datum of BR-WM-01.
The nearby fibre has dimension two in $H^1$, whereas the singular
fibre has dimension one.  Thus the extra quotient cannot be
ordinary $H^1$ of the singular curve.
\end{proof}

The already established decomposition theorem concerns the pure
\emph{complex} pushed forward by a proper morphism.  It does not
assert that every ordinary stalk cohomology group of that complex
is pure of its cohomological degree at a singular stratum.  This
distinction prevents a false purity claim for the boundary curves
of $\overline{\mathcal M}_{g,n}$.  The supplied entrance-category
claims H.3--H.6 remain background at the level at which their intake
checked them; their missing original proofs have not been silently
promoted to checks made in this run.

\subsection{The standard conjectures: which additional gate remains}

Grothendieck's \emph{Standard conjectures on algebraic cycles}
(1969), \S2, p.~195, asks for algebraicity of the K\"unneth
projectors (C); \S3, pp.~195--197, asks for the Lefschetz
standard conjectures (A and B); \S4, pp.~197--198, asks for
positivity of the primitive algebraic intersection form (Hdg).
The implication to the finite-field Weil conjectures is stated in
\S4, Remark (2), p.~197; the relation to numerical equivalence
is in Remarks (3)--(4), pp.~197--198.
There is an essential positive consequence of finite-field purity:
Katz--Messing (1974), Theorem 2(1), pp.~76--77, constructs algebraic
K\"unneth projectors over finite fields.  The distinct weights make the
Frobenius characteristic polynomials in different degrees coprime;
Chinese-remainder polynomials in the algebraic Frobenius correspondence
then act as the cohomological projectors.  Theorem 1 of that paper
transfers the required Frobenius characteristic polynomials to crystalline
cohomology under its stated axioms.  These are already established
consequences of the finite-field theorem, and a complete new proof of
that theorem would give another route to them.  They do not, by
themselves, construct the inverse Lefschetz correspondence or prove
positivity on algebraic primitive cycles in general.  The classical
number-field Dirichlet assertion does not add those correspondences.

For the abelian Artin--Tate subcategory the relevant structures are
already explicit: Artin projectors are obtained from finite
correspondences and character idempotents, and Tate objects have
their standard duality.  BR-AT-01 settles the global zero-location
assertion for their $L$-functions.  For arbitrary smooth projective
varieties the current gate is still an algebraic correspondence
realizing the inverse Lefschetz operator, together with the
required positive pairing on algebraic primitive classes.  This
work supplies neither that correspondence nor a counterexample to
the standard conjectures.  The unresolved conjectures are recorded
as such, rather than marked ``refuted'' from the absence of a proof.

\section{D\'erivateurs, localization, and information forgotten}
\label{sec:root-derivators}

The supplied DER\_LOCALIZERS intake identifies the linear
``Thm(?)'' at cote 134-8/048--054 and its remaining acyclicity
condition C.10.  Its original scans and the full proof of its
general counterexample C.26 are not supplied.  A specific motivic
instance is nevertheless established in the literature:
Voevodsky (2000), Theorem 3.1.12, gives strict homotopy invariance
for Nisnevich cohomology of homotopy-invariant presheaves with
transfers over a perfect field; Proposition 3.2.3 identifies the
Suslin-complex localization; Theorem 3.2.6 gives the embedding.
Those are geometric inputs to the localization theorem, independent
of RH.

\paragraph{Finding BR-DER-01 (constant rational coefficients detect exactly
rational homology).}
Let $\mathcal C$ be a nonzero presentable stable
$\mathbb Q$-linear category, with a specified nonzero object
$\mathbf1$ and the colimit-preserving exact action
$D(\mathbb Q)\to\mathcal C$, $V\mapsto V\otimes\mathbf1$.
For a functor $f:A\to B$ between small categories, the induced map
on homotopy colimits of the constant diagram $\mathbf1$ is an
equivalence if and only if $Nf$ is a rational homology equivalence.
Thus this particular constant-object localizer does not change
when one equips $\mathcal C$ with additional Tate twists,
Frobenius actions, or $L$-functions while retaining the specified
rational action and a nonzero unit.

\begin{proof}
The tensoring of a presentable stable category over spaces, followed
by the $\mathbb Q$-linear action, identifies the homotopy colimit
of a constant diagram on $A$ with
$C_*(NA;\mathbb Q)\otimes\mathbf1$.  Let $K$ be the cone in
$D(\mathbb Q)$ of the chain map induced by $Nf$.  Exactness shows
that the homotopy-colimit map is an equivalence exactly when
$K\otimes\mathbf1=0$.
Every complex of vector spaces splits as its homology together
with a contractible complex: choose complements to boundaries
inside cycles and to cycles inside each chain group.  The
differential identifies the latter complement with the
boundaries one degree lower, giving the contractible summands.
This works for unbounded complexes degree by degree.  Hence
$K\simeq\bigoplus_i H_i(K)[i]$ in $D(\mathbb Q)$.
If some $H_i(K)\ne0$, choose a nonzero vector and a linear
projection onto its span.  This exhibits $\mathbb Q[i]$ as a
retract of $K$.  Its image $\mathbf1[i]$ is then a nonzero
retract of $K\otimes\mathbf1$, a contradiction.  Conversely,
$K=0$ plainly implies $K\otimes\mathbf1=0$.
Thus the latter vanishing is equivalent to zero rational homology
of the cone, which is exactly the asserted condition on $Nf$.
\end{proof}

The quantifiers are deliberately explicit.  This proves the
constant-unit statement needed for rational motivic categories;
it does not identify all coefficient-dependent localizers, and
does not assert that every reversible localizer is of this form.
It also avoids the intake's unverified blanket full-faithfulness
assertion for arbitrary coefficient rings.

\paragraph{Finding BR-DER-02 (the reduced Tate object is essential).}
In geometric motives over a field, with the rational or usual
integral coefficients for which the cited motivic identities hold,
\[
 M(\mathbb G_m)\simeq\mathbf1\oplus\mathbf1(1)[1],\qquad
 \widetilde M(\mathbb G_m)[-1]\simeq\mathbf1(1).
\]
The formula in the DER\_LOCALIZERS intake omits the reduction if
$\mathbb Z_{\mathrm{tr}}(\mathbb G_m)$ means the full representable
presheaf with transfers.  Its full localized motive shifted by
$-1$ has an additional $\mathbf1[-1]$ summand.

\begin{proof}
Cover $\mathbb P^1$ by its two affine lines.  The motivic
Mayer--Vietoris triangle is
\[
 M(\mathbb G_m)\longrightarrow\mathbf1\oplus\mathbf1
 \longrightarrow M(\mathbb P^1)\longrightarrow M(\mathbb G_m)[1].
\]
The projective-line decomposition is
$M(\mathbb P^1)=\mathbf1\oplus\mathbf1(1)[2]$.
Each affine-line inclusion has degree-one component to
$\mathbf1$ and zero component to the Tate summand; the latter
lies in $\operatorname{Hom}(\mathbf1,\mathbf1(1)[2])
=\operatorname{Pic}(\operatorname{Spec}k)=0$.
After an invertible change of basis in $\mathbf1\oplus\mathbf1$,
the middle map is the identity on one unit summand and zero on
the other, with zero map to $\mathbf1(1)[2]$.  Taking its fibre
gives $M(\mathbb G_m)=\mathbf1\oplus\mathbf1(1)[1]$.
The rational point $1\in\mathbb G_m$ splits the augmentation;
its reduced fibre is exactly the second summand.  Shifting proves
the assertion.  The projective-line definition and reduction are
Voevodsky (2000), \S2.1 following Proposition 2.1.3; the Picard
calculation is Corollary 3.4.3.  This proof spells out the corrected
form rather than reproving the construction of motives.
\end{proof}

\section{A concrete boundary for the anabelian and Teichm\"uller bridges}
\label{sec:root-anabelian}

\paragraph{Finding BR-ANA-01 (identical abelian cohomological data do not
reconstruct a four-punctured sphere).}
The two hyperbolic curves over $\mathbb Q$
\[
 U_2=\mathbb P^1\setminus\{0,1,\infty,2\},\qquad
 U_3=\mathbb P^1\setminus\{0,1,\infty,3\}
\]
are not isomorphic even over $\overline{\mathbb Q}$, although
their geometric $\ell$-adic cohomology representations are
isomorphic in every degree.  For every prime $p\notin\{2,3\}$
and integer $r\geq1$ they have the same point count and local zeta
function:
\[
 \#U_i(\mathbb F_{p^r})=p^r-3,\qquad
 Z(U_i/\mathbb F_p,T)=\frac{(1-T)^3}{1-pT}.
\]
Thus the global good-prime cohomological zeta functions agree,
even after imposing RH on their Dirichlet/Tate factors.

\begin{proof}
An isomorphism of the affine smooth curves extends uniquely to
their smooth projective compactifications and permutes their
four missing points.  The six possible cross-ratios of an
unordered quadruple with parameter $t$ are
\[
 t,\quad 1-t,\quad t^{-1},\quad(1-t)^{-1},\quad
 t/(t-1),\quad(t-1)/t.
\]
For $t=2$ the resulting set is $\{2,-1,1/2\}$; for $t=3$ it
is $\{3,-2,1/3,-1/2,3/2,2/3\}$.  They are disjoint, so no
projective-linear transformation carries one unordered boundary
to the other.  This proves nonisomorphism.

For either boundary $D$ the localization sequence on
$\mathbb P^1_{\overline{\mathbb Q}}$ is Galois equivariant.
Since its four points are rational, the residue sequence gives
\[
 H^0(U_i,\mathbb Q_\ell)=\mathbb Q_\ell,\qquad
 H^1(U_i,\mathbb Q_\ell)
 =\ker\bigl(\mathbb Q_\ell(-1)^4
             \xrightarrow{\text{sum}}\mathbb Q_\ell(-1)\bigr)
 \simeq\mathbb Q_\ell(-1)^3,
\]
and higher ordinary cohomology vanishes.  This determines the
representations independently of $t$.  Alternatively compactly
supported cohomology has $H_c^1=\mathbb Q_\ell^3$ and
$H_c^2=\mathbb Q_\ell(-1)$.
At each stated prime the four rational points remain distinct,
so subtracting them from $\mathbb P^1(\mathbb F_{p^r})$ gives
$p^r+1-4$.  Exponentiating
$\sum_{r\geq1}(p^r-3)T^r/r$ yields the claimed zeta function.
\end{proof}

This is an explicit counterexample to reconstructing the
Teichm\"uller parameter from the indicated cohomological or
$L$-function data.  It does not contradict Grothendieck's
anabelian reconstruction proposals, whose input is a nonabelian
arithmetic fundamental group with its Galois action and peripheral
structure.  It shows why the abelian Artin--Tate result, even with
effective recovery of all its finite characters, is not itself
that nonabelian input.  A comparison with $\widehat{\mathrm{GT}}$
would have to specify and prove compatibility of a construction
with the noncommutative fundamental groups, their inertia, and
the tower maps.  None is constructed by the coincident words
``hexagon'' and ``tower'' in the supplied proof sketch.


% END root_core.tex

% BEGIN root_weight_criterion.tex
\section{A finite exterior-power gate for weight--monodromy}
\label{sec:finite-weight-gate}

Scholze (2012), arXiv:1111.4914, Lemma 9.5, printed p.~49,
records Deligne's tensor-power criterion: purity of the monodromy
filtration is equivalent to upper weight bounds on the inertia
invariants of every tensor power and every dual tensor power.
The following finite exterior-power formulation is proved here in
full. It reduces that infinite collection of representations; it
does not assert the needed bounds for an arbitrary geometric
representation. No historical novelty claim is made.

\paragraph{Finding BR-WM-02 (finite exterior-power criterion).}
Let $V$ be a complex vector space of dimension $d\geq1$, let
$q>1$, let $w\in\mathbb R$, and let $F$ be invertible and $N$
nilpotent with $FNF^{-1}=q^{-1}N$. Write $M$ for the monodromy
filtration centered at zero. On exterior powers, $N$ denotes the
induced derivation; on the dual it is $-{}^tN$. The following
three conditions are equivalent.
\begin{enumerate}
\item On $\operatorname{gr}_j^M V$, every eigenvalue of $F$
has absolute value $q^{(w+j)/2}$.
\item For each $1\leq m\leq d$, all eigenvalues on
$\ker(N:\bigwedge^m V\to\bigwedge^m V)$ have absolute value
at most $q^{mw/2}$, and all eigenvalues on
$\ker(N:\bigwedge^m V^\vee\to\bigwedge^m V^\vee)$ have absolute
value at most $q^{-mw/2}$.
\item For every $1\leq m\leq d$ there are constants
$C_m,C'_m>0$ and $A_m,A'_m\geq0$ such that, for every
positive integer $r$,
\[
 \left|\operatorname{Tr}\left(F^r\mid
             \ker N_{\wedge^m V}\right)\right|
       \leq C_m r^{A_m}q^{rmw/2},
 \quad
 \left|\operatorname{Tr}\left(F^r\mid
             \ker N_{\wedge^m V^\vee}\right)\right|
       \leq C'_m r^{A'_m}q^{-rmw/2}.
\]
\end{enumerate}
When these conditions hold, condition (3) holds with $A_m=A'_m=0$
and the respective dimensions in place of $C_m,C'_m$ (or any
positive constants when a space is zero).

\begin{proof}
We first justify replacing $F$ by its semisimple part $F_s$.
If $V_\lambda$ is the generalized $\lambda$-eigenspace of $F$,
the covariance identity implies
$N(V_\lambda)\subset V_{\lambda/q}$.
Since $F_s$ is scalar $\lambda$ on $V_\lambda$, it follows
that $F_sNF_s^{-1}=q^{-1}N$. The unipotent part
$F_u=F_s^{-1}F$ commutes with $N$ and $F_s$.
The same statements hold on exterior powers and duals.
Their kernels of $N$ and all monodromy-graded spaces are
preserved by both parts. On any such subquotient, multiplying
the semisimple operator by its commuting unipotent part
does not change its eigenvalues or the traces of its powers.
Thus all three conditions are unchanged by this replacement.

There is a decomposition into $F_s,N$-stable Jordan strings.
Here is a direct proof, including the invariant complement.
Let $l$ be the largest Jordan length of $N$ on a nonzero
remaining summand. Choose an $F_s$-eigenvector $v$ of
eigenvalue $\mu$ with $N^{l-1}v\ne0$, and put
$e_j=N^jv$ for $0\leq j<l$.
Then $F_se_j=\mu q^{-j}e_j$. Choose an eigenfunctional
$a$ of $F_s$ with
$a(F_sx)=\mu q^{-(l-1)}a(x)$ and $a(e_{l-1})=1$.
Such a functional exists by semisimplicity. It vanishes
on $e_j$ for $j<l-1$, because those eigenvalues are distinct.
The map
\[
 P(x)=\sum_{j=0}^{l-1}a(N^{l-1-j}x)e_j
\]
is the identity on the span $W$ of the $e_j$.
The equality $N^l=0$ shows $PN=NP$, including the endpoint
terms. Also $N^rF_s=q^rF_sN^r$, so the $j$th coefficient
of $P(F_sx)$ is $\mu q^{-j}$ times that of $P(x)$.
Therefore $PF_s=F_sP$. Its kernel is an invariant
complement. Induction on dimension gives the asserted
decomposition.

For a string of length $l$, assign
$h(e_j)=l-1-2j$. This is its monodromy grading:
$N$ lowers it by two, and $N^a$ is an isomorphism from
grade $a$ to grade $-a$ whenever those grades occur.
Thus condition (1) on this string says
\[
       |\mu|=q^{(w+l-1)/2}.                         \tag{WM.1}
\]
The monodromy filtration of the direct sum is the sum
of these filtrations, by its defining properties.

Assume condition (2). The determinant line
$\bigwedge^lW$ is contained in $\ker N$ in
$\bigwedge^lV$, and its eigenvalue is
$\delta=\mu^l q^{-l(l-1)/2}$.
The dual determinant line is an invariant line in
$\bigwedge^lV^\vee$, because the direct-sum decomposition
also splits the dual representation. The two bounds with
$m=l\leq d$ give
$|\delta|\leq q^{lw/2}$ and
$|\delta|^{-1}\leq q^{-lw/2}$, hence equality.
Solving it gives (WM.1) for every string, proving (1).

Assume condition (1). On each string define
\[
 H e_j=(l-1-2j)e_j,\qquad
 E e_j=j(l-j)e_{j-1},\quad E e_0=0.
\]
Direct multiplication gives $[E,N]=H$.
Make the strings orthogonal, and choose their basis norms
recursively by
\[
 \|e_{j+1}\|^2=(j+1)(l-j-1)\|e_j\|^2
       \quad(0\leq j<l-1).
\]
All these factors are positive, and these norms give
$N^*=E$. On $\bigwedge^mV$ with its induced positive
Hermitian form, the induced derivations again satisfy
$N^*=E$ and $[E,N]=H$.
If $x$ is in $\ker N$ and has $H$-weight $h$, then
\[
 h\|x\|^2=\langle [E,N]x,x\rangle=-\|Ex\|^2\leq0.
\]
Thus $h\leq0$ for every nonzero such vector.
The operators $F_s$ and $H$ commute and are semisimple.
By (WM.1), on $\bigwedge^mV$ each simultaneous
$F_s,H$-eigenvector of $H$-weight $h$ has Frobenius
absolute value $q^{(mw+h)/2}$: multiply the eigenvalue
absolute values of the chosen wedge basis.
The kernel of $N$ is stable under $F_s$ and $H$, so it
has a simultaneous eigenbasis. Its eigenvalues therefore
obey the first bound of (2). The dual strings have the
same lengths and central weight $-w$. Applying the
same argument to them gives the second bound.

Condition (2) implies (3), with no power of $r$,
by summing the eigenvalues of $F^r$ with their
algebraic multiplicities. For the converse we give
the required elementary power-sum argument. For a
finite-dimensional invertible operator with distinct
eigenvalues $\alpha$ and multiplicities $b_\alpha>0$,
\[
 \sum_{r\geq1}\operatorname{Tr}(F^r)z^r
       =\sum_\alpha\frac{b_\alpha\alpha z}{1-\alpha z}.
\]
A bound $|\operatorname{Tr}(F^r)|\leq Cr^AR^r$
makes the left side holomorphic for $|z|<R^{-1}$.
If $|\alpha|>R$, the right side has a pole at
$z=\alpha^{-1}$ inside that disc. The other terms
are holomorphic there, and $b_\alpha\ne0$, so the
pole cannot cancel. The identity initially holds
in a smaller disc and hence contradicts analytic
continuation. Therefore $|\alpha|\leq R$.
Apply this to the two invariant spaces in (3),
with $R=q^{mw/2}$ and $R=q^{-mw/2}$ respectively.
Zero-dimensional invariant spaces require no argument.
This proves (2) and completes all implications.
\end{proof}

\paragraph{Geometric use and the remaining gate.}
For an $\ell$-adic representation with unipotent inertia,
inertia invariants equal $\ker N$. Apply BR-WM-02
after every complex embedding of the algebraic Frobenius
eigenvalues. For $V=H^i(X_{\overline K},\mathbb Q_\ell)$,
take $w=i$ and $d=\dim V$. Weight--monodromy for this
specified $V$ is therefore equivalent to the two
families of trace bounds in (3), for only
$1\leq m\leq d$. A finite extension of $K$ first
makes the inertia unipotent; the weight assertion
is invariant under that extension.

This is a strictly finite list of representations,
although the trace bounds still quantify over all
positive Frobenius powers $r$. It uses actual
inertia invariants and both dual directions. Bounds
only for alternating total traces, only for
$\det(1-TF\mid V)$, or only in one dual direction
are not substituted for these hypotheses.
In a geometric realization, exterior projectors
place these representations inside tensor powers
of the cohomology, so a prospective positivity
argument has explicit targets. The adopted
Dirichlet zero-location theorem alone does not
establish these local trace bounds.

% END root_weight_criterion.tex

% BEGIN consequences_criteria_WM_formality.tex
\section{Weight--monodromy, Hyodo--Kato cohomology, and rigid analytic
formality}\label{weightmonodromy-hyodokato-cohomology-and-rigid-analytic-formality}

\textbf{Scope and provenance.} Literature cutoff: 20 November 2026.
Checked on 21 November 2026 UTC by consequences\_criteria, GPT-6 Astra
Pro; reasoning effort is not exposed. The adopted RH and Dirichlet-GRH
inputs are not needed for these local geometric results. The sources
below are additional geometric inputs. The 2026 preprints are used as
cited theorems, with their statements and indicated dependency arguments
inspected; their complete foundational proofs have not been
independently reconstructed.

\subsection{CC-WM01 --- Exact known classes of
weight--monodromy}\label{cc-wm01-exact-known-classes-of-weightmonodromy}

Fix the convention that geometric Frobenius \(F\) and monodromy satisfy
\(FNF^{-1}=q^{-1}N\), where \(q\) is the residue-field cardinality.
Monodromy purity of weight \(i\) means that eigenvalues on
\(\operatorname{gr}^M_j\) are algebraic and have absolute value
\(q^{(i+j)/2}\) under every complex embedding. The monodromy filtration
is centered at zero.

\begin{longtable}[]{@{}
  >{\raggedright\arraybackslash}p{(\columnwidth - 4\tabcolsep) * \real{0.3333}}
  >{\raggedright\arraybackslash}p{(\columnwidth - 4\tabcolsep) * \real{0.3333}}
  >{\raggedright\arraybackslash}p{(\columnwidth - 4\tabcolsep) * \real{0.3333}}@{}}
\toprule\noalign{}
\begin{minipage}[b]{\linewidth}\raggedright
Primary result and inspected locator
\end{minipage} & \begin{minipage}[b]{\linewidth}\raggedright
Exact relevant hypotheses
\end{minipage} & \begin{minipage}[b]{\linewidth}\raggedright
Conclusion and boundary
\end{minipage} \\
\midrule\noalign{}
\endhead
\bottomrule\noalign{}
\endlastfoot
Ito and Zou (2026), arXiv:2610.07890v1, \textbf{Theorem B}, p.~2;
\textbf{Theorem 4.4}, pp.~12--13 & \(K/\mathbf Q_p\) finite; \(X/K\)
hyper-Kähler: smooth projective, geometrically simply connected, with
\(H^0(X,\Omega^2_{X/K})=K\sigma\) for a nowhere-degenerate two-form;
definition, p.~1 & Weight--monodromy for every \(0\le i\le2\dim X\) and
every prime \(\ell\), including \(p\), via the associated Weil--Deligne
representation. \textbf{No \(b_2\ge4\) hypothesis occurs in this
theorem}; that restriction concerns separate compatibility results. \\
Scholze (2012), arXiv:1111.4914v1, \textbf{Theorem 9.6}, p.~49;
Conjecture 9.3, p.~48 & \(k\) a characteristic-zero local field; \(Y/k\)
geometrically connected, proper and smooth; \(Y\) a
\textbf{set-theoretic} complete intersection in a projective smooth
toric variety & Weight--monodromy in all degrees for \(\ell\ne p\). \\
Binda, Kato, and Vezzani (2025), arXiv:2207.00369v2, \textbf{Theorem
1.2}, p.~2; \textbf{Theorem 5.6}, p.~28 & \(K/\mathbf Q_p\) finite;
\(Y/K\) smooth and a \textbf{scheme-theoretic} complete intersection in
a projective smooth toric variety & The p-adic conjecture in every
degree. Theorem 5.6 more generally treats such a variety over
\(\mathbf C_p\), asserting quasi-purity of
\(H^i_{\mathrm{HK}}(Y^{\mathrm{an}})\). The
set-theoretic/scheme-theoretic distinction is explicit on p.~2. \\
Wear (2023), arXiv:2303.05610v1, \textbf{Theorem 1.1.2}, p.~3;
coefficient convention, pp.~1--2 & \(k\) a local field of residue
characteristic \(p\); \(A/k\) abelian; \(Y\) geometrically connected,
proper and smooth, and a \textbf{set-theoretic} complete intersection in
\(A\) & Weight--monodromy in all degrees for \(\ell\ne p\). This permits
equal or mixed characteristic and does not provide the p-adic case. \\
\end{longtable}

Scholze's \textbf{Theorem 9.4}, p.~48, states Deligne's
equal-characteristic result for \(k\) the local field at
\(x\in C(\mathbf F_q)\), where \(C/\mathbf F_q\) is a curve and the
variety comes from a proper smooth scheme over \(C\setminus\{x\}\). This
exactly scoped globally defined case suffices in his reduction.

The inspected Ito--Zou dependency chain is \textbf{Lemma 3.8}, p.~8
(purity in degrees at most two), \textbf{Lemma 4.2}, p.~12 (the
higher-degree tensor square is a summand of tensors of \(H^2\) and its
dual), and \textbf{Lemma 4.3}, p.~12 (tensor-square detection), followed
by Theorem 4.4. The complete last argument was read; its earlier
semisimplicity and Mumford--Tate inputs remain cited results.

In a semistable Hyodo--Kato model with \(q=p^a\), use the linear
operator \(\Phi^a\), rather than semilinear \(\Phi\), for traces. The
relation \(N\Phi=p\Phi N\) yields \(\Phi^aN(\Phi^a)^{-1}=q^{-1}N\).
Binda--Kato--Vezzani's p.~2 comparison discussion identifies the
pertinent potentially semistable p-adic realization; their Theorem 5.6's
final splitting argument was read, retaining the preceding motivic
comparison theorems as inputs.

\subsection{CC-WM02 --- The cochain-algebra
consequence}\label{cc-wm02-the-cochain-algebra-consequence}

Petrov and Zavyalov (2026), arXiv:2607.14517v2, \textbf{Theorem 1.2},
p.~2, is \textbf{Theorem 5.1 and Lemma 5.2}, p.~18. For \(X/K\) smooth
proper rigid analytic, \(K/\mathbf Q_p\) finite, weight--monodromy in
all degrees for a fixed prime \(\ell\), possibly \(p\), implies

\[
 R\Gamma_{\mathrm{et}}(X_{\mathbf C_p},\mathbf Q_\ell)
 \simeq \bigoplus_i H^i_{\mathrm{et}}(X_{\mathbf C_p},\mathbf Q_\ell)[-i]
\]

as underlying \(E_\infty\)-algebras. For \(\ell=p\), de Rham cochains
are formal too. Definitions 5.5--5.6, pp.~19--20, specify monodromy
purity; Examples 5.4, p.~19, use \(D_{\mathrm{pst}}\) at \(p\). No
canonical or Galois-equivariant splitting is asserted. The final proof
is on pp.~25--26.

Smooth proper \textbf{algebraic} varieties already have formal cochains
unconditionally, as the source explains on p.~2. Thus the substantive
additional application is to genuinely rigid analytic spaces. Its
\textbf{Construction 3.1 and Proposition 3.2}, p.~5, provide a smooth
proper rigid surface with nonformal cochains at every prime and in de
Rham cohomology. Smoothness and properness alone therefore cannot give
the needed hypothesis.

\subsection{CC-WM03 --- A finite family of exterior trace conditions
implies rigid
formality}\label{cc-wm03-a-finite-family-of-exterior-trace-conditions-implies-rigid-formality}

\textbf{Proposition.} Let \(K/\mathbf Q_p\) be finite and \(X/K\)
smooth, proper and rigid analytic. Fix a prime \(\ell\), including the
possibility \(\ell=p\). For each nonzero cohomology group let

\[
 V_i=D\bigl(H^i_{\mathrm{et}}(X_{\mathbf C_p},\mathbf Q_\ell)\bigr),
 \qquad d_i=\dim V_i,
\]

be the Weil--Deligne representation in Petrov--Zavyalov's Examples 5.4.
Choose a finite extension \(L/K\) killing the finite inertial actions in
these finitely many representations. Put \(q=\#k_L\), and write
\((F_i,N_i)\) for their Frobenius and monodromy over \(L\). Extend
scalars to an algebraic closure of the coefficient field and fix an
embedding of \(\overline{\mathbf Q}\) into that closure. Suppose that
all \(F_i\)-eigenvalues lie in its image, so they can be regarded as
algebraic numbers. For \(1\le m\le d_i\), define

\[
 T_{i,m}(r)=\operatorname{Tr}\left(F_i^r\mid
       \ker N_{\wedge^m V_i}\right),\qquad
 T^\vee_{i,m}(r)=\operatorname{Tr}\left((F_i^\vee)^r\mid
       \ker N_{\wedge^m V_i^\vee}\right).
\]

Here \(F_i^\vee\) denotes inverse-transpose action, and \(N\) acts as a
derivation on exterior powers, with \(N^\vee=-{}^tN\) on the dual. For
every \(i,m\) and embedding
\(\sigma:\overline{\mathbf Q}\hookrightarrow\mathbf C\), suppose that
positive constants \(C_{i,m,\sigma},C^\vee_{i,m,\sigma}\) and
nonnegative exponents \(A_{i,m,\sigma},A^\vee_{i,m,\sigma}\) exist such
that, for \textbf{every integer \(r\ge1\)},

\[
 \begin{split}
 |\sigma(T_{i,m}(r))|&\le C_{i,m,\sigma}r^{A_{i,m,\sigma}}q^{rmi/2},\\
 |\sigma(T^\vee_{i,m}(r))|&\le C^\vee_{i,m,\sigma}r^{A^\vee_{i,m,\sigma}}q^{-rmi/2}.
 \end{split} \tag{CC-WM.1}
\]

Then weight--monodromy holds in every degree over \(K\), and the
underlying \(E_\infty\)-algebra
\(R\Gamma_{\mathrm{et}}(X_{\mathbf C_p},\mathbf Q_\ell)\) is formal. For
\(\ell=p\), \(R\Gamma_{\mathrm{dR}}(X/K)\) is formal as well.
Conversely, weight--monodromy implies algebraicity and (CC-WM.1), with
exponents zero and constants equal to the respective kernel dimensions
when these dimensions are positive.

\textbf{Complete proof of the deduction.} There are finitely many
nonzero finite-dimensional cohomology groups. Intersect their open
inertia kernels and pass to a finite extension to obtain \(L\). Each
pair satisfies \(F_iN_iF_i^{-1}=q^{-1}N_i\).

Fix \(i\). Replacing \(F_i\) by its semisimple part leaves all relevant
traces and monodromy-graded eigenvalues unchanged, by the first part of
BR-WM-02. The same proof decomposes the pair into invariant Jordan
strings, with bases \(e_0,\ldots,e_{s-1}\) such that \(Ne_j=e_{j+1}\),
\(Ne_{s-1}=0\), and \(Fe_j=\mu q^{-j}e_j\). Since every \(\mu\) is
algebraic, these matrices define a pair over a number field. Thus every
\(\sigma\) can be applied to the pair, without requiring an algebraic
model for the full original Weil action. Exterior powers, duals and
kernels commute with this characteristic-zero field extension; their
power traces are precisely the algebraic numbers in (CC-WM.1).

Apply BR-WM-02 to this complex pair with dimension \(d_i\) and weight
\(w=i\). Its trace criterion gives

\[
 |\sigma(\alpha)|=q^{(i+j)/2}
 \quad\text{for every eigenvalue }\alpha
 \text{ on }\operatorname{gr}^M_jV_i.
\]

Doing this for every \(\sigma\) proves monodromy purity over \(L\),
including algebraicity. This is the hypothesis in Petrov--Zavyalov's
Definition 5.6.

For completeness, purity descends from \(L\) to \(K\). Write
\(q_0=\#k_K\). For any chosen geometric Frobenius \(\phi\) over \(K\),
some positive power \(\phi^f\) belongs to the open subgroup \(W_L\);
replacing \(L\) by a further finite extension if necessary does not
affect purity, since its Frobenius is a positive power after the finite
inertial action has been killed. The element \(\phi^f\) has degree \(f\)
over \(K\), so its eigenvalues have absolute value \(q_0^{f(i+j)/2}\) on
\(\operatorname{gr}^M_j\). Indeed, over \(L\) the finite inertial action
is trivial and its degree is \(f/[k_L:k_K]\); hence it acts by that
power of an \(L\)-Frobenius. If \(\alpha\) is a \(\phi\)-eigenvalue,
\(\alpha^f\) is such an eigenvalue. Algebraicity descends from
\(\alpha^f\) by the equation \(t^f-\alpha^f=0\), and every complex
absolute value of \(\alpha\) is \(q_0^{(i+j)/2}\). Scaling \(N\) under a
change of monodromy normalization does not change its filtration. This
proves descent; ascent follows by the same power calculation after
killing finite inertia.

Petrov--Zavyalov's Theorem 5.1 now gives étale formality. At \(p\),
Lemma 5.2 gives de Rham formality. The \(p\)-adic pair is obtained from
\(D_{\mathrm{pst}}\); when a semistable Hyodo--Kato description is
available, use its linear Frobenius and \(N\). No equality with ordinary
\(p\)-adic Galois inertia invariants is needed.

Conversely, monodromy purity gives algebraicity and BR-WM-02's first
condition under every complex embedding. That theorem bounds eigenvalues
on the two invariant exterior spaces by \(q^{mi/2}\) and \(q^{-mi/2}\).
A power trace is the sum of the eigenvalue powers, with algebraic
multiplicities, also for a nonsemisimple Frobenius. The triangle
inequality proves (CC-WM.1) with zero exponents and the stated
constants. A zero cohomology group requires no conditions; a zero kernel
admits any positive constant. Every assertion follows. \(\square\)

\textbf{Exact gain and limit.} There are at most \(2\sum_i d_i\) tested
representations; the quantifier over Frobenius powers remains infinite.
This is a criterion and its formality consequence, not a proof of the
bounds for arbitrary \(X\). For \(\ell\ne p\), after unipotent Galois
inertia is achieved, the kernels are actual inertia invariants of the
exterior representations. At \(p\), they belong to the Weil--Deligne or
Hyodo--Kato realization. Alternating traces on total cohomology do not
supply the separate inequalities. Both dual bounds and all complex
embeddings are required. The adopted Dirichlet GRH does not, by itself,
establish these geometric bounds.

\subsection{Exact primary bibliography and version
record}\label{exact-primary-bibliography-and-version-record}

\begin{itemize}
\tightlist
\item
  Binda, F., Kato, H., \& Vezzani, A. (2025). On the \(p\)-adic
  weight-monodromy conjecture for complete intersections in toric
  varieties. \emph{Inventiones Mathematicae, 241}, 559--603.
  \url{https://doi.org/10.1007/s00222-025-01344-x}. Inspected version:
  \url{https://arxiv.org/pdf/2207.00369v2,} dated 10 June 2025, 36 pages;
  locators above refer to this version. Published online 17 June 2025.
\item
  Ito, K., \& Zou, H. (2026). \emph{Semisimplicity, purity and
  Mumford-Tate conjecture for hyper-Kähler varieties}
  (arXiv:2610.07890v1) {[}Preprint{]}.
  \url{https://arxiv.org/abs/2610.07890v1}. Submitted 6 October 2026, 07:38:56
  UTC. Inspected PDF: \url{https://arxiv.org/pdf/2610.07890v1,} 23 pages.
\item
  Petrov, A., \& Zavyalov, B. (2026). \emph{Formality for rigid-analytic
  spaces satisfying the weight-monodromy conjecture}
  (arXiv:2607.14517v2) {[}Preprint{]}.
  \url{https://arxiv.org/abs/2607.14517v2}. First submitted 16 July 2026; v2
  revised 17 August 2026, 21:05:01 UTC. Inspected PDF:
  \url{https://arxiv.org/pdf/2607.14517v2,} 28 pages.
\item
  Scholze, P. (2012). Perfectoid spaces. \emph{Publications
  Mathématiques de l'IHÉS, 116}, 245--313.
  \url{https://doi.org/10.1007/s10240-012-0042-x}. Inspected author version:
  \url{https://arxiv.org/pdf/1111.4914v1,} submitted 21 November 2011,
  16:14:23 UTC, 51 pages; numbered locators above refer to that author
  version.
\item
  Wear, P. (2023). \emph{Perfectoid covers of abelian varieties and the
  weight-monodromy conjecture} (arXiv:2303.05610v1) {[}Preprint{]}.
  \url{https://arxiv.org/abs/2303.05610v1}. Submitted 9 March 2023, 22:41:32
  UTC. Inspected PDF: \url{https://arxiv.org/pdf/2303.05610v1,} 54 pages.
\end{itemize}

% END consequences_criteria_WM_formality.tex

% BEGIN consequences_criteria_stable_curve.tex
\section{Universal stable curves: exact stack purity and
decomposition}\label{universal-stable-curves-exact-stack-purity-and-decomposition}

Research module \nolinkurl{consequences_criteria}; GPT-6 Astra Pro; effort
setting not exposed; 20 November 2026. This module uses established
finite-field weight theory independently of the adopted Alpöge RH and
Dirichlet GRH. New deductions below are supplied with complete proofs.
No claim is made about independent reconstruction of the foundational
BBD or stack six-functor theorems.

\subsection{CC-ST01. Verified stack theorem and its
application}\label{cc-st01.-verified-stack-theorem-and-its-application}

\textbf{Primary theorem.} Sun (2012), arXiv:1009.4398v3, Theorem 1.2,
p.~3, treats proper morphisms of finite diagonal between finite-type
stacks over a finite field with affine stabilizers. For a pure input, it
gives a decomposition into shifted semisimple perverse cohomology
\textbf{after extension to the algebraic closure}. The convention
removing the subscript 0 is explicit in §1.3.3, p.~3. The weight
calculation uses Proposition 3.9(iii), p.~8, and Theorem 3.5, p.~6;
geometric semisimplicity and splitting are Theorems 3.11--3.12, p.~8,
with the application proved in §3.13, p.~9. These are arXiv printed-page
locators, not guessed journal-page translations.
\href{https://arxiv.org/pdf/1009.4398v3}{Primary PDF}.

\textbf{Geometric hypotheses.} Let \(g,n\geq0\) satisfy \(2g-2+n>0\),
let \(q\) be a power of \(p\), and let \(\ell\ne p\). Set \[
Y_0=\overline{\mathcal M}_{g,n,\mathbf F_q},\qquad
X_0=\overline{\mathcal M}_{g,n+1,\mathbf F_q},\qquad
\pi_0:X_0\longrightarrow Y_0,
\] where the last marking is forgotten and the curve stabilized. Put
\(d=3g-3+n\) and \(D=d+1\). Knudsen (1983), Corollary 2.6 and Theorem
2.7, p.~179, identify contraction and stabilization as inverse
constructions and give the smooth proper moduli stacks with
normal-crossing boundary. Knudsen (2012), introduction, pp.~1--2,
explicitly verifies that \(\pi_0\) is representable and is the universal
stable curve, hence is proper and flat with geometrically connected
nodal curve fibres. The local expression \(t_i\mapsto xy\) is also
printed in the proof of Knudsen's Theorem 2.7 on p.~179.
\href{https://journals.msp.org/mscand/article/download/1622/1621/1653}{Original
journal PDF}; \href{https://arxiv.org/pdf/1106.1588v2}{author's later
clarification}.

With \(\Lambda=\overline{\mathbf Q}_\ell\) and a fixed embedding
\(\iota:\Lambda\hookrightarrow\mathbf C\), define \[
K_0=\Lambda_{X_0}[D]=\mathrm{IC}_{X_0},\quad
L_0=R\pi_{0*}K_0,\quad
P_{i,0}={}^p\!H^i(L_0).
\] Then \(L_0\) is pure of weight \(D\), and \(P_{i,0}\) is pure of
weight \(D+i\). After base change to \(\overline{\mathbf F}_q\), there
is a noncanonical isomorphism \[
L\simeq\bigoplus_{i\in\mathbf Z}P_i[-i],
\] and every \(P_i\) is a semisimple perverse sheaf, hence a finite
direct sum of intersection complexes of irreducible lisse sheaves on
smooth locally closed substacks.

\textbf{Complete application proof.} Stability makes all geometric
stabilizers finite, so they are affine. The moduli stacks are of finite
type. Since \(\pi_0\) is representable and proper, its relative diagonal
is a closed immersion and therefore finite. Thus the hypotheses on the
morphism in Sun's theorem hold; arbitrary Artin stacks would not
suffice.

The dimension formula may be checked on the dense smooth-curve locus.
The tangent space of a smooth pointed curve is
\(H^1(C,T_C(-\sum p_i))\). Stability gives \(H^0(C,T_C(-\sum p_i))=0\),
and Riemann--Roch gives its dimension \(3g-3+n\). The normal-crossing
boundary has empty interior, and the stack is smooth, so this is its
dimension on every component. The total space has dimension \(D=d+1\).

The constant sheaf on the smooth stack \(X_0\) has weight zero. With
cohomological shifts, \(K_0\) is perverse and pure of weight \(D\);
equivalently, its unique ordinary cohomology sheaf sits in degree \(-D\)
and its Verdier dual is \(K_0(D)\). Properness gives
\(R\pi_{0!}=R\pi_{0*}\) in this finite-diagonal setting. The upper and
lower weight bounds of Sun's Proposition 3.9(iii) therefore give both
inequalities needed for purity of \(L_0\). His Theorem 3.5 gives the
weights \(D+i\). Finally, Theorems 3.11--3.12 give the asserted
geometric decomposition and semisimplicity. The description of simple
perverse objects is precisely the one used in Sun's Theorem 3.11, citing
Laszlo--Olsson, Theorem 8.2(ii). This proves the stated application.
\(\square\)

\textbf{Scope of the conclusion.} A pure complex need not have every
ordinary stalk cohomology group pure of its cohomological degree. In
particular, the statement about \(R\pi_{0*}\Lambda\) being pure of
weight zero does not assert that the \(H^1\) of every nodal fibre is
pure of weight one. The latter generally has both graph and
normalization contributions. Neither arithmetic semisimplicity nor a
Frobenius-equivariant choice of the displayed splitting follows merely
by retaining Sun's conclusion over \(\mathbf F_q\); the theorem
explicitly takes the geometric base change. Decomposition is not a
smooth-local property that can simply be descended from scheme charts:
Sun gives a counterexample in his introduction, pp.~1--2.

\subsection{CC-ST02. Three perverse degrees and their endpoint
summands}\label{cc-st02.-three-perverse-degrees-and-their-endpoint-summands}

In the preceding situation, the only possibly nonzero \(P_{i,0}\) have
\(i=-1,0,1\). Their weights are respectively \(d,d+1,d+2\).
Geometrically, the endpoints are \[
P_{-1}\simeq\Lambda_Y[d],\qquad
P_1\simeq\Lambda_Y[d](-1).
\] Consequently the constant-coefficient direct image admits a
noncanonical geometric decomposition \[
R\pi_*\Lambda_X\simeq
\Lambda_Y\;\oplus\;P_0[-d-1]\;\oplus\;\Lambda_Y(-1)[-2].
\tag{CC-ST.1}
\] If \(U=\mathcal M_{g,n,\overline{\mathbf F}_q}\subset Y\) is the
smooth-curve locus, then \[
P_0|_U\simeq (R^1\pi_*\Lambda)|_U[d].
\] The middle perverse sheaf can include boundary-supported summands;
equation (CC-ST.1) does not discard them.

\textbf{Complete proof.} Proper base change and the cohomological
dimension of a proper curve give \(R^j\pi_*\Lambda=0\) for
\(j\notin\{0,1,2\}\). Thus the ordinary cohomology of \(L\) is zero
outside degrees \(-d-1,-d,1-d\). For each of these degrees \(k\), the
support has dimension at most \(d\leq1-k\). The support criterion for
the middle perverse \(t\)-structure therefore gives
\(L\in{}^pD^{\leq1}\). The criterion is checked on an étale atlas;
relative dimension zero means no additional stack shift occurs.

Because \(X\) is smooth of dimension \(D\), Verdier duality gives
\(\mathbb D_XK=K(D)\). Proper duality then gives \(\mathbb D_YL=L(D)\).
Perverse duality interchanges \({}^pD^{\leq1}\) and \({}^pD^{\geq-1}\),
and Tate twisting preserves these categories, so \(L\in{}^pD^{\geq-1}\).
Hence \(P_i=0\) for \(|i|>1\). The same vanishing holds before geometric
base change, since that pullback is conservative and perverse exact. The
weight assertions now follow from CC-ST01.

To identify \(P_1\), use its geometric semisimplicity. A simple
constituent with support of dimension \(s\) restricts on a dense smooth
part of that support to a nonzero lisse sheaf shifted by \(s\). Its
contribution through the direct summand \(P_1[-1]\subset L\) therefore
has nonzero ordinary cohomology in degree \(1-s\). But \(L\) has no
ordinary cohomology above \(1-d\). Consequently \(1-s\leq1-d\), so
\(s=d\): every simple constituent of \(P_1\) has full support on a
component of \(Y\).

Over \(U\), smooth proper base change identifies \(P_1|_U\) with
\(R^2\pi_*\Lambda[d]|_U=\Lambda_U[d](-1)\), since the fibres are
geometrically connected smooth proper curves. A semisimple perverse
sheaf whose constituents all have full support is the intermediate
extension of its restriction to a suitable dense open. Since \(Y\) is
smooth and the local system just obtained is constant, this intermediate
extension is \(\Lambda_Y[d](-1)\), component by component. This gives
\(P_1\).

Perverse duality also gives \(\mathbb D_YP_1=P_{-1}(D)\). On the smooth
\(d\)-dimensional stack \(Y\), \[
\mathbb D_Y(\Lambda_Y[d](-1))=\Lambda_Y[d](d+1)=\Lambda_Y[d](D),
\] so \(P_{-1}=\Lambda_Y[d]\). Substitute these two identifications in
\(L\simeq P_{-1}[1]\oplus P_0\oplus P_1[-1]\) and shift by \(-D\) to
obtain (CC-ST.1). The same calculation of perverse degree on \(U\)
identifies the middle term with \(R^1\pi_*\Lambda[d]\). Every asserted
step has now been proved. \(\square\)

\subsection{CC-ST03. Why geometric semisimplicity does not imply
arithmetic
semisimplicity}\label{cc-st03.-why-geometric-semisimplicity-does-not-imply-arithmetic-semisimplicity}

Over \(\operatorname{Spec}\mathbf F_q\), let Frobenius act on
\(\Lambda^2\) by \[
F=\begin{pmatrix}1&1\\0&1\end{pmatrix}.
\] This is a continuous \(\ell\)-adic representation: the homomorphism
\(\widehat{\mathbf Z}\to\mathbf Z_\ell\) followed by
\(a\mapsto\left(\begin{smallmatrix}1&a\\0&1\end{smallmatrix}\right)\)
provides it, with the generator chosen as geometric Frobenius. All
eigenvalues are one, so the associated sheaf is punctually pure of
weight zero. It is not arithmetically semisimple, since its only
Frobenius-fixed line has no invariant complement. After pullback to
\(\operatorname{Spec}\overline{\mathbf F}_q\), the Galois action
disappears and the sheaf is a sum of two constant rank-one sheaves. This
fully explicit example shows why purity alone does not license the
stronger arithmetic conclusion. It does not assert that the particular
universal-curve decomposition has no arithmetic refinement.

\subsection{CC-ST04. Complete geometric support formula, including the
separating
boundary}\label{cc-st04.-complete-geometric-support-formula-including-the-separating-boundary}

Write \(j:U=\mathcal M_{g,n,\overline{\mathbf F}_q}\hookrightarrow Y\)
and \(\mathbb V=(R^1\pi_*\Lambda)|_U\). For each irreducible separating
boundary divisor \(D\subset Y\), let \(D^\circ\) be its open locus of
curves with exactly one node. Let \(\mathcal C_D\to D^\circ\) be the
finite étale two-element cover that chooses one of the two irreducible
components of the fibre. Define the rank-one local system \[
\epsilon_D=\Lambda[\mathcal C_D]/\Lambda\mathbf 1.
\] Here \(\Lambda\mathbf1\) is the diagonal constant sub-local system of
the component permutation local system. For a locally closed immersion
\(k_D:D^\circ\hookrightarrow Y\), the notation \((k_D)_{!*}\) means
intermediate extension to its closure followed by the closed pushforward
to \(Y\).

There is a geometric isomorphism of semisimple perverse sheaves \[
P_0\simeq j_{!*}(\mathbb V[d])\;\oplus\!
\bigoplus_{D\text{ separating}}(k_D)_{!*}\bigl(\epsilon_D[d-1]\bigr)(-1).
\tag{CC-ST.2}
\] Thus substituting (CC-ST.2) into (CC-ST.1) determines every perverse
support and its coefficient system. There are no simple constituents
whose support has codimension at least two, and no constituent supported
on a nonseparating boundary divisor. Intersection complexes in the
formula can nevertheless have nonzero stalks on deeper boundary strata.

The local system \(\epsilon_D\) is the sign of the permutation of the
two components. If those components are distinguished by their genera or
by the labelled markings, it is constant. In the standard labelled
moduli problem, the only boundary type allowing exchange is \(n=0\),
\(g\) even, with genera \(g/2\) and \(g/2\); there one must retain the
component-exchange local system on the stack. A twist of weight two is
essential: the coefficient system on \(D^\circ\) is \(\epsilon_D(-1)\),
so its intermediate extension has perverse weight \(2+(d-1)=d+1\).

\textbf{Complete proof.} A simple constituent of \(P_0\) supported in
dimension \(s\) contributes ordinary cohomology to \(L\) in degree
\(-s\) on a dense smooth part of its support. The upper ordinary degree
bound \(1-d\) from CC-ST02 forces \(-s\leq1-d\), hence \(s\geq d-1\).
Every proper support is therefore a divisor. Since
\(P_0|_U=\mathbb V[d]\) is lisse in one ordinary degree, no such proper
support meets \(U\). The full-support sum is \(j_{!*}(\mathbb V[d])\);
all other possible constituents lie on boundary divisors.

Fix such a divisor and remove its intersections with other boundary
branches and any smaller loci needed to make its restrictions lisse. On
the resulting open set, write \(i:D^\circ\hookrightarrow Y\) for the
closed smooth divisor and \(A=j_{!*}(\mathbb V[d])\). The
intermediate-extension property gives \({}^pH^0(i^*A)=0\). Also
\(i^*A\in{}^pD^{\leq0}\). On the smooth \((d-1)\)-dimensional divisor
the perverse structure on lisse complexes is the ordinary structure
shifted by \(d-1\), so \[
\mathcal H^{1-d}(i^*A)=0.
\] The required intermediate-extension vanishing is stated for stacks in
Laszlo--Olsson (2009), Lemma 6.1, author-preprint p.~10; only this
vanishing implication is used here.

Taking ordinary cohomology in degree \(1-d\) in the geometric
decomposition of \(L\) now gives, on \(D^\circ\), \[
(R^2\pi_*\Lambda)|_{D^\circ}\simeq
\Lambda_{D^\circ}(-1)\oplus E_D,
\tag{CC-ST.3}
\] where \(E_D[d-1]\) is the restriction of the sum of the
divisor-supported constituents of \(P_0\). The first term comes from
\(P_1[-1]\), and \(P_{-1}[1]\) contributes nothing in this degree. No
full-support middle term contributes by the vanishing just proved.

For a connected proper nodal curve, normalization identifies \(H^2\)
with the sum of \(\Lambda(-1)\) over its irreducible components: the
difference between the curve and its normalization is supported at
finitely many nodes and has no cohomology in degree two. This
identification is functorial, so over the one-node locus
\(R^2\pi_*\Lambda\) is exactly the component permutation local system
with a Tate twist. A nonseparating node leaves the curve irreducible,
giving a rank-one system; (CC-ST.3) then forces \(E_D=0\).

A separating node gives precisely two components. Their monodromy
permutation group is either trivial or the order-two group. In the
trivial case, every rank-one quotient by a constant line is constant. In
the order-two case, the constant subrepresentation is the unique
invariant line, and the quotient is the sign representation. Since the
coefficients have characteristic zero, these permutation representations
are semisimple. Equation (CC-ST.3) therefore gives
\(E_D\simeq\epsilon_D(-1)\) in both cases. Intermediate extension
uniquely determines each semisimple divisor-supported sum from this
dense-open restriction, which proves (CC-ST.2).

Finally, a component exchange preserves genus and each labelled marking.
It is impossible when the genera differ or when any labelled marking
lies on either component. If \(n=0\) and the two genera agree, the two
components are unordered and their exchange is exactly the finite
permutation character retained in \(\epsilon_D\). This proves the last
description and completes the proof. \(\square\)

\subsection{References}

Knudsen, F. F. (1983). The projectivity of the moduli space of stable
curves, II: The stacks \(M_{g,n}\). \emph{Mathematica Scandinavica, 52},
161--199. \url{https://doi.org/10.7146/math.scand.a-12001} . Original article,
Corollary 2.6 and Theorem 2.7, p.~179, visually inspected at
\url{https://journals.msp.org/mscand/article/download/1622/1621/1653} .

Knudsen, F. F. (2012). \emph{A closer look at the stacks of stable
pointed curves} (arXiv:1106.1588v2). \url{https://arxiv.org/abs/1106.1588v2} .
Introduction, pp.~1--2, directly inspected. The version inspected
supplies the author's clarification of the contraction/stabilization
construction.

Laszlo, Y., \& Olsson, M. (2009). Perverse \(t\)-structure on Artin
stacks. \emph{Mathematische Zeitschrift, 261}(4), 737--748.
\url{https://doi.org/10.1007/s00209-008-0348-z} . Author preprint, Theorem
8.2(ii), p.~12:
\url{https://www.cmls.polytechnique.fr/perso/laszlo/articleweb/faisceaux-pervers.pdf}
. Statement directly inspected; its whole foundational proof is not
reconstructed here.

Sun, S. (2012). Decomposition theorem for perverse sheaves on Artin
stacks over finite fields. \emph{Duke Mathematical Journal, 161}(12),
2297--2310. \url{https://doi.org/10.1215/00127094-1723657} . Inspected
version: \url{https://arxiv.org/pdf/1009.4398v3} . Exact theorem and page
locators are given above.

% END consequences_criteria_stable_curve.tex

% BEGIN consequences_criteria_RESIDUE.tex
\section{RESIDUE H.3, H.5, H.6: independently established categorical
content}\label{residue-h.3-h.5-h.6-independently-established-categorical-content}

Research module \nolinkurl{consequences_criteria}; GPT-6 Astra Pro; effort
setting not exposed; 20 November 2026. The available file
\nolinkurl{sources/INTAKE_RESIDUE_20261120.md} reports H.3, H.5 and H.6
at statement level and explicitly says their proofs and cited sources
were not opened by that intake. The underlying \nolinkurl{sections/H.tex}
is absent from the supplied source tree. The results below therefore
have their own complete proofs and precise models; they are not
retrospective certifications of unseen proofs.

\subsection{CC-RES01. Direct image for the universal curve is right Kan
extension}\label{cc-res01.-direct-image-for-the-universal-curve-is-right-kan-extension}

Let \(g,n\geq0\) with \(2g-2+n>0\), and work with the complex analytic
topological Deligne--Mumford stacks \[
X=\overline{\mathcal M}_{g,n+1}^{\mathrm{top}},\qquad
Y=\overline{\mathcal M}_{g,n}^{\mathrm{top}},\qquad
\pi:X\longrightarrow Y.
\] Stratify both by stable dual graph. Let \(\mathcal E\) be a compactly
generated presentable \(\infty\)-category, such as spaces, spectra, or
\(D(R)\) for an associative ring \(R\). Write \(\Pi_X\) and \(\Pi_Y\)
for the actual exit-path \(\infty\)-categories of these stratified
stacks. Then:

\begin{enumerate}
\def\labelenumi{\arabic{enumi}.}
\tightlist
\item
  The forgetful map induces an exit functor
  \(F=\Pi(\pi):\Pi_X\to\Pi_Y\).
\item
  The sheaf direct image \(\pi_*\) takes graph-constructible
  \(\mathcal E\)-valued sheaves to graph-constructible sheaves.
\item
  Under exodromy equivalences \(E_X,E_Y\), there is a natural
  equivalence \[
  E_Y\pi_*\simeq\operatorname{Ran}_{F}E_X.
  \tag{CC-RES.1}
  \] In particular, \[
  (E_Y\pi_*\mathcal F)(y)\simeq
  \lim_{(y\downarrow F)} E_X\mathcal F.
  \tag{CC-RES.2}
  \] The indexing category consists of pairs \((x,\alpha:y\to F(x))\).
  It is not replaced by the strict fibre of \(F\) without a further
  cofinality theorem. For \(\mathcal E=D(R)\), \(\pi_*\) denotes the
  derived sheaf direct image.
\end{enumerate}

\textbf{Inspected primary inputs.} Jansen (2023), Corollary 6.6 and
Theorem 6.7, p.~34, identify these exit categories with Charney--Lee
categories and give their exodromy equivalences; Theorem 6.8, p.~35,
gives the derived functor-category formulation. Jansen (2024), Theorem
2.17, pp.~10--11, gives proper base change for étale topological stacks;
Proposition 3.24, p.~21, gives compactly generated coefficients. Haine,
Porta, and Teyssier (2026 version), Theorem 3.2.3, pp.~34--35, states
that every morphism between exodromic stratified \(\infty\)-topoi
induces the requisite exit functor; Observation 2.3.6, p.~20, identifies
pullback with precomposition. Their Proposition 5.2.11, p.~53, supplies
the product formula used below. These are explicit locators in
arXiv:2308.09551v2, arXiv:2308.09550v3, and arXiv:2401.12825v2
respectively.

\textbf{Complete proof.} Forgetting the last marked point induces a map
of dual-graph strata: remove its leg and stabilize every now-unstable
rational vertex. Contraction of edges, corresponding to smoothing nodes,
is compatible with this operation because stable contraction of families
commutes with base change. Thus \(\pi\) is a stratified map. Jansen's
theorem says that both stratified stacks are exodromic. The
functoriality theorem of Haine--Porta--Teyssier gives \(F\), and its
pullback compatibility gives \[
E_X\pi^*\simeq F^*E_Y,
\] first for space-valued sheaves and then, by the compactly generated
coefficient construction, for \(\mathcal E\)-valued sheaves.

We prove the constructibility claim, which is necessary to identify the
actual sheaf direct image with the right adjoint inside constructible
sheaves. Fix a dual-graph stratum \(Y_G\). On a sufficiently small
contractible analytic chart \(B\to Y_G\), label the finitely many
components, nodes, branches and old marked points. The family over \(B\)
has constant nodal type. Its normalization is a disjoint union of proper
smooth families of compact surfaces, with disjoint labelled sections for
the old marked points and the two preimages of every node. One can
trivialize these families preserving all those sections: choose a
horizontal distribution tangent to the section images, using disjoint
neighbourhoods of the sections and a partition of unity; lift the radial
paths in a small coordinate ball \(B\). Properness makes the lifts exist
throughout these paths after shrinking the ball. They give smooth
trivializations depending continuously on the endpoint. Gluing each pair
of node sections then gives a fibre-preserving homeomorphism \[
\pi^{-1}(B)\simeq C\times B,
\tag{CC-RES.3}
\] where \(C\) is a fixed compact nodal curve, with its nodes and old
marked points distinguished.

This is a stratified trivialization for a refinement of the induced
graph stratification: an extra point in a component's smooth unmarked
part adds a leg to that vertex; at an old marked point it inserts a
rational tail; at a node it inserts a rational bridge. Each of these
conditions is preserved by the trivialization. Thus a
graph-constructible sheaf on \(X\) restricts to a constructible sheaf
for this finite point-and-component stratification on \(C\times B\).

Both \(C\) with this stratification and the ball \(B\) are conically
stratified. The exit product formula gives \[
\Pi(C\times B)\simeq\Pi(C)\times\Pi(B)\simeq\Pi(C),
\] since \(B\) is contractible. Exodromy therefore shows that every such
constructible sheaf on \(C\times B\) is pulled back from a constructible
sheaf \(\mathcal G\) on \(C\). These conical, finite-dimensional spaces
use the usual constructible sheaf categories; the equivalent hypersheaf
formulation of the product theorem yields the same objects here.

The map \(C\to *\) is proper. Proper base change for the product square
now identifies the direct image to \(B\) of the pullback of
\(\mathcal G\) with the constant sheaf having value
\(\Gamma(C,\mathcal G)\). Proper base change for \(\pi\) along
\(Y_G\hookrightarrow Y\), followed by the chart \(B\to Y_G\), identifies
this with \((\pi_*\mathcal F)|_B\). The base-change theorem holds also
for compactly generated coefficients: express \(\mathcal E\)-valued
sheaves through finite-limit-preserving functors from the opposite of the category of compact objects to
space-valued sheaves; evaluation detects the required equivalence,
reducing it to the space-valued statement. Hence \(\pi_*\mathcal F\) is
locally constant on each graph stratum. This proves (2).

It remains to identify the functor. For every
\(A\in\operatorname{Fun}(\Pi_Y,\mathcal E)\) and
\(\mathcal F\in\operatorname{Shv}_{\mathrm{cbl}}(X;\mathcal E)\),
exodromy and the sheaf adjunction give naturally \[
\begin{aligned}
\operatorname{Map}(A,E_Y\pi_*\mathcal F)
&\simeq\operatorname{Map}(E_Y^{-1}A,\pi_*\mathcal F)\\
&\simeq\operatorname{Map}(\pi^*E_Y^{-1}A,\mathcal F)\\
&\simeq\operatorname{Map}(F^*A,E_X\mathcal F).
\end{aligned}
\] Thus \(E_Y\pi_*E_X^{-1}\) is right adjoint to precomposition by
\(F\). By uniqueness of right adjoints, it is right Kan extension. Its
usual pointwise limit formula gives (CC-RES.2). This proves all
assertions. \(\square\)

\textbf{Relation to the intake.} This independently establishes a
precise geometric version of H.6 for the canonical exit categories.
Jansen identifies \(\Pi_{g,n}\) with
\(\mathrm{SC}_{g,n}^{\mathrm{op}}\); consequently the statement can be
expressed using presheaves on entrance categories, with right Kan
extension along the opposite \(u^{\mathrm{op}}\) of the entrance functor
\(u\). What remains uninspected is the equality of the unseen archive's
chosen, explicitly presented forgetful functor with this canonical
functor. The assertion is complex topological. No Frobenius action, Tate
twist, or finite-field weight filtration is created by exodromy; the
separate finite-field statement is CC-ST01--CC-ST04.

\subsection{CC-RES02. Complete lifting-category calculation for the
nodal local
model}\label{cc-res02.-complete-lifting-category-calculation-for-the-nodal-local-model}

Let \(\mathcal C\) be the ordinary category with objects \(0,1\), with
\[
\operatorname{End}(0)=\{1_0\},\quad
\operatorname{End}(1)=\mathbf Z,\quad
\operatorname{Hom}(0,1)=\{a\},\quad
\operatorname{Hom}(1,0)=\varnothing,
\] where composition in \(\mathbf Z\) is addition and \(m\circ a=a\).
Define \[
F:\mathcal C\times\mathcal C\longrightarrow\mathcal C
\] by \(F(1,1)=1\) and \(F(x)=0\) on the other three objects. On the
endomorphisms of \((1,1)\), set \(F(r,s)=r+s\); all arrows between
objects mapped to 0 go to \(1_0\), and all arrows from those objects to
\((1,1)\) go to \(a\).

For \(x\in\mathcal C^2\) and \(\alpha:F(x)\to y\), define the strict
lifting category \(\mathcal L(x,\alpha)\) as follows. Its objects are
arrows \(\beta:x\to z\) such that \(F(z)=y\) and \(F(\beta)=\alpha\). A
morphism from \(\beta\) to \(\beta'\) is an arrow \(\delta:z\to z'\)
with \(F(\delta)=1_y\) and \(\delta\beta=\beta'\). Then \[
|N\mathcal L(x,\alpha)|\simeq
\begin{cases}
S^1,&x=(0,0),\ \alpha=a:0\to1,\\
*,&\text{all other possible pairs }(x,\alpha),
\end{cases}
\tag{CC-RES.4}
\] In the exceptional case the lifting category itself is the one-object
groupoid \(B\mathbf Z\). In the remaining groupoid cases it is
equivalent to the terminal category; the cases with an initial object
need only have contractible nerve.

\textbf{Complete proof.} If \(F(x)=0\) and \(\alpha=1_0\), the lifting
category is the undercategory of \(x\) in the full subcategory
\(F^{-1}(0)\). The identity arrow of \(x\) is initial, so its nerve is
contractible.

Suppose \(F(x)=0\) and \(\alpha=a\). The target must be \(z=(1,1)\). Its
endomorphisms over \(1_1=0\in\mathbf Z\) form \[
K=\ker(\mathbf Z^2\xrightarrow{+}\mathbf Z)
 =\{(r,-r):r\in\mathbf Z\}\simeq\mathbf Z.
\] For \(x=(0,0)\) there is exactly one arrow \(\beta:x\to(1,1)\), and
every element of \(K\) fixes it. Hence the lifting category is exactly
the one-object groupoid \(BK\).

For \(x=(1,0)\) the possible arrows \(\beta\) are indexed by
\(\mathbf Z\), coming from the first coordinate. The element \((r,-r)\)
acts by translation by \(r\). This action is free and transitive, so its
action groupoid is equivalent to the terminal category. For \(x=(0,1)\)
the same argument uses translation by \(-r\) in the second coordinate.

Finally, if \(F(x)=1\), then \(x=(1,1)\), \(y=1\), and \(\alpha\) is an
arbitrary integer \(m\). Its lifts are the pairs \((r,s)\) with
\(r+s=m\). The group \(K\) acts freely and transitively on this set by
addition, so the lifting category again is equivalent to the terminal
category. No arrow \(1\to0\) exists, so these cases exhaust every
possible \((x,\alpha)\). This proves (CC-RES.4). \(\square\)

\textbf{Geometric interpretation with its precise limit.} The exit
category of a complex disk stratified by its origin and punctured disk
is \(\mathcal C\): the punctured stratum has fundamental group
\(\mathbf Z\) and the origin is the cone's initial object. The product
stratification of a bidisk has exit category \(\mathcal C^2\). The
multiplication map \((u,v)\mapsto uv\) induces the displayed functor
because the winding numbers add on \((\mathbf C^*)^2\). Knudsen's
universal-curve local model is exactly \(t=uv\) (1983, proof of Theorem
2.7, p.~179); the analytic nodal coordinates are also used explicitly in
Jansen (2023), Appendix B.2, pp.~43--44. At \((u,v)=(0,0)\) the extra
marked point lies at the node of the contracted curve; stabilization
inserts a rational bridge. Thus the single \(B\mathbf Z\) case is
precisely the local bridge-contraction followed by smoothing described
in the intake. This verifies that local mechanism without asserting a
global equivalence of all the archive's lifting categories with these
local ones.

\subsection{CC-RES03. Exact remaining scope of
H.5}\label{cc-res03.-exact-remaining-scope-of-h.5}

The independent results above resolve the actual sheaf/right-Kan
identification and compute the exceptional nodal lifting category in a
fully specified model. They do not by themselves prove that an arbitrary
categorical base change of the archive's entrance functor preserves
\(W_\infty\)-equivalences, nor its claimed global smoothness
classification. Properness of \(\pi\) supplies geometric proper base
change; that statement quantifies over geometric base changes and cannot
be silently replaced by the universal quantifier over functors in a
definition of \(W_\infty\)-properness. Likewise, a noncontractible local
lifting category does not exclude a possible contraction caused by
additional global morphisms unless the local-to-global comparison is
proved. The missing input is the archive's full functor, the definitions
it uses, and its comparison of global lifting categories with the local
models. H.5 remains an uninspected archive assertion; no contradiction
has been established.

\subsection{References}

Haine, P. J., Porta, M., \& Teyssier, J.-B. (2026). \emph{Exodromy
beyond conicality} (arXiv:2401.12825v2; revised 24 June 2026).
\url{https://arxiv.org/abs/2401.12825v2} . Inspected: Observation 2.3.1,
p.~19; Observation 2.3.6, p.~20; Theorem 3.2.3 and its proof,
pp.~34--35; Proposition 5.2.11, p.~53.
\url{https://arxiv.org/pdf/2401.12825v2} .

Jansen, M. Ø. (2023). \emph{Moduli stack of stable curves from a
stratified homotopy viewpoint} (arXiv:2308.09551v2).
\url{https://arxiv.org/abs/2308.09551v2} . Inspected: Corollary 6.6 and
Theorem 6.7, p.~34; Theorem 6.8, p.~35; Appendix B, local coordinates
and trivializations, pp.~40--45. \url{https://arxiv.org/pdf/2308.09551v2} .

Jansen, M. Ø. (2024). Stratified homotopy theory of topological
\(\infty\)-stacks: A toolbox. \emph{Journal of Pure and Applied Algebra,
228}(11), Article 107710. \url{https://doi.org/10.1016/j.jpaa.2024.107710} .
Inspected preprint arXiv:2308.09550v3: Theorem 2.17 and proof,
pp.~10--11; Definition 3.12 and following discussion, pp.~16--17; Remark
3.15, p.~17; Proposition 3.24, p.~21. \url{https://arxiv.org/pdf/2308.09550v3}
.

Knudsen, F. F. (1983). The projectivity of the moduli space of stable
curves, II: The stacks \(M_{g,n}\). \emph{Mathematica Scandinavica, 52},
161--199. \url{https://doi.org/10.7146/math.scand.a-12001} . Theorem 2.7 and
its displayed formal local map, p.~179, visually inspected:
\url{https://journals.msp.org/mscand/article/download/1622/1621/1653} .

% END consequences_criteria_RESIDUE.tex

% BEGIN root_linear_archive.tex
\section{Cote 134-8: the linear lemmas, normalization, and a proved range of the proposed theorem}

\subsection{The recovered question and the scope of its predicates}
The primary source is Grothendieck's \emph{Pursuing stacks}, cote
134-8, manuscript footers 47--54, in the official Montpellier PDF
\url{https://grothendieck.umontpellier.fr/134-8.pdf}.  This PDF has a
cover: its physical page number is the manuscript footer plus one.
The displayed \emph{Thm(?)} is on footer 48, physical page 49.
The original uses a connective chain-complex target, denoted with a
subscript on \(D\); we do not replace it by the unbounded derived
category.  The statements proved below use the unambiguous modern
notation \(D_{\geq0}(k)\), with homological differentials of degree
\(-1\).

On footer 47 the scalar functor \(\kappa(U)=\mathbf1\otimes_k U\)
is required to commute with inverse limits and is given a left adjoint
\(\varphi\).  The separately mentioned right adjoint of \(\kappa\)
is \(\operatorname{Hom}_{\mathcal M}(\mathbf1,-)\).  The setup treats
\(\kappa\) as exact.  There are enough projectives, the canonical
map \(L\varphi(\mathbf1)\to k\) is an isomorphism, and a monomorphism
\(\alpha:\mathbf1\to J\) is sought with
\[
 X\longrightarrow J\otimes X\text{ monic},\qquad
 L\varphi(J\otimes X)=0\quad(X\in\mathcal M).
 \tag{134.1}
\]
The vanishing includes degree zero.  The proposed conclusion is that
localizing \(\mathcal M\) at the maps inverted by \(L\varphi\)
gives the connective derived category of scalar modules.  Footer 48
then adds projective generators of finite presentation and sketches a
chain-complex functor and its right adjoint.  Its right margin also
records a projectivity condition on the scalar modules of morphisms.
We state this extra condition explicitly in BR-13405.

The supplied DER\_LOCALIZERS intake reports results C.1--C.29,
including a counterexample outside the finitely presented-generator
case.  Their original proofs are absent.  The present proofs therefore
establish the claims they state directly; they do not certify the
unseen counterexample or silently identify its hypotheses with these
handwritten predicates.  No RH or GRH input is needed for this section.

\subsection{BR-13401. The weakened form of Lemme 1 is valid}
Let \(\mathcal M,\mathcal A\) be abelian categories, with enough
projectives in \(\mathcal M\), and let
\(\Lambda:\mathcal M\to\operatorname{Ch}_{\geq0}(\mathcal A)\)
be an exact additive functor.  Put \(\varphi=H_0\Lambda\).
Assume that, for every projective \(P\in\mathcal M\),
\[
 H_j(\Lambda P)=0\quad(j>0).
 \tag{134.2}
\]
Then \(\varphi\) is right exact and there is a natural isomorphism
\[
 L\varphi(X)\simeq\Lambda X
 \quad\text{in }D_{\geq0}(\mathcal A),\qquad X\in\mathcal M.
 \tag{134.3}
\]
It is unnecessary to require \(\Lambda P\) to be a projective object
of the chain-complex category.  This is the weakened condition written
beside Lemme 1 on manuscript footer 53, physical page 54.

\paragraph{Complete proof.}
Exactness of \(\Lambda\) and right exactness of \(H_0\) imply right
exactness of \(\varphi\).  Choose a projective resolution
\(P_\bullet\to X\) in nonnegative homological degrees.  Form the
first-quadrant double complex
\(B_{a,b}=(\Lambda P_a)_b\), with the usual sign in one
differential so that the differentials anticommute.  Its total complex
uses only the finite direct sum \(\bigoplus_{a+b=n}B_{a,b}\) in each
degree \(n\).

The augmentations \(P_\bullet\to X\) give a map
\(\operatorname{Tot}B\to\Lambda X\).  In every fixed internal degree
\(b\), its augmented horizontal complex is exact, since the degree
\(b\) functor of \(\Lambda\) is exact.  The first-quadrant filtration
therefore proves this map is a quasi-isomorphism.  In the other
direction, the natural augmentations
\(\Lambda P_a\to H_0(\Lambda P_a)[0]\) give a map
\(\operatorname{Tot}B\to\varphi(P_\bullet)\).  Its vertical homology
is an isomorphism in degree zero and zero in positive degree by
(134.2).  The other first-quadrant filtration proves that it too is a
quasi-isomorphism.  These assertions can equivalently be checked by
successive filtration quotients in each total degree; no convergence
of an infinite diagonal is being assumed.  The complex
\(\varphi(P_\bullet)\) computes \(L\varphi(X)\), giving (134.3).
Comparison maps between projective resolutions are unique up to chain
homotopy; applying the additive functor and totalization preserves
these homotopies.  The resulting isomorphism is consequently natural
and independent of the resolution. \(\square\)

\subsection{BR-13402. Projective objects of the connective chain category}
Suppose \(\mathcal A\) is abelian with enough projectives and
coproducts, and coproducts of projectives are projective.  For
\(n\geq1\), let \(D^n(Q)\) be the complex with \(Q\) in degrees
\(n,n-1\), identity differential, and zero elsewhere.  A complex
\(P\) is projective as an object of
\(\operatorname{Ch}_{\geq0}(\mathcal A)\) if and only if it is
isomorphic to
\[
 U[0]\oplus\bigoplus_{n\geq1}D^n(Q_{n-1}),
 \tag{134.4}
\]
where \(U,Q_j\) are projective in \(\mathcal A\).  The sum here is
locally finite in the chain degree.  In particular, a projective
complex is not the same as an arbitrary complex with projective terms.

\paragraph{Complete proof.}
There are natural identifications
\[
 \operatorname{Hom}_{\mathrm{Ch}}(Q[0],C)
   =\operatorname{Hom}_{\mathcal A}(Q,C_0),\qquad
 \operatorname{Hom}_{\mathrm{Ch}}(D^nQ,C)
   =\operatorname{Hom}_{\mathcal A}(Q,C_n).
\]
Consequently the indicated building blocks are projective.  Their
coproducts are projective, proving one implication.

For the converse, choose projective epimorphisms \(Q_n\to P_n\)
for all \(n\geq0\).  The preceding identifications give a degreewise
epimorphism
\[
 Q_0[0]\oplus\bigoplus_{n\geq1}D^nQ_n\longrightarrow P.
\]
It splits because \(P\) is projective.  Thus the terms of \(P\)
are projective, \(H_jP=0\) for \(j>0\), and \(H_0P\) is
projective, all by taking the corresponding direct summands of the
source.  Write \(B_j=\operatorname{im}(d:P_{j+1}\to P_j)\).
The sequence \(0\to B_0\to P_0\to H_0P\to0\) splits, so
\(B_0\) is projective.  Inductively,
\(0\to B_j\to P_j\to B_{j-1}\to0\) splits for every \(j\geq1\),
because \(H_jP=0\) and \(B_{j-1}\) is projective.  These splittings
identify \(P_0\) with \(H_0P\oplus B_0\) and \(P_j\) with
\(B_j\oplus B_{j-1}\), with differential the identity on the second
factor and zero on the first.  This is (134.4), with \(U=H_0P\).
\(\square\)

\subsection{BR-13403. Unit normalization and the exact remaining counit test}
Let \(k\) be a commutative ring.  Assume the following precise data:
\(\mathcal M\) is a cocomplete \(k\)-linear abelian category with
enough projectives; \(\kappa(U)=\mathbf1\otimes_k U\) is an exact
scalar functor with left adjoint \(\varphi\); and there is an exact
\(k\)-linear adjunction
\[
 \Lambda:\mathcal M\rightleftarrows
 \operatorname{Ch}_{\geq0}(k):G
 \tag{134.5}
\]
with both functors exact, \(H_0\Lambda=\varphi\), and (134.2).
Assume also that the canonical counit
\(\Lambda G(k[0])\to k[0]\) is a quasi-isomorphism.  Equivalently,
under the identification below, this is the canonical normalization
\(L\varphi(\mathbf1)\to k\), including degree zero.

Set \(W=\Lambda^{-1}(\mathrm{qis})\).  Then (134.5) induces
mutually inverse equivalences on the localizations, with the induced
unit and counit, if and only if
\[
 L\varphi(GP)=0
 \quad\text{for every acyclic projective object }
 P\in\operatorname{Ch}_{\geq0}(k).
 \tag{134.6}
\]
Thus (134.6) is an exact residual test for this specified adjunction,
not merely a sufficient list of examples.  It is the question left at
the bottom of manuscript footer 54, physical page 55.

\paragraph{Complete proof.}
For every scalar module \(U\) and every \(X\in\mathcal M\),
the adjunctions give natural isomorphisms
\[
 \operatorname{Hom}_{\mathcal M}(X,G(U[0]))
 =\operatorname{Hom}_{\mathrm{Ch}}(\Lambda X,U[0])
 =\operatorname{Hom}_k(\varphi X,U)
 =\operatorname{Hom}_{\mathcal M}(X,\kappa U).
\]
Yoneda identifies \(G(U[0])\) with \(\kappa U\), and identifies
the two counits.  If \(U\) is projective, it is a summand of a free
module.  Both \(\kappa\) and \(\Lambda\) preserve coproducts:
the first by its scalar-tensor construction, the second because it is
a left adjoint.  Therefore the normalization for \(k\) implies that
\(\Lambda G(U[0])\to U[0]\) is a quasi-isomorphism for every
projective \(U\), by coproducts and retracts.  The coproduct of these
quasi-isomorphisms is a quasi-isomorphism in modules, since direct
sums are exact.

Every projective chain object splits by BR-13402 as \(U[0]\oplus P\)
with \(U\) projective and \(P\) acyclic projective.  Additivity of
both adjoints means that its counit is a quasi-isomorphism exactly when
the counit on \(P\) is one.  Since \(P\) is acyclic, the latter
condition is precisely \(\Lambda GP\simeq0\), equivalently (134.6)
by BR-13401.  The exact-adjunction criterion F-SM17 proves that counits
on all projective chain objects are quasi-isomorphisms if and only if
all counits are, and that this is equivalent to the stated localized
adjunction equivalence.  This proves both implications. \(\square\)

\subsection{BR-13404. The proposed equivalence holds over fields and over \(\mathbf Z\)}
Retain all the data of BR-13403.  Assume \(\mathcal M\) has enough
injectives and the tensor interval (134.1) exists.  Suppose there is
an integer \(r\geq0\) such that every projective \(k\)-module has
injective dimension at most \(r\).  Then (134.6) holds and the exact
adjunction induces an equivalence
\[
 W^{-1}\mathcal M\simeq D_{\geq0}(k).
 \tag{134.7}
\]
In particular this conclusion holds for a field \(k\) and for
\(k=\mathbf Z\), once the specified exact adjunction has been
constructed.  These are the rational and integral scalar rings used
in the programme.  The result does not assert an unbounded derived
equivalence or eliminate the construction hypotheses of BR-13403.

\paragraph{Complete proof.}
If \(I\) is injective in \(\mathcal M\), its monomorphism
\(I\to J\otimes I\) splits.  Applying \(L\varphi\) and (134.1)
shows \(L\varphi(I)=0\), since it is a retract of zero.  Any object
\(X\) having a finite injective resolution consequently satisfies
\(L\varphi(X)=0\).  Indeed, in a short exact sequence
\(0\to X\to I\to X'\to0\) with \(I\) injective, the triangle for
the left derived functor gives
\(L\varphi(X')\simeq L\varphi(X)[1]\).  Starting at the final
injective term of a finite resolution and working backwards proves
the assertion, including degree-zero vanishing.

Let \(P\) be an acyclic projective chain object.  BR-13402 gives
\(P=\bigoplus_{n\geq1}D^n(Q_{n-1})\) for projective scalar
modules \(Q_j\).  Choose length-\(r\) injective resolutions of the
\(Q_j\).  Apply each disk functor and take the locally finite sum
over \(j\).  In each chain degree only two of these resolutions
occur, so the resulting sequence is exact; it is a length-\(r\)
resolution of \(P\).

Each of its terms is injective as a chain object.  To see this without
an unsupported assertion about infinite sums of injectives, note that
\(D^{j+1}(I)\) is right adjoint, in the appropriate variable, to
the exact evaluation \(C\mapsto C_j\):
\[
 \operatorname{Hom}_{\mathrm{Ch}}(C,D^{j+1}I)
    =\operatorname{Hom}_k(C_j,I).
\]
It is therefore injective if \(I\) is injective.  The locally finite
sum of these disks equals their product degree by degree.  Products
of injective objects are injective, since maps into the product are
products of maps and products of surjections of abelian groups are
surjective.  This proves the asserted injectivity.

Because \(\Lambda\) is exact, its right adjoint \(G\) sends
injectives to injectives.  Because \(G\) is exact, applying it to
the length-\(r\) resolution gives a finite injective resolution of
\(GP\) in \(\mathcal M\).  The first paragraph now gives
\(L\varphi(GP)=0\).  BR-13403 proves (134.7).

Over a field, every module is injective, so \(r=0\) works.  Over
\(\mathbf Z\), every free module \(F\) has
\(0\to F\to F\otimes\mathbf Q\to(F\otimes\mathbf Q)/F\to0\),
whose two latter groups are divisible and hence injective.  Every
projective abelian group is a summand of a free one; taking summands
shows injective dimension at most one.  For completeness, the
divisibility assertion uses the elementary injectivity criterion:
a homomorphism \(n\mathbf Z\to D\) extends to \(\mathbf Z\to D\)
exactly when the prescribed image of \(n\) is divisible by \(n\).
The usual ideal criterion for injectivity then applies.  Thus \(r=1\)
works over \(\mathbf Z\). \(\square\)

\subsection{BR-13405. An explicit construction of the exact adjunction in the corollary's locally projective case}
Let \(R\) be a small \(k\)-linear category such that every
\(R(x,y)\) is a projective \(k\)-module, and let
\(\mathcal M=\operatorname{Fun}_k(R^{\mathrm{op}},\operatorname{Mod}k)\).
Assume \(\kappa(U)=T\otimes_k U\), for a fixed right \(R\)-module
\(T\), is exact and has a left adjoint \(\varphi\).  Then there
is an exact \(k\)-linear functor \(\Lambda\) representing
\(L\varphi\) on objects and an exact right adjoint \(G\), as in
(134.5).  This supplies a fully specified construction when the
finite-presentation projective generators and the additional scalar
projectivity condition on footer 48 are interpreted by such a ringoid.

\paragraph{Complete proof.}
Write \(P_x=R(-,x)\) for the right representable and set
\(D(x)=\varphi(P_x)\).  The dependence on \(x\) makes \(D\) a
left \(R\)-module.  Since \(\kappa\) is exact, its left adjoint
\(\varphi\) sends projectives to projectives, so each \(D(x)\)
is a projective scalar module.  In fact it is finitely generated:
\[
 \operatorname{Hom}_k(D(x),U)
   =\operatorname{Hom}_{\mathcal M}(P_x,\kappa U)
   =T(x)\otimes_k U
\]
commutes with arbitrary coproducts.  A projective module with this
property is finitely generated.  To verify this last assertion, write
it as a summand of \(k^{(I)}\).  Its inclusion in that direct sum
has, by the coproduct property, only finitely many nonzero coordinate
maps, so it is a summand of a finite free module.

There is a natural isomorphism \(\varphi(X)=D\otimes_R X\).
Both sides preserve coproducts and cokernels, and they agree naturally
on each representable.  Every right module has a presentation by
coproducts of representables; applying the two functors to such a
presentation proves the isomorphism and its naturality.

Take a projective resolution \(B_\bullet\to D\) by left
\(R\)-modules, choosing each \(B_a\) to be a coproduct of left
representables \(R(x,-)\).  Put
\[
 \Lambda(X)=B_\bullet\otimes_R X.
 \tag{134.8}
\]
Every degree of this functor is a sum of evaluation functors, and is
therefore exact; hence \(\Lambda\) is exact.  Its zeroth homology
is \(D\otimes_R X=\varphi X\).  For \(X=P_x\), the complex is
\(B_\bullet(x)\), whose positive homology vanishes and whose zeroth
homology is \(D(x)\).  Coproducts and retracts give the same
positive-homology vanishing for every projective \(X\).
BR-13401 proves that (134.8) represents \(L\varphi\).

Define a right \(R\)-module by
\[
 (GC)(x)=\operatorname{Hom}_{\operatorname{Ch}_{\geq0}(k)}
                    (B_\bullet(x),C).
 \tag{134.9}
\]
The ordinary tensor--Hom adjunction, degree by degree with the chain
map compatibility, gives \(\Lambda\dashv G\).  Every term
\(B_a(x)\) is a coproduct of projective scalar modules \(R(y,x)\),
and is projective.  The complex \(B_\bullet(x)\) resolves the
projective module \(D(x)\).  Its successive boundary sequences
therefore split, exactly as in BR-13402; it is a projective object of
the chain category.  It follows that (134.9) is an exact functor of
\(C\), for every \(x\).  Exactness in a functor category is
objectwise, so \(G\) is exact.  This completes the construction.
\(\square\)

\subsection{BR-13406. Positive-degree interval acyclicity is insufficient}
The degree-zero clause in (134.1) cannot be removed.  More precisely,
replace its vanishing by
\(L_i\varphi(J\otimes X)=0\) only for \(i>0\), keeping the scalar
unit normalization and all the exact-adjunction data.  The proposed
canonical equivalence need no longer hold.

\paragraph{Complete counterexample.}
Take a nonzero field \(k\), \(\mathcal M=\operatorname{Mod}k\)
with its ordinary tensor product, \(\kappa=\varphi=\mathrm{id}\),
and \(J=\mathbf1=k\), \(\alpha=\mathrm{id}\).  The interval map
\(X\to J\otimes X\) is the identity, all positive derived functors
of \(\varphi\) vanish, and \(L\varphi(k)=k\).  Let
\(\Lambda(X)=X[0]\) and \(G(C)=C_0\).  These are exact adjoint
functors, since a chain map from \(X[0]\) into a nonnegative complex
is exactly a module map \(X\to C_0\).  They satisfy (134.2).
The class \(W\) consists precisely of isomorphisms of modules, so
the canonical localized functor still has only complexes concentrated
in degree zero in its essential image.  The object \(k[1]\) of
\(D_{\geq0}(k)\) has nonzero first homology and cannot lie in that
image.  Thus the asserted canonical equivalence fails.  Exactly the
missing degree-zero condition fails, since
\(\varphi(J\otimes X)=X\) need not be zero. \(\square\)

\paragraph{What is settled and what remains.}
The original Lemme 1 is proved under its weaker, explicitly annotated
hypothesis.  F-SM17 proves the exact-adjunction/counit lemma without an
infinite-coproduct assumption on its right adjoint.  BR-13403 identifies
the residual acyclic-projective test, and BR-13404 discharges it over
fields and \(\mathbf Z\) under the specified exact-adjunction data.
BR-13405 constructs those data in the locally projective ringoid case.
For a general scalar ring without the finite-injective-dimension
condition, (134.6) remains the precise test for the raw adjunction.
An equivalence constructed by a different derived or model-category
adjunction would require its own comparison with (134.5).  The much
stronger assertion that every ringoid presentation automatically
satisfies the scalar projectivity hypothesis has not been made.


% END root_linear_archive.tex

% BEGIN source_cote1348_counit.tex
\section{Cote 134-8, Lemma 2: a complete counit criterion}
\label{sec:SM-cote1348-counit}

The manuscript at handwritten pages 52--53 considers an adjunction from
an abelian category to a category of complexes, the inverse image of
quasi-isomorphisms, and tests on the adjunction counit.  Its conditions
labelled $(ii')$ and $(ii'')$ have revisions that should not be
silently identified with a modern assertion.  The following gives a
precise version of the clear counit/projective-object criterion.  All
complexes here are homological and concentrated in degrees at least
zero.  The exactness hypothesis on the right adjoint is explicit.

\paragraph{Finding F-SM-17 (localization equivalence and projective counit tests).}
Let $\mathcal A$ and $\mathcal M$ be abelian categories, with
$\mathcal A$ having enough projectives.  Put
$\mathcal B=\operatorname{Ch}_{\geq0}(\mathcal A)$ and let
\[
 \Lambda:\mathcal M\rightleftarrows\mathcal B:G
 \qquad(\Lambda\dashv G)
\]
be an additive adjunction such that both functors are exact.  Write
$\eta$ and $\varepsilon$ for its unit and counit, let
$\mathcal Q$ be the quasi-isomorphisms in $\mathcal B$, and put
$W=\Lambda^{-1}(\mathcal Q)$.  Fix universes in which the following
localizations exist.  Then these conditions are equivalent:
\begin{enumerate}
\item For every $C\in\mathcal B$, the counit
      $\varepsilon_C:\Lambda G(C)\to C$ is a quasi-isomorphism.
\item The same assertion holds for every projective object $P$ of
      the abelian category $\mathcal B$.
\item $G$ sends $\mathcal Q$ into $W$, and the induced functor
      \[
       \overline\Lambda:\mathcal M[W^{-1}]
          \longrightarrow \mathcal B[\mathcal Q^{-1}]
          =D_{\geq0}(\mathcal A)
      \]
      is an equivalence.
\end{enumerate}
Whenever these hold, $G^{-1}(W)=\mathcal Q$, the induced functors
$\overline\Lambda\dashv\overline G$ are mutually inverse equivalences,
and every $\eta_X$ belongs to $W$.

The equivalence of the first and third conditions does not require
abelian categories or exactness; it uses the displayed adjunction and
the saturation of quasi-isomorphisms.  For the implication from the
second condition to the first, exactness of the composite
$H=\Lambda G$ suffices in place of separate exactness of both functors.
No preservation of infinite coproducts by $G$ is assumed.

\begin{proof}
\emph{Enough projectives in the category of complexes.}
For a projective $P\in\mathcal A$, let $S^0(P)$ be $P$ in degree zero.
For $n\geq1$, let $D^n(P)$ have $P$ in degrees $n$ and $n-1$,
with the identity as its only nonzero differential.  There are natural
identifications
\[
 \operatorname{Hom}_{\mathcal B}(S^0(P),C)
       =\operatorname{Hom}_{\mathcal A}(P,C_0),
 \qquad
 \operatorname{Hom}_{\mathcal B}(D^n(P),C)
       =\operatorname{Hom}_{\mathcal A}(P,C_n).
\]
An epimorphism of complexes is an epimorphism in every degree.
Consequently these stalk and disk complexes are projective objects of
$\mathcal B$.

For any $C$, choose epimorphisms $P_n\to C_n$ with each $P_n$
projective.  The associated chain maps give an epimorphism
\[
 Q=S^0(P_0)\oplus\bigoplus_{n\geq1}D^n(P_n)\longrightarrow C.
\]
This coproduct exists without any assumption about arbitrary
coproducts in $\mathcal A$: in degree $j$ its value is the finite
biproduct $P_j\oplus P_{j+1}$.  A map from $Q$ lifts across any
epimorphism of complexes by lifting separately its maps from the
$P_n$ into degree $n$.  Hence $Q$ is projective.  Iterating on kernels
gives a projective resolution of every object of $\mathcal B$.

\emph{The finite-diagonal total-complex argument.}
We record the elementary convergence fact used below.  For a first
quadrant bicomplex $X_{r,s}$ with commuting differentials $d_h,d_v$,
give
\[
 \operatorname{Tot}(X)_N=\bigoplus_{r+s=N}X_{r,s},
 \qquad d=d_h+(-1)^r d_v.
\]
If every column is acyclic, the total complex is acyclic.  Indeed, for
a total cycle choose its component with largest first index $r$.
The cycle equation makes that component a vertical cycle, since there
is no component with first index $r+1$.  A vertical primitive, with
the appropriate sign, removes it by subtracting a total boundary;
the only newly introduced component has first index $r-1$.  Repeating
terminates after finitely many steps.  The same proof, starting with
the largest second index and using horizontal primitives, works if
every row is acyclic.  This argument uses only the finitely many
positions on a diagonal.  In a general $\mathcal A$ one may perform
this cycle calculation after applying
$\operatorname{Hom}_{\mathcal A}(P,-)$ for each projective $P$.
That functor is exact and commutes with the finite sums in question;
enough projectives makes these functors detect zero homology.

Two consequences will be used.  A morphism of first quadrant
bicomplexes that is a quasi-isomorphism in each column induces a
quasi-isomorphism of total complexes: take its columnwise mapping
cone, apply the column argument, and identify its total complex with
the mapping cone of the total map, using the usual signs.  Also, if
$X_{\bullet,s}\to C_s$ is an exact augmented horizontal complex for
each $s$, the augmentation
$\operatorname{Tot}(X)\to C$ is a quasi-isomorphism.  For this, insert
$C_s$ as the zeroth column and shift the old columns right by one;
the resulting first quadrant bicomplex has exact rows, and its total
complex is the cone of the augmentation, up to the corresponding
sign convention.  Thus its cone is acyclic.  Neither consequence
uses a totalization with infinite sums in a fixed degree.

\emph{Projective tests imply all tests.}
Assume the second condition, and write $H=\Lambda G$.  Choose a
projective resolution in the abelian category $\mathcal B$,
\[
 \cdots\longrightarrow P_2\longrightarrow P_1
       \longrightarrow P_0\longrightarrow C\longrightarrow0 .
\]
Regard $(P_r)_s$ as a first quadrant bicomplex.  Since $H$ is exact,
the augmented horizontal complex
\[
 \cdots\longrightarrow H(P_1)\longrightarrow H(P_0)
       \longrightarrow H(C)\longrightarrow0
\]
is also exact in $\mathcal B$, hence in each internal degree.
The preceding augmentation argument gives quasi-isomorphisms
$\operatorname{Tot}(P_\bullet)\to C$ and
$\operatorname{Tot}(H(P_\bullet))\to H(C)$.

Naturality of the counit gives a morphism between these bicomplexes.
On its $r$th column it is $\varepsilon_{P_r}$, a quasi-isomorphism by
assumption. The column argument therefore makes the top arrow in
the commutative square
\[
\begin{array}{ccc}
 \operatorname{Tot}(H(P_\bullet))&
       \xrightarrow{\operatorname{Tot}(\varepsilon)}&
                         \operatorname{Tot}(P_\bullet)\\
 \downarrow&&\downarrow\\
 H(C)&\xrightarrow{\varepsilon_C}&C
\end{array}
\]
a quasi-isomorphism.  Both vertical arrows are quasi-isomorphisms.
On homology the square then shows that $\varepsilon_C$ is an
isomorphism in every degree.  This proves the first condition.  The
converse implication to the second is immediate.

Notice that $H$ was applied separately to the resolution objects
$P_r$.  We neither equated $H(\operatorname{Tot}P_\bullet)$ with
$\operatorname{Tot}H(P_\bullet)$ nor applied $G$ to any infinite
coproduct.

\emph{All counits imply the localization equivalence.}
Suppose the first condition holds.  For a morphism $v:C\to C'$,
naturality gives
\[
 \varepsilon_{C'}\circ\Lambda G(v)=v\circ\varepsilon_C.
\]
Since both counits are quasi-isomorphisms, the two-out-of-three
property gives
\[
 \Lambda G(v)\in\mathcal Q
       \quad\Longleftrightarrow\quad v\in\mathcal Q.
\]
This is $G^{-1}(W)=\mathcal Q$ and, in particular, allows $G$ to
induce $\overline G$ on the localizations.  The triangular identity
\[
 \varepsilon_{\Lambda X}\circ\Lambda(\eta_X)
       =1_{\Lambda X}
\]
gives $\Lambda(\eta_X)\in\mathcal Q$, hence $\eta_X\in W$.
The unit and counit therefore descend to natural isomorphisms on the
localized categories.  Their triangular identities remain valid,
so $\overline\Lambda$ and $\overline G$ are inverse equivalences.

\emph{The converse with the localized adjunction made explicit.}
Assume the third condition.  Because both functors preserve the
respective classes being inverted, their original unit and counit
descend to natural transformations on the localizations: naturality
on original arrows implies naturality on their formal inverses.
Their triangular identities give an adjunction
$\overline\Lambda\dashv\overline G$.

The unit of an adjunction whose left adjoint is fully faithful is an
isomorphism: for each $Z$, its map on $\operatorname{Hom}(Z,-)$ is,
under adjunction, the fully faithful map on morphisms induced by the
left adjoint, so the conclusion follows from Yoneda.  Apply this to
$\overline\Lambda$.  The triangular identity makes the counit an
isomorphism on objects in its essential image, and essential
surjectivity and naturality make it an isomorphism on every object.
Thus the image of each $\varepsilon_C$ in
$\mathcal B[\mathcal Q^{-1}]$ is invertible.

Finally, a chain map whose image is invertible in this localization
is a quasi-isomorphism.  Each homology functor factors through the
localization, so it sends that map to an isomorphism.  This proves
the first condition and completes all implications.
\end{proof}

\paragraph{Scope of the projective test and of the source reading.}
The phrase \emph{projective object of
$\operatorname{Ch}_{\geq0}(\mathcal A)$} means projectivity in the
abelian category.  It is stronger than being projective in every
degree: a nonzero projective object placed only in degree one is
termwise projective but is not a projective complex.  For example the
evident epimorphism $D^1(P)\to S^1(P)$ admits no chain section when
$P\ne0$.  Testing the counit on \emph{all} termwise projective
complexes would be a stronger assumption than condition two and
therefore also sufficient; the two classes must nevertheless not be
identified.  Testing only projectives of $\mathcal A$ placed in
degree zero is not what has been proved here.

The source's arrow from $\mathcal M$ to the category of complexes is
the left adjoint; its reverse arrow is the right adjoint.
Handwritten page~53 explicitly discusses the right adjoint preserving
epimorphisms, which, together with its automatic left exactness,
gives the exactness used here.  Page~52 displays the all-counits
quasi-isomorphism condition and an equivalence after localization.
F-SM-17 establishes this precise modern formulation completely.
It does not assert that every overwritten formulation of
$(ii')$ or $(ii'')$ has been deciphered, nor does it infer a localization
equivalence from the left functor alone without the stated compatibility
of its right adjoint.

The primary locator is Grothendieck, A. (n.d.),
\emph{Pursuing stacks: The Modelizing Story, Chapter VII,
Linearization of homotopy types}, cote~134-8, handwritten
pp.~52--53, physical PDF pp.~53--54:
\url{https://grothendieck.umontpellier.fr/134-8.pdf}.
The argument above is a new complete presentation of the stated
criterion, with no claim that its categorical mechanism is historically
new or that the overwritten variants have been recovered verbatim.

% END source_cote1348_counit.tex

% BEGIN source_cote1348_purity.tex
\section{Cote 134-8: two extensions using purity}
\label{sec:SM-cote1348-purity}

Retain the precise connective setup of BR-13403 and the ringoid
construction of BR-13405. Thus $k$ is commutative,
$R$ is a small $k$-linear category with projective scalar modules
$R(x,y)$, $\mathcal M=\operatorname{Mod}\!-\!R$, and
\[
 \varphi(X)=D\otimes_R X,\qquad
 \Lambda(X)=B_\bullet\otimes_R X,\qquad
 (GC)(x)=\operatorname{Hom}_{\mathrm{Ch}_{\geq0}(k)}(B_\bullet(x),C).
 \tag{SM24.1}
\]
Here $B_\bullet\to D$ is the chosen resolution by coproducts of
left representables. The canonical unit normalization and the full
interval condition, including degree zero, remain in force.
For precision we also write the scalar compatibility used below:
\[
 J\otimes_{\mathcal M}(X\otimes_k U)
   \simeq (J\otimes_{\mathcal M}X)\otimes_k U,\qquad
 \alpha_{X\otimes_k U}=\alpha_X\otimes_k U.
 \tag{SM24.2}
\]
It holds for a $k$-linear tensor product preserving colimits in each
variable: present $U$ by free modules and use coproducts and cokernels.
We make it explicit rather than infer it from an unspecified tensor
operation. It is the scalar compatibility of the tensor setup being
used.

\paragraph{Finding F-SM-24 (the ordinary injective-dimension hypothesis can be replaced).}
Under these hypotheses, either one of the following two conditions
suffices for the specified adjunction to induce
\[
 \mathcal M[W^{-1}]\simeq D_{\geq0}(k),
 \qquad W=\Lambda^{-1}(\mathrm{qis}).
 \tag{SM24.3}
\]
\begin{enumerate}
\item There is a finite $r$ such that every projective $k$-module
      has a pure-injective resolution of length at most $r$.
\item Every $R(x,y)$ is finitely generated projective over $k$.
\end{enumerate}
The second condition imposes no injective or pure-injective dimension
condition on $k$. The first, by the known cardinality theorem stated
below, applies to every countable scalar ring and more generally to
$|k|=\aleph_t$ for a finite integer $t\geq0$.
In the first condition, $R(x,y)$ may be infinitely generated.

\begin{proof}
\emph{Purity conventions and elementary facts.}
A monomorphism of right modules is pure if tensoring with every
left module keeps it injective. Equivalently, the associated short
exact sequence remains exact after applying $\operatorname{Hom}(F,-)$
for every finitely presented right module $F$.
This equivalence also holds for ringoid modules. Indeed, a presentation
of $F$ by finite sums of representables makes a lift of a map
$F\to C$ precisely a solution of a finite system of module equations
whose constants are in the submodule. Tensoring with the cokernel
of the transposed finite matrix tests solvability of that same system.
Arbitrary tensor tests reduce to finitely presented tests because
every module is a filtered colimit of finitely presented modules.
In particular, products of pure short exact sequences are pure:
maps from a finitely presented module into a product are tuples of
maps, and products of surjections of abelian groups are surjective.

A module $E$ is pure injective if maps into $E$ extend across pure
monomorphisms. Products of pure-injective modules are pure injective,
by extending the coordinate maps. Every scalar module $Q$ has a
pure monomorphism into a pure-injective module. Here is the needed
construction, avoiding any envelope assumption. Write
$Q^+=\operatorname{Hom}_{\mathbb Z}(Q,\mathbb Q/\mathbb Z)$.
For every module $V$, the character module $V^+$ is pure injective:
use
\[
 \operatorname{Hom}_k(M,V^+)
   =\operatorname{Hom}_{\mathbb Z}(V\otimes_k M,\mathbb Q/\mathbb Z)
\]
and injectivity of $\mathbb Q/\mathbb Z$.
The evaluation map $e_Q:Q\to Q^{++}$ is pure.
Indeed $e_Q^+:Q^{+++}\to Q^+$ has the section $e_{Q^+}$.
For every $V$, the character dual of $V\otimes e_Q$ is therefore
surjective. Since $\mathbb Q/\mathbb Z$ is an injective cogenerator
of abelian groups, $V\otimes e_Q$ is injective. Thus $Q^{++}$ is
the required pure-injective target.

We will also use the fact that all functors
$\operatorname{Tor}_n^R(D,-)$ preserve pure short exact sequences.
For ringoid modules, the character dual is taken at every object
and reverses the variance; the Ext groups below are in left
$R$-modules.
For completeness, suppose $0\to X\to Y\to Z\to0$ is pure.
The evaluation pairing
$X^+\otimes_R X\to\mathbb Q/\mathbb Z$ extends across
$X^+\otimes_R X\to X^+\otimes_R Y$, which is injective by purity.
The resulting map $X^+\to Y^+$ splits restriction
$Y^+\to X^+$. Thus the character-dual sequence of left
$R$-modules splits.
The natural identity
\[
 \operatorname{Tor}_n^R(D,X)^+
   \simeq \operatorname{Ext}_R^n(D,X^+)
 \tag{SM24.4}
\]
follows by using the projective resolution $B_\bullet$ and the
exactness of character duality.
Applying $\operatorname{Ext}_R^n(D,-)$ to the split dual sequence
and dualizing back proves the assertion. Exactness is reflected
by the character dual because it is an injective cogenerator.
In particular, an $L\varphi$-acyclic module has only
$L\varphi$-acyclic pure submodules.

\emph{Relative pure injectives are killed by the actual interval.}
For $x\in R$ and a scalar module $E$, put
\[
 C_x(E)(y)=\operatorname{Hom}_k(R(x,y),E).
 \tag{SM24.5}
\]
This is the right adjoint to evaluation at $x$:
\[
 \operatorname{Hom}_{\mathcal M}(X,C_x(E))
     =\operatorname{Hom}_k(X(x),E).
 \tag{SM24.6}
\]
Call a monomorphism objectwise pure if its evaluation at every
$x$ is pure over $k$. When $E$ is pure injective, (SM24.6)
shows that $C_x(E)$ is injective with respect to objectwise pure
monomorphisms. So is every product of such modules.

By (SM24.2), the interval monomorphism
$\alpha_X:X\to J\otimes_{\mathcal M}X$ is objectwise pure:
after tensoring its evaluation with any $U$, it is the evaluation
of the assumed monomorphism $\alpha_{X\otimes_k U}$.
Consequently, if
\[
 I=\prod_{a\in A} C_{x_a}(E_a),\qquad E_a\text{ pure injective},
\]
the identity on $I$ extends across $\alpha_I$.
Thus $I$ is a retract of $J\otimes_{\mathcal M}I$.
The full interval acyclicity gives
\[
 L\varphi(I)=0.
 \tag{SM24.7}
\]
No assertion that $I$ is an ordinary injective object is needed.

\emph{The counit test is a product of coinduced scalar coefficients.}
Let $P$ be an acyclic projective object of
$\mathrm{Ch}_{\geq0}(k)$. By BR-13402,
\[
 P=\bigoplus_{n\geq0}D^{n+1}(Q_n)
   =\prod_{n\geq0}D^{n+1}(Q_n),
 \qquad Q_n\text{ projective over }k.
\]
The equality with the product holds degreewise, since only two
disks contribute in any degree. Since $G$ is a right adjoint,
\[
 GP=\prod_{n\geq0}\operatorname{Hom}_k(B_n(-),Q_n).
 \tag{SM24.8}
\]
Writing $B_n=\bigoplus_{a\in A_n}R(x_a,-)$ gives
\[
 GP=\prod_{n\geq0}\prod_{a\in A_n} C_{x_a}(Q_n).
 \tag{SM24.9}
\]
The chain-map formula behind (SM24.8) is
$\operatorname{Hom}_{\mathrm{Ch}}(B,D^{n+1}Q)
=\operatorname{Hom}_k(B_n,Q)$.
In particular, no preservation of coproducts by $G$ was assumed.

\emph{First sufficient condition.}
Choose pure-injective resolutions of length $r$ for all $Q_n$.
The functor $C_x$ is exact, since $R(x,y)$ is projective over $k$
for each $y$. Apply these functors and take the products in
(SM24.9). Products are exact in this module category, so the
result is a finite exact sequence
\[
 0\longrightarrow GP\longrightarrow I^0\longrightarrow
       I^1\longrightarrow\cdots\longrightarrow I^r
       \longrightarrow0
\]
whose terms have the form appearing in (SM24.7).
Every $L\varphi(I^j)$ vanishes.
The long exact sequences of the derived functors, or successive
distinguished triangles, now give $L\varphi(GP)=0$.
This induction is finite; no infinite totalization is involved.

\emph{Second sufficient condition.}
Embed each $Q_n$ purely into $E_n=Q_n^{++}$.
For any finitely presented right $R$-module $F$, the scalar
module $F(x)$ is finitely presented over $k$: evaluate a finite
presentation of $F$ by representables and use the present
finitely generated projectivity of all $R(x,y)$.
If $0\to Q_n\to E_n\to V_n\to0$ is the resulting pure sequence,
(SM24.6) shows that
\[
 0\to C_x(Q_n)\to C_x(E_n)\to C_x(V_n)\to0
\]
is pure as a sequence of right $R$-modules. Exactness uses
projectivity of $R(x,y)$; the purity test uses finite presentation
of $F(x)$ and scalar purity.
Taking products gives a pure monomorphism
\[
 GP\longrightarrow I=
       \prod_{n\geq0}\prod_{a\in A_n}C_{x_a}(E_n).
\]
By (SM24.7), $L\varphi(I)=0$.
By (SM24.4), every $\operatorname{Tor}_n^R(D,GP)$ injects
into the corresponding zero group for $I$.
Thus $L\varphi(GP)=0$ in this case also.

In either case the exact residual condition BR-13403 is satisfied.
Together with the unchanged canonical unit normalization,
F-SM-17 gives (SM24.3) for the specified connective adjunction.
\end{proof}

\paragraph{The countable-ring consequence uses a known theorem.}
Chach\'olski, Neeman, Pitsch, and Scherer (2018),
Theorem~9.2, p.~926, explicitly state the
Kie{\l}pi\'nski--Simson/Gruson--Jensen theorem:
a ring of cardinality $\aleph_t$, $t\in\mathbb N$, has pure
global dimension at most $t+1$. Their preceding paragraph
states this in terms of finite pure-injective dimensions.
Thus F-SM-24(1) applies with $r=t+1$, and in particular with
$r=1$ for every countably infinite scalar ring.
Finite rings are covered as well: apply that theorem to
$k\times\mathbb Z$ and take the first idempotent component of
a pure-injective resolution. A pure sequence and a pure-injective
module over a product ring are exactly pairs of such sequences
and modules over the factors, as follows from the componentwise
tensor and Hom formulas. Hence the resulting resolution over
$k$ has the same length. This proves the stated countable-ring
consequence without claiming a new proof of the cardinality theorem.

\paragraph{Finding F-SM-25 (why the second proof does not automatically extend to arbitrary projective morphism modules).}
Coinduction need not carry a scalar pure short exact sequence to a
pure sequence of ringoid modules when its morphism modules are
infinitely generated projective. The following is a complete
counterexample to that particular inference. It is not a
counterexample to the interval theorem.

\begin{proof}
Take $k=\mathbb Z$ and the one-object ringoid $R=\mathbb Z[t]$,
which is free, but infinitely generated, over $k$.
Let
$E=\widehat{\mathbb Z}=\prod_p\mathbb Z_p$ and $V=E/\mathbb Z$,
where $\mathbb Z$ embeds diagonally.
The exact sequence
\[
 0\longrightarrow\mathbb Z\longrightarrow E
       \longrightarrow V\longrightarrow0
 \tag{SM25.1}
\]
is pure over $\mathbb Z$. For every $n>0$, the map
$\mathbb Z/n\to E/nE$ is the Chinese-remainder isomorphism;
tensoring with $\mathbb Z$ also preserves the injection.
The structure theorem for finitely presented abelian groups,
followed by filtered colimits, gives all tensor tests for purity.

Choose $a=\sum_{j\geq2}2^{j!}\in\mathbb Z_2$.
It is not an integer: its binary expansion has infinitely many
zeros and ones, whereas an ordinary integer has an eventually
zero or eventually one $2$-adic expansion.
Let $e\in E$ have $2$-adic component $a$ and all other components
zero. For each $n\geq0$, choose
$0\leq r_n<2^n$ with $r_n\equiv a\pmod{2^n}$, taking $r_0=0$.
Then
\[
 q_n=(e-r_n)/2^n\in E,\qquad
 2q_{n+1}-q_n=(r_n-r_{n+1})/2^n\in\mathbb Z.
\]
Division is legitimate in the $2$-adic component by the chosen
congruence and in every other component because $2$ is a unit.
Thus $2^{-n}\mapsto q_n+\mathbb Z$ defines a homomorphism
\[
 h:\mathbb Z[1/2]\longrightarrow V.
\]
It has no lift to $E$. The image of $1$ under any such lift would
belong to $\bigcap_n2^n E$, so its $2$-adic component would be
zero. But being a lift would make its difference from $e$ a
diagonal integer, forcing $a$ to be an integer, a contradiction.

Now $F=R/(2t-1)$ is a finitely presented $R$-module whose
underlying abelian group is $\mathbb Z[1/2]$.
For $C(U)=\operatorname{Hom}_{\mathbb Z}(R,U)$, the coinduction
adjunction identifies
\[
 \operatorname{Hom}_R(F,C(U))
       =\operatorname{Hom}_{\mathbb Z}(\mathbb Z[1/2],U).
\]
The coinduced version of (SM25.1) is short exact because
$R$ is free over $\mathbb Z$. It is not pure over $R$:
the displayed Hom map from $C(E)$ to $C(V)$ is not surjective,
by the constructed $h$. This proves the claim.
\end{proof}

\paragraph{The residual generality is stated exactly.}
F-SM-24 removes the ordinary scalar injective-dimension condition
for all countable scalar rings, for the other rings of finite
pure global dimension just described, and for every scalar ring
when the ringoid is locally finitely generated projective.
The remaining case has neither available finite pure-injective
dimension nor the finite-presentation property used in the second
proof. For that case the exact unresolved assertion is still
\[
 \operatorname{Tor}_j^R
 \left(D,\prod_{n\geq0}\prod_{a\in A_n}C_{x_a}(Q_n)\right)=0
 \quad(j\geq0)
 \tag{SM25.2}
\]
for every sequence of projective scalar modules $Q_n$.
F-SM-25 blocks one proposed automatic purity argument for
(SM25.2); it does not provide data satisfying the original
interval and normalization assumptions and violating (SM25.2).
The ordinary finite-injective-dimension proof, the relative-purity
proofs here, and a hypothetical infinite resolution must therefore
not be conflated.

\subsection*{Source for the cardinality input}
Chach\'olski, W., Neeman, A., Pitsch, W., \& Scherer, J. (2018).
Relative homological algebra via truncations.
\emph{Documenta Mathematica, 23}, 895--937.
\url{https://ems.press/content/serial-article-files/26419?nt=1}.
Exact input: \S9.2 and Theorem~9.2, p.~926.
The source attributes this result to
Kie{\l}pi\'nski--Simson (1975) and Gruson--Jensen (1973);
their original proofs are not claimed to have been reconstructed here.


% END source_cote1348_purity.tex

% BEGIN source_cote1348_fp.tex
\section{Further progress on the arbitrary-ring case of cote 134-8}

\subsection{F-SM-26: the full interval forces infinite relative Tor amplitude}
Assume the data of BR-13403 and the full interval condition (134.1).
Write \(T=\mathbf1\), and suppose \(k\ne0\).  For every object \(X\), put
\[
 S(X)=\operatorname{coker}\bigl(X\xrightarrow{\alpha_X}J\otimes X\bigr).
\]
Then
\[
 L\varphi(S^nT)\simeq k[n]\qquad(n\ge0).
 \tag{SM26.1}
\]
In the ringoid situation of BR-13405, where
\(\varphi(X)=D\otimes_R X\), this implies
\[
 \operatorname{Tor}^{R}_n(D,S^nT)\cong k\ne0
 \qquad(n\ge0).
 \tag{SM26.2}
\]
Consequently \(D\) has infinite flat dimension, and hence infinite
projective dimension, as a left \(R\)-module.  A finite flat-dimension
hypothesis on \(D\) cannot close the remaining counit test while retaining
the actual nonzero normalized interval hypotheses.

\paragraph{Complete proof.}
The defining sequence
\(0\to X\to J\otimes X\to S(X)\to0\) is exact.
The exact complex-valued representative \(\Lambda\) sends it to a short
exact sequence of complexes.  Its associated triangle, together with
\(\Lambda(J\otimes X)\simeq0\), gives
\[
 L\varphi(SX)\simeq L\varphi(X)[1].
\]
Here homological shift \([1]\) raises the homology degree by one.
Iteration and the canonical normalization \(L\varphi(T)\simeq k[0]\)
give (SM26.1).  Taking degree \(n\) homology gives (SM26.2).
Finite flat dimension \(r\) would make every such Tor group zero for
\(n>r\), a contradiction.  Finite projective dimension implies finite
flat dimension, completing the proof. \(\square\)

\subsection{F-SM-27: an arbitrary-scalar result when \(D\) is of type \(\mathrm{FP}_\infty\)}
Retain the ringoid, scalar, exact-adjunction and full interval hypotheses
of BR-13403--05.  Thus \(R\) is a small \(k\)-linear category, each
\(R(x,y)\) is a projective scalar module,
\(\mathcal M=\operatorname{Fun}_k(R^{\mathrm{op}},\operatorname{Mod}k)\),
\(\varphi=D\otimes_R-\), and the specified adjunction is formed from
a resolution \(B_\bullet\to D\) by coproducts of left representables.
Retain the canonical unit normalization, including degree zero.

Assume the scalar compatibility used in F-SM-24:
\[
 J\otimes(X\otimes_kU)\cong(J\otimes X)\otimes_kU,\qquad
 \alpha_{X\otimes_kU}=\alpha_X\otimes_kU.
 \tag{SM27.1}
\]
This holds for a \(k\)-linear tensor product preserving colimits in each
variable.  Finally suppose that \(D\), as a left \(R\)-module, has a
resolution
\[
 \cdots\longrightarrow A_2\longrightarrow A_1
 \longrightarrow A_0\longrightarrow D\longrightarrow0
 \tag{SM27.2}
\]
in which every \(A_n\) is a finitely generated projective left
\(R\)-module: equivalently, a retract of a finite coproduct of left
representables.  This is the additional condition
\(D\in\mathrm{FP}_\infty(R)\); no finite length for (SM27.2) is required.

Then the exact adjunction specified in BR-13405 induces
\[
 W^{-1}\mathcal M\simeq D_{\ge0}(k).
 \tag{SM27.3}
\]
There is no cardinality, ordinary injective-dimension or pure
injective-dimension restriction on \(k\), and the scalar modules
\(R(x,y)\) need not be finitely generated.  More strongly, with
\[
 C_x(U)(y)=\operatorname{Hom}_k(R(x,y),U),
 \tag{SM27.4}
\]
\(L\varphi\) vanishes on every \(C_x(U)\) and on every product of such
objects, for arbitrary scalar modules \(U\).

\paragraph{Exact known input.}
We use Benson and Goodearl's (2000) periodic-flat-module theorem:
a flat module \(V\) in an exact sequence
\(0\to V\to P\to V\to0\), with \(P\) projective, is projective.
The precise primary locator is Theorem 2.5, printed p.~50
(also stated as Theorem 1.1 on p.~46), in
\emph{Periodic flat modules, and flat modules for finite groups}.
Its proof occupies pp.~50--51, after Lemmas 2.1--2.4.
The complete deduction from that known theorem to the contractibility
needed below is given here.  In particular, ordinary acyclicity of a
complex of projectives, unbounded to the right cohomologically, is
not substituted for pure acyclicity.

\paragraph{Complete proof.}
We use the convention that a pure monomorphism of scalar modules
remains injective after tensoring with every scalar module.
A scalar module \(E\) is pure-injective if maps to \(E\) extend across
all such monomorphisms.  For every scalar module \(V\), its character
module
\[
 V^+=\operatorname{Hom}_{\mathbf Z}(V,\mathbf Q/\mathbf Z)
\]
is pure-injective: tensor--Hom adjunction identifies extension across
a pure monomorphism with extension to the injective abelian group
\(\mathbf Q/\mathbf Z\).  We will use this elementary fact both to
produce tests and to detect purity.

\emph{1. Coinduction with pure-injective coefficients is annihilated.}
There is a natural adjunction
\[
 \operatorname{Hom}_{\mathcal M}(X,C_x(E))
   =\operatorname{Hom}_k(X(x),E).
 \tag{SM27.5}
\]
Consequently, when \(E\) is pure-injective, \(C_x(E)\) is injective
with respect to all monomorphisms that are scalar-pure at every
object of \(R\).  Condition (SM27.1) and the interval's monomorphism
condition show that \(\alpha_X:X\to J\otimes X\) is objectwise
scalar-pure: tensoring it with any \(U\) gives the monomorphism
\(\alpha_{X\otimes_k U}\).  Thus \(\alpha_{C_x(E)}\) splits in
\(\mathcal M\), by extending the identity using (SM27.5).
The interval condition gives
\[
 L\varphi(C_x(E))=0\qquad(E\text{ scalar pure-injective}).
 \tag{SM27.6}
\]
This argument does not declare \(C_x(E)\) ordinarily injective.

\emph{2. Dualizing the finite-projective resolution.}
Fix \(x\), and form the cochain complex
\[
 K_x^n=\operatorname{Hom}_{R\text{-}\mathrm{Mod}}
                (A_n,R(x,-)),\qquad
 \delta^n(f)=f\circ d_{n+1}.
 \tag{SM27.7}
\]
Every \(K_x^n\) is a projective scalar module: for
\(A_n=R(y,-)\) it is \(R(x,y)\), and the general case follows by
finite coproducts and retracts.  For every scalar module \(U\),
evaluation and finite-projective duality give natural isomorphisms
\[
 A_n\otimes_R C_x(U)
   \cong\operatorname{Hom}_k(K_x^n,U).
 \tag{SM27.8}
\]
For a representable this is the defining formula (SM27.4).
Finite coproducts convert to finite products, which coincide, and
retracts preserve the isomorphism.  This also proves naturality in
\(A_n\); hence the isomorphisms commute with the differentials.
The homological complex \(A_\bullet\otimes_R C_x(U)\), computing
\(L\varphi(C_x(U))\), is therefore the transpose complex
\(\operatorname{Hom}_k(K_x^\bullet,U)\).

By (SM27.6), this transpose complex is acyclic for every pure-injective
\(U\).  In particular take the injective cogenerator
\(E_0=\operatorname{Hom}_{\mathbf Z}(k,\mathbf Q/\mathbf Z)\).
Then
\(\operatorname{Hom}_k(K_x^\bullet,E_0)
 =\operatorname{Hom}_{\mathbf Z}(K_x^\bullet,\mathbf Q/\mathbf Z)\);
exactness and the cogenerator property imply that \(K_x^\bullet\)
is acyclic.

Let \(Z^n=\ker(\delta^n)\).  As \(K_x^\bullet\) starts in degree zero,
acyclicity implies \(Z^0=0\) and gives exact sequences
\[
 0\longrightarrow Z^n\longrightarrow K_x^n
   \xrightarrow{\delta^n}Z^{n+1}\longrightarrow0.
 \tag{SM27.9}
\]
These sequences are pure.  To see this directly, fix a pure-injective
\(E\) and a map \(f:Z^n\to E\), with \(n\ge1\).
The map \(f\delta^{n-1}:K_x^{n-1}\to E\) is a cycle in the
transpose complex.  Its acyclicity gives \(g:K_x^n\to E\) with
\(g\delta^{n-1}=f\delta^{n-1}\).  Surjectivity onto \(Z^n\)
implies \(g|_{Z^n}=f\).  The assertion for \(n=0\) is immediate.
In particular this extension property holds for every \(E=V^+\).
Tensor--Hom adjunction and the fact that
\(\mathbf Q/\mathbf Z\) detects nonzero elements imply that
\(V\otimes_k Z^n\to V\otimes_k K_x^n\) is injective for every
\(V\).  This proves purity of (SM27.9).

\emph{3. Contractibility, using the genuine periodicity input.}
A pure submodule of a flat module is flat.  In the present instance
one can check this without another theorem: for an injection
\(V'\to V\), put \(Z^n\to K_x^n\) in the tensor square.
Both horizontal maps are injective by purity, and the right vertical
map is injective by flatness of \(K_x^n\).  Therefore the left
vertical map is injective.  This is the flatness criterion for
\(Z^n\).

Take the direct sum of (SM27.9) for \(n\ge0\).
Since \(Z^0=0\), the two end sums are canonically isomorphic after
shifting their indices.  We obtain
\[
 0\longrightarrow Z\longrightarrow P\longrightarrow Z
 \longrightarrow0,\qquad
 Z=\bigoplus_{n\ge0}Z^n,\quad P=\bigoplus_{n\ge0}K_x^n.
 \tag{SM27.10}
\]
Here \(Z\) is flat and \(P\) is projective.
The cited Benson--Goodearl Theorem 2.5 makes \(Z\) projective.
Each \(Z^n\), as a direct summand of \(Z\), is consequently
projective.  Thus every sequence (SM27.9) splits.
Choosing splittings identifies
\(K_x^n=Z^n\oplus Z^{n+1}\) and
\(\delta^n(z,w)=(w,0)\); the maps
\(h^n(z,w)=(0,z)\) give a contracting homotopy
(with \(h^0=0\)).  This proves contractibility of \(K_x^\bullet\),
including its unbounded right-hand end.

For every scalar module \(U\), transpose this homotopy.
Then \(\operatorname{Hom}_k(K_x^\bullet,U)\) is contractible.
By (SM27.8),
\[
 L\varphi(C_x(U))=0\qquad\text{for every }x,U.
 \tag{SM27.11}
\]

\emph{4. Products and the specified adjunction.}
Each functor \(A_n\otimes_R-\) commutes with arbitrary products,
because \(A_n\) is a retract of a finite sum of representables,
whose tensor functors are evaluation functors.
Therefore for any family \(X_i\) there is a natural isomorphism of
complexes
\[
 A_\bullet\otimes_R\prod_iX_i
   \cong\prod_i(A_\bullet\otimes_R X_i).
 \tag{SM27.12}
\]
Products of scalar modules are exact, so a product of acyclic
complexes is acyclic.  Formula (SM27.11) now proves the stronger
vanishing assertion for every product of coinduced objects.
Notice that this product argument computes \(L\varphi\) with \(A\);
it does not assume that the originally chosen \(B_n\) is finite.

Let \(P\) be an acyclic projective object of
\(\operatorname{Ch}_{\ge0}(k)\).  BR-13402 writes
\[
 P=\bigoplus_{n\ge0}D^{n+1}(Q_n)
   =\prod_{n\ge0}D^{n+1}(Q_n),
 \tag{SM27.13}
\]
where each \(Q_n\) is projective; the equality of sum and product is
degreewise, since only two summands occur in any chain degree.
The right adjoint \(G\) preserves products.  For every \(y\),
\[
 G(D^{n+1}Q_n)(y)
   =\operatorname{Hom}_{\mathrm{Ch}}(B_\bullet(y),D^{n+1}Q_n)
   =\operatorname{Hom}_k(B_n(y),Q_n).
\]
The latter formula is the disk adjunction, obtained by sending the
degree-\(n\) map freely to \(Q_n\); its degree-\(n+1\) companion is
determined by the differential.  Write the chosen resolution term as
\(B_n=\bigoplus_{a\in I_n}R(x_{n,a},-)\).
It follows naturally as right \(R\)-modules that
\[
 GP\cong\prod_{n\ge0}\prod_{a\in I_n}C_{x_{n,a}}(Q_n).
 \tag{SM27.14}
\]
The preceding product vanishing gives \(L\varphi(GP)=0\).
BR-13403 and F-SM-17 conclude (SM27.3).  This last invocation uses
the canonical unit normalization to handle projective complexes'
degree-zero summands; the coinduction vanishing itself did not need
that normalization. \(\square\)

\paragraph{Exact remaining scope.}
Condition \(\mathrm{FP}_\infty\) is an added hypothesis on \(D\),
not a consequence asserted here of the existence of finitely
presented projective generators of \(\mathcal M\).
It supplies finite-projective duality in (SM27.8) and product
commutation in (SM27.12).  Neither identity has been asserted for an
arbitrary infinite coproduct of representables.
Thus F-SM-27 adds a third proved range to F-SM-24:
arbitrary scalar rings and arbitrary projective scalar morphism
modules, provided \(D\) has this finite-generation property in every
resolution degree.  The case without any of these extra hypotheses
remains the exact test (134.6).

A countable-subring argument has not removed that last condition.
Although a tensor element uses finitely many tensors, its coefficients
in \(GP\) are global linear maps on the modules \(B_n(y)\).
A boundary found after restriction to a smaller ringoid can require
extending such maps back to all of \(B_n(y)\), compatibly with the
ringoid action.  Finite support of the outer tensor expression does
not provide that extension.  No automatic descent or compatible
extension assertion is used in the proved result.

\paragraph{Primary reference.}
Benson, D. J., \& Goodearl, K. R. (2000).
Periodic flat modules, and flat modules for finite groups.
\emph{Pacific Journal of Mathematics, 196}(1), 45--67.
\url{https://doi.org/10.2140/pjm.2000.196.45}.
Theorem 2.5, p.~50; proof pp.~50--51.
Primary publisher PDF:
\url{https://msp.org/pjm/2000/196-1/pjm-v196-n1-p03-s.pdf}.


% END source_cote1348_fp.tex

% BEGIN source_cote1348_compact_unit.tex
\subsection{F-SM-28: a compact unit and compact interval do not force \(D\in\mathrm{FP}_\infty\)}
For every nonzero commutative ring \(k\), there is an example satisfying
the full unit normalization and interval hypotheses of BR-13403--05,
with the following additional properties:
\begin{itemize}
\item the monoidal unit and the interval are finitely presented
projective objects;
\item the category has a countable set of finitely presented projective
generators, and all scalar modules of morphisms between these generators
are either \(0\) or \(k\);
\item the left ringoid module \(D\) representing \(\varphi\) is not
finitely generated, so in particular is not of type
\(\mathrm{FP}_\infty\).
\end{itemize}
The desired connective localization equivalence nevertheless holds
in this example.  The example therefore disproves the proposed
automatic-\(\mathrm{FP}_\infty\) inference, not the original
equivalence.

\paragraph{Complete construction and proof.}
Let \(\mathbb N\) denote the category with one morphism \(i\to j\)
when \(i\le j\).  Put
\[
 \mathcal M=\operatorname{Fun}
     \bigl(\mathbb N,\operatorname{Ch}_{\ge0}(k)\bigr),
 \qquad
 T=\operatorname{const}_{\mathbb N}(k[0]).
 \tag{SM28.1}
\]
This is a cocomplete \(k\)-linear abelian category.  Give it the
pointwise tensor product of chain complexes, with the usual signed
differential; its monoidal unit is \(T\).  Define
\[
 \kappa(U)=\operatorname{const}_{\mathbb N}(U[0]),\qquad
 \Lambda(X)=\operatorname*{colim}_{i\in\mathbb N}X_i,\qquad
 \varphi(X)=H_0\Lambda(X).
 \tag{SM28.2}
\]
The constant functor \(G(C)=\operatorname{const}_{\mathbb N}(C)\)
is right adjoint to \(\Lambda\), by the universal property of colimits.
Both are exact: exactness of \(G\) is pointwise, and filtered colimits
of modules are exact, so \(\Lambda\) is exact in every chain degree.
The latter familiar exactness can also be read elementwise: a kernel
or relation in a filtered colimit is witnessed at one sufficiently
large index; a finite collection of such witnesses has a common
upper index.  Thus kernels and cokernels of short exact sequences
commute with this colimit.

For a scalar module \(U\), a chain map \(X_i\to U[0]\) is precisely
a map \(H_0X_i\to U\).  Taking compatible families gives
\[
 \operatorname{Hom}_{\mathcal M}(X,\kappa U)
   \cong\operatorname{Hom}_k
       \left(\operatorname*{colim}_{i}H_0X_i,U\right)
   =\operatorname{Hom}_k(\varphi X,U).
 \tag{SM28.3}
\]
Hence \(\varphi\dashv\kappa\).
The scalar functor \(\kappa\) is exact, preserves all inverse limits,
and is compatible with the tensor products.  It also has the required
right adjoint: since \(0\) is initial in \(\mathbb N\), a map
\(T\to X\) is determined by an arbitrary element of \((X_0)_0\);
thus \(\operatorname{Hom}_{\mathcal M}(T,X)=(X_0)_0\), and
scalar tensor--Hom gives \(\kappa\dashv\operatorname{Hom}_{\mathcal M}(T,-)\).

For \(a\ge1\), let \(C_a=D^a(k)\) be the disk with identity
differential from degree \(a\) to degree \(a-1\), and set
\(C_0=k[0]\).  For each \(i,a\ge0\), define
\[
 (P_{i,a})_j=
 \begin{cases}0,&j<i,\\ C_a,&j\ge i,\end{cases}
 \tag{SM28.4}
\]
with identity transition maps after \(i\).  These free extensions
satisfy
\[
 \operatorname{Hom}_{\mathcal M}(P_{i,a},X)
   =\operatorname{Hom}_{\mathrm{Ch}}(C_a,X_i)
   =(X_i)_a.
 \tag{SM28.5}
\]
The formula for \(a=0\) is the sphere adjunction and the formula for
\(a\ge1\) is the left disk adjunction.  Evaluation is exact and
commutes with filtered colimits.  Therefore every \(P_{i,a}\) is
projective and finitely presented.  They generate: for every element
in every \((X_i)_a\), the corresponding map \(P_{i,a}\to X\) contributes
that element, and the sum of these maps is a degreewise epimorphism.
In particular
\[
 T=P_{0,0}
\]
is finitely presented and projective.

We have
\[
 \Lambda(P_{i,0})=k[0],\qquad
 \Lambda(P_{i,a})=D^a(k)\quad(a\ge1).
\]
Thus \(\Lambda\) has zero positive homology on projectives, first on
the generators and then on their coproducts and retracts.
The finite-diagonal proof BR-13401 gives
\[
 L\varphi(X)\simeq\Lambda(X).
 \tag{SM28.6}
\]
In particular, \(L\varphi(T)\to k[0]\) is the canonical identity
normalization.

Take
\[
 J=P_{0,1}=\operatorname{const}_{\mathbb N}(D^1(k)),
 \qquad
 \alpha:T\longrightarrow J
\]
to be the identity inclusion in degree zero.  The interval \(J\)
is also finitely presented and projective.
For any \(X\), at every diagram index the tensor complex has
\[
 ((J\otimes X)_i)_n=(X_i)_n\oplus(X_i)_{n-1},\qquad
 d(x,y)=(d_Xx+y,-d_Xy),
 \tag{SM28.7}
\]
where the second summand is zero for \(n=0\).
The map \(\alpha_X\) is \(x\mapsto(x,0)\), hence monic.
The maps \(h_n(x,y)=(0,x)\) satisfy \(dh+hd=1\).
They commute with the diagram structure maps, so give a natural
contraction in \(\mathcal M\).  Applying the exact colimit functor
preserves this contraction.  Equation (SM28.6) therefore gives
\(L\varphi(J\otimes X)=0\) in every degree.
The scalar compatibility (SM27.1) also holds directly.
This verifies the full interval, including its degree-zero clause.

Now take \(R\) to be the full small \(k\)-linear category of the
objects \(P_{i,a}\), with
\[
 R((i,a),(j,b))=
 \operatorname{Hom}_{\mathcal M}(P_{i,a},P_{j,b}).
\]
Restricted Yoneda identifies \(\mathcal M\) with right \(R\)-modules.
For completeness, any object has a presentation by sums of the
\(P_{i,a}\), as just proved; projectivity and finite presentation make
restricted Yoneda exact and make it preserve those sums.  Such a
presentation proves full faithfulness, and a presentation of any
right \(R\)-module by representables gives essential surjectivity
by taking the corresponding cokernel in \(\mathcal M\).

By (SM28.5), a scalar morphism module is zero when \(i<j\), and
when \(i\ge j\) it is
\(\operatorname{Hom}_{\mathrm{Ch}}(C_a,C_b)=(C_b)_a\).
For \(b=0\) this is \(k\) only when \(a=0\).
For \(b\ge1\) it is \(k\) precisely when \(a=b\) or \(a=b-1\).
Thus each scalar morphism module is \(0\) or \(k\).
All hypotheses required for the exact-adjunction construction
BR-13405 are satisfied.

The left \(R\)-module \(D\) representing \(\varphi\) is
\[
 D(i,a)=\varphi(P_{i,a})=
 \begin{cases}k,&a=0,\\0,&a\ge1.\end{cases}
 \tag{SM28.8}
\]
Furthermore
\[
 R((i,0),(j,0))=
 \begin{cases}k,&j\le i,\\0,&j>i,\end{cases}
 \tag{SM28.9}
\]
and the nonzero arrows act on (SM28.8) by their scalar coefficients.
If \(D\) were finitely generated, it would be a quotient of a finite
sum of left representables.  A map from a representable based at
\((i,a)\) to \(D\) is specified by an element of \(D(i,a)\).
For \(a\ge1\) it is zero.  For the remaining finitely many indices
\((i,0)\), choose an upper bound \(N\).  At every \((j,0)\) with
\(j>N\), all those representables evaluate to zero by (SM28.9).
Their image cannot equal \(D(j,0)=k\ne0\).
This contradiction proves that \(D\) is not finitely generated.

Finally, this construction still models the desired connective
derived category.  The counit
\(\Lambda G(C)\to C\) is an identity for every complex \(C\).
The unit \(X\to G\Lambda X\) becomes an identity after \(\Lambda\),
hence belongs to \(W=\Lambda^{-1}(\mathrm{qis})\).
Also \(G\) sends quasi-isomorphisms to \(W\), since \(\Lambda G\)
is the identity.  The ordinary adjunction consequently descends
with invertible unit and counit, and gives
\[
 W^{-1}\mathcal M\simeq D_{\ge0}(k).
\]
For an adjunction formed from any representable resolution as in
BR-13405, the same conclusion, including the raw counit criterion,
also follows from F-SM-24's finitely generated scalar-morphism case.
Thus the failure of finite generation for \(D\) creates no
counterexample to the original claim in this example. \(\square\)

\paragraph{Source scope and consequence for the search.}
A fresh visual reading of manuscript footers 47 and 48 found the
conditions (a)--(d), followed by the corollary's finitely presented
projective generators and its scalar-projectivity marginal condition.
Those lines do not state finite presentation of the monoidal unit.
The example above supplies that stronger property anyway, together
with finite presentation and projectivity of the interval itself.
Therefore neither adding a compact unit nor adding this compact
interval makes F-SM-27 an automatic solution of the full general-ring
question.  This is an explicit construction, not an inference from
an unread part of the manuscript.


% END source_cote1348_compact_unit.tex

% BEGIN source_cote1348_realization.tex
\section{F-SM-29: the underlying connective localization over every scalar ring}

\subsection{Precise statement and distinction from the ordinary right adjoint}
Let \(k\) be an arbitrary commutative ring and \(R\) a small
\(k\)-linear category.  Set
\(\mathcal M=\operatorname{Fun}_k(R^{\mathrm{op}},\operatorname{Mod}k)\).
Assume the following data from BR-13403:
\begin{enumerate}
\item an exact scalar functor \(\kappa(U)=T\otimes_kU\), preserving
coproducts, with left adjoint \(\varphi\);
\item an exact additive left adjoint
\(\Lambda:\mathcal M\to\operatorname{Ch}_{\ge0}(k)\), with a right
adjoint \(G\), such that \(H_0\Lambda=\varphi\) and
\(H_j\Lambda(P)=0\) for \(j>0\) and projective \(P\);
\item the canonical map
\(\Lambda\kappa(k)\to k[0]\), obtained from the counit of
\(\varphi\dashv\kappa\), is a quasi-isomorphism;
\item for every \(X\) there is the interval monomorphism
\(X\to J\otimes X\), and
\(L\varphi(J\otimes X)=0\), including degree zero.
\end{enumerate}
Put \(W=\Lambda^{-1}(\mathrm{qis})\).  Then the functor induced by
\(\Lambda\), equivalently by \(L\varphi\), is an equivalence
\[
 W^{-1}\mathcal M\xrightarrow{\ \sim\ }D_{\ge0}(k).
 \tag{SM29.1}
\]
There is no restriction on the cardinality or homological dimensions
of \(k\), no finite-generation restriction on the scalar morphism
modules, and no \(\mathrm{FP}_\infty\) hypothesis on \(D\).
In particular BR-13405 supplies these exact-adjunction data under
the corollary's stated locally projective ringoid hypotheses.

The source identification here is precise.  On manuscript footer 47,
condition (a) supplies the left adjoint \(\varphi\) of \(\kappa\)
and identifies its right adjoint as
\(X\mapsto\operatorname{Hom}_{\mathcal M}(\mathbf1,X)\).
Thus exactness and preservation of coproducts of \(\kappa\) follow
from its two adjoints; footer 46 also explicitly writes
``\(\kappa\) exact.''  Condition (b) displays the derived adjunction,
and (c) displays \(L\varphi(\kappa(k))\to k\) as an isomorphism.
We use this arrow as the displayed adjunction's counit, not just an
unspecified isomorphism between these objects.  Footer 48 adds the
finitely presented projective generators, and its right margin adds
scalar projectivity of the morphism modules.  These are the
module-category corollary's hypotheses.  We do not assert that the
broader abstract conditions (a)--(d), without this presentability
and exact-adjunction construction, have been settled by this proof.

The inverse constructed below uses a homotopy colimit.
The proof does not assert that the ordinary right adjoint \(G\)
preserves every quasi-isomorphism, nor that
\(\Lambda G(C)\to C\) is a quasi-isomorphism for every ordinary
complex \(C\).  That stronger assertion remains precisely the
acyclic-projective test of BR-13403.  The right \emph{derived}
adjoint is a quasi-inverse by (SM29.1).

\subsection{Exact background inputs}
We use Smith's model-structure recognition theorem as proved in
Beke (2000), Theorem 1.7, author-PDF p.~8.
Its hypotheses are checked below.  The accessibility facts used
are Beke's Proposition 1.15, p.~12, and Proposition 1.18, p.~13.
For a cofibrantly generated model category, the projective model
structure on diagrams over a small category is the standard
diagram-model theorem: see Dugger, current author version,
Theorem 19.4, p.~91, and the discussion of its generating
cofibrations in Beke, author-PDF p.~20.
We use Dold--Kan for abelian categories in the precise form of
the Stacks Project, Theorem 14.24.3, Tag \nolinkurl{019G}.

The homotopy-colimit totalization formula needed in the argument
is proved below, including its naturality and augmentation
compatibility.  It also agrees with Dugger's Proposition 19.9,
p.~93 of the current 109-page author PDF.
None of these known background theorems is attributed to Alp\"oge.

\subsection{Complete proof}

\paragraph{1. All scalar modules satisfy canonical normalization.}
Write
\[
 e_U:\Lambda\kappa(U)\longrightarrow U[0]
 \tag{SM29.2}
\]
for the natural composite through
\(H_0\Lambda\kappa(U)[0]=\varphi\kappa(U)[0]\) and the adjunction
counit.  The hypothesis says \(e_k\) is a quasi-isomorphism.
Both \(\Lambda\) and \(\kappa\) preserve coproducts.
Consequently \(e_U\) is a quasi-isomorphism for free \(U\), and
then for projective \(U\) by retracts.

For arbitrary \(U\), take a projective scalar resolution
\(Q_\bullet\to U\) in nonnegative homological degrees.
Apply the exact functor \(\Lambda\kappa\) to obtain the
first-quadrant double complex
\[
 E_{a,b}=(\Lambda\kappa Q_a)_b.
\]
The augmentation
\(\operatorname{Tot}E\to\Lambda\kappa U\) is a quasi-isomorphism:
in each internal degree \(b\), the augmented horizontal sequence
is exact.  The maps \(e_{Q_a}\) give a quasi-isomorphism
\(\operatorname{Tot}E\to Q_\bullet\), by taking vertical homology.
The augmentation \(Q_\bullet\to U[0]\) is a quasi-isomorphism.
These maps form a commutative square with \(e_U\), so \(e_U\)
is a quasi-isomorphism.  Each total degree has finitely many
summands; both filtration arguments therefore have finite
diagonals.  We have proved the \emph{canonical}, natural assertion
\[
 \Lambda\kappa(U)\simeq U[0]\quad\text{for every scalar }U.
 \tag{SM29.3}
\]

\paragraph{2. The model structure on \(\mathcal M\).}
The module category of a small ringoid is locally finitely
presentable.  Concretely, presentations use coproducts of the
representables \(P_x=R(-,x)\), and every presentation is the filtered
colimit of its finite sets of generators and relations.  Limits,
colimits and exactness are computed objectwise.
In particular \(\mathcal M\) is a Grothendieck abelian category.

There is an explicit set of generating monomorphisms:
\[
 \mathcal I=\{\,U\hookrightarrow P_x:
       x\in R,\ U\text{ a submodule of }P_x\,\}.
 \tag{SM29.4}
\]
It is a set, because \(R\) is small and each representable has a
set of subobjects.  We verify that every monomorphism is a
transfinite composition of pushouts of maps in \(\mathcal I\).
Identify a given monomorphism with \(A\subset B\), and well-order
all the elements of all modules \(B(x)\).  At a successor stage,
adjoin the next element \(b\in B(x)\) to the submodule \(A_\beta\)
already constructed.  It defines \(P_x\to B\).  Let
\(U\subset P_x\) be the inverse image of \(A_\beta\).
Then the pushout
\[
 A_\beta\amalg_U P_x
\]
maps isomorphically onto \(A_\beta+\operatorname{im}(P_x\to B)\):
the kernel of \(P_x\to B\) is contained in \(U\), so the pushout
imposes exactly the needed identifications.  At a limit stage
take the union.  The construction eventually gives \(B\).
Conversely pushouts, transfinite compositions and retracts of
monomorphisms remain monomorphisms in modules.  Thus
\(\operatorname{cof}(\mathcal I)\) is exactly the monomorphisms.

The weak-equivalence class \(W\) is accessible and accessibly
embedded in the arrow category.  Indeed, \(\Lambda\), being a
left adjoint, preserves colimits; homology of module complexes
commutes with filtered colimits.  Hence
\[
 H_*\Lambda:\mathcal M\longrightarrow
       \prod_{n\ge0}\operatorname{Mod}k
\]
is accessible.  The class of isomorphisms in the arrow category
of graded modules is accessible: evaluation on the source is an
equivalence from that class to graded modules, and its inclusion
preserves filtered colimits.  The accessible inverse-image
criterion, Beke Proposition 1.18, applies to \(H_*\Lambda\).
Proposition 1.15 then gives Smith's solution-set condition at
the set \(\mathcal I\).

We check the other recognition conditions individually.
Closure under retracts and two-out-of-three follow from the
same properties of quasi-isomorphisms.
A monomorphism lies in \(W\) if and only if its cokernel has
acyclic \(\Lambda\)-image, because \(\Lambda\) is exact.
A pushout of a monomorphism has the same cokernel, so pushouts
preserve monomorphisms in \(W\).
Transfinite compositions also preserve this class: \(\Lambda\)
preserves the relevant colimits, filtered colimits of modules
are exact, and all the induced maps on homology in the
transfinite sequence are isomorphisms.

Finally, a map with the right lifting property against
\(\mathcal I\) has it against every monomorphism.
Lifting against \(0\to Y\), where \(Y\) is its target, supplies
a section.  Its kernel \(I\) is injective: a map from a subobject
to \(I\) extends across the containing object by applying the
same lifting property with zero map to \(Y\).
Thus the original map is a split epimorphism with injective
kernel.  The interval monomorphism
\(I\to J\otimes I\) splits, by injectivity, so
\[
 L\varphi(I)=0
\]
as a retract of \(L\varphi(J\otimes I)=0\).
BR-13401 identifies \(L\varphi(I)\) with \(\Lambda I\).
The split epimorphism consequently lies in \(W\).

All four conditions of Beke Theorem 1.7 are now verified.
There is a cofibrantly generated model structure on
\(\mathcal M\) with weak equivalences \(W\) and all monomorphisms
as cofibrations.  Every object is cofibrant.  Its homotopy
category is the localization \(W^{-1}\mathcal M\).

The same recognition argument gives a model structure on
\(\operatorname{Ch}_{\ge0}(k)\) with quasi-isomorphisms as weak
equivalences and all monomorphisms as cofibrations.
This category is itself a module category: its finitely presented
projective generators are \(k[0]\) and the disks \(D^n(k)\),
\(n\ge1\), by the evaluation adjunctions in BR-13402.
For its last recognition condition, an injective chain object
embeds into \(D^1(k)\otimes I\), a contractible complex, via
the degree-zero endpoint; the embedding splits and therefore
the injective chain object is acyclic.  All the other
conditions are exactly the preceding module-category argument.
This construction is within nonnegative complexes; no
unbounded-model restriction is being assumed.
The target homotopy category is \(D_{\ge0}(k)\).

Exactness of \(\Lambda\) and the definition of \(W\) show that
\(\Lambda\) preserves cofibrations and trivial cofibrations.
Together with its given right adjoint \(G\), it is therefore
a Quillen adjunction.  Its left derived functor is represented
by \(\Lambda\) itself, since every source object is cofibrant.

\paragraph{3. Homotopy colimits and the normalized total complex.}
We record the precise realization fact used below.
For a simplicial object
\(Y:\Delta^{\mathrm{op}}\to\operatorname{Ch}_{\ge0}(k)\),
there is a natural isomorphism in \(D_{\ge0}(k)\)
\[
 \operatorname*{hocolim}_{\Delta^{\mathrm{op}}}Y
       \simeq\operatorname{Tot}N_{\mathrm{simp}}Y.
 \tag{SM29.5}
\]
The totalization is the direct-sum totalization, with the usual
sign on one differential.  Its diagonals are finite.

Here is a proof in the model structures just specified.
Use the projective model structure on simplicial diagrams and
a functorial projectively cofibrant replacement \(QY\to Y\).
The colimit/constant-diagram adjunction is Quillen, since the
constant functor preserves the objectwise fibrations and trivial
fibrations of the projective diagram structure.
By definition the homotopy colimit is \(\operatorname{colim}QY\).
There is a natural augmentation
\[
 \theta_Z:\operatorname{Tot}N_{\mathrm{simp}}Z
              \longrightarrow\operatorname{colim}Z
\]
for every such diagram \(Z\), induced by its degree-zero
augmentation and zero on positive simplicial degrees.

Let \(F_{[m]}\) be left adjoint to evaluation at \([m]\) on
simplicial diagrams.  It is the copower by the representable
simplicial set \(\Delta[m]\).  For every chain complex \(C\),
\[
 \operatorname{Tot}N_{\mathrm{simp}}F_{[m]}(C)
       =N(k[\Delta[m]])\otimes_k C,\qquad
 \operatorname{colim}F_{[m]}(C)=C.
\]
The normalized chains of \(\Delta[m]\) contract onto \(k[0]\).
Explicitly, on a nondegenerate ordered simplex prepend the
vertex \(0\); if this creates a repeated vertex, its normalized
value is zero.  The boundary identity is
\(dh+hd=1-\iota\epsilon\).
Tensoring this contraction with \(C\), with the standard total
sign, proves that \(\theta_{F_{[m]}C}\) is a quasi-isomorphism
for arbitrary \(C\).

Projective diagram cofibrations are generated by the maps
\(F_{[m]}(i)\), with \(i\) in a set of generating target
cofibrations.  Both normalized totalization and colimit
preserve colimits; normalization does so by Dold--Kan,
and totalization does so by direct sums.
For a pushout adjoining a cell \(F_{[m]}(A)\to F_{[m]}(B)\),
both resulting maps from the preceding stage are monomorphisms.
Indeed, before the pushout these maps are the tensor of \(A\to B\)
with the degreewise free complex \(N(k[\Delta[m]])\), and
\(A\to B\) itself, respectively.
Their cokernels are respectively
\(N(k[\Delta[m]])\otimes(B/A)\) and \(B/A\).
The comparison on those cokernels is a quasi-isomorphism by
the preceding contraction.  The long exact homology sequences
therefore propagate the assertion about \(\theta\) across the
cell pushout.  Coproducts of cells, transfinite compositions
and retracts preserve it, since direct sums and filtered
colimits of modules are exact.  Every projectively cofibrant
diagram is a retract of such a cell diagram.  Hence
\(\theta_{QY}\) is a quasi-isomorphism.

The objectwise quasi-isomorphism \(QY\to Y\) also induces a
quasi-isomorphism on normalized total complexes.
Indeed normalization is exact in the simplicial direction,
so internal homology commutes with it.  The vertical-homology
comparison and the first-quadrant filtration prove the claim.
We have the natural zigzag
\[
 \operatorname{colim}QY
 \xleftarrow{\ \sim\ }\operatorname{Tot}N_{\mathrm{simp}}QY
 \xrightarrow{\ \sim\ }\operatorname{Tot}N_{\mathrm{simp}}Y,
\]
which proves (SM29.5).  All maps commute with an augmentation
of \(Y\) to a constant diagram; this gives the augmentation
compatibility needed later.

For any small diagram shape \(I\), the left Quillen functor
\(\Lambda\) preserves homotopy colimits in the derived sense:
\[
 \Lambda\operatorname*{hocolim}_I Z
       \simeq\operatorname*{hocolim}_I\Lambda Z.
 \tag{SM29.6}
\]
To verify this precisely, use projective model structures on
the diagram categories.  The objectwise adjunction is
Quillen, since its right adjoint preserves the objectwise
fibrations and trivial fibrations.  Thus \(\Lambda\) sends a
projectively cofibrant replacement to a projectively
cofibrant diagram, and it preserves the objectwise weak
equivalences involved.  It commutes with ordinary colimit
because it is a left adjoint.  Computing both derived
colimits with these cofibrant diagrams gives (SM29.6),
naturally in \(Z\).

\paragraph{4. Construction of a homotopical inverse.}
Let \(\Gamma\) be the Dold--Kan inverse from nonnegative
scalar chain complexes to simplicial scalar modules.
Thus \(N\Gamma C=C\), and every \((\Gamma C)_n\) is a
finite direct sum of the modules \(C_j\), \(0\le j\le n\).
Choose functorial projective-diagram cofibrant replacements
in \(\mathcal M^{\Delta^{\mathrm{op}}}\), and let \(h\) be
the resulting ordinary functor representing homotopy colimit.
Define an ordinary functor on chain complexes by
\[
 \mathcal L(C)=h\bigl(\kappa\Gamma C\bigr)\in\mathcal M.
 \tag{SM29.7}
\]
Equations (SM29.3), (SM29.5), and (SM29.6) give a natural
isomorphism in the derived category
\[
 \begin{split}
 \Lambda\mathcal L(C)
 &\simeq\operatorname*{hocolim}_{\Delta^{\mathrm{op}}}
               \Lambda\kappa\Gamma C\\
 &\simeq\operatorname*{hocolim}_{\Delta^{\mathrm{op}}}
               (\Gamma C)[0]\\
 &\simeq N\Gamma C=C.
 \end{split}
 \tag{SM29.8}
\]
The middle comparison is induced by the actual natural maps
\(e_U\), which are quasi-isomorphisms for every coefficient
module.  No projectivity of the modules in \(C\) is required.

If \(f:C\to C'\) is a quasi-isomorphism, naturality of
(SM29.8) shows that \(\Lambda\mathcal L(f)\) is an
isomorphism in the derived category, hence is a
quasi-isomorphism of complexes.
By the definition of \(W\), \(\mathcal L(f)\in W\).
Thus \(\mathcal L\) descends to a functor
\(D_{\ge0}(k)\to W^{-1}\mathcal M\), with
\(\Lambda\mathcal L\cong1\).
This argument proves homotopicality of \(\mathcal L\);
it does not incorrectly assume that \(\kappa\Gamma\)
takes quasi-isomorphisms to objectwise weak equivalences.

\paragraph{5. The other natural isomorphism.}
There is a functorial projective resolution \(P_\bullet X\to X\)
in the module category \(\mathcal M\).
One explicit choice starts with the coproduct of \(P_x\)
indexed by every element of every \(X(x)\), maps it onto
\(X\), and iterates the same construction on its kernel.
Maps of modules send the indexing elements to their images,
so the construction is functorial.

For a projective object \(P\), the adjunction unit
\[
 \eta_P:P\longrightarrow\kappa\varphi P
\]
belongs to \(W\).  The source \(\Lambda P\) has no positive
homology.  The target has no positive homology by (SM29.3).
On zeroth homology the map is
\(\varphi\eta_P:\varphi P\to\varphi\kappa\varphi P\);
the triangle identity gives
\(\epsilon_{\varphi P}\varphi\eta_P=1\).
The map \(\epsilon_{\varphi P}\) is an isomorphism by
(SM29.3), so \(\varphi\eta_P\) is an isomorphism too.

Apply Dold--Kan inside the abelian category \(\mathcal M\)
to \(P_\bullet X\), and put
\[
 A_X=h\Gamma(P_\bullet X).
\]
Every simplicial term \(\Gamma(P_\bullet X)_n\) is a finite
sum of projectives.  Their units consequently give an
objectwise \(W\)-equivalence of simplicial diagrams
\[
 \Gamma(P_\bullet X)
       \longrightarrow\kappa\Gamma(\varphi P_\bullet X).
\]
Here \(\varphi\) and \(\kappa\) commute with the finite sums
and differential matrices defining \(\Gamma\).
Taking homotopy colimits gives a natural \(W\)-equivalence
\[
 A_X\longrightarrow\mathcal L(\varphi P_\bullet X).
 \tag{SM29.9}
\]

The augmentation of the projective resolution gives a natural
map \(A_X\to X\).  Applying \(\Lambda\) and (SM29.5)--(SM29.6)
identifies its derived image with
\[
 \operatorname{Tot}_{a,b}(\Lambda P_aX)_b
               \longrightarrow\Lambda X.
 \tag{SM29.10}
\]
The identification is compatible with the augmentation by
Step 3.  Additivity of \(\Lambda\) makes it commute with the
finite-sum formula for \(\Gamma\), so normalizing that simplicial
direction gives \(\Lambda P_aX\) in degree \(a\).
For every \(b\), the augmented horizontal sequence
is exact because \(\Lambda\) is exact.
The first-quadrant filtration shows that (SM29.10) is a
quasi-isomorphism.  Therefore \(A_X\to X\) belongs to \(W\).

Finally define
\(T_X=\operatorname{Tot}_{a,b}(\Lambda P_aX)_b\).
The natural augmentations already used in BR-13401 give
\[
 \varphi P_\bullet X
       \xleftarrow{\ \sim\ }T_X
       \xrightarrow{\ \sim\ }\Lambda X.
 \tag{SM29.11}
\]
The left arrow is a quasi-isomorphism by the positive-degree
acyclicity of \(\Lambda\) on projectives, and the right arrow
is (SM29.10).  All comparisons are natural in \(X\).
Since \(\mathcal L\) is homotopical by Step 4, the zigzag
\[
 X\xleftarrow{\ \sim\ }A_X
   \xrightarrow{\ \sim\ }\mathcal L(\varphi P_\bullet X)
   \xleftarrow{\ \sim\ }\mathcal L(T_X)
   \xrightarrow{\ \sim\ }\mathcal L(\Lambda X)
 \tag{SM29.12}
\]
consists of maps in \(W\).
After localization it supplies
\(1\cong\mathcal L\Lambda\), naturally.
Together with (SM29.8), this proves (SM29.1).
Every bicomplex in the proof lies in the first quadrant;
no infinite-diagonal convergence claim has entered.
\(\square\)

\subsection{What this settles and what remains separate}
This establishes the original \emph{underlying} localization
conclusion for the finitely presented-projective-generator
corollary once the local scalar projectivity supplies the
exact adjunction BR-13405.  It works over every scalar ring.
The proof only needs the interval to annihilate injective
objects for the Smith recognition step; its supplied full
hypothesis guarantees this.

The distinction concerning \(G\) is substantive.
A right Quillen functor is applied to a fibrant replacement
when its right derived functor is formed.
The equivalence just proved therefore provides
\(\Lambda G(C^{\mathrm{fib}})\simeq C^{\mathrm{fib}}\)
at that derived level, with the usual cofibrant conventions,
but does not eliminate the test (134.6) for every ordinary
acyclic projective chain object \(P\).
F-SM-24 and F-SM-27 remain positive results for that stronger
ordinary-adjoint assertion.
The general handwritten theorem without the module-category
or comparable presentability assumptions is not claimed here.

\paragraph{References and version-specific locators.}
Beke, T. (2000). Sheafifiable homotopy model categories.
\emph{Mathematical Proceedings of the Cambridge Philosophical
Society, 129}(3), 447--475.
\url{https://doi.org/10.1017/S0305004100004722}.
Author PDF:
\url{https://faculty.uml.edu/tbeke/sheafi.pdf}.
Theorem 1.7, p.~8; Propositions 1.15 and 1.18, pp.~12--13;
diagram-model discussion p.~20.  These are author-PDF pages,
not asserted journal-page translations.

Dugger, D. (2008). \emph{A primer on homotopy colimits}
[Unpublished notes; current 109-page author version retrieved
November 21, 2026].
\url{https://pages.uoregon.edu/ddugger/hocolim.pdf}.
Theorem 19.4 and the following derived-colimit discussion,
p.~91; Proposition 19.9, p.~93.
Earlier shorter versions have different proposition numbers.

The Stacks Project Authors. (2026). \emph{Dold--Kan}.
Theorem 14.24.3, Tag \nolinkurl{019G}; Section 14.24, Tag
\nolinkurl{019D}.
\url{https://stacks.math.columbia.edu/tag/019G}.


% END source_cote1348_realization.tex

% BEGIN source_cote1348_grothendieck.tex
\section{F-SM-30: the abstract interval theorem for Grothendieck categories}

\subsection{Statement and exact historical scope}
Let \(k\) be a commutative ring and let \(\mathcal M\) be a
\(k\)-linear Grothendieck abelian category with enough projectives.
Suppose that the scalar functor
\[
 \kappa:\operatorname{Mod}k\longrightarrow\mathcal M,
 \qquad U\longmapsto\mathbf1\otimes_k U,
\]
is exact, preserves coproducts, and has a left adjoint \(\varphi\).
Write \(F(X)=L\varphi(X[0])\), in nonnegative homological degrees.
Assume the following two conditions:
\begin{enumerate}
\item the canonical derived counit
\(L\varphi(\kappa k)\to k[0]\) is a quasi-isomorphism;
\item every \(X\) has a monomorphism \(X\to A_X\) with
\(L\varphi(A_X)=0\), including degree zero.
\end{enumerate}
The original full interval gives the second condition with
\(A_X=J\otimes X\).  Set
\[
 W=\{f\in\operatorname{Mor}(\mathcal M):F(f)\text{ is an equivalence}\}.
\]
Then the ordinary localization has the required equivalence
\[
 W^{-1}\mathcal M\simeq D_{\ge0}(k),                         \tag{SM30.1}
\]
induced by \(L\varphi\).  In fact the corresponding
\(\infty\)-categories are equivalent.

No exact chain-valued realization \(\Lambda\), no finite-presentation
assumption on projective generators, and no scalar projectivity
assumption on their morphism modules is required.  The additional
hypothesis relative to the bare handwritten conditions (a)--(d) is
that \(\mathcal M\) is Grothendieck: it has a generator and exact
filtered colimits.  Enough projectives is already condition (b).
The exact scalar functor and coproduct preservation follow from its
two adjoints in the original setup, as explained in F-SM-29.
We do not infer AB5 or existence of a generator from the handwriting.
Thus (SM30.1) closes this broad and precisely stated case, while
leaving a genuinely more general abelian-category assertion distinct.

\subsection{Known inputs and coherence convention}
We use the following known theorems in their stated forms.
\begin{itemize}
\item Beke (2000), Proposition 3.10, author-PDF pp.~24--25:
Grothendieck abelian categories are locally presentable.
Proposition 1.12, p.~11, gives a set generating all monomorphisms
under pushouts and transfinite composition in this setting.
Theorem 1.7 and Propositions 1.15, 1.18 provide Smith recognition
and the accessibility/solution-set implications.
\item Lurie, \emph{Higher Algebra}, Example 1.3.3.4, p.~98,
Theorem 1.3.4.4 and Proposition 1.3.4.6, p.~109:
a right-exact functor between
abelian categories with enough projectives has a genuine derived
\(\infty\)-functor, computed by projective complexes, and the
derived \(\infty\)-category is the localization at quasi-isomorphisms.
\item Lurie, \emph{Higher Algebra}, Proposition 1.3.4.22,
p.~118, and Propositions 1.3.4.24--25, p.~119: a combinatorial
model category presents a presentable \(\infty\)-category;
model homotopy colimits compute its colimits, and diagrams
indexed by ordinary categories admit strict diagram representatives.
\item Lurie, \emph{Higher Topos Theory}, Proposition 4.4.2.7,
p.~297: a functor preserving pushouts and small coproducts
preserves small colimits.  The theorem is applied for each
regular size bound, within a fixed ambient universe.
\end{itemize}
We also use Dold--Kan, Stacks Project Theorem 14.24.3,
Tag \nolinkurl{019G}, and the projective diagram model theorem.
The normalized-total-complex calculation needed below has the
complete proof in F-SM-29, Step 3.  Its extension from modules to
a Grothendieck abelian category is explained explicitly below.
These background theorems are inputs, not claimed new proofs.

All references to \(F\) as a functor of \(\infty\)-categories
mean the actual derived functor followed by relative localization.
No functor defined only on homotopy categories is used to infer
preservation of coherent colimits.

\subsection{Complete proof}

\paragraph{1. Projectives, filtered colimits, and accessibility.}
Choose a generator \(G_0\) of \(\mathcal M\) and a projective
epimorphism \(P\to G_0\).  Then \(P\) is a projective generator:
\(\operatorname{Hom}(P,-)\) is exact, and if an object receives
no nonzero map from \(P\), it receives none from \(G_0\) and
is zero.  The evaluation map
\(\bigoplus_{\operatorname{Hom}(P,X)}P\to X\) is epic;
otherwise projectivity and generation give a map to its nonzero
cokernel which would lift and contradict evaluation.
This construction and successive kernels supply functorial
projective resolutions.  Additivity of this resolution functor
is not required.  The adjunction and exactness of \(\kappa\)
show that \(\varphi\) takes projectives to projective scalar
modules.  Coproducts of projectives are projective, since a
family of maps to an epimorphism can be lifted separately.

We verify, rather than assume, that all \(L_j\varphi\) commute
with filtered colimits.  Let \(I\) be a small filtered category.
For \(i\in I\) and projective \(Q\in\mathcal M\), the free
diagram on \(Q\) at \(i\) has value
\[
 (i_!Q)(j)=\bigoplus_{I(i,j)}Q.
\]
It is projective in \(\mathcal M^I\), by its adjunction to
exact evaluation.  These diagrams give enough projectives.
Every projective diagram is a retract of a coproduct of them.
Evaluation therefore takes projective diagrams to projectives,
and so does colimit: \(\operatorname{colim}_I i_!Q=Q\),
because the category \(i/I\) has an initial object.

Resolve an arbitrary diagram \(X:I\to\mathcal M\) by such
projective diagrams \(Q_\bullet\to X\).
Evaluation gives a projective resolution of each \(X_i\).
AB5 says colimit is exact, so \(\operatorname{colim}Q_\bullet\)
is a projective resolution of \(\operatorname{colim}X\).
The left adjoint \(\varphi\) commutes with colimits, and
filtered colimits of scalar modules are exact.  Consequently
\[
 L_j\varphi(\operatorname{colim}_I X_i)
   \simeq\operatorname{colim}_I L_j\varphi(X_i)
       \quad(j\ge0).                                    \tag{SM30.2}
\]
The isomorphism is canonical: it is induced by the diagram's
colimit maps, and comparison of projective resolutions preserves
it.  The same argument for a discrete diagram, using exact
coproducts in a Grothendieck category, proves
\[
 F\Bigl(\bigoplus_i X_i\Bigr)\simeq\bigoplus_i F(X_i).
                                                               \tag{SM30.3}
\]

The functor \(X\mapsto(L_j\varphi X)_{j\ge0}\) to graded
scalar modules is accessible, by (SM30.2) and local
presentability.  The isomorphisms in the target arrow category
form an accessible, accessibly embedded subcategory.
Beke Proposition 1.18 makes its inverse image \(W\) accessible
and accessibly embedded; Proposition 1.15 gives the solution-set
condition for any set of candidate generating cofibrations.

\paragraph{2. The model structure is constructed directly from \(L\varphi\).}
Use the set \(\mathcal I\) furnished by Beke Proposition 1.12,
so \(\operatorname{cof}(\mathcal I)\) is precisely all
monomorphisms.  Its hypotheses hold: subobjects have effective
unions in an abelian category, and AB5 preserves monomorphisms
under transfinite composition.  Closure of \(W\) under retracts
and two-out-of-three follows from its definition.

For a monomorphism \(u:A\to B\), the derived triangle of
\(0\to A\to B\to\operatorname{coker}u\to0\) gives
\[
 u\in W\quad\Longleftrightarrow\quad F(\operatorname{coker}u)=0.
                                                               \tag{SM30.4}
\]
Pushout retains this cokernel.  A transfinite composition of
monomorphisms in \(W\) stays in \(W\): use (SM30.2) at every
limit ordinal and transfinite induction on the maps from the
initial stage.  Hence the trivial cofibrations are stable under
the operations required in Smith's condition (c2).

Suppose \(u:B\to C\) has the right lifting property against
\(\mathcal I\), hence against every monomorphism.  Lifting
against \(0\to C\) gives a section, and its kernel \(K\) is
injective: a map into \(K\) extends across any monomorphism
by a lifting square with bottom map zero into \(C\).
The assumed monomorphism \(K\to A_K\) splits by injectivity.
Thus \(F(K)\) is a retract of \(F(A_K)=0\).
Since \(B\simeq C\oplus K\), the split epimorphism \(u\)
belongs to \(W\).  This is Smith's condition (c1).
All conditions now hold, so there is a combinatorial model
structure on \(\mathcal M\) with weak equivalences \(W\)
and cofibrations all monomorphisms.  Every object is cofibrant.
Let \(\mathcal C\) be its underlying \(\infty\)-category.

We use the analogous all-monomorphism, quasi-isomorphism model
on \(\operatorname{Ch}_{\ge0}(\mathcal M)\).
For clarity its existence does not result from restricting an
unbounded model.  The same recognition applies directly:
this chain category is Grothendieck, homology preserves filtered
colimits, and a monomorphism is a quasi-isomorphism exactly
when its cokernel is acyclic.  An injective chain object \(I\)
embeds into the contractible complex \(D^1(\mathbf Z)\otimes I\)
by the degree-zero inclusion; this embedding splits, so \(I\)
is acyclic.  The preceding injective-kernel argument proves
the remaining lifting condition.  This also gives the model
for \(\operatorname{Ch}_{\ge0}(k)\).
Their underlying \(\infty\)-categories are the connective
derived categories, denoted \(\mathcal D_{\ge0}(\mathcal M)\)
and \(\mathcal D_{\ge0}(k)\).

\paragraph{3. Coherent derived functors and all-scalar normalization.}
The right-exact functor \(\varphi\) gives the genuine derived
\(\infty\)-functor
\[
 \mathbb L\varphi:\mathcal D_{\ge0}(\mathcal M)
                     \longrightarrow\mathcal D_{\ge0}(k).
\]
It is the restriction of the exact right-t-exact functor of
\emph{Higher Algebra}, Example 1.3.3.4, and is computed by
applying \(\varphi\) to projective complexes.  This also
identifies its value on \(X[0]\) with the usual derived functor.
Exactness implies preservation of pushouts, which stay connective.
Projective resolutions and exact coproducts give preservation
of coproducts, with the actual comparison maps of the derived
functor.  HTT Proposition 4.4.2.7 therefore shows that
\(\mathbb L\varphi\) preserves all small connective colimits.

The exact functor \(\kappa\), preserving coproducts, induces
a functor on connective derived \(\infty\)-categories.
It too preserves pushouts and coproducts, hence all small
colimits by the same proposition.  There is the canonical
derived counit \(\mathbb L\varphi\kappa\to\mathrm{id}\).
Concretely, on a complex it is represented by
\(\varphi P(\kappa C)\to\varphi\kappa C\to C\), using
a functorial projective replacement and the original counit.
These are natural maps of complexes before localization, so
the displayed counit is a coherent natural transformation.
It is an equivalence on \(k[0]\) by hypothesis, hence on
free scalar modules and their projective retracts.
Every scalar module \(U\) is the realization of
\(\Gamma(Q_\bullet)[0]\), where \(Q_\bullet\to U\) is a
projective scalar resolution and \(\Gamma\) is Dold--Kan's
inverse normalization.  Each \(\Gamma(Q)_n\) is a finite sum
of projective modules.  Both functors in this counit preserve
that realization.  Thus the \emph{canonical} map
\[
             F(\kappa U)\longrightarrow U[0]
 \quad\hbox{is an equivalence for every }U.                \tag{SM30.5}
\]

Here and below the realization identity used is
\[
 \operatorname*{colim}_{\Delta^{\rm op}}\Gamma(C)_n[0]
                          \simeq C                       \tag{SM30.6}
\]
in the appropriate connective derived category, naturally
including its augmentation.  The proof in F-SM-29, Step 3,
works for a Grothendieck category without change to its logical
steps: normalization is an exact equivalence; coproducts and
filtered colimits are exact; the normalized chains of a simplex
are finite free abelian in each degree; their vertex contraction
acts by finite sums of identity maps on any coefficient object.
The free-diagram cell proof then gives
\(\operatorname{hocolim}Y\simeq\operatorname{Tot}NY\).
Cell cokernels, transfinite stages and retracts are handled as
there, using exact sums and AB5, and every total degree has
finitely many bidegrees.  Applying this to \(Y=\Gamma(C)[0]\)
proves (SM30.6), with its natural augmentation.

Compose \(\mathcal M\to\mathcal D_{\ge0}(\mathcal M)\),
\(X\mapsto X[0]\), with \(\mathbb L\varphi\).
By definition this coherent functor takes \(W\) to
equivalences.  The universal property of relative
\(\infty\)-localization therefore gives
\[
 \overline F:\mathcal C\longrightarrow\mathcal D_{\ge0}(k).
                                                               \tag{SM30.7}
\]
This supplies the required coherence before any colimit claim
about \(\overline F\) is made.

\paragraph{4. The localized functor preserves all colimits.}
Since every object of \(\mathcal M\) is cofibrant, its ordinary
coproducts compute model homotopy coproducts.  Equation (SM30.3)
therefore proves that \(\overline F\) preserves coproducts,
including the empty one.

Represent any pushout diagram in \(\mathcal C\) by a
cofibrant span in the model category; the strictification and
colimit comparisons are HA Propositions 1.3.4.24--25.
We may thus suppose that both maps in
\(B\leftarrow A\to C\) are monomorphisms and its ordinary
pushout is \(P=B\amalg_A C\).  There is a short exact sequence
\[
 0\longrightarrow A\xrightarrow{(u,-v)}B\oplus C
                     \longrightarrow P\longrightarrow0.
                                                               \tag{SM30.8}
\]
It exhibits the degree-zero square as a pushout in
\(\mathcal D_{\ge0}(\mathcal M)\).  Applying the actual
derived functor \(\mathbb L\varphi\) gives the pushout
square in \(\mathcal D_{\ge0}(k)\), with the same comparison
maps as (SM30.7).  Hence \(\overline F\) preserves pushouts.
Its source is presentable by HA Proposition 1.3.4.22.
HTT Proposition 4.4.2.7 now gives preservation of all small
colimits, in particular simplicial realizations.

\paragraph{5. Construction of both inverse comparisons.}
For a nonnegative scalar complex \(C\), set
\[
 L(C)=\operatorname*{colim}_{\Delta^{\rm op}}
                        \kappa\Gamma(C)_n\quad\text{in }\mathcal C.
                                                               \tag{SM30.9}
\]
Functorial projective diagram replacement may be used to choose
an ordinary representative of this colimit.  By Step 4,
(SM30.5), and (SM30.6), there is a natural equivalence
\[
                      \overline F L(C)\simeq C.            \tag{SM30.10}
\]
If \(f:C\to C'\) is a quasi-isomorphism, this comparison
shows that the chosen map representing \(L(f)\) lies in \(W\).
Thus \(L\) descends to a functor
\(\mathcal D_{\ge0}(k)\to\mathcal C\).
This argument does not assume that \(\kappa\Gamma(f)\)
is an objectwise weak equivalence.

Take the functorial projective resolution \(P_\bullet X\to X\)
from Step 1.  For a projective object \(Q\), the adjunction
unit \(\eta_Q:Q\to\kappa\varphi Q\) belongs to \(W\).
Indeed its source under \(F\) is \(\varphi Q[0]\); its
target is identified with the same object by (SM30.5); the
triangle identity says the resulting map is the identity.
In every simplicial degree \(\Gamma(P_\bullet X)_n\) is
projective.  The levelwise units consequently give equivalences
in \(\mathcal C\), and their realization gives
\[
 A_X:=\operatorname*{colim}_{\Delta^{\rm op}}
                    \Gamma(P_\bullet X)_n
       \xrightarrow{\ \sim\ }L(\varphi P_\bullet X).
                                                               \tag{SM30.11}
\]
The augmentation also gives \(a_X:A_X\to X\) in \(\mathcal C\).
We check this particular map, rather than just comparing its
source and target abstractly.  Applying \(\overline F\) and
Step 4 identifies it with the augmentation
\[
 \operatorname*{colim}_{\Delta^{\rm op}}
       \mathbb L\varphi(\Gamma(P_\bullet X)_n[0])
                   \longrightarrow\mathbb L\varphi(X[0]).
\]
Since \(\mathbb L\varphi\) preserves colimits, (SM30.6)
identifies this map with the image of the quasi-isomorphism
\(P_\bullet X\to X[0]\); hence it is an equivalence.
On ordinary representatives the augmentation map therefore
lies in \(W\), so \(a_X\) itself is an equivalence in
\(\mathcal C\).  Finally \(\varphi P_\bullet X\) is the
functorial representative of \(F(X)\).  We have natural
equivalences
\[
              X\xleftarrow{\ \sim\ } A_X
                    \xrightarrow{\ \sim\ }L(\overline F X).
                                                               \tag{SM30.12}
\]
The augmentations, units and resolution maps used to construct
this span are natural.  Equations (SM30.10)--(SM30.12) prove
that \(L\) and \(\overline F\) are mutually inverse
\(\infty\)-functors.  Passing to homotopy categories gives
(SM30.1). \(\square\)

\subsection{What the proof removes and what it does not remove}
The proof constructs accessibility and the model structure
directly from \(L\varphi\), rather than adding an exact
chain-valued realization as an assumption.  It works for every
commutative scalar ring and for projective generators which
need not be finitely presented.  For the underlying localization,
the full tensor interval is used only through the stated
acyclic-monomorphism consequence.

The AB5 step is substantive.  Merely observing that an
endofunctor such as \(V\mapsto V^{**}/V\) is exact over a field
does not make its square-zero module category Grothendieck.
For the intake's three-sort proposal, an asserted periodic
formula \(L_{3n}\varphi(\kappa V)=(V^{**}/V)^{\circ n}\)
would contradict (SM30.5) for infinite-dimensional \(V\),
so all of the present hypotheses cannot simultaneously hold
there.  This is a diagnostic conditional implication, not a
completed reconstruction of that counterexample: its symmetric
tensor coherence and full interval still require their absent
proof.  No conclusion about that missing proof is inferred.

There is also no ordinary right adjoint \(G\) claimed here.
Whenever one chooses an exact chain-valued adjunction separately,
its raw counit remains the stronger and separately identified
test in BR-13403.  The present theorem concerns the underlying
localization in the explicit Grothendieck range.

\subsection{Primary references}
\begin{itemize}
\item Beke, T. (2000). Sheafifiable homotopy model categories.
\emph{Mathematical Proceedings of the Cambridge Philosophical
Society, 129}(3), 447--475.
\url{https://doi.org/10.1017/S0305004100004722}.
Author PDF: \url{https://faculty.uml.edu/tbeke/sheafi.pdf}.
\item Lurie, J. (2009). \emph{Higher topos theory}
(Annals of Mathematics Studies, Vol.~170). Princeton University
Press. \url{https://www.math.ias.edu/~lurie/papers/HTT.pdf}.
Exact input: Proposition 4.4.2.7, printed p.~297.
\item Lurie, J. (2017, September 18). \emph{Higher algebra}.
Author manuscript.
\url{https://people.math.harvard.edu/~lurie/papers/HA.pdf}.
Exact inputs: Example 1.3.3.4, p.~98; Theorem 1.3.4.4 and
Proposition 1.3.4.6, p.~109; Proposition 1.3.4.22, p.~118; and Propositions
1.3.4.24--25, p.~119.
\item The Stacks Project Authors. (2026). \emph{The Stacks
Project}. Theorem 14.24.3, Tag \nolinkurl{019G}.
\url{https://stacks.math.columbia.edu/tag/019G}.
\end{itemize}

% END source_cote1348_grothendieck.tex

% BEGIN source_cote1348_no_ab5.tex
\section{F-SM-31: the full interval hypotheses do not force AB5 or presentability}

\subsection{Statement}
For every field \(k\), there is a complete and cocomplete
\(k\)-linear abelian symmetric monoidal category \(\mathcal M\)
with enough projectives satisfying all the precise adjoint,
unit-normalization, and full tensor-interval conditions of
cote 134-8, footers 47--48, which is neither AB5 nor locally
presentable.  Its tensor product is exact and preserves colimits
in both variables,
and the scalar functor is exact and strong symmetric monoidal.
Thus the missing Grothendieck hypothesis in F-SM-30 cannot be
deduced from those original conditions.

This is a counterexample to that proposed implication, not to
the original localization conclusion.  The latter holds in the
example, as is proved at the end.  No double-dual, ultrapower,
or unreconstructed tensor-coherence assertion is used.

\subsection{Construction and complete coherence check}
Let \(\mathcal V=\operatorname{Vect}_k\) and set
\[
 \mathcal M=\operatorname{Ch}_{\ge0}(k)\times\mathcal V^{\rm op}.
                                                               \tag{SM31.1}
\]
We denote its objects by \((C,A)\), where \(A\) is an ordinary
vector space but belongs to the opposite-category factor.
Thus a morphism \((C,A)\to(D,B)\) consists of a chain map
\(f:C\to D\) and a linear map \(g:B\to A\).
The category is abelian, complete and cocomplete, componentwise.
Every object of \(\mathcal V^{\rm op}\) is projective because
every monomorphism of vector spaces splits.  The chain category
has the usual projective generators \(k[0],D^n(k)\), \(n\ge1\).
Consequently \(\mathcal M\) has enough projectives.
It even has a projective generator: combine those chain
generators with \((0,k)\), and take their coproduct.
The latter generates the opposite factor because linear
functionals to \(k\) separate nonzero maps of vector spaces.

For a vector space \(U\), the canonical scalar action on
\(\mathcal V^{\rm op}\) is
\[
             U\odot A=\operatorname{Hom}_k(U,A),
             \quad\text{regarded as an object of }\mathcal V^{\rm op}.
                                                               \tag{SM31.2}
\]
It is covariant in \(U\) after the opposite-category convention.
The canonical currying isomorphism gives
\((U\otimes V)\odot A\simeq U\odot(V\odot A)\),
and \(k\odot A\simeq A\).
Define the tensor product by
\[
 (C,A)\boxtimes(D,B)=
 \left(C\otimes_kD,
    \operatorname{Hom}_k(C_0,B)\oplus
    \operatorname{Hom}_k(D_0,A)\right).                    \tag{SM31.3}
\]
The finite sum in the second component is the same vector
space for products and coproducts.  For morphisms
\((f,g):(C,A)\to(C',A')\) and
\((h,j):(D,B)\to(D',B')\), the map on this component,
in the ordinary vector-space direction, is
\[
 \begin{aligned}
 \operatorname{Hom}(C'_0,B')&\longrightarrow\operatorname{Hom}(C_0,B),
                 &q&\longmapsto j\,q\,f_0,\\
 \operatorname{Hom}(D'_0,A')&\longrightarrow\operatorname{Hom}(D_0,A),
                 &q&\longmapsto g\,q\,h_0.
 \end{aligned}                                             \tag{SM31.4}
\]
This defines a bifunctor additive in each variable.
Its unit is \(\mathbf1=(k[0],0)\).
In fact the tensor is exact in each variable: evaluation at
chain degree zero is exact, all vector-space short exact
sequences split, and tensor and either variance of Hom over
a field preserve their exactness, with the opposite variance
accounted for as in (SM31.4).  It preserves coproducts as well.
For example, their second component satisfies
\[
 \operatorname{Hom}\Bigl(\bigoplus_i(C_i)_0,B\Bigr)
 \oplus\operatorname{Hom}\Bigl(D_0,\prod_iA_i\Bigr)
 \simeq\prod_i\bigl(\operatorname{Hom}((C_i)_0,B)
                 \oplus\operatorname{Hom}(D_0,A_i)\bigr).
\]
This is the required coproduct identity in the opposite
factor.  Since any colimit in an abelian category is obtained
from coproducts and a cokernel, tensor preserves all colimits
in each variable.  In particular the example also satisfies
the stronger right-exactness conventions for a tensor product.

We specify and check all coherence maps.  In the first component
use the standard associator, unit maps, and signed symmetry of
nonnegative chain complexes.  In the second component an
\(m\)-fold tensor of \((C_i,A_i)\), in any parenthesization,
has the canonical normal form
\[
 \bigoplus_{i=1}^m
  \operatorname{Hom}_k\left(
      \bigotimes_{j\ne i}(C_j)_0,\ A_i\right),             \tag{SM31.5}
\]
where the tensor factors are in their original order.
This follows by induction from (SM31.3), distributivity over
finite sums, and currying.  Define every associator by the
canonical identifications with this normal form.  The two
pentagon composites coincide because on each labelled summand
they are the same currying map on the same ordered tensor
factors.  Define the symmetry by swapping the two components
in (SM31.3) and using the usual interchange of degree-zero
tensor factors in (SM31.5).  The hexagon composites induce the
same permutation of these factors and the same permutation
of the summand labels, and hence coincide.  The triangle and
unit axioms reduce to \(\operatorname{Hom}_k(k,A)=A\).
The first-component signed symmetry contributes no sign to
the second component, which uses only chain degree zero.
These identifications are natural, by (SM31.4).
Thus this is a symmetric monoidal category with all the
required coherence, not merely a formula on objects.

\subsection{Scalar adjoints and derived normalization}
The actual scalar copower of the tensor unit is
\[
 \kappa(U)=(U[0],0).
\]
Indeed the second component of any coproduct of copies of the
zero object remains zero, and scalar presentations give the
displayed first component.  Formula (SM31.3) makes \(\kappa\)
strong symmetric monoidal.  It is exact and preserves all
small limits and colimits, by the componentwise constructions.
Its two adjoints are
\[
 \varphi(C,A)=H_0(C),\qquad
 \psi(C,A)=C_0,
 \qquad \varphi\dashv\kappa\dashv\psi.                    \tag{SM31.6}
\]
For the left adjunction, a chain map \(C\to U[0]\) is exactly
a linear map \(H_0(C)\to U\); for the right adjunction, a
chain map \(U[0]\to C\) is exactly a linear map \(U\to C_0\).
The opposite-vector component contributes only its unique
zero map in either calculation.  In particular
\(\psi=\operatorname{Hom}_{\mathcal M}(\mathbf1,-)\),
as required in the original setup.

Let \(\Lambda(C,A)=C\).  This is an exact functor to
\(\operatorname{Ch}_{\ge0}(k)\), with \(H_0\Lambda=\varphi\).
A projective object of \(\mathcal M\) has projective first
component in the chain category, whose positive homology
vanishes.  The fully proved finite-diagonal lemma BR-13401
therefore gives, naturally,
\[
                     L\varphi(C,A)\simeq C.                \tag{SM31.7}
\]
The canonical map on the unit is \(k[0]\to k[0]\), the
identity.  This verifies the original unit normalization.

\subsection{The complete tensor interval, including degree zero}
Take \(J=(D^1(k),0)\), where \(D^1(k)\) has identity
differential from degree one to degree zero, and let
\(\alpha:\mathbf1\to J\) be the degree-zero inclusion.
For every \(X=(C,A)\), there is the natural identification
\[
                 J\boxtimes X=(D^1(k)\otimes C,A).          \tag{SM31.8}
\]
Under it, \(\alpha\boxtimes1_X\) has second component the
identity of \(A\).  Its first component is the degreewise
monomorphism \(C\to D^1(k)\otimes C\).
In degree \(n\), write the latter tensor complex as
\(C_n\oplus C_{n-1}\).  Its differential and contraction are
\[
 d(x,y)=(d_Cx+y,-d_Cy),\qquad h(x,y)=(0,x).
\]
Direct substitution gives \(dh+hd=1\).
Consequently \(\alpha\boxtimes1_X\) is monic for every
\(X\), and (SM31.7) gives
\[
                    L\varphi(J\boxtimes X)=0
\]
in every degree, including degree zero.  This verifies the
full original interval condition, without a finiteness
restriction on \(C\) or \(A\).

\subsection{Failure of AB5 and of local presentability}
We give explicit arguments for both failures.
Let \(I=\mathbb N\), let \(V=k^{(I)}\), and for a finite
subset \(S\subset I\) put
\(V_S=k^{(I\setminus S)}\subset V\).
There is an inverse system of short exact vector-space
sequences
\[
          0\longrightarrow V_S\longrightarrow V
                  \longrightarrow V/V_S\longrightarrow0.
\]
Its limits have
\(\lim_S V_S=0\), \(\lim_S V=V\), and
\(\lim_S V/V_S=k^I\).
The resulting map \(V\to k^I\) is not surjective, since
the constant family \((1)_{i\in I}\) is not finitely supported.
Passing to opposite categories turns this into a filtered
system of short exact sequences whose colimit is not exact.
Thus \(\mathcal V^{\rm op}\), and hence \(\mathcal M\),
fails AB5.

For local presentability fix any infinite regular cardinal
\(\lambda\) and any nonzero vector space \(A\).
Let \(I\) be a set of cardinality \(\lambda\), and now
index the same tail subspaces \(V_S=k^{(I\setminus S)}\)
by all subsets \(S\subset I\) of cardinality less than
\(\lambda\).  This indexing poset is \(\lambda\)-filtered,
and its colimit in \(\mathcal V^{\rm op}\) is again zero,
because the intersection of the tail subspaces is zero.
Choose \(0\ne a\in A\).  The map \(V\to A\) taking every
basis vector to \(a\) restricts nontrivially to every \(V_S\).
It therefore defines a nonzero element of
\[
 \operatorname*{colim}_{|S|<\lambda}
          \operatorname{Hom}_k(V_S,A),
\]
whereas \(\operatorname{Hom}_{\mathcal V^{\rm op}}(A,0)=0\).
Thus no nonzero object of \(\mathcal V^{\rm op}\) is
\(\lambda\)-presentable, for any infinite regular \(\lambda\).

If \((C,A)\) is \(\lambda\)-presentable in \(\mathcal M\),
apply the preceding filtered diagram in its second factor,
with zero first factor.  The same calculation forces \(A=0\).
Filtered colimits of objects with zero second factor still
have zero second factor.  Hence no family of
\(\lambda\)-presentable objects can generate \(\mathcal M\)
under \(\lambda\)-filtered colimits: for example \((0,k)\)
cannot be so obtained.  Since this holds for every regular
\(\lambda\), \(\mathcal M\) is not locally presentable.

\subsection{The localization conclusion nevertheless holds here}
By (SM31.7), \(W\) consists exactly of the maps whose first
components are quasi-isomorphisms.  Projection
\(\Lambda(C,A)=C\) has a section \(G(C)=(C,0)\), with
\(\Lambda G=\mathrm{id}\).  The natural map
\((C,A)\to(C,0)\) belongs to \(W\), since its first component
is the identity.  These functors and this natural weak-equivalence
map give mutually inverse functors after localization.  Therefore
\[
                W^{-1}\mathcal M\simeq D_{\ge0}(k).
\]
The example consequently rules out only the proposed deduction
of AB5 or presentability from the original predicates.  It does
not decide the bare abstract localization assertion in general.
\(\square\)

\paragraph{Primary-source scope.}
The original predicates checked here are those of
Grothendieck, A. (n.d.), \emph{Pursuing stacks: Linearization
of homotopy types}, cote 134-8, manuscript footers 47--48,
\url{https://grothendieck.umontpellier.fr/134-8.pdf}.
The construction and all checks above are supplied in full;
the derived-functor identification invokes only the already
proved BR-13401 and BR-13402.

% END source_cote1348_no_ab5.tex

% BEGIN source_cote1348_no_ab4.tex
\section{F-SM-43: the full original interval hypotheses do not force AB4}

\subsection{Statement and exact scope}
For every field \(k\), there exists a complete and cocomplete
\(k\)-linear abelian symmetric monoidal category \(\mathcal M\)
with enough projectives satisfying the scalar adjunction,
canonical unit-normalization, and full tensor-interval assumptions
of cote~134-8, footers~46--48, for which countable coproducts of
monomorphisms are not always monic.  The tensor is right exact
and preserves colimits in each variable; it is even closed.
The distinguished interval and its quotient by the unit are flat.
Thus none of these extra features forces AB4.

Nevertheless, the canonical derived scalar functor in this example
preserves arbitrary coproducts, its canonical counit is an
equivalence on every scalar module, and its ordinary localization
is \(D_{\ge0}(k)\).  This is a counterexample only to automatic
AB4, not to the original localization conclusion.

\subsection{Construction and complete proof}
Let \(\mathcal A\) be any complete and cocomplete \(k\)-linear
abelian category with enough projectives.  We shall at the end
take \(\mathcal A\) to be the explicit category of F-SM-34,
which fails countable AB4.  No tensor on \(\mathcal A\) is assumed.

For a vector space \(U\) and \(A\in\mathcal A\), let
\(U\odot A\) denote the canonical scalar copower, characterized by
\[
 \operatorname{Hom}_{\mathcal A}(U\odot A,B)
  =\operatorname{Hom}_k\bigl(U,
               \operatorname{Hom}_{\mathcal A}(A,B)\bigr).
                                                               \tag{SM43.1}
\]
It exists: after choosing a basis of \(U\), take the corresponding
coproduct of copies of \(A\).  The displayed universal property
makes the construction independent of that choice and functorial.
Currying in vector spaces and Yoneda give the canonical maps
\[
 (U\otimes V)\odot A\simeq U\odot(V\odot A),
       \qquad k\odot A\simeq A.                              \tag{SM43.2}
\]
Every iterated such map is the same iterated currying map under
(SM43.1), proving its associativity and unit coherence.
The action preserves colimits in each variable: in \(U\) this
follows from (SM43.1); for fixed \(U\), its basis description is
a coproduct functor on \(\mathcal A\), which preserves colimits.
No claim of exactness in \(A\) is made.

Set
\[
 \mathcal M=\operatorname{Ch}_{\ge0}(k)\times\mathcal A,
 \qquad
 (C,A)\boxtimes(D,B)=
  \bigl(C\otimes_k D, C_0\odot B\oplus D_0\odot A\bigr).
                                                               \tag{SM43.3}
\]
The product category is abelian, complete and cocomplete.
The tensor formula is a bifunctor, using the functorial copower
maps on its second component.  It is additive and preserves
colimits in each variable, since degree-zero evaluation of
connective chain complexes preserves colimits.  Hence it is
right exact in each variable.

Here are all its coherence maps.  The first component uses the
usual chain tensor, associator and Koszul symmetry.  The second
component of any parenthesization of an \(m\)-fold product has
the canonical normal form
\[
 \bigoplus_{i=1}^{m}
       \left(\bigotimes_{j\ne i}(C_j)_0\right)\odot A_i.
                                                               \tag{SM43.4}
\]
Define associators through this normal form, using (SM43.2).
Both pentagon composites are the same currying map on each
labelled summand.  Define the symmetry by permuting the summand
labels and the degree-zero scalar factors; its hexagon is the
equality of the corresponding permutations.  These maps are
natural by (SM43.1).  Their unit is
\(\mathbf1=(k[0],0)\), and the triangle reduces to
\(k\odot A=A\).  No Koszul sign enters the second component,
whose scalar factors all have degree zero.  This verifies the
full symmetric monoidal structure without a tensor on
\(\mathcal A\).

For completeness, closedness also follows directly.  The scalar
cotensor \(U\pitchfork B\) exists by products in \(\mathcal A\),
with
\[
 \operatorname{Hom}_{\mathcal A}(A,U\pitchfork B)
  =\operatorname{Hom}_k\bigl(U,
                \operatorname{Hom}_{\mathcal A}(A,B)\bigr).
\]
Let \([C,D]_{\ge0}\) be the usual connective internal Hom of
chain complexes: its positive-degree terms are those of the
full Hom complex, and its degree-zero term is the kernel of
the differential out of degree zero.  Tensor--Hom currying
identifies maps into this complex with chain maps from a tensor
product.  If \(D^1(U)=(U\xrightarrow{1}U)\) in degrees one
and zero, then a chain map \(E\to D^1(U)\) is determined
by an arbitrary map \(E_0\to U\); its degree-one component
is that map composed with \(d_E\).  Consequently an internal
Hom for (SM43.3) is
\[
 [(C,A),(D,B)]_{\mathcal M}=
 \left([C,D]_{\ge0}\oplus
          D^1\bigl(\operatorname{Hom}_{\mathcal A}(A,B)\bigr),
                         C_0\pitchfork B\right).             \tag{SM43.5}
\]
Indeed maps from \((E,F)\) into its three indicated summands
are exactly the chain map \(E\otimes C\to D\) and the two
maps \(E_0\odot A\to B\), \(C_0\odot F\to B\) specifying
a map \((E,F)\boxtimes(C,A)\to(D,B)\).

The chain category has enough projectives, generated by
\(k[0]\) and the projective disks \(D^n(k)\), \(n\ge1\).
Together with enough projectives in \(\mathcal A\),
componentwise projective epimorphisms prove that
\(\mathcal M\) has enough projectives.
The actual scalar copower of \(\mathbf1\) is
\[
 \kappa U=(U[0],0),\qquad
 \varphi(C,A)=H_0(C),\qquad \psi(C,A)=C_0,
                   \qquad\varphi\dashv\kappa\dashv\psi.
                                                               \tag{SM43.6}
\]
The two adjunctions follow by the ordinary factorization of
chain maps to degree zero through \(H_0\), and by evaluation
for maps from a degree-zero complex.  Maps involving the zero
second component contribute the one zero map.  The scalar
functor is exact, preserves all limits and colimits, and is
strong symmetric monoidal by (SM43.3).

Let \(\Lambda(C,A)=C\).  There is a natural canonical
identification
\[
                  L\varphi(C,A)\simeq C.                    \tag{SM43.7}
\]
To check it, resolve \((C,A)\) by projective objects
\((P_r,Q_r)\) of \(\mathcal M\).  Each \(P_r\) is
projective in the chain category, hence has zero positive
homology and projective \(H_0\).  The bicomplex \((P_r)_n\)
has a first-quadrant total complex.  Its resolution-direction
augmentation to \(C\) is a quasi-isomorphism, because it is
exact in that direction in every chain degree.  Its other
augmentation to \(H_0(P_\bullet)\) is a quasi-isomorphism,
because each \(P_r\) has zero positive homology.  Every total
degree is a finite sum, so this is precisely the finite-diagonal
argument of BR-13401, with no AB4 assumption on
\(\mathcal A\).  Applying \(\varphi\) to the projective
resolution gives \(H_0(P_\bullet)\), which proves (SM43.7)
with its natural comparisons.  On \(\kappa U\), the canonical
counit is consequently the identity of \(U[0]\), for every
vector space \(U\).

Take \(J=(D^1(k),0)\), and let
\(\alpha:\mathbf1\to J\) be the degree-zero inclusion.
Then
\[
 J\boxtimes(C,A)=(D^1(k)\otimes C,A),\qquad
 (J/\mathbf1)\boxtimes(C,A)=(k[1]\otimes C,0).              \tag{SM43.8}
\]
The map \(\alpha\boxtimes1\) is the identity in the second
component and degreewise monic in the first.  Write the latter
chain complex in degree \(n\) as \(C_n\oplus C_{n-1}\).
Its differential and a contraction are
\[
 d(x,y)=(d_Cx+y,-d_Cy),\qquad h(x,y)=(0,x),
                     \qquad dh+hd=1.                         \tag{SM43.9}
\]
Thus (SM43.7) gives \(L\varphi(J\boxtimes X)=0\) in
every degree, including zero, for every \(X\).  Moreover both
functors in (SM43.8) are exact.  This verifies the full interval
and even the optional flat-quotient feature, independently of
exactness of arbitrary coproducts in \(\mathcal A\).

We now choose the category \(\mathcal A=\mathcal N\) of
F-SM-34, whose abelian structure, all limits and colimits, and
enough projectives were proved there.  Recall the exact
countable non-AB4 witness.  With \(E(V)=V^{**}/V\), put
\[
 P_0(V)=(V,V,1,0),\qquad
 P_1(V)=(EV,V,0,1),\qquad V=k^{(\mathbb N)}.
\]
Since \(Ek=0\), each map
\(P_1(k)\to P_0(k)\) with components \((0,1)\) is monic.
Their countable coproduct is \(P_1(V)\to P_0(V)\), because
both \(P_i\) are free functors and preserve coproducts.
This map has kernel \((EV,0,0,0)\ne0\).  Nonvanishing follows
by pulling a nonzero linear functional on
\(k^{\mathbb N}/k^{(\mathbb N)}\) back to \(k^{\mathbb N}\):
it gives a bidual class which no evaluation vector represents.
Embedding these maps as \((0,P_1(k))\to(0,P_0(k))\) in
\(\mathcal M\) preserves this coproduct and its nonzero kernel.
Thus countable AB4 fails in \(\mathcal M\).

Finally, by (SM43.7), the weak equivalences are exactly the
maps whose first components are quasi-isomorphisms.
Projection \(\Lambda\) has the section \(G(C)=(C,0)\), with
\(\Lambda G=1\), and the natural map
\((C,A)\to(C,0)\) is a weak equivalence.  After localization
these maps yield mutually inverse functors, proving
\[
                     W^{-1}\mathcal M\simeq D_{\ge0}(k).
                                                               \tag{SM43.10}
\]
All coproducts in \(\mathcal M\) have the usual chain
coproduct as first component.  Hence (SM43.7) also identifies
each canonical map
\(\bigoplus_iL\varphi(X_i)\to
L\varphi(\bigoplus_iX_i)\) with the identity of that chain
coproduct.  Thus the derived coproduct condition is strictly
weaker than AB4 in this full interval setting.  This example
does not decide whether the original assumptions always force
that weaker condition or all-free scalar normalization.
\(\square\)

\paragraph{Source and dependence scope.}
The original predicates are from Grothendieck, A. (n.d.),
\emph{Pursuing stacks: Linearization of homotopy types},
cote~134-8, footers~46--48,
\url{https://grothendieck.umontpellier.fr/134-8.pdf}.
The abelian non-AB4 category is exactly the previously
fully proved F-SM-34; the derived identification uses the
already proved finite-diagonal lemma BR-13401.  The tensor,
all its coherence, closedness, the interval, and the
localization calculation are verified above.

% END source_cote1348_no_ab4.tex

% BEGIN source_cote1348_low_degree.tex
\section{F-SM-44: automatic low-degree compatibility and the precise higher obstruction}

\subsection{Statement}
Let \(k\) be a commutative ring and \(\mathcal M\) a complete
and cocomplete \(k\)-linear abelian category with enough
projectives.  Suppose there are \(k\)-linear adjunctions
\[
 \varphi:\mathcal M\rightleftarrows\operatorname{Mod}_k:\kappa,
       \qquad \varphi\dashv\kappa\dashv\psi,
       \qquad \kappa\text{ exact}.                           \tag{SM44.1}
\]
For every set-indexed family \((X_s)_{s\in S}\), the canonical
maps
\[
 b_n:\bigoplus_{s\in S}L_n\varphi(X_s)
           \longrightarrow L_n\varphi\!\left(\bigoplus_{s\in S}X_s\right)
                                                               \tag{SM44.2}
\]
are monomorphisms for all \(n\ge0\), and are isomorphisms for
\(n=0,1\).  No AB4, tensor exactness, closedness, or interval
assumption is needed for this assertion.

More precisely, fix an injective cogenerator \(I\) of
\(\operatorname{Mod}_k\), put \(Y=\kappa I\), and let
\(C_n=\operatorname{coker}b_n\).  Restriction to the summands
gives a natural exact sequence
\[
\begin{split}
 0\longrightarrow\operatorname{Hom}_k(C_n,I)
 &\longrightarrow
 \operatorname{Ext}^n_{\mathcal M}\!\left(\bigoplus_sX_s,Y\right)\\
 &\longrightarrow\prod_s\operatorname{Ext}^n_{\mathcal M}(X_s,Y)
       \longrightarrow0 .
\end{split}                                                    \tag{SM44.3}
\]
For \(n=0,1\), the left term vanishes.  Thus the remaining
derived-coproduct condition is precisely the vanishing, in
degrees \(n\ge2\), of the classes in the middle group whose
restriction to every summand is zero.

If in addition the canonical unit counit
\[
                      L\varphi(\kappa k)\longrightarrow k[0]
                                                               \tag{SM44.4}
\]
is an equivalence, then for every scalar module \(U\) the
canonical maps satisfy
\[
         \varphi\kappa U\simeq U,
                    \qquad L_1\varphi(\kappa U)=0.            \tag{SM44.5}
\]
In particular, any failure of full scalar normalization under
the original assumptions begins in degree at least two.
The exact scalar embedding also induces
\[
 \operatorname{Ext}^1_k(U,V)
       \simeq\operatorname{Ext}^1_{\mathcal M}(\kappa U,\kappa V)
                                                               \tag{SM44.6}
\]
for all \(U,V\).

\subsection{Complete proof}
\paragraph{1. Products are exact, and scalar duality computes the relevant Ext.}
Products in \(\mathcal M\) preserve kernels.  To see that they
preserve epimorphisms, take a family of epimorphisms
\(A_s\to B_s\), and choose a projective epimorphism
\(P\to\prod_sB_s\).  Its component maps lift to \(A_s\)
because \(P\) is projective.  Their product gives a map
\(P\to\prod_sA_s\) whose composite to \(\prod_sB_s\)
is epic.  Hence \(\prod_sA_s\to\prod_sB_s\) is epic.
This proves exactness of products without an AB4 assumption.

The functor \(\varphi\) takes projective objects to projective
scalar modules: \(\operatorname{Hom}_k(\varphi P,-)
=\operatorname{Hom}_{\mathcal M}(P,\kappa(-))\) is exact.
For a projective resolution \(P_\bullet\to X\), the
adjunction and injectivity of \(I\) therefore give, naturally,
\[
 \operatorname{Ext}^n_{\mathcal M}(X,\kappa I)
   =H^n\operatorname{Hom}_k(\varphi P_\bullet,I)
   =\operatorname{Hom}_k(L_n\varphi X,I).                     \tag{SM44.7}
\]
An injective cogenerator always exists in the scalar category;
one choice is \(I=\operatorname{Hom}_{\mathbf Z}(k,\mathbf Q/\mathbf Z)\).
The adjunction with the underlying abelian group identifies its
Hom functor with \(\operatorname{Hom}_{\mathbf Z}(-,\mathbf Q/\mathbf Z)\),
which is exact and detects nonzero abelian groups.

\paragraph{2. Every family of extension classes is the restriction of one class.}
We prove surjectivity of
\[
 R_n:\operatorname{Ext}^n_{\mathcal M}\!\left(\bigoplus_sX_s,Y\right)
              \longrightarrow\prod_s\operatorname{Ext}^n_{\mathcal M}(X_s,Y)
                                                               \tag{SM44.8}
\]
for \(n\ge1\).  The empty family is immediate, so assume
\(S\ne\varnothing\).  Represent prescribed classes \(e_s\)
by exact Yoneda \(n\)-extensions
\[
 0\longrightarrow Y\longrightarrow E_{s,n-1}
      \longrightarrow\cdots\longrightarrow E_{s,0}
      \longrightarrow X_s\longrightarrow0.
                                                               \tag{SM44.9}
\]
Their termwise product is exact by Step~1 and represents a class
\[
             e_{\Pi}\in
 \operatorname{Ext}^n_{\mathcal M}\!\left(\prod_sX_s,\prod_sY\right).
\]
Since \(\kappa\) is a right adjoint, \(\prod_sY=\kappa(I^S)\).
In the scalar module category the canonical inclusion
\(I^{(S)}\hookrightarrow I^S\) is monic.  The finite-support
codiagonal \(I^{(S)}\to I\) extends by injectivity of \(I\)
to a map
\[
                   t:I^S\longrightarrow I,
                   \qquad t\iota_s=1_I\quad(s\in S),         \tag{SM44.10}
\]
where \(\iota_s\) is the coordinate inclusion.
Let \(j:\bigoplus_sX_s\to\prod_sX_s\) be the canonical
map.  It need not be monic, and no such assertion is used.
Pull back the right end of \(e_{\Pi}\) along \(j\), and
push out its left end along \(\kappa t\).  These endpoint
operations on Yoneda extensions give
\[
                   e=(\kappa t)_*\,j^*e_{\Pi}
              \in\operatorname{Ext}^n_{\mathcal M}
                         \!\left(\bigoplus_sX_s,Y\right).     \tag{SM44.11}
\]

We check each restriction, including its coefficient map.
Denote the coproduct inclusion by \(\mu_s\), the coordinate
inclusion into \(\prod_sX_s\) by \(v_s\), and the coordinate
inclusion \(Y\to\prod_sY\) by \(w_s=\kappa\iota_s\).
There is a morphism from the \(s\)-th extension (SM44.9) to
its product extension, using coordinate inclusions in every
term, with endpoint maps \(w_s,v_s\).  Functoriality of
Yoneda pushout and pullback gives
\[
                 v_s^*e_{\Pi}=(w_s)_*e_s.
\]
Because \(j\mu_s=v_s\) and \((\kappa t)w_s=1_Y\),
\[
 \mu_s^* e
   =(\kappa t)_*v_s^*e_{\Pi}
   =((\kappa t)w_s)_*e_s=e_s.                               \tag{SM44.12}
\]
This proves the claimed surjectivity in every degree.  The
choice of \(t\) need not be natural and does not affect the
surjectivity assertion; \(R_n\) itself is canonical.

\paragraph{3. In degree one, restriction is also injective.}
This step holds in every abelian category with the relevant
coproduct, for every coefficient object \(Y\).  An extension
\[
                0\longrightarrow Y\longrightarrow E
                   \xrightarrow{q}\bigoplus_sX_s\longrightarrow0
\]
whose pullback to each \(X_s\) splits admits maps
\(u_s:X_s\to E\) with \(qu_s=\mu_s\).  By the coproduct
universal property, these assemble to \(u:\bigoplus_sX_s\to E\)
with \(qu=1\).  The extension splits, proving injectivity
of \(R_1\).  In degree zero the ordinary Hom coproduct
identity is already an isomorphism.

\paragraph{4. Dualize the restriction maps.}
Under (SM44.7) and
\(\operatorname{Hom}_k(\bigoplus_sL_n\varphi X_s,I)
=\prod_s\operatorname{Hom}_k(L_n\varphi X_s,I)\),
the map \(R_n\) is precisely \(\operatorname{Hom}_k(b_n,I)\).
The exact faithful contravariant functor
\(\operatorname{Hom}_k(-,I)\) detects monomorphisms by
surjectivity, and detects isomorphisms.  Step~2 therefore
makes every \(b_n\) monic, and Step~3 makes \(b_1\) an
isomorphism.  The case \(n=0\) also follows directly from
preservation of coproducts by the left adjoint \(\varphi\).
Dualizing
\(0\to\bigoplus_sL_n\varphi X_s\xrightarrow{b_n}
L_n\varphi(\bigoplus_sX_s)\to C_n\to0\)
gives exactly (SM44.3).

Equivalently, for the canonical derived comparison
\[
 b:\bigoplus_s L\varphi(X_s)
                \longrightarrow L\varphi\!\left(\bigoplus_sX_s\right),
                                                               \tag{SM44.13}
\]
the homology of its cone in degree \(n\) is \(C_n\): the
long exact sequence of the cone and the injectivity of
\(b_n,b_{n-1}\) give this assertion.  Thus this cone has
zero homology in degrees zero and one, and full derived
coproduct compatibility is equivalent to \(C_n=0\) for
all \(n\ge2\).  Ordinary scalar-module coproducts are exact,
so their homology commutes with the coproduct in (SM44.13);
no derived coproduct assertion in \(\mathcal M\) is used here.

\paragraph{5. Canonical scalar normalization through degree one.}
Assume (SM44.4).  Its degree-zero map says that the ordinary
counit \(\varphi\kappa k\to k\) is an isomorphism.
Both \(\varphi\) and \(\kappa\) preserve coproducts and
cokernels, since both are left adjoints in (SM44.1).
The same canonical counit is therefore an isomorphism on
free modules; applying it to a free presentation and taking
cokernels proves
\(\varphi\kappa U\xrightarrow{\sim}U\) for every \(U\).
In particular \(\kappa\), as the right adjoint of
\(\varphi\), is fully faithful.

For \(F=k^{(S)}\), one has
\(\kappa F=\bigoplus_s\kappa k\).  Applying the isomorphism
\(b_1\) and (SM44.4) gives
\(L_1\varphi(\kappa F)=0\).  For any scalar \(U\), choose
an exact sequence \(0\to K\to F\to U\to0\) with \(F\)
free.  Exactness of \(\kappa\) and the derived long exact
sequence give
\[
 0=L_1\varphi(\kappa F)\longrightarrow L_1\varphi(\kappa U)
       \longrightarrow\varphi\kappa K\longrightarrow\varphi\kappa F.
\]
Under the canonical counit identifications, the last map is
the monomorphism \(K\to F\).  Thus
\(L_1\varphi(\kappa U)=0\), proving (SM44.5).

Finally, given an extension
\(0\to\kappa V\to Z\to\kappa U\to0\), the vanishing
just proved makes its image under \(\varphi\) the exact
sequence \(0\to V\to\varphi Z\to U\to0\).
The unit \(Z\to\kappa\varphi Z\) gives a morphism from
the original extension to the image of this scalar extension.
Its endpoint maps are isomorphisms by the adjunction identities,
so its middle map is an isomorphism.  Applying \(\varphi\)
and \(\kappa\) to extensions therefore gives inverse maps
on their classes, including the Baer operations, which proves
(SM44.6).  \(\square\)

\subsection{Sharp scope and the remaining original step}
The tensor and interval hypotheses have not been used to prove
this result.  Thus it applies in particular to the original
abstract setup.  It reduces the unresolved coproduct question
to \emph{surjectivity} of (SM44.2) in degrees \(n\ge2\),
or equivalently to injectivity of the corresponding higher
restriction maps (SM44.8).  It does not prove that the tensor
or interval kills those higher kernels.

The bound cannot be improved by the additive assumptions alone:
in F-SM-34, for the family of countably many copies of
\(\kappa k\), the degree-two map is
\(0\to E(k^{(\mathbb N)})\ne0\).  This is a precise additive
sharpness example, not a full symmetric-tensor counterexample.

\subsection{Primary reference and check scope}
Argud\'in-Monroy, A. (2022). The Yoneda Ext and arbitrary
coproducts in abelian categories. \emph{Glasgow Mathematical
Journal, 64}(2), 277--291.
\url{https://doi.org/10.1017/S0017089521000070}.
Version inspected: arXiv:1904.12182v2, September~8, 2020,
\url{https://arxiv.org/pdf/1904.12182}.
Proposition~3.11(b), p.~12, proves the degree-one injection by
assembling sections; the discussion after Theorem~3.13, p.~14,
explicitly states its validity for AB3 and separates the
higher-degree problem.  Only these precise source passages
and the publisher's bibliographic metadata were inspected for
this lemma.  The product-extension construction and its
application to the two scalar adjoints are proved in full above;
no general AB4 theorem is being applied to \(\mathcal M\).

% END source_cote1348_low_degree.tex

% BEGIN source_cote1348_ab4.tex
\section{F-SM-32: the interval theorem with exact coproducts}

\subsection{Statement and remaining historical scope}
Let \(k\) be a commutative ring.  Let \(\mathcal M\) be a
\(k\)-linear abelian category with enough projectives and all small
coproducts.  Suppose that
\(\varphi:\mathcal M\rightleftarrows\operatorname{Mod}k:\kappa\)
is an adjunction, with \(\kappa\) exact and coproduct preserving.
Write \(\mathcal D_{\mathcal M}=D_{\ge0}(\mathcal M)_\infty\),
\(\mathcal D_k=D_{\ge0}(k)_\infty\), and
\[
 F(X)=\mathbb L\varphi(X[0]),\qquad
 W=\{f:F(f)\text{ is an equivalence}\}.
\]
Assume:
\begin{enumerate}
\item the \emph{canonical} derived counit
\(F(\kappa k)\to k[0]\) is an equivalence;
\item each \(X\in\mathcal M\) admits a monomorphism
\(a_X:X\hookrightarrow A_X\) with \(F(A_X)=0\), including
degree-zero homology;
\item for every small family \((X_i)\), the canonical map
\[
       \bigoplus_i F(X_i)\longrightarrow
                         F\!\left(\bigoplus_iX_i\right)
                         \quad\text{is an equivalence}.       \tag{SM32.1}
\]
\end{enumerate}
Then \(F\) induces an equivalence
\[
        \mathcal M[W^{-1}]_\infty\simeq\mathcal D_k.            \tag{SM32.2}
\]
Consequently its homotopy category is the ordinary localization
\(W^{-1}\mathcal M\simeq D_{\ge0}(k)\).

In particular, condition (SM32.1) holds if \(\mathcal M\)
satisfies \emph{AB4}: its small coproduct functors are exact.
Thus the full interval theorem holds for every such AB4 category
with enough projectives, without AB5, a generator, local
presentability, an exact chain-valued realization, or a
functorial choice of projective resolutions.

The original interval supplies the second assumption by
\(A_X=J\otimes X\) and \(a_X=\alpha\otimes1_X\).
Fresh visual reading of cote~134-8, footers~46--48, confirms
the scalar adjoints, exactness of \(\kappa\), and the bilinear
tensor convention.  These pages do not explicitly require
exactness of arbitrary copowers on \emph{every} object of
\(\mathcal M\), or state AB4.  Bilinearity and the scalar
adjoints alone are not treated as such a requirement here.
The result therefore removes the AB5/presentability gate and
reduces the remaining abstract case to the precise coproduct
condition (SM32.1); it does not silently add AB4 to the original
handwritten conditions (a)--(d).

Fix universes \(\mathbb U\in\mathbb V\), with the categories
\(\mathbb V\)-small and the specified coproducts
\(\mathbb U\)-small.  Localizations are initially formed in
\(\mathbb V\).  The proof does not assume that their mapping
spaces are \(\mathbb U\)-small; the final equivalence supplies
that conclusion when \(\operatorname{Mod}k\) is formed in
\(\mathbb U\).  Geometric realizations have the countable
indexing category \(\Delta^{\mathrm{op}}\).

\subsection{Known inputs}
We use the following primary results, with their hypotheses
kept distinct from the argument below.
\begin{itemize}
\item Cisinski, \emph{Higher categories and homotopical algebra},
author version dated July~30, 2026, Proposition~7.5.6,
printed pp.~372--373, and its dual: the localization of a
category of cofibrant objects has finite colimits, and a
functor preserving its initial object and cofibration pushouts
induces a finite-colimit-preserving functor.  Proposition~7.7.1,
p.~396, handles products of such relative categories.
\item Ramzi, Sosnilo, and Winges (2025), Definition~2.5 and
Theorem~2.6(1),(4), pp.~5--6: in the same finite cofibration
setting, coproducts stable under weak equivalences pass to the
Dwyer--Kan localization.  The proof of Theorem~2.6, together
with Remark~2.3 and Example~2.4, gives the left-fraction
description used below.  Neither a sequential-colimit axiom
nor accessibility is a hypothesis here.
\item Lurie, \emph{Higher algebra} (September~18, 2017),
Theorem~1.3.3.8, p.~101: for an abelian category with enough
projectives, restriction to its projectives identifies functors
from its connective derived \(\infty\)-category preserving
geometric realizations with arbitrary functors on those
projectives.  This is an equivalence of \emph{functor
\(\infty\)-categories}, including natural transformations.
Lemma~1.3.3.11(2), pp.~103--104, and Proposition~1.3.3.16,
pp.~106--107, show that connective derived functors coming
from right \(t\)-exact exact functors preserve these
realizations.  Lemma~1.3.3.17 and its proof, p.~107, identify
the realization of the Dold--Kan simplicial object of a
connective projective complex with that complex, including
the augmentation comparison.  Theorem~1.3.4.4, p.~109,
identifies the ordinary chain-category localization.
\item Lurie, \emph{Higher topos theory},
Propositions~4.4.2.6--7, pp.~296--297: small coproducts and
pushouts generate small colimits, including the corresponding
preservation statement.
\end{itemize}
These are existing background theorems.  The complete deduction
of (SM32.2), especially its two coherent inverse comparisons,
is given next.

\subsection{Complete proof}

\paragraph{1. Exact coproducts imply (SM32.1).}
Suppose AB4.  For each \(i\), choose a projective resolution
\(P_{i,\bullet}\to X_i\).  A coproduct of projectives is
projective: maps from each summand lift separately through an
epimorphism.  AB4 makes \(\bigoplus_iP_{i,\bullet}\) a
projective resolution of \(\bigoplus_iX_i\).
As a left adjoint, \(\varphi\) preserves coproducts, so applying
it to this resolution gives
\[
 \varphi\!\left(\bigoplus_iP_{i,\bullet}\right)
                  =\bigoplus_i\varphi(P_{i,\bullet}).
\]
Coproducts in \(\mathcal D_k\) are computed by these ordinary
coproducts of complexes.  The comparison maps induced by the
summand inclusions are precisely those in (SM32.1).  This proves
the claimed sufficient condition, without any filtered-colimit
argument.  The remainder of the proof only uses (SM32.1).

\paragraph{2. A finite cofibration structure and its localization.}
Take every monomorphism of \(\mathcal M\) to be a cofibration.
Initial maps are cofibrations, and their pushouts exist and are
monomorphisms.  The coherent derived functor
\(\mathbb L\varphi\) sends a short exact sequence to its
derived cofiber sequence.  In particular, for a monomorphism
\(u\),
\[
       u\in W\quad\Longleftrightarrow\quad
                       F(\operatorname{coker}u)=0.             \tag{SM32.3}
\]
A pushout retains this cokernel, so trivial cofibrations are
stable under arbitrary pushout.  More generally, pushing a
weak equivalence along a monomorphism compares two short
exact sequences with the same quotient; their derived
triangles show that the resulting map is again in \(W\).

Every map \(f:X\to Y\) has the factorization
\[
 X\xrightarrow{(a_X,f)}A_X\oplus Y
                          \xrightarrow{\operatorname{pr}_Y}Y. \tag{SM32.4}
\]
The first map is monic, and the second is in \(W\), because
\(F\) is additive and \(F(A_X)=0\).  Choices in this
factorization need not be functorial.  Two-out-of-three and
two-out-of-six follow from the definition of \(W\).
These verifications give all the cofibrant-object axioms in
Ramzi--Sosnilo--Winges Definition~2.5.

Write \(\gamma:\mathcal M\to\mathcal C=
\mathcal M[W^{-1}]_\infty\).  The cited localization theorem
gives finite colimits in \(\mathcal C\).  By (SM32.1), small
coproducts of maps in \(W\) stay in \(W\); Theorem~2.6(4)
therefore gives small coproducts in \(\mathcal C\), preserved
by \(\gamma\).  HTT Proposition~4.4.2.6 gives all small
colimits in \(\mathcal C\).

The actual coherent functor
\(F=\mathbb L\varphi\circ(X\mapsto X[0])\) factors as
\[
             \overline F:\mathcal C\longrightarrow\mathcal D_k.
                                                               \tag{SM32.5}
\]
It preserves finite colimits.  Indeed a pushout along a
monomorphism \(A\to B\), with other map \(A\to D\), has
the exact sequence
\[
 0\longrightarrow A\longrightarrow B\oplus D
                   \longrightarrow B\amalg_A D\longrightarrow0,
\]
so \(F\) sends this pushout to a derived pushout.  Apply the
last assertion of the dual of Cisinski Proposition~7.5.6.
Equation (SM32.1) and preservation of coproducts by
\(\gamma\) show that \(\overline F\) preserves coproducts.
HTT Proposition~4.4.2.7 now shows that it preserves all small
colimits, in particular geometric realizations.

We will also need \emph{conservativity}.  The fraction
description in the cited cofibration theorem represents every
morphism of \(h\mathcal C\) by
\[
                  X\xrightarrow{f}Z\xleftarrow{s}Y,
                       \qquad s\in W.                        \tag{SM32.6}
\]
If its image under \(\overline F\) is invertible, then
\(F(f)\) is invertible because \(F(s)\) is.  Thus
\(f\in W\), and (SM32.6) is already invertible in
\(h\mathcal C\).  An \(\infty\)-morphism is an equivalence
exactly when it is invertible in its homotopy category.
Consequently \(\overline F\) is conservative.  This argument
does not infer conservativity merely from the definition of
\(W\); it uses the fraction theorem.

\paragraph{3. Scalar normalization and projective units.}
Exactness of \(\kappa\) gives, for projective \(P\in\mathcal M\),
\[
 \operatorname{Hom}_k(\varphi P,-)
                  =\operatorname{Hom}_{\mathcal M}(P,\kappa(-))
\]
exact.  Thus \(\varphi P\) is projective over \(k\).
There is the canonical derived adjunction
\(\mathbb L\varphi\dashv\kappa:\mathcal D_k\to
\mathcal D_{\mathcal M}\), with \(\kappa\) applied termwise;
its counit on a complex is represented by
\(\varphi P(\kappa C)\to\varphi\kappa C\to C\).
This notation denotes a projective representative, not a
functorial choice of one.

For a free scalar module \(k^{(I)}\), (SM32.1), preservation
of coproducts by \(\kappa\), and the normalization on \(k\)
show that the canonical counit
\[
          F(\kappa k^{(I)})\longrightarrow k^{(I)}[0]
\]
is an equivalence.  Naturality passes this assertion to every
projective scalar retract \(Q\).  These identifications are
the actual counit identifications.  Both derived functors
\(\mathbb L\varphi\kappa\) and the identity on
\(\mathcal D_k\) preserve geometric realizations, by HA
Lemma~1.3.3.11(2).  Theorem~1.3.3.8 therefore promotes the
projective identifications to the equivalence of the canonical
counit on all of \(\mathcal D_k\).

In particular, for projective \(P\in\mathcal M\), the original
adjunction unit satisfies
\[
                 \eta_P:P\longrightarrow\kappa\varphi P
                            \quad\in W.                       \tag{SM32.7}
\]
Its image under \(F\), followed by the just-established
counit identification, is the identity of \(\varphi P[0]\)
by the adjunction triangle identity.

\paragraph{4. Coherent extension from projectives.}
By HA Theorem~1.3.3.8, the functor
\(Q\mapsto\gamma\kappa Q\) on projective scalar modules
extends to a geometric-realization-preserving functor
\[
                         L:\mathcal D_k\longrightarrow\mathcal C.
                                                               \tag{SM32.8}
\]
The same theorem, applied to \(P\mapsto\gamma P\), gives
\[
                    T:\mathcal D_{\mathcal M}\longrightarrow\mathcal C.
                                                               \tag{SM32.9}
\]
All comparisons here are comparisons of \(\infty\)-functors.
Since \(\overline F\) preserves realizations, the functors
\(\overline F L\) and \(\mathrm{id}_{\mathcal D_k}\)
preserve them and have the canonically equivalent restrictions
to projective scalar modules supplied in Step~3.
Full faithfulness in HA Theorem~1.3.3.8 gives
\[
                      \overline F L\simeq\mathrm{id}_{\mathcal D_k}.
                                                               \tag{SM32.10}
\]
Likewise, \(T\) and \(L\mathbb L\varphi\) preserve
realizations.  Their restrictions to projective \(P\) are
\(\gamma P\) and \(\gamma\kappa\varphi P\), and (SM32.7)
gives a natural equivalence between them.  Thus
\[
                             T\simeq L\mathbb L\varphi.       \tag{SM32.11}
\]

For clarity, the following step constructs the required natural
augmentation; an objectwise choice of resolutions would not
by itself supply this comparison.

\paragraph{5. The comma-category augmentation and its identification.}
For a functor \(g:\mathcal A\to\mathcal E\), where both
categories admit geometric realizations, let
\((\mathrm{id}_{\mathcal E}\downarrow g)\) have objects
\((e,a,u:e\to g(a))\).  This comma \(\infty\)-category has
geometric realizations computed in its two components, even
if \(g\) does not preserve them.  Indeed a simplicial diagram
has colimit arrow
\[
 \operatorname*{colim}e_n\longrightarrow
 \operatorname*{colim}g(a_n)\longrightarrow
                                g(\operatorname*{colim}a_n).
                                                               \tag{SM32.12}
\]
To check its universal property, use the mapping-space formula
\[
 \operatorname{Map}((e,a,u),(e',a',u'))
 =\operatorname{Map}_{\mathcal E}(e,e')
  \mathop{\times}_{\operatorname{Map}_{\mathcal E}(e,g(a'))}
                 \operatorname{Map}_{\mathcal A}(a,a').        \tag{SM32.13}
\]
Mapping out of each componentwise colimit is a limit, and
limits of spaces commute with this pullback.  Hence (SM32.12)
satisfies the asserted universal property.  It also follows
that a functor into this comma category preserves realizations
whenever both its component functors do.

Apply this with \(g=\gamma:\mathcal M\to\mathcal C\),
viewing \(\mathcal M\) as its nerve.  An ordinary abelian
category has geometric realizations, computed by its reflexive
coequalizers.  Extend
\[
       P\longmapsto(\gamma P,P,\mathrm{id}_{\gamma P})
                                                               \tag{SM32.14}
\]
from the projectives by HA Theorem~1.3.3.8, now with this
comma category as target.  Its first component is \(T\).
Its second component is \(H_0:\mathcal D_{\mathcal M}
\to\mathcal M\): this functor is left adjoint to the heart
inclusion \(\iota:X\mapsto X[0]\), preserves realizations,
and agrees with the second component on projectives.
We obtain a coherent natural transformation
\[
                             b:T\longrightarrow\gamma H_0.  \tag{SM32.15}
\]

We identify its image under \(\overline F\), rather than
assuming that it is the desired augmentation.  The functor
\[
 (\mathrm{id}_{\mathcal C}\downarrow\gamma)
       \longrightarrow(\mathrm{id}_{\mathcal D_k}\downarrow F)
\]
induced by \(\overline F\) preserves realizations by
(SM32.12).  On the other hand, the functor
\[
 C\longmapsto
 \bigl(\mathbb L\varphi(C),H_0(C),
       \mathbb L\varphi(C\to\iota H_0(C))\bigr)               \tag{SM32.16}
\]
also preserves realizations: its components
\(\mathbb L\varphi\) and \(H_0\) do, and (SM32.13)
checks the comma-category claim.  Both resulting functors from
\(\mathcal D_{\mathcal M}\) to this last comma category
restrict on projective \(P\) to
\((\varphi P[0],P,\mathrm{id})\).  The full faithfulness
part of HA Theorem~1.3.3.8 identifies them as arrow-valued
functors.  In particular, when \(C=X[0]\),
\(\overline F(b_{X[0]})\) identifies with the image of the
identity truncation \(X[0]\to H_0(X[0])[0]\), and is an
equivalence.  Conservativity from Step~2 proves the natural
equivalence
\[
                             T\iota\xrightarrow{\sim}\gamma. \tag{SM32.17}
\]

Equations (SM32.11) and (SM32.17) now give a natural equivalence
\(L\overline F\gamma\simeq\gamma\).  The universal property
of relative \(\infty\)-localization is fully faithful on
functor \(\infty\)-categories; it therefore descends this
comparison to
\[
                         L\overline F\simeq\mathrm{id}_{\mathcal C}.
                                                               \tag{SM32.18}
\]
Together with (SM32.10), this proves (SM32.2).
Passing to homotopy categories gives the asserted ordinary
localization equivalence. \(\square\)

\subsection{Concrete realization and exact limitations}
If \(Q_\bullet\) is a connective complex of projective scalar
modules representing \(C\in\mathcal D_k\), the preceding
inverse is concretely
\[
               L(C)\simeq
                  \operatorname*{colim}_{\Delta^{\mathrm{op}}}
                                      \gamma\kappa\Gamma(Q)_n.
                                                               \tag{SM32.19}
\]
For a projective resolution \(P_\bullet\to X\),
\(T(X[0])\simeq\operatorname*{colim}\gamma\Gamma(P)_n\),
and (SM32.15) is its augmentation to \(\gamma X\).
These identifications follow from HA Lemma~1.3.3.17; the
comma-category argument proves that their choices fit together
coherently.  No ordinary filtered colimit in \(\mathcal M\)
is substituted for either realization.

The new scope strictly includes the AB4, non-AB5,
nonpresentable example F-SM-31.  What remains in the bare
abstract archive claim is whether its actual remaining
assumptions force (SM32.1), or whether a full tensor-and-interval
counterexample violates it.  A merely additive three-sort
construction is insufficient for that negative conclusion.
The theorem also makes no claim that an ordinary right adjoint
\(G\) to a separately chosen exact realization is already a
quasi-inverse.  Its raw counit test remains distinct.

\subsection{Primary references}
\begin{itemize}
\item Cisinski, D.-C. (2019). \emph{Higher categories and
homotopical algebra}. Cambridge University Press.
Author version dated July~30, 2026:
\url{https://cisinski.app.uni-regensburg.de/CatLR.pdf}.
Exact locators used here refer to that version:
Proposition~7.5.6, pp.~372--373; Proposition~7.7.1, p.~396;
Definition~7.7.2 and Proposition~7.7.4, p.~397.
\item Grothendieck, A. (n.d.). \emph{Cote 134-8}
[Unpublished manuscript], footers~46--48.
Archives Grothendieck, Universit\'e de Montpellier.
\url{https://grothendieck.umontpellier.fr/134-8.pdf}.
\item Lurie, J. (2009). \emph{Higher topos theory}
(Annals of Mathematics Studies, Vol.~170).
Princeton University Press.
\url{https://www.math.ias.edu/~lurie/papers/HTT.pdf}.
\item Lurie, J. (2017, September 18). \emph{Higher algebra}
[Author manuscript].
\url{https://people.math.harvard.edu/~lurie/papers/HA.pdf}.
Exact inputs: Theorem~1.3.3.8, p.~101;
Lemma~1.3.3.11(2), pp.~103--104;
Proposition~1.3.3.16 and Lemma~1.3.3.17, pp.~106--107;
Theorem~1.3.4.4, p.~109.
\item Ramzi, M., Sosnilo, V., \& Winges, C. (2025).
\emph{Every motive is the motive of a stable
\(\infty\)-category} (arXiv:2503.11338v1).
\url{https://arxiv.org/abs/2503.11338}.
Exact input: Definition~2.5 and Theorem~2.6(1),(4),
pp.~5--6, with Remark~2.3 and Example~2.4, p.~4.
\end{itemize}

% END source_cote1348_ab4.tex

% BEGIN source_cote1348_countable.tex
\section{F-SM-33: a countable-coproduct refinement}

\subsection{Statement}
Keep the abelian category, enough projectives, scalar adjunction,
exactness and coproduct preservation of \(\kappa\), and full
acyclic-monomorphism assumption from F-SM-32.  Replace its
coproduct and scalar-normalization assumptions by:
\begin{enumerate}
\item for every countable family \((X_n)\), the canonical map
\[
       \bigoplus_n F(X_n)\longrightarrow F(\bigoplus_nX_n)
                         \quad\text{is an equivalence};       \tag{SM33.1}
\]
\item for every small set \(I\), the canonical derived counit
\[
                  F(\kappa k^{(I)})\longrightarrow k^{(I)}[0]
                         \quad\text{is an equivalence}.       \tag{SM33.2}
\]
\end{enumerate}
Then \(F\) again induces
\[
                   \mathcal M[W^{-1}]_\infty\simeq D_{\ge0}(k)_\infty.
                                                               \tag{SM33.3}
\]
Exact countable coproducts in \(\mathcal M\) suffice for
(SM33.1), by the projective-resolution argument in
F-SM-32, Step~1.  No assertion that the scalar adjoints imply
(SM33.2) is included.

\subsection{Complete proof}
The finite cofibrant-object structure of F-SM-32, Step~2,
uses only short exact sequences and the factorization
\(X\hookrightarrow A_X\oplus Y\twoheadrightarrow Y\).
It therefore remains valid here.  By (SM33.1), \(W\) is
closed under countable coproducts.  Apply
Ramzi--Sosnilo--Winges Theorem~2.6(4) with the cardinal bound
\(\aleph_1\): the localization \(\mathcal C\) has countable
coproducts and its localization functor \(\gamma\) preserves
them.  It also has finite colimits.  The actual coherent
\(\overline F:\mathcal C\to\mathcal D_k\) preserves finite
colimits by the exact-triangle argument in F-SM-32, and
preserves countable coproducts by (SM33.1).

HTT Propositions~4.4.2.6--7, with the same cardinal bound,
give countable colimits and their preservation.  The nerve
of \(\Delta^{\mathrm{op}}\) is countable: the simplex
category has countably many objects and morphisms, and its
simplices are finite composable strings.  Hence
\(\mathcal C\) admits geometric realizations and
\(\overline F\) preserves them.  No uncountable colimit in
\(\mathcal C\) is needed below.

The left-fraction argument in F-SM-32, equation (SM32.6),
does not use any coproduct hypothesis and proves that
\(\overline F\) is conservative.  Naturality of the
canonical counit in (SM33.2) passes its equivalence from
free scalar modules to their projective retracts.
Exactness of \(\kappa\) again makes \(\varphi P\)
projective for projective \(P\in\mathcal M\).
Consequently the triangle identity makes
\(P\to\kappa\varphi P\) a member of \(W\).

Apply HA Theorem~1.3.3.8 to extend
\(Q\mapsto\gamma\kappa Q\) on projective scalars to a
realization-preserving functor \(L:\mathcal D_k\to\mathcal C\),
and \(P\mapsto\gamma P\) on projectives of \(\mathcal M\)
to \(T:\mathcal D_{\mathcal M}\to\mathcal C\).
All required targets admit geometric realizations by the
preceding paragraph.  The projective counits and units give,
by the full faithfulness of the same theorem,
\[
                 \overline F L\simeq\mathrm{id},\qquad
                         T\simeq L\mathbb L\varphi.           \tag{SM33.4}
\]
Here \(\mathbb L\varphi\) preserves realizations by
HA Lemma~1.3.3.11(2).

The completely proved comma-category lemma (SM32.12)--(SM32.13)
requires only geometric realizations in its two component
categories.  Apply it to
\((\mathrm{id}_{\mathcal C}\downarrow\gamma)\).
Extending the functor \(P\mapsto(\gamma P,P,\mathrm{id})\)
gives the natural transformation \(b:T\to\gamma H_0\).
After applying \(\overline F\), the arrow-valued uniqueness
in HA Theorem~1.3.3.8 identifies this with
\[
          \mathbb L\varphi(C)\longrightarrow
                          \mathbb L\varphi(H_0(C)[0]).
\]
This is precisely the comparison from F-SM-32, Step~5;
its proof uses only realization preservation and is
applicable with the countable colimits just established.
On \(C=X[0]\), the last arrow is the identity truncation.
Conservativity therefore gives \(T\iota\simeq\gamma\).
Combine this with (SM33.4) and descend the resulting natural
equivalence \(L\overline F\gamma\simeq\gamma\) through
the relative localization.  It yields
\(L\overline F\simeq\mathrm{id}_{\mathcal C}\), proving
(SM33.3). \(\square\)

\subsection{F-SM-35: when an exact tensor convention already forces AB4}
Let \(\mathcal A\) be an abelian category with small products
and coproducts.  Assume that for every small set \(I\), the
copower functor \(X\mapsto X^{(I)}\) preserves monomorphisms.
Then \(\mathcal A\) is AB4.

\paragraph{Proof.}
Let \((f_i:A_i\hookrightarrow B_i)_{i\in I}\) be any
family of monomorphisms.  The map
\(f=\prod_i f_i:\prod_iA_i\to\prod_iB_i\) is monic,
because products preserve kernels.  Therefore
\(f^{(I)}\) is monic by hypothesis.  Put
\(A=\prod_iA_i\), \(B=\prod_iB_i\).  The coordinate
inclusions \(e_i:A_i\to A\), with every other coordinate
zero, and projections \(p_i:A\to A_i\) define maps
\[
 \bigoplus_iA_i\xrightarrow{e}A^{(I)}
                   \xrightarrow{p}\bigoplus_iA_i,
 \qquad pe=1.
\]
Here \(e\) sends the \(i\)-th summand into the \(i\)-th
copy of \(A\), and \(p\) sends that copy through \(p_i\)
into the \(i\)-th summand.  The identical construction for
\(B\) commutes with \(\bigoplus_i f_i\) and \(f^{(I)}\).
Thus \(\bigoplus_i f_i\) is a retract of the monomorphism
\(f^{(I)}\) in the arrow category, so is monic.
Coproducts already preserve cokernels, being left adjoints
to diagonal functors.  They are therefore exact. \(\square\)

In particular, if \(\mathcal A\) has an exact tensor and
natural identifications
\[
              X\otimes\kappa(k^{(I)})\simeq X^{(I)},
                                                               \tag{SM35.1}
\]
the lemma applies.  A closed symmetric monoidal structure
supplies (SM35.1), because tensoring preserves coproducts
and \(\kappa(k^{(I)})=\mathbf1^{(I)}\).
This is a conditional interpretation result: the source
pages inspected say bilinear and exhibit the scalar
adjoints; they do not explicitly supply this full closed,
exact tensor convention or all products.  We consequently
do not use the lemma to declare bare (a)--(d) proved.

\subsection{Primary references}
The known inputs, with their exact primary locators and URLs,
are those listed in F-SM-32: Ramzi, Sosnilo, and Winges
(2025), Definition~2.5 and Theorem~2.6, pp.~5--6,
\url{https://arxiv.org/abs/2503.11338};
Cisinski (2019, author version July~30, 2026),
Proposition~7.5.6, pp.~372--373,
\url{https://cisinski.app.uni-regensburg.de/CatLR.pdf};
Lurie (2017), Theorem~1.3.3.8, p.~101, and
Lemma~1.3.3.11(2), pp.~103--104,
\url{https://people.math.harvard.edu/~lurie/papers/HA.pdf};
and Lurie (2009), Propositions~4.4.2.6--7, pp.~296--297,
\url{https://www.math.ias.edu/~lurie/papers/HTT.pdf}.
F-SM-35 is proved directly above.

% END source_cote1348_countable.tex

% BEGIN source_cote1348_additive_counterexample.tex
\section{F-SM-34: the scalar adjoints do not force free-scalar normalization}

\subsection{Exact conclusion}
There is a complete and cocomplete \(k\)-linear abelian
category \(\mathcal N\), with enough projectives, and an exact
functor \(\kappa:\operatorname{Vect}_k\to\mathcal N\) having
both adjoints, such that the canonical counit
\[
              L\varphi(\kappa k)\longrightarrow k[0]
\]
is an equivalence, whereas the same counit is not an equivalence
on the countably generated free space \(k^{(\mathbb N)}\).
Here \(\varphi\dashv\kappa\), and even the underived
counit \(\varphi\kappa V\to V\) is an isomorphism for every
\(V\).  Countable coproducts of monomorphisms fail to be
monomorphisms in an explicitly exhibited family.

Passing to \(\operatorname{Ch}_{\ge0}(\mathcal N)\) supplies
an exact scalar functor with both adjoints, the same canonical
normalization on \(k\), and a natural monomorphism of every
object into a fully derived-acyclic object.  Even its ordinary
localization conclusion fails.  This is a counterexample to
the \emph{additive, acyclic-effacement version} without a
coproduct condition.  No symmetric monoidal tensor on this
category satisfying the original archive assumptions is
constructed here.  Consequently this result does not refute
the full tensor-and-interval assertion in cote~134-8.

\subsection{Complete construction and proof}

\paragraph{1. The coefficient functor.}
Fix a field \(k\), and set
\[
                       E(V)=V^{**}/V,
\]
where \(V\to V^{**}\) is the canonical evaluation embedding.
The functor \(E\) is \(k\)-linear, additive and exact:
every short exact sequence of vector spaces splits, and an
additive functor preserves a split exact sequence.
It vanishes on finite-dimensional spaces, in particular on
\(k\).

For \(V=k^{(\mathbb N)}\), it is nonzero.  Identify
\(V^*=k^{\mathbb N}\).  The quotient
\(k^{\mathbb N}/k^{(\mathbb N)}\) is nonzero, as witnessed
by the all-ones sequence.  Choose a nonzero linear functional
on that quotient and pull it back to \(k^{\mathbb N}\).
It is a nonzero element of \(V^{**}\) vanishing on the
finite-support subspace.  No nonzero evaluation functional
from \(V\) vanishes on that subspace, since its coordinate
values recover the vector.  Its class in \(E(V)\) is therefore
nonzero.  This uses ordinary vector-space choice only.

\paragraph{2. The two-sort abelian category.}
An object of \(\mathcal N\) is
\[
 X=(X_0,X_1,b:X_0\to X_1,a:E(X_1)\to X_0),
                 \qquad ba=0,\quad aE(b)=0.                  \tag{SM34.1}
\]
A morphism is a pair \((u_0,u_1)\) satisfying
\(u_1b=b'u_0\) and \(u_0a=a'E(u_1)\).
Kernels and cokernels are computed on the two vector spaces:
exactness of \(E\) gives the induced structure maps and
preserves both relations.  The comparison from coimage to
image is an isomorphism componentwise.  Thus \(\mathcal N\)
is abelian, and the forgetful functor to
\(\operatorname{Vect}_k^2\) is exact and faithful.

For vector spaces \(U,V\), define
\[
 \mathsf P(U,V)=
       (U\oplus E(V),\ U\oplus V),
 \quad b(u,e)=(u,0),\quad a(e_U,e_V)=(0,e_V).                \tag{SM34.2}
\]
Here \(E(U\oplus V)=E(U)\oplus E(V)\).  The two relations
are immediate.  A map \(\mathsf P(U,V)\to X\) is uniquely
specified by maps \(u:U\to X_0\), \(v:V\to X_1\):
its components are
\((u,aE(v))\) and \((bu,v)\).  Hence
\[
 \operatorname{Hom}_{\mathcal N}(\mathsf P(U,V),X)
     =\operatorname{Hom}_k(U,X_0)\times\operatorname{Hom}_k(V,X_1).
                                                               \tag{SM34.3}
\]
This is a free-object adjunction.  Since the forgetful functor
is exact, each \(\mathsf P(U,V)\) is projective.  The counit
\(\mathsf P(X_0,X_1)\to X\) is componentwise epic: its
components contain the identity of \(X_0\) and of \(X_1\),
respectively.  There are therefore enough projectives.

For completeness, all small limits and colimits exist.
Products are formed componentwise; the map
\(E(\prod X_{i,1})\to\prod X_{i,0}\) has components
\(a_iE(\operatorname{pr}_i)\).  Both relations follow
after projection to each factor.  Together with the already
constructed kernels this gives all limits.

Here is an explicit construction of coproducts, without
assuming that \(E\) preserves them.  Write \(U\) for the
forgetful functor, \(\mathsf P\) for its left adjoint, and
\(T=U\mathsf P\).  For a family \(X_i\), put
\(L=\bigoplus_iUX_i\), \(R=\bigoplus_iTUX_i\) in
\(\operatorname{Vect}_k^2\).  Form the coequalizer in
\(\mathcal N\)
\[
             \mathsf P(R)\rightrightarrows\mathsf P(L)
                                      \longrightarrow Q.     \tag{SM34.4}
\]
The first arrow is induced by the sum of the algebra action
maps \(TUX_i\to UX_i\).  The second is the free extension
of the map \(R\to T L\) whose restriction to the \(i\)-th
summand is \(T\) of the summand inclusion \(UX_i\to L\).
These action maps are simply the free-adjunction counits.
One may also verify them by the square-zero monad
\(T=\mathrm{id}\oplus H\), \(H(U,V)=(E(V),U)\): its
multiplication adds the two \(H\)-summands and kills
\(H^2\).  Associativity adds the three linear \(H\)-terms
and kills all higher terms on either side; its algebra
relation is exactly (SM34.1).

Coequalizers exist as cokernels of differences.  A map from
\(\mathsf P(L)\) to an object \(Z\) corresponds to a family
of underlying linear maps \(UX_i\to UZ\).  Equalizing the
two arrows in (SM34.4) says precisely that each map commutes
with the two structure maps of \(X_i\) and \(Z\).
Thus \(\operatorname{Hom}(Q,Z)=\prod_i
\operatorname{Hom}(X_i,Z)\), proving the coproduct property.
Coproducts and cokernels give all small colimits.

\paragraph{3. Both scalar adjoints.}
Define \(\kappa V=(V,0,0,0)\).  Directly from the
morphism equations,
\[
 \begin{split}
 \operatorname{Hom}_{\mathcal N}(X,\kappa V)
       &=\operatorname{Hom}_k(\operatorname{coker}a,V),\\
 \operatorname{Hom}_{\mathcal N}(\kappa V,X)
       &=\operatorname{Hom}_k(V,\ker b).
 \end{split}                                                 \tag{SM34.5}
\]
Thus \(\varphi X=\operatorname{coker}a\) is left adjoint
to \(\kappa\), and \(\psi X=\ker b\) is its right
adjoint.  The functor \(\kappa\) is exact by componentwise
exactness, preserves all coproducts and products by these
adjoints, and is the vector-space copower of its distinguished
object \(\kappa k\).  The underived counit is literally
\(\varphi\kappa V=V\).

Put
\[
 P_0(V)=\mathsf P(V,0)=(V,V,\mathrm{id},0),\qquad
 P_1(V)=\mathsf P(0,V)=(E(V),V,0,\mathrm{id}).                \tag{SM34.6}
\]
The epimorphism \(P_0(V)\to\kappa V\) has kernel
\(S(V)=(0,V,0,0)\).  The epimorphism
\(P_1(V)\to S(V)\) has kernel \(\kappa E(V)\).
Repeating gives a projective resolution of \(\kappa V\)
with terms
\[
       P_{2n}=P_0(E^nV),\qquad P_{2n+1}=P_1(E^nV).
                                                               \tag{SM34.7}
\]
The odd differential has components \((0,\mathrm{id})\);
the positive even differential has components
\((\mathrm{id},0)\).  The preceding kernel calculations
prove exactness at every term, not merely that the
differentials compose to zero.

Applying \(\varphi\) gives \(E^nV\) in degree \(2n\)
and zero in odd degrees.  Therefore
\[
 L_{2n}\varphi(\kappa V)=E^nV,\qquad
 L_{2n+1}\varphi(\kappa V)=0\quad(n\ge0).                    \tag{SM34.8}
\]
For \(V=k\), this is \(k\) in degree zero and zero
elsewhere, and its canonical counit is the identity in
degree zero.  For \(V=k^{(\mathbb N)}\), degree two is
the nonzero space \(E(V)\).  This proves failure of the
free-scalar counit despite both scalar adjoints and the
canonical unit normalization.

There is also an explicit failure of exact countable sums.
Since \(E(k)=0\), the maps
\[
                S(k)=P_1(k)\hookrightarrow P_0(k)
\]
with components \((0,\mathrm{id})\) are monomorphisms.
Their countable coproduct is
\[
                    P_1(V)\longrightarrow P_0(V),
                        \qquad V=k^{(\mathbb N)},            \tag{SM34.9}
\]
because the free functor \(\mathsf P\) preserves
coproducts.  Its components are again
\((0,\mathrm{id})\), by the free-map formula (SM34.3).
Its kernel is \(\kappa E(V)\ne0\), so it is not monic.
Likewise \(\bigoplus_nL\varphi(\kappa k)\to
L\varphi(\bigoplus_n\kappa k)\) fails in degree two.

\paragraph{4. Actual acyclic effacement after adjoining chain degrees.}
Let \(\mathcal M=\operatorname{Ch}_{\ge0}(\mathcal N)\).
It is complete and cocomplete abelian.  It has enough
projectives: for projective \(Q\in\mathcal N\), the
complex \(Q[0]\) is projective in the chain category, as
is each disk \(D^n(Q)\), \(n\ge1\); their Hom functors
are evaluation in degrees \(0\) and \(n\), respectively,
followed by \(\operatorname{Hom}(Q,-)\).  Choose projective
epimorphisms onto each term of a complex and use the
corresponding disks and degree-zero complex.  Their
coproduct maps degreewise epimorphically to the complex.
There are at most two disk summands in any fixed degree
in this chosen epimorphism, so its construction does not
assume exact infinite coproducts.  A projective chain
object is a retract of such a coproduct.  The disk part
is explicitly contractible, and the remaining degree-zero
part is projective.

The exact scalar functor is
\(\kappa_{\mathcal M}V=(\kappa V)[0]\).  Its left and
right adjoints are, respectively,
\[
       \varphi_{\mathcal M}(C)=\varphi(H_0C),\qquad
                     \psi_{\mathcal M}(C)=\psi(C_0).           \tag{SM34.10}
\]
The adjunctions follow by factoring chain maps to a
degree-zero complex through \(H_0\), and by evaluating
chain maps from a degree-zero complex in degree zero.

There is a canonical identification
\[
       L\varphi_{\mathcal M}(C)=\mathbb L\varphi(C)
                                  \quad\text{in }D_{\ge0}(k), \tag{SM34.11}
\]
where the right-hand side regards \(C\) as a complex in
\(\mathcal N\).  To verify it, resolve \(C\), as an
object of the chain category, by projective chain objects
\(R_{r,\bullet}\).  Each has projective terms,
no positive chain homology, and projective \(H_0\), since
it is a retract of a coproduct of the disks and degree-zero
projectives just described.  The first-quadrant total
complex \(\operatorname{Tot}R\to C\) is a quasi-isomorphism
by exactness in the resolution direction.  The row
augmentations also give a quasi-isomorphism
\(\operatorname{Tot}R\to H_0(R_{\bullet,\bullet})\).
Both assertions use finite diagonals in each total degree.
The latter complex has projective \(\mathcal N\)-terms,
and applying \(\varphi\) gives exactly the resolution
formula for \(L\varphi_{\mathcal M}(C)\).  This proves
(SM34.11), with its canonical comparisons.  No infinite
sum of short exact sequences in \(\mathcal N\) is used.

Define \(A(C)_n=C_n\oplus C_{n-1}\), with \(C_{-1}=0\),
and
\[
        d(x,y)=(dx+y,-dy),\quad
        a_C(x)=(x,0),\quad h(x,y)=(0,x).                     \tag{SM34.12}
\]
These finite-matrix formulas make sense in any abelian
category.  The map \(a_C:C\to A(C)\) is a natural
degreewise monomorphism, and \(dh+hd=1\).
Thus \(A(C)\) is contractible and (SM34.11) implies
\(L\varphi_{\mathcal M}(A(C))=0\), including degree zero.
The canonical scalar counit on \(k\) still is an
equivalence by (SM34.8), and the one on
\(k^{(\mathbb N)}\) still is not.

\paragraph{5. Why the additive localization conclusion actually fails.}
Put
\[
 W_{\mathcal M}=(L\varphi_{\mathcal M})^{-1}
                         (\text{isomorphisms}).
\]
Suppose, contrary to the claim,
that the induced ordinary functor
\(W_{\mathcal M}^{-1}\mathcal M\to D_{\ge0}(k)\)
were an equivalence.  Since quasi-isomorphisms belong to
\(W_{\mathcal M}\), localization in stages and (SM34.11)
would make
\[
          \mathbb L\varphi:D_{\ge0}(\mathcal N)
                                         \longrightarrow D_{\ge0}(k)
\]
a localization functor.  It has right adjoint the exact
termwise functor \(\kappa\), since
\(\varphi\dashv\kappa\) and \(\varphi\) takes
projectives to projectives.

A localization functor having a right adjoint has fully
faithful right adjoint.  One can check this elementary fact
directly: for \(q:\mathcal A\to\mathcal B\),
\(q\dashv r\), each \(rY\) is local for the inverted
maps, by the adjunction.  A zigzag into \(rY\) can be
uniquely reduced to an actual map by successively extending
across each inverted map; hence
\(\operatorname{Hom}_{\mathcal A}(X,rY)=
\operatorname{Hom}_{\mathcal B}(qX,qrY)\).
The adjunction identifies composition with the counit
\(qrY\to Y\) with a bijection on these Hom sets.
Essential surjectivity of \(q\) and Yoneda make the
counit an isomorphism.  This is also the ordinary instance
of Cisinski Proposition~7.1.17.

In the present example the counit on
\(V=k^{(\mathbb N)}\) has nonzero degree-two source
homology \(E(V)\), by (SM34.8).  It cannot be an
isomorphism in the derived category.  This contradiction
proves failure of the asserted additive localization,
and therefore also of an \(\infty\)-categorical
equivalence induced by this functor. \(\square\)

\subsection{What this counterexample does not certify}
The complete construction certifies the two adjoints of the
scalar functor, its exactness, enough projectives, all limits and colimits,
canonical unit normalization, actual monic acyclic
effacement, and failure of the localization conclusion.
It does not supply the associative, symmetric, unital
tensor of the handwritten original, or identify (SM34.12)
with tensoring by an object of such a monoidal category.
That additional construction remains a substantive gate.
In particular, an exact closed tensor compatible with
the distinguished scalar unit would be impossible here:
F-SM-35 would imply AB4, contradicting the explicit
monomorphism calculation (SM34.9), which also persists in
degree zero in the chain category.

\subsection{Sources}
All example-specific assertions are proved above.  The
derived-category and localization background is the same
as in F-SM-32: Lurie (2017), \emph{Higher algebra},
Theorem~1.3.4.4, p.~109,
\url{https://people.math.harvard.edu/~lurie/papers/HA.pdf};
and Cisinski (2019, author version July~30, 2026),
\emph{Higher categories and homotopical algebra},
Proposition~7.1.17, p.~331,
\url{https://cisinski.app.uni-regensburg.de/CatLR.pdf}.
The scalar/tensor historical distinction refers to
Grothendieck's cote~134-8, footers~46--48,
\url{https://grothendieck.umontpellier.fr/134-8.pdf}.

% END source_cote1348_additive_counterexample.tex

% BEGIN source_cote1348_tensor_obstructions.tex
\section{Two obstructions to the direct tensor repairs of F-SM34}

\paragraph{Scope.}
These are obstructions to specified constructions, not a proof that
the additive category in F-SM34 admits no symmetric monoidal
structure.  In particular, they do not turn that additive
counterexample into a counterexample to all the handwritten
conditions of cote~134-8.  The full tensor, its scalar unit, and the
interval remain separate requirements.

\subsection{F-SM36: no normalized contraction of the double-dual quotient}

Let \(k\) be a field, let \(V^*=\operatorname{Hom}_k(V,k)\), and
put \(E(V)=V^{**}/\operatorname{ev}_V(V)\).  There is no family of
linear maps
\[
 c_{V,W}:E(V\otimes_k W)\longrightarrow E(V)\otimes_k W
                                                               \tag{SM36.1}
\]
natural in both vector spaces and equal to the identity for
\(W=k\), under the usual unit identifications.  In fact, for any
fixed \(V\) with \(E(V)\ne0\), naturality in \(W\) alone is
already impossible.

\begin{proof}
First such a \(V\) exists.  Take \(V=k^{(\mathbb N)}\), so that
\(V^*=k^{\mathbb N}\).  The quotient
\(k^{\mathbb N}/k^{(\mathbb N)}\) is nonzero, as witnessed by
the constant sequence \(1\).  Choose a nonzero linear functional
on this quotient and pull it back to \(f\in V^{**}\).
If \(f\) were evaluation at \(v\in V\), its vanishing on every
finite-support coordinate vector of \(V^*\) would force every
coordinate of \(v\) to vanish.  This would give \(f=0\), a
contradiction.  Thus \(e=[f]\in E(V)\) is nonzero.

The rest of the argument applies to any \(V\) and any
\(f\in V^{**}\) with nonzero class \(e\).  Set
\(W=k^{(\mathbb N)}\), with basis \((w_n)\), and identify
\[
        (V\otimes W)^*\simeq\prod_{n\ge0}V^*.
\]
On the subspace \(\bigoplus_{n\ge0}V^*\) of this product,
define the functional
\[
             (\lambda_n)\longmapsto\sum_n f(\lambda_n).
\]
The sum is finite on this subspace.  Extend this functional
linearly to the whole product, obtaining
\(F\in(V\otimes W)^{**}\), and write \(u=[F]\).
For the coordinate functional \(\epsilon_n:W\to k\), the
dual of \(1_V\otimes\epsilon_n\) is the inclusion of the
\(n\)-th copy of \(V^*\) into \(\prod_mV^*\).  Consequently
\[
       E(1_V\otimes\epsilon_n)(u)=e\qquad(n\ge0).             \tag{SM36.2}
\]
Naturality and the stated normalization of \(c\) would give
\[
      (1_{E(V)}\otimes\epsilon_n)c_{V,W}(u)=e
                      \qquad(n\ge0).                        \tag{SM36.3}
\]
Every element of \(E(V)\otimes k^{(\mathbb N)}\), however,
has only finitely many nonzero coordinates.  Equations
(SM36.3) contradict \(e\ne0\).  The linear extensions used
here are ordinary algebraic extensions of linear functionals;
no topology or continuity is involved.
\end{proof}

\paragraph{Consequence for the actual additive category.}
Write \(\mathcal N\) for F-SM34's category of
\[
 (X_0,X_1,b:X_0\to X_1,a:E(X_1)\to X_0),
 \qquad ba=0,\quad aE(b)=0,
\]
and \(\kappa(W)=(W,0)\).  There is no unital bifunctor
\(\star:\mathcal N\times\mathcal N\to\mathcal N\) with
unit \(\kappa(k)\) which has natural identifications
\[
 U_i(X\star\kappa(W))\simeq U_i(X)\otimes_kW
                  \quad(i=0,1),                             \tag{SM36.4}
\]
compatible with the canonical vector-space unit identifications.
Here \(U_i\) extracts the indicated component.  This conclusion
does not even require associativity or symmetry of \(\star\).

\begin{proof}
For every \(V\) the object
\(P_1(V)=(E(V),V,0,\mathrm{id})\) belongs to \(\mathcal N\),
and depends functorially on \(V\).  Under (SM36.4), the
\(a\)-operation of \(P_1(V)\star\kappa(W)\) is a map
exactly of the form (SM36.1).  Functoriality makes it natural
in \(V,W\).  Compatibility of the unit isomorphism, which is
an isomorphism in \(\mathcal N\), makes its value at \(W=k\)
the identity on \(E(V)\).  F-SM36 rules this out.
\end{proof}

This leaves open a tensor whose underlying vector spaces are
enlarged.  In particular, (SM36.4) is not a consequence of the
mere bilinearity written in the original scan.  A tensor
preserving scalar copowers in \(\mathcal N\) would involve
the actual coproducts of \(\mathcal N\), whose components
are not generally ordinary direct sums.

\subsection{F-SM37: the arrow monad has no usual tensorial strength}

Let \(\mathcal A=\operatorname{Fun}([0\to1],\operatorname{Vect}_k)\).
Regard \(C=(C_0\xrightarrow{b}C_1)\) as a cochain complex
in degrees \(0,1\), and give \(\mathcal A\) its truncated
cochain tensor
\[
 \begin{aligned}
 (C\odot D)_0&=C_0\otimes D_0,\\
 (C\odot D)_1&=(C_1\otimes D_0)\oplus(C_0\otimes D_1),\\
 b_{C\odot D}&=(b_C\otimes1,\ 1\otimes b_D).
 \end{aligned}                                               \tag{SM37.1}
\]
Its unit is \(I=(k\to0)\).  Associativity and symmetry are
the restrictions of the cochain tensor: in an arbitrary finite
product the degree-zero term has only degree-zero factors, and
the degree-one term has precisely one degree-one factor.

For any nonzero additive exact endofunctor \(E\) of vector
spaces, define
\[
 H(C)=(E(\operatorname{coker}b_C)\to0),\qquad
 T=\mathrm{id}_{\mathcal A}\oplus H.                         \tag{SM37.2}
\]
Then \(H^2=0\), so \(T\) is the square-zero monad whose unit
is the first inclusion and whose multiplication adds the two
\(H\)-summands.  For \(E(V)=V^{**}/V\), its algebra category
is precisely \(\mathcal N\) above.  Nevertheless there is no
natural family
\[
       s_{C,D}:T(C)\odot D\longrightarrow T(C\odot D)
                                                               \tag{SM37.3}
\]
which is the identity when \(D=I\).  Thus \(T\) cannot be
a strong monad for (SM37.1); in particular this model cannot
obtain a tensor by invoking a commutative strong monad on that
base category.

\begin{proof}
Since the cokernel of \(E(\operatorname{coker}b)\to0\) is
zero, \(H^2=0\).  The monad identities follow by adding the
three copies of \(H(C)\) in \(T^3(C)\); all composites of two
\(H\)'s vanish.  An algebra structure has the form
\(1_C\oplus\bar a:C\oplus H(C)\to C\).  The requirement
that it be an arrow morphism is \(b_C\bar a=0\), and the
monad multiplication condition is automatic because \(H^2=0\).
Exactness of \(E\) identifies a map
\(E(\operatorname{coker}b_C)\to C_0\) with a map
\(a:E(C_1)\to C_0\) satisfying \(aE(b_C)=0\).
This gives the stated algebra-category identification, including
its morphisms.

Choose \(V\) with \(E(V)\ne0\), and write
\[
 S=(0\to V),\qquad D(V)=(V\xrightarrow{1}V),\qquad
 D=D(k).
\]
Let \(j:S\to D(V)\) be zero in degree zero and the identity
in degree one.  Let \(q:D\to I\) be the identity in degree
zero and zero in degree one.  Directly from (SM37.1),
\[
 \begin{aligned}
 S\odot D&=S, & S\odot q&=1_S,\\
 T(S)&=(E(V)\xrightarrow{0}V), & T(D(V))&=D(V).
 \end{aligned}
                                                               \tag{SM37.4}
\]
Naturality of \(s\) with respect to \(q\), and the
normalization at \(I\), force
\[
                    (s_{S,D})_0=1_{E(V)}.                    \tag{SM37.5}
\]
On the other hand,
\[
 D(V)\odot D=(V\xrightarrow{v\mapsto(v,v)}V\oplus V).
\]
The map \(j\odot1_D\) sends the degree-one term \(V\)
to the first summand of \(V\oplus V\).  Its induced map
on cokernels is the isomorphism
\[
       V\longrightarrow(V\oplus V)/\operatorname{diag}(V).
\]
It follows that the degree-zero map of \(T(j\odot1_D)\)
is an isomorphism from \(E(V)\) onto the \(H\)-summand of
\(T(D(V)\odot D)_0\), and hence is nonzero.
But \(T(j)_0:E(V)\to V\) is zero.  The naturality square
for \(j\) and \(D\) would therefore give
\[
 T(j\odot1_D)_0(s_{S,D})_0
   =(s_{D(V),D})_0(T(j)\odot1_D)_0=0,
\]
contradicting (SM37.5).  No associativity axiom for a strength
was used: naturality and its unit axiom already contradict
each other.
\end{proof}

\paragraph{What is still open.}
F-SM36 excludes an ordinary componentwise scalar tensor, and
F-SM37 excludes the indicated strong-monad construction.
Neither excludes a different tensor on \(\mathcal N\), a
different monoidal base for an infinitary theory, or another
category satisfying all the original assumptions.  Conversely,
neither proves that the original interval forces the missing
derived normalization on infinite scalar modules.  These are
the exact remaining alternatives.

\paragraph{Bounded literature check on a different proposed route.}
An ordinary Hopf-algebra structure does not automatically supply
the needed tensor on a contramodule category.  Johnston (2026,
Section~2.3) explicitly constructs an action involving a comodule
and a contramodule, using \(\operatorname{Hom}_k\), rather than
a tensor of two arbitrary contramodules.  Consequently a
completed Hopf-algebra candidate still needs its own construction
of the symmetric tensor, its trivial-module unit, and both
scalar adjoints.  This observation is not a nonexistence theorem
for all such completed constructions.

\paragraph{Reference.}
Johnston, D. (2026). Contramodules for algebraic groups: The
existence of mock projectives. \emph{Canadian Mathematical
Bulletin, 69}(1), 184--200. Published online August~4, 2025.
\url{https://doi.org/10.4153/S0008439525100933}.
Section~2.3, introductory paragraph and displayed
\(\operatorname{Hom}_k\) bifunctor.  The elementary proofs
F-SM36--37 above do not use that paper.

\paragraph{Two further repairs tested; no new original-status conclusion.}
The intake's three-sort proposal refers to an unavailable full
\nolinkurl{C.tex} and an unrun coherence script.  For the radical-square-zero
functor \(H(U_0,U_1,U_2)=(E(U_2),U_0,U_1)\), its free projectives are
\(P_0(V)=(V,V,0)\), \(P_1(V)=(0,V,V)\), and
\(P_2(V)=(EV,0,V)\), with the corresponding identity structure maps.
Merely substituting these into the standard truncated-cochain tensor
does not work.  For a nonzero ghost
\(a:P_0(k)\to P_2(W)\) and the map
\(b:P_2(W)\to P_1(W)\) that is the identity in the last component,
one has \(ba=0\).  Tensoring with \(P_0(k)\) under the proposed
standard disk rules gives \((a,0)\) followed by
\(\binom{b}{\pm1}\).  Their composite has nonzero second component
\(\pm a\).  This excludes only these specific rules, not the
unavailable intake construction.

There is a separate defect in the one-sort square-zero repair.
For an exact linear functor \(E\) on vector spaces, the monad
\(T=\mathrm{id}\oplus E\) with zero \(E^2\)-multiplication has
the usual commutative square-zero strengths.  An algebra is a map
\(a:EV\to V\) with \(aE(a)=0\).  Suppose \(E\) vanishes on
finite-dimensional spaces and is nonzero on some \(V\).
The trivial algebra \(\kappa(k)\) is free, since \(Ek=0\).
If the trivial-algebra embedding \(\kappa\) had a right adjoint,
it would preserve coproducts; its value on every free vector space
\(V\) would therefore have to be the free algebra \(T(V)\).
But \(T(V)\) has nonzero structure map when \(EV\ne0\), while
\(\kappa(V)\) has zero structure map.  Thus this repair lacks
one of the required scalar adjoints.

Its usual algebra tensor also fails scalar compatibility.  The
tensor of two trivial algebras is the quotient of the free algebra
on \(V\otimes W\) by the images of the two strengths; explicitly
its underlying vector space is
\[
 (V\otimes W)\oplus E(V\otimes W)/R,\qquad
 R=\operatorname{im}(EV\otimes W)
       +\operatorname{im}(V\otimes EW).
\]
The structure map on the second summand is the quotient map from
\(E(V\otimes W)\); the square-zero condition adds no further
relation there.  For \(W=V\), choose a basis \((v_i)\) and define
\(\delta(v_i)=v_i\otimes v_i\) and
\(m(v_i\otimes v_j)=\delta_{ij}v_i\).  Then \(m\delta=1\).
For each fixed \(v\in V\), the maps \(m(-\otimes v)\) and
\(m(v\otimes-)\) have finite-dimensional image.  Thus \(E(m)\)
kills both strength images, whereas
\(E(m)E(\delta)=1_{EV}\).  Consequently
\(E(V\otimes V)/R\ne0\) if \(EV\ne0\), and the algebra tensor
cannot be \(\kappa(V\otimes V)\).  These two failed repairs do
not justify adding tensor exactness, closedness or AB4 to the
original finite-additivity assumptions.

% END source_cote1348_tensor_obstructions.tex

% BEGIN source_cote1348_ext_tensor.tex
\section{A tensor obstruction detected by the scalar two-extensions}

\subsection{F-SM38: a countable scalar tensor cannot be exact in F-SM34}

Let \(\mathcal N\), \(E(V)=V^{**}/V\), and
\(\kappa(V)=(V,0)\) be as in F-SM34.  Suppose that
\(\mathcal N\) has an additive monoidal product with unit
\(\kappa(k)\), and that \(\kappa\) is strong monoidal:
\[
   \kappa(V)\otimes\kappa(W)\simeq\kappa(V\otimes_kW),
                                                               \tag{SM38.1}
\]
with the usual unit compatibility.  Then the functor
\[
       -\otimes\kappa(k^{(\mathbb N)}):
                   \mathcal N\longrightarrow\mathcal N       \tag{SM38.2}
\]
cannot be exact.  This uses only (SM38.1), not an
identification of the tensor of an arbitrary object with a
scalar object as a copower or as a componentwise tensor.
Closedness and symmetry are not needed.

The same conclusion holds for the external connective chain
category \(\mathcal M=\operatorname{Ch}_{\ge0}(\mathcal N)\)
and \(\kappa_{\mathcal M}(V)=\kappa(V)[0]\), for any monoidal
product with this unit and the analogous scalar compatibility.
Thus the actual additive interval example of F-SM34 cannot be
completed by a tensor which is exact even in this one countable
free scalar factor.  If the proposed tensor is right exact,
that scalar object is necessarily nonflat.

\begin{proof}
\emph{1. The canonical two-extension.}
Retain the functorial projective objects
\[
 P_0(V)=(V,V,1,0),\qquad P_1(V)=(E(V),V,0,1),
 \qquad S(V)=(0,V)
\]
from F-SM34.  There is a natural exact two-extension
\[
 \delta_V:\quad
 0\longrightarrow\kappa E(V)\longrightarrow P_1(V)
   \longrightarrow P_0(V)\longrightarrow\kappa(V)
   \longrightarrow0.                                      \tag{SM38.3}
\]
The middle map is zero in component zero and the identity
in component one.  Its kernel is \(\kappa E(V)\), its image
is \(S(V)\), and its cokernel is \(\kappa(V)\).

Dimension shifting through the two short exact sequences in
(SM38.3) gives natural identifications
\[
 \begin{aligned}
 \operatorname{Ext}^2_{\mathcal N}(\kappa V,\kappa U)
  &\simeq\operatorname{Ext}^1_{\mathcal N}(S(V),\kappa U)\\
  &\simeq\operatorname{Hom}_{\mathcal N}(\kappa E(V),\kappa U)
   \simeq\operatorname{Hom}_k(E(V),U).
 \end{aligned}                                               \tag{SM38.4}
\]
For the second identification,
\(\operatorname{Hom}_{\mathcal N}(P_1(V),\kappa U)=0\),
because \(P_1(V)\) represents the component-one functor,
whereas that component of \(\kappa U\) is zero.
The higher Ext groups of \(P_0(V)\) and \(P_1(V)\) vanish
by their projectivity.  These observations justify both
dimension-shifting isomorphisms, without any condition on
injective objects or infinite coproducts.  Under (SM38.4),
\(\delta_V\) is the identity of \(E(V)\): the connecting
map sends a homomorphism \(E(V)\to U\) to the pushout of
(SM38.3) along that homomorphism.  All identifications are
natural, contravariantly in \(V\) and covariantly in \(U\).

\emph{2. Exact tensoring would produce the prohibited contractions.}
Put \(W=k^{(\mathbb N)}\), and assume that (SM38.2) is
exact.  Tensoring (SM38.3) with \(\kappa W\) then produces
a genuine exact two-extension.  Using (SM38.1) and (SM38.4),
its class corresponds to a linear map
\[
          c:E(V\otimes W)\longrightarrow E(V)\otimes W.
                                                               \tag{SM38.5}
\]
For each coordinate functional \(\epsilon_n:W\to k\),
the natural map given by tensoring with \(\kappa(\epsilon_n)\)
is a morphism from this two-extension to \(\delta_V\).
The right unit isomorphism identifies the latter extension
with (SM38.3).  On the two endpoints these morphisms are
\(\kappa(1_{E(V)}\otimes\epsilon_n)\) and
\(\kappa(1_V\otimes\epsilon_n)\), respectively.
Naturality of (SM38.4), or equivalently equality of the
pushout and pullback classes of this morphism of extensions,
therefore gives
\[
    (1_{E(V)}\otimes\epsilon_n)c
                    =E(1_V\otimes\epsilon_n)\qquad(n\ge0).
                                                               \tag{SM38.6}
\]
Choose \(V\) and \(u\in E(V\otimes W)\) as in the complete
proof of F-SM36, so that the right side of (SM38.6), evaluated
at \(u\), is the same nonzero \(e\in E(V)\) for every
\(n\).  Then \(c(u)\) has infinitely many nonzero
coordinates in \(E(V)\otimes k^{(\mathbb N)}\), an
impossibility.  This proves the first assertion.  Only the
single functor (SM38.2) was assumed exact; tensoring with the
unit is automatically exact because it is isomorphic to the
identity.

\emph{3. Passage to the actual chain-category interval example.}
For a projective \(P\in\mathcal N\), the complex \(P[0]\)
is projective in \(\operatorname{Ch}_{\ge0}(\mathcal N)\):
maps from it are \(\operatorname{Hom}_{\mathcal N}(P,C_0)\),
and degree-zero evaluation is exact.  Applying \([0]\) to a
projective resolution in \(\mathcal N\) is a projective
resolution in this chain category, with exactness checked
degree by degree.  Consequently
\[
 \operatorname{Ext}^2_{\mathcal M}(\kappa V[0],\kappa U[0])
          \simeq\operatorname{Hom}_k(E(V),U),                 \tag{SM38.7}
\]
with the same universal class \(\delta_V[0]\).  The
argument of Step~2 applies verbatim to any tensor on
\(\mathcal M\) having the stipulated unit and scalar
compatibility.  No tensor inherited from \(\mathcal N\)
was assumed in this last assertion.
\end{proof}

\paragraph{Relation to the original question.}
The handwritten assumptions do not explicitly require
exactness of (SM38.2).  Thus F-SM38 is a stronger obstruction
to constructing a tensor on F-SM34, and it does not decide
the remaining abstract bilinear-tensor question.  Its input
is scalar strong-monoidal compatibility alone; it is distinct
from F-SM35's argument using exact copowers and products.

% END source_cote1348_ext_tensor.tex

% BEGIN source_cote1348_chain_counit.tex
\section{The raw counit for categories already of connective-chain form}

\subsection{F-SM39: coinduced disks remove the scalar-dimension restriction}

Let \(k\) be an arbitrary commutative ring, and let \(R\) be a
small \(k\)-linear category such that every \(R(a,b)\) is a
projective \(k\)-module; finite generation is not assumed.
Put \(\mathcal N=\operatorname{Mod}\!-\!R\), with right-module
convention, and suppose that
\[
 \varphi_{\mathcal N}\dashv\kappa_{\mathcal N}\dashv
 \psi_{\mathcal N},\qquad
 \kappa_{\mathcal N}:\operatorname{Mod}(k)\longrightarrow\mathcal N
                                                               \tag{SM39.1}
\]
are \(k\)-linear adjunctions with \(\kappa_{\mathcal N}\) exact.
Assume the canonical derived counit at the scalar ring is a
quasi-isomorphism:
\[
       L\varphi_{\mathcal N}(\kappa_{\mathcal N}k)
                         \longrightarrow k[0].              \tag{SM39.2}
\]
Set
\[
 \begin{aligned}
 \mathcal M&=\operatorname{Ch}_{\ge0}(\mathcal N),\\
 \kappa_{\mathcal M}(U)&=\kappa_{\mathcal N}(U)[0],\\
 \varphi_{\mathcal M}(X)&=\varphi_{\mathcal N}(H_0X).
 \end{aligned}                                               \tag{SM39.3}
\]
All complexes in this statement are homological.

Use the standard compact projective generators of \(\mathcal M\):
if \(P_a=R(-,a)\), take
\[
 Q_{a,0}=P_a[0],\qquad Q_{a,m}=D_m(P_a)\quad(m\ge1),
                                                               \tag{SM39.4}
\]
where \(D_m(A)\) has \(A\) in degrees \(m,m-1\) and identity
differential.  Let \(S\) be their full small \(k\)-linear
category, identify \(\mathcal M\) with right \(S\)-modules,
and let \(D(s)=\varphi_{\mathcal M}(Q_s)\) be the left
coefficient module.  Choose any resolution
\(B_\bullet\to D\) by projective left \(S\)-modules.
Form the exact adjunction of BR-13405:
\[
 \begin{aligned}
 \Lambda(X)&=B_\bullet\otimes_S X,\\
 (GC)(s)&=\operatorname{Hom}_{\operatorname{Ch}_{\ge0}(k)}
                                  (B_\bullet(s),C).
 \end{aligned}                                               \tag{SM39.5}
\]
Then the \emph{ordinary, unreplaced} counit
\[
                    \epsilon_C:\Lambda G(C)\longrightarrow C
                                                               \tag{SM39.6}
\]
is a quasi-isomorphism for every connective complex \(C\).
Consequently this very adjunction induces
\[
 \mathcal M[W^{-1}]\simeq D_{\ge0}(k),\qquad
          W=\Lambda^{-1}(\mathrm{qis}).                     \tag{SM39.7}
\]
No finite injective or pure-injective dimension for scalar
projectives, no finitely generated scalar morphism modules,
and no condition \(D\in\mathrm{FP}_\infty\) is imposed.

\begin{proof}
\emph{1. Generators, exactness, and the derived functor.}
The generators (SM39.4) satisfy
\[
       \operatorname{Hom}_{\mathcal M}(Q_{a,m},X)=X_m(a).
                                                               \tag{SM39.8}
\]
These evaluation functors are exact and preserve coproducts,
so the \(Q_{a,m}\) are compact projective generators.
Presentations by their coproducts give the stated restricted
Yoneda equivalence with \(\operatorname{Mod}\!-\!S\).
Formula (SM39.8) also shows that every \(S(s,t)\) is either
zero or a scalar morphism module of \(R\); in particular it
is projective over \(k\).

The scalar adjunctions of \(\mathcal M\) are
\[
 \varphi_{\mathcal M}\dashv\kappa_{\mathcal M}\dashv
        (\psi_{\mathcal N}\circ\operatorname{ev}_0).
                                                               \tag{SM39.9}
\]
For the first adjunction, a chain map into an object in degree
zero factors uniquely through \(H_0\); for the second, a
chain map from an object in degree zero is an arbitrary map
to the degree-zero term.  The middle functor is exact.
It follows that \(\varphi_{\mathcal M}\) sends projectives
to projective scalar modules.  Thus each \(D(s)\) is
projective over \(k\).

The value \(B_\bullet(s)\to D(s)\) is a projective
scalar resolution of a projective scalar module.  Splitting
the augmentation and then its successive kernels, starting
in degree zero, decomposes \(B_\bullet(s)\) into its
projective degree-zero homology and projective disks.
It is therefore a projective object of
\(\operatorname{Ch}_{\ge0}(k)\).  This proves that \(G\)
is exact, by (SM39.5).  Each \(B_n\) is flat over \(S\),
so \(\Lambda\) is exact as a functor to the abelian category
of complexes.  The right exact coproduct-preserving functor
\(\varphi_{\mathcal M}\) is \(D\otimes_S-\): the two
functors agree on the compact projective generators and
their coproducts, and hence on their cokernel presentations.
Consequently
\[
                 \Lambda(X)\simeq L\varphi_{\mathcal M}(X)
                    \quad\hbox{in }D_{\ge0}(k).              \tag{SM39.10}
\]

We will need the further identification
\[
       L\varphi_{\mathcal M}(X)
                         \simeq L\varphi_{\mathcal N}(X),     \tag{SM39.11}
\]
where the right side applies the derived functor to the
\(\mathcal N\)-complex \(X\).  Here is a finite-diagonal
proof.  Resolve \(X\), as an object of \(\mathcal M\), by
projective chain objects \(A_r\), \(r\ge0\).  A projective
chain object is a retract of a coproduct of the generators
(SM39.4); hence its terms and its \(H_0\) are projective
in \(\mathcal N\), and its positive homology is zero.
The bicomplex \((A_r)_m\) has finite total diagonals.
Exactness of the resolution in every fixed degree gives
\(\operatorname{Tot}(A_\bullet)\to X\) a quasi-isomorphism.
The row augmentations give a second quasi-isomorphism
\[
          \operatorname{Tot}(A_\bullet)
                            \longrightarrow H_0(A_\bullet).
\]
This follows by taking row homology, which is concentrated
in degree zero.  Thus \(H_0(A_\bullet)\) is a connective
projective \(\mathcal N\)-complex representing \(X\), and
applying \(\varphi_{\mathcal N}\) to it computes both sides
of (SM39.11).  The augmentations make the identification
compatible with the canonical scalar counits.  In particular,
if \(X\) is contractible as an outer \(\mathcal N\)-complex,
then
\[
                            L\varphi_{\mathcal M}(X)=0.       \tag{SM39.12}
\]

\emph{2. The coinduced objects are outer disks.}
For \(a\in R\) and a scalar module \(U\), let
\[
             I_a(U)(b)=\operatorname{Hom}_k(R(a,b),U).
\]
It is the right adjoint to evaluation at \(a\) in
\(\mathcal N\).  For every \(m\ge0\), the right adjoint
to degree-\(m\) evaluation on connective homological
complexes is \(D_{m+1}\): a map
\(X\to D_{m+1}(A)\) is uniquely determined by its degree-\(m\)
component \(X_m\to A\).  Therefore the right adjoint to
(SM39.8) is
\[
                        D_{m+1}(I_a(U)).                    \tag{SM39.13}
\]
Equivalently, under the right \(S\)-module description,
\[
 \operatorname{Hom}_k(S((a,m),-),U)
                   \simeq D_{m+1}(I_a(U)).                  \tag{SM39.14}
\]
Every object in (SM39.14) is contractible in the outer chain
direction.  Products of such objects are contractible:
take the product of their contracting homotopies degree by
degree.  This uses products of maps and their homotopy
identities, not exactness of products.  A retract of a
contractible complex is contractible as well: if \(rs=1\)
and \(h\) contracts the ambient complex, then \(rhs\)
contracts the retract.

If \(L\) is any projective left \(S\)-module, it is a
retract of a coproduct of left representables.  Hence
\(\operatorname{Hom}_k(L(-),U)\) is a retract of a product
of (SM39.14), so it is outer contractible.  In particular,
for every \(n\ge1\) and every scalar module \(U\),
\[
 \begin{aligned}
 G(D_n(U))
   &=\operatorname{Hom}_{\mathrm{Ch}}(B_\bullet,D_n(U))\\
   &\simeq\operatorname{Hom}_k(B_{n-1}(-),U)
                 \quad\hbox{is outer contractible}.
 \end{aligned}                                               \tag{SM39.15}
\]
The chain-map identity follows because a map into a disk
is determined by its degree-\((n-1)\) component; the
degree-\(n\) component is its composite with the source
differential.  There is no injectivity assumption on \(U\).

\emph{3. The canonical scalar normalization holds on projectives.}
The functors \(\varphi_{\mathcal N}\) and
\(\kappa_{\mathcal N}\) preserve coproducts, each being a
left adjoint in (SM39.1).  Coproducts are exact in the module
category \(\mathcal N\).  Taking a coproduct of projective
resolutions therefore computes \(L\varphi_{\mathcal N}\)
on a coproduct, and shows that this derived functor preserves
coproducts.  The direct sum of the canonical maps (SM39.2)
gives the canonical quasi-isomorphism for every free scalar
module \(k^{(A)}\).  Taking retracts gives, for every
projective scalar module \(Q\),
\[
       L\varphi_{\mathcal M}(\kappa_{\mathcal M}Q)
                \simeq L\varphi_{\mathcal N}(\kappa_{\mathcal N}Q)
                         \longrightarrow Q[0]               \tag{SM39.16}
\]
as a quasi-isomorphism.

On the other hand, chain maps into \(Q[0]\) factor through
the degree-zero homology of \(B_\bullet\).  Thus
\[
          G(Q[0])=\operatorname{Hom}_k(D(-),Q)
                         \simeq\kappa_{\mathcal M}Q.         \tag{SM39.17}
\]
The counit (SM39.6) for \(Q[0]\) is exactly (SM39.16):
it is evaluation after the augmentation \(B_\bullet\to D\).
In particular it is the canonical counit, not merely some
isomorphism between its source and target in the derived
category.

\emph{4. Test every projective target, without commuting \(G\) with sums.}
Let \(P\) be a projective object of
\(\operatorname{Ch}_{\ge0}(k)\).  The projective-chain
decomposition BR-13402 gives projective scalar modules
\(Q_0,Q_1,Q_2,\ldots\) and an isomorphism
\[
          P\simeq Q_0[0]\oplus\prod_{n\ge1}D_n(Q_n).
                                                               \tag{SM39.18}
\]
The displayed product of disks is also their coproduct:
only two disks contribute in any fixed degree.
Because \(G\) is a right adjoint, it preserves this product.
By (SM39.15), its value on the product of disks is an outer
contractible object \(K\in\mathcal M\).  Therefore
\[
                 G(P)\simeq\kappa_{\mathcal M}Q_0\oplus K.
                                                               \tag{SM39.19}
\]
Equations (SM39.10)--(SM39.12) show that \(\Lambda K\) is
acyclic.  The projection in (SM39.18) and its image under
\(\Lambda G\) are consequently quasi-isomorphisms.  The
counit square for that projection commutes, and its bottom
map is the quasi-isomorphism (SM39.16).  Two-out-of-three
proves that \(\epsilon_P\) is a quasi-isomorphism.
Only a finite biproduct was passed through \(\Lambda\) in
(SM39.19); no product-preservation assertion about
\(\Lambda\) or coproduct-preservation assertion about \(G\)
has entered.

\emph{5. Extend the test to all targets and descend the adjunction.}
For completeness, the projective-target criterion F-SM17
applies directly as follows.  Resolve an arbitrary \(C\),
as an object of the abelian category of connective
complexes, by projective chain objects \(P_r\).
Both the identity functor and \(\Lambda G\) are exact.
The augmented first-quadrant bicomplexes therefore give
quasi-isomorphisms
\[
 \operatorname{Tot}(\Lambda G(P_\bullet))\longrightarrow\Lambda G(C),
 \qquad
 \operatorname{Tot}(P_\bullet)\longrightarrow C.
\]
The map between the total complexes is a quasi-isomorphism
because its map on each row is \(\epsilon_{P_r}\), already
proved to be one.  All total diagonals are finite, so no
infinite-totalization issue or continuity assumption is
needed.  Naturality and two-out-of-three give (SM39.6) for
every \(C\).

If \(f\) is a quasi-isomorphism of scalar complexes,
naturality of these counits shows that \(\Lambda G(f)\)
is a quasi-isomorphism.  Thus \(G\) is homotopical for the
specified weak equivalences.  The triangle identity
\(\epsilon_{\Lambda X}\Lambda(\eta_X)=1\) shows that
each unit \(\eta_X:X\to G\Lambda X\) is a weak
equivalence.  Both ordinary functors therefore descend,
and their unit and counit become isomorphisms, proving
(SM39.7).
\end{proof}

\paragraph{Scope and comparison with the earlier sufficient conditions.}
This proves the literal counit for the resolution adjunction
(SM39.5), not only the existence of some inverse to the
localized realization.  It applies to all locally
\(k\)-projective ringoids \(R\) in (SM39.1), including those
with infinitely generated scalar morphism modules, because
the extra connective-chain structure supplies the explicit
coinduced contractions.  It does not establish this
structure for every category in the original corollary.
The general raw-counit question remains outside this class.

The earlier example F-SM28 is a consistency check: its
\(\operatorname{Fun}(\mathbb N,\operatorname{Ch}_{\ge0}(k))\)
can be written in the form (SM39.3).  Its raw counit was
already covered by F-SM24's finitely generated scalar
morphism case, so it is not counted here as a newly solved
example.  The present proof imposes none of that condition,
F-SM27's \(\mathrm{FP}_\infty\) condition, or F-SM24's
scalar pure-injective-dimension bound.

% END source_cote1348_chain_counit.tex

% BEGIN source_cote1348_product_interval.tex
\section{A product-preserving interval and the general module-category counit}

\subsection{F-SM40: the interval functor, rather than the resolution, can supply finiteness}

Let \(k\) be a commutative ring, let \(R\) be a small
\(k\)-linear category with projective scalar morphism modules,
and put \(\mathcal M=\operatorname{Mod}\!-\!R\).
Retain the scalar adjunctions, exact scalar functor, canonical
normalization, and exact resolution adjunction of BR-13403--05:
\[
 \begin{gathered}
 \varphi\dashv\kappa\dashv\psi,\qquad
 \kappa:\operatorname{Mod}(k)\longrightarrow\mathcal M
       \text{ exact},\qquad
 L\varphi(\kappa k)\longrightarrow k[0]\text{ a quasi-isomorphism},\\
 \varphi=D\otimes_R-,\qquad B_\bullet\longrightarrow D
       \text{ a projective left-}R\text{ resolution},\\
 \Lambda X=B_\bullet\otimes_R X,\qquad
 (GC)(x)=\operatorname{Hom}_{\operatorname{Ch}_{\ge0}(k)}
                                  (B_\bullet(x),C).
 \end{gathered}                                                \tag{SM40.1}
\]
The normalization map is the canonical derived counit, not
an unspecified isomorphism of objects.

Suppose there is a \(k\)-linear endofunctor \(K\) of
\(\mathcal M\) and a natural transformation
\(\alpha:\mathrm{Id}\to K\) such that
\[
 \alpha_X\text{ is monic and }L\varphi(KX)=0
                 \quad\text{for every }X\in\mathcal M.       \tag{SM40.2}
\]
Assume in addition that \(K\) preserves coproducts,
cokernels, and products.  In the original interval setup,
\(K=J\otimes_{\mathcal M}-\); thus this is an explicit
condition on the actual interval functor.  Exactness of
\(K\) is not required.

Then the ordinary counit \(\Lambda G(C)\to C\) is a
quasi-isomorphism for every connective scalar complex.
Consequently this very adjunction induces
\[
 \mathcal M[W^{-1}]\simeq D_{\ge0}(k),\qquad
                  W=\Lambda^{-1}(\mathrm{qis}).              \tag{SM40.3}
\]
More strongly, for arbitrary objects \(x_i\in R\) and
arbitrary scalar modules \(U_i\),
\[
 L\varphi\!\left(\prod_i C_{x_i}(U_i)\right)=0,
 \qquad C_x(U)(y)=\operatorname{Hom}_k(R(x,y),U).              \tag{SM40.4}
\]
No scalar dimension or cardinality bound, no finite generation
of \(R(x,y)\), and no \(\mathrm{FP}_\infty\) hypothesis on
\(D\) is imposed.

\begin{proof}
\emph{1. A finite-presentation criterion for tensor functors.}
We first prove the elementary ringoid version needed below.
For a left \(R\)-module \(L\), the functor
\(L\otimes_R-\) preserves products if and only if \(L\)
is finitely presented.  Here a finite presentation uses
finite coproducts of left representables \(R(a,-)\).

For the forward implication, index a family by all pairs
\(i=(a_i,\ell_i)\) with \(\ell_i\in L(a_i)\), and put
\(P_i=R(-,a_i)\).  Product preservation makes
\[
 L\otimes_R\prod_iP_i\longrightarrow
                \prod_i(L\otimes_RP_i)=\prod_iL(a_i)          \tag{SM40.5}
\]
surjective.  A preimage of the tuple \((\ell_i)_i\) is a
finite sum \(\sum_{j=1}^m\ell'_j\otimes w_j\), where
\(\ell'_j\in L(b_j)\) and
\(w_j\in\prod_iR(b_j,a_i)\).  Writing its coordinates
as \(r_{ji}:b_j\to a_i\) gives
\[
                  \ell_i=\sum_{j=1}^m L(r_{ji})(\ell'_j).
\]
Thus finitely many \(\ell'_j\) generate \(L\).  Choose a
finite coproduct \(F\) of left representables and an
epimorphism \(F\to L\), and write \(H\) for its kernel.

Now index a second family by every pair
\(i=(a_i,h_i)\) with \(h_i\in H(a_i)\), again using
\(P_i=R(-,a_i)\).  The tuple \((h_i)_i\), viewed in
\(\prod_i F(a_i)\), comes from a unique element
\(z\in F\otimes_R\prod_iP_i\): tensoring with a finite
coproduct of left representables is a finite coproduct of
evaluation functors, and therefore preserves products.
The image of \(z\) in \(L\otimes_R\prod_iP_i\) has zero
image under (SM40.5).  Injectivity of that comparison map
makes this image zero.  Right exactness of tensoring in
the left-module variable now writes
\[
                  z=\sum_{j=1}^s h'_j\otimes v_j
       \quad\text{in }F\otimes_R\prod_iP_i,
\]
with finitely many \(h'_j\in H(b_j)\).  Taking coordinates
shows that these \(h'_j\) generate every element of \(H\).
Consequently \(H\) is finitely generated and \(L\) is
finitely presented.

Conversely, tensor a finite presentation of \(L\) with
an arbitrary product of right modules.  The two finite
representable sums commute with the product.  Products
of scalar modules are exact, so taking the product of
the corresponding right exact presentations and comparing
cokernels proves the converse.  All index sets used above
are sets because \(R\) is small and its modules have
set-valued components.

\emph{2. Each evaluation of the interval has a finitely presented coefficient.}
For \(x\in R\) define a left \(R\)-module
\[
                   L_x(a)=(K(R(-,a)))(x).                   \tag{SM40.6}
\]
Its left action is induced by maps between right
representables.  There is a natural evaluation pairing
\[
                  L_x\otimes_R X\longrightarrow(KX)(x).
                                                               \tag{SM40.7}
\]
For an element of \(X(a)\), use its associated map
\(R(-,a)\to X\), apply \(K\), and evaluate at \(x\).
The \(k\)-linearity and naturality of this operation make
the pairing balanced.  It is an isomorphism on right
representables.  Since \(K\) preserves coproducts and
cokernels, presentations by coproducts of representables
make (SM40.7) an isomorphism for every \(X\).
This gives a direct proof of the required tensor
description, without an additional representation theorem.

Because \(K\) preserves products, (SM40.7) shows that
\(L_x\otimes_R-\) preserves products.  Step 1 makes
\(L_x\) finitely presented.  Under (SM40.7), the evaluated
natural transformation \(\alpha_X(x)\) is tensoring with
the left-module map
\[
                       u_x:R(x,-)\longrightarrow L_x        \tag{SM40.8}
\]
obtained by evaluating \(\alpha\) on right representables.
The assertion follows first on these representables and
then on their coproduct presentations.

\emph{3. The evaluated interval inclusions split.}
By (SM40.2), tensoring (SM40.8) with every right module
gives a monomorphism.  Thus \(u_x\) is a pure monomorphism
of left \(R\)-modules.  Its quotient \(V_x\) is finitely
presented.  Indeed, take a finite presentation of \(L_x\).
The map from the projective left representable \(R(x,-)\)
lifts to the finite free term of that presentation; adding
this one representable to its relation term gives a finite
presentation of \(V_x\).

A pure short exact sequence with finitely presented
quotient splits.  In this case the finite-equation argument
can be stated explicitly.  Choose finitely many generators
and relations for \(V_x\), and lift the generators to
\(L_x\).  Each relation on the lifts gives an element of
\(R(x,-)\).  Those finitely many linear equations, with
constants in \(R(x,-)\), have a solution in \(L_x\).
Purity gives a solution in \(R(x,-)\).  Subtracting that
solution from the chosen lifts makes all relations vanish
and defines a section \(V_x\to L_x\).  The finite-equation
criterion for purity is the one proved in F-SM24 by
tensoring the transposed finite presentation matrix.
We obtain a left-module retraction
\(r_x:L_x\to R(x,-)\) of \(u_x\).

Tensoring \(r_x\) with \(X\) now splits
\[
                     \alpha_X(x):X(x)\longrightarrow(KX)(x)
                                                               \tag{SM40.9}
\]
as a scalar map, naturally in \(X\).  The choices \(r_x\)
need not be natural in \(x\); no retraction \(K\to\mathrm{Id}\)
of functors into \(\mathcal M\) has been asserted.
This distinction is essential.

\emph{4. Arbitrary products of coinduced coefficients are killed.}
Put \(I=\prod_i C_{x_i}(U_i)\).  Coinduction gives
\[
 \operatorname{Hom}_{\mathcal M}(X,I)
             =\prod_i\operatorname{Hom}_k(X(x_i),U_i).       \tag{SM40.10}
\]
For every monomorphism whose evaluations split over \(k\),
the restriction map in (SM40.10) is surjective: extend each
coordinate by composing with its chosen scalar retraction.
Apply this to \(\alpha_I:I\to KI\), whose evaluations
split by Step 3.  The identity of \(I\) extends to a map
\(KI\to I\).  Hence \(I\) is a retract of \(KI\).
Equation (SM40.2) and functoriality of \(L\varphi\) show
that \(L\varphi(I)\) is a retract of zero, proving
(SM40.4).  Neither injectivity nor pure injectivity of
the coefficients \(U_i\) is involved.

\emph{5. The unreplaced counit.}
If \(P\) is an acyclic projective object of
\(\operatorname{Ch}_{\ge0}(k)\), BR-13402 writes it as
\[
           P=\prod_{n\ge1}D_n(Q_n),\qquad Q_n\text{ projective over }k.
                                                               \tag{SM40.11}
\]
The product is also the coproduct degree by degree, but
we use that \(G\), as a right adjoint, preserves the
product.  The disk adjunction gives
\[
            GP=\prod_{n\ge1}\operatorname{Hom}_k(B_{n-1}(-),Q_n).
                                                               \tag{SM40.12}
\]
Each projective left module \(B_{n-1}\) is a retract of
a coproduct of left representables.  Its scalar Hom
module is therefore a retract of a product of coinduced
modules \(C_x(Q_n)\).  Taking products of these retractions
makes (SM40.12) a retract of one product of such modules.
Step 4 and preservation of retracts give
\(L\varphi(GP)=0\).

For clarity, the canonical normalization in (SM40.1)
also handles every degree-zero projective scalar module.
The adjunction identifies \(G(Q[0])\) with \(\kappa Q\)
and its counit with the scalar derived counit.  The
functors \(\kappa\) and \(\Lambda\) preserve coproducts,
and scalar direct sums preserve quasi-isomorphisms.
The normalization for \(k\) therefore passes to free
modules and then to their projective retracts.

Every projective chain object is a finite biproduct of
such a degree-zero object and an acyclic projective
complex.  Its counit is consequently a quasi-isomorphism.
The exact-adjunction criterion F-SM17 extends this to
every target: resolve the target by projective chain
objects, apply the two exact functors \(\Lambda G\) and
\(\mathrm{Id}\), and compare their first-quadrant total
complexes.  The diagonals are finite, the augmentations
are quasi-isomorphisms, and the counit on every resolution
term is a quasi-isomorphism.  Naturality and the triangle
identities then descend the actual unit and counit to
inverse isomorphisms, proving (SM40.3).
\end{proof}

\paragraph{Exact remaining scope.}
F-SM40 imposes finiteness on the evaluated interval
coefficients \(L_x\), not on \(D\), its resolution terms,
or the scalar modules \(R(x,y)\).  Equivalently, among
coproduct- and cokernel-preserving interval functors,
Step 1 identifies the extra product-preservation condition
with finite presentation of all the \(L_x\).
The split-evaluation conclusion is stronger than the
objectwise purity used in F-SM24.

The handwritten bilinear tensor convention has not been
shown to imply that \(J\otimes-\) preserves products.
Thus this result removes another explicit class from
the general raw-counit gate; it does not declare that
gate closed under the bare original assumptions.
In particular, no passage from an arbitrary left adjoint
to product preservation is used.

% END source_cote1348_product_interval.tex

% BEGIN source_cote1348_split_interval.tex
\section{Split evaluations suffice without product preservation}

\subsection{F-SM42: the relative-injectivity criterion for the actual interval}

Retain precisely the ringoid, scalar normalization, and exact
resolution adjunction (SM40.1).  Thus
\(\mathcal M=\operatorname{Mod}\!-\!R\), the scalar
\(R(x,y)\) are projective, and the canonical scalar
derived counit at \(k\) is a quasi-isomorphism.
Let \(K\) be a \(k\)-linear endofunctor with a natural
map \(\alpha:\mathrm{Id}\to K\).  Suppose
\[
 \begin{gathered}
       L\varphi(KX)=0\quad(X\in\mathcal M),\\
       \alpha_X(x):X(x)\longrightarrow(KX)(x)
          \text{ is a split scalar monomorphism}
                    \quad(X\in\mathcal M,\ x\in R).
 \end{gathered}                                                \tag{SM42.1}
\]
The scalar retractions need not be natural in \(x\), or
even in \(X\).  No preservation of products, coproducts,
cokernels, or exact sequences is required of \(K\).
Then the ordinary counit \(\Lambda G(C)\to C\) is a
quasi-isomorphism for every connective scalar complex.
In particular the specified adjunction induces
\(\mathcal M[W^{-1}]\simeq D_{\ge0}(k)\), where
\(W=\Lambda^{-1}(\mathrm{qis})\), with its actual unit and counit.
Moreover, for arbitrary scalar modules \(U_i\),
\[
                  L\varphi\!\left(\prod_i C_{x_i}(U_i)\right)=0,
 \qquad C_x(U)(y)=\operatorname{Hom}_k(R(x,y),U).              \tag{SM42.2}
\]

\begin{proof}
Put \(I=\prod_iC_{x_i}(U_i)\).  The coinduction adjunction
identifies
\[
 \operatorname{Hom}_{\mathcal M}(X,I)
          =\prod_i\operatorname{Hom}_k(X(x_i),U_i).
                                                               \tag{SM42.3}
\]
The identity of \(I\) therefore corresponds to a family
of scalar maps \(q_i:I(x_i)\to U_i\).  Choose scalar
retractions \(r_i:(KI)(x_i)\to I(x_i)\) of
\(\alpha_I(x_i)\).  The maps \(q_i r_i\) correspond,
by the same adjunction, to an actual \(R\)-module map
\(r:KI\to I\).  Its composite with \(\alpha_I\)
corresponds to the original family \(q_i\), so
\(r\alpha_I=1_I\).  Thus \(I\) is a retract of
\(KI\), and (SM42.1) proves (SM42.2).

For an acyclic projective scalar chain object
\(P=\prod_{n\ge1}D_n(Q_n)\), the disk adjunction gives
\[
               GP=\prod_{n\ge1}\operatorname{Hom}_k(B_{n-1}(-),Q_n).
                                                               \tag{SM42.4}
\]
Each projective \(B_{n-1}\) is a retract of a coproduct
of left representables.  Hence \(GP\) is a retract of
a product of coinduced modules.  Formula (SM42.2) gives
\(L\varphi(GP)=0\).  The canonical scalar normalization
extends to free scalar modules and then projective
retracts, because \(\kappa\) and \(\Lambda\) are left
adjoints and scalar direct sums are exact.  Thus the
counit is a quasi-isomorphism on every projective chain
object, including its degree-zero projective summand.
The exact-adjunction criterion F-SM17, whose extension
uses a first-quadrant resolution with finite diagonals,
gives the result for every target and descends the actual
adjunction.  This is also precisely Step 5 of F-SM40;
its proof uses no further property of \(K\).
\end{proof}

\paragraph{Application to pointwise tensor products.}
Let \(\mathcal C\) be a small ordinary category,
\(R=k[\mathcal C]\), and
\(\mathcal M=\operatorname{Fun}(\mathcal C^{\mathrm{op}},
\operatorname{Mod}k)\), with pointwise tensor and
constant tensor unit \(\underline{k}\).  The scalar
functor is the exact constant functor; its left and
right adjoints are colimit and limit.  Retain the
canonical normalization and full interval acyclicity.
If the interval map
\(\alpha:\underline{k}\to J\) splits over \(k\) at
every object of \(\mathcal C\), then F-SM42 applies
to \(KX=J\otimes_kX\).  Indeed a scalar retraction
\(r_c:J(c)\to k\) gives
\(r_c\otimes1_{X(c)}\) as a retraction of
\(\alpha_X(c)\), for every \(X\).

For example, this scalar splitting holds whenever the
underlying interval is a reduced free presheaf
\[
             J(c)=k[I(c)]/k[s_0(c)],\qquad
                     \alpha_c(1)=[s_1(c)],                  \tag{SM42.5}
\]
where \(I\) has two natural sections \(s_0,s_1\) which
are distinct at every object.  Projecting onto the basis
element \([s_1(c)]\) gives the required scalar
retraction.  This projection is not asserted to be
natural in \(c\).  The bases \(I(c)\) can be infinite;
no finite morphism-set or finite interval hypothesis
has entered.  The normalization and full acyclicity
still have to be verified for a proposed interval;
the free-module description alone does not imply them.

\paragraph{The added range is strict: an actual interval which does not preserve products.}
For every nonzero commutative ring \(k\), take
\(\mathcal M=\operatorname{Ch}_{\ge0}(k)\), with its
usual homological tensor product and scalar unit
\(k[0]\).  In the standard ringoid presentation,
evaluation at a projective generator is evaluation
of a chain degree, and all scalar morphism modules
are zero or \(k\).  The scalar left adjoint is
\(\varphi=H_0\), with
\(L\varphi(X)\simeq X\) by the finite-diagonal
identification (SM39.11) for \(\mathcal N=\operatorname{Mod}k\).
Thus the canonical normalization holds.

Let \(V=k^{(\mathbb N)}\), with basis \(e_0,e_1,\ldots\),
put \(J=D_1(V)\), and take
\(\alpha:k[0]\to J\) to send \(1\) to \(e_0\)
in degree zero.  For every \(X\),
\[
  (J\otimes X)_n=(V\otimes X_n)\oplus(V\otimes X_{n-1}),
        \qquad \alpha_X(x)=(e_0\otimes x,0),                \tag{SM42.6}
\]
with \(X_{-1}=0\).  Projection onto the \(e_0\)
coefficient in the first summand splits this map over
\(k\) in each degree.  It is not claimed to be a chain
retraction.  If \(h:J_0\to J_1\) is the identity of
\(V\), extended by zero in the other degrees, then
\(dh+hd=1_J\).  The map \(h\otimes1_X\) contracts
\(J\otimes X\): the two cross terms involving \(d_X\)
cancel with their Koszul signs.  Therefore the full
interval acyclicity in (SM42.1) holds, including degree
zero.

Nevertheless \(K=J\otimes-\) does not preserve countable
products.  Apply its product comparison to the family
\(X_j=k[0]\), \(j\in\mathbb N\), and inspect degree
zero.  The comparison is
\[
 k^{(\mathbb N)}\otimes_k k^{\mathbb N}
                 \longrightarrow (k^{(\mathbb N)})^{\mathbb N}.
                                                               \tag{SM42.7}
\]
Every element on the left has nonzero coordinates in
only finitely many of the \(V\)-basis directions,
uniformly across all product factors.  The tuple
\((e_j)_{j\in\mathbb N}\) on the right has no such
finite support, and is not in the image.  This proves
failure of product preservation.  Thus F-SM42 genuinely
weakens F-SM40's hypothesis on the actual interval
functor.  This illustrative category was already
covered by F-SM39 and is not counted as a newly solved
example of the original question.

\paragraph{Residual scope.}
The split-evaluation condition is a sufficient condition,
not a converse assertion about the raw counit.  A case
still outside these proofs must have no verified interval
satisfying (SM42.1), in addition to escaping the earlier
scalar, resolution-finiteness, and connective-chain
conditions.  No derivation of scalar splitting from the
bare handwritten bilinearity convention is claimed.

% END source_cote1348_split_interval.tex

% BEGIN source_cote1348_projective_tensor_failure.tex
\section{A composition obstruction to the standard projective tensor proposal}

\subsection{F-SM41: the truncated-disk rules do not define a bifunctor}

Work over a field \(k\) in the two-sort abelian category
\(\mathcal N\) of F-SM34, with
\(E(V)=V^{**}/V\).  Its projectives include
\[
 P_0(V)=(V,V,1,0),\qquad P_1(V)=(E(V),V,0,1).
\]
This finding concerns only the following proposed rules,
motivated by the tensor product of two-term cochain disks:
\[
 \begin{aligned}
 P_0(V)\star P_0(W)&=P_0(V\otimes W)\oplus P_1(V\otimes W),\\
 P_0(V)\star P_1(W)&=P_1(V\otimes W),\\
 P_1(V)\star P_0(W)&=P_1(V\otimes W),\\
 P_1(V)\star P_1(W)&=0.
 \end{aligned}                                                \tag{SM41.1}
\]
Use the standard disk splitting in which the two degree-one
coordinates are \((u,u+v)\).  On a cross map
\(a:P_0(V)\to P_1(W)\), defined by a linear map
\(V\to E(W)\), the proposed rule tensors by the forward
strength \(E(W)\otimes Z\to E(W\otimes Z)\).
On the cross map \(b:P_1(W)\to P_0(W)\) defined by
\(1_W\), use the standard first-factor disk rule.
These rules do not preserve composition, so they do not
define a bifunctor even on the indicated projective objects.

\begin{proof}
Choose \(W\) with \(E(W)\ne0\), for example
\(W=k^{(\mathbb N)}\), as proved in F-SM34 and F-SM36.
Fix \(0\ne e\in E(W)\).  Let
\[
 a:P_0(k)\longrightarrow P_1(W),\qquad
 b:P_1(W)\longrightarrow P_0(W)
\]
have component maps \(a=(1\mapsto e,0)\) and
\(b=(0,1_W)\), respectively.  They satisfy the two-sort
morphism relations, and their composite is zero in both
components:
\[
                              ba=0.                         \tag{SM41.2}
\]

Put \(T=P_0(k)\).  The normalized forward strength at
the scalar unit is the canonical identification
\(E(W)\otimes k=E(W\otimes k)\).  Hence the proposed
map on the first cross arrow is
\[
      a\star1_T=(a,0):P_0(k)\oplus P_1(k)\longrightarrow P_1(W).
                                                               \tag{SM41.3}
\]
For the second cross arrow, its degree-one coordinates
in the unsplit disk tensor are \((w,0)\).  In the chosen
coordinates \((u,u+v)\), these are \((u,v)=(w,-w)\).
Thus the standard rule is
\[
      b\star1_T=\binom{b}{-1}:
                  P_1(W)\longrightarrow P_0(W)\oplus P_1(W).
                                                               \tag{SM41.4}
\]
Multiplying these matrices gives
\[
 (b\star1_T)(a\star1_T)
       =\begin{pmatrix}0&0\\-a&0\end{pmatrix}\ne0.           \tag{SM41.5}
\]
The lower-left component is nonzero because \(a\ne0\);
this remains true in characteristic two.  A bifunctor
additive in the first variable would instead send
\((ba)\star1_T\) to zero by (SM41.2), contradicting
(SM41.5).  Neither an associativity nor a symmetry
coherence check is needed to exhibit the failure.
\end{proof}

\paragraph{Scope.}
The obstruction excludes these particular standard
truncated-disk maps together with the proposed object
rules and normalized forward strength.  It is not a
nonexistence theorem for every possible tensor product
on \(\mathcal N\).  In particular, it does not complete
the additive example into a counterexample to the full
handwritten symmetric-tensor hypotheses.

% END source_cote1348_projective_tensor_failure.tex

% BEGIN root_methods.tex
\section{The geometric method: a usable arithmetic determinant reduction}
\label{sec:root-method}

There are two distinct multiplicities in the prompt.  An automorphic
restriction multiplicity is the dimension of a space of equivariant
functionals.  A Stepanov multiplicity is an order of vanishing or an
integral divisibility order.  A theorem relating them must include an
integral lattice, a valuation, and compatible maps; a complex
representation-theoretic dimension alone does not supply those data.
The following finite-dimensional form isolates the actual numerical
inequality needed by the arithmetic determinant method.

\paragraph{Finding BR-DET-01 (divisibility versus archimedean size).}
Let $K$ be a number field with ring of integers $\mathcal O_K$.
Let $A$ be an $r\times r$ matrix over $\mathcal O_K$.  Suppose its
$j$th column has every entry in a nonzero integral ideal $\mathfrak a_j$,
and for each embedding $\sigma:K\hookrightarrow\mathbb C$ the
Euclidean norm of that column is at most $C_{\sigma,j}>0$.
Count all embeddings, including both members of a complex pair.
If
\[
 \sum_{j=1}^r\log N\mathfrak a_j
   >\sum_{\sigma:K\hookrightarrow\mathbb C}
                     \sum_{j=1}^r\log C_{\sigma,j},             \tag{D.1}
\]
then $\det A=0$.  If $\det A\ne0$, the reverse weak inequality
is necessary.

\begin{proof}
Expansion by permutations shows
$\det A\in\mathfrak a_1\cdots\mathfrak a_r$.
For $a=\det A\ne0$, inclusion of integral ideals
$(a)\subseteq\mathfrak a_1\cdots\mathfrak a_r$ gives
\[
 \prod_jN\mathfrak a_j
 \leq N(a)=|N_{K/\mathbb Q}(a)|
 =\prod_\sigma|\sigma(a)|.
\]
Hadamard's determinant inequality, applied separately at every
embedding, gives
$|\sigma(a)|\leq\prod_j C_{\sigma,j}$.
Multiplying and taking logarithms proves the reverse weak
inequality.  Its contradiction with (D.1) proves vanishing.
\end{proof}

\paragraph{Finding BR-DET-02 (the exact accumulation law for iterated lifts).}
Let $L_j$ be free finite-rank $\mathcal O_K$-lattices with Hermitian
norms at the embeddings, and let
$\theta_j:L_{j-1}\to L_j$ be $\mathcal O_K$-linear maps.
Suppose
\[
 \theta_j(L_{j-1})\subseteq\mathfrak b_jL_j,
 \qquad \|\theta_{j,\sigma}\|\leq c_{j,\sigma}.
\]
Then
\[
 \theta_m\cdots\theta_1(L_0)
   \subseteq(\mathfrak b_1\cdots\mathfrak b_m)L_m,
 \qquad
 \|\theta_{m,\sigma}\cdots\theta_{1,\sigma}\|
   \leq\prod_{j=1}^m c_{j,\sigma}.                            \tag{D.2}
\]
After an integral evaluation map $e:L_m\to\mathcal O_K^r$, columns
obtained from chosen vectors $v_k\in L_0$ obey BR-DET-01 with
the bound
\[
 C_{\sigma,k}=\|e_\sigma\|
       \left(\prod_{j=1}^m c_{j,\sigma}\right)\|v_{k,\sigma}\|,
\]
or any strictly positive upper bound if a column is zero.
Thus a valid application of ``one order
per lift'' requires the ideal-inclusion assertion in (D.2), and
a valid use of arbitrarily many lifts requires their accumulated
logarithmic gain to exceed their accumulated norm cost.

\begin{proof}
For any ideal $\mathfrak b$ and any $\mathcal O_K$-linear map
$\theta$, one has $\theta(\mathfrak bL)=\mathfrak b\theta(L)$:
both sides consist of finite sums of coefficients from
$\mathfrak b$ times images of elements of $L$.
Induction proves the lattice inclusion.  Submultiplicativity of
operator norms proves the other inequality.  An integral
evaluation map preserves the ideal inclusions; applying the
composite operator to $v_k$ bounds its norm by the product of
the operator norms and $\|v_{k,\sigma}\|$. The preceding
finding applies to the resulting matrix.
\end{proof}

These findings are elementary reductions, not a replacement for
Alp\"oge's arithmetic estimates.  They explain exactly which
content the four-page method would have to export in order to
transfer to another motivic family.  CM-period divisibility and
the integral theta structures supply the ideals.  A Rallis
formula and period estimates may control the archimedean norms.
The Iwasawa averaging and relative trace formula must preserve
the necessary independence and integrality while improving the
cost.  The order-$m$ cancellation in the sketch matters because
it changes the sign of the strict inequality (D.1).  Neither
that cancellation nor a uniform tower estimate is a formal
consequence of the five exported RH declarations.

The exact external comparison points are Gan--Gross--Prasad
(2009), Conjectures 17.1 and 17.3 (local distinguished members),
Theorem 14.1 and the reductions in Theorems 15.1 and 16.1
(multiplicity-one statements in their stated cases), and
Conjecture 24.1 (global periods and central values).  These are
the statuses and numbering of the supplied work; we do not call
all of its conjectures theorems from its presence in the ZIP.
Gan--Qiu--Takeda (2014), Theorem 11.4, is a precise Rallis
inner-product formula, with range, cuspidality, and local
conditions.  Applying it repeatedly requires checking those
conditions at every step.  Its special values generally belong
to higher-rank automorphic $L$-functions, outside the stated
Dirichlet scope.  GRH also allows zeros at the centre; it does
not assert that a central period is nonzero.

\subsection{What this does for crystals and for towers}

For a rank-one abelian Artin--Tate object unramified at a
realization prime $p$, its crystalline realization has Frobenius
eigenvalue $p^{-n}$ times the unramified finite character, in the
convention of BR-AT-01.  Its Newton slope is $-n$ and its
Hodge--Tate weight is $-n$.  The period $t$ with $\varphi(t)=pt$ gives this normalization
(Berger, 2004, \S II.3.1 and II.4.3); an unramified finite twist
changes the eigenvalue only by a root of unity.  This local rank-one
compatibility is already known; GRH is a global assertion about the continued
Euler product.  For a ramified finite character one must first
state a potentially crystalline or other appropriate local
realization, rather than silently use the unramified formula.
For the CM curve $y^2=x^3+1$, the separate argument below does
give a genuine new use of Dirichlet GRH: it controls the supply
of inert, trace-zero primes used by the divisibility mechanism.
Its two-dimensional $H^1$ remains outside the diagonal
Artin--Tate Hodge types.

In a cyclotomic tower of conductor $p^a$, the finite Frobenius
recognition theorem below applies at every fixed level, with
$Q=\operatorname{lcm}(p^a,4)$ and bound
$(\varphi(Q)\log Q)^2$.  This is an explicit effective
consequence for the tower's finite abelian layers.  An inverse
limit of integral Selmer groups, a statement about its
characteristic ideal, or a Vandiver-type vanishing assertion
requires additional integral information.  BR-MIX-01 supplies
an explicit example of information lost on passing to the
weight-graded $L$-function; it prevents substituting a
semisimple spectral calculation for that missing integral
calculation.

\subsection{Arithmetic duality is a precise analogy, not a missing construction}

Artin--Verdier duality is a pairing for constructible torsion
\'etale coefficients on number rings, with compact support and
the real-place conventions.  A standard exact source is Milne
(2006), \emph{Arithmetic Duality Theorems}, Chapter II,
Theorem 3.1(b), p.~177.  Its degree-three duality explains the arithmetic
three-dimensional analogy.  It does not equip
$\operatorname{Spec}\mathbb Z$ with a positive Hilbert-space
cohomology whose determinant is $\xi$.  In fact the infinitely
many zeros of $\xi$ already exclude an ordinary characteristic
polynomial of a single finite-dimensional complex operator.
If such a finite-dimensional operator $A$ existed, its
determinant $\det(sI-A)$ would be a polynomial with at most
$\dim A$ zeros, contradicting the known zero count.
An infinite-dimensional space, a trace class transformation
of the variable as in BR-TR-01, or an appropriate regularized
determinant is additional structure.  This is the precise
analytic point at which the arithmetic-duality analogy and a
Hilbert--P\'olya or Deninger construction must be connected by
a new theorem.


% END root_methods.tex

% BEGIN consequences_criteria.tex
\section{Consequences and criteria under the adopted RH and Dirichlet
GRH}\label{consequences-and-criteria-under-the-adopted-rh-and-dirichlet-grh}

Research module \nolinkurl{consequences_criteria}. Author/checker: GPT-6
Astra Pro; effort setting not exposed. Date: 20 November 2026 UTC. The
mathematical input is the RH and Dirichlet-GRH return adopted by the
task. This module does not claim an independent proof or a kernel check
of that return. Known conditional results are used with their actual
hypotheses; an implication asserted here is relative to the adopted
input, and hence has no remaining RH condition within the task's adopted
theory.

The propositions labelled CC-F are deductions and counterexamples
established in this run; no claim of historical priority or exhaustive
novelty is made.

The main distinction is between \textbf{theorems whose hypotheses are
now discharged}, \textbf{equivalent criteria whose asserted positivity
or approximation now follows}, and \textbf{additional geometric or
statistical structures that an equivalence does not supply}. An
input-scope obstruction does not establish logical independence: it
identifies the extra hypothesis in the available theorem and prevents an
unjustified application of that theorem.

\subsection{1. Highest-priority
outcomes}\label{highest-priority-outcomes}

\begin{longtable}[]{@{}
  >{\raggedright\arraybackslash}p{(\columnwidth - 6\tabcolsep) * \real{0.2500}}
  >{\raggedright\arraybackslash}p{(\columnwidth - 6\tabcolsep) * \real{0.2500}}
  >{\raggedright\arraybackslash}p{(\columnwidth - 6\tabcolsep) * \real{0.2500}}
  >{\raggedright\arraybackslash}p{(\columnwidth - 6\tabcolsep) * \real{0.2500}}@{}}
\toprule\noalign{}
\begin{minipage}[b]{\linewidth}\raggedright
ID
\end{minipage} & \begin{minipage}[b]{\linewidth}\raggedright
Outcome
\end{minipage} & \begin{minipage}[b]{\linewidth}\raggedright
Exact operative source
\end{minipage} & \begin{minipage}[b]{\linewidth}\raggedright
What is discharged or remains
\end{minipage} \\
\midrule\noalign{}
\endhead
\bottomrule\noalign{}
\endlastfoot
CC-F01 & Every smooth proper Borger \(\Lambda\)-scheme over
\(\mathbb Z[1/M]\) has cohomological Dirichlet-completed factors
satisfying RH on their weight-shifted lines & Borger (2009), Theorem
0.1, p.~3; Theorem 5.3, p.~17; Corollary 6.2 and (6.2.1), p.~20 &
Dirichlet GRH suffices. This is a statement about realizations and their
\(L\)-data; no Chow-motive decomposition or cycle-surjectivity assertion
is inferred \\
CC-F01b & The rational compactly supported cohomology of every separated
finite-type \(\Lambda\)-scheme over \(\mathbb Z\) actually splits into
finite abelian characters and Tate twists & Borger (2009), §6 setup and
Corollary 6.2, equation (6.2.1), followed by the weight-space descent
proved here & No torsion-freeness hypothesis or semisimplification is
inserted. Dirichlet GRH supplies the resulting primitive factors;
bad-fiber model zeta functions require separate treatment \\
CC-C01 & Effective Chebotarev with the Lagarias--Odlyzko error term for
\(L/K\) whenever the top field \(L/\mathbb Q\) is abelian &
Lagarias--Odlyzko (1977), Theorem 1.1, equation (1.3) & Their exact
analytic hypothesis is RH for \(\zeta_L\), discharged by Dirichlet
factorization. Merely assuming \(L/K\) abelian does not do this \\
CC-C02 & A prime in every reduced class \(a\pmod q\), \(q>3\), below
\((\varphi(q)\log q)^2\) & Lamzouri--Li--Soundararajan (LLS), corrected
arXiv v3 (2016), Corollary 1.2, p.~3 & Dirichlet GRH suffices; a
representative for every class supports finite character tests,
including higher-rank semisimple representations \\
CC-F02 & Siegel's lower bound for primitive real Dirichlet \(L(1,\chi)\)
has a computable constant for every positive rational exponent & LLS,
Theorem 1.5, pp.~3--4, plus the full effective deduction below &
Dirichlet GRH suffices. This is the analytic Siegel theorem; it is not
an asserted effective integral-point theorem \\
CC-F03 & The Hilbert-class-field hypothesis is a real scope boundary
even for an imaginary quadratic base & Explicit splitting field of
\(x^3-x-1\), proved below & The required ideal-class-character
\(L\)-function can be an irreducible degree-two Artin \(L\)-function
over \(\mathbb Q\), outside Dirichlet factorization \\
CC-P01--CC-P04 & Weil positivity, Li positivity, the strong
Nyman--Beurling approximation, and the bare spectral realization all
hold in their precise senses & Exact locators in §4 & Stronger
constructions, support ranges, approximation asymptotics, and geometric
functoriality are separate questions \\
CC-P07 & The concrete de Branges sufficient positivity conditions
analyzed by Conrey--Li are false and remain false & Conrey--Li (2000),
§§3.1--3.2 and §4 Remark & A sufficient condition is not automatically
an equivalent criterion. RH does not imply these conditions \\
\end{longtable}

\subsection{2. The Borger bridge, with a full
proof}\label{the-borger-bridge-with-a-full-proof}

\subsubsection{\texorpdfstring{CC-F01. Smooth proper
\(\Lambda\)-cohomology satisfies its shifted analytic
RH}{CC-F01. Smooth proper \textbackslash Lambda-cohomology satisfies its shifted analytic RH}}\label{cc-f01.-smooth-proper-lambda-cohomology-satisfies-its-shifted-analytic-rh}

\textbf{Statement.} Let \(X\) be smooth and proper over
\(\mathbb Z[1/M]\), of finite type, and equipped with a
\(\Lambda\)-structure in Borger's sense. Fix a prime \(\ell\), an
algebraic closure \(\overline{\mathbb Q}_\ell\), and an identification
of the algebraic roots of unity in it with complex roots of unity. Put
\[
 V^j=H^j_{\text{ét}}(X_{\overline{\mathbb Q}},\mathbb Q_\ell).
\] Then \(V^{2m+1}=0\), and there is a finite family of primitive
Dirichlet characters, counted with nonnegative integer multiplicities
\(r_{m,\chi}\), such that \[
 V^{2m}\otimes\overline{\mathbb Q}_\ell
 \simeq\bigoplus_\chi(\chi_\ell\otimes\overline{\mathbb Q}_\ell(-m))^{\oplus r_{m,\chi}}.
 \tag{CC-F01.1}
\] The convention for \(\chi_\ell\) is that its value on geometric
Frobenius at an unramified prime \(p\) is \(\chi(p)\). Define the
primitive Dirichlet completion of the cohomological factor by \[
 L^{\mathrm{Dir}}_{2m}(X,s)=\prod_\chi L(s-m,\chi)^{r_{m,\chi}}.
 \tag{CC-F01.2}
\] Every nontrivial zero of this factor lies on \(\Re s=m+\tfrac12\).
The product has a pole at \(s=m+1\) of order precisely \(r_{m,1}\) and
has no other poles. The corresponding completed factors have the usual
gamma and conductor factors of these primitive Dirichlet characters. At
all sufficiently good primes, (CC-F01.2) has exactly the cohomological
Euler factor of \(X\).

If \(S\) contains \(\ell\), the bad primes of \(X\), and every character
conductor prime, then the incomplete factor satisfies the exact identity
\[
 L^S_{2m}(X,s)
 =L^{\mathrm{Dir}}_{2m}(X,s)
   \prod_{p\in S}\prod_\chi
       (1-\chi(p)p^{m-s})^{r_{m,\chi}},
 \tag{CC-F01.3}
\] where a primitive character is assigned value zero at a prime
dividing its conductor. The explicitly deleted Euler factors can
introduce zeros on \(\Re s=m\). They are excluded from the
nontrivial-zero assertion. It would be incorrect to assert that every
zero of an incomplete Euler product lies on the shifted critical line.

\textbf{Proof.} Borger's Theorem 0.1(c) says that the Hodge
decomposition of \(H^j(X(\mathbb C),\mathbb C)\) has no components of
type \((a,b)\) with \(a\ne b\). Smoothness and properness give purity of
this Hodge structure of weight \(j\); its nonzero components therefore
satisfy \(a+b=j\). Consequently odd cohomology vanishes. Comparison with
étale cohomology proves \(V^{2m+1}=0\).

Borger's Corollary 6.2, equation (6.2.1), supplies a cyclotomic field
\(F=\mathbb Q(\zeta_N)\) over which the representation is a sum of Tate
twists, with multiplicities the diagonal Hodge numbers. In degree \(2m\)
the only possible Hodge type is \((m,m)\). Therefore \[
 V^{2m}|_{G_F}\simeq\mathbb Q_\ell(-m)^{b_{2m}}.
\] The twist \(W_m=V^{2m}(m)\) is thus trivial on \(G_F\) and factors
through the finite abelian group \(\mathrm{Gal}(F/\mathbb Q)\).
Characteristic-zero representations of a finite group are semisimple, by
averaging a projection over the group. Every irreducible representation
of a finite abelian group over an algebraically closed field is
one-dimensional: a common eigenspace of the commuting operators is
invariant, and irreducibility forces it to be the whole space. It
follows that \(W_m\) splits into finite-order characters. Characters of
\(\mathrm{Gal}(\mathbb Q(\zeta_N)/\mathbb Q)\) are Dirichlet characters;
replacing an imprimitive character by its primitive character changes
the modulus description, not the Galois character. This proves
(CC-F01.1).

On \(\mathbb Q_\ell(-m)\), geometric Frobenius has eigenvalue \(p^m\).
Hence the good-prime determinant is \[
 \det(1-\mathrm{Frob}_p p^{-s}\mid V^{2m})
   =\prod_\chi(1-\chi(p)p^{m-s})^{r_{m,\chi}}.
\] Its inverse is the Euler factor of (CC-F01.2). Euler products agree
in a right half-plane of absolute convergence, and their analytic
continuations consequently agree wherever defined. Removing the finitely
many factors in \(S\) gives (CC-F01.3).

By the adopted Dirichlet GRH, every nontrivial zero of \(L(s-m,\chi)\)
has real part \(m+1/2\). A finite product with nonnegative
multiplicities has no other nontrivial zeros. Only the primitive trivial
character contributes a pole, and it contributes \(\zeta(s-m)\). At
\(s=m+1\), all nonprincipal Dirichlet factors are nonzero by the
classical nonvanishing of \(L(1,\chi)\); thus the pole order is exactly
\(r_{m,1}\). Finally, if \(\chi(p)\ne0\), it is a root of unity, so
\(1-\chi(p)p^{m-s}=0\) implies \(p^{m-\Re s}=1\), or \(\Re s=m\). If
\(\chi(p)=0\), the factor is identically one. This proves every
assertion. \(\square\)

\textbf{Scope and significance.} This supplies analytic RH to an entire
geometric class identified inside the supplied \(\mathbf F_1\)
literature. The argument proves a decomposition of realizations and
their Dirichlet-completed \(L\)-functions. It does not identify a Chow
motive with a direct sum of Artin--Tate Chow motives, prove a Hodge or
Tate cycle map surjective, or construct the infinite-dimensional
cohomology envisaged in Borger's introduction, p.~3. Borger's Theorem
6.1 applies much more generally to separated finite-type
\(\Lambda\)-schemes; Corollary 6.2 gives their rational compactly
supported cohomology after cyclotomic restriction. For arbitrary mixed
objects one must separately control the actual extensions and bad-prime
inertia. CC-F01b proves that Borger's particular rational cohomology
class actually splits, so no such extension is being silently discarded
there; its asserted arithmetic-model comparison is restricted to
good-prime data.

\subsubsection{\texorpdfstring{CC-F01b. Actual rational splitting for
every separated finite-type
\(\Lambda\)-scheme}{CC-F01b. Actual rational splitting for every separated finite-type \textbackslash Lambda-scheme}}\label{cc-f01b.-actual-rational-splitting-for-every-separated-finite-type-lambda-scheme}

\textbf{Statement.} Let \(X\) be a separated \(\Lambda\)-scheme of
finite type over \(\mathbb Z\), let \(n\ge0\), and let \(\ell\) be
prime. For \[
 V=H^n_{\text{ét},c}(X_{\overline{\mathbb Q}},\mathbb Q_\ell),
\] there is an actual isomorphism, after extension of coefficient field,
\[
 V\otimes\overline{\mathbb Q}_\ell
 \simeq\bigoplus_{m,\chi}
 (\chi_\ell\otimes\overline{\mathbb Q}_\ell(-m))^{\oplus r_{m,n,\chi}}.
 \tag{CC-F01b.1}
\] The characters have finite abelian image, the multiplicities are
nonnegative integers, and only finitely many occur. In particular this
rational representation is semisimple. This assertion is about
cohomology with \(\mathbb Q_\ell\) coefficients: it does not assume that
integral or finite-coefficient cohomology is torsion-free, and asserts
neither integral splitting nor absence of torsion. The primitive
completion of its good-prime \(L\)-data is \[
 L_n^{\mathrm{Dir}}(X,s)=\prod_{m,\chi}L(s-m,\chi)^{r_{m,n,\chi}},
 \tag{CC-F01b.2}
\] whose nontrivial zeros lie on the lines \(\Re s=m+1/2\) for the
weights present. Deleted Euler factors are explicit finite products and
can add zeros on \(\Re s=m\). For a nonproper or singular model, no
equality with the zeta function of its bad fibers is asserted.

\textbf{Proof.} In the setup of Borger §6, \(X\) is exactly a separated
finite-type \(\Lambda\)-scheme over \(\mathbb Z\). Corollary 6.2, p.~20,
equation (6.2.1), gives an isomorphism of the actual representation \[
 V|_H\simeq\bigoplus_m\mathbb Q_\ell(-m)^{r_{m,n}},
 \qquad H=G_{\mathbb Q(\zeta_N)}.
\] The source does not insert a semisimplification superscript here. Let
\(\epsilon_\ell:G_\mathbb Q\to\mathbb Z_\ell^\times\) be the arithmetic
cyclotomic character, so that \(\mathbb Q_\ell(-m)\) carries
\(\epsilon_\ell^{-m}\). Its restriction to \(H\) has infinite image:
\(H\) has finite index in \(G_\mathbb Q\), while the image of the full
cyclotomic character is infinite, and the image of a finite-index
subgroup has finite index in that image. The powers
\(\epsilon_\ell^{-m}|_H\) are distinct for distinct integers \(m\),
since the torsion subgroup of \(\mathbb Z_\ell^\times\) is finite.

Define the corresponding isotypic subspace \[
 V_m=\{v\in V:\rho(h)v=\epsilon_\ell(h)^{-m}v\text{ for every }h\in H\}.
\] The given direct-sum decomposition gives \(V=\bigoplus_m V_m\). Since
\(H\) is normal and \(\epsilon_\ell\) is a character of the whole group,
for \(g\in G_\mathbb Q\), \(h\in H\), and \(v\in V_m\), \[
 \rho(h)\rho(g)v
 =\rho(g)\rho(g^{-1}hg)v
 =\epsilon_\ell(g^{-1}hg)^{-m}\rho(g)v
 =\epsilon_\ell(h)^{-m}\rho(g)v.
\] Thus every \(V_m\) is \(G_\mathbb Q\)-stable. The twist \(V_m(m)\) is
trivial on \(H\), hence factors through the finite abelian quotient
\(G_\mathbb Q/H\). By characteristic-zero semisimplicity and the
character decomposition of a finite abelian group, it splits over
\(\overline{\mathbb Q}_\ell\) into one-dimensional finite-order
characters. Undoing the twists proves (CC-F01b.1). The determinant
computation and primitive-character identification in CC-F01 apply to
every summand and give (CC-F01b.2) at good primes. The adopted Dirichlet
GRH and the same elementary computation of zeros of deleted Euler
factors prove the analytic assertions. No passage through a nonsplit
extension or an exactness claim for inertia invariants occurs.
\(\square\)

The nontrivial zero contributions in (CC-F01b.2) cannot be canceled by
poles of other factors: their real parts are half-integers, whereas the
only possible poles occur at the integers \(m+1\). Trivial zeros and
possible cancellation among trivial zeros and poles are separate from
this assertion. The importance of this refinement is that Borger's
actual rational splitting eliminates an extension obstruction inside
this specific geometric class. General mixed Artin--Tate objects outside
this class can still have nontrivial extensions. Borger's Remark 6.3
explicitly cautions against an integral-coefficient analogue.

\subsubsection{CC-F01a. The hexagonal CM curve does not itself furnish
Borger's finite-type
descent}\label{cc-f01a.-the-hexagonal-cm-curve-does-not-itself-furnish-borgers-finite-type-descent}

Borger (2009), §2.9, pp.~10--11, proves that a connected smooth proper
curve of genus at least one over \(\mathbb Q\), with a smooth proper
model over some \(\mathbb Z[1/M]\), has no \(\Lambda\)-structure. In
particular the genus-one CM curve with hexagonal period lattice in the
adopted note cannot be identified with a smooth proper
\(\Lambda\)-scheme over \(\mathbb Z[1/M]\). CM endomorphisms at split
primes do not supply the commuting Frobenius lifts at all good primes
required in Borger's definition. This refutes that particular
identification; it does not refute the note's use of a CM curve or its
separate reference to \(\mathbf F_1\).

\subsection{3. Classical consequences and exact hypothesis
map}\label{classical-consequences-and-exact-hypothesis-map}

\subsubsection{CC-C01. Lagarias--Odlyzko for abelian top
fields}\label{cc-c01.-lagariasodlyzko-for-abelian-top-fields}

For a finite Galois extension \(L/K\), let \(G=\mathrm{Gal}(L/K)\), let
\(C\) be a conjugacy class, let \(\pi_C(x,L/K)\) count unramified prime
ideals of \(K\) of norm at most \(x\) with Frobenius in \(C\), and let
\(n_L=[L:\mathbb Q]\), \(\Delta_L=|\mathrm{disc}(L)|\). Lagarias and
Odlyzko (1977), Theorem 1.1, equation (1.3), require \textbf{RH for
\(\zeta_L\)} and conclude, for \(x>2\), \[
 \left|\pi_C(x,L/K)-\frac{|C|}{|G|}\operatorname{Li}(x)\right|
 \le c_1\left\{\frac{|C|}{|G|}\sqrt{x}\log(\Delta_Lx^{n_L})+\log\Delta_L\right\},
\] with an absolute effectively computable \(c_1\). The inspected
witness is the full retypeset primary text, Theorem 1.1 on its p.~4; the
original article occupies pp.~409--464. Their Corollary 1.2 gives an
effective least-prime norm bound
\(c_2(\log\Delta_L)^2(\log\log\Delta_L)^4\); later sharpenings should
not be silently substituted for this displayed corollary.

If \(L/\mathbb Q\) is abelian, its Dedekind zeta function is a product
of primitive Dirichlet \(L\)-functions, including \(\zeta\); the adopted
statements therefore discharge the analytic hypothesis. Any intermediate
\(K\) is allowed. If only \(L/K\) is abelian, this factorization into
Dirichlet \(L\)-functions over \(\mathbb Q\) is unavailable in general.
For unrestricted \(L\), the missing input is RH for that top-field
Dedekind zeta function, not Artin holomorphy as an additional hypothesis
of Theorem 1.1.

\subsubsection{CC-C02. Least primes and finite character
tests}\label{cc-c02.-least-primes-and-finite-character-tests}

LLS, corrected arXiv v3 (2016), Corollary 1.2, p.~3, gives \[
 P(a,q)\le (\varphi(q)\log q)^2\qquad(q>3,\ (a,q)=1)
\] under Dirichlet GRH. The small moduli are finite explicit cases.
Theorem 1.1 gives stronger subgroup-escape bounds for \(q\ge3000\);
Corollary 1.1 gives least quadratic nonresidue \(<(\log q)^2\) for prime
\(q\ge5\). These need only Dirichlet GRH. The corrected v3 constants
must be used: Theorem 1.2 has \(\alpha(2)=0.8\), \(\alpha(3)=0.7\), and
\(\alpha(h)=0.66\) for \(h>3\). Earlier-version constants are not
interchangeable.

For a rank-one character, its values on group generators determine it.
For an arbitrary representation, the trace at the chosen generators need
not determine the character. The safe finite certificate tests every
residue class (or every class of the actual quotient), as in the
separate archive module.

The earlier explicit bound is also now checked against its original
source: Bach and Sorenson (1996), Theorem 5.3 (ERH), printed p.~1731,
gives \(P(a,q)<2(q\log q)^2\) for positive moduli \(q\ge2\) and
\((a,q)=1\). Their §2, p.~1720, specifies RH for \(\zeta_E\) and
\(\zeta_{\mathbb Q}\); in this application \(E=\mathbb Q(\zeta_q)\), so
the adopted input suffices. Their Corollary 3.5, pp.~1722--1723,
concerns a prime in every ideal class of a quadratic field and
explicitly sets \(E\) equal to its Hilbert class field. That application
has the wider top-field hypothesis discussed in CC-F03, even though its
base is quadratic. The stronger LLS bound is the one used for the finite
certificate.

\subsubsection{CC-C03. Miller's deterministic
algorithm}\label{cc-c03.-millers-deterministic-algorithm}

Miller (1976), Theorem 2, p.~301, proves under his ERH a deterministic
primality algorithm taking \(O(b^4\log\log b)\) Turing-machine steps for
an integer of binary length \(b\). His Appendix, pp.~316--317,
formulates the relevant ERH for Dirichlet \(L\)-functions. The adopted
Dirichlet GRH therefore discharges it. This validates that algorithm and
its complexity analysis; deterministic primality testing was already
unconditional by other algorithms, so one must not describe the
consequence as the first proof that primality belongs to P. The exact
bound on p.~301 was inspected visually as well as in extracted text.

\subsubsection{CC-C04. Bach's explicit witnesses and the class-group
boundary}\label{cc-c04.-bachs-explicit-witnesses-and-the-class-group-boundary}

Bach (1990), \emph{Explicit bounds for primality testing and related
problems}, has distinct applications. Its abstract gives a Miller or
Solovay--Strassen compositeness witness at most \(2(\log n)^2\), under
the relevant ERH. LLS §1.1, p.~1, attributes to Bach a prime outside a
proper subgroup below \(2(\log q)^2\), but that prime is allowed to
divide \(q\); imposing \(p\nmid q\) changes that cited bound to
\(3(\log q)^2\). This distinction must be preserved when applying a
generator theorem.

Belabas, Diaz y Diaz, and Friedman (2008), \emph{Small generators of the
ideal class group}, introduction pp.~2--3 of the inspected author
preprint, explicitly distinguish the two zero sums used by Bach: one for
a nontrivial class character and one for the Dedekind zeta function.
Their own Theorem 1.1, preprint p.~3, assumes RH for every nontrivial
ideal-class-character \(L\)-function. This supplies a directly inspected
primary locator for the hypothesis boundary.

The commonly cited class-group corollary is generation by prime ideals
of norm at most \(12(\log|\Delta_K|)^2\), under GRH for the relevant
ideal-class-character Hecke \(L\)-functions. This is identified as a
consequence of Bach's Theorem 4, p.~376, by subsequent primary
algorithmic literature. \textbf{Access/check status:} the original AMS
article was not successfully retrieved in this run (HTTP 403); its
theorem number is not represented here as directly checked against the
original. Its exact original wording is a bibliographical gate, while
the stronger operative LLS bounds and the mathematical hypothesis
boundary are directly checked. The following explicit example explains
why ``\(K/\mathbb Q\) abelian'' does not by itself discharge the
class-character hypothesis.

\subsubsection{CC-F03. An abelian base with a non-Dirichlet
class-character
factor}\label{cc-f03.-an-abelian-base-with-a-non-dirichlet-class-character-factor}

\textbf{Statement.} Put \(K=\mathbb Q(\sqrt{-23})\), and let \(H\) be
the splitting field of \(f(x)=x^3-x-1\). Then \(H/K\) is a nontrivial
unramified cyclic cubic extension. A nontrivial character of
\(\mathrm{Gal}(H/K)\) is an unramified class character of \(K\), but its
induction to \(G_\mathbb Q\) is an irreducible two-dimensional
representation. It is not a direct sum, or a virtual integer
combination, of abelian one-dimensional representations. Thus its Hecke
\(L\)-function is not supplied by an abelian-Artin factorization into
Dirichlet \(L\)-functions.

\textbf{Proof.} A reducible monic cubic in \(\mathbb Z[x]\) has an
integer root dividing its constant term. Neither \(1\) nor \(-1\) is a
root of \(f\), so \(f\) is irreducible. Its discriminant is
\(-4(-1)^3-27(-1)^2=-23\), a nonsquare. A transitive subgroup of \(S_3\)
is \(A_3\) or \(S_3\); the discriminant being a nonsquare forces
\(\mathrm{Gal}(H/\mathbb Q)=S_3\). The fixed field of \(A_3\) is \(K\),
since the product of the root differences squares to \(-23\).
Consequently \(\mathrm{Gal}(H/K)=A_3\cong C_3\).

For \(p\ne23\), the polynomial has separable reduction. Its roots lie in
an unramified extension of \(\mathbb Q_p\): choose a finite residue
extension splitting the reduction, lift it to the corresponding
unramified local extension, and lift each simple root by Hensel's lemma.
Hence the splitting field is unramified at all such primes.

At \(p=23\), direct multiplication gives \[
 x^3-x-1\equiv(x-10)^2(x-3)\pmod{23}.
\] The simple root \(3\) lifts to a root \(a\in\mathbb Z_{23}\), and
polynomial division writes \(f=(x-a)g\) with \(g\) monic quadratic over
\(\mathbb Z_{23}\). The resultant of \(x-a\) and \(g\) is a unit,
because \(3\) differs from \(10\) modulo 23. The discriminant product
formula \(\mathrm{disc}(f)=\mathrm{disc}(g)\mathrm{Res}(x-a,g)^2\)
therefore gives valuation one for \(\mathrm{disc}(g)\). The splitting
field of \(g\) is a ramified quadratic extension of \(\mathbb Q_{23}\);
the linear factor adds nothing. Its inertia group acts as a
transposition on the three roots. The relative inertia group for \(H/K\)
is its intersection with \(A_3\), which is trivial. Thus \(H/K\) is
unramified also above 23. There are no real places of \(K\), so there is
no archimedean ramification. Global class field theory identifies this
cyclic unramified extension with a quotient of the ideal class group.

Let \(\eta\) be a nontrivial complex character of \(A_3\). Its conjugate
by a transposition is \(\eta^{-1}\ne\eta\). The induced representation
has character values \(2,0,-1\) on the identity, a transposition, and a
three-cycle, respectively. Its character norm is
\((4+3\cdot0+2\cdot1)/6=1\), so it is irreducible. Every one-dimensional
character of \(S_3\) kills \(A_3=[S_3,S_3]\), and every virtual
combination of them has equal character values on the identity and on a
three-cycle. The induced character does not. More generally, every
abelian character of \(G_\mathbb Q\) is trivial on its commutator
subgroup; the image of this commutator subgroup in the quotient \(S_3\)
is \(A_3\). A virtual combination of abelian \(G_\mathbb Q\)-characters
likewise cannot equal this induced character. Artin induction of
\(L\)-functions identifies \(L_K(s,\eta)\) with the degree-two induced
Artin \(L\)-function, so the claimed abelian factorization is
unavailable. \(\square\)

This is a counterexample to the proposed \emph{factorization route}, not
a proof that Dirichlet GRH could never imply this Hecke GRH by a
different theorem. In this example the missing hypothesis is the
zero-line theorem for that class-character factor, or sufficiently
strong Dedekind GRH for the corresponding Hilbert class extension. One
need not identify \(H\) with the entire Hilbert class field for this
obstruction.

\subsubsection{CC-F03b. A sufficient condition that does discharge the
class-character
hypothesis}\label{cc-f03b.-a-sufficient-condition-that-does-discharge-the-class-character-hypothesis}

\textbf{Statement.} Let \(K/\mathbb Q\) be a number field whose Hilbert
class field \(H_K\) is abelian over \(\mathbb Q\). Then every
ideal-class-character Hecke \(L\)-function of \(K\) is a product of
primitive Dirichlet \(L\)-functions, and its GRH follows from the
adopted input. In particular this applies when \(K\) is imaginary
quadratic and its ideal class group has exponent dividing two.

\textbf{Proof.} Write \(G=\mathrm{Gal}(H_K/\mathbb Q)\),
\(U=\mathrm{Gal}(H_K/K)\), and let \(\eta\) be a character of \(U\). The
Hilbert class field correspondence identifies \(\eta\) with the chosen
ideal-class character. Because \(G\) is finite abelian,
\(\mathrm{Ind}_U^G\eta\) is a direct sum of one-dimensional characters
of \(G\). Induction compatibility of Artin \(L\)-functions therefore
writes \(L_K(s,\eta)\) as the product of the corresponding finite
abelian \(L\)-functions over \(\mathbb Q\), which are primitive
Dirichlet factors. This proves the first claim.

For imaginary quadratic \(K\), its Hilbert class field is stable under
complex conjugation, and conjugation acts on the class group by
inversion: an ideal \(\mathfrak a\) satisfies
\([\overline{\mathfrak a}]=[\mathfrak a]^{-1}\), since
\(\mathfrak a\overline{\mathfrak a}\) is generated by the ideal norm.
The exact sequence for \(\mathrm{Gal}(H_K/\mathbb Q)\) splits by complex
conjugation, which has order two and restricts nontrivially to \(K\).
Thus this group is the semidirect product of the class group by the
inversion action. If the class group has exponent dividing two,
inversion is trivial and the semidirect product is abelian. The first
part applies. \(\square\)

\subsubsection{CC-C05. Littlewood's
consequences}\label{cc-c05.-littlewoods-consequences}

Titchmarsh, revised by Heath-Brown (1986), Theorem 14.2 and equations
(14.2.5)--(14.2.6), pp.~336--337, gives the Littlewood bounds leading
from RH to the Lindelöf hypothesis. Theorem 14.25(A)--(C), pp.~369--370,
gives convergence of \(\sum\mu(n)n^{-s}=1/\zeta(s)\) for \(\Re s>1/2\),
and the equivalence \[
 \mathrm{RH}\quad\Longleftrightarrow\quad
 M(x)=O_\varepsilon(x^{1/2+\varepsilon})
 \quad\text{for every }\varepsilon>0.
\] These are now available with the RH hypothesis discharged. This does
not yield the much stronger inequality \(|M(x)|\le\sqrt{x}\).

For the edge of the Dirichlet critical strip, the exact effective form
used here is LLS Theorem 1.5. Write \[
 B_1(q)=\log\log q-\log2+\tfrac12+\frac1{\log\log q}
       +\frac{14\log\log q}{\log q}.
\] For a primitive character of conductor \(q\ge10^{10}\), \[
 |L(1,\chi)|\ge\frac{\pi^2}{12e^\gamma B_1(q)}.
\] LLS Corollary 1.3 then gives, for negative fundamental discriminant
\(-q\), \[
 h(-q)\ge\frac{\pi}{12e^\gamma}\frac{\sqrt q}{B_1(q)}.
\] These require only Dirichlet GRH. Bounds for Hecke or higher-rank
automorphic \(L\)-functions require their corresponding analytic
hypotheses.

\subsubsection{CC-F02. A fully effective analytic Siegel
bound}\label{cc-f02.-a-fully-effective-analytic-siegel-bound}

\textbf{Statement.} For every positive rational \(\varepsilon\), there
is an algorithm producing a positive rational \(c_\varepsilon\) such
that every primitive real nonprincipal Dirichlet character \(\chi\) of
conductor \(q\) satisfies \[
 L(1,\chi)\ge c_\varepsilon q^{-\varepsilon}.
\]

\textbf{Proof.} For \(q\ge10^{10}\), put \(t=\log q\) and \(u=\log t\).
Then \(t>23\) and \(u>3\). Since \(-\log2+1/2<0\), \(1/u<1/3\), and
\(14/t<1\), \[
 B_1(q)\le2u+\tfrac13<3u.
\] The displayed LLS bound yields \[
 L(1,\chi)\ge \frac{\pi^2}{36e^\gamma\log\log q}.
\] The value is positive: the real Euler product is positive on \(s>1\),
continuity and nonvanishing at \(s=1\) give positivity at 1. Choose a
computable integer \(Q\ge10^{10}\) with
\(\log Q\ge\max(2,2/\varepsilon^2)\). For \(q\ge Q\), the elementary
estimate \(e^v\ge v^2/2\) gives \[
 q^\varepsilon=e^{\varepsilon t}
 \ge\frac{\varepsilon^2t^2}{2}\ge t\ge\log t=\log\log q.
\] Consequently \(q^\varepsilon L(1,\chi)\ge\pi^2/(36e^\gamma)\)
throughout that range.

Only finitely many conductors remain. Primitive real characters of
conductor below \(Q\) can be enumerated by finite arithmetic. For each
such nonprincipal \(\chi\), periodicity and zero sum on a period imply
\(A(x)=\sum_{n\le x}\chi(n)\) has \(|A(x)|\le q\). Abel summation gives
the certified tail estimate \[
 \left|L(1,\chi)-\sum_{n\le N}\frac{\chi(n)}n\right|
 \le \frac{2q}{N}.
\] Increase \(N\) until the rational number
\(b_\chi=\sum_{n\le N}\chi(n)/n-2q/N\) is positive. This terminates
because the partial sums converge to the positive number \(L(1,\chi)\).
Since \(q^\varepsilon\ge1\), one may choose a positive rational
\(c_\varepsilon\) below both \(\pi^2/(36e^\gamma)\) and every resulting
\(b_\chi\). The finite minimum and the computable constant permit such a
choice. For \(q<Q\),
\(L(1,\chi)\ge b_\chi\ge c_\varepsilon\ge c_\varepsilon q^{-\varepsilon}\);
for \(q\ge Q\), the preceding tail-range estimate applies. This
algorithm establishes effectivity, without claiming practical
efficiency. \(\square\)

\subsubsection{CC-C06. Cramér's RH theorem, as distinct from his
stronger
conjecture}\label{cc-c06.-cramuxe9rs-rh-theorem-as-distinct-from-his-stronger-conjecture}

Cramér's classical RH consequence is
\(p_{n+1}-p_n=O(\sqrt{p_n}\log p_n)\). An inspected primary modern
source is Dudek (2014), \emph{On the Riemann hypothesis and the
difference between primes}, arXiv:1402.6417v2: Theorem 1, p.~2, states
Cramér's result; Theorem 2, p.~2, proves the explicit version \[
 (x-\tfrac4\pi\sqrt x\log x,x]\quad\text{contains a prime for every }x\ge2;
\] Theorem 3 replaces \(4/\pi\) by \(1+\varepsilon\) for sufficiently
large \(x\). RH is the sole unproved analytic hypothesis. Cramér's
conjectural \(O((\log p_n)^2)\) gap bound is not the theorem being
discharged. The original 1920 theorem number was not directly inspected;
the exact operative locators here are Dudek's numbered statements, with
the original work retained in the bibliography.

\subsubsection{CC-C07. Selberg's mean square for short
intervals}\label{cc-c07.-selbergs-mean-square-for-short-intervals}

The usable RH consequence is \[
 \int_1^X|\psi(t+h)-\psi(t)-h|^2\,dt
 \ll Xh(\log X)^2\qquad(1\le h\le X).
\] An inspected modern primary statement is Bhowmik--Schlage-Puchta,
\emph{Mean representation number of integers as the sum of primes},
Lemma 1, which explicitly attributes the input to Selberg (1943),
equation (13). The original Selberg article is \emph{On the normal
density of primes in small intervals, and the difference between
consecutive primes}, pp.~87--105; it is paper 12 in \emph{Collected
papers}, vol.~I. The original equation was not directly inspected in
this run, so that locator is marked as transmitted by the later primary
paper.

The elementary consequence needed here requires no further zero
hypothesis. Apply the mean square with upper endpoint \(2X\). For fixed
\(\delta>0\), the measure of \(t\in[X,2X]\) for which
\(|\psi(t+h)-\psi(t)-h|>\delta h\) is at most
\(O(X(\log X)^2/(\delta^2h))\), by integrating over that set. Thus if
\(h/(\log X)^2\to\infty\), the relative exceptional measure tends to
zero. This is an almost-all-interval statement, not a bound in every
interval and not the Selberg eigenvalue conjecture.

\subsubsection{CC-C08. Montgomery's restricted pair
correlation}\label{cc-c08.-montgomerys-restricted-pair-correlation}

Montgomery (1973), §1, unnumbered Theorem, p.~181, equations (1)--(2),
defines \[
 F(\alpha,T)=\left(\frac{T}{2\pi}\log T\right)^{-1}
 \sum_{0<\gamma,\gamma'\le T}
 T^{i\alpha(\gamma-\gamma')}
 \frac4{4+(\gamma-\gamma')^2}.
\] Under RH, \[
 F(\alpha,T)=(1+o(1))T^{-2\alpha}\log T+\alpha+o(1)
\] uniformly on \(0\le\alpha\le1-\varepsilon\), for each fixed
\(\varepsilon>0\). RH suffices for this theorem. The unrestricted
pair-correlation prediction, discussed separately in the same paper at
equations (11)--(12), is not obtained by dropping a hypothesis from the
displayed theorem: its Fourier-support obstruction remains. RH
determines the horizontal locations, not all vertical correlations or
simplicity.

\subsection{4. Criteria and programmes}\label{criteria-and-programmes}

\begin{longtable}[]{@{}
  >{\raggedright\arraybackslash}p{(\columnwidth - 6\tabcolsep) * \real{0.2500}}
  >{\raggedright\arraybackslash}p{(\columnwidth - 6\tabcolsep) * \real{0.2500}}
  >{\raggedright\arraybackslash}p{(\columnwidth - 6\tabcolsep) * \real{0.2500}}
  >{\raggedright\arraybackslash}p{(\columnwidth - 6\tabcolsep) * \real{0.2500}}@{}}
\toprule\noalign{}
\begin{minipage}[b]{\linewidth}\raggedright
ID
\end{minipage} & \begin{minipage}[b]{\linewidth}\raggedright
Precise conclusion now available
\end{minipage} & \begin{minipage}[b]{\linewidth}\raggedright
Exact locator
\end{minipage} & \begin{minipage}[b]{\linewidth}\raggedright
Remaining mathematical task
\end{minipage} \\
\midrule\noalign{}
\endhead
\bottomrule\noalign{}
\endlastfoot
CC-P01 & The global Weil quadratic form is nonnegative on its specified
test-function space & Connes--Consani, supplied source, Appendix C,
Proposition C.1, p.~51; explicit formula Appendix B & The local
archimedean or semilocal comparison with a particular cutoff/Sonin
projection is an additional identity or inequality, not an automatic
termwise consequence of the global sum \\
CC-P02 & Every classical Li coefficient is nonnegative; indeed CC-F04
proves strict positivity for the nonzero Dirichlet systems & Li (1997),
Theorem 1, p.~325; generalized criterion Lagarias (2007), Theorem 2.4,
p.~1702 (preprint Theorem 2.2, p.~11) & Higher-rank automorphic
positivity requires RH for the corresponding \(L\)-function \\
CC-P03 & \(\mathbf1_{(0,1]}\) belongs to the closed span of
\(x\mapsto\{1/(nx)\}\), \(n\ge1\), in \(L^2(0,\infty)\) & Báez-Duarte
(2002), Theorem 1.1, p.~1 & Convergence of a particular proposed
coefficient sequence or optimal distance asymptotics needs its own
theorem; the ordinary Möbius truncations diverge, as discussed at (1.1),
p.~2 \\
CC-P04 & There is an abstract self-adjoint operator whose spectral
ordinates are the zeros, counting multiplicity & Elementary construction
CC-F05 below & Arithmetic origin, a prescribed trace formula, canonical
domain, and geometric functoriality do not follow from the diagonal
construction \\
CC-P05 & Connes's RH-equivalent trace assertion for the \(\mathbb Q\)
version has its zero hypothesis discharged & Connes (1999), §VIII,
Theorem 5 in positive characteristic, and the explicit \(\mathbb Q\)
analogue on preprint pp.~44--47 & The literally numbered theorem
concerns positive characteristic. General number-field Hecke
\(L\)-functions exceed the adopted Dirichlet statement \\
CC-P06 & Deninger's proposed eigenvalue location is correct for the
adopted abelian Artin--Tate factors & Deninger, \emph{The Hilbert--Pólya
strategy and height pairings}, §2, Conjectures 1--2, pp.~2--5 & The
functorial cohomology, determinant formula, Hodge star, and positive
height pairing remain construction/positivity questions; CC-F06 refutes
a general converse from eigenvalue location to the adjoint identity \\
CC-P07 & The specific sufficient de Branges conditions tested by
Conrey--Li remain refuted & Conrey--Li (2000), Theorems 1--2;
§§3.1--3.2, equations (3.1)--(3.4); §4 Remark & Any successful revised
Hilbert-space construction must avoid those disproved positivity
assumptions \\
CC-P08 & The \(\Lambda\)-geometric class in CC-F01 has its analytic RH &
Borger Theorem 0.1 and Corollary 6.2 & Infinite-dimensional
\(\mathbf F_1\) cohomology, an archimedean theory, and a geometric proof
of RH are separate parts of the programme \\
CC-P09 & The already proved restricted Katz--Sarnak density results can
now be interpreted with real ordinates in the Dirichlet families &
Rubinstein (2001), Theorem 3.1, pp.~150--151 & Full support and
higher-rank families do not follow. Theorem 3.1 itself was already
unconditional, allowing complex ordinates \\
\end{longtable}

\subsubsection{CC-P01. Exact global positivity and the archimedean
boundary}\label{cc-p01.-exact-global-positivity-and-the-archimedean-boundary}

For \(g\in C_c^\infty(\mathbb R_{>0})\), use
\(\widetilde g(s)=\int_0^\infty g(x)x^{s-1}dx\) and impose
\(\widetilde g(0)=\widetilde g(1)=0\). On RH the zero-side form is \[
 Q(g)=\sum_\rho\widetilde g(\rho)
       \overline{\widetilde g(1-\overline\rho)}
     =\sum_\rho|\widetilde g(\rho)|^2\ge0.
\] The sums converge because the Mellin transform decays rapidly on a
vertical strip and the zero count is \(O(T\log T)\). Proposition C.1 of
the supplied Connes--Consani paper expresses the equivalent
prime-and-archimedean inequality with their sign convention as
\(\sum_v W_v(g*\overline g^{\sharp})\le0\); it allows finitely many
additional Mellin vanishing conditions as well. The global inequality is
now a theorem in the adopted theory. Their Theorem 1 and Theorem 6.11
concern a more specific comparison of the archimedean term with a
Sonin-space trace, with support and vanishing restrictions. Those
results remain their existing proved theorems; RH alone does not justify
enlarging their support class or declaring each separate local term
positive.

\subsubsection{CC-F04. Strict real-part positivity of Dirichlet Li
coefficients}\label{cc-f04.-strict-real-part-positivity-of-dirichlet-li-coefficients}

\textbf{Statement.} For the zeros of \(\zeta\), or of any primitive
Dirichlet \(L\)-function after the usual completion and removal of the
trivial-character poles, define the Li coefficients by the standard
symmetrically ordered zero sum \[
 \lambda_n=\sum_\rho\left[1-\left(1-\frac1\rho\right)^n\right],\qquad n\ge1.
\] Under the adopted statements, \(\Re\lambda_n>0\) for every \(n\ge1\).
For a real character, or for \(\zeta\), \(\lambda_n\) itself is real and
strictly positive. The criterion is the known theorem of Li and its
generalized form; the short proof here records the strictness and the
treatment of complex characters.

\textbf{Proof.} Every zero has the form \(\rho=1/2+i\gamma\), and \[
 z_\rho=1-\frac1\rho=\frac{-1/2+i\gamma}{1/2+i\gamma}
\] has modulus one. Write \(z_\rho=e^{i\theta_\rho}\). Each real-part
summand is \(1-\cos(n\theta_\rho)\ge0\). For \(|\gamma|\to\infty\),
\(\theta_\rho=O(1/|\gamma|)\); hence these real parts are
\(O_n(1/\gamma^2)\). The usual zero count \(O(T\log(q(T+2)))\), with
\(q\) fixed, makes the real-part series absolutely convergent, so taking
real parts of the prescribed symmetric limit gives the sum of these
nonnegative terms.

There are infinitely many zeros with arbitrarily large \(|\gamma|\). For
finite \(\gamma\), \(z_\rho\ne1\), while \(z_\rho\to1\) as
\(|\gamma|\to\infty\). For a fixed \(n\), choose a sufficiently large
zero so that the nonzero principal argument satisfies
\(0<|\theta_\rho|<2\pi/n\). Then \(1-\cos(n\theta_\rho)>0\), so the
entire sum is strictly positive. In the real-character and zeta cases,
conjugate zeros occur with equal multiplicities; their contributions are
conjugates, making the standard sum real. \(\square\)

No assertion of reality is made for the coefficient attached to a
complex Dirichlet character. The generalized Li criterion uses its real
part (Lagarias, published Theorem 2.4; preprint Theorem 2.2).

\subsubsection{CC-F05. The bare Hilbert--Pólya existence
statement}\label{cc-f05.-the-bare-hilbertpuxf3lya-existence-statement}

\textbf{Statement and proof.} List the nontrivial zeros of \(\zeta\),
with multiplicity, as \(1/2+i\gamma_j\), and choose an enumeration with
\(|\gamma_j|\to\infty\). On \(\mathcal H=\ell^2(\mathbb N)\), put \[
 D(a_j)=(\gamma_j a_j),\qquad
 \operatorname{Dom}(D)=\{a\in\ell^2:\sum_j\gamma_j^2|a_j|^2<\infty\}.
\] Finite sequences give a dense domain. Since the \(\gamma_j\) are
real, \(D\) is symmetric. If \(b\in\operatorname{Dom}(D^*)\), testing
the adjoint relation against the coordinate vectors forces
\((D^*b)_j=\gamma_jb_j\); membership of \(D^*b\) in \(\ell^2\) then
forces \(b\in\operatorname{Dom}(D)\). Conversely every such \(b\)
satisfies the adjoint relation by Cauchy--Schwarz. Thus \(D^*=D\). The
resolvent \((D-i)^{-1}\) is diagonal with entries
\((\gamma_j-i)^{-1}\to0\), and its truncations are finite-rank operators
converging in operator norm. Therefore it is compact. Each coordinate
vector is an eigenvector with the prescribed ordinate, with
multiplicities retained. If a complex number \(z\) is not an ordinate,
discreteness and escape to infinity give \(\inf_j|\gamma_j-z|>0\), so
the diagonal operator \(R_z\) with entries \((\gamma_j-z)^{-1}\) is
bounded. The coordinate identity \(DR_z=I+zR_z\) shows that its range is
in \(\operatorname{Dom}(D)\), bounded in the graph norm; also
\((D-z)R_z=I\) and \(R_z(D-z)=I\) on their respective domains. Hence
there is no additional spectrum. This proves the abstract existence
statement. \(\square\)

This construction uses the zeros themselves as input. It supplies none
of the independent arithmetic explanation sought in a stronger
Hilbert--Pólya programme. The same construction works for each primitive
Dirichlet function. No zero-simplicity assumption is required.

\subsubsection{CC-P05. Connes: exact theorem
scope}\label{cc-p05.-connes-exact-theorem-scope}

Connes's \emph{Trace formula in noncommutative geometry and the zeros of
the Riemann zeta function} has several assertions that should not be
merged. Its §III Theorem 1, preprint p.~13, gives a spectral realization
of critical zeros (with the multiplicity cutoff stated there), with
noncritical zeros entering a resonance description. The adopted RH
removes the noncritical alternatives for \(\zeta\) and the relevant
Dirichlet twists. In §VIII, Theorem 5 is explicitly formulated for
global fields of positive characteristic. The text then constructs and
discusses the corresponding \(\mathbb Q\) cutoff using the real place
and prolate spaces on preprint pp.~44--47, ending with the stated
analogues of Lemma 1, Theorem 5 and Lemma 3. For \(\mathbb Q\), unitary
Hecke characters are Dirichlet characters times imaginary norm powers;
these shift ordinates without changing the real zero line. Thus the
analytic hypothesis in that \(\mathbb Q\) criterion is covered. The
statement for general number-field Hecke characters requires additional
GRH outside the adopted family. A proof of an equivalent trace assertion
does not, by itself, provide a new geometric derivation of its required
cutoff limit.

\subsubsection{CC-F06. Eigenvalue purity does not imply Deninger's
positive adjoint
structure}\label{cc-f06.-eigenvalue-purity-does-not-imply-deningers-positive-adjoint-structure}

\textbf{Statement.} For every real \(w\), there is a complex linear
endomorphism with spectrum on \(\Re s=w/2\) that admits no
positive-definite Hermitian form for which \(A+A^*=wI\).

\textbf{Proof.} On \(\mathbb C^2\), put \(A=(w/2)I+N\), where
\(Ne_2=e_1\) and \(Ne_1=0\). Its only eigenvalue is \(w/2\), and
\(N^2=0\ne N\). If a positive Hermitian form made \(A+A^*=wI\), then
\(N^*=-N\). Taking \(x=e_2\) gives \[
 \|Nx\|^2=\langle x,N^*Nx\rangle
          =-\langle x,N^2x\rangle=0,
\] contradicting positive definiteness and \(Nx=e_1\ne0\). \(\square\)

Deninger's proposed identity in §2, pp.~4--5, is
\(\nu(f,f')=(\theta f,f')+(f,\theta f')\), obtained from a trace, cup
product and Hodge star. It implies the desired line. The example proves
that a general converse from line location to that positive adjoint
identity is false. It does not refute the existence of the particular
geometric objects Deninger conjectures. Semisimplicity, a pairing and
the functorial geometry are additional content.

\subsubsection{CC-P07. What the de Branges counterexamples
settle}\label{cc-p07.-what-the-de-branges-counterexamples-settle}

Conrey and Li's Theorems 1--2 formulate sufficient positivity conditions
in de Branges spaces. In §3.1 they show that the proposed conditions
fail for \(E(z)=\xi(1-iz)\), and in §3.2 they give analogous failures
for the character modulo 4. Their §4 Remark attributes to Sarnak a
nonnumerical proof of failure of the kernel-shift condition (3.3), via
the value distribution of \(\log\zeta\), with a Dirichlet analogue for
(3.8). The derivative-test counterexamples to (3.1) in §3 are cited as
the published numerical results, not as computations independently
certified in this run. The nonnumerical kernel-shift refutation already
suffices to disprove the blanket assertion that every proposed
sufficient positivity condition follows from RH. The accurate outcome is
a refutation of those particular sufficient conditions. It would be a
logical error to deduce them from RH.

\subsubsection{CC-P09. Katz--Sarnak and the restriction that
survives}\label{cc-p09.-katzsarnak-and-the-restriction-that-survives}

Rubinstein's Theorem 3.1, pp.~150--151, treats primitive quadratic
characters \(\chi_d\), averaged over fundamental discriminants with
\(X/2\le|d|<X\). For product Schwartz tests
\(f(x_1,\ldots,x_n)=\prod_i f_i(x_i)\), each \(f_i\) even, with Fourier
support \(\sum_i|u_i|<1\), the normalized \(n\)-level density is the
symplectic determinant with kernel \[
 K_{-1}(x,y)=\frac{\sin\pi(x-y)}{\pi(x-y)}
            -\frac{\sin\pi(x+y)}{\pi(x+y)}.
\] The sentence immediately preceding the theorem explicitly allows
complex ordinates and says RH is not assumed. Thus the theorem itself
was already unconditional. Dirichlet GRH now supplies real ordinates but
does not remove its Fourier-support restriction. Full number-field
family statistics, and analogues involving higher-rank \(L\)-functions,
are not consequences of just the adopted zero-line statements. The
finite-field monodromy and equidistribution part of Katz--Sarnak
likewise was not waiting for number-field RH.

\subsection{5. Explicit boundary
catalogue}\label{explicit-boundary-catalogue}

\begin{longtable}[]{@{}
  >{\raggedright\arraybackslash}p{(\columnwidth - 2\tabcolsep) * \real{0.5000}}
  >{\raggedright\arraybackslash}p{(\columnwidth - 2\tabcolsep) * \real{0.5000}}@{}}
\toprule\noalign{}
\begin{minipage}[b]{\linewidth}\raggedright
Proposed extension
\end{minipage} & \begin{minipage}[b]{\linewidth}\raggedright
Status and actual additional gate
\end{minipage} \\
\midrule\noalign{}
\endhead
\bottomrule\noalign{}
\endlastfoot
General effective Chebotarev from Dirichlet GRH & Theorem 1.1 of
Lagarias--Odlyzko needs RH for \(\zeta_L\). It is discharged when
\(L/\mathbb Q\) is abelian; arbitrary top fields need more \\
Bach class-group generators for every abelian \(K/\mathbb Q\) & The
class-character Hecke \(L\)-functions need their own zero theorem, or a
suitable Hilbert-class-extension Dedekind GRH; CC-F03 shows why abelian
base is insufficient for the offered factorization argument \\
Full generalized Li positivity & Lagarias Theorem 2.4 (preprint 2.2) is
family by family. General unitary cuspidal \(GL_N\) functions require
their own RH; Dirichlet GRH covers the finite-order \(GL_1/\mathbb Q\)
family and its norm shifts \\
Full GUE / Katz--Sarnak statistics / simplicity of all zeros & The cited
restricted statistics keep their test-function restrictions. No theorem
here upgrades zero location to arbitrary correlations or simplicity \\
Cramér's logarithmic-square gap conjecture & Distinct from Cramér's
RH-conditional square-root-times-log theorem \\
Universal local positivity in an explicit formula & Positivity of a sum
of local terms does not imply positivity of each summand.
Connes--Consani's extra local comparison has its own hypotheses \\
Deninger's positive cohomology or a canonical Hilbert--Pólya operator &
The analytic target is supplied for the adopted factors; functorial
spaces, adjoints, pairings and trace compatibility remain. CC-F06
disproves an unrestricted linear-algebra converse \\
A \(\Lambda\)-structure on the hexagonal elliptic curve & Refuted by
Borger §2.9; split-prime CM Frobenius lifts are weaker data \\
A Chow-motive Artin--Tate decomposition from the cohomological
\(L\)-factorization & Not proved by the factorization; a realization is
not automatically fully faithful or conservative in the required motivic
sense \\
Effective integral points merely from the phrase ``effective Siegel'' &
No such implication is asserted. CC-F02 proves only the analytic lower
bound for real Dirichlet \(L(1,\chi)\) and its class-number
consequences \\
\end{longtable}

\subsection{6. Source register and exact-reading
status}\label{source-register-and-exact-reading-status}

The following are primary works or mirrors of primary texts. Web
references in the companion TSV are provenance aids; the paper should
cite authors, dates and the locators above. No original statement number
is guessed when an original was not inspected.

\begin{enumerate}
\def\labelenumi{\arabic{enumi}.}
\item
  \textbf{Borger, J. (2009).} \emph{\(\Lambda\)-rings and the field with
  one element} (arXiv:0906.3146v1). Theorems 0.1, 5.3, 6.1; Corollary
  6.2, equation (6.2.1); Corollaries 6.5--6.6; §2.9.
  \url{https://arxiv.org/pdf/0906.3146} . Supplied TeX and external PDF both
  inspected. The PDF's regenerated date is not used as the publication
  date.
\item
  \textbf{Lamzouri, Y., Li, X., \& Soundararajan, K. (2015; corrected
  preprint 2016).} Conditional bounds for the least quadratic
  non-residue and related problems. \emph{Mathematics of Computation,
  84}(295), 2391--2412. \url{https://doi.org/10.1090/S0025-5718-2015-02925-1}
  . Operative corrected text: \url{https://arxiv.org/pdf/1309.3595} , v3, 8
  December 2016, internal date 12 December 2016. Theorems 1.1, 1.4, 1.5;
  Corollaries 1.1--1.3. PDF and extracted text inspected.
\item
  \textbf{Lagarias, J. C., \& Odlyzko, A. M. (1977).} Effective versions
  of the Chebotarev density theorem. In A. Fröhlich (Ed.),
  \emph{Algebraic number fields: L-functions and Galois properties}
  (pp.~409--464). Academic Press. Theorem 1.1, equation (1.3); Corollary
  1.2. Inspected retypeset primary text:
  \url{https://aareyanmanzoor.github.io/assets/articles/lagarias-odlyzko.pdf}
  .
\item
  \textbf{Miller, G. L. (1976).} Riemann's hypothesis and tests for
  primality. \emph{Journal of Computer and System Sciences, 13}(3),
  300--317. Theorem 2, p.~301; Appendix, pp.~316--317.
  \url{https://www.cs.cmu.edu/~glmiller/Publications/Papers/Mi76.pdf}
  . Original PDF inspected, including rendered p.~301.
\item
  \textbf{Bach, E. (1990).} Explicit bounds for primality testing and
  related problems. \emph{Mathematics of Computation, 55}(191),
  355--380. \url{https://doi.org/10.1090/S0025-5718-1990-1023756-8} . Original
  PDF access failed; author abstract and LLS attribution inspected.
  Theorem 4 / p.~376 class-group locator remains to be directly checked.
\item
  \textbf{Titchmarsh, E. C. (1986).} \emph{The theory of the Riemann
  zeta-function} (2nd ed., D. R. Heath-Brown, rev.). Oxford University
  Press. Theorems 14.2 and 14.25(A)--(C). Original book scan:
  \url{https://sites.math.rutgers.edu/~zeilberg/EM18/TitchmarshZeta.pdf}
  . Inspected extracted text; equations crosschecked against theorem
  context.
\item
  \textbf{Cramér, H. (1920).} Some theorems concerning prime numbers.
  \emph{Arkiv för Matematik, Astronomi och Fysik, 15}(5), 1--33.
  Historical attribution; original statement number not inspected.
  Operative source follows.
\item
  \textbf{Dudek, A. W. (2014/2015).} On the Riemann hypothesis and the
  difference between primes. Preprint arXiv:1402.6417v2; published in
  \emph{International Journal of Number Theory, 11}(3), 771--778.
  Theorems 1--3 in the inspected preprint (published numbering
  1.1--1.3). \url{https://arxiv.org/pdf/1402.6417} .
\item
  \textbf{Selberg, A. (1943).} On the normal density of primes in small
  intervals, and the difference between consecutive primes. \emph{Archiv
  for Mathematik og Naturvidenskab, 47}(6), 87--105. Equation (13), as
  cited by the next primary work; original not directly inspected.
\item
  \textbf{Bhowmik, G., \& Schlage-Puchta, J.-C. (2010).} Mean
  representation number of integers as the sum of primes. \emph{Nagoya
  Mathematical Journal, 200}, 27--33.
  \url{https://doi.org/10.1215/00277630-2010-010} . Lemma 1, p.~2 in arXiv v3,
  with explicit attribution to Selberg equation (13). Primary PDF
  inspected: \url{https://arxiv.org/pdf/0806.3295} ; initial mirror:
  \url{https://paperzz.com/doc/9064999/mean-representation-number-of-integers-as-the-sum}
  .
\item
  \textbf{Montgomery, H. L. (1973).} The pair correlation of zeros of
  the zeta function. In \emph{Analytic number theory} (Proceedings of
  Symposia in Pure Mathematics, Vol. 24, pp.~181--193). American
  Mathematical Society. \url{https://doi.org/10.1090/pspum/024/9944} . §1
  unnumbered Theorem, p.~181, equations (1)--(2); conjectural equations
  (11)--(12). Primary-text mirror inspected:
  \url{https://studylib.net/doc/7952799/the-pair-correlation-of-zeros-of-the-zeta-function}
  .
\item
  \textbf{Li, X.-J. (1997).} The positivity of a sequence of numbers and
  the Riemann hypothesis. \emph{Journal of Number Theory, 65}(2),
  325--333. \url{https://doi.org/10.1006/jnth.1997.2137} . Theorem 1, p.~325;
  equation (1.4). Primary-text mirror:
  \url{https://paperzz.com/doc/7086946/--science-direct} .
\item
  \textbf{Lagarias, J. C. (2007).} Li coefficients for automorphic
  \(L\)-functions. \emph{Annales de l'Institut Fourier, 57}(5),
  1689--1740. \url{https://doi.org/10.5802/aif.2311} . Published Theorem 2.4,
  p.~1702:
  \url{https://www.numdam.org/article/AIF_2007__57_5_1689_0.pdf} . In
  arXiv v4 (2005), \url{https://arxiv.org/pdf/math/0404394} , the
  corresponding statement is Theorem 2.2, p.~11. Both primary statements
  inspected; their different numbering is retained.
\item
  \textbf{Báez-Duarte, L. (2002/2003).} A strengthening of the
  Nyman--Beurling criterion for the Riemann hypothesis. \emph{Rendiconti
  Lincei. Matematica e Applicazioni, 14}, 5--11. Theorem 1.1 in
  arXiv:math/0202141v2: \url{https://arxiv.org/pdf/math/0202141} . PDF
  inspected. The quantitative continuation is arXiv:math/0205003; not
  used here without checking its exact numbered statement.
\item
  \textbf{Connes, A. (1999).} Trace formula in noncommutative geometry
  and the zeros of the Riemann zeta function. \emph{Selecta Mathematica,
  5}, 29--106. \url{https://arxiv.org/pdf/math/9811068} . §VIII, Theorem 5,
  and preprint pp.~44--47. PDF inspected; characteristic scope checked
  against the theorem's opening words.
\item
  \textbf{Connes, A., \& Consani, C. (2020).} Weil positivity and trace
  formula, the archimedean place. arXiv:2006.13771.
  \url{https://arxiv.org/pdf/2006.13771} . Supplied TeX
  \nolinkurl{weil-compo.tex}; Appendix C Proposition C.1; Theorem 1;
  Theorem 6.11; Appendix B. Supplied text and PDF locators inspected.
\item
  \textbf{Deninger, C. (2010).} The Hilbert--Pólya strategy and height
  pairings. In \emph{Casimir force, Casimir operators and the Riemann
  hypothesis} (pp.~275--283). Walter de Gruyter. §2, Conjectures 1--2
  and the adjoint calculation, preprint pp.~2--5. Author/institute PDF:
  \url{https://www.uni-muenster.de/SFB878/publications/files/php7aMxKR3029.pdf}
  ; arXiv: \url{https://arxiv.org/pdf/1001.1621} . PDF inspected; publication
  metadata checked on the author's publications page.
\item
  \textbf{Conrey, J. B., \& Li, X.-J. (2000).} A note on some positivity
  conditions related to zeta- and \(L\)-functions. \emph{International
  Mathematics Research Notices, 2000}(18), 929--940.
  \url{https://arxiv.org/pdf/math/9812166} . Theorems 1--2; §§3.1--3.2; §4
  Remark. Primary PDF inspected.
\item
  \textbf{Rubinstein, M. (2001).} Low-lying zeros of \(L\)-functions and
  random matrix theory. \emph{Duke Mathematical Journal, 109}(1),
  147--181. Theorem 3.1, pp.~150--151.
  \url{https://web.williams.edu/Mathematics/sjmiller/public_html/ntprob17/handouts/Lfns/Rubinstein_ThesisPaper_Duke.pdf}
  . Original PDF inspected. The theorem's explicit absence of an RH
  assumption was checked.
\item
  \textbf{Belabas, K., Diaz y Diaz, F., \& Friedman, E. (2008).} Small
  generators of the ideal class group. \emph{Mathematics of Computation,
  77}(262), 1185--1197. \url{https://doi.org/10.1090/S0025-5718-07-02003-0} .
  Inspected author preprint dated 5 December 2005:
  \url{https://www.math.u-bordeaux.fr/~kbelabas/research/OnBach.pdf}
  . Introduction pp.~2--3 and Theorem 1.1, p.~3.
\item
  \textbf{Bach, E., \& Sorenson, J. (1996).} Explicit bounds for primes
  in residue classes. \emph{Mathematics of Computation, 65}(216),
  1717--1735. \url{https://doi.org/10.1090/S0025-5718-96-00763-6} . Original
  author-hosted PDF inspected:
  \url{https://www.researchgate.net/profile/Jonathan-Sorenson-2/publication/220576951_Explicit_bounds_for_primes_in_residue_classes/links/55100ea70cf2ac2905afbceb/Explicit-bounds-for-primes-in-residue-classes.pdf}
  . Theorem 5.3, p.~1731; §2, p.~1720; Corollaries 3.4--3.5,
  pp.~1722--1723. The numerical computations reported by the authors
  were not rerun here.
\end{enumerate}

\subsection{7. Check record and unresolved bibliographical
gates}\label{check-record-and-unresolved-bibliographical-gates}

Every new proof in this module must be read as a whole after
composition. The companion
\nolinkurl{consequences_criteria_CHECK_EVENTS.tsv} records the actual
completion of those reads and any correction; it does not call a
same-author proof read an independent re-proof. The adopted Alpöge
statements receive an event for their use, not for a claimed
verification of the unseen full archive.

The unresolved bibliographical items are narrowly delimited: direct
inspection of Bach's original Theorem 4 and the original Selberg
equation (13). These are not presented as completed direct-source
checks. The operative least-prime, cohomological, effective-Siegel and
criterion results above have explicit inspected primary locators. No
mathematical claim of global logical independence is made from a failure
of the current factorization route.

% END consequences_criteria.tex

% BEGIN consequences_criteria_GT_addendum.tex
\section{Current Grothendieck--Teichmüller scope: 20 November
2026}\label{current-grothendieckteichmuxfcller-scope-20-november-2026}

Research addendum \nolinkurl{consequences_criteria}; GPT-6 Astra Pro;
effort setting not exposed. This updates the literature status without
enlarging the adopted Alpöge input or claiming a kernel check of another
proof archive.

\subsection{CC-GT01. The 2026 rational claim has substantive primary
support}\label{cc-gt01.-the-2026-rational-claim-has-substantive-primary-support}

It would be outdated to label rational Deligne--Drinfeld freeness simply
open. OpenAI's 23 September 2026 preprint, Theorem 1.1, asserts freeness
of the rational graded solution Lie algebra on one generator in each odd
weight at least three, together with its weight-completed version. Its
mirrored introduction, the statement of its all-weight upper bound, and
parts of the argument were inspected; the entire 41-page original proof
and its Lean archive were not independently checked here. Plectis Atlas
itself is an index of other parties' reports, not our verification
method. The Valency page could not be retrieved directly.

Two subsequent primary mathematical texts supply stronger evidence than
an index entry:

\begin{longtable}[]{@{}
  >{\raggedright\arraybackslash}p{(\columnwidth - 4\tabcolsep) * \real{0.3333}}
  >{\raggedright\arraybackslash}p{(\columnwidth - 4\tabcolsep) * \real{0.3333}}
  >{\raggedright\arraybackslash}p{(\columnwidth - 4\tabcolsep) * \real{0.3333}}@{}}
\toprule\noalign{}
\begin{minipage}[b]{\linewidth}\raggedright
Source
\end{minipage} & \begin{minipage}[b]{\linewidth}\raggedright
Exact locator inspected
\end{minipage} & \begin{minipage}[b]{\linewidth}\raggedright
Mathematical contribution
\end{minipage} \\
\midrule\noalign{}
\endhead
\bottomrule\noalign{}
\endlastfoot
Willwacher (2026b), arXiv:2610.15592v4, 28 October & Theorem 1.1, p.~2;
Theorem 2.2, p.~3, proved in §§3.1--3.3, pp.~4--6 & The required upper
bound \(\dim_{\mathbb Q}\mathfrak{grt}_1(\mathbb Q)_n\le d_n\), where
\(d_n\) is the dimension of the free Lie algebra on odd generators \\
Hirose (2026), arXiv:2610.23719v1, 16 October & Theorem 4 and §1.2, p.~2
& An integral Hoffman spanning theorem modulo confluence relations,
combined there with Brown and Furusho to identify the motivic group with
\(\mathrm{GRT}\) \\
\end{longtable}

For the first paper, the full chain proving the upper bound was read:
integral rank comparison, the characteristic-two projection, its
explicit reconstruction, the image restriction, and the free-Lie
dimension calculation. The key derivation intertwining and square-zero
identities were checked algebraically. This is a targeted reading of an
existing proof, not an independent reproof or a formal verification.
Theorem 1.3 and §4 give further characteristic-two structure, but they
are unnecessary when Brown's injection is retained. Hirose's theorem and
stated motivic consequence were inspected; his full technical proof was
not checked in this run.

Neither proof uses RH or Dirichlet GRH. They are external 2026
mathematical results, not achievements attributable to the adopted
analytic input.

\subsection{CC-GT02. What Drinfeld and Brown had already
supplied}\label{cc-gt02.-what-drinfeld-and-brown-had-already-supplied}

Drinfeld (1991), Proposition 6.3, p.~859, gives a homogeneous element of
\(\mathfrak{grt}_1\) in every odd degree \(n\ge3\) with nonzero
depth-one term. The unnumbered ``Questions'' immediately following ask
about precisely one generator in each odd degree and freeness. The
primary proposition, proof and questions were visually inspected on the
rendered original scan; its existence assertion should not be conflated
with freeness or exhaustion.

Brown (2012), Theorem 1.1, p.~950, gives the motivic Hoffman basis; the
paragraph immediately following explains why it proves Conjecture 1 on
p.~949, namely faithfulness of the motivic action. The structure of the
motivic prounipotent Lie algebra is recalled on p.~950. An exact
numbered modern formulation of its free-subalgebra consequence is
Willwacher (2026a), Theorem 1.1, p.~1. His Theorem 1.2, also p.~1,
extends the assertion to any homogeneous odd-degree family with nonzero
depth-one coefficients. These are injectivity statements. The upper
bound in CC-GT01 supplies the distinct exhaustion step.

\subsection{CC-GT03. Combining injection and upper bound, including
completion}\label{cc-gt03.-combining-injection-and-upper-bound-including-completion}

\textbf{Statement.} Let \[
 \mathfrak f=\operatorname{Lie}_{\mathbb Q}\langle e_3,e_5,e_7,\ldots\rangle,
 \qquad \deg e_j=j.
\] Write
\(\mathfrak g^{\mathrm{gr}}=\bigoplus_n\mathfrak{grt}_{1,n}(\mathbb Q)\)
for the graded algebra of polynomial solutions, and
\(\widehat{\mathfrak g}=\prod_n\mathfrak{grt}_{1,n}(\mathbb Q)\) for its
weight completion; some sources use \(\mathfrak{grt}_1\) for the latter.
Take the weight-preserving injection
\(\iota:\mathfrak f\hookrightarrow\mathfrak g^{\mathrm{gr}}\) supplied
by Brown's result and the upper bound of Willwacher (2026b), Theorem
1.1. Then \(\iota\) is an isomorphism, and its extension
\(\prod_n\mathfrak f_n\to\widehat{\mathfrak g}\) is an isomorphism of
complete Lie algebras.

\textbf{Proof.} In every weight \(n\), the space \(\mathfrak f_n\) is
finite-dimensional: only finitely many generators can occur, and every
generator has degree at least three. Injectivity gives \[
 \dim\mathfrak f_n\le\dim\mathfrak{grt}_{1,n}.
\] The cited upper bound reverses this inequality. Therefore the
injective linear map \(\iota_n\) is bijective in every weight. Their
direct sum is a bijective graded Lie homomorphism, hence a Lie algebra
isomorphism. Their product is a continuous bijection for the product,
equivalently weight-filtration, topologies; its inverse is the product
of the degreewise inverses and is continuous. A bracket component of
fixed weight is a finite sum of brackets of lower-weight components, so
the product map and its inverse preserve the completed bracket. This
proves the completed assertion. \(\square\)

This deduction identifies exactly which result completed the rational
theorem. It supplies no theorem about the full profinite absolute Galois
group.

\subsection{CC-GT04. The full profinite question is a different
assertion}\label{cc-gt04.-the-full-profinite-question-is-a-different-assertion}

Fix Drinfeld's original profinite group \(\widehat{\mathrm{GT}}\), with
its commutator, invertibility, symmetry, hexagon and pentagon
conditions. Bortnovskyi, Dolgushev, Holikov, and Pashkovskyi (2026
revision), §1.3, equations (1.5)--(1.8), pp.~4--5, gives the fixed Ihara
embedding and cyclotomic compatibility. In that paper's coordinates, \[
 \operatorname{Ih}(g)=\bigl((\chi_{\mathrm{cyc}}(g)-1)/2,f_g\bigr),
 \qquad
 x\mapsto x^{\chi_{\mathrm{cyc}}(g)},\quad
 y\mapsto f_g^{-1}y^{\chi_{\mathrm{cyc}}(g)}f_g.
\] Its cyclotomic projection is \((m,f)\mapsto2m+1\), so the composite
equals \(\chi_{\mathrm{cyc}}\). The paper's introduction distinguishes
surjectivity of this whole embedding from its proved surjectivity onto
particular quotients; its Theorems 5.3 and 5.7 concern such quotient
results. Rational Lie-algebra freeness does not establish full profinite
surjectivity.

The safe status statement for this report is precise: \textbf{the full
profinite identification is not settled by the adopted RH/Dirichlet GRH
or by the rational Deligne--Drinfeld theorem; the inspected 2026 primary
profinite work treats full surjectivity as unresolved.} No exhaustive
claim about all unpublished assertions is needed. The gentle,
pro-\(\ell\), prounipotent, and higher-genus variants must not be
substituted for the fixed original profinite embedding.

\subsection{CC-GT05. The exact positive role of the cyclotomic
quotient}\label{cc-gt05.-the-exact-positive-role-of-the-cyclotomic-quotient}

The cyclotomic character and Kronecker--Weber identify finite abelian
characters of \(G_{\mathbb Q}\) with Dirichlet characters. Consequently
the adopted GRH controls the analytic factors of finite abelian Artin
representations and their Tate twists, including the Borger realizations
of CC-F01 and CC-F01b. Those results also apply to arithmetic
representations known to factor through the cyclotomic projection in the
preceding diagram.

This analytic statement concerns those representations and their
\(L\)-functions. It does not determine the kernel of the cyclotomic
projection, the associator component \(f_g\), or surjectivity of
\(\operatorname{Ih}\). Abstractly, surjectivity after projection cannot
establish surjectivity before projection: for any profinite abelian
group \(A\), the map \(A\to A\times C_2\), \(a\mapsto(a,1)\), has
surjective composite with the projection to \(A\) and is a proper
embedding. This elementary observation is a counterexample to that
general inference, not a counterexample to any conjecture about the
actual group \(\widehat{\mathrm{GT}}\).

\subsection{Exact bibliography and access
record}\label{exact-bibliography-and-access-record}

\begin{enumerate}
\def\labelenumi{\arabic{enumi}.}
\tightlist
\item
  \textbf{Brown, F. (2012).} Mixed Tate motives over \(\mathbb Z\).
  \emph{Annals of Mathematics, 175}(2), 949--976.
  \url{https://doi.org/10.4007/annals.2012.175.2.10} . Primary PDF:
  \url{https://annals.math.princeton.edu/wp-content/uploads/annals-v175-n2-p10-p.pdf}
  . Theorem 1.1 and its stated consequence directly inspected.
\item
  \textbf{Drinfeld, V. G. (1991).} On quasitriangular quasi-Hopf
  algebras and on a group that is closely connected with
  \(\operatorname{Gal}(\overline{\mathbb Q}/\mathbb Q)\).
  \emph{Leningrad Mathematical Journal, 2}(4), 829--860. Original
  Russian publication: \emph{Algebra i Analiz, 2}(4), 149--181 (1990).
  Scan: \url{https://www.lamfa.u-picardie.fr/marin/docs/drinfeld.pdf} .
  Proposition 6.3 and the following unnumbered Questions, p.~859,
  visually inspected on the rendered original page.
\item
  \textbf{Willwacher, T. (2026a).} \emph{On free generators in the
  Grothendieck--Teichmüller Lie algebra} (arXiv:2607.18793v2, 6 August
  2026). \url{https://arxiv.org/pdf/2607.18793v2} . Theorems 1.1--1.2, p.~1,
  directly inspected.
\item
  \textbf{Willwacher, T. (2026b).} \emph{The Deligne--Drinfeld
  conjecture after OpenAI} (arXiv:2610.15592v4, 28 October 2026).
  \url{https://arxiv.org/pdf/2610.15592v4} . Theorem 1.1 and its entire proof
  chain in §§1--3 read; key algebraic identities checked. Theorem 1.3
  and §4 are additional, not required inputs to CC-GT03.
\item
  \textbf{Hirose, M. (2026).} \emph{Confluence relations, Hoffman basis
  over \(\mathbb Z_{(2)}\), and the Deligne--Drinfeld conjecture}
  (arXiv:2610.23719v1, 16 October 2026).
  \url{https://arxiv.org/pdf/2610.23719v1} . Theorem 4 and §1.2, p.~2,
  directly inspected; full technical proof not independently checked
  here.
\item
  \textbf{Bortnovskyi, I., Dolgushev, V. A., Holikov, B., \&
  Pashkovskyi, V. (2024/2026).} \emph{Accessing non-abelian quotients of
  the Grothendieck--Teichmueller group via elementary tools}
  (arXiv:2405.11725v3, 6 August 2026).
  \url{https://arxiv.org/pdf/2405.11725v3} . Introduction, §1.3, equations
  (1.5)--(1.8), pp.~4--5, and the described quotient scope directly
  inspected.
\item
  \textbf{OpenAI. (2026, September 23).} \emph{The Deligne--Drinfeld
  conjecture}. OpenAI Math Release preprint.
  \url{https://github.com/openai/math/blob/fd4aeeb2ee4fc729c18d98444fed42fd0529eeeb/preprints/The-Deligne-Drinfeld-conjecture-September-23-2026/paper.pdf}
  . Inspected mirror:
  \url{https://wcook04.github.io/plectis-atlas/papers/The-Deligne-Drinfeld-conjecture-September-23-2026.html}
  . Theorem 1.1 and dimension-bound route inspected; the complete
  original proof and Lean archive were not checked in this run.
\end{enumerate}

% END consequences_criteria_GT_addendum.tex

% BEGIN source_and_masullo.tex
% Candidate section. This file is intended for inclusion in the root paper.
% IDs F-SM-* identify results proved in this run, not claims of historical novelty.
\section{The two supplied routes: exact scope and mathematical bridges}
\label{sec:source-masullo}

\subsection{What is adopted, what has been checked, and what is not exported}

We use Alp\"oge's five declarations in \nolinkurl{ALPOGE_Main.lean} as the research
inputs specified in the task.  The check performed here is an inspection of the
declarations and of \nolinkurl{ALPOGE_Statements.lean}, not a fresh Lean kernel
verification: the import \nolinkurl{Supplements.Kernels.MainTheoremsAll} is absent
from the supplied ZIP, as is the remainder of the formal development.  Neither
\nolinkurl{lean} nor \nolinkurl{lake} is installed in the working environment.  The
supplied scan is a report about another source tree; no claim is made that the
175,187 files were rescanned here.  This distinction does not change the adopted
mathematical premises, but it determines the correct check-event label.

The declaration names and their exact targets are as follows.
\begin{enumerate}
\item \texttt{riemannHypothesis : RiemannHypothesis := rhOfAll}; the target is
Mathlib's named proposition.
\item \texttt{riemannZeta\_eq\_zero\_re : RiemannZetaZerosOnLine := rhOfAll};
the supplied definition excludes the negative even integers and $s=1$.
\item \texttt{LFunction\_ne\_zero\_of\_half\_lt\_re :
DirichletLNonvanishingRightOfHalf := allChi}; this quantifies over every positive
modulus, every complex Dirichlet character, and $\Re s>1/2$, with the exception
$(\chi,s)=(1,1)$.
\item \texttt{LFunction\_eq\_zero\_re : DirichletLZerosOnLine := claimAll.2.1};
this concerns zeros in the open strip $0<\Re s<1$.
\item \texttt{LFunction\_eq\_zero\_re\_of\_isPrimitive :
PrimitiveDirichletLZerosOnLine := critAll}; this is the primitive-character
specialization, again in the open strip.
\end{enumerate}
The pinned Mathlib revision is
\nolinkurl{0df444a360eaa60ab8c11dca51a86af692955474}, in
\nolinkurl{ALPOGE_lakefile.toml}.  None of these five declarations states a
weight-monodromy theorem, a theorem on all Hecke characters of all number
fields, a nonabelian Artin holomorphy theorem, or a theorem specifying an
integral algebraic $K$-group of $\mathbb Z$.

Two corrections to the prompt's reading of the four-page note are material.
The theta-lift chain has length of order $\sqrt X$, both in the paragraph on
page~1 and in step~8 on page~4.  The radical is visible in the PDF.  The claim
about $K$-groups in step~17, footnote~6, is explicitly tentative and gives no
degrees or quantifiers.  It cannot be turned into a specified $K$-group
vanishing theorem merely by citing the note.  The relevant export, if one
exists in the full development, must first be identified.

\subsection{Masullo's different target and its exact remaining gate}

Masullo, \emph{The Music of the Primes}, page~6, specifies all number fields
and all unitary Hecke characters.  This is a larger target than the five
Alp\"oge declarations above.  The analytic route can be located precisely:
\begin{enumerate}
\item Section~1.2.1, pages~13--15: a divisor exact sequence of nuclear
Fr\'echet spaces, then Hardy--Sobolev Hilbert completions whose spectral
quotients retain zeros in a prescribed vertical strip.
\item Sections~1.2.2--1.2.3, pages~15--16, and Chapter~11: orbital bases
and a decomposition of a lifted geometric kernel into local terms.
\item Theorem~10.4.2.8, pages~463--471: the spectral trace after the
relevant truncation/projection operations.
\item Theorem~11.2.2.1, pages~490--491: orbital nuclearity on the
zero line for Sobolev degree $s\geq 2[F:\mathbb Q]$.
\item Theorem~11.3.2.1, pages~495--496: orbital trace calculations;
Theorem~12.4.1.1, pages~514--519: the $S$-local trace calculation.
\item Theorem~12.5.1.1, pages~519--520: the proposed global trace
identity with only the zeros inside the chosen strip on its spectral side.
\item Sections~1.4.2--1.4.5, pages~28--30: the proposed deduction of
zero location by holomorphic families of test functions.
\end{enumerate}
The end of Chapter~13 is not a usable theorem statement: on page~522,
Theorem~13.3.1.2 consists of a heading with instructions to insert the
statement and prove the main results.  Theorem~13.3.1.1 on the same page is
marked for correction.  More substantially, Section~12.3 on page~514, devoted
to the global-to-$S$-local deconstruction, consists only of its heading and an
instruction to include the section.  The proof of Theorem~12.5.1.1 on
page~520 uses the resulting decomposition and exchange of diagonal integration
with the local sum.  These are concrete missing links in the supplied written
argument.  The following proposition supplies the zero-location deduction
once the appropriate trace identity has been established; it does not supply
the missing kernel identity.

\paragraph{Finding F-SM-01 (a shrinking-strip trace criterion).}
Let $\mathcal Z$ be a locally finite multiset in
$\{\lambda\in\mathbb C:|\Re\lambda|\leq B\}$, with positive integral
multiplicities $m(\lambda)$, and suppose
\[
 \sum_{\lambda\in\mathcal Z}
       \frac{m(\lambda)}{1+|\Im\lambda|^2}<\infty.
\]
For $f\in C_c^\infty(0,\infty)$ put
$\widehat f(\lambda)=\int_0^\infty f(t)e^{\lambda t}\,dt$.
Assume that for some sequence $a_j\downarrow0$,
\begin{equation}\label{eq:sm-strip}
 \sum_{\lambda\in\mathcal Z}m(\lambda)\widehat f(\lambda)
  =\sum_{\substack{\lambda\in\mathcal Z\\|\Re\lambda|<a_j}}
           m(\lambda)\widehat f(\lambda)
 \qquad(f\in C_c^\infty(0,\infty)).
\end{equation}
Then every $\lambda\in\mathcal Z$ is imaginary.  Conversely, if every
$\lambda$ is imaginary, \eqref{eq:sm-strip} holds for every $a_j>0$.

\begin{proof}
Fix $a=a_j$ and let $\mathcal E$ be the submultiset with
$|\Re\lambda|\geq a$.  The hypothesis gives
$\sum_{\mathcal E}m(\lambda)\widehat f(\lambda)=0$ for every such test
function.  Fix a real number $w>B$ and a complex number $z$ with
$\Re z>B$.  On $t\geq0$ define
\[
 h_{z,w}(t)=e^{-zt}-e^{-wt}.
\]
Choose smooth functions $\eta,\zeta$ on $[0,\infty)$ with $0\leq\eta,\zeta
\leq1$, $\eta=0$ on $[0,1]$, $\eta=1$ on $[2,\infty)$, $\zeta=1$ on
$[0,1]$, and $\zeta=0$ on $[2,\infty)$.  Set
$f_{\epsilon,R}(t)=\eta(t/\epsilon)\zeta(t/R)h_{z,w}(t)$, where
$0<\epsilon\leq1$ and $R\geq1$.  These are admissible compactly supported
smooth test functions on $(0,\infty)$.

For $\beta\in[-B,B]$, the $L^1$ norms of
$g_{\epsilon,R,\beta}=e^{\beta t}f_{\epsilon,R}(t)$ and its second
derivative are bounded independently of $\epsilon,R,\beta$.  Here is the
estimate at the only potentially singular endpoint.  Since $h_{z,w}(0)=0$,
$h_{z,w}(t)=O(t)$, $h'_{z,w}(t)=O(1)$ and $h''_{z,w}(t)=O(1)$ for
$0\leq t\leq2$.  On $[\epsilon,2\epsilon]$ the terms containing
$\eta''(t/\epsilon)/\epsilon^2$ therefore have $L^1$ norm
$O(\epsilon^{-2}\epsilon\epsilon)=O(1)$, and those containing
$\eta'(t/\epsilon)/\epsilon$ also have $L^1$ norm $O(1)$.
The remaining terms near zero are bounded.  At infinity all terms are
bounded by a polynomial in $t$ times $e^{-ct}$, with
$c=\min(\Re z,w)-B>0$; the derivatives of $\zeta(t/R)$ have size
$O(R^{-1})$ and $O(R^{-2})$.  This proves the uniform bounds.
Twice integrating by parts in the Fourier variable, together with the
$L^1$ bound when $|\gamma|\leq1$, now gives
\[
 |\widehat f_{\epsilon,R}(\beta+i\gamma)|
 \leq\frac{C}{1+\gamma^2}
 \qquad(\beta\in[-B,B]).
\]
Dominated convergence in $t$, followed by dominated convergence in the
absolutely summable zero series, gives
\[
 0=\sum_{\lambda\in\mathcal E}m(\lambda)
       \left(\frac1{z-\lambda}-\frac1{w-\lambda}\right)
       \qquad(\Re z>B).
\]
The series on the right is normally convergent on every compact subset of
$\mathbb C\setminus\mathcal E$: outside a finite set each summand is bounded
by $C_Km(\lambda)/(1+|\Im\lambda|^2)$.  It therefore defines a
meromorphic function on $\mathbb C$.  At any $\lambda_0\in\mathcal E$
its residue is the strictly positive integer $m(\lambda_0)$; local finiteness
ensures that the other terms are holomorphic near $\lambda_0$.
The identity theorem, applied to the connected complement of the discrete
set $\mathcal E$, makes the meromorphic function identically zero.  A zero
meromorphic function has no pole, so $\mathcal E$ is empty.  This holds for
every $a_j$, proving $\Re\lambda=0$ for every $\lambda$.  The converse
is immediate because no terms are omitted.
\end{proof}

\paragraph{Finding F-SM-02 (the precise theorem-level comparison).}
For one fixed unitary Hecke character, suppose the established full explicit formula
and a proposed shrinking-strip trace formula have exactly the same
geometric distribution on test functions supported at idelic norm $>1$.
After taking its character isotype and setting $t=\log|g|$, their equality
is \eqref{eq:sm-strip} for $\lambda=\rho-1/2$.  The usual zero-counting
bound $N(T)=O(T\log(T+2))$ implies the summability required in
F-SM-01, so the trace formula implies GRH for this character.  Conversely,
GRH implies this numerical trace identity by the full explicit formula.
For the Dirichlet characters covered by Alp\"oge's adopted theorem, the
numerical identity is therefore available without assuming Masullo's
operator-theoretic construction.

\begin{proof}
The trace distributions subtract to exactly \eqref{eq:sm-strip}.  By a
dyadic decomposition,
\[
 \sum_{2^k<|\Im\lambda|\leq2^{k+1}}
 \frac{m(\lambda)}{1+|\Im\lambda|^2}
 \ll \frac{2^k(k+1)}{2^{2k}},
\]
whose sum is finite; local finiteness handles bounded ordinates.
Apply F-SM-01.  Conversely, all zeros lie in every open strip of positive
width centered on the critical line, so restricting the spectral sum removes
no terms.  This proves the equality of distributions.  No assertion about
the existence of a particular kernel, basis, or geometric realization is
used or obtained in this converse.
\end{proof}

The full idelic explicit formula is recorded in Masullo's Theorem~B.4.4.7,
page~623, whereas the restricted one is Theorem~12.5.1.1,
pages~519--520.  To apply F-SM-02 to the manuscript one must first align the
Haar-measure and multiplicity conventions in those statements, and prove
the missing deconstruction.  F-SM-01 is independent of those normalization
choices.  It also avoids passing separately to divergent archimedean
integrals: the difference of exponentials vanishes at the identity.

\subsection{A larger Artin bridge, conditional on Masullo's larger target}

\paragraph{Finding F-SM-03 (off-line Artin holomorphy from Hecke GRH).}
Assume GRH for every finite-order Hecke character of every number field.
Let $K$ be a number field and $V$ a finite-dimensional complex
representation of $\operatorname{Gal}(\overline K/K)$ with finite image.
Then its Artin $L$-function is holomorphic and nonzero on
\[
 U=\{s:0<\Re s<1,\quad\Re s\ne1/2\}.
\]
The same assertion for $V(n)$ holds on $U-n$ under the convention
$L(V(n),s)=L(V,s+n)$.

\begin{proof}
Choose a finite Galois extension $E/K$ through which $V$ factors and put
$G=\operatorname{Gal}(E/K)$.  Brauer induction gives
\[
 [V]=\sum_{j=1}^r n_j\operatorname{Ind}_{H_j}^{G}\psi_j,
 \qquad n_j\in\mathbb Z,\qquad\dim\psi_j=1.
\]
The $\psi_j$ have finite order.  With $K_j=E^{H_j}$, Artin reciprocity
and inductivity of Artin $L$-functions give the meromorphic identity
\[
 L_K(V,s)=\prod_{j=1}^r L_{K_j}(\psi_j,s)^{n_j}.
\]
Each Hecke factor is holomorphic in $0<\Re s<1$, including the trivial
character's Dedekind zeta function, whose possible pole is at $1$.
The assumed GRH says that no such factor vanishes on $U$.  An integral
power, positive, zero, or negative, of a nowhere-zero holomorphic function
is again nowhere-zero and holomorphic.  Their finite product has the same
property.  Substituting $s+n$ proves the twisted assertion.
\end{proof}

The exact representation-theoretic input is Serre (1977),
\emph{Linear Representations of Finite Groups}, \S10.5, Theorem~20,
page~78.  The $L$-function factorization is equation~(2) in Cogdell,
\emph{On Artin $L$-functions}, \S2.1, page~10; the compatibility of the
ramified local factors is explained in \S2.2, page~12.  This is a conditional
bridge to all finite-image Artin motives, substantially beyond the
abelian-over-$\mathbb Q$ scope of the supplied Alp\"oge theorem.
It does not prove holomorphy on the critical line: factors in the denominator
can vanish there.  Nor does the proof imply general automorphy, the Artin
conjecture, or the existence of a motivic realization functor.

\subsection{Two explicit corrections to the geometric analogies}

\paragraph{Finding F-SM-04 (the norm does not recover the unit strata).}
The assertion on Masullo page~493 that a saturated all-local-units stratum
in $\mathbb A_F^\times$ is the inverse image of the corresponding unit
stratum of $\mathbb A_{\mathbb Q}^\times$ under the norm map is false.
Here ``saturated'' has its explicit meaning from page~492: a set of places
contains all places over a rational prime whenever it contains one.
The displayed stratum on page~493 is a subset of the ideles themselves,
defined by $g_v\in\mathcal O_{F_v}^{\times}$ at each listed place;
it is not a subset enlarged by multiplication by $F^{\times}$.
The correct condition at a rational prime $p$ for the norm to be a unit is
\[
 \sum_{v\mid p} f(v/p)v_v(g_v)=0,
\]
where $v_v$ assigns value $1$ to a uniformizer and $f(v/p)$ is the residue
degree.  The stronger all-local-units condition is $v_v(g_v)=0$ for
every $v\mid p$.

\begin{proof}
For a finite extension $F_v/\mathbb Q_p$, a uniformizer has norm of
$p$-adic valuation $f(v/p)$, and the norm of a unit is a unit; by factoring
an arbitrary element into a uniformizer power times a unit,
$v_p(N_{F_v/\mathbb Q_p}(g_v))=f(v/p)v_v(g_v)$.  Multiplication of the
local norm factors proves the displayed formula.
Take $F=\mathbb Q(i)$ and $p=5$, which splits into two degree-one places
$v_1,v_2$.  Thus $F_{v_j}\simeq\mathbb Q_5$.  Let the idele $g$ have
$g_{v_1}=5$, $g_{v_2}=5^{-1}$ and every other component $1$.
Its norm is the identity idele because the two $5$-adic components multiply
to $1$.  Neither of these components is a local unit, so $g$ does not lie
in the asserted unit stratum.  This is the required counterexample.
The all-local-units stratum remains well-defined directly; the error in
identifying it with a norm inverse image does not by itself disprove
the later trace identity.
\end{proof}

\paragraph{Finding F-SM-05 (a maximal curve need not have gonality $g+1$).}
The equality for maximal curves asserted on Masullo page~535 is false.
Over $\mathbb F_9$, the smooth plane curve
\[
 H:\quad Y^3Z+YZ^3=X^4
\]
is maximal of genus $3$ and has gonality at most $3$, hence strictly
less than $g+1=4$.

\begin{proof}
In characteristic $3$, the three partial derivatives of
$Y^3Z+YZ^3-X^4$ are $-X^3,Z^3,Y^3$, so they have no common projective
zero.  The curve is smooth.  Over an algebraic closure, any two positive-degree
plane-curve components intersect by B\'ezout's theorem, creating a singular
point; a repeated component would also contradict smoothness.  Thus this
smooth plane quartic is geometrically integral.  Its genus is
$(4-1)(4-2)/2=3$ (the plane-curve genus formula, Stacks Project,
Lemma~53.9.3).  For each $x\in\mathbb F_9$, $x^4$
is the norm of $x$ and lies in $\mathbb F_3$; the trace
$y\mapsto y^3+y$ from $\mathbb F_9$ onto $\mathbb F_3$ has three
elements in each fibre.  There are therefore $9\cdot3=27$ affine
points.  At $Z=0$ the equation forces $X=0$, leaving the unique point
$[0:1:0]$.  Thus $\#H(\mathbb F_9)=28=9+1+2\cdot3\cdot3$,
which is the maximal value.  The nonconstant rational function $x=X/Z$
has function-field degree at most $3$, because $y=Y/Z$ satisfies
$y^3+y=x^4$ over $\mathbb F_9(x)$.  It extends to a morphism to
$\mathbb P^1$ on the smooth projective curve; consequently the gonality
is at most $3$.
\end{proof}

The correction strengthens rather than defeats the point being made in
Masullo \S14.2, pages~534--536: the proposed number-field multiplicity
mechanism cannot simply be transported to function fields.  For this
curve the Frobenius eigenvalue multiplicity in $H^1$ is $2g=6$, whereas
the gonality is at most $3$: the finite-field Weil theorem gives modulus
$3$ for each of its six eigenvalues, and their sum is
$9+1-28=-18$, so equality in the triangle inequality makes all six
equal to $-3$.

\subsection{The categorical and positivity information still required}

The stack diagrams are especially relevant to Grothendieck's programme,
but their stated status matters.  Masullo's equations~(8.4.0.6) and
(8.4.0.7), page~390, and (8.4.0.8)--(8.4.0.9), page~391, relate
Fr\'echet-space divisor sequences to periodic-orbit quotient stacks and
boundary divisors.  Remark~8.4.0.1 on page~391 explicitly says that the
vertical arrows are not morphisms.  They express an interpretation, not
a constructed sheaf, realization functor, or equivalence of categories.
Consequently these diagrams do not supply an arithmetic cohomology theory
of $\operatorname{Spec}\mathbb Z$, a six-functor formalism, or a comparison
with the \emph{D\'erivateurs}.  A literal categorical bridge would require
those arrows to be constructed and their compatibility with the trace,
tensor products, Tate twists, and local terms proved.

\paragraph{Finding F-SM-06 (spectral location alone does not supply
semisimplicity or a positive unitary realization).}
There exists a complex finite-dimensional operator with purely imaginary
spectrum that is not skew-adjoint for any positive definite Hermitian
inner product.

\begin{proof}
On $\mathbb C^2$ take $A=i\gamma I+N$, where $\gamma\in\mathbb R$,
$Ne_2=e_1$ and $Ne_1=0$.  Its spectrum is the singleton $\{i\gamma\}$.
If $A$ were skew-adjoint for a positive definite Hermitian form, then
$e^{tA}$ would preserve its norm for every real $t$: differentiating
$\langle e^{tA}v,e^{tA}v\rangle$ gives zero.  But
$e^{tA}e_2=e^{i\gamma t}(e_2+te_1)$, whose norm is at least
$|t|\|e_1\|-\|e_2\|$ by the triangle inequality and therefore tends
to infinity.  This contradicts norm preservation.
\end{proof}

This refutes a general inference, not RH and not the possibility of a
particular Hilbert--P\'olya realization.  Masullo's spectral discussion
expressly permits Jordan blocks (\S1.2.1, page~14), so zero location
must be separated from a theorem killing nilpotent parts or proving
zero simplicity.

\paragraph{Finding F-SM-07 (a bounded biorthogonal lift is not automatically
a Schauder basis).}
The proof step in Masullo's Lemma~10.4.2.5, page~462, that obtains
the Schauder basis property from the displayed boundedness and
biorthogonality alone is invalid.

\begin{proof}
Let $E=\ell^2(\{0,1,2,\ldots\})$, $E_0=\mathbb Ce_0$, $F=E$,
and let $x_n$ be the image of $e_n$ in $E/E_0$ for $n\geq1$.
These $x_n$ form the usual orthonormal basis.  Set $u_n=e_0+e_n$
and let $\phi_n$ be the $n$th coordinate functional.  Then
$\phi_n(u_m)=\delta_{nm}$ and $\|u_n\|=\sqrt2$.
The collar conditions hold with $r=1$, $r'=11/10$ and $\epsilon=1/2$.
Yet
\[
 \frac1N\sum_{n=1}^{N}u_n
   =e_0+\frac1N\sum_{n=1}^{N}e_n\longrightarrow e_0
\]
in $E$.  Thus $e_0$ belongs to the closed span of the $u_n$, while
every $\phi_n(e_0)=0$.  If the $u_n$ were a Schauder basis of that
closed span, applying the continuous biorthogonal coordinate functionals
to the expansion of $e_0$ would give every coefficient zero, forcing
$e_0=0$, a contradiction.
\end{proof}

This counterexample concerns the inference used in the written proof.
It does not refute the existence of suitable lifts in the Hilbert-space
case.  The following construction supplies that repair in the form relevant
to the Hardy--Sobolev Hilbert spaces.

\paragraph{Finding F-SM-08 (Hilbert-space basis lifting with a controlled
graph).}
Let $E$ be a Hilbert space, $E_0$ a closed subspace, $q:E\to E/E_0$ the
quotient, and $F\subset E$ a linear subspace such that $F\cap E_0$ is
dense in $E_0$.  Let $(x_n)$ be a Schauder basis of $E/E_0$ with
coefficient functionals $(\phi_n)$, or a Markushevich basis with
its total biorthogonal functionals $(\phi_n)$, and assume each
$x_n\in q(F)$.
For any positive numbers $\delta_n$ satisfying
$\sum_n\|\phi_n\|\delta_n<\infty$, there are $u_n\in F$ with
$q(u_n)=x_n$, $\|u_n-Sx_n\|<\delta_n$, where
$S:E/E_0\to E_0^\perp$ is the isometric orthogonal section, such that
$(u_n)$ is respectively a Schauder or Markushevich basis of its closed span.
That span maps
isomorphically to $E/E_0$ under $q$.

\begin{proof}
Choose $v_n\in F$ with $q(v_n)=x_n$.  Then $Sx_n-v_n\in E_0$.
Density of $F\cap E_0$ provides $w_n\in F\cap E_0$ with
$\|v_n+w_n-Sx_n\|<\delta_n$; put $u_n=v_n+w_n$ and
$d_n=u_n-Sx_n\in E_0$.  The series
\[
 Kx=\sum_n\phi_n(x)d_n
\]
converges absolutely for every $x\in E/E_0$ and defines a bounded
operator to $E_0$, with norm at most
$\sum_n\|\phi_n\|\delta_n$.  Set $T=S+K$.
Then $qT=\operatorname{id}$, so $\|Tx\|\geq\|x\|$ and $T$
is a bounded linear isomorphism onto a closed subspace of $E$.
The inverse is the restriction of $q$.  Since $Tx_n=u_n$, in the
Schauder case the image under $T$ of the basis expansion of $x$ is a
convergent basis expansion in $(u_n)$.  In the Markushevich case,
density of the span and totality of the biorthogonal functionals
$\phi_n\circ q$ on the image are transported by the isomorphism $T$.
Their closed span is $T(E/E_0)$ in both cases, proving every assertion.
\end{proof}

Thus one lift-related gap has an explicit repair.  This does not repair
all of Theorem~10.4.2.8 or the missing Section~12.3: the estimates for
the resulting integral kernels and the interchange of diagonal
integration with deconstruction are separate requirements.

\subsection{Comparison of methods and what each can actually transfer}

\paragraph{Finding F-SM-09 (the exact power-sum bridge to purity).}
Let $\alpha_1,\ldots,\alpha_r$ be distinct nonzero complex numbers,
let $c_1,\ldots,c_r$ be nonzero complex coefficients, and let $R>0$.
Suppose there are constants $C>0$ and $A\geq0$ such that
\[
 \left|\sum_{j=1}^r c_j\alpha_j^n\right|\leq C n^A R^n
 \qquad(n\geq1).
\]
Then $|\alpha_j|\leq R$ for every $j$.  If additionally the set of
$\alpha_j$ is invariant under $\alpha\mapsto R^2/\alpha$, then
$|\alpha_j|=R$ for every $j$.

\begin{proof}
The power-series
\[
 F(z)=\sum_{n\geq1}\left(\sum_{j=1}^r c_j\alpha_j^n\right)z^n
\]
converges absolutely on $|z|<R^{-1}$ by the assumed bound.  Near $0$ it
is the rational function
\[
 \sum_{j=1}^r\frac{c_j\alpha_jz}{1-\alpha_jz}.
\]
For any $j$, the latter has a genuine pole at $z=1/\alpha_j$:
the $j$th summand has residue $-c_j/\alpha_j\ne0$ and the other
summands are holomorphic there because the $\alpha_j$ are distinct.
If $|\alpha_j|>R$, this pole would lie in the disk on which the power
series is holomorphic, contradicting equality by analytic continuation.
Thus $|\alpha_j|\leq R$.  Under the stated reciprocal symmetry,
$R^2/\alpha_j$ is another listed number, so
$R^2/|\alpha_j|\leq R$, which gives the reverse inequality.
\end{proof}

Taking $R=q^{w/2}$ identifies the precise missing inputs for a purity
argument of the form suggested by Alp\"oge's footnote~3: bounds over
\emph{all} finite extensions, a nonzero net coefficient at each eigenvalue,
and the appropriate duality.  A bound for a single finite field is
insufficient for this deduction.  For the trace of an honest finite-dimensional
space the coefficients are the positive algebraic multiplicities, so
they cannot cancel.  For alternating cohomological traces this must
be proved separately.

\paragraph{Finding F-SM-10 (why alternating traces alone do not give purity).}
For any $R>0$ there is a graded vector space with Frobenius whose
alternating traces satisfy every bound in F-SM-09, while a Frobenius
eigenvalue has modulus greater than $R$.

\begin{proof}
Take $V^0=V^1=\mathbb C$, all other graded pieces zero, and let
Frobenius act by the same scalar $a>R$ on both nonzero pieces.
Then $\sum_i(-1)^i\operatorname{tr}(F^n\mid V^i)=a^n-a^n=0$
for every $n\geq1$, while both pieces have eigenvalue $a$.
\end{proof}

This is the exact point of a no-cancellation or positivity input;
it supplies spectral information that a virtual determinant can hide.
Even F-SM-09 determines eigenvalue moduli, not the nilpotent part of
Frobenius, as F-SM-06 illustrates.

Alp\"oge's step~5 (\emph{Remark}, page~3 and footnote~3) seeks a
noncancelling sign in cohomological power sums; his steps~7--10 and~12--16
(page~4) obtain arithmetic multiplicity and norm estimates from unitary
theta lifts, periods, and integral structures.  Masullo's primary gain is
different: a coarser completion retains only part of the zero divisor,
while the geometric trace is asserted to remain unchanged.  F-SM-01
shows precisely why preservation of that trace forces the omitted part
to be empty.  The first route needs a positivity/divisibility estimate;
the second needs trace compatibility through completion, lifting,
localization, and a limiting process.

Neither gain automatically gives algebraic cycles, crystalline
Frobenius, a realization functor, a weight filtration on an arbitrary
motive, or reconstruction of nonabelian fundamental groups.  In the
supplied text the Iwasawa tower is an arithmetic amplification device
(Alp\"oge steps~11 and~19), and the hexagonal CM curve supplies a
specific theta-geometric setting (step~3).  Neither is identified by a
theorem with the Teichm\"uller tower or the profinite GT group.
The testable bridges are the explicit $L$-function identities and the
trace/norm compatibilities just described.  The correspondence between
the words ``tower'', ``hexagon'', or ``multiplicity'' by itself does not
constitute any such compatibility.

The Masullo route also does not eliminate the logical possibility of a
geometric or foliated realization of a different construction.  Its
Sections~1.2.7 and~1.5 compare its own analytic objects with Deninger's
expectations; the absence of a three-manifold in this construction is
not a nonexistence theorem for all other constructions.

% END source_and_masullo.tex

% BEGIN source_masullo_numeric.tex
% Supplement to source_and_masullo.tex. Original proofs of the stated deductions;
% no claim of historical novelty. Background inputs are explicitly identified.
\section{The actual Hilbert-space sheets and a numerical completion over \texorpdfstring{$\mathbb Q$}{Q}}
\label{sec:masullo-numeric}

\subsection{What cote 1 actually says}

The supplied transcription of cote~1, scan~079, gives Corollaries~4 and~5:
coefficientwise domination of the Fredholm determinant of $T$ by that of
$|T|$, and hence $|\det(1+zT)|\leq\det(1+|z||T|)$.  The following
eigenvalue/product theorem and its trace corollary are crossed out, as is the
draft proof continuing on scan~080.  Scans~111--112 state and attempt to prove
the genus-zero lemma used by that draft.  These supplied sheets contain no
Riemann-$\xi$ function, arithmetic spectral realization, or unfinished question
asserting that such a realization should exist.  Thus BR-TR-01 and BR-TR-02
are arithmetic applications of this analytic framework.  The next two
findings also give complete proofs of the actual analytic assertions on the
supplied sheets, including the crossed-out product assertion.  ``Crossed
out'' is a description of this source, not a claim that the classical theorem
is false or remains open.

\paragraph{Finding F-SM-11 (the proposition on scans 111--112, corrected).}
Let $f(z)=\sum_{n\geq0}a_nz^n\not\equiv0$ and
$h(z)=\sum_{n\geq0}b_nz^n\not\equiv0$ be entire, with
$|a_n|\leq|b_n|$ for all $n$.  Suppose that $h$ has genus zero, and that the
nonzero zeros $z_j$ of $f$, repeated with multiplicity, satisfy
$\sum_j|z_j|^{-1}<\infty$.  Then $f$ has genus zero.  Here genus zero means
a constant times a power of $z$ times a normally convergent product of
factors $1-z/z_j$, with no nonconstant exponential factor.

\begin{proof}
Write $h(z)=c z^r\prod_j(1-z/w_j)$, where
$\sum_j|w_j|^{-1}<\infty$.  The function
$H(z)=|c|z^r\prod_j(1+z/|w_j|)$ has nonnegative coefficients that majorize
the absolute values of the coefficients of $h$.  To see this, first use the
finite products, then pass to their locally uniform limits; Cauchy's
coefficient formula preserves each coefficient inequality.  Therefore
$|f(z)|\leq H(|z|)$.

Every genus-zero product $P$ satisfies, for each $\varepsilon>0$,
$|P(z)|\leq C_\varepsilon e^{\varepsilon|z|}$.  Indeed, split its factors
after an index $N$ for which the sum of the reciprocal moduli in the tail is
less than $\varepsilon/2$.  The tail is bounded by
$\exp((\varepsilon/2)|z|)$, using $\log(1+x)\leq x$; the finite product
and the power of $z$ are bounded by a constant times
$\exp((\varepsilon/2)|z|)$.  Applying this to $H$ proves the same bound
for $f$.  Hadamard factorization for an entire function of order at most one
now gives
\[
 f(z)=e^{A+bz}P(z),\qquad P(z)=z^m\prod_j(1-z/z_j).
\]
More explicitly, the genus-one factors in the usual Hadamard formula may
be replaced by genus-zero factors, because $\sum_j1/z_j$ converges
absolutely; the resulting linear exponential is absorbed into $b$.

Suppose $b\ne0$.  Replace $z$ by $\overline b z/|b|$ to make $b>0$;
the preceding radial bounds are unchanged.  Fix
$\pi/3<\theta<\pi/2$, and put $\delta=b\cos\theta/2>0$.  On the sector
$|\arg z|\leq\theta$, the bound for $f$ with $\varepsilon=\delta$ gives
\[
 |P(z)|\leq C e^{-\delta|z|}.
\]
Globally the genus-zero bound for $P$ gives
$|P(z)|\leq C'e^{\delta|z|/3}$.  Set $\omega=e^{2\pi i/3}$.
For every nonzero $z$, at least one of $z,\omega z,\omega^2z$ belongs
to this sector, because its half-width exceeds $\pi/3$.  Hence the entire
function $G(z)=P(z)P(\omega z)P(\omega^2z)$ satisfies
\[
 |G(z)|\leq C(C')^2e^{-\delta|z|/3}.
\]
Liouville's theorem makes $G$ constant, and the displayed bound makes that
constant zero.  A finite product of nonzero entire functions cannot vanish
identically.  This contradicts $P\not\equiv0$, so $b=0$.
\end{proof}

Three transcription-level points matter.  Scan~112 literally specifies
$\pi/3<\theta<\pi$ while subsequently requiring
$\Re z\geq c|z|$ throughout the sector, for $c>0$.
The latter assertion requires $\theta<\pi/2$; the proof above makes the
compatible choice $\pi/3<\theta<\pi/2$.  Scan~111 writes a rotation by
$b/|b|$; with the substitution $f(\Theta z)$ the correcting rotation is
$\Theta=\overline b/|b|$.  Finally, coefficient domination gives
$|f(z)|\leq H(|z|)$, not generally $|f(z)|\leq|h(z)|$ at the same complex
argument, as written on scan~112.  For example, $f=1$ and $h=1+z$
satisfy the coefficient inequality but violate the latter inequality at
$z=-1$.  These corrections repair the proof of the stated proposition.
The uncertain sum/product sign in the transcription of scan~111 is not
used as a premise: the canonical-product hypothesis is stated explicitly.

\paragraph{Finding F-SM-12 (the Fredholm assertions on scans 079--080).}
Let $T$ be a trace-class operator on a complex Hilbert space.  Put
$D_T(z)=\sum_{n\geq0}\alpha_n(T)z^n$, where
$\alpha_n(T)=\operatorname{Tr}(\bigwedge^nT)$ and $\alpha_0(T)=1$.
If $s_j$ are the singular values and $\lambda_j$ the nonzero eigenvalues,
each repeated with its multiplicity (algebraic multiplicity for
$\lambda_j$), then
\[
 |\alpha_n(T)|\leq\alpha_n(|T|),\quad
 |D_T(z)|\leq D_{|T|}(|z|),\quad
 \sum_j|\lambda_j|\leq\|T\|_1,
\]
\[
 D_T(z)=\prod_j(1+z\lambda_j),\qquad
 \operatorname{Tr}T=\sum_j\lambda_j.
\]
In particular $D_T$ has genus zero.  The elementary trace-class ideal,
exterior-power, and compact-operator spectral theorems are the analytic
background here; no arithmetic hypothesis is used.

\begin{proof}
The singular values of $\bigwedge^nT$ are the products
$s_{j_1}\cdots s_{j_n}$ with $j_1<\cdots<j_n$.  Thus
\[
 |\operatorname{Tr}(\bigwedge^nT)|
 \leq\|\bigwedge^nT\|_1
 =\sum_{j_1<\cdots<j_n}s_{j_1}\cdots s_{j_n}
 =\alpha_n(|T|).
\]
The last quantity is at most $\|T\|_1^n/n!$, proving entire convergence
and the two domination inequalities.  Diagonalizing the positive compact
operator $|T|$ gives directly
$D_{|T|}(z)=\prod_j(1+zs_j)$, a genus-zero product.

For any finite collection of nonzero eigenvalues, include their full
generalized eigenspaces and let $E$ be their finite-dimensional invariant
sum.  A unitary triangularization of $T|_E$ gives an orthonormal basis
$e_1,\ldots,e_d$ with diagonal entries the selected eigenvalues.  There is
a diagonal unitary on $E$, extended by zero on $E^\perp$, whose product
with $T$ has trace $\sum_{j=1}^d|\lambda_j|$.  Its operator norm is at
most one.  The trace ideal inequality therefore bounds this sum by
$\|T\|_1$.  Increasing the finite collection proves the summability
assertion.

For completeness, the determinant zero information needed next does not
assume the desired product formula.  The exterior-power definition is
continuous under trace-norm approximation, uniformly for $z$ in compact
sets.  For fixed $n$, this follows by telescoping the tensor product
$A^{\otimes n}-B^{\otimes n}$ and restricting by the bounded
antisymmetrizing projection.  The bounds
$|\alpha_n(A)|\leq\|A\|_1^n/n!$ make the tails uniformly small on
bounded trace-norm sets, proving the asserted locally uniform continuity.
Finite-rank determinant identities consequently imply similarity
invariance and
$\det((1+A)(1+B))=\det(1+A)\det(1+B)$ for trace-class $A,B$.
In particular, if $1+A$ is invertible, its inverse has the form $1+B$
with $B$ trace class, so $\det(1+A)\ne0$.
For a nonzero eigenvalue $\lambda$, the compact spectral theorem gives
the decomposition by its finite-dimensional Riesz subspace and its
invariant closed complement.  Similarity invariance and the direct-sum
identity factor $D_T$ into the ordinary determinant on the Riesz subspace
and the determinant on the complement.  At $z=-1/\lambda$, the latter
factor is nonzero, and the former has zero of order exactly the algebraic
multiplicity of $\lambda$.  At every other point where $1+zT$ is
invertible the determinant is nonzero.  These are all its zeros by the
compact spectral theorem.

Consequently the nonzero zeros of $D_T$ have summable reciprocal moduli.
Its coefficients are dominated by those of the genus-zero function
$D_{|T|}$.  Finding F-SM-11 gives genus zero for $D_T$.  Its zero divisor
and $D_T(0)=1$ identify it as $\prod_j(1+z\lambda_j)$.  Normal convergence
and $\sum_j|\lambda_j|<\infty$ allow differentiation at zero, giving
$\operatorname{Tr}T=\sum_j\lambda_j$.
\end{proof}

This proves the classical Hilbert-space product and trace assertions
written in the deleted passage without using them to prove the
genus-zero lemma.  It should not be described as an RH-dependent solution
of an arithmetic question on those sheets.

\subsection{A fully normalized numerical form of Masullo's main identity}

We now separate a numerical consequence that can be completed from the
adopted Dirichlet theorem from the geometric deconstruction not supplied
in Masullo's Section~12.3.  No kernel-decomposition claim is needed for
the following result.

\paragraph{Finding F-SM-13 (Masullo 12.5.1.1 over $\mathbb Q$, as a numerical distribution).}
Adopt Alp\"oge's Dirichlet zero-location theorem.  Let $\chi$ be a
primitive Dirichlet character of conductor $q$, including the trivial
character of conductor one; write $\chi(-1)=(-1)^a$ with $a\in\{0,1\}$.
Let $h\in C_c^\infty(\mathbb R_{>0})$ satisfy $h(1)=0$, and put
\[
 \widehat h(z)=\int_0^\infty h(t)t^z\,\frac{dt}{t},
 \qquad\delta_\chi=\begin{cases}1,&q=1,\\0,&q>1.\end{cases}
\]
Define the following concrete linear functionals:
\[
 A_a(h)=\frac12\int_0^\infty h(t)t^{1/2}
 \left(\frac1{|1-t|}+\frac{(-1)^a}{1+t}\right)\frac{dt}{t}
 =\int_0^\infty
 \frac{e^{-(a+1/2)x}}{1-e^{-2x}}
 \bigl(h(e^x)+h(e^{-x})\bigr)\,dx,
\]
\[
 B_\chi(h)=\sum_{n\geq1}\frac{\Lambda(n)}{\sqrt n}
 \bigl(\chi(n)h(n)+\overline\chi(n)h(n^{-1})\bigr).
\]
The integrals are ordinary absolutely convergent integrals, and the prime
power sum is finite.  For every $\alpha>0$ one has
\begin{equation}\label{eq:SM-numeric-strip}
 A_a(h)+B_\chi(h)
 =\delta_\chi\bigl(\widehat h(1/2)+\widehat h(-1/2)\bigr)
 -\!\sum_{\substack{\rho:\,L(\rho,\chi)=0,\ 0<\Re\rho<1\\
                         |\Re\rho-1/2|<\alpha}}
       m_\chi(\rho)\widehat h(\rho-1/2).
\end{equation}
The zero sum is absolutely convergent.  Thus this is an equality of
continuous linear functionals on the closed test-function subspace
$\{h\in C_c^\infty(\mathbb R_{>0}):h(1)=0\}$, with its usual inductive
limit topology.

Under the canonical identification
$C_{\mathbb Q}\simeq\widehat{\mathbb Z}^{\times}\times\mathbb R_{>0}$,
take Haar measure $du\,dt/t$, where $du$ has total mass one, and set
\[
 \phi_{\chi,h}(u,t)=\widehat\chi(u)h(t),\qquad
 c_{\chi,s}(u,t)=\widehat\chi(u)^{-1}t^s.
\]
Here $\widehat\chi$ denotes the finite character on
$\widehat{\mathbb Z}^{\times}$ induced by $\chi$.
Here and above $L(s,\chi)$ denotes the finite Dirichlet series.  Masullo's
notation $L(c)$ includes the archimedean factor (Definition~3.2.1.2,
page~123).  With these conventions his finite factor satisfies
$L_f(c_{\chi,s})=L(s,\chi)$, and his completed function is
\[
 L_{\mathrm M}(c_{\chi,s})=
 \pi^{-(s+a)/2}\Gamma((s+a)/2)L(s,\chi).
\]
It has exactly the nontrivial zero divisor appearing above.  Moreover
\[
 M(\phi_{\chi,h})(c_{\chi,s}|\cdot|^{-1/2})
       =\widehat h(s-1/2).
\]
Equation~\eqref{eq:SM-numeric-strip} is exactly the $\chi$-component of
Masullo's Theorem~12.5.1.1, pages~519--520, for $F=\mathbb Q$ and
$I=[-\alpha,\alpha]$.  Finite Fourier decomposition consequently proves
that numerical theorem for every $\phi\in\mathcal D(C_{\mathbb Q})$
with $\phi(C_{\mathbb Q}^1)=0$, in Masullo's own Haar normalization.

\begin{proof}
We give the contour calculation as well as the local normalization checks.
The classical analytic inputs are the continuation and functional equation
of primitive Dirichlet $L$-functions, the order-one completed function,
and its zero count $N_\chi(T)=O_\chi(T\log(2+T))$.  These are background
theorems; this proof does not claim to re-prove the analytic continuation
or the functional equation.  They are the standard inputs behind the
explicit formula (Iwaniec and Kowalski, 2004, Sections~4.6, 5.3, and~5.5).

Let
\[
 \mathcal L_\chi(s)=(q/\pi)^{(s+a)/2}
            \Gamma((s+a)/2)L(s,\chi),\qquad H(s)=\widehat h(s-1/2).
\]
The completed function has precisely the nontrivial zeros, and has simple
poles at $0,1$ exactly when $\delta_\chi=1$.  Its functional equation is
$\mathcal L_\chi(s)=\varepsilon_\chi\mathcal L_{\bar\chi}(1-s)$,
where the nonzero constant $\varepsilon_\chi$ disappears after logarithmic
differentiation.  Integration by parts after $t=e^x$ shows that $H(s)$
decays faster than any inverse power of $|\Im s|$, uniformly on every
fixed vertical strip.  This and the zero count prove absolute convergence
of $\sum_\rho m_\chi(\rho)H(\rho)$.

Integrate $(\mathcal L'_\chi/\mathcal L_\chi)(s)H(s)$ around rectangles
with vertical sides $c=3/2$ and $1-c=-1/2$.  Horizontal sides can be chosen
to contribute zero in the limit.  One explicit justification uses the
Hadamard logarithmic derivative of
$[s(s-1)]^{\delta_\chi}\mathcal L_\chi(s)$: in each interval $[j,j+1]$
choose a height at distance at least $c_0/(j\log(2+j))$ from the ordinates
of all zeros of modulus at most $2j+4$; the zero count leaves such a height
for a sufficiently small fixed $c_0>0$.  Do the same below the real axis.
The nearby terms of the logarithmic derivative are then polynomially
bounded in $j$, and its tail is bounded using
$\sum_{|\rho|>R}|\rho|^{-2}=O_\chi(\log(2+R)/R)$.
The rapid decay of $H$ kills both horizontal integrals.

The residue theorem and the functional equation therefore give
\begin{align*}
 \sum_\rho m_\chi(\rho)H(\rho)-\delta_\chi(H(0)+H(1))
 =\frac1{2\pi i}\int_{(c)}\biggl[
 \frac{\mathcal L'_\chi}{\mathcal L_\chi}(s)H(s)
 +\frac{\mathcal L'_{\bar\chi}}{\mathcal L_{\bar\chi}}(s)H(1-s)
 \biggr]\,ds.
\end{align*}
The constant $\tfrac12\log(q/\pi)$ in each logarithmic derivative
contributes $\log(q/\pi)h(1)=0$, by Mellin inversion.  On this line the
Dirichlet series for $L'/L$ converges absolutely.  Mellin inversion gives
\[
 \frac1{2\pi i}\int_{(c)}H(s)n^{-s}\,ds=n^{-1/2}h(n),\qquad
 \frac1{2\pi i}\int_{(c)}H(1-s)n^{-s}\,ds=n^{-1/2}h(n^{-1}).
\]
Thus the two $L'/L$ terms contribute $-B_\chi(h)$.

For the gamma terms use, for $\Re z>0$,
\[
 \psi(z)=-\gamma+\int_0^\infty
                 \frac{e^{-v}-e^{-zv}}{1-e^{-v}}\,dv.
\]
First truncate the $v$-integral to $[\epsilon,R]$.  All constant terms
pair with $h(1)=0$.  Fubini and Mellin inversion give for the remaining
term
\[
 -\frac12\int_\epsilon^R
 \frac{e^{-(a/2+1/4)v}}{1-e^{-v}}
       \bigl(h(e^{v/2})+h(e^{-v/2})\bigr)\,dv.
\]
The truncated digamma integrals on the line $\Re s=c$ are bounded by
$C(1+|\Im s|)$ uniformly in $\epsilon,R$: near zero use the mean-value
bound on $e^{-v}-e^{-(s+a)v/2}$, and away from zero use exponential
decay.  This justifies passage to the limit against rapidly decreasing
$H(s)+H(1-s)$.  The resulting real integral is absolutely convergent,
because $h(1)=0$ and $h$ is smooth, and equals $-A_a(h)$ after $v=2x$.
The expression for $A_a$ in the statement follows by splitting its
$t$-integral at one and substituting $t=e^x,e^{-x}$; the elementary
identity
\[
 \frac12\left(\frac1{1-r}+\frac{(-1)^a}{1+r}\right)
       =\frac{r^a}{1-r^2},\qquad 0<r<1,
\]
gives exactly its displayed coefficient.  We have proved the full
explicit formula, with the sum over all nontrivial zeros.
The adopted Dirichlet theorem puts every such zero on $\Re\rho=1/2$.
For every $\alpha>0$ its restriction to the indicated open strip therefore
omits no zero, proving~\eqref{eq:SM-numeric-strip}.

We now identify the idelic terms.  For an idele $b=(b_\infty,(b_p)_p)$,
put $r=\operatorname{sgn}(b_\infty)\prod_p p^{v_p(b_p)}\in\mathbb Q^\times$.
Its class has the unique representative
$(t,(u_p)_p)=(b_\infty/r,(b_p/r)_p)$ with $t>0$ and all $u_p$ units;
moreover $t=|b|$.  This proves the stated group identification.
Masullo's Definitions~A.3.2.1--A.3.2.2, page~560, give
$\mu_\infty=dx/|x|$, $\mu_p(\mathbb Z_p^\times)=1$, and
\[
 \mu_1(C_{\mathbb Q}^1)=\frac{2^{r_1}(2\pi)^{r_2}h_{\mathbb Q}R_{\mathbb Q}}
 {\sqrt{|d_{\mathbb Q}|}w_{\mathbb Q}}=\frac{2}{2}=1.
\]
Hence the quotient measure is precisely $du\,dt/t$, without an omitted
volume factor.

For $p\nmid q$, the local idele $b_p=p^kz$ with
$z\in\mathbb Z_p^\times$ has
$\phi_{\chi,h}(b_p)=\chi(p)^{-k}h(p^{-k})$.  The $k=0$ term is zero;
on $k>0$ one has $|1-b_p|_p=1$, and on $k=-m<0$ one has
$|1-b_p|_p=p^m$.  Each shell has multiplicative measure one.  Its local
integral, multiplied by $\log p$, is therefore
\[
 (\log p)\sum_{m\geq1}p^{-m/2}
 \bigl(\chi(p)^m h(p^m)+\overline\chi(p)^m h(p^{-m})\bigr).
\]
For $p\mid q$, primitivity makes the local unit character nontrivial;
unit averaging annihilates each shell with $k\ne0$, while $k=0$ is again
identically zero since $h(1)=0$.  This gives exactly $B_\chi(h)$.
At the real place,
$\phi_{\chi,h}(x)=\operatorname{sgn}(x)^a h(|x|)$; splitting
$\mathbb R^\times$ into its two components and retaining Masullo's
factor $1/2$ gives exactly $A_a(h)$.  The potential singularity at $x=1$
is integrable because $h(1)=0$, so Weil's principal value agrees with the
ordinary integral here.

Orthogonality on the mass-one compact group gives the asserted Mellin
transform, gives zero for all other finite-character components, and gives
$M(\phi_{\chi,h})(|\cdot|^{\pm1/2})=
\delta_\chi\widehat h(\pm1/2)$.  Every Bruhat test function on
$C_{\mathbb Q}$ factors through a finite quotient of its compact factor
and is smooth with compact support on its positive-real factor.  Its
finite Fourier expansion is therefore
$\phi(u,t)=\sum_\chi\widehat\chi(u)h_\chi(t)$, with each character
indexed once by its primitive conductor.  The condition
$\phi(u,1)=0$ for all $u$ is equivalent to $h_\chi(1)=0$ for every
component.  Adding the finitely many component identities proves the
stated assertion for all $\phi$.  Continuity follows on each fixed
compact-support/finite-quotient test space from the local integral bounds,
the finite prime-power sum, and the rapid-decay estimate for the absolutely
convergent zero sum; these estimates are compatible with the inductive
limit topology.
\end{proof}

\paragraph{What is completed and what remains geometric.}
Finding F-SM-13 identifies the numerical distribution in Masullo's
restricted-strip formula for every allowed test function over
$\mathbb Q$, and gives all conductor, parity, Haar, and Mellin
normalizations.  The conductor term vanishes here for the precise reason
$h(1)=0$; it has not been silently discarded in a formula for unrestricted
test functions.  The statement indexes each idelic character by its
primitive conductor; treating an imprimitive Dirichlet series as though it
were that same completed idelic $L$-function would introduce extra deleted
Euler factors.

This numerical identity does not prove Masullo's geometric assertion
\[
 \int_{C_F}k(x,x)\,d\mu(x)
 =\int_{C_F}k_S(x,x)\,d\mu(x)
   +\sum_{p\nmid S}\int_{C_F}k_p(x,x)\,d\mu(x),
\]
nor identify the kernels in the missing Section~12.3, justify an exchange
of the local sum with diagonal restriction/integration, or establish its
Sobolev, nuclearity, and quotient compatibility hypotheses.  Equality of
two already identified numerical distributions does not establish that
a proposed kernel construction realizes them.  It does, however, complete
the exact numerical statement of Theorem~12.5.1.1 for $\mathbb Q$ from
the adopted Alp\"oge theorem and classical analytic inputs.  No extension
to arbitrary number fields or their Hecke characters is made by this
argument.

% END source_masullo_numeric.tex

% BEGIN source_masullo_kernel.tex
\section{The actual Masullo kernel: a diagonal bound, an indentation defect, and a concrete repair}
\label{sec:SM-actual-kernel}

We use the operators actually named in the manuscript:
\[
 H=W_I^s(C_F)=L^2(C_F,w_{I,s}\,d\mu),\quad
 S_\phi=A(\phi)P_a^{\mathrm{spec}},\quad
 T_{\rho,\phi}=M_\rho A(\phi)P_a^{\mathrm{geom}}.
\]
The projections are Definition~10.4.1.1, p.~456; nuclearity for
$S_\phi$ is asserted in Theorem~10.4.2.8, pp.~463--464.
Definition~10.4.3.7, pp.~473--474, gives
$k^{\mathrm{geom}}_\phi(x,y)=\phi(xy^{-1})-k^{\mathrm{spec}}_\phi(x,y)$.
Equation~(10.4.3.48), p.~473, assigns to $T_{\rho,\phi}$ the
weighted-measure kernel
$w_{I,s}(y)^{-1}\rho(x)k^{\mathrm{geom}}_\phi(x,y)$.
The first result is conditional on the indicated actual $S_\phi$
being trace class and on this kernel identification.  It does not
replace that operator by a separately constructed zero spectrum or
recheck the entire earlier nuclearity argument.

\subsection{An explicit diagonal convergence estimate}

\paragraph{Finding F-SM-18 (trace-density domination and valid-indentation shell convergence).}
Let $\nu=w\mu$, where $\mu$ is a positive Radon measure and $w>0$.
Every trace-class operator $S$ on $L^2(X,\nu)$ has a unique
$d_S^\nu\in L^1(X,\nu)$ with
\[
 \operatorname{Tr}(M_fS)=\int_X f\,d_S^\nu\,d\nu
 \quad(f\in L^\infty),\qquad \|d_S^\nu\|_1\leq\|S\|_1.
 \tag{SM.18.1}
\]
Its density relative to $\mu$ is $d_S^\mu=w\,d_S^\nu$, with
the same $L^1$ norm.  For the actual $S_\phi$, if $\phi(1)=0$,
the trace-density geometric diagonal is $d=-d_{S_\phi}^\mu$, so
\[
 \|d\|_{L^1(\mu)}\leq\|S_\phi\|_1.
 \tag{SM.18.2}
\]
When the manuscript's averaging procedure gives the trace
diagonal of its spectral kernel, this agrees with its smoothened
geometric diagonal, since the convolution kernel has continuous
diagonal $\phi(1)$.  Its cutoff diagonal is $\rho d$;
for a sequence $|\rho_j|\leq1$ with $\rho_j\to1$ almost
everywhere, $\rho_jd\to d$ in $L^1$.

Suppose now $X=G/H$, with $G$ second countable, locally compact
abelian and $H$ closed and discrete with counting measure.
Take $\mu$ to be the quotient Haar measure and $D$ a measurable
fundamental domain carrying its corresponding measure $\mu_D$.
For a nonnegative measurable normalized indentation $b$, put
\[
 \sum_{r\in H}b(r^{-1}x)=1,\qquad
 \omega_b(x,h)=\sum_{r\in H}b(r^{-1}x)b(r^{-1}h^{-1}x).
 \tag{SM.18.3}
\]
Let measurable subsets $G_N$ increase to $G$, set
$\alpha_N=1_{G_N}$, and define on $D$
\[
 \beta_N(x)=\sum_{h\in H}\omega_b(x,h)
                       (1-\alpha_N(h^{-1}x)).
 \tag{SM.18.4}
\]
Then $0\leq\beta_N\leq1$, $\beta_N\to0$, and the omitted
absolute diagonal mass is bounded, uniformly for $|\rho|\leq1$, by
\[
 \int_D\sum_h|\rho(\pi x)d(\pi x)\omega_b(x,h)|
                 (1-\alpha_N(h^{-1}x))\,d\mu_D(x)
 \leq \int_D|d(\pi x)|\beta_N(x)\,d\mu_D(x)\longrightarrow0.
 \tag{SM.18.5}
\]
The initial-stratum and successive-shell diagonal densities
are absolutely summable in $L^1$, with sum of their norms equal
to $\|d\|_1$.  Thus diagonal integration and bounded cutoff
limits commute with this shell sum.  This conclusion does not
identify the pieces with independently specified local operators.

\begin{proof}
Write a singular-value expansion
$S=\sum_j\sigma_j u_j\otimes v_j^*$ with unit vectors $u_j,v_j$.
The series $\sum_j\sigma_j u_j\overline{v_j}$ converges in $L^1$,
since $\|u_j\overline{v_j}\|_1\leq1$ and
$\sum_j\sigma_j=\|S\|_1$.  Its sum is $d_S^\nu$.
The rank-one trace identity, followed by trace-norm and $L^1$
limits, proves (SM.18.1).  Testing against all bounded functions
gives uniqueness; the same estimate proves contractivity of
$S\mapsto d_S^\nu$.  Change of measure gives $d_S^\mu$.
The defining difference for the geometric kernel, with $\phi(1)=0$,
gives (SM.18.2), and dominated convergence gives the cutoff limit.

The discrete subgroup $H$ is countable.  All terms in $\omega_b$
are nonnegative, and Tonelli gives
\[
 \sum_h\omega_b(x,h)
 =\sum_r b(r^{-1}x)\sum_h b(r^{-1}h^{-1}x)=1.
 \tag{SM.18.6}
\]
Consequently $\beta_N$ lies in $[0,1]$.  For fixed $x,h$,
$\alpha_N(h^{-1}x)\to1$; dominated convergence with the
probability weights $\omega_b(x,h)$ gives $\beta_N(x)\to0$.
Another dominated convergence argument against $|d|$ proves
(SM.18.5), including uniformity in $\rho$.

Explicitly, let $c_0=\sum_h\omega_b(x,h)\alpha_0(h^{-1}x)$
and $c_j=\sum_h\omega_b(x,h)(\alpha_j-\alpha_{j-1})(h^{-1}x)$
for $j\geq1$.  Then $c_j\geq0$ and $\sum_jc_j=1$.
The shell densities are $d(\pi x)c_j(x)$, and Tonelli yields
$\sum_j\|dc_j\|_1=\|d\|_1$.  This proves every stated interchange.
\end{proof}

For a common-Hilbert-space realization of a proposed trace-class
operator remainder $E_N$, an equivalent trace-norm target is
\[
 \|d_{E_N}\|_1\leq\|E_N\|_1
   =\inf_{E_N=U_NV_N}\|U_N\|_2\|V_N\|_2\longrightarrow0.
 \tag{SM.18.7}
\]
The infimum runs over Hilbert--Schmidt factorizations.
The inequality $\|UV\|_1\leq\|U\|_2\|V\|_2$ is
Cauchy--Schwarz with trace duality.  Conversely, for
$E=W|E|$, take $U=W|E|^{1/2}$ and $V=|E|^{1/2}$;
both squared Hilbert--Schmidt norms equal $\operatorname{Tr}|E|$.
This proves the infimum formula.  Its vanishing is equivalent
to trace-norm convergence, and is sufficient for diagonal
convergence; it is not necessary for a scalar trace equality alone.
Hilbert--Schmidt convergence alone is insufficient: $N^{-1}$
times a rank-$N$ orthogonal projection has Hilbert--Schmidt norm
$N^{-1/2}$, trace norm one, and trace one.

Theorem~11.1.2.1, pp.~480--483, states a cofinal intersection
of closed carving subspaces and an identification by special
theta functions.  It does not state a common-space realization
and remainder estimate (SM.18.7).  On the other hand,
(SM.18.5) solves diagonal summation once the actual spectral
trace-class input and a normalized indentation are available.
Identifying those diagonal pieces with the local kernels remains
an additional assertion in the unwritten Section~12.3.

\subsection{The displayed orbital indentation vanishes}

\paragraph{Finding F-SM-19 (literal formula defect and nonzero test).}
Interpret equation~(11.3.2.2), p.~495, as the scalar-valued
function it is called, with the ordinary Haar measure
$d\mu_1(g)$ of its adjacent equations:
\[
 U(x,y)=\sum_{\xi\in F^\times}\int_{E_F^1}
   1_{\tau(\mathbb R_{>0})}(\xi gx)
   1_{\tau(\mathbb R_{>0})}(\xi gy)\,d\mu_1(g).
 \tag{SM.19.1}
\]
For $F=\mathbb Q$ this is identically zero, is not a normalized
indentation, and cannot yield the stated orbital trace formula
(11.3.2.1) under this literal interpretation.

\begin{proof}
Definition~A.2.3.6, p.~549, equation~(A.2.3.14), defines
$E_F^1$ as an actual subset of $\mathbb A_F^{\times,1}$.
For $\mathbb Q$ it is
$E_{\mathbb Q}^1=\{1\}_\infty\times\widehat{\mathbb Z}^{\,\times}$.
Definitions~A.3.2.1--2, pp.~559--560, use ordinary Haar measure
here, of volume one, with counting measure assigned only to
$\mathbb Q^\times$.  This infinite compact group has atomless
Haar measure: a positive singleton mass $c$ would give mass
$Nc$ to $N$ distinct points for every $N$, a contradiction.

Fix $x$ and $\xi$.  If $\xi gx=\tau(t)$ for $g\in E_{\mathbb Q}^1$,
the product formula and $|\tau(t)|=t$ force $t=|x|$.
Thus $g=\xi^{-1}\tau(|x|)x^{-1}$ is the only possible value.
The first indicator in (SM.19.1) is supported on at most one
point of $E_{\mathbb Q}^1$.  Each integral is zero and the
countable sum is zero, for all $x,y$.

For a nonzero right side of the orbital formula, choose
$F=\mathbb Q$, $\chi=1$, $S=\{\infty\}$, $a=1/2$, $s=4$,
and any $\alpha>0$.  Then $\operatorname{Pic}(\mathbb Z)=1$
and all stated ramification assumptions hold.
Choose $h\in C_c^\infty((3/2,5/2))$ with $h(2)=1$ and set
$\phi(u,t)=h(t)$ on
$C_{\mathbb Q}\simeq\widehat{\mathbb Z}^{\,\times}\times\mathbb R_{>0}$.
This is a Bruhat test function, finite-unit invariant and
supported over norms greater than one.  For the finite uniformizer
idele $\pi_2$, $|\pi_2|=1/2$, whence
\[
 \sum_{m\in\mathbb Z}\phi(\pi_2^m)
       =\sum_mh(2^{-m})=1.
\]
For the shell containing the place $2$, the stated right side is
$\log2/\mu_2(\mathbb Z_2^\times)>0$; the stated normalization
makes this $\log2$.  The literal indented kernel on the left
is zero, including after cutoffs, smoothening, and any integration
in the remaining norm variable.  The claimed equality therefore
fails with formula (SM.19.1) interpreted as printed.
\end{proof}

The introductory discussion, pp.~21--22, explicitly speaks of
scalar-valued measurable functions and allows indicators of
fundamental domains.  It introduces no Dirac interpretation
for the indicators in (SM.19.1).  A distributional replacement
or independently normalized disintegration would be a new
formula requiring its own operator and diagonal estimates.
This finding does not rule out such a corrected realization.

\subsection{A valid rational indentation and compact averaging}

\paragraph{Finding F-SM-20 (two explicit repairs and the remaining identification).}
For $\mathbb Q$, put
$Z=\mathbb R_{>0}\times\widehat{\mathbb Z}^{\,\times}
\subset\mathbb A_{\mathbb Q}^\times$ and $b=1_Z$.
Then $b$ is a continuous normalized indentation and
$b^\Delta(x,y)=\sum_\xi b(\xi^{-1}x)b(\xi^{-1}y)$ is valid.
For $x\in Z$ and $q\in\mathbb Q^\times$,
\[
 b^\Delta(x,qx)=1_{\{q=1\}}.
 \tag{SM.20.1}
\]
Hence this lift preserves the global diagonal but does not itself
recover the nontrivial prime-orbital formulas.

There is also a legitimate compact averaging for the actual operators.
Write $K=\widehat{\mathbb Z}^{\,\times}$ with probability Haar
measure, and assume $\phi(u,t)=h(t)$, $\rho(u,t)=\rho(t)$.
For $H_0=L^2(\mathbb R_{>0},w(t)\,dt/t)$, put
\[
 (Jf)(u,t)=f(t),\quad (J^*F)(t)=\int_KF(u,t)\,du,\quad E=JJ^*.
\]
Then $E$ is the orthogonal projection onto $K$-invariants, and
\[
 T_{\rho,\phi}=ET_{\rho,\phi}E,\qquad S_\phi=ES_\phi E.
 \tag{SM.20.2}
\]
For either operator $T$ in (SM.20.2), if $T$ is trace class,
\[
 \operatorname{Tr}_H T=\operatorname{Tr}_{H_0}(J^*TJ),
 \qquad \|T\|_1=\|J^*TJ\|_1.
 \tag{SM.20.3}
\]
If $K_T$ is its weighted Hilbert--Schmidt kernel, the compressed
kernel is the actual compact average
\[
 K_0(t,s)=\int_K\int_KK_T((u,t),(v,s))\,du\,dv,\qquad
 \|K_0\|_2\leq\|K_T\|_2.
 \tag{SM.20.4}
\]

\begin{proof}
For an idele $x$, let $r=\prod_p p^{v_p(x_p)}>0$, a finite
product, and take $q=\operatorname{sgn}(x_\infty)r$.
Then $q^{-1}x\in Z$.  Two such rational multipliers differ
by a rational number with all finite valuations zero and
positive sign, hence by one.  Thus the translates of $Z$
are disjoint and cover the ideles.  This subset is open
and closed and maps homeomorphically onto $C_{\mathbb Q}$.
Its intersection with the inverse image of a compact subset
of the quotient is compact.  Therefore $b$ has the stated
indentation properties and $\sum_qb(q^{-1}x)=1$.
For $x\in Z$ the first factor in $b^\Delta(x,qx)$ forces
$\xi=1$, and the second then forces $q=1$, proving (SM.20.1).

In particular, for the initial subgroup
$G_0=\mathbb R^\times\times\widehat{\mathbb Z}^{\,\times}$,
all outside-shell diagonal pieces from this lift are zero:
the only surviving $q$ is one and $x\in Z\subset G_0$.
The part labelled $S$ still contains the original global
spectral diagonal.  Its identification with a separately
defined $S$-local Paley--Wiener operator has not followed.

For averaging, Fubini and the volume-one normalization give
the displayed adjoint and $J^*J=1$, so $E$ is the claimed
orthogonal projection.  The actual weight $w_{I,s}$ depends
only on the norm.  Thus $K$ acts unitarily.
The closed geometric theta-function subspace of
Definition~10.4.1.1 is $C_F$-invariant, as explicitly used
in Lemma~10.4.1.2, p.~456.  That subspace and its orthogonal
complement are therefore $K$-invariant; their projections
commute with every $K$-translation and hence with $E$.
Convolution by $\phi$ satisfies
$EA(\phi)=A(\phi)=A(\phi)E$, by integrating the compact
variable, and $M_\rho$ commutes with $E$.  This proves
(SM.20.2).

Relative to $H=JH_0\oplus(JH_0)^\perp$, such a $T$ is
$(J^*TJ)\oplus0$.  The traces and singular values of this
block operator prove (SM.20.3).  Fubini first for simple
tensors and then $L^2$ density proves (SM.20.4); Jensen's
inequality for probability Haar measure on $K\times K$
gives its norm bound.
\end{proof}

These repairs permit a genuine trace calculation on the norm
coordinate for the indicated invariant part.  They do not
identify the compressed operator with the $S$-local operator
in Theorem~12.4.1.1, pp.~514--519.  That needs a specified
intertwiner carrying the actual kernels, projections and cutoff
limits to those local operators.  The zero indentation in
(SM.19.1) supplies none, and the closed-subspace carving
statement supplies no trace-norm estimate for that identification.

\subsection{Two precise obstructions to naive replacements}

\paragraph{Finding F-SM-21 (the unchanged archimedean local kernel misses a nonzero prime term).}
There exists a real $h\in C_c^\infty((3/2,5/2))$ such that
$h(2)=1$ and
\[
 P(h):=\widehat h(1/2)+\widehat h(-1/2)=0,\qquad A_0(h)=0,
 \tag{SM.21.1}
\]
where $A_0$ is the real, even archimedean distribution in F-SM-13.
Under the adopted Dirichlet theorem, its full zero distribution is
\[
 -\sum_\rho m(\rho)\widehat h(\rho-1/2)
       =\frac{\log2}{\sqrt2}>0.
 \tag{SM.21.2}
\]
Consequently the following two manuscript identifications cannot
be joined by simply replacing the zero indentation with $1_Z$
and leaving the $S=\{\infty\}$ local operator unchanged:
the global spectral trace identification of
Theorem~10.4.2.8 and equation~(12.5.1.4), and the local
Paley--Wiener trace limit of Theorem~12.4.1.1.
Under these explicitly stated operator inputs, the former
geometric cutoff trace tends to $\log2/\sqrt2$, whereas the
latter trace tends to zero, including the combined removed
moment correction $P(h)$.

More precisely, if trace-preserving transports put those
trace-class cutoffs on a common Hilbert space, their difference
$E_\lambda$ must satisfy
\[
 \liminf_{\lambda\to\infty}\|E_\lambda\|_1
       \geq \frac{\log2}{\sqrt2}.
 \tag{SM.21.3}
\]
This rules out that unchanged-kernel identification; it does not
rule out an assembly with the required nonzero prime contribution.

\begin{proof}
On the indicated interval, write
\[
 P(h)=\int h(t)p(t)\,dt,\quad
 A_0(h)=\int h(t)a(t)\,dt,\quad
 p(t)=\frac{t+1}{t^{3/2}},\quad
 a(t)=\frac{\sqrt t}{t^2-1}.
\]
The ratio
$R(t)=a(t)/p(t)=t^2/((t-1)(t+1)^2)$ is strictly decreasing for
$t>1$: its logarithmic derivative is
$(-t^2+t-2)/(t(t^2-1))<0$.
Choose nonzero nonnegative bumps $v_1,v_2$ supported,
respectively, in $(1.55,1.7)$ and $(2.3,2.45)$.
The quotients $A_0(v_i)/P(v_i)$ are weighted averages of
$R$ over these disjoint ordered intervals, so the first
is strictly larger than the second.  The matrix with
columns $(P(v_i),A_0(v_i))$ is therefore invertible.
Choose a real bump $u$ supported in $(1.9,2.1)$ with
$u(2)=1$, solve that two-by-two system for real $c_1,c_2$,
and set $h=u-c_1v_1-c_2v_2$.  This proves (SM.21.1) and
$h(2)=1$.

The only nonzero term in the finite distribution of F-SM-13 is
$\Lambda(2)h(2)/\sqrt2=\log2/\sqrt2$; the reciprocal terms
are outside the support.  Its explicit formula reads
$A_0(h)+B_1(h)=P(h)-\sum_\rho m(\rho)\widehat h(\rho-1/2)$,
so (SM.21.2) follows.  All zero sums converge absolutely.
The adopted zero-location input makes the same statement
hold in every positive-width strip used by the manuscript.

The source's global operator identification is exactly the
left side of (SM.21.2).  For $S=\{\infty\}$, its stated
local trace limit is $A_0(h)$, since $h(1)=0$ removes the
cutoff logarithm.  The combined moment correction also
vanishes by (SM.21.1).  Thus their trace difference tends
to $\log2/\sqrt2$.  The elementary inequality
$|\operatorname{Tr}E_\lambda|\leq\|E_\lambda\|_1$
proves (SM.21.3).
\end{proof}

\paragraph{Finding F-SM-22 (multiplying by a compact diagonal Dirac kernel destroys compactness).}
For the actual compact-invariant operators of F-SM-20, write
$T_{\rho,\phi}=E_K\otimes B_\rho$ on
$L^2(K)\otimes H_0$, where $E_K$ projects onto constants
and $B_\rho=J^*T_{\rho,\phi}J$.
If its radial kernel is $K_0(t,s)$, the proposed
multiplicative replacement
\[
 K_0(t,s)\longmapsto \delta_K(uv^{-1})K_0(t,s)
 \tag{SM.22.1}
\]
defines $1_{L^2(K)}\otimes B_\rho$ on tensor test functions.
If $B_\rho\ne0$, this operator is not compact and hence
is not trace class.  Under the operator inputs of F-SM-21,
$B_\rho\ne0$ for all sufficiently advanced cutoffs in that
example.  Thus (SM.22.1) cannot serve as a trace-class repair
of the ordinary-function formula.

\begin{proof}
The normalized compact diagonal distribution acts by
$\int_K\delta_K(uv^{-1})f(v)\,dv=f(u)$, so integrating
the product kernel gives $1\otimes B_\rho$ on simple
tensors and then on the Hilbert completion.
The group $K=\widehat{\mathbb Z}^{\,\times}$ has infinitely
many distinct continuous characters, forming an infinite
orthonormal set $(\chi_j)$ in $L^2(K)$.
Choose a unit vector $v\in H_0$ with $B_\rho v\ne0$.
The unit vectors $\chi_j\otimes v$ have images
$\chi_j\otimes B_\rho v$, which are mutually orthogonal
and have the same positive norm.  Their images have no
convergent subsequence.  The operator is not compact.
For the example in F-SM-21, its asserted global trace
limit is positive, so the trace-class $B_\rho$ has
nonzero trace and is nonzero for all sufficiently
advanced cutoffs.  This proves the final assertion.
\end{proof}

This is an obstruction to the particular \emph{multiplicative}
Dirac interpretation.  A subsequent compact-isotypic
projection changes the operator again: compression by $E_K$
restores $E_K\otimes B_\rho$.  Such a repair is possible
as an operation, but it returns to the global operator
and still does not supply the missing local identification.

\paragraph{Primary source and verification scope.}
Masullo, A. (2026). \emph{The music of the primes}
[supplied manuscript].
\url{https://math216.wordpress.com/wp-content/uploads/2026/11/the-music-of-the-primes.pdf}.
Additional exact locators: pp.~21--22; Lemma~12.2.1.1,
pp.~504--505; Proposition~12.2.2.2, pp.~508--510;
the missing Section~12.3, p.~514;
Theorem~11.3.2.1 and equation~(11.3.2.2), p.~495;
Definition~A.2.3.6, p.~549; Definitions~A.3.2.1--2,
pp.~559--560.  The p.~495 formula was inspected visually.
F-SM-18--22 prove the statements written here, without
certifying the earlier full trace-class construction or the
absent local-operator identification.

% END source_masullo_kernel.tex

% BEGIN source_K_bridge.tex
\section{Integral $K$-groups: concrete known consequences and the missing export}
\label{sec:source-K-bridge}

The final footnote in Alp\"oge's four-page remark suggests consequences for
$K$-groups of $\mathbb Z$, but the supplied declarations do not identify a
degree, a torsion prime, an eigenspace, or a comparison map.  There are
nevertheless precise existing integral results with which such an export
can be compared.  We keep ordinary integral $K$-groups, their reductions
modulo $p$, and their $p$-primary subgroups distinct.

\subsection{The exact cyclotomic comparison}

\paragraph{Finding F-SM-14 (the mod-$p$ comparison and its vanishing consequence).}
Let $p$ be an odd prime, let $n\geq1$, put
$F_p=\mathbb Q(\zeta_p)$, and let $C_p$ be the $p$-primary subgroup of
$\operatorname{Cl}(F_p)$.  Let $\omega$ be the Teichm\"uller character of
$\Delta_p=\operatorname{Gal}(F_p/\mathbb Q)$.  There is an isomorphism of
$\mathbb F_p$-vector spaces
\begin{equation}\label{eq:SM-Kclass}
 K_{2n}(\mathbb Z)/p
 \simeq H^2_{\text{ét}}
       (\mathbb Z[1/p],\mu_p^{\otimes(n+1)})
 \simeq (C_p/pC_p)(\omega^{-n}).
\end{equation}
Here the last expression denotes a character eigenspace.  In particular,
\[
 K_{2n}(\mathbb Z)\{p\}=0
 \quad\Longleftrightarrow\quad C_p(\omega^{-n})=0.
\]
The comparison with \emph{mod-$p$} \'etale cohomology in the first
isomorphism is a known theorem, not a consequence of RH: see Weibel
(2013), Chapter~VI, Corollary~8.3, chapter page~48, following
Theorem~8.2, chapter page~47.  The following proof spells out its
class-group specialization and the vanishing equivalence.

\begin{proof}
Let $R=\mathbb Z[\zeta_p,1/p]$.  The map
$\operatorname{Spec}R\to\operatorname{Spec}\mathbb Z[1/p]$ is finite
\'etale Galois with group $\Delta_p$, of order $p-1$.  Taking invariants
on $\mathbb F_p[\Delta_p]$-modules is exact, by averaging; the
Hochschild--Serre spectral sequence therefore identifies cohomology over
$\mathbb Z[1/p]$ with the invariants of cohomology over $R$.

The Kummer sequence gives
\[
 0\longrightarrow\operatorname{Pic}(R)/p
 \longrightarrow H^2_{\text{ét}}(R,\mu_p)
 \longrightarrow\operatorname{Br}(R)[p]\longrightarrow0.
\]
There is exactly one prime of $F_p$ above $p$, generated by
$1-\zeta_p$.  Thus localizing at $p$ leaves the ideal class group
unchanged.  The field $F_p$ is totally imaginary, and the local-invariant
sequence for the Brauer group of its ring of $S$-integers has only this
one finite place; the sum-of-invariants map is then injective.  Hence
$\operatorname{Br}(R)[p]=0$.  The Kummer identification is
$\Delta_p$-equivariant and reads
\[
 H^2_{\text{ét}}(R,\mu_p)
       \simeq C_p/pC_p.
\]
Because $R$ contains $\mu_p$, the higher coefficient twist is constant
over $R$.  With its descent action retained, it gives
\[
 H^2_{\text{ét}}(R,\mu_p^{\otimes(n+1)})
 \simeq (C_p/pC_p)\otimes_{\mathbb F_p}\mu_p^{\otimes n}.
\]
The last tensor factor has character $\omega^n$ modulo $p$.
Its invariant subspace is exactly the $\omega^{-n}$ eigenspace of
$C_p/pC_p$, proving~\eqref{eq:SM-Kclass} from the cited comparison.

The finite $p$-group $C_p$ splits into its Teichm\"uller eigenspaces,
since $p-1$ is invertible in $\mathbb Z_p$; reduction modulo $p$ commutes
with the corresponding idempotents.  A finite $p$-group $A$ has
$A/pA=0$ if and only if $A=0$: if $A=pA$, iteration gives
$A=p^rA=0$ for $p^r$ annihilating $A$.  Apply this first to each class
component, and then to the $p$-primary component of the finite group
$K_{2n}(\mathbb Z)$.  The latter finiteness is the usual integral
$K$-theory finiteness theorem, also used in the cited comparison.
\end{proof}

The assertion is a comparison of reductions modulo $p$ and hence a
vanishing criterion.  It does not identify the entire finite
$p$-primary groups with one another.  In positive-rank odd degrees the
correct comparison uses $p$-adic coefficients/completion; an
uncompleted finitely generated $K$-group cannot surject onto a copy of
$\mathbb Z_p$.  This distinction is material when reading informal
Chern-map summaries.

\subsection{A complete known vanishing and its class-group consequence}

Dutour Sikiri\'c, Elbaz-Vincent, Kupers, and Martinet (2026),
Theorem~6.4, prove
\[
 K_8(\mathbb Z)=0.
\]
The theorem is on page~19 of the author-posted updated manuscript;
Corollary~7.1 on page~20 concludes
$C_p(\omega^{p-5})=0$ for every prime $p>3$.  The journal article is in
\emph{Journal of the Institute of Mathematics of Jussieu},
\emph{25}(1), 565--590, DOI~10.1017/S1474748025101394;
its online publication date is 24 November 2025.  This is a known result
independent of the adopted Alp\"oge input.

The final deduction has a short, exact audit trail.  Their
Proposition~6.1 puts $H_9(BQ;\mathbb Z)$ in the class of finite groups
whose prime divisors are at most seven, and Proposition~6.3 puts the
kernel of the relevant Hurewicz map in the analogous class for five.
Since $K_8(\mathbb Z)=\pi_9(BQ)$, its possible prime divisors are at
most seven.  The odd primes $3,5,7$ are regular (their relevant Bernoulli
numerators are those of $B_2=1/6$ and $B_4=-1/30$), so Weibel VI,
Proposition~10.5, chapter page~63, removes their torsion.  The
$2$-primary part is zero by Weibel VI, Corollary~9.8, chapter page~55.
This last input makes explicit the prime-two step that is not written in
the updated manuscript's very short proof of Theorem~6.4.  The Voronoi
cell enumerations and the preceding homology computations are carried
as published inputs; they have not been rerun here.

Applying F-SM-14 with $n=4$ gives
$(C_p/pC_p)(\omega^{-4})=0$, hence
$C_p(\omega^{-4})=0$.  As $-4\equiv p-5\pmod{p-1}$, this is the stated
cyclotomic class-group consequence.  The argument works at the level of
finite quotients without an incorrect uncompleted Chern map.  It also
does not need the assertion, as printed in the manuscript's Section~7,
that $H^1_{\text{ét}}(\mathbb Z[1/p],\mathbb Z_p(j))$
always vanishes when $p>j+1$.  For the positive twists relevant here,
that assertion requires even $j$: for odd $j>1$ this group has a
rank-one summand, as
Weibel VI, Exercise~8.3, chapter page~50, specifies.

\subsection{A density-one family with all degrees specified}

Chen et al. (2026), arXiv:2602.05090v1, Theorem~1.1 on page~2, prove
the following quantitative partial regularity statement.  For each
$\alpha>1/2$, let
\[
 M_\alpha(p)=\left\lfloor\frac{\sqrt p}{(\log p)^\alpha}\right\rfloor.
\]
The number of primes $p\leq X$ for which some even index
$2\leq2k\leq\min(M_\alpha(p),p-3)$ has
$p\mid\operatorname{num}(B_{2k})$ is
$O_\alpha(X/(\log X)^{2\alpha})$.  Their Corollary~1.4, page~3,
gives the corresponding vanishing of the even class-group components;
Corollary~1.7(4), page~5, gives the $K_{4k-2}$ consequence.
The formal verification reported by that paper concerns its main
elementary estimate; no Lean build of that archive has been performed
in this run, and the deep comparison inputs remain cited mathematics.

\paragraph{Finding F-SM-15 (both families of integral torsion consequences).}
For a fixed $\alpha>1/2$, outside
$O_\alpha(X/(\log X)^{2\alpha})$ odd primes $p\leq X$, the following
hold simultaneously for every integer $k$ with
$2\leq2k\leq\min(M_\alpha(p),p-3)$ and every integer $b\geq0$:
\begin{align*}
 K_{4k-2+2b(p-1)}(\mathbb Z)\{p\}&=0,\\
 K_{2p-2-4k+2b(p-1)}(\mathbb Z)\{p\}&=0.
\end{align*}
In particular, each fixed $\alpha$ gives a density-one set of primes
with these simultaneous vanishings.  The first line for $b=0$ is the
published consequence just cited; the second line makes explicit the
$4r$-degree consequence of its partial Vandiver components.

\begin{proof}
For a nonexceptional prime the partial regularity theorem gives
$p\nmid\operatorname{num}(B_{2k})$ throughout the range.  The
Herbrand--Ribet theorem identifies this with
$C_p(\omega^{1-2k})=0$; its exponent is also written $p-2k$ modulo
$p-1$.  Leopoldt reflection implies
$C_p(\omega^{2k})=0$.  These are exactly the two cited cyclotomic
inputs in Chen et al., equation~(4) and the deduction of
Corollary~1.4, page~3.

For the first line of the conclusion put
$n=2k-1+b(p-1)\geq1$ in F-SM-14.  Then
$-n\equiv1-2k\pmod{p-1}$.  For the second put
$n=p-1-2k+b(p-1)$.  The bound $2k\leq p-3$ ensures $n\geq2$, and
$-n\equiv2k\pmod{p-1}$.  Thus F-SM-14 proves both assertions, for
every $b$ at once.  They introduce no additional exceptional primes.
Finally, the prime number theorem gives
\[
 \frac{O_\alpha(X/(\log X)^{2\alpha})}{\pi(X)}
       =O_\alpha((\log X)^{1-2\alpha})\longrightarrow0.
\]
This proves the density assertion.
\end{proof}

The degrees in the second line are divisible by four because $p$ is
odd.  The degrees in the first line are congruent to two modulo four.
Both are precise $p$-primary assertions; neither line is an integral
vanishing assertion for an entire $K$-group at every prime.  Also, these
are pre-existing partial-regularity and comparison inputs, independent
of complex RH; their proof does not manufacture a missing exported
theorem from Alp\"oge's final footnote.

\subsection{The degree-one boundary and what an export must state}

\paragraph{Finding F-SM-16 (the elementary degree-one calculation).}
$K_1(\mathbb Z)\simeq\{\pm1\}\simeq\mathbb Z/2$.

\begin{proof}
Use the classical definition
$K_1(\mathbb Z)=\operatorname{GL}_\infty(\mathbb Z)^{\mathrm{ab}}$.
The determinant is a surjection onto $\{\pm1\}$.  Its kernel is
$\operatorname{SL}_\infty(\mathbb Z)$, which is generated by elementary
matrices: the Euclidean algorithm reduces the first column of a
unimodular matrix to a unit vector using integer row additions and signed
row interchanges; column additions clear the rest of its first row, and
induction applies to the remaining unimodular block.  Sign corrections
are allowed because
$\left(\begin{smallmatrix}0&1\\-1&0\end{smallmatrix}\right)
=e_{12}(1)e_{21}(-1)e_{12}(1)$ and its square is
$\operatorname{diag}(-1,-1)$.  In determinant one the final block also
has determinant one.  Stabilization supplies an extra index when needed.

For distinct indices $i,j,r$,
$e_{ij}(a)=[e_{ir}(a),e_{rj}(1)]$, so every elementary matrix is a
commutator in the stable group.  Conversely every commutator has
determinant one.  Thus the determinant kernel is precisely the
commutator subgroup, proving the asserted abelianization.
\end{proof}
The exact standard locator is Weibel III, Example~1.3.5, chapter
page~4, following Whitehead's Lemma~1.3.3.  This elementary nonzero
group already prevents a reading of the tentative footnote as blanket
vanishing in all positive degrees.

To make a further integral consequence of the supplied analytic method
reviewable, an export must specify its degree set, its coefficient or
torsion convention, its prime range, and the exact arithmetic map it
controls.  For example, an exported vanishing of the component
$C_p(\omega^{-n})$ would give the corresponding $p$-primary vanishing
in $K_{2n}(\mathbb Z)$ by F-SM-14.  To conclude an entire finite
$K_{2n}(\mathbb Z)$ is zero, one must cover every prime, including the
separate prime-two calculation.  Zero location for complex Dirichlet
$L$-functions alone supplies no such integral eigenspace or
divisibility statement in the supplied five declarations.

\subsection{Primary references and checked locators}

Chen, E., Cummins, C., Eltschig, B., Grubisic, D., Haller, L.,
Hong, L., Kurghinyan, A., Lau, K., Leather, H., Lee, S.,
Markosyan, A., Ono, K., Patel, M., Pendharkar, G., Rathi, V.,
Schneidman, A., Seeker, V., Sengupta, S., Sinha, I., \ldots\
Zhang, J. (2026).
\emph{Almost all primes are partially regular}
(arXiv:2602.05090v1).
\url{https://arxiv.org/abs/2602.05090}.
The used locators are Theorem~1.1, p.~2; equation~(4),
the reflection implication and Corollary~1.4, p.~3; and
Corollary~1.7(4), p.~5.  The natural-language proof of the main
quantitative theorem in Sections~2--3, pp.~6--10, was read throughout;
the formalization scope is stated in Section~4, pp.~10--11.

Dutour Sikiri\'c, M., Elbaz-Vincent, P., Kupers, A., \& Martinet, J.
(2026). Voronoi complexes in higher dimensions, cohomology of
$\operatorname{GL}_N(\mathbb Z)$ for $N\geq8$ and the triviality of
$K_8(\mathbb Z)$.
\emph{Journal of the Institute of Mathematics of Jussieu, 25}(1),
565--590.
\url{https://doi.org/10.1017/S1474748025101394}.
Updated author manuscript:
\url{https://jamartin.perso.math.cnrs.fr/Publications/DEKM.pdf}.
The used locators in this manuscript are Proposition~6.1, p.~18;
Proposition~6.3 and Theorem~6.4, p.~19; and Corollary~7.1, p.~20.
The complete final deductions were read; upstream Voronoi
computations were not reproduced.

Weibel, C. A. (2013).
\emph{The K-book: An introduction to algebraic K-theory}
(Graduate Studies in Mathematics, Vol.~145).
American Mathematical Society.
\url{https://doi.org/10.1090/gsm/145}.
Author's chapter files:
\url{https://sites.math.rutgers.edu/~weibel/Kbook/Kbook.III.pdf}
and
\url{https://sites.math.rutgers.edu/~weibel/Kbook/Kbook.VI.pdf}.
The exact locators are III, Example~1.3.5, chapter p.~4;
VI, equation~(8.1.1) and Theorem~8.2, chapter p.~47;
VI, Corollary~8.3, chapter p.~48; VI, Exercise~8.3, chapter p.~50;
VI, Corollary~9.8, chapter p.~55; VI, Proposition~10.5,
chapter p.~63; and VI, Theorem~10.9, chapter p.~65.
Chapter pagination differs from the printed volume; theorem
numbering is unchanged.  These primary statements and their
adjacent arguments were read directly.  The underlying motivic
comparison theorems are carried as established inputs.

% END source_K_bridge.tex
\clearpage
\part{Unresolved questions and unchecked source assertions}
\section{Current gates and approaches still in progress}
No mathematical assertion in this part is promoted to a theorem by
its inclusion. The executable question register is \nolinkurl{QUESTIONS.tsv};
the dated approach record is \nolinkurl{LEDGER.md}. A missing historical
proof package and an open mathematical conjecture are different statuses.

\begin{description}[style=nextline,leftmargin=0pt,labelwidth=0pt,labelsep=0pt]
\item[Cote 14, scans 25--38: original quotient identification.]
The generated quotient has full proofs in AQ-14-FORWARD and
AQ-14-CONVERSE. Original Definition1.1, Proposition1.2 and
Corollary2.3 now identify that quotient under the explicitly bracketed
projective-parameter convention. AQ-14-QUASIPROJECTIVE now removes the projective qualification
in characteristic zero for smooth quasiprojective parameters, by an
integral equality of the generated relations. AQ-14-NONPROPER now
proves the cohomological implication in every characteristic.
AQ-14-ALTERATION proves that the relation quotient in characteristic
$p>0$ is $p$-primary torsion, hence proves equality after inverting $p$
and on the prime-to-$p$ torsion in the resulting cycle quotients.
CC-ALB01--02 now remove that residual integral gate: the ordinary
Albanese kernel is divisible, while one fixed alteration degree kills
its entire image for a fixed correspondence. Thus the projective and
quasiprojective parameter relations agree integrally in every
characteristic. BR-14-SEP further reduces every smooth connected
separated finite-type parameter to a dense affine open, proving
the same integral relation without quasiprojectivity. Supplied scan38 is physical PDFpage38,
handwritten page37; both locators are retained.
\item[Cote 15, scans 100--106: full simple-motive classification.]
Weil-number occurrence is proved. The cited 2026 CM Hodge theorem,
with the inspected Milne implications, discharges the rational Tate and
standard-conjecture inputs for abelian varieties and their products.
AQ-ABELIAN-RECOGNITION proves full faithfulness after scalar extension
and Frobenius-polynomial recognition in the abelian--Artin numerical
category. CC-AG01--03 further classify its idempotent completion over
$\overline{\mathbf Q}_\ell$ by finite Weil-number gradings and
compute its diagonalizable group, including its Artin torsion quotient. CC-AG04--06 now also prove the
actual base-field transition: finite correspondences and numerical
idempotents descend, and the category over the algebraic closure is
classified by Weil germs. Its protorus is the identity component of
the fixed-field group, with quotient $D_E(\mu_\infty)$.
CC-AG07--08 give the exact finite-transition sequence and classify
rational object descent by Galois-invariant multiplicities divisible
by their explicit Schur indices; this is no longer an unspecified gate.
CC-AG09--10 give the exact rational base-field essential image and an
even-multiplicity obstruction over $\mathbb F_3$. CC-AG11 gives a
finite root test and the cyclic multiplicity monoid for integral-slope
germs; it does not assert the same cyclic formula for fractional slopes.
CC-AG12 removes the unbounded root-order search for each specified
rational object, including fractional slopes: rank bounds give a finite
polynomial and finite positive multiplicity equations, enumerating all
rational preimage isomorphism classes over a fixed finite field.
CC-AG13 computes the entire geometric multiplicity monoid over the
prime field in every slope sector: it is generated by the least
representative exponent divided by the geometric Schur index.
CC-AG14 refutes a general cyclic extension of that formula: an explicit
weight-one orbit over $\mathbb F_8$ has exactly the geometric and
specified finite-target multiplicities $\langle2,3\rangle$. The
corresponding geometrically simple abelian threefold has its square
and cube, but not itself, descend up to isogeny, by classical
Honda--Tate alone. CC-AG15 shows that the same representative exponent,
index and stabilized centre can give the different exact monoid
$3\mathbf N$ over $\mathbb F_{163^3}$; a local splitting obstruction
survives every root-of-unity twist. BR-15-PRIME-SUPPORT reduces generators
of the geometric multiplicity monoid to root-of-unity towers at the
finitely many primes dividing the least representative exponent. A
uniform bound on those tower heights is now proved by CC-AG16: an
explicit finite norm polynomial gives generators for the entire
positive geometric monoid, its exact gaps and conductor, over every
fixed finite source field. Its cutoff is independent of the target
multiplicity. The separate CC-AG12 procedure still enumerates every
preimage class of a specified target. CC-AG17 further proves
$H_a(\xi)=H_{\gcd(a,h_\xi)}(\xi)$ and an explicit label-orbit
bijection between exponents differing by $h_\xi$. Thus one finite
divisor table computes the entire geometric field-of-definition
profile; a specified finite-sum object has a computable minimal
divisor antichain describing all its source fields.
CC-XC01--04 now reduce membership of a specified smooth projective
variety to a determinant condition on one projected finite family
of cross-cycle maps. This condition forces semisimplicity and yields
fixed-prime Tate and numerical--homological comparison on the stated
products. CC-XC05--07 replace the positive witness by finite rational
cycle certificates with an intersection-trace determinant verifier.
An accepted certificate supplies the same rational splitting for
every prime different from the characteristic. An enumerable search
halts exactly on these positive instances. CC-XC08--10 prove that this
search terminates for all smooth Fermat and diagonal hypersurfaces of
degree prime to the characteristic, and for the stated sums of
squarefree binary forms with an optional pure-power term. Membership
also descends across finite fields and is preserved by the proved
trace and codimension-two blowup operations. CC-XC11--12 further
remove the rational-target choice: crystalline Frobenius forces every
finite Schur-index divisibility, and hard Lefschetz gives the required
real-place parity for the actual cohomology of any smooth projective
variety. Its Frobenius polynomials uniquely specify a rational
abelian-sector spectral target. CC-XC14 constructs its concrete untwisted presentation by a terminating
algorithm using only cycles on known abelian products. CC-XC15 also
computes a fixed rational duality and its inverse. Thus only one family
of forward rational cross-cycles varies in the remaining positive
search; its reverse maps are fixed adjoints, and the computable
intersection determinant detects an isomorphism at every prime.
For arbitrary $X$, this alone does not construct those cycles or
assert semisimplicity. CC-XC16--17 now produce the needed cycles when
twisted products of curves generate the algebraically trivial
$\mathrm{CH}_i$ for $i<\lfloor\dim X/2\rfloor$ over a universal
domain. Vial's self-dual splitting leaves a residual variety of
dimension at most one; the proved finite-field descent retains actual
Chow correspondences and their rational-equivalence witnesses. This
includes universal rational $\mathrm{CH}_0=\mathbf Q$ in dimensions
two and three, hence all smooth projective geometrically rationally
chain connected surfaces and threefolds. Testing Chow groups only over
$\overline{\mathbf F}_p$ does not verify the universal-domain hypothesis.
CC-XC18 gives a further geometric reduction for flat fibrations whose
geometric fibres have the stated small rational Chow groups: the motive
is the base summands plus one projector of dimension $\dim X-2n$.
For quadric fibrations the residual dimension is $\dim B$ for even
relative dimension, and $\dim B-1$ for odd relative dimension.
Thus every such fibration over a curve is certified. Over a surface,
odd relative dimension leaves exactly the base condition; even relative
dimension retains the base and one specified surface projector.
All geometric fibres, including the singular ones, are part of the
hypothesis; generic-fibre information alone has not been substituted.
CC-XC19 additionally records Vial's published rational--numerical
comparison on quadric fibrations over bases of dimension at most two.
The equivalence sandwich proves the standard conjecture $D$ and
cycle-class injectivity on every finite extension of each such total
space. This statement does not supply Tate surjectivity or the
comparison on all powers and products.
CC-XC20--21 also settle actual correspondence existence for geometrically
Tate Chow motives. Their finite-field forms are Artin--Tate, with the
Galois action retained. The published supersingular K3 theorem supplies
this input for all supersingular K3 surfaces, their Hilbert schemes and
the stated primitive stable-sheaf moduli. All finite products then have
rational--homological--numerical equality and Tate surjectivity.
This discharges an XC18 residual projector when it is actually identified
with a summand of these generators; a general such identification remains
unproved.
CC-XC22 extends these all-product conclusions to geometrically supersingular
abelian Chow motives and proves semisimple Frobenius; its fixed-target
algorithm uses only supersingular generators and needs no all-abelian input.
CC-XC23 applies the inspected October2026 lifting theorem to excellent
supersingular varieties of the four known symplectic deformation types,
with $p>\dim X$. The $K3^{[n]}$ case and the OG10/OG6 cases with Artin
invariant different from12/4 are geometrically Tate; the generalized
Kummer case with invariant different from3 is supersingular abelian.
Excellence includes unchanged Hodge numbers and an extended primitive
polarization. Missing lifting hypotheses and excluded invariants remain
unresolved by that theorem.
BR-15-QUADRIC-DISCRIMINANT makes the even-quadric condition concrete
when the fibration has a line-bundle quadratic presentation, positive
even relative dimension, odd characteristic, corank at most one and
smooth discriminant. Actual Chow membership, and separately homological
membership, is then equivalent to that of its base and its canonical
discriminant double cover. Finite orbit-category retractions prove both
directions. The $\mathbf P^2$ base with supersingular-K3 cover is
therefore fully settled by BR-15-QUADRIC-TATE, including all products.
CC-XC13 extends the rational spectral target to proper smooth schemes
using the cited crystalline comparison and equivariant pairing theorems.
The concrete cycle-verification algorithm retains its projective scope. Extending the classification
to arbitrary varieties still requires proving that this search always
finds such a certificate; no cycle-degree bound or negative-halting
test is established. The
quaternion example AQ-RATIONAL-FIBRE also shows why a rational neutral
fibre functor and arbitrary rational multiplicities cannot be silently
inserted. AQ-WEIL-NOTATION recovers the historical totally real and CM
closures on the earlier original sheets and proves the norm-condition
identification; that historical notation is no longer a gate.
The hypotheses of Milne (1994), \S2, have therefore been
discharged in a specified sector, not deleted globally.
\item[Cote 18, scans 11--13 and 59: period density.]
The overly strong formulation is refuted, and generation for one
elliptic curve has a precise known answer. Non-CM transcendence degree
four and the non-CM algebraicity question for the Hodge-splitting
scalar remain the specified transcendence gates. Group generation
does not imply density of a particular period locus; AQ-COCYCLE
gives the exact distinction. The separate four numbered comparison
problems on physical page 43/footer 42 now have scoped classical
answers KN-COMP01--04. BR-18-AD01 proves the separate adelic
genericity assertion on physical page 49/footer 48, retaining the
compact-open-image hypothesis for its Galois application.
Physical page 11/footer 10 explicitly defines $K=k(\alpha,\beta)$,
but the complete intervening-sheet reread has not established that
physical page 45/footer 44 uses the same field. That later symbol
assignment remains a source gate.
\item[Cote 134-8, footers 47--54: abstract scope and ordinary right adjoint.]
F-SM-29 now proves the underlying connective localization conclusion
of the finitely presented-projective-generator corollary over every
scalar ring, with the original marginal scalar-projectivity condition.
It checks all Smith recognition hypotheses and constructs an inverse
by Dold--Kan realization. There is no remaining coefficient-cardinality
or $\mathrm{FP}_\infty$ restriction on this conclusion.
The ordinary right adjoint is a separate stronger assertion:
F-SM-24 and F-SM-27 establish it in the stated pure-injective and
finite-generation ranges. F-SM39 now proves the literal unreplaced
counit for every category already of the form $\operatorname{Ch}_{\ge0}(\operatorname{Mod}\!-\!R)$
with the stipulated scalar adjoints, canonical unit normalization and
projective scalar $R(a,b)$, without finite-generation or scalar-dimension
restrictions. Coinduced objects are explicit outer disks. Outside these
classes its exact Tor test SM25.2 remains. F-SM40 also settles the raw
counit when the actual interval functor preserves coproducts,
cokernels and products: its evaluation coefficients are finitely
presented, yielding scalar-split interval inclusions and annihilation
of every product of arbitrary coinduced modules. This product
condition is not inferred from the handwritten bilinearity. F-SM42
further proves the raw counit whenever the actual acyclic interval
inclusion is split over the scalar ring at each ringoid object. No
preservation of products, coproducts, cokernels or exact sequences
by the interval functor is required. Arbitrary coinduced products
are relatively injective for these inclusions. The scalar splitting
condition itself remains additional. The right derived adjoint is already a quasi-inverse by
F-SM-29. F-SM-30 proves the Grothendieck-category case.
F-SM-32 now removes AB5, a generator and presentability: it suffices
that the canonical derived-coproduct comparison (SM32.1) hold, as it
does in every AB4 abelian category with enough projectives. The
coherent inverse and its actual heart augmentation are proved without
functorial projective resolutions. The remaining abstract gate is
whether the handwritten scalar/tensor/interval assumptions force
(SM32.1), or admit a full counterexample. The extra compatibility
condition is not inferred from bilinearity alone. F-SM-33 further replaces arbitrary coproduct compatibility by its
countable version and canonical counit normalization on all free scalar
modules. F-SM-35 proves AB4 from exact copowers when products exist;
this conditional convention is not read into the manuscript.
F-SM44 now proves that the canonical derived-coproduct map is
injective in every degree and is an isomorphism in degrees zero and
one, without using tensor or interval. Its remaining cokernel in
degree $n\ge2$ is dual, under an injective scalar cogenerator, to
higher extension classes that vanish on every summand. Unit
normalization also forces $\varphi\kappa U=U$ and
$L_1\varphi(\kappa U)=0$ for every scalar module, and the scalar
embedding preserves and reflects all degree-one extension classes.
Thus only the higher surjectivity and higher scalar-normalization
conditions remain at this compatibility gate. The original tensor
and interval have not yet been shown to kill those higher kernels.
F-SM-31 refutes automatic AB5 and presentability,
using a complete symmetric monoidal category with an exact tensor,
a projective generator and the full interval, but without AB5 or
local presentability. Its localization conclusion is nevertheless
true, so this is not a counterexample to the general theorem.
F-SM-34 supplies a fully proved additive counterexample with both
scalar adjoints, exact scalars, unit normalization and monic acyclic
effacement. It fails the free-scalar counit and localization conclusion,
but no original symmetric tensor is supplied.
F-SM36--37 rule out the specified componentwise scalar tensor and
the usual strong-monad tensor repair of that example. F-SM38 further
excludes every scalar-compatible tensor on this example that is exact
in the single countable free scalar factor, using its canonical scalar
two-extension. These are
failed-route obstructions, not a counterexample with all original
tensor hypotheses.
F-SM43 shows that even closedness and a flat interval with flat
quotient do not force AB4: its complete symmetric monoidal example
has the original adjunctions, canonical normalization and interval,
yet a countable coproduct of monomorphisms fails to be monic.
Its derived scalar functor still preserves coproducts and gives the
claimed localization, so the automatic derived-normalization question
remains distinct from the disproved automatic-AB4 route.
The intake's earlier three-sort construction has not been certified
with all original monoidal and adjunction hypotheses.
\item[Cote 136, physical 111--121: the integral non-strict questions.]
The original mathematical sheets and their coefficient conventions
are recovered. AQ-136-COMPARISON proves the weight-tail comparison;
AQ-136-DP-MODEL constructs an admissible theory over every ring.
AQ-136-STRICT refutes the strict equalizing-source assertion over
every nonzero ring, including $\mathbf Z$.
AQ-136-RATIONAL proves common sources and terminal ordinary
all-arrow localization over $\mathbf Q$-algebras.
AQ-136-PLOCAL and AQ-136-PLOCAL-LOCALIZATION extend both non-strict
assertions to every $\mathbf Z_{(p)}$-algebra, finite products of such
rings, and hence every positive-characteristic ring.
The coefficient-preserving simplicial cotensor path has a full proof;
it does not require a characteristic-zero cochain interval.
AQ-136-SEMILOCAL now proves terminality of the ordinary all-arrow
groupoid whenever all but finitely many rational primes are units in
the coefficient ring, by arithmetic path gluing and compatible finite
cones. AQ-136-COSPAN proves that every admissible cospan over every
coefficient ring has an admissible path pullback with strict projections.
AQ-136-CONNECTED-SOURCES therefore identifies the common-source relation
with connectedness, and AQ-136-SEMILOCAL-SOURCES supplies strict common
sources for every finite family over all these mixed semilocal rings.
Over $\mathbf Z$, common-source existence is now exactly connectedness;
connectedness and triviality of all loops remain open and distinct. The one-weight integral substitution fails by arbitrarily large
prime-weight classes. Admissibility does not force the weight-one
row to be acyclic, so the standard weight-one theory is not an
automatic source for every target. AQ-136-MULTIWEIGHT also proves
that independent coefficient disks in every large weight leave nonzero
Koszul $H_1$ in arbitrarily large weights, already rationally.
AQ-136-DERHAM-CELLS rules out repairing this by arbitrary coherent
higher simplicial zero-form cells of weights at least two followed
by the full relative special-PD de Rham complex: coefficient
boundaries force nonzero one-form homology at arbitrarily large
nonunit primes. AQ-136-FREE-ZERO-FORMS strengthens this obstruction
to arbitrary weight-one zero-form cells whenever the weight-one
coefficient summand has a simplicial retraction. Every map to the
actual admissible truncated standard target $C_N$, $N\geq2$, forces
that retraction. Thus the full relative free-zero-form de Rham route
cannot supply all integral common sources even with weight-one cells.
Nonfree positive-form relations and the proved coskeleton pullbacks
remain outside this obstruction.
AQ-136-TENSOR gives the actual derived-module calculation for
$B\otimes_S B$: its zero-form row has first homology $k/pk$
in every prime-power weight, despite simplexwise freeness and vertical
cohomology $S$. The tensor target therefore fails over every nonzero
non-$\mathbf Q$-algebra. These failed routes are recorded
separately from the actual existence questions.
AQ-136-TAME additionally excludes a prescribed fully vertically
acyclic odd-disk closure, while constructing an admissible target
with the relevant torsion class. It retains the need for compatible
whole-object transitions and uniform fixed-degree weight bounds. AQ-136-PATH-LIFT refutes the
general source-lifting property for the actual arithmetic path arrow,
already over a nonzero rational algebra and over mixed semilocal rings.
It does not refute either original non-strict assertion. AQ-136-RELATIVE-H1 and AQ-136-PARTIAL-SOURCE
now supply, over every nonzero integer-torsion-free coefficient ring,
an actual strict partial source with maps to any finite target family,
all horizontal tail conditions, rational vertical cohomology $S$,
and integral vertical $H^0=S,H^1=0$ uniformly in high weights.
The higher vertical cohomology is integer-torsion.
AQ-136-HIGHER-COMPLETION proves that these partial properties alone
do not guarantee an admissible relative completion: a specified
partial source with two maps to the same standard theory has an
explicit circle Bockstein obstruction. The next construction must
maintain an additional compatibility invariant or replace the partial
source/maps. The target pair already has a common admissible source,
so neither original non-strict assertion is refuted.
AQ-136-DERIVED-COMPATIBILITY identifies a necessary additional
condition: after sections and derived coefficient change by
$\mathbf Z/n$, all prescribed maps must induce the same canonical
comparison map, in a weight tail uniform in every modulus and every
finite simplicial set. Its explicit cone class recovers the earlier
completion obstruction without flatness or base-change assumptions.
Sufficiency of these compatibilities is not asserted.
AQ-136-MULTIPLICATIVE proves that the canonical ordinary cohomology
comparison always preserves the weighted cup product, without
flatness. AQ-136-PD-OPERATIONS makes each divided-power operation
well defined by its simplicial map of cycle spaces.
AQ-136-PD-RIGIDITY then identifies these operations between every
pair of theories, uniformly in the weight tail, on finite simplicial
sets over every ring and on all simplicial sets when the additive
coefficient group is finitely generated, including $\mathbf Z$.
Finitely many prime-local comparisons and a universal
Eilenberg--Mac Lane space remove the factorial-torsion ambiguity in
that scope. These operations therefore give no distinction between
integral components. The infinite-space coefficient-localization
counterexample explains why the universal-space argument is needed;
none of these operation results supplies a strict common source.
AQ-136-FIRST-SKELETON now gives actual admissible replacements of
any finite target family with one shared free integral first skeleton:
$F_0=S$, and $F_1$ has even zero-form PD disks in every weight above
one bound. The coskeleton pullback proof verifies both uniform tails.
Thus higher-simplex compatibility is the remaining comparison task
for these replacements; integral connectedness itself is not asserted.
\item[Cote 123: the original apposition questions are now closed.]
The dimension, rational-singularity, differential-rank and
subregular-uniqueness questions are answered in their original scope.
CC-APG01--05 now gives the complete apposition quadruple and
generator-enriched classifications in every simple complex adjoint
type, exact free orbit counts, and a finite root-data test for each
full automorphism group. An outer-invariant class exists exactly
outside type $A_{2m}$; existence does not select a unique canonical
class. CC-APG03 proves opposite choices exist for arbitrary pairs of
maximal tori. The all-choice answer is now complete in every classical simple
type: $A_{n-1}$ exactly for prime $n$, $B_r$ and $C_r$ exactly
for $r\leq2$, and no $D_r$ with $r\geq4$. CC-APG07--08
supply the explicit odd/even orthogonal counterexamples. The
BR-123-WEDGE, BR-123-OPPF4 and BR-123-OPPE8 now supply
explicit nonopposite choices in $F_4$ and $E_8$.  The formerly remaining
exceptional cases are also complete: BR-123-OPPG2 proves every
choice opposite in $G_2$, and CC-APG09--11 gives nonopposite
choices in $E_6,E_7$. There is no remaining type in the original
complex some/all-choice question.
\item[Weight--monodromy and universal-curve categories.]
The local monodromy operator and its invariant exterior powers must
be retained. BR-WM-02 supplies a finite gate and has passed a
recorded independent whole-proof reading. CC-ST01--04 give the geometric
stack decomposition with its explicit separating-divisor summands.
CC-RES01--02 prove the canonical exit/right-Kan statement and the nodal
lifting-category calculation. Identifying an unseen chosen archive
functor with that canonical functor, or deducing all categorical base
changes from geometric proper base change, requires another argument.
Purity of a proper pushforward
complex is not purity of every ordinary singular-fibre cohomology group.
\item[The analytic method as a transferable theorem.]
The lattice ideals, repeated theta-lift estimates, period nonvanishing,
and tower bounds must be exported with exact hypotheses from the
full argument. BR-DET-01/02 identify the numerical determinant gate.
The footnote on $K$-groups contains no degree-indexed theorem to apply.
F-SM-14--16 instead give exact comparisons and explicit families of
torsion consequences of separately cited integral arithmetic theorems.
\item[Masullo's operator realization.]
The fully normalized numerical shrinking-strip identity is proved in
the Dirichlet case in F-SM-13, including the finite Fourier expansion
of the tests and the archimedean parity term. F-SM-18 proves diagonal summation with a valid normalized indentation,
and F-SM-20 gives the actual compact-invariant compression. F-SM-19
refutes the printed indentation: it vanishes for $\mathbf Q$.
F-SM-21--22 rule out two specific naive replacements, by a nonzero
prime-trace discrepancy and by noncompactness. The remaining
global-to-$S$-local kernel assembly in \S12.3, p.~514, must supply
a valid intertwiner and the required prime corrections; the
diagonal domination step alone no longer describes the gate. All-number-field Hecke scope and
zero-multiplicity assertions require additional proofs.
\item[Profinite GT, anabelian reconstruction and geometric spectra.]
The full fixed Ihara embedding's surjectivity, arithmetic associator
data, and a spectrum constructed geometrically from primes are distinct
targets. The current rational GT theorem and the adopted Dirichlet
theorem do not identify those objects. Explicit counterexamples here
refute only the stated routes through abelian cohomology or bare spectra.
\end{description}

\section{Checks and current status}
Every completed new proof is associated with its actual reading events
in \nolinkurl{CHECK_EVENTS.tsv}. No unproved candidate in this part is
used as an input to a completed proof. The timed checkpoints and the
three-hour, then hourly movement reviews are recorded in
\nolinkurl{LEDGER.md}; the two-no-movement stopping rule remains in force.
\clearpage
\section{References for the root geometric and analytic sections}
References in the consequences and GT sections additionally give their
own exact source records. All archive locators refer to the numbered
scan, excluding the unnumbered cover, at the Montpellier catalogue:
\url{https://grothendieck.umontpellier.fr/archives-grothendieck/}.
\begin{enumerate}[leftmargin=0pt,labelsep=0pt,label={}]
\item Beilinson, A. A., Bernstein, J., \& Deligne, P. (1982).
\emph{Faisceaux pervers}. Ast\'erisque, 100. Operative locators:
Stabilit\'es 5.1.14, pp.~128--129; Corollary 5.4.4,
Theorem 5.4.5 and Corollary 5.4.6, p.~142.
\url{https://publications.ias.edu/sites/default/files/Faisceaux%20pervers.pdf}
\item Berger, L. (2004). \emph{An introduction to the theory of
$p$-adic representations} [Author preprint, arXiv:math/0210184].
Operative locators: \S II.3.1 and II.4.1--II.4.3.
\url{https://arxiv.org/abs/math/0210184}
\item Deligne, P. (1980). La conjecture de Weil II.
\emph{Publications Math\'ematiques de l'IH\'ES, 52}, 137--252.
Operative locators: Theorem 1.8.4; Problem 1.8.15,
pp.~178--179; 1.9.5, p.~181.
\url{https://publications.ias.edu/sites/default/files/Number40.pdf}
\item Deligne, P., \& Goncharov, A. B. (2005). Groupes fondamentaux
motiviques de Tate mixte. \emph{Annales Scientifiques de l'\'Ecole
Normale Sup\'erieure, 38}(1), 1--56.
Operative locators in supplied source: \S1.1, equations
(1.1.2), (1.1.4), Proposition 1.7 and its proof.
\url{https://arxiv.org/abs/math/0302267}
\item Gan, W. T., Gross, B. H., \& Prasad, D. (2009).
\emph{Symplectic local root numbers, central critical $L$-values,
and restriction problems} [Preprint, arXiv:0909.2999].
Operative locators: Theorems 14.1, 15.1, 16.1;
Conjectures 17.1, 17.3, 24.1. These labels refer to the supplied version.
\url{https://arxiv.org/abs/0909.2999}
\item Gan, W. T., Qiu, Y., \& Takeda, S. (2014).
\emph{The regularized Siegel--Weil formula (the second term identity)
and the Rallis inner product formula} [Preprint, arXiv:1207.4709].
Operative locator: Theorem 11.4.
\url{https://arxiv.org/abs/1207.4709}
\item Grothendieck, A. (1969). Standard conjectures on algebraic
cycles. In \emph{Algebraic geometry, Bombay 1968} (pp.~193--199).
Oxford University Press. Operative locators: \S2, p.~195;
\S3, pp.~195--197; \S4 and Remarks 2--4, pp.~197--198.
\url{https://download.uni-mainz.de/mathematik/Algebraische%20Geometrie/Lehre/WS23.Padische.1968.Grothendieck.StandardConjs.pdf}
\item Kashiwara, M. (1986). A study of variation of mixed Hodge
structure. \emph{Publications of the Research Institute for Mathematical
Sciences, 22}(5), 991--1024.
\url{https://doi.org/10.2977/PRIMS/1195177264}.
Operative locators: \S1.8--1.9; Theorems 4.4.1 and 4.5.2;
Propositions 4.5.1 and 5.2.1.
\item Katz, N. M., \& Messing, W. (1974). Some consequences of the
Riemann hypothesis for varieties over finite fields.
\emph{Inventiones Mathematicae, 23}, 73--77.
Operative locators: Theorem 1; Theorem 2(1), pp.~76--77.
\url{https://web.math.princeton.edu/~nmk/old/katzmessing.pdf}
\item Milne, J. S. (2006). \emph{Arithmetic duality theorems} (2nd ed.).
BookSurge. Operative locator: Chapter II, Theorem 3.1(b), p.~177;
Corollaries 3.2--3.3 specify the ordinary and compact-support variants.
\url{https://jmilne.org/math/Books/ADTnot.pdf}
\item Pearlstein, G. (2000). \emph{Degenerations of mixed Hodge structure}
(arXiv:math/0002030v3, 6 October). Operative locator: Theorem 5.1(i)--(iii),
p.~30 of this preprint. \url{https://arxiv.org/abs/math/0002030v3}
\item Scholze, P. (2012). \emph{Perfectoid spaces}
[Preprint, arXiv:1111.4914]. Operative locators in the inspected
preprint: Conjecture 9.3 and Theorem 9.4, p.~48;
Lemma 9.5 and Theorem 9.6, p.~49.
\url{https://arxiv.org/abs/1111.4914}
\item Titchmarsh, E. C. (1986). \emph{The theory of the Riemann
zeta-function} (2nd ed., revised by D. R. Heath-Brown).
Oxford University Press. Operative locators: \S2.12,
pp.~29--30; Theorems 9.3--9.4 and equation (9.4.3), p.~214.
\item Voevodsky, V. (2000). Triangulated categories of motives
over a field. In V. Voevodsky, A. Suslin, \& E. M. Friedlander,
\emph{Cycles, transfers, and motivic homology theories}
(Annals of Mathematics Studies, Vol.~143, pp.~188--238).
Princeton University Press. Operative locators in the author PDF:
\S2.1 after Proposition 2.1.3; Theorem 3.1.12; Corollary 3.4.3.
\url{https://www.math.ias.edu/vladimir/sites/math.ias.edu.vladimir/files/s5_0.pdf}
\item Deligne, P., \& Katz, N. (Eds.). (1973).
\emph{Groupes de monodromie en g\'eom\'etrie alg\'ebrique,
SGA 7 II} (Lecture Notes in Mathematics, Vol.~340). Springer.
Operative locator: Expos\'e XV, Th\'eor\`eme 3.4(ii)--(iii),
pp.~195--196 (specialization, vanishing cycles and Picard--Lefschetz).
\end{enumerate}

\section{Primary bibliography and source audit for the AQ
module}\label{primary-bibliography-and-source-audit-for-the-aq-module}

All mathematical attributions below distinguish the cited theorem from
the deductions proved in \nolinkurl{archive_questions.tex}. Page numbers
for Milne, Achter--Casalaina-Martin--Vial, KSV and Caradot--Lin refer to
the specified author/arXiv version, not a different published
pagination. The original archive quotations come from the supplied
transcriptions or the explicitly identified official scans recovered in
this run; the intake reports are secondary records of earlier work and
do not substitute for their absent proof files.

\subsection{References in APA style}\label{references-in-apa-style}

Achter, J. D., Casalaina-Martin, S., \& Vial, C. (2021). \emph{On the
image of the second ℓ-adic Bloch map} (arXiv:2007.07180v2).
\url{https://arxiv.org/pdf/2007.07180v2}

\begin{itemize}
\tightlist
\item
  Relevant primary statements: Appendix A, Propositions A.23--A.26
  (functoriality and specialization), A.28--A.29 (codimension one and
  zero-cycles/Roitman). The appendix explicitly treats arbitrary
  codimension where used. The entire paper is not claimed independently
  reproved.
\end{itemize}

Bloch, S. (1979). Torsion algebraic cycles and a theorem of Roitman.
\emph{Compositio Mathematica, 39}(1), 107--127.
\url{https://www.numdam.org/article/CM_1979__39_1_107_0.pdf}

\begin{itemize}
\tightlist
\item
  Used: Proposition 3.5, p. 115, correspondence compatibility;
  Proposition 3.8, p. 117, specialization; Theorem 4.2, p. 119,
  prime-to-characteristic Roitman. The modern Appendix A above supplies
  a convenient explicit formulation of the same operations.
\end{itemize}

Boavida de Brito, P., Horel, G., \& Robertson, M. (2019). Operads of
genus zero curves and the Grothendieck--Teichmüller group.
\emph{Geometry \& Topology, 23}, 299--346.
\url{https://doi.org/10.2140/gt.2019.23.299}

\begin{itemize}
\tightlist
\item
  Primary PDF: \url{https://msp.org/gt/2019/23-1/gt-v23-n1-p07-s.pdf}
\item
  Used: Definition 7.2(II), p. 324, the profinite
  Grothendieck--Teichmüller hexagon. This does not define the real
  hexagonal elliptic lattice of the archive.
\end{itemize}

Broe, M. (2026). \emph{The Tate conjecture for abelian fourfolds over
finite fields} (arXiv:2608.28651v1, August 19).
\url{https://arxiv.org/pdf/2608.28651v1}

\begin{itemize}
\tightlist
\item
  Theorem 1.1, p. 1, gives the codimension-two cycle-map surjection for
  every abelian fourfold over a finite field. Corollary 1.2, p. 2, and
  Corollary 8.2, p. 20, give standard conjecture D in dimension four
  over arbitrary fields. Corollary 8.1(i), p. 20, gives rational equals
  numerical equivalence with rational coefficients over a finite field.
  Exact statements and their scope were inspected; the complete preprint
  proof was not independently checked here. The later all-abelian result
  below is stronger for the present realization application.
\end{itemize}

Caradot, A., \& Lin, Z. (2026). \emph{Differential graded algebras with
divided powers and homotopy Lie algebras} (arXiv:2511.22614v4, June 20).
\url{https://arxiv.org/pdf/2511.22614v4}

\begin{itemize}
\tightlist
\item
  Checked directly: §2.1, pp. 5--6, the termwise-free convention;
  Proposition 3.10 and its proof, pp. 19--20; Theorem 3.15, pp. 23--24.
  AQ-SA08-LIFT refutes the degreewise-surjective lifting statement in
  this exact version. Its map is not a quasi-isomorphism, so it does not
  refute Theorem 3.15. No assertion is made about changes in a later
  journal version.
\end{itemize}

Chudnovsky, G. V. (1980). Algebraic independence of values of
exponential and elliptic functions. In \emph{Proceedings of the
International Congress of Mathematicians (Helsinki, 1978)} (Vol. 1, pp.
339--350). Academia Scientiarum Fennica.

\begin{itemize}
\tightlist
\item
  Official proceedings:
  \url{https://www.mathunion.org/fileadmin/ICM/Proceedings/ICM1978.1/ICM1978.1.ocr.pdf}
\item
  Exact primary locator: §2A, Theorem 1 and Corollary 2, printed p. 343.
  The algebraic-invariants setting is fixed on p. 339. Theorem 1 gives a
  transcendence-degree lower bound for periods and quasi-periods;
  Corollary 2 gives algebraic independence of a nonzero CM elliptic
  period and π. The meeting year is 1978; publication is 1980. The
  relevant twelve-page extract was read locally.
\end{itemize}

Duren, P. L. (1983). \emph{Univalent functions} (Grundlehren der
mathematischen Wissenschaften, Vol. 259). Springer.
\url{https://link.springer.com/book/9780387907956}

\begin{itemize}
\tightlist
\item
  Standard locator: Theorem 2.3, p. 31, Koebe quarter theorem. The book
  locator was corroborated indirectly; the AQ module supplies the entire
  needed area-theorem/odd-function/Möbius proof, so the bound does not
  depend on an unexamined book proof.
\end{itemize}

Grothendieck, A. (n.d.). \emph{Théorie des cycles algébriques:
Conjectures de Hodge, Tate, Lefschetz, Weil} {[}Unpublished manuscripts,
cote 14{]}. Archives Grothendieck, Université de Montpellier.
\url{https://grothendieck.umontpellier.fr/14.pdf}

\begin{itemize}
\tightlist
\item
  Supplied source: scan 38, manuscript p. 37. It states the rational
  ℓ-adic conclusion, labels Théorème 3.1 ``dém. pas faite,'' and asks
  for a converse. The original quotient is now identified:
  Definition1.1/Proposition1.2(i), physical PDF page24/footer23;
  Proposition1.2(v), physical25/footer24; Theorem2.1,
  physical31/footer30; Corollary2.3, physical34/footer33. The word
  projectif is bracketed in Definition1.1; the proofs cover its
  projective-parameter reading. Physical page38 corresponds to
  footer/manuscript37, matching the supplied scan038 convention.
\item
  Definition 1.1 says only a field k. No explicit algebraically closed
  convention was found in the preceding relevant handwritten pages
  inspected. Our theorem retains algebraic closure as an explicit
  hypothesis; the finite-field rational-point counterexample shows why
  it matters.
\end{itemize}

Grothendieck, A. (n.d.). \emph{Théorie arithmétique. Théorie de Galois
des motifs} {[}Unpublished manuscripts, cote 15{]}. Archives
Grothendieck, Université de Montpellier.
\url{https://grothendieck.umontpellier.fr/15.pdf}

\begin{itemize}
\tightlist
\item
  Supplied sources: scans 100, 102, 104 and 106, manuscript pp. 99, 101,
  103 and 105. Official physical pages 90--106 were additionally
  retrieved where relevant, and the defining sheets visually inspected:
  physical94/footer93 identifies \texttt{Q′} as the totally real
  algebraic closure of Q; physical98/footer97 writes the norm and
  norm-one normalization; physical100/footer99 gives the boxed
  all-conjugate purity correction. Physical102/footer101,104/footer103
  and106/footer105 were also visually compared with the supplied
  transcriptions. The exhaustion question on physical104 remains in a
  deleted passage. AQ-WEIL-NOTATION proves the exact field/group
  identification with those norm/purity conditions retained. The
  counterexample to a valuation-only paraphrase does not attribute that
  reading to the manuscript term \emph{unité absolue}. Uncertain
  superscripts and crossed revisions remain source qualifications.
\end{itemize}

Grothendieck, A. (n.d.). \emph{Motivic, period and realization notes}
{[}Unpublished manuscripts, cote 18{]}. Archives Grothendieck,
Université de Montpellier.
\url{https://grothendieck.umontpellier.fr/18.pdf}

\begin{itemize}
\tightlist
\item
  Supplied sources: scans 9, 11, 13, 38 and 59; manuscript pp. 8, 10,
  12, 37 and 58. These contain respectively the Tate-filtration
  conventions, strong period triangle, three elliptic questions,
  real-involution claim and cocycle-image genericity reformulation. The
  title here is descriptive, not asserted to be a manuscript title.
\end{itemize}

Grothendieck, A. (1983). \emph{Courbes elliptiques réelles; capping and
nodal-propagation notes} {[}Unpublished manuscripts, cote 67{]}.
Archives Grothendieck, Université de Montpellier.
\url{https://grothendieck.umontpellier.fr/67.pdf}

\begin{itemize}
\tightlist
\item
  Supplied sources: scans 74--82, archive foliation 76--84, and scan 91,
  archive foliation 93. Scan 91 is explicitly dated ``Déc 83.'' This
  date is not assigned to every sheet of the cote. Scan 75 specifies a
  nonzero tangent tensor; scan 78 uses its inverse dual pairing and
  mutually polar tangent disks; scans 76--77 impose no norm restriction
  on that tensor. These are the precise source conventions used in
  AQ-67-RANGE.
\end{itemize}

Grothendieck, A. (n.d.). \emph{Pseudoline-relative-position surface
notes} {[}Unpublished manuscripts, cote 154{]}. Archives Grothendieck,
Université de Montpellier.
\url{https://grothendieck.umontpellier.fr/154.pdf}

\begin{itemize}
\tightlist
\item
  Supplied source: scan 120. Its Euler characteristic is
  cellular/topological. No finite-field scheme or Frobenius structure on
  that surface is specified. The title is descriptive.
\end{itemize}

Habegger, P. (2014). \emph{Singular moduli that are algebraic units}
(arXiv:1402.1632v1). \url{https://arxiv.org/pdf/1402.1632v1}

\begin{itemize}
\tightlist
\item
  Used: §1, p. 1, the classical integrality of CM j-invariants, to
  certify that j = −6912/23 is not CM. The AQ proof does not use the
  paper's finiteness theorem for singular units.
\end{itemize}

Honda, T. (1968). Isogeny classes of abelian varieties over finite
fields. \emph{Journal of the Mathematical Society of Japan, 20}(1--2),
83--95. \url{https://doi.org/10.2969/jmsj/02010083}

\begin{itemize}
\tightlist
\item
  Main Theorem, p. 83: existence for weight-one Weil integers. Primary
  scan:
  \url{https://people.math.rochester.edu/faculty/doug/otherpapers/Honda-Ab.pdf}
\item
  AQ-WEIL-EXIST derives the specified weight-i cohomological occurrence
  using an i-th root and Künneth. This is an unconditional known
  consequence, not a new RH application.
\end{itemize}

Kreutz, T., Shen, M., \& Vial, C. (2026). \emph{De Rham--Betti classes
with coefficients} (arXiv:2206.08618v3).
\url{https://arxiv.org/pdf/2206.08618v3}

\begin{itemize}
\tightlist
\item
  Also: \url{https://arxiv.org/html/2206.08618v3}
\item
  Read: Proposition 3.1, Theorem 5.13(ii), Proposition 5.14. The
  coefficient/CM-field restriction in Theorem 5.13(ii) is satisfied for
  one elliptic curve. The result identifies the generated period torsor,
  answering 18/013(b); it does not settle maximal transcendence degree
  for a non-CM period point.
\end{itemize}

Lagarias, J. C., \& Odlyzko, A. M. (1977). Effective versions of the
Chebotarev density theorem. In A. Fröhlich (Ed.), \emph{Algebraic number
fields: L-functions and Galois properties} (pp. 409--464). Academic
Press.

\begin{itemize}
\tightlist
\item
  Full retype of the original article:
  \url{https://aareyanmanzoor.github.io/assets/articles/lagarias-odlyzko.pdf}
\item
  Used: Theorem 1.1, equation (1.3), retype p. 4. For the quadratic
  extension Q(√−3)/Q, degree 2 and discriminant 3, its hypothesis
  follows from ζ\_K(s)=ζ(s)L(s,χ\_−3) and the adopted Dirichlet GRH. Its
  effective absolute constant is left as c1, not invented numerically.
\end{itemize}

Lamzouri, Y., Li, X., \& Soundararajan, K. (2015). Conditional bounds
for the least quadratic non-residue and related problems.
\emph{Mathematics of Computation, 84}(295), 2391--2412.
\url{https://arxiv.org/pdf/1309.3595v3}

\begin{itemize}
\tightlist
\item
  Operative source: corrected arXiv v3, dated December 8, 2016;
  Corollary 1.2, p. 3. Under Dirichlet GRH, q\textgreater3 and (a,q)=1,
  the least prime in the class a mod q is at most (φ(q) log q)².
  AQ-FINITE uses Q=lcm(N,4), so the q\textgreater3 condition is
  automatic.
\end{itemize}

Lamzouri, Y., Li, X., \& Soundararajan, K. (2017). Corrigendum to
``Conditional bounds for the least quadratic non-residue and related
problems.'' \emph{Mathematics of Computation, 86}(307), 2551--2554.
\url{https://doi.org/10.1090/mcom/3261}

\begin{itemize}
\tightlist
\item
  Bibliographic/version control: the operative theorem was checked in
  corrected v3. The standalone corrigendum's complete text was not
  independently opened in this AQ run; no unverified early-version
  constants are used.
\end{itemize}

Milne, J. S. (1994). Motives over finite fields. In U. Jannsen, S. L.
Kleiman, \& J.-P. Serre (Eds.), \emph{Motives} (Proceedings of Symposia
in Pure Mathematics, Vol. 55, Part 1, pp. 401--459). American
Mathematical Society. \url{https://jmilne.org/math/articles/1994aP.pdf}

\begin{itemize}
\tightlist
\item
  Exact author-PDF locators: Conjecture 1.14 and Proposition 1.15, p.
  11; Corollary 1.16, pp. 11--12; start of §2, p. 12; Proposition 2.6
  and Remark 2.7, p. 15.
\item
  The injectivity/classification and abelian/Artin generation claims use
  Tate. The eigenvalue-existence argument is unconditional via Honda.
  These two uses are separated explicitly in AQ-WEIL-EXIST.
\item
  Author's corrections page:
  \url{https://jmilne.org/math/xnotes/Seattle1.html}
\end{itemize}

Milne, J. S. (1999). Lefschetz motives and the Tate conjecture.
\emph{Compositio Mathematica, 117}(1), 45--76.
\url{https://doi.org/10.1023/A:1000776613765}

\begin{itemize}
\tightlist
\item
  Author manuscript:
  \url{https://www.jmilne.org/math/articles/1999bP.pdf}
\item
  Theorem 7.1, journal p. 72, proves CM Hodge implies the Tate statement
  (0.1) for all abelian varieties over an algebraic closure of a finite
  field. Introduction (0.1), journal p. 46, includes numerical equals
  homological and algebraicity of all Tate classes; the paragraph
  immediately following explains finite-model descent and the full
  pole-order statement. Remark 7.6(b), journal p. 76, extends the
  consequence to homological motives of abelian type. The CM-Hodge
  premise is supplied by the 2026 result below.
\end{itemize}

Milne, J. S. (2002). Polarizations and Grothendieck's standard
conjectures. \emph{Annals of Mathematics, 155}(2), 599--610.
\url{https://annals.math.princeton.edu/2002/155-2/p07}

\begin{itemize}
\tightlist
\item
  Operative final author version, September 19, 2001:
  \url{https://www.jmilne.org/math/articles/2002aP.pdf}
\item
  Theorem 3.3(a),(b), author pp. 8--9, gives numerical versus ℓ-adic
  homological equivalence and Hodge standard for all abelian varieties
  over algebraically closed fields, conditional on CM Hodge. Remark 3.5,
  author p. 9, treats good Weil cohomology theories. Earlier
  long-version Theorem 4.6 is not the numbering of this final author
  PDF. The primitive Hodge pairing has the sign (−1)\^{}r; the
  abbreviated 2026 report\textquotesingle s displayed pairing omits that
  sign.
\end{itemize}

Milne, J. S. (2026). \emph{On the Hodge, Tate, and standard conjectures
for abelian varieties} (Hexagon:2611.00244v1, November 9).
\url{https://www.jmilne.org/math/xnotes/HTS.pdf}

\begin{itemize}
\tightlist
\item
  Theorem 2.1 states CM Hodge; Theorem 3.1, p. 6, states the full Tate
  conjecture for all abelian varieties over finite fields; Theorem 3.5,
  p. 8, states Hodge standard for all abelian varieties. Section 3.7
  explains reduction to finite fields. This author report confirms the
  current status and cites the underlying proofs; it is not substituted
  for a new full proof check of those papers.
\end{itemize}

Mohajer, M. (2026). \emph{Rational elliptic period relations across
realizations: Complex lattices, local p-adic periods, and Drinfeld
modules} (arXiv:2611.25399v1, November 19).
\url{https://arxiv.org/pdf/2611.25399v1}

\begin{itemize}
\tightlist
\item
  Section 1.1, equation (1.3), pp. 2--3, retains the classical full
  complex-period conjecture; Proposition 2.1 concerns holomorphic E2.
  Theorems 4.2 and 4.4 concern local B\_dR and normalized C\_p periods.
  These exact scopes were inspected; the complete 39-page preprint proof
  is not certified here. Its local statements do not settle the complex
  E2* splitting-scalar problem.
\end{itemize}

Sebbar, A., \& Besrour, K. (2026). \emph{Zeros of Katz's p-adic
Eisenstein series E2: Hypergeometric congruences and isogenies}
(arXiv:2611.14502v1, November 10).
\url{https://arxiv.org/pdf/2611.14502v1}

\begin{itemize}
\tightlist
\item
  Section 8.4, Remark 8.11, Proposition 8.12, and Question 8.13, pp.
  32--33, leave algebraic non-CM zeros of the p-adic Katz form as a
  question and obtain nonvanishing under the crystalline period
  conjecture. Algebraicity over Q\_p is explicitly distinguished from
  algebraicity over Q. The abstract and these exact statements were
  read; the whole preprint is not claimed checked.
\end{itemize}

Tate, J. (1966). Endomorphisms of abelian varieties over finite fields.
\emph{Inventiones Mathematicae, 2}, 134--144.
\url{https://doi.org/10.1007/BF01404549}

\begin{itemize}
\tightlist
\item
  Primary scan:
  \url{https://pazuki.perso.math.cnrs.fr/index_fichiers/Tate66.pdf}
\item
  Main Theorem, p. 134, is the Tate-module endomorphism theorem. Theorem
  2(a), p. 140, identifies Q{[}π{]} with the centre of the semisimple
  endomorphism algebra, giving the semisimple Frobenius input. Theorem
  2(d), p. 140, identifies the quaternion algebra in the
  scalar-Frobenius supersingular case. AQ-RATIONAL-FIBRE verifies the
  explicit model over F\_9 before invoking this theorem.
\end{itemize}

van Hoften, P., \& Sempliner, J. (2026). \emph{The Hodge conjecture for
CM abelian varieties} (Hexagon:2610.00259v2).
\url{https://hexagonmath.org/2610.00259v2}

\begin{itemize}
\tightlist
\item
  Author-hosted PDF:
  \url{https://polvanhoften2.github.io/Pol\%20van\%20Hoften\%20and\%20Jack\%20Sempliner\%20-\%20The\%20Hodge\%20conjecture\%20for\%20CM\%20abelian\%20varieties.pdf}
\item
  Theorem I, p. 1, proves rational Hodge for every complex CM abelian
  variety. Sections 1.1.2--1.1.3, pp. 1--2, state the consequences for
  standard conjectures and finite-field Tate using Milne. The exact
  theorem and downstream hypotheses were inspected. The complete proof
  is a cited research input, not independently certified in the AQ
  module. Version 2 is retained explicitly to avoid relying on an
  unspecified first posting.
\end{itemize}

Waldschmidt, M. (2004). Open Diophantine problems. \emph{Moscow
Mathematical Journal, 4}(1), 245--305.
\url{https://webusers.imj-prg.fr/~michel.waldschmidt/articles/pdf/ODP-MoscowMathJ.pdf}

\begin{itemize}
\tightlist
\item
  Exact locator: p. 277, question immediately after Conjecture 3.22
  (Katz), about the nonholomorphic weight-two Eisenstein invariant for a
  non-CM lattice with algebraic invariants. The AQ module writes its
  exact normalized Hodge-splitting scalar.
\end{itemize}

Walker, M. E. (2007). The morphic Abel--Jacobi map. \emph{Compositio
Mathematica, 143}, 909--944.
\url{https://doi.org/10.1112/S0010437X07002278}

\begin{itemize}
\tightlist
\item
  Primary PDF:
  \url{https://www.cambridge.org/core/services/aop-cambridge-core/content/view/4A44146CE680F4F3C857048FDEDD9F48/S0010437X07002278a.pdf/the-morphic-abel-jacobi-map.pdf}
\item
  Read: Example 5.5, p. 928 (the zero-cycle morphic Jacobian/Albanese);
  Theorem 5.7(2),(4), p. 930 (cycle-map and every-Chow-correspondence
  compatibility); Theorem 6.1, p. 933, and the complete proof of Theorem
  6.2, pp. 933--934 (homologically trivial cycles with nonzero Lawson
  cylinder after an elliptic product); §7 (the separate ordinary
  regular-universality question).
\item
  No unchecked effective hypersurface degree is used. ``Sufficiently
  general of sufficiently large degree'' is the precise retained theorem
  hypothesis.
\end{itemize}

\subsection{Additional primary sources for the exact capping and
compactification
results}\label{additional-primary-sources-for-the-exact-capping-and-compactification-results}

Brock, J. F., \& Margalit, D. (2007). Weil--Petersson isometries via the
pants complex. \emph{Proceedings of the American Mathematical Society,
135}(3), 795--803. \url{https://arxiv.org/abs/math/0412499}

\begin{itemize}
\tightlist
\item
  Operative author version: December 21, 2004, arXiv:math/0412499v1;
  \url{https://jeffbrock.net/home/text/papers/dan/www/wp.pdf} . Main
  Theorem, p. 2, explicitly includes the four-punctured sphere and
  states that every Weil--Petersson isometry is induced by the extended
  mapping class group. The finite kernel in that case does not affect
  the countability argument.
\end{itemize}

Fırat, A. H. (2021). Hyperbolic three-string vertex. \emph{Journal of
High Energy Physics, 2021}(8), Article 35.
\url{https://doi.org/10.1007/JHEP08(2021)035}

\begin{itemize}
\tightlist
\item
  Operative primary version: arXiv:2102.03936v2, August 18, 2021,
  \url{https://arxiv.org/pdf/2102.03936} . Equations (3.11), (3.13),
  (3.14a), (3.23), and (3.30)--(3.34), pp. 14--17 and 24--25, provide
  the hypergeometric projective connection and the arclength
  flat-cylinder cap coordinate. AQ-67-EXACT supplies its own cusp-limit
  and continuous phase proof; no unproved numerical integer-branch
  prescription is used.
\end{itemize}

Kahn, B., \& Sujatha, R. (2008). \emph{Birational geometry and
localisation of categories} (author manuscript, May 22).
\url{https://webusers.imj-prg.fr/~bruno.kahn/preprints/Birat.pdf}

\begin{itemize}
\tightlist
\item
  Theorem 1.6.3(a), p. 11, states the characteristic-zero resolution and
  smooth compactification consequence in the exact form used: a smooth
  quasiprojective variety admits a smooth projective compactification.
  AQ-14-QUASIPROJECTIVE proves the correspondence/Albanese comparison
  from that cited input. This is an inspected primary theorem statement;
  the Hironaka proof is not independently reproduced.
\end{itemize}

Mondello, G. (2011). Riemann surfaces with boundary and natural
triangulations of the Teichmüller space. \emph{Journal of the European
Mathematical Society, 13}(3), 635--684.
\url{https://doi.org/10.4171/JEMS/263}

\begin{itemize}
\tightlist
\item
  Primary PDF:
  \url{https://ems.press/content/serial-article-files/31778} . Section
  2.2, p. 640, recalls Fenchel--Nielsen coordinates with zero boundary
  lengths interpreted as cusps. Section 5.4, Theorem 5.4, p. 660, treats
  the real-analytic infinite-grafting parametrization; the canonical cap
  is its punctured flat cylinder in a disk coordinate.
\end{itemize}

National Institute of Standards and Technology. (n.d.). \emph{Digital
Library of Mathematical Functions}. \url{https://dlmf.nist.gov/}

\begin{itemize}
\tightlist
\item
  Exact inspected identities: (15.4.21),
  \url{https://dlmf.nist.gov/15.4.E21} (zero-balanced hypergeometric
  logarithmic asymptotic); (5.5.5), \url{https://dlmf.nist.gov/5.5.E5}
  (gamma duplication); (5.9.16), \url{https://dlmf.nist.gov/5.9.E16}
  (digamma integral); (5.11.1), \url{https://dlmf.nist.gov/5.11.E1}
  (analytic-log Stirling expansion), and (5.11.13),
  \url{https://dlmf.nist.gov/5.11.E13} (gamma-ratio asymptotic).
  Accessed November 21, 2026. The positivity and range arguments are
  proved in full in AQ-67-EXACT.
\end{itemize}

Wolpert, S. A. (1983). On the symplectic geometry of deformations of a
hyperbolic surface. \emph{Annals of Mathematics, 117}(2), 207--234.
\url{https://doi.org/10.2307/2007075}

\begin{itemize}
\tightlist
\item
  Original scan:
  \url{https://math.umd.edu/~swolpert/preprints/On_the_Symplectic.pdf} .
  Lemma 4.1, printed p. 229, gives the exact twist Hamiltonian identity
  Ω(t\_α)=−dℓ\_α. The finite-area surface scope is stated in §1, pp.
  208--209. AQ-67-NONHOLO uses Brock--Margalit for the exceptional
  four-punctured-sphere isometry assertion, rather than an unqualified
  invocation of Royden.
\end{itemize}

\subsection{Additional primary sources for the all-characteristic
nonproper
implication}\label{additional-primary-sources-for-the-all-characteristic-nonproper-implication}

Deligne, P. (1974). La conjecture de Weil: I. \emph{Publications
Mathématiques de l'IHÉS, 43}, 273--307.
\url{https://doi.org/10.1007/BF02684373}

\begin{itemize}
\tightlist
\item
  Original primary scan:
  \url{https://www.numdam.org/article/PMIHES_1974__43__273_0.pdf} .
  Theorem (1.6), printed p. 276, is the smooth-projective Frobenius
  purity theorem. Twisting H\^{}\{2i−1\} by i gives geometric weight −1;
  a torus Tate module has geometric weight −2. The candidate proves the
  resulting vanishing without requiring semisimplicity.
\end{itemize}

Geisser, T. (2015). Rojtman's theorem for normal schemes.
\emph{Mathematical Research Letters, 22}(4), 1129--1144.
\url{https://arxiv.org/abs/1402.1831}

\begin{itemize}
\tightlist
\item
  Operative directly inspected primary version: arXiv:1402.1831v2,
  February 25, 2015, \url{https://arxiv.org/pdf/1402.1831} . Theorem
  1.1, preprint p. 2, gives prime-to-characteristic torsion
  Suslin--Albanese isomorphism and H₁\^{}S(T,Z)⊗Q\_l/Z\_l=0 for normal
  separated finite-type schemes, without assuming smooth
  compactification. Theorem 1.3, preprint p. 3, supplies the
  smooth-curve special case H₀\^{}S(C,Z)\^{}0≅Alb\_sa(C)(k). Lemma 2.2,
  p. 5, records the Albanese map on Suslin H₀. The exact theorem
  statements were read; the entire Geisser proof is a cited input, not a
  new whole-proof certification.
\end{itemize}

Geisser, T., \& Schmidt, A. (2016). Tame class field theory for singular
varieties over algebraically closed fields. \emph{Documenta Mathematica,
21}, 91--123. \url{https://doi.org/10.4171/DM/528}

\begin{itemize}
\tightlist
\item
  Original journal PDF:
  \url{https://ems.press/content/serial-article-files/26306?nt=1} .
  Theorem 1.1, printed p. 92, gives the finite perfect
  pairing/reciprocity isomorphism for coefficients prime to the
  characteristic. Theorem 7.1, printed p. 118, identifies the pairing
  with the Suslin--Voevodsky comparison. The candidate passes through
  direct and inverse coefficient limits explicitly; it does not claim
  that every finite-coefficient H¹ of a variety equals that of its
  Albanese.
\item
  The published PDF lists authors Thomas Geisser and Alexander Schmidt
  in that order; the publisher metadata happens to reverse them. The
  citation follows the paper itself.
\end{itemize}

Kahn, B. (2014). Algebraic tori as Nisnevich sheaves with transfers.
\emph{Annales de la Faculté des Sciences de Toulouse: Mathématiques,
23}(3), 699--715. \url{https://doi.org/10.5802/afst.1421}

\begin{itemize}
\tightlist
\item
  Directly inspected final journal PDF:
  \url{https://webusers.imj-prg.fr/~bruno.kahn/preprints/tores-nisnevich.pdf}
  . Lemma 4.1(a),(b), printed p. 707, proves that a pointed morphism of
  semiabelian varieties is a homomorphism and that a semiabelian variety
  is its own Serre Albanese. Equivalent earlier locator: July 4, 2013
  author version,
  \url{https://webusers.imj-prg.fr/~bruno.kahn/preprints/tores-nisnevich8.pdf}
  , Lemma 3, p. 8. Both exact statements and the short proof were read.
\end{itemize}

Sun, S. (2017). Generic base change, Artin's comparison theorem, and the
decomposition theorem for complex Artin stacks. \emph{Journal of
Algebraic Geometry, 26}(3), 513--555.
\url{https://doi.org/10.1090/JAG/683}

\begin{itemize}
\tightlist
\item
  Operative inspected primary version: arXiv:1511.01393v1, November 4,
  2015, \url{https://arxiv.org/pdf/1511.01393} . Section 1.3.1, p. 2,
  permits Λ=Z\_l. Section 2.1, p. 3, gives condition (LO) and the
  integral stratifiable adic category. Definition 2.2, p. 3, specifies
  one dense open valid for all subsequent base changes. Lemma 2.3(iii),
  pp. 4--5, and its proof explicitly obtain one dense open for all
  finite coefficient levels by finite Jordan--Hölder dévissage. Theorem
  2.5(i), p. 7, states generic base change; Remark 2.5.1(ii), p. 8,
  recalls preservation of stratifiability. The candidate uses only
  finite-type F\_p bases with l≠p, so (LO) applies. A family of separate
  opens, one for each l\^{}n, would not suffice and is not the argument
  used.
\end{itemize}

Bloch correspondence and specialization inputs are those already
precisely cited in AQ-14-FORWARD: Bloch (1979), Propositions 3.5 and
3.8, and Achter--Casalaina-Martin--Vial, arXiv:2007.07180v2,
Propositions A.23, A.24 and A.26. The latter paper\textquotesingle s
(0.2)--(0.3) and §A.3 also identify the rational Tate module of the
torsion-coefficient cohomology with ordinary Q\_l-cohomology.

Kahn, B., \& Otsubo, N. (2026). \emph{Periods and CM motives of abelian
type} (arXiv:2611.13723v1; manuscript dated November 11).
\url{https://arxiv.org/pdf/2611.13723}

\begin{itemize}
\tightlist
\item
  Exact targeted locators: §2.4, pp. 6--7; Definition 2.5, p. 8;
  Corollaries 3.21, 3.25 and 3.35, pp. 24, 26 and 30--31; Remark 3.22,
  p. 25; Theorems 4.36, 4.39 and 4.44, pp. 46--47 and 49. The new
  comparison statements concern rank-one coefficient motives and period
  classes. AQ-KO-RANK corrects the ancillary rank sentence in Lemma 4.6,
  p. 34, by a complete elementary proof. This targeted reading is not
  certification of the full 64-page preprint.
\end{itemize}

Illusie, L., \& Temkin, M. (2014). Exposé X: Gabber's modification
theorem (log smooth case). In L. Illusie, Y. Laszlo, \& F. Orgogozo
(Eds.), \emph{Travaux de Gabber sur l'uniformisation locale et la
cohomologie étale des schémas quasi-excellents} (\emph{Astérisque,
363--364}, pp. 167--212).
\url{https://www.numdam.org/item/AST_2014__363-364__167_0.pdf}

\begin{itemize}
\tightlist
\item
  Theorem 3.5(ii)--(iii), printed p. 202, gives the separable projective
  prime-to-ℓ alteration over a perfect field used in AQ-14-ALTERATION.
  Definition 3.3.1, p. 197, specifies prime-to-ℓ degree on every generic
  component. The introduction, p. 167, and §3.5.4, p. 207, distinguish
  this separable strengthening from Theorem 2.1. Regular finite-type
  sources over the perfect base are smooth, and projectivity over a
  projective compactification supplies projective parameters. The exact
  theorem was directly inspected; its external proof remains a cited
  input.
\end{itemize}

Milne, J. S. (2008, March 16). \emph{Abelian varieties} (Version 2.0).
\url{https://www.jmilne.org/math/CourseNotes/AV.pdf}

\begin{itemize}
\tightlist
\item
  Chapter I, Theorem 3.2, printed p. 17, extends a rational map from a
  nonsingular variety to an abelian variety over the whole source. This
  exact numbering belongs to the specified course-note version. It is
  the extension input for the fibre-sum trace in AQ-14-ALTERATION.
\end{itemize}

\subsection{Check and scope record}\label{check-and-scope-record}

\nolinkurl{archive_check_events.tsv} records the AQ author's complete
readings of the written proofs and the separate symbolic/finite checks.
Independent readings by the main agent and \nolinkurl{source_and_masullo}
are recorded by those reviewers, not manufactured here. The full prior
SA06/07/08/11/22/23 proof returns are absent from the supplied packet.
Their available intake rows are preserved verbatim in
\nolinkurl{archive_intake_inventory.tsv}; indexing does not certify a
proof.

\subsection{Cote 66 source recovery and affine
comparison}\label{cote-66-source-recovery-and-affine-comparison}

Grothendieck, A. (n.d.). \emph{List of subjects} {[}Unpublished
manuscript, cote 66, physical PDF p. 4, archive footer 3, subject 21{]}.
Archives Grothendieck, Université de Montpellier.
\url{https://grothendieck.umontpellier.fr/66.pdf}

The Stacks Project Authors. (2026). \emph{The Stacks Project}.
\url{https://stacks.math.columbia.edu}

\begin{itemize}
\tightlist
\item
  The parent module \nolinkurl{root_henselian.tex}, BR-66-01--03, uses
  Theorem59.82.7/09ZI (Gabber comparison), Lemma15.11.6/09XI (idempotent
  criterion), Theorem59.91.11/095T (proper base change),
  Lemma15.11.4/0ALJ (complete pairs), Lemma15.12.1/0A02 (henselization),
  and Lemmas21.4.2--3/03AI--03AJ (torsor sections and H1). The complete
  deductions received an independent whole-proof reading by this
  module\textquotesingle s author; the exact Stacks statements were
  directly inspected on21November2026. The parent alone recovered and
  visually checked the original66 sheet; this cross-review does not
  relabel that as a separate source inspection here.
\end{itemize}

\subsection{Cote 136 recovered source and exact simplicial
inputs}\label{cote-136-recovered-source-and-exact-simplicial-inputs}

Grothendieck, A. (n.d.). \emph{Notes on de Rham functors and admissible
theories} {[}Unpublished manuscript, cote 136; descriptive title{]}.
Archives Grothendieck, Université de Montpellier.
\url{https://grothendieck.umontpellier.fr/136.pdf}

\begin{itemize}
\tightlist
\item
  Directly visually inspected: physical111/footer110 (functors and
  simplex representation),113/footer112 and115/footer114 (cohomological
  comparison),117/footer116 (vertical conditions3a/3b and
  simplex-uniformity question),119/footer118 (horizontal flasqueness
  condition4),121/footer120 (common-source a, strict equalizing-source
  b, and generated groupoid). The intervening reverses are unrelated
  printed material. The exact coefficient, grading, strict PD and
  uniform-cutoff conventions used in the proofs are written in AQ-136;
  no unseen prior SA08/SA09 proof is certified.
\end{itemize}

The Stacks Project Authors. (2026). \emph{The Stacks Project}.
\url{https://stacks.math.columbia.edu}

\begin{itemize}
\tightlist
\item
  Exact inspected inputs: Theorem14.24.3, Dold--Kan,
  \url{https://stacks.math.columbia.edu/tag/019G}; Lemma14.29.3, a
  normalized chain homotopy gives a simplicial homotopy,
  \url{https://stacks.math.columbia.edu/tag/01A4}; Remark14.26.4, every
  functor preserves simplicial homotopy,
  \url{https://stacks.math.columbia.edu/tag/019M}; Lemma23.5.1, the PD
  polynomial universal property,
  \url{https://stacks.math.columbia.edu/tag/07H5}. The AQ module proves
  the contractions, PD square-zero construction, equalizer and
  rational/localization deductions in full. These known input statements
  were directly inspected on21November2026; the phrase whole-proof check
  refers to the written AQ deductions, not to an independent
  reconstruction of all Stacks inputs.
\end{itemize}

\subsection{Integral Albanese completion and the adelic
assertion}\label{integral-albanese-completion-and-the-adelic-assertion}

Katz, N. M. (1993). Affine cohomological transforms, perversity, and
monodromy. \emph{Journal of the American Mathematical Society, 6}(1),
149--222. \url{https://web.math.princeton.edu/~nmk/old/ACT.pdf}

\begin{itemize}
\tightlist
\item
  Corollary3.4.3, printed p.182, gives smooth irreducible affine
  sections preserving irreducibility of a specified finite étale cover;
  Corollary3.4.2, p.181, gives the connected version. Directly inspected
  for the complete independent reading of CC-ALB01--02. One curve is
  chosen for each fixed n; no simultaneous finite curve for all covers
  is inferred.
\end{itemize}

Suslin, A., \& Voevodsky, V. (1996). Singular homology of abstract
algebraic varieties. \emph{Inventiones Mathematicae, 123}, 61--94.
\url{https://doi.org/10.1007/BF01232367} . Author copy:
\url{https://www.math.ias.edu/vladimir/sites/math.ias.edu.vladimir/files/singular_homology_published.pdf}

\begin{itemize}
\tightlist
\item
  Theorem3.1, printed p.65, and its proof pp.65--66 concern the integral
  finite-correspondence complex of a smooth affine curve and identify H0
  with the relative Picard group of the smooth compactification and its
  reduced boundary. This is the integral curve input in CC-ALB01, not
  the separate symmetric-power comparison in the introduction after
  inversion of the characteristic.
\end{itemize}

Milne, J. S. (2017). \emph{Algebraic groups: The theory of group schemes
of finite type over a field}. Cambridge University Press.
\url{https://www.jmilne.org/math/Books/iAG2017.pdf}

\begin{itemize}
\tightlist
\item
  Directly inspected for CC-ALB01: Theorem5.28 and
  Corollaries5.26--5.27, p.105 (quotients, smoothness and torsors);
  Proposition5.58, p.114 (connected--étale sequence);8.21, p.152 (finite
  faithfully flat multiplication on an abelian variety). The proof
  retains finite nonreduced kernels and universal homeomorphisms in
  characteristic p.
\end{itemize}

The Stacks Project Authors. (2026). \emph{The Stacks Project}.
\url{https://stacks.math.columbia.edu}

\begin{itemize}
\tightlist
\item
  Directly inspected for the prescribed-point curve argument:
  Lemma33.47.1, Tag0FD5,
  \url{https://stacks.math.columbia.edu/tag/0FD5}, and Lemma37.32.3,
  Tag0G4F, \url{https://stacks.math.columbia.edu/tag/0G4F}. The stated
  immersion, codimension-two intersection and base-locus hypotheses are
  all checked in the CC proof.
\end{itemize}

Poonen, B. (2017). \emph{Rational points on varieties} (Graduate Studies
in Mathematics, Vol.186). American Mathematical Society.
\url{https://math.mit.edu/~poonen/papers/Qpoints.pdf}

\begin{itemize}
\tightlist
\item
  Proposition3.5.75 and Remark3.5.76, p.87, give local Zariski density
  and avoidance of a countable union of proper closed subvarieties in
  any neighborhood of a smooth point. Proposition5.2.4, pp.121--122, and
  Corollary5.2.18, p.125, supply the connected-group geometric
  irreducibility and characteristic-zero smoothness. Directly inspected
  with the original18 physical49/footer48 image and the entire
  BR-18-AD01 deduction. No compact-open Galois image theorem is inferred
  from the adopted RH/GRH inputs.
\end{itemize}

The Stacks Project Authors. (2026). \emph{The Stacks Project},
Remark21.39.11 (Simplicial modules), Tag08QD, following Lemma21.39.10.
\url{https://stacks.math.columbia.edu/tag/08QD} . Exact inspected input
for AQ-136-TENSOR: flat simplicial tensor products are computed by the
total tensor of associated complexes; split normalization gives the
normalized version. The binomial and Tor calculations are proved in full
here. The supplied SA08 intake statement A-S09-26 is explicitly credited
as an earlier report of the tensor obstruction; its absent full proof
and stronger rational-Hadamard PD-closure claim are not certified.

\begin{itemize}
\tightlist
\item
  The Stacks Project Authors. (2026). \emph{Coskeleton functors},
  Section14.19, Lemmas14.19.2--14.19.4 (Tag0AMA).
  \url{https://stacks.math.columbia.edu/tag/0AMA}. Directly inspected
  the finite-limit formula, unchanged truncation and their proofs for
  AQ-136-FIRST-SKELETON; the graph and cohomology arguments are given in
  full in the module.
\end{itemize}

\subsection{Integral cohomology-theory source audit and operation
comparison}\label{integral-cohomology-theory-source-audit-and-operation-comparison}

Cartan, H. (1976). Théories cohomologiques. \emph{Inventiones
Mathematicae, 35}, 261--271.
\url{https://gdz.sub.uni-goettingen.de/download/pdf/PPN356556735_0035/LOG_0030.pdf}

\begin{itemize}
\tightlist
\item
  The entire eleven-page article was visually read from the official
  image PDF (physical2--12, after its repository cover). Exact locators:
  §2 p263; Theorem1 p264; §4 and Theorem2 p266; Theorem3 p267; §6 p268;
  §7 pp269--271. The integration theorem uses linear simplicial cochain
  maps, explicitly without requiring multiplicative compatibility. Its
  additional simplex-dimension condition and the flat-cycle hypothesis
  of Theorem3 are retained. The article does not state the original
  strict PD common-source or terminal-localization result.
\end{itemize}

Hatcher, A. (2002). \emph{Algebraic topology}. Cambridge University
Press. Corrected author manuscript:
\url{https://pi.math.cornell.edu/~hatcher/AT/AT.pdf}

\begin{itemize}
\tightlist
\item
  Directly inspected inputs: Theorem3.2 p195 (integral universal
  coefficients), relative connecting/product calculation p218,
  Theorem3B.6 pp275--276 and reduced smash form p276 (integral Kunneth),
  Theorem4.32 p366 (Hurewicz), Theorem4.57 pp393--394 (Eilenberg--Mac
  Lane representability and fundamental class). The new proof uses free
  integral chains of underlying spaces rather than a tensor exactness
  assertion about potentially nonflat cycle modules.
\end{itemize}

Hatcher, A. (2004, January 5). \emph{Spectral sequences in algebraic
topology: Chapter 1, The Serre spectral sequence}. Author manuscript.
\url{https://pi.math.cornell.edu/~hatcher/SSAT/SSch1.pdf}

\begin{itemize}
\tightlist
\item
  Lemma1.10 is stated on p15; its complete proof spans pp15--16 and was
  directly read. For a finitely generated abelian group A it gives
  finite generation of every integral homology group of K(A,n). Lemma1.9
  p15 supplies the fibration step. This finite-generation input,
  together with UCT, is exactly what permits coefficient-localization
  detection on the universal space in AQ-136-PD-RIGIDITY.
\end{itemize}

\subsection{Source-route references}\label{source-route-references}

\begin{itemize}
\tightlist
\item
  Alpöge, L. (2026, November 20). \emph{Remark} {[}Four-page research
  announcement{]}. \url{https://alpo.ge/remark.pdf}. Adopted theorem
  locators are the five named declarations in
  \nolinkurl{ALPOGE_Main.lean}; definitions in
  \nolinkurl{ALPOGE_Statements.lean}.
\item
  Masullo, A. M. (2026). \emph{The music of the primes} {[}Unfinished
  manuscript released November 20{]}.
  \url{https://math216.wordpress.com/wp-content/uploads/2026/11/the-music-of-the-primes.pdf}.
  Exact internal locators listed above.
\item
  Vakil, R. (2026, November 20). \emph{Hurried announcement: Alessandro
  Masullo GRH}. Math 216: Foundations of Algebraic Geometry.
  \url{https://math216.wordpress.com/2026/11/20/hurried-announcement-alessandro-masullo-grh/}.
\item
  Serre, J.-P. (1977). \emph{Linear representations of finite groups}
  (L. L. Scott, Trans.). Springer.
  \url{https://doi.org/10.1007/978-1-4684-9458-7}. Exact input: §10.5,
  Theorem 20, p.78 (Brauer induction). Authoritative publisher record
  retrieved; theorem text was inspected in a digitized book hosted by
  Rochester, search ref \nolinkurl{}.
\item
  Cogdell, J. W. (n.d.). \emph{On Artin L-functions}. Author manuscript,
  \url{https://people.math.osu.edu/cogdell.1/artin-www.pdf}. Exact
  inputs: §2.1, p.10, equation (2); §2.2, p.12, ramified Euler factors
  and inductivity. Retrieved ref \nolinkurl{}; full
  31-page PDF saved as background. Use n.d. because this inspected
  manuscript has no explicit date. The author homepage now links a
  distinct 33-page PDF,
  \url{https://people.math.osu.edu/cogdell.1/Artin-www.pdf}, under the
  published title \emph{L-functions and non-abelian class field theory,
  from Artin to Langlands} (2015), in D. Dumbaugh and J. Schwermer,
  \emph{Emil Artin and Beyond---Class Field Theory and L-Functions}
  (pp.129--161), EMS, DOI 10.4171/146. Author page and publisher
  contents verified (\nolinkurl{},
  \nolinkurl{}); publication date 31 March 2015
  verified on publisher page (\nolinkurl{}). Do not
  transfer the inspected 31-page manuscript page numbers to that
  differently paginated chapter.
\item
  The Stacks Project Authors. (2026). \emph{Plane curves}, §53.9, Lemma
  53.9.3. \url{https://stacks.math.columbia.edu/tag/0BYA}. Retrieved ref
  \nolinkurl{}. This supplies the exact plane-curve genus
  formula; no new proof of the known formula is claimed.
\end{itemize}

The F-SM-02 zero-counting input can be cited to Iwaniec and Kowalski,
\emph{Analytic Number Theory}, Theorem 5.8 (p.104). This locator was
corroborated in search but the actual theorem page was not yet read
during this subrun; it should remain a citation-check item until
inspected or replaced by an already verified equivalent input. F-SM-01
itself assumes only the explicit summability condition and does not
depend on this unresolved bibliographic detail.

\end{document}
